\section*{目 录}

第 1 章 集合与逻辑

考点 1 集合间的基本关系 1

考点 2 集合的运算 2

考点 3 充分条件与必要条件 4

考点 4 命题综合问题 7

第 2 章 不等式

考点 1 等式与不等式的性质 9

考点 2 不等式的求解 10

考点 3 基本不等式及其应用 12

第 3 章 函数

考点 1 函数的定义域与值域 15

考点 2 函数单调性的判断及其应用 16

考点 3 函数的奇偶性及其应用 17

考点 4 函数图像 19

考点 5 函数零点及其应用 20

考点 6 函数模型 22

考点 7 函数综合问题 26

第 4 章 幂函数、指数函数与对数函数

考点 1 幂函数及其应用 29

考点 2 指数函数及其应用 31

考点 3 对数函数及其应用 33

考点 4 幂指对函数值大小比较 36

第 5 章 三角函数

考点 1 任意角的正弦、余弦、正切、余切及同角三角比的基本关系 41

考点 2 诱导公式及其应用 41

考点 3 两角和与差的正弦、余弦、正切公式及二倍角公式 42

考点 4 简单的三角恒等变换 43

考点 5 解三角形 44

考点 6 解三角形之实际应用 51

考点 7 三角函数的图像与性质 59

考点 8 三角函数图像的伸缩平移变换 65

第 6 章 数列

考点 1 数列的增减性 69

考点 2 数列的递推公式 70

考点 3 等差数列及其前 $n$ 项和 72

考点 4 等比数列及其前 $n$ 项和 75

考点 5 数列中的数学文化 77

考点 6 数列求和 78

考点 7 数列与不等式、最值及范围问题 80

考点 8 数列与其他知识点的关联问题 81

第 7 章 平面解析几何

考点 1 直线方程与圆的方程 87

考点 2 直线与圆的位置关系及其应用 88

考点 3 椭圆方程及其性质 90

考点 4 双曲线方程及其性质 106

考点 5 抛物线方程及其性质 116

考点 6 直线与圆锥曲线的位置关系及其应用 123

第 8 章 平面向量

考点 1 向量的概念及其线性运算 128

考点 2 向量的基本定理及其应用 130

考点 3 向量的数量积与向量的坐标表示 130

考点 4 向量的夹角与向量的坐标表示 133

第 9 章 复数

考点 1 复数的概念与运算 135

考点 2 复数的几何意义 137

第 10 章 空间向量与立体几何

考点 1 空间几何体的结构特征、表面积和体积 140

考点 2 空间向量及其应用 148

考点 3 直线、平面平行的判定与性质 149

考点 4 直线、平面垂直的判定与性质 152

考点 5 空间角与距离 156

考点 6 空间几何体的截面、交线及动态问题 179

第 11 章 导数及其应用

考点 1 导数的基本计算及其应用 182

考点 2 求切线方程及其应用 183

考点 3 公切线问题 180

考点 4 利用导数判断函数单调性及其应用 191

考点 5 求极值、最值及其应用 197

考点 6 利用导数研究函数的极值点及其应用 201

考点 7 导数与函数的基本性质结合问题 202

考点 8 利用导数研究函数的零点及其应用 203

第 12 章 排列、组合与二项式定理

考点 1 排列、组合 207

考点 2 二项式定理 209

第 13 章 概率与统计

考点 1 随机抽样 213

考点 2 图表类统计图综合 213

考点 3 样本的数字特征 217

考点 4 变量间的相关关系 224

考点 5 互斥事件的概率计算 228

考点 6 古典概率 231

考点 7 条件概率 237

考点 8 全概率公式 241

考点 9 正态分布指定区间的概率 243

考点 10 独立性检验为载体及其应用 245

考点 11 离散型分布及期望、方差 252

第 14 章 数学建模问题

考点 1 圆锥曲线模型 266

考点 2 函数模型 268

考点 3 不等式模型 270

第 15 章 新颖题型

考点 1 数列新定义 273

考点 2 函数新定义 274

考点 3 集合新定义 279

考点 4 其他新定义 280

\section*{第 1 章 集合与逻辑}

\begin{center}

\begin{tabular}{|p{4cm}|p{12cm}|}
\hline
考点 & 课标要求与命题趋势 \\
\cline{1-2}
1. 集合间的基本关系 & \ding{172}  通过实例，了解集合的含义，理解元素与集合的属于关系  \ding{173}  针对具体问题，能在自然语言和图形语言的基础上，用符号语言刻画集合. 在具体情境中，了解全集与空集的含义  \ding{174}  理解集合之间包含与相等的含义，能识别给定集合的子集 \ding{175}  能够在现实情境或数学情境中，概括出数学对象的一般特征，并用集合语言予以表达  \ding{176}  学会用三种语言(自然语言、图形语言、符号语言)表达数学研究对象，并能进行转换 \\
\cline{1-2}
2. 集合的运算 & \ding{172}  理解两个集合的并集与交集的含义，能求两个集合的并集与交集 \ding{173}  理解在给定集合中一个子集的补集的含义，能求给定子集的补集 \ding{174}  能使用 Venn 图表达集合的基本关系与基本运算，体会图形对理解抽象概念的作用 \ding{175}  掌握集合的基本关系与基本运算 \ding{176}  高考一般给出两个集合，要求通过解不等式求出集合，然后通过集合的运算得出答案 \\
\cline{1-2}
3. 充分条件与必要条件 & \ding{172}  通过对典型数学命题的梳理，理解必要条件的意义，理解性质定理与必要条件的关系 \ding{173}  通过对典型数学命题的梳理，理解充分条件的意义，理解判定定理与充分条件的关系 \ding{174}  通过对典型数学命题的梳理，理解充要条件的意义，理解数学定义与充要条件的关系 \ding{175}  常以关联的知识点作为命题背景，考查充分条件与必要条件，难度随载体而定 \\
\cline{1-2}
4. 命题综合问题 & \ding{172}  通过已知的数学实例，理解全称量词与存在量词的意义 \ding{173}  能正确使用存在量词对全称量词命题进行否定 \ding{174}  能正确使用全称量词对存在量词命题进行否定 \ding{175}  能够借助常用逻辑用语进行数学表达、论证和交流，体会常用逻辑用语在数学中的作用 \ding{176}  全称量词命题和存在量词命题的否定及参数求解，是高考复习与考查的重点 \\
\cline{1-2}
\hline
\end{tabular}
\end{center}
\clearpage
\subsection*{考点1 集合间的基本关系}

\subsubsection*{一、填空题}
\begin{enumerate}
    \item (2023$\cdot$上海春季卷$\cdot$1)已知集合 $A = \{ 1,2\} ,B = \{ 1,a\}$ ，且 $A = B$ ，则 $a =$ \blkda[2]{2}.
    \item (2024$\cdot$静安区二模$\cdot$1)中国国旗上所有颜色组成的集合为\blkda[2]{$\{$红，黄$\}$}.
    \item (2024$\cdot$闵行区一模$\cdot$1)已知集合 $M = \{ 0,1,a + 1\}$ ，若 $- 1 \in  M$ ，则实数 $a =$ \blkda[2]{$-2$}.
    \item (2025$\cdot$杨浦区一模$\cdot$1)已知集合 $A = \{ a,b\}$ ，则 $A$ 的子集个数为\blkda[2]{4}.
    \item (2025$\cdot$金山区二模$\cdot$1)已知集合 $A = \left\{  {x \mid  {x}^{2} - {5x} + 6 = 0}\right\}  ,B = \{ 2,4,6\}$ ，则 $A \cap  B =$ \blkda[2]{$\{ 2\}$}.
    \begin{jieda}
        因为 $A = \left\{  {x \mid  {x}^{2} - {5x} + 6 = 0}\right\}   = \{ 2,3\}$ ,所以 $A \cap  B = \{2\}$.
    \end{jieda}
    \item (2025$\cdot$嘉定区二模$\cdot$1)已知集合 $A = \{ x \mid   - 1 < x \leq  3\}$ ，集合 $B = \{ x \mid  1 \leq  x < 4\}$ ，则 $A \cap  B =$ \blkda[2]{$\{ x \mid  1 \leq  x \leq  3\}$}.
    \item (2024$\cdot$长宁区二模$\cdot$1)已知集合 $A = \{ 1,2\} ,B = \{ 1,a,3\}$ ，且 $A \subseteq  B$ ，则 $a =$ \blkda[2]{2}.
    \item (2024$\cdot$宝山区二模$\cdot$7)已知集合 $A = \left\{  {2,\left| {a + 1}\right| ,a + 3}\right\}$ ，且 $1 \in  A$ ，则实数 $a$ 的值为\blkda[2]{0}.
    \begin{jieda}
        $|a+1| = 1$，解得$a = -2$或$0$，代入检验，$a = -2$时违背互异性.
    \end{jieda}
    \item (2025$\cdot$普陀区二模$\cdot$10)设: ${i\text{ 、 }j}$ ， $n$ 为正整数，集合 $M = \{ 1,3,5,\cdots ,{{2n} - 1}\}$ ，若集合 $A$ 满足 $A \subseteq  M$ ， 且对 $A$ 中任意的两个元素 $i\text{ 、 }j$ ，皆有 $\left| {i - j}\right|  > 4$ 成立，记满足条件的集合 $A$ 的个数为 ${a}_{n}$ ，则 ${a}_{8} =$ \blkda[2]{28}.
    \begin{jieda}
        \begin{dingautolist}{172}
            \item 枚举
            
            当 $n = 8$ 时, $M = \{ 1,3,5,\cdots ,{15}\}$. 
            
            若 $A$ 为空集，符合题意,只有 1 种; 
            
            若 $A$ 为一元集: 共有 8 种; 
            
            若 $A$ 为二元集: 如 $\{ 1,7\} ,\{ 1,9\} ,\{ 1,{11}\} ,\{ 1,{13}\} ,\{ 1,{15}\}, \\
            \{ 3,9\} ,\{ 3,{11}\}, \{ 3,{13}\} ,\{ 3,{15}\} ,\{ 5,{11}\} ,\{ 5,{13}\} ,\{ 5,{15}\} ,\{ 7,{13}\} ,\{ 7,{15}\} ,  \{ 9,{15}\}$ ,共有 15 种; 
            
            若 $A$ 为三元集: 如 $\{ 1,7,{13}\} ,\{ 1,7,{15}\}$ , $\{ 1,9,{15}\} ,\{ 3,9,{15}\}$ ,共有 4 种
            
            所以总共有: ${a}_{8} = 1 + 8 +  {15} + 4 = {28}$ 种

            \item 寻找递推关系
            
            $a_n = a_{n-1}+a_{n-3}$，$a_1 = 2, a_2 = 3, a_3 = 4, a_4 = 6, a_5 = 9, a_6 = 13, a_7 = 19, a_8 = 28$
        \end{dingautolist}
    \end{jieda}
    \item (2024$\cdot$嘉定区二模$\cdot$12)若规定集合 $E = \{ 0,1,2,\cdots ,n\}$ 的子集 $\left\{  {{a}_{1},{a}_{2},{a}_{3},\cdots ,{a}_{m}}\right\}$ 为 $E$ 的第 $k$ 个子集，其中 $k = {2}^{{a}_{1}} + {2}^{{a}_{2}} + {2}^{{a}_{3}} + \cdots  + {2}^{{a}_{m}}$ ，则 $E$ 的第 211 个子集是\blkda[2]{$\{0,1,4,6,7\}$}.
    \begin{jieda}
        $211 = 128 + 64 + 16 + 2 + 1$
    \end{jieda}
\end{enumerate}

\subsubsection*{二、选择题}
\begin{enumerate}
    \item (2024$\cdot$徐汇区一模$\cdot$15)已知集合 $M = \{ \left( {x,y}\right)  \mid  y = f\left( x\right) \}$ ，若对于任意 $\left( {x,y}\right)  \in  M$ ，总存在与之相应的 $\left( {{x}^{\prime },{y}^{\prime }}\right)  \in  M$ (其中 $x \neq  {x}^{\prime }$ ),使得 $\left| {x{x}^{\prime } + y{y}^{\prime }}\right|  = \sqrt{{x}^{2} + {y}^{2}} \cdot  \sqrt{{\left( {x}^{\prime }\right) }^{2} + {\left( {y}^{\prime }\right) }^{2}}$ 成立,则称集合 $M$ 是“$\Omega$ 集合”. 下列选项为“$\Omega$ 集合”的是 \dotfill{(\kh{D})}
    \xx{$M = \left\{  {\left( {x,y}\right)  \mid  y = \dfrac{1}{x},x > 0}\right\}$}{$M = \left\{  {\left( {x,y}\right)  \mid  y = {{e}}^{x} - 2}\right\}$}{$M = \{ \left( {x,y}\right)  \mid  y = \cos x\}$}{$M = \left\{  {\left( {x,y}\right)  \mid  y = {x}^{3}}\right\}$} 
    \begin{jieda}
        记点 $P\left( {x,y}\right)$ ,点 $Q\left( {{x}^{\prime },{y}^{\prime }}\right)$ ,则 $\left| {x{x}^{\prime } + y{y}^{\prime }}\right|  = \sqrt{{x}^{2} + {y}^{2}} \cdot  \sqrt{{\left( {x}^{\prime }\right) }^{2} + {\left( {y}^{\prime }\right) }^{2}} \Longleftrightarrow  \left| {\vv{OP} \cdot  \vv{OQ}}\right|  = \left| \vv{OP}\right|  \cdot  \left| \vv{OQ}\right|$ , 即 $P\text{ 、 }O\text{ 、 }Q$ 三点共线.
        
        由于 $P\text{ 、 }Q$ 不重合,故问题转化为直线 ${OP}$ 与函数 $y = f\left( x\right)$ 的图像至少还有另一个不同的交点. 对于 ${A}$ 项，取 $P\left( {1,1}\right)$ ，排除；对于 ${B}$ 项，取 $P\left( {0, - 1}\right)$ ，排除；对于 $C$ 项，取 $P\left( {0,1}\right)$ ，排除. 而函数 $y = {x}^{3}$ 为奇函数，对于任意非原点的点 $P$ ,均存在与之关于原点对称的点 $Q$ 符合题意,故选 D.
    \end{jieda}
    
    \item (2024$\cdot$宝山区一模$\cdot$16)已知集合 $S$ 是由某些正整数组成的集合，且满足:若 $a \in  S$ ，则当且仅当 $a = m + n$ (其中 $m\text{ 、 }n \in  S,m \neq  n$ )，或 $a = p + q$ (其中正整数 $p\text{ 、 }q \notin  S$ ,且 $p \neq  q$ ).
    
    现有如下两个命题: \ding{172}  $4 \in  S$ ;  \ding{173}  集合 $\{ x \mid  x = {3n} + 5,n \in  \mathbf{N}\}  \subseteq  S$，则下列选项中正确的是 \dotfill{(\kh{C})}
    \xx{\ding{172}  \ding{173} 都是真命题}{\ding{172} 是真命题， \ding{173} 是假命题}{\ding{172} 是假命题， \ding{173} 是真命题}{\ding{172}  \ding{173} 都是假命题}
    \begin{jieda}
        $1 \notin S, 2 \notin S, 3=1+2 \in S, 4 \notin S, 5 = 1+4 \in S$.

        当$n = 0$时，满足${3n} + 5 \in S$，假设$n = k$时${3n} + 5 \in S$，则当$n = k+1$时，${3n} + 5 = 3 + (3k+5)$，因为$3 \in S, 3k+5 \in S$，因此${3n} + 5 \in S$. 根据数学归纳法可知${3n} + 5 \in S$对$n \in \mathbf{N}$恒成立.
    \end{jieda}
\end{enumerate}

\subsection*{考点2 集合的运算}

\subsubsection*{一、填空题}
\begin{enumerate}
    \item (2025$\cdot$上海秋季卷$\cdot$1)已知全集 $U = \{ x \mid  2 \leq  x \leq  5,x \in  \mathbf{R}\}$ ，集合 $A = \{ x \mid  2 \leq  x < 4,x \in  \mathbf{R}\}$ ，则 $\overline{A} =$ \blkda[2]{$\left\lbrack  {4,5}\right\rbrack$}.

    \item (2025$\cdot$杨浦区二模$\cdot$1)已知集合 $A = \{ 1,2,3,4\} ,B = \{ x \mid  1 < x < 4\}$ ，则 $A \cap  B =$ \blkda[2]{$\{ 2,3\}$}.

    \item (2024$\cdot$上海秋季卷$\cdot$1)设全集 $U = \{ 1,2,3,4,5\}$ ，若集合 $A = \{ 2,4\}$ ，则 $\overline{A} =$ \blkda[2]{$\{ 1,3,5\}$}.

    \item (2025$\cdot$徐汇区二模$\cdot$1)已知全集 $U = \{ x\left| \right| x - 1 \mid   \leq  2,x \in  \mathbf{R}\} ,A = \left\lbrack  {1,3}\right\rbrack$ ，则 $\overline{A} =$ \blkda[2]{$\lbrack  - 1,1)$}.

    \item (2025$\cdot$闵行区二模$\cdot$1)设全集 $U = \{  - 1,0,1,2\}$ ，若集合 $A = \{ 0,2\}$ ，则 $\overline{A} =$ \blkda[2]{$\{  - 1,1\}$}.

    \item (2025$\cdot$青浦区二模$\cdot$1)已知集合 $A = \{ x \mid  {2x} \leq  1\} ,B = \{  - 1,0,1,2\}$ ，则 $A \cap  B =$ \blkda[2]{ $\{  - 1,0\}$}.

    \item (2025$\cdot$静安区二模$\cdot$1)已知全集为 $\mathbf{R}$，集合 $A = \{ x \mid   - 2 < x \leq  1\}$ ，则 $\overline{A} =$ \blkda[3]{$(-\infty-2]\cup(1,+\infty)$}.

    \item (2025$\cdot$虹口区二模$\cdot$1)已知全集 $U = \{ 1,2,3,4\} ,A = \{ 1,2\}$ ，则 $\overline{A} =$ \blkda[2]{$\{ 3,4\}$}.

    \item (2025$\cdot$上海春季卷$\cdot$1)已知集合 $A = \{ x \mid  x > 0\} ,B = \{  - 1,0,1,2\}$ ，则 $A \cap  B =$ \blkda[2]{$\{ 1,2\}$}.

    \item (2025$\cdot$闵行区一模$\cdot$1)设集合 $A = \{ 0,1,2,3,4\} ,B = \{ x \mid  0 < x < 3\}$ ，则 $A \cap  B =$ \blkda[2]{$\{ 1,2\}$}.

    \item (2025$\cdot$宝山区二模$\cdot$2)已知集合 $A = \{ 1,2\} ,B = \{ 2,4\}$ ，则 $A \cup  B =$ \blkda[2]{$\{ 1,2,4\}$}.

    \item (2025$\cdot$金山区一模$\cdot$1)已知全集 $U = \{  - 1,0,1,2\}$ ，集合 $A = \{ x \mid   - 1 < x \leq  1,x \in  \mathbf{N}\}$ ，则 $\overline{A} =$ \blkda[2]{$\{  - 1,2\}$}.

    \item (2025$\cdot$宝山区一模$\cdot$1)已知集合 $A = \left( {-1,1}\right) ,B = \mathbf{Z}$ ，则 $A \cap  B =$ \blkda[2]{$\{ 0\}$}.

    \item (2025$\cdot$崇明区一模$\cdot$1)已知集合 $A = \{ 1,2,3,4,5\} ,B = \{ x \mid  x > \sqrt{5}\}$ ，则 $A \cap  B =$ \blkda[2]{$\{ 3,4,5\}$}.

    \item (2025$\cdot$虹口区一模$\cdot$1)已知集合 $A = \{ x \mid  \left| x\right|  < 2\} ,B = \{ 1,2,3\}$ ，则 $A \cap  B =$ \blkda[2]{$\{ 1\}$}.

    \item (2025$\cdot$黄浦区一模$\cdot$1)若集合 $A = \{ 1,2\} ,B = \{ 1,3\}$ ，则 $A \cup  B =$ \blkda[2]{$\{ 1,2,3\}$}.

    \item (2025$\cdot$奉贤区一模$\cdot$1)设全集 $U = \{ 1,2,3,4\}$ ，集合 $A = \{ 2,4\}$ ，则 $\overline{A} =$ \blkda[2]{$\{ 1,3\}$}.

    \item (2025$\cdot$静安区一模$\cdot$1)设集合 $A = \{ 1,3,5,7\} ,B = \{ 2,3,4,5\}$ ，则 $A \cap  B =$ \blkda[2]{$\{ 3,5\}$}.

    \item (2024$\cdot$崇明区二模$\cdot$1)若集合 $A = \{  - 2,0,1\} ,B = \{ x \mid  x <  - 1$ 或 $x > 0\}$ ，则 $A \cap  B =$ \blkda[2]{$\{  - 2,1\}$}.

    \item (2024$\cdot$黄浦区二模$\cdot$1)若集合 $A = \left\lbrack  {1,4}\right\rbrack  ,B = \left\lbrack  {2,5}\right\rbrack$ ，则 $A \cup  B =$ \blkda[2]{$\left\lbrack  {1,5}\right\rbrack$}.

    \item (2024$\cdot$嘉定区二模$\cdot$1)设集合 $A = \left\lbrack  {1,3}\right\rbrack$ ， $B = \left( {2,4}\right)$ ，则 $A \cup  B =$ \blkda[2]{$\lbrack 1,4)$}.

    \item (2024$\cdot$杨浦区一模$\cdot$1)已知全集为 $\mathbf{R}$ ，集合 $A = \left( {2, + \infty }\right)$ ，则 $A$ 的补集可用区间表示为 $\overline{A} =$ \blkda[2]{$( - \infty ,2\rbrack$}.

    \item (2024$\cdot$长宁区一模$\cdot$1)已知集合 $A = ( - \infty ,4\rbrack ,B = \{ 1,3,5,7\}$ ，则 $A \cap  B =$ \blkda[2]{$\{ 1,3\}$}.

    \item (2025$\cdot$普陀区一模$\cdot$1)设全集 $U = \{  - 1,0,1,2,4\}$ ，若集合 $A = \{  - 1,2,4\}$ ，则 $\overline{A} =$ \blkda[2]{$\{ 0,1\}$}.

    \item (2025$\cdot$松江区一模$\cdot$1)已知集合 $A = \lbrack 4, + \infty ),B = \{ 2,4,6,8\}$ ，则 $A \cap  B =$ \blkda[2]{$\{ 4,6,8\}$}.

    \item (2024$\cdot$杨浦区二模$\cdot$1)已知集合 $A = \left( {0,4}\right)$ ， $B = \left( {1,5}\right)$ ，则 $A \cap  B =$ \blkda[2]{$\left( {1,4}\right)$}.

    \item (2024$\cdot$黄浦区一模$\cdot$1)已知集合 $A = \{ x \mid  x \leq  2\}$ ， $B = \{ x \mid  x \geq   - 1\}$ ，则 $A \cap  B =$ \blkda[2]{$[-1,2]$}.

    \item (2024$\cdot$金山区一模$\cdot$1)已知集合 $A = \{ 1,2,3\}$ ， $B = \{ 3,4,5\}$ ，则 $A \cap  B =$ \blkda[2]{$\{ 3\}$}.

    \item (2024$\cdot$清东新区一模$\cdot$1)已知全集 $U = \{ 1,2,3,4\}$ ，集合 $A = \{ 1,3\}$ ，则 $\overline{A} =$ \blkda[2]{$\{ 2,4\}$}.

    \item (2024$\cdot$青浦区一模$\cdot$1)已知集合 $A = \lbrack  - 2,3),B = \{ x \mid   - 1 < x < 6\}$ ，则 $A \cap  B =$ \blkda[2]{$\left( {-1,3}\right)$}.

    \item (2024$\cdot$松江区一模$\cdot$1)已知全集为 $\mathbf{R}$ ，集合 $P = \{ x \mid  x \geq  1\}$ ，则集合 $\overline{P} =$ \blkda[2]{$\{ x \mid  x < 1\}$}.

    \item (2024$\cdot$金山区二模$\cdot$1)已知集合 $M = \{ 1,3,5,7,9\}$ ， $N = \left\{  {x \mid  {2}^{x} = 8}\right\}$ ，则 $M \cap  N =$ \blkda[2]{$\{ 3\}$}.

    \item (2024$\cdot$闵行区二模$\cdot$1)已知集合 $A = \{ x \mid  {2x} + 1 \leq  0\} ,B = \{  - 2, - 1,0\}$ ，则 $A \cap  B =$ \blkda[2]{$\{  - 2, - 1\}$}.

    \item (2024$\cdot$浦东新区二模$\cdot$1)已知集合 $A = \{ 0,1,2\}$ ，集合 $B = \left\{  {x \mid  {2}^{x} > 3}\right\}$ ，则 $A \cap  B =$ \blkda[2]{$\{ 2\}$}.

    \item (2024$\cdot$徐汇区二模$\cdot$1)已知集合 $A = \left\{  {y \mid  y = {x}^{2} + 2}\right\}$ ，集合 $B = \left\{  {x \mid  {x}^{2} - {4x} + 3 \geq  0}\right\}$ ，那么 $A \cap  B =$ \blkda[2]{$\lbrack 3, + \infty )$}.
    \begin{jieda}
        描述法表示的集合注意表示元素的符号.
    \end{jieda}

    \item (2024$\cdot$虹口区一模$\cdot$1)已知集合 $A = \{ 0,1,2,3,4,5\} ,B = \{ x \mid  \left| {x - 2}\right|  \leq  1\}$ ，则 $A \cap  B =$ \blkda[2]{$\{ 1,2,3\}$}.

    \item (2024 $\cdot$徐汇区一模$\cdot$1)已知全集 $U = \mathbf{R}$ ，集合 $M = \{ x \mid  \left| x\right|  > 2\}$ ，则 $\overline{M} =$ \blkda[2]{$\left\lbrack  {-2,2}\right\rbrack$}.

    \item (2025$\cdot$长宁区一模$\cdot$1)设全集为 $\mathbf{R}$，集合 $A = \left\{  {x \mid  {x}^{2} - {2x} - 3 \geq  0}\right\}$ ，则 $\overline{A} =$ \blkda[2]{$\left( {-1,3}\right)$}.

    \item (2024$\cdot$普陀区二模$\cdot$2)已知 $a \in  \mathbf{R}$ ，设集合 $A = \{ 1,a,4\}$ ，集合 $B = \{ 1,a + 2\}$ ，若 $A \cap  B = B$ ，则 $a =$ \blkda[2]{2}.
    \begin{jieda}
        $A \cap  B = B \Longleftrightarrow B \subseteq A$
        
        当 $a + 2 = a$ 时,等式不成立,舍去; 当 $a +  2 = 4$ 时,解得 $a = 2$ ,此时 $A = \{ 1,2,4\} ,B = \{ 1,4\}$ ,满足题意，故 $a = 2$.
    \end{jieda}

    \item (2024$\cdot$奉贤区一模$\cdot$2)设集合 $A = \{  - 2, - 1,0,5,{10},{20}\} ,B = \{ x \mid  \lg x < 1\}$ ，则 $A \cap  B =$ \blkda[2]{$\{5\}$}.

    \item (2025$\cdot$青浦区一模$\cdot$2)已知集合 $A = \{ x \mid  x = {{2k} - 1},k \in  \mathbf{N}\} ,B = \{  - 1,0,1,2,3\}$ ，则 $A \cap  B =$ \blkda[2]{$\{  - 1,1,3\}$}.

    \item (2025$\cdot$黄浦区二模$\cdot$2)设 $a \in  \mathbf{R}$ ，集合 $A = \left\lbrack  {1,3}\right\rbrack$ ， $B = \left\lbrack  {a,4}\right\rbrack$ ，若 $A \cap  B = \left\lbrack  {2,3}\right\rbrack$ ，则 $a =$ \blkda[2]{2}.

    \item (2025$\cdot$崇明区二模$\cdot$3)已知全集 $U = \mathbf{R}$ ，集合 $A = \{ 1,2,4\}$ ， $B = \{ 2,4,5\}$ ，则 $A \cap  \overline{B} =$ \blkda[2]{$\{1\}$}.
    \begin{jieda}
        寻找在$A$中，不在$B$中的元素.
    \end{jieda}

    \item (2024$\cdot$虹口区二模$\cdot$4)已知集合 $A = \{ x \mid  \tan x < 0\}$ ， $B = \left\{  {x\left| {\dfrac{x - 2}{x} \leq  0}\right. }\right\}$ ，则 $A \cap  B =$ \blkda[2]{$\left( {\dfrac{\pi }{2},2}\right\rbrack$}.
    \begin{jieda}
        $B = (0, 2]$，结合正切图像可解.
    \end{jieda}

    \item (2024$\cdot$普陀区一模$\cdot$7)设集合 $M = \{ 2,0, - 1\} ,N = \{ x \mid  \left| {x - a}\right|  < 1\}$ ，若 $M \cap  N$ 的真子集的个数是 1，则正实数 $a$ 的取值范围为\blkda[2]{$\left( {0,1}\right)  \cup  \left( {1,3}\right)$}.
    \begin{jieda}
        由题意可知：$M,N$有一个公共元素，$N = (a-1, a+1)$，画数轴分类讨论即可.
    \end{jieda}
\end{enumerate}
\subsubsection*{二、选择题}
\begin{enumerate}
    \item (2023$\cdot$上海秋季卷$\cdot$13)已知集合 $P = \{ 1,2\}$ ， $Q = \{ 2,3\}$ ，若 $M = \{ x \mid  x \in  P$ 且 $x \notin  Q\}$ ，则 $M =$ \dotfill{(\kh{A})}
    \xx{$\{ 1\}$}{$\{2\}$}{$\{ 3\}$}{$\{ 1,2,3\}$}
    \item (2025$\cdot$浦东新区二模$\cdot$13)已知集合 $P = \{  - 1,1,3,5,7\}$ ，集合 $Q = \{ x \mid  \left| {x - 3}\right|  > 3\}$ ，全集为 $\mathbf{R}$ ，则 $P \cap  \overline{Q} =$ \dotfill{(\kh{D})}
    \xx{$\{  - 1,1\}$}{$\{  - 1,7\}$}{$\{ 1,3,5,7\}$}{$\{ 1,3,5\}$}

    \item (2024$\cdot$崇明区一模$\cdot$13)已知集合 $A = \{ x \mid   - 2 \leq  x \leq  3\} ,B = \{ x \mid  x > 0\}$ ，则 $A \cup  B =$ \dotfill{(\kh{D})}
    \xx{$\left\lbrack  {-2,3}\right\rbrack$}{$\left\lbrack  {0,3}\right\rbrack$}{$\left( {0, + \infty }\right)$}{$\lbrack  - 2, + \infty )$}

    \item (2024$\cdot$松江区二模$\cdot$13)已知集合 $A = \{ x \mid  0 \leq  x \leq  4\} ,B = \{ x \mid  x = {2n},n \in  \mathbf{Z}\}$ ，则 $A \cap  B =$ \dotfill{(\kh{D})}
    \xx{$\{ 1,2\}$}{$\{ 2,4\}$}{$\{ 0,1,2\}$}{$\{ 0,2,4\}$}
\end{enumerate}
\subsection*{考点3  充分条件与必要条件}

\subsubsection*{一、填空题}
\begin{enumerate}
    \item (2025$\cdot$徐汇区一模$\cdot$6)已知 $m$ 、 $n$ 为空间中两条不同的直线， $\alpha$ 、 $\beta$ 为两个不同的平面，若 $m \subset  \alpha$ ， $\alpha  \cap  \beta  = n$ ，则 $m//n$ 是 $m//\beta$ 的\blkda[2]{充要}条件. (填:“充分非必要”，“必要非充分”，“充要”，“既非充分又非必要条件”中的一个)
    \item (2025$\cdot$长宁区一模$\cdot$9)已知 $\alpha  : {2}^{x} + {\log }_{2}x \leq  2,\beta  : x < m$ ，若 $\alpha$ 是 $\beta$ 的充分条件，则实数 $m$ 的取值范围是\blkda[2]{$(1, +\infty)$}.
    \begin{jieda}
        $y = {2}^{x} + {\log }_{2}x$严增，$\alpha : x \in (-\infty, 1]$，$(-\infty, 1] \subseteq (-\infty, m)$，故$m > 1$.
    \end{jieda}
\end{enumerate}

\subsubsection*{二、选择题}
\begin{enumerate}
    \item (2025$\cdot$闵行区一模$\cdot$13)在空间中，“ $a$ 、 $b$ 为异面直线”是“$a$ 、 $b$ 不相交” 的 \dotfill{(\kh{A})}
    \xx{充分非必要条件}{必要非充分条件}{充要条件}{既非充分又非必要条件}
    \item (2025$\cdot$静安区一模$\cdot$13)设 $a$ 、 $b \in  \mathbf{R}$ ，则“$a + b > 0$”是“$a > 0$ 且 $b > 0$”的 \dotfill{(\kh{B})}
    \xx{充分非必要条件}{必要非充分条件}{充要条件}{既非充分又非必要条件}

    \item (2025$\cdot$宝山区一模$\cdot$13)设 ${ab} > 0$ ，则“$a > b$”是“$\dfrac{1}{a} < \dfrac{1}{b}$”的 \dotfill{(\kh{C})}
    \xx{充分非必要条件}{必要非充分条件}{充要条件}{既非充分又非必要条件}

    \item (2025$\cdot$嘉定区一模$\cdot$13)已知 $a$ 为正数，则“$a > 3$”是“${a}^{a} > {a}^{3}$”的 \dotfill{(\kh{A})}
    \xx{充分非必要条件}{必要非充分条件}{充要条件}{既非充分又非必要条件}
    \begin{jieda}
        根据幂函数的单调性可知充分性显然
        
        当$a \in (0, 1)$时，指数函数$y = a^x$严减，故$a^a > a^3$依然成立，故必要性不成立.
    \end{jieda}

    \item (2025$\cdot$闵行区二模$\cdot$14)已知 $a > 0,b > 0$ ，则“ $a + b > 2$ ”是“ ${ab} > 1$ ”的 \dotfill{(\kh{B})}
    \xx{充分非必要条件}{必要非充分条件}{充要条件}{既非充分又非必要条件}
    \begin{jieda}
        在平面直角坐标系上找出对应区域即可知.
    \end{jieda}

    \item (2024$\cdot$闵行区二模$\cdot$13)设 $a \in  \mathbf{R}$ ，则“${a}^{2} > 1$”是“${a}^{3} > 1$”的 \dotfill{(\kh{B})}
    \xx{充分非必要条件}{必要非充分条件}{充要条件}{既非充分又非必要条件}

    \item (2024$\cdot$长宁区二模$\cdot$13)设 $z \in  \mathbf{C}$ ，则“$z = \overline{z}$”是“$z \in  \mathbf{R}$”的 \dotfill{(\kh{C})}
    \xx{充分非必要条件}{必要非充分条件}{充要条件}{既非充分又非必要条件}

    \item (2024$\cdot$宝山区一模$\cdot$13)“ $x > 1$ ”是“ $\left| x\right|  > 1$ ”的 \dotfill{(\kh{A})}
    \xx{充分非必要条件}{必要非充分条件}{充要条件}{既非充分又非必要条件}

    \item (2024$\cdot$黄浦区一模$\cdot$13)设 $x \in  \mathbf{R}$ ，则“ ${x}^{3} > 8$ ”是“ $\left| x\right|  > 2$ ”的 \dotfill{(\kh{A})}
    \xx{充分非必要条件}{必要非充分条件}{充要条件}{既非充分又非必要条件}

    \item (2024$\cdot$金山区一模$\cdot$13)已知 $a$ 、 $b$ 、 $c$ 是实数，则“ $a > b$ ”是“ $a{c}^{2} > b{c}^{2}$ ”的 \dotfill{(\kh{B})}
    \xx{充分非必要条件}{必要非充分条件}{充要条件}{既非充分又非必要条件}
    \begin{jieda}
        充分性反例：$c = 0$.
    \end{jieda}

    \item (2024$\cdot$静安区一模$\cdot$13)已知 $\alpha  : x > 1,\beta  : \dfrac{1}{x} < 1$ ，则 $\alpha$ 是 $\beta$ 的 \dotfill{(\kh{A})}
    \xx{充分非必要条件}{必要非充分条件}{充要条件}{既非充分又非必要条件}

    \item (2024$\cdot$浦东新区一模$\cdot$13)如果 $a > 0 > b$ ，则下列不等式中一定成立的是 \dotfill{(\kh{D})}
    \xx{$\sqrt{a} > \sqrt{-b}$}{${a}^{2} > {b}^{2}$}{${a}^{2} < {ab}$}{${a}^{3} > {b}^{3}$}

    \item (2024$\cdot$青浦区一模$\cdot$13)已知 $a$ 、 $b \in  \mathbf{R}$ ，则“$a > b$”是“${a}^{3} > {b}^{3}$”的 \dotfill{(\kh{C})}
    \xx{充分非必要条件}{必要非充分条件}{充要条件}{既非充分又非必要条件}

    \item (2024$\cdot$徐汇区一模$\cdot$13)设 ${z}_{1}$ 、 ${z}_{2} \in  \mathbf{C}$ ，则“${z}_{1}$ 、 ${z}_{2}$ 中至少有一个虚数”是“${z}_{1} - {z}_{2}$ 为虚数”的 \dotfill{(\kh{B})}
    \xx{充分非必要条件}{必要非充分条件}{充要条件}{既非充分又非必要条件}
    \begin{jieda}
        当 ${z}_{1} = 1 + {i},{z}_{2} = {i}$ 时, ${z}_{1} - {z}_{2} = 1$ 不为虚数，充分性不成立.

        当 ${z}_{1} - {z}_{2}$ 为虚数时, ${z}_{1}\text{ 、 }{z}_{2}$ 的虚部不全为 0，故 ${z}_{1}$ 、 ${z}_{2}$ 中至少有一个虚数，必要性成立. 
    \end{jieda}

    \item (2025$\cdot$虹口区二模$\cdot$13)若 $a$ 是实数，则“${a}^{2} > 1$”是“$a > 2$”的\dotfill{(\kh{B})}
    \xx{充分非必要条件}{必要非充分条件}{充要条件}{既非充分又非必要条件}

    \item (2024$\cdot$浦东新区二模$\cdot$13)“$a = 1$” 是 “直线 $ax - {2y} - 2 = 0$ 与直线 $x - \left( {a + 1}\right) y + 1 = 0$ 平行”的 \dotfill{(\kh{C})}
    \xx{充分非必要条件}{必要非充分条件}{充要条件}{既非充分又非必要条件}
    \begin{jieda}
        $\vv{n_1} = (a, -2), \vv{n_2} = (1, -(a+1))$，$\vv{n_1} // \vv{n_2} \Longleftrightarrow a^2+a-2 = 0$，解得$a = 1$或$-2$，代回检验可知$a = 1$时两直线平行，$a = -2$时两直线重合.
    \end{jieda}

    \item (2025$\cdot$静安区二模$\cdot$13)“$m<2$”是“一元二次不等式 ${x}^{2} + {mx} + 1 > 0$ 的解集为 $\mathbf{R}$ ”的 \dotfill{(\kh{B})}
    \xx{充分非必要条件}{必要非充分条件}{充要条件}{既非充分又非必要条件}

    \item (2025$\cdot$杨浦区一模$\cdot$13)已知实数 $a \neq  0$ ，则“ $a > 2$ ”是“ $\dfrac{1}{a} < \dfrac{1}{2}$ ”的\dotfill{(\kh{A})}
    \xx{充分非必要条件}{必要非充分条件}{充要条件}{既非充分又非必要条件}

    \item (2025$\cdot$奉贤区一模$\cdot$13)在 $\bigtriangleup  {ABC}$ 中，“ $C = \dfrac{\pi }{2}$ ”是“ ${\sin }^{2}A + {\sin }^{2}B = 1$ ”的 \dotfill{(\kh{A})}
    \xx{充分非必要条件}{必要非充分条件}{充要条件}{既非充分又非必要条件}
    \begin{jieda}
        ${\sin }^{2}A + {\sin }^{2}B = 1$即$sinB = |cosA|$
        
        当 $C = \dfrac{\pi }{2}$ 时,则 ${\sin }^{2}A + {\sin }^{2}B = {\sin }^{2}A + {\sin }^{2}\left( {\dfrac{\pi }{2} - A}\right)  = {\sin }^{2}A + {\cos }^{2}A = 1$ ,充分性成立; 
        
        当 $A = \dfrac{2\pi }{3},B = \dfrac{\pi }{6}$ 时 ${\sin }^{2}A + {\sin }^{2}B = 1$ ,但是 $C = \dfrac{\pi }{6}$ ,则必要性不成立
    \end{jieda}

    \item (2025$\cdot$普陀区一模$\cdot$13)已知 $\alpha$ 为任意角，则“$\cos {2\alpha } = \dfrac{1}{3}$”是“$\sin \alpha  = \dfrac{\sqrt{3}}{3}$ ”的 \dotfill{(\kh{B})}
    \xx{充分非必要条件}{必要非充分条件}{充要条件}{既非充分又非必要条件}

    \item (2025$\cdot$浦东新区二模$\cdot$14)“$a>b$”是“$\lg a > \lg b$ ”的 \dotfill{(\kh{B})}
    \xx{充分非必要条件}{必要非充分条件}{充要条件}{既非充分又非必要条件}
    \begin{jieda}
        $a, b$未必在定义域内.
    \end{jieda}

    \item (2025$\cdot$金山区一模$\cdot$14)若“ $x > a$ ”是“ ${x}^{2} - {2x} - 3 < 0$ ”的必要不充分条件，则实数 $a$ 的取值范围是 \dotfill{(\kh{B})}
    \xx{$\left( {-\infty , - 1}\right)$}{$\left( {-\infty , - 1}\right)$}{$\left( {-1, + \infty }\right)$}{$\lbrack  - 1, + \infty )$}

    \item (2025$\cdot$宝山区二模$\cdot$14)“ $a > b$ ”的一个必要非充分条件是 \dotfill{(\kh{C})}
    \xx{$\ln \left( {a - b}\right)  > 0$}{${2}^{a} > {2}^{b}$}{${\left( \dfrac{1}{2}\right) }^{a - b} \leq  1$}{${a}^{3} > {b}^{3}$}

    \item (2024$\cdot$长宁区一模$\cdot$14)“$P(A \cap B) = P(A) P(B)$”是“事件$A$ 与事件 $\overline{B}$ 互相独立”的\dotfill{(\kh{C})}
    \xx{充分非必要条件}{必要非充分条件}{充要条件}{既非充分又非必要条件}

    \item (2025$\cdot$青浦区一模$\cdot$16)对于数列 $\left\{  {a}_{n}\right\}$ ，设数列 $\left\{  {a}_{n}\right\}$ 的前 $n$ 项和为 ${S}_{n}$ ，给出下列两个命题: \ding{172}  存在函数 $y = f\left( x\right)$ ，使得 ${S}_{n} = f\left( {a}_{n}\right)$ ； \ding{173}  存在函数 $y = g\left( x\right)$ ，使得 $n = g\left( {a}_{n}\right)$. 则 \ding{172} 是 \ding{173} 的 \dotfill{(\kh{B})}
    \xx{充分非必要条件}{必要非充分条件}{充要条件}{既非充分又非必要条件}
    \begin{jieda}
        取 ${S}_{n} = {a}_{n} = 0$ ,存在 $y = f\left( x\right)  = 0$ ,使得 ${S}_{n} = f\left( {a}_{n}\right)$ 成立,此时由函数定义知,不存在函数 $y = g\left( x\right)$ ,使得 $n =  g\left( {a}_{n}\right)$
        
        当存在函数 $y = g\left( x\right)$ ,使得 $n = g\left( {a}_{n}\right)$ 成立时,由于 $n$ 与 ${a}_{n}$ 为一一对应关系，所以$\{a_n\}$中每一项各不相同，所以存在函数 $y = f\left( x\right)$ ，使得 ${S}_{n} = f\left( {a}_{n}\right)$.
    \end{jieda}

    \item (2024$\cdot$上海秋季卷$\cdot$15)定义一个集合 $\Omega$ ，集合中的元素是空间内的点集，任取 ${P}_{1}$ 、 ${P}_{2}$ 、 ${P}_{3} \in  \Omega$ ， 存在不全为 0 的实数 ${\lambda }_{1}$ 、 ${\lambda }_{2}$ 、 ${\lambda }_{3}$ ，使得 ${\lambda }_{1}\vv{O{P}_{1}} + {\lambda }_{2}\vv{O{P}_{2}} + {\lambda }_{3}\vv{O{P}_{3}} = \vv{0}$. 已知 $\left( {1,0,0}\right)  \in  \Omega$ ，则 $\left( {0,0,1}\right)  \notin  \Omega$ 的充分条件是 \dotfill{(\kh{C})}
    \xx{$\left( {0,0,0}\right)  \in  \Omega$}{$\left( {-1,0,0}\right)  \in  \Omega$}{$\left( {0,1,0}\right)  \in  \Omega$}{$\left( {0,0, - 1}\right)  \in  \Omega$}
    \begin{jieda}
        存在不全为 0 的实数 ${\lambda }_{1}$ 、 ${\lambda }_{2}$ 、 ${\lambda }_{3}$ ，使得 ${\lambda }_{1}\vv{O{P}_{1}} + {\lambda }_{2}\vv{O{P}_{2}} + {\lambda }_{3}\vv{O{P}_{3}} = \vv{0}$即$\vv{O{P}_{1}}, \vv{O{P}_{2}}, \vv{O{P}_{3}}$不构成一组基. 故应有$\left( {1,0,0}\right)，\left( {0,0,1}\right)$与选项构成一组基.
    \end{jieda}
\end{enumerate}

\subsubsection*{三、解答题}
\begin{enumerate}
    \item (2024$\cdot$普陀区一模$\cdot$19)设函数 $y = f\left( x\right)$ 的表达式为 $f\left( x\right)  = a{{e}}^{x} + {{e}}^{-x}$.
    \begin{enumerate}
        \item 求证:“$a = 1$”是 “函数 $y = f\left( x\right)$ 为偶函数” 的充要条件
        \item 若 $a = 1$ ，且 $f\left( {m + 2}\right)  \leq  f\left( {{2m} - 3}\right)$ ，求实数 $m$ 的取值范围.
    \end{enumerate}
    \begin{jieda}
        \begin{enumerate}
            \item 充分性:当 $a = 1$ 时， $f\left( x\right)  = {{e}}^{x} + {{e}}^{-x}$ ，定义域为 $\left( {-\infty , + \infty }\right)$ 关于原点对称,则 $f\left( {-x}\right)  = {{e}}^{-x} + {{e}}^{x} = f\left( x\right)$ , 即函数 $y = f\left( x\right)$ 为偶函数. 
            
            必要性: 当函数 $y = f\left( x\right)$ 为偶函数时,即 $f\left( {-x}\right)  = f\left( x\right)$ ,则 $a{{e}}^{-x} + {{e}}^{x} = a{{e}}^{x} + {{e}}^{-x}$ ,即 $\left( {a - 1}\right) \left( {{{e}}^{-x} - {{e}}^{x}}\right)  = 0$ 对 $x \in  \mathbf{R}$ 恒成立,则 $a - 1 = 0$ ,即 $a = 1$. 
            
            综上所述，“ $a = 1$ ”是“函数 $y = f\left( x\right)$ 为偶函数”的充要条件.
            \item 当 $a = 1$ 时， $f\left( x\right)  = {{e}}^{x} + {{e}}^{-x}$为偶函数，在$(-\infty, 0)$严减，在$(0, +\infty)$严增，故$f\left( {m + 2}\right)  \leq  f\left( {{2m} - 3}\right) \Longleftrightarrow |m+2| \leq |2m-3| \Longleftrightarrow \left( {{3m} - 1}\right) \left( {m - 5}\right)  \geq  0$，故实数 $m$ 的取值范围是 $\left( {-\infty ,\dfrac{1}{3}}\right\rbrack   \cup  \lbrack 5, + \infty )$.
        \end{enumerate}
    \end{jieda}
    \item (2024$\cdot$上海春季卷$\cdot$21)对于定义在 $\mathbf{R}$ 上的函数 $y = f\left( x\right)$ ，记集合 ${M}_{a} = \left\{  {t \mid  t = f\left( x\right)  - f\left( a\right) ,x \geq  }\right.  a\} ,{L}_{a} = \left\{  {t \mid  t = f\left( x\right)  - f\left( a\right) ,x \leq  a}\right\}$.
    \begin{enumerate}
        \item 若 $f\left( x\right)  = {x}^{2} + 1$ ，求 $M$ (1) 和 $L$ (1)
        \item 若 $f\left( x\right)  = {x}^{3} - 3{x}^{2}$ ，求证:对任意 $a \in  \mathbf{R}$ ，都有 $M\left( a\right)  \subseteq  \lbrack  - 4, + \infty )$ ，且存在 $a$ ，使得 $- 4 \in  M\left( a\right)$
        \item 已知定义在 $\mathbf{R}$ 上的函数 $y = f\left( x\right)$ 有最小值，证明:“$y = f\left( x\right)$ 是偶函数”的充要条件为 “对任意 $c \in  \mathbf{R}$ ,都有 ${M}_{\left( -c\right) } = {L}_{c}$”
    \end{enumerate}
    \begin{jieda}
        \begin{enumerate}
            \item  $M\left( 1\right)  = \left\{  {t \mid  t = {x}^{2} + 1 - 2,x \geq  1}\right\}   = \lbrack 0, + \infty )$
            
            $L\left( 1\right)  = \left\{  {t\left| {t = {x}^{2} + 1 - 2 = {x}^{2} - 1}\right. ,x \leq  1}\right\}   = \lbrack  - 1, + \infty )$
            \item 证明:由题意知， ${M}_{a} = \left\{  {t \mid  t = f\left( x\right)  - f\left( a\right) ,x \geq  a}\right\}   = \{ t \mid  \left. {t = {x}^{3} - 3{x}^{2} - \left( {{a}^{3} - 3{a}^{2}}\right) ,x \geq  a}\right\}$ ,记 $g\left( x\right)  = {x}^{3} - 3{x}^{2} - {a}^{3} +  3{a}^{2}$，$g(a) = 0$
            
            ${g}^{\prime }\left( x\right)  = 3{x}^{2} - {6x} = 0 \Longrightarrow  x = 0$ 或 $x = 2$ 
            
            $g(x)$在$(-\infty, 0), (2, +\infty)$严增，在$(0, 2)$严减，故
            
            \begin{dingautolist}{172}
                \item  当 $a \geq  2,t = {x}^{3} - 3{x}^{2} - {a}^{3} + 3{a}^{2},x \geq  a$为严格增函数,因为 $g\left( a\right)  = 0$ ,所以此时 $M\left( a\right)  = \lbrack 0, + \infty ) \subseteq  \lbrack  - 4, + \infty )$ ,符合条件;

                \item  当 $0 \leq  a < 2$ 时， $t = {x}^{3} - 3{x}^{2} - {a}^{3} + 3{a}^{2},x \geq  a$ 先增后减， ${t}_{\min } = g\left( 2\right)  =  - {a}^{3} - 3{a}^{2} - 4$ ,因为 $- {a}^{3} - 3{a}^{2} - 4 = {a}^{2}\left( {3 - a}\right)  \geq  0\left( {a\text{ 取等号 }}\right)$ ,所以 ${t}_{\min } = g\left( 2\right)  =  - {a}^{3} - 3{a}^{2} - 4 \geq   - 4$ ,此时 $M\left( a\right)  = \left\lbrack  {-{a}^{3} - 3{a}^{2} - 4, + \infty }\right)  \subseteq  \left\lbrack  {-4, + \infty }\right)$ ,符合条件;

                \item 当 $a < 0$ 时， $t = {x}^{3} - 3{x}^{2} - {a}^{3} + 3{a}^{2},x \geq  a$ ，在 $\lbrack a,0)$ 严格增，在 $\left\lbrack  {0,2}\right\rbrack$ 严格减，在 $\lbrack 2, + \infty )$ 严格增， ${t}_{\min } = \min \{ g\left( a\right), g\left( 2\right) \}  = \min \left\{  {0, - {a}^{3} + 3{a}^{2} - 4}\right\}$ ,因为 $h\left( a\right)  =  - {a}^{3} + 3{a}^{2} - 4$ , 当 $a < 0$ 时， ${h}^{\prime }\left( a\right)  =  - 3{a}^{2} + {6a} < 0$ ，则此时 $M\left( a\right)  = \left\lbrack  {t}_{\min }\right.$ ， $+ \infty ) \subseteq  \lbrack  - 4, + \infty )$ 成立
            \end{dingautolist}
            综上所述，对任意 $a \in  \mathbf{R}$ ，都有 $M\left( a\right)  \subseteq  \lbrack  - 4, + \infty )$ ,且存在 $a=0$ ,使得 $- 4 \in  M\left( a\right)$
            \item 必要性: ${M}_{\left( -c\right) } = \left\{  {t \mid  t = f\left( x\right)  - f\left( {-c}\right) ,x \geq   - c}\right\} = \left\{  {t \mid  t = f\left( -x\right)  - f\left( {c}\right) , -x \geq - c}\right\} \\
            = \left\{ t \mid  t = f\left(-x\right)  - f\left( {c}\right) , x \leq c\right\}$
            
            因为 $y = f\left( x\right)$ 为偶函数，$f(x) = f(-x)$，所以${M}_{\left( -c\right) } = \left\{  t \mid  t = f\left(x\right)  - f\left( {c}\right) , x \leq c\right\} = {L}_{c}$

            充分性：假设$f(x)$不是偶函数，即存在$x_0 > 0$使得$f(x_0) \neq f(-x_0)$

            设$f(x)$在$x = m$处取得最小值，若$m = x_0$或$-x_0$，取$c = 0$，${M}_{\left( -c\right) }$最小值为$f(m)-f(0)$，$L_c$最小值必然大于$f(m) - f(0)$，违背${M}_{\left( -c\right) } = {L}_{c}$. 故必有$f(m) = f(-m)$，故不妨设$m > 0$

            若$m > x_0$，则取$-c = x_0$，${M}_{\left( -c\right) }$最小值为$f(m)-f(x_0)$，$L_c$最小值为$f(m) - f(-x_0)$，违背${M}_{\left( -c\right) } = {L}_{c}$.

            若$0 < m < x_0$，则取$c = x_0$，${M}_{\left( -c\right) }$最小值为$f(m)-f(-x_0)$，$L_c$最小值为$f(m) - f(x_0)$，违背${M}_{\left( -c\right) } = {L}_{c}$.

            综上所述，“ $y = f\left( x\right)$ 是偶函数”的充要条件为“对任意 $c \in$ R,都有 ${M}_{\left( -c\right) } = {L}_{c}$ ”.
        \end{enumerate}
    \end{jieda}
\end{enumerate}

\subsection*{考点4 命题综合问题}

\subsubsection*{一、选择题}
\begin{enumerate}
    \item (2024$\cdot$崇明区一模$\cdot$15)已知点 $M$ 为正方体 ${ABCD} - {A}_{1}{B}_{1}{C}_{1}{D}_{1}$ 内部(不包含表面)的一点. 给出下列两个命题:

    ${q}_{1}$ : 过点 $M$ 有且只有一个平面与 $A{A}_{1}$ 和 ${B}_{1}{C}_{1}$ 都平行;

    ${q}_{2}$ : 过点 $M$ 至少可以作两条直线与 $A{A}_{1}$ 和 ${B}_{1}{C}_{1}$ 所在的直线都相交.

    则以下说法正确的是 \dotfill{(\kh{A})}

    \xx{命题 ${q}_{1}$ 是真命题,命题 ${q}_{2}$ 是假命题}{命题 ${q}_{1}$ 是假命题,命题 ${q}_{2}$ 是真命题}{命题 ${q}_{1}\text{ 、 }{q}_{2}$ 都是真命题}{命题 ${q}_{1}\text{ 、 }{q}_{2}$ 都是假命题}
    \begin{jieda}
        已知点 $M$ 为正方体 ${ABCD} - {A}_{1}{B}_{1}{C}_{1}D$ 内 (不包含表面) 的一点,过点 $M$ 的平面为 $\alpha$ ,如图所示. 对于 ${q}_{1}$ ,在平面 $A{A}_{1}{D}_{1}D$ 与平面 $B{B}_{1}{C}_{1}C$ 之间与平面 $A{A}_{1}{D}_{1}D$ 与平面 $B{B}_{1}{C}_{1}C$ 平行的平面均与 $A{A}_{1}$ 和 ${B}_{1}{C}_{1}$ 平行,如平面 $\alpha$ ,当点 $M$ 为正方体 ${ABCD} - {A}_{1}{B}_{1}{C}_{1}D$ 内(不包含表面)的一点， 满足要求的平面有且只有一个,故命题 ${q}_{1}$ 是真命题;        
        \begin{center}
        \includegraphics[width=0.2\textwidth]{bodyxb/src/8_1383_1050_190_187_0.jpg}
        \end{center}
        
        对于 ${q}_{2}$，直线外一点和直线确定一个平面，$M$和$B_1C_1$确定平面$\gamma_1$，$M$和$AA_1$确定平面$\gamma_2$，$\gamma_1 \cap \gamma_2 = l$，$M \in l$，过$M$至多能做出$l$一条直线和$A{A}_{1}$ 和 ${B}_{1}{C}_{1}$ 所在的直线都相交，所以命题 ${q}_{2}$ 是假命题. 

    \end{jieda}
    \item (2023$\cdot$上海秋季卷$\cdot$16)对于一段曲线C，若存在点 $M$ ，使得对于任意的点 $P \in  C$ ，都存在点 $Q \in  C$ ,使得 $\left| {PM}\right|  \cdot  \left| {QM}\right|  = 1$ ,则称曲线 $C$ 为“自相关曲线”. 现有如下两个命题：\ding{172}  任何椭圆都是“自相关曲线”;  \ding{173}  存在双曲线是“自相关曲线”. 则下列正确的是 \dotfill{(\kh{B})}
    \xx{\ding{172} 真命题， \ding{173} 真命题}{\ding{172} 真命题， \ding{173} 假命题}{\ding{172} 假命题， \ding{173} 真命题}{\ding{172} 假命题， \ding{173} 假命题}
    \begin{jieda}
        曲线 $C$ 为“自相关曲线”即存在点 $M$ ，使得$C$上点到$M$的距离取值范围为$\left[\dfrac{1}{t}, t\right]$，显然对于椭圆肯定存在$M$，而对于双曲线肯定不存在$M$.
    \end{jieda}
    
    \item (2024$\cdot$奉贤区二模$\cdot$16)如图，在等腰梯形 ${ABCD}$ 中， ${AD}//{BC},{AD} = 1,{BC} = m(m>1), \angle {ABC} = \dfrac{\pi }{3}$. 点 $E$ 是线段 ${AB}$ 上的一点,点 $F$ 在线段 ${DC}$ 上, $\dfrac{DF}{DC} = t$.
    \begin{center}
    \includegraphics[width=0.2\textwidth]{bodyxb/src/11_1169_698_222_236_0.jpg}
    \end{center}

    命题 \ding{172} :若 $\vv{AE} = \dfrac{1}{2}\vv{EB}$ ，则 $\vv{EF} \cdot \vv{AD}$ 随着 $t$ 的增大而减小

    命题 \ding{173} :设 $\dfrac{AE}{AB} = x$ ，若存在线段 ${EF}$ 把梯形 ${ABCD}$ 的面积分成上下相等的两个部分,那么 $x \geq  \dfrac{m - 1}{2m},t = f\left( x\right)$ 随着 $x$ 的增大而减小.

    则下列选项正确的是 \dotfill{(\kh{A})}
    
    \xx{命题 \ding{172} 不正确，命题 \ding{173} 正确}{命题 \ding{172} 、命题 \ding{173} 都不正确}{命题 \ding{172} 正确，命题 \ding{173} 不正确}{命题 \ding{172} 、命题 \ding{173} 都正确}
    \begin{jieda}
        命题 \ding{172} 错误,理由如下: $\vv{EF} \cdot  \vv{AD}$ 的大小等价于 $\vv{EF}$ 在 $\vv{AD}$ 上的投影，由图可知，投影是不断增大的；
        
        命题 \ding{173}  正确,理由如下: 
        延长$BA, CD$交于$G$，易知$\triangle GBC$为正三角形，$\dfrac{\dfrac{m^2+1}{2}}{m^2} = \dfrac{S_{\triangle GEF}}{S_{\triangle GBC}} = \dfrac{|GE|\cdot|GF|}{|GB| \cdot |GC|} = \dfrac{[x(m-1)+1][t(m-1)+1]}{m^2}$，故$\dfrac{m^2+1}{2} = [x(m-1)+1][t(m-1)+1]$

        故$\left\{\begin{array}{l}
            2x = \dfrac{(m+1) - 2t}{t(m-1)+1} \\
            2t = \dfrac{(m+1) - 2x}{x(m-1)+1} 
        \end{array}\right.$，根据反比例函数性质可知均为严减，代入边界值可知$2x \geq \dfrac{m-1}{m}$
        \begin{center}
        \includegraphics[width=0.2\textwidth]{bodyxb/src/8_1407_1815_169_253_0.jpg}
        \end{center}
    \end{jieda}
    \item (2025$\cdot$浦东新区二模$\cdot$16)已知圆锥曲线 $\Gamma$ 的对称中心为原点 $O$ ，若对于 $\Gamma$ 上的任意一点 $A$ ，均存在 $\Gamma$ 上两点 $B$ 、 $C$ ，使得原点 $O$ 到直线 ${AB}$ 、 ${AC}$ 和 ${BC}$ 的距离都相等，则称曲线 $\Gamma$ 为“完美曲线”. 现有如下两个命题: \ding{172}  任意椭圆都是“完美曲线”； \ding{173}  存在双曲线是“完美曲线”. 下列判断正确的是 \dotfill{(\kh{A})}
    \xx{\ding{172} 是真命题， \ding{173} 是假命题}{\ding{172} 是假命题， \ding{173} 是真命题}{\ding{172}  \ding{173} 都是真命题}{\ding{172}  \ding{173} 都是假命题}
    \begin{jieda}
        对于命题 \ding{172} : 已知过椭圆上任意一点 $A$ 作以原点为圆心的圆的切线,分别交椭圆于 $B\text{ 、 }C$ 两点,连接 ${BC}$ . 根据直线与圆的位置关系,当 ${BC}$ 与圆相切时,满足给定条件. 当 ${BC}$ 与圆相交时,因为圆的圆心是固定的原点,我们可以通过缩小圆的半径,使得圆逐渐靠近 ${BC}$ ,直到 ${BC}$ 与圆相切; 同理,当 ${BC}$ 与圆相离时,扩大圆的半径,也能使圆靠近 ${BC}$ 直至相切. 所以从直线与圆位置关系的动态调整角度可知, 一定能找到合适的圆半径使得 ${BC}$ 与圆相切,故  \ding{172}  正确.
        
        对于命题 \ding{173} :当 $A$ 在双曲线顶点时，过 $A$ 作圆的切线，交双曲线于另外两点 $B$ 、 $C$ . 由双曲线的性质可知，双曲线在顶点附近的形状特点决定了，过顶点作圆的切线与双曲线相交得到的线段 ${BC}$ ,其整体位置与以原点为圆心的圆是相离的,这是因为双曲线的渐近线性质以及顶点处的曲线走向, 使得从顶点出发的切线与双曲线相交形成的线段不会与圆相切, 所以  \ding{173} 不正确.

        \begin{center}
        \includegraphics[width=0.6\textwidth]{bodyxb/src/9_228_769_513_194_0.jpg}
        \end{center}
    \end{jieda}
    
    \item (2024$\cdot$青浦区一模$\cdot$16)定义:如果曲线段 $C$ 可以一笔画出，那么称曲线段 $C$ 为单轨道曲线，比如圆、椭圆都是单轨道曲线；如果曲线段 $C$ 由两条单轨道曲线构成，那么称曲线段 $C$ 为双轨道曲线. 对于曲线 $\Gamma  : \sqrt{{\left( x + 1\right) }^{2} + {y}^{2}} \cdot  \sqrt{{\left( x - 1\right) }^{2} + {y}^{2}} = m\left( {m > 0}\right)$ 有如下命题:

    $p$ : 存在常数 $m$ ,使得曲线 $\Gamma$ 为单轨道曲线; $q$ : 存在常数 $m$ ,使得曲线 $\Gamma$ 为双轨道曲线.

    下列判断正确的是 \dotfill{(\kh{A})}
    \xx{$p$ 和 $q$ 均为真命题}{$p$ 和 $q$ 均为假命题}{$p$ 为真命题， $q$ 为假命题}{$p$ 为假命题， $q$ 为真命题}
    \begin{jieda}
        记 $f\left( {x,y}\right)  = \sqrt{{\left( x + 1\right) }^{2} + {y}^{2}} \cdot  \sqrt{{\left( x - 1\right) }^{2} + {y}^{2}} - m$ , 易得 $f\left( {-x,y}\right)  = f\left( {x, - y}\right)  = f\left( {-x, - y}\right)  = f\left( {x,y}\right)  = 0$ , 因此曲线 $f\left( {x,y}\right)  = 0$ 关于 $x$ 轴、 $y$ 轴成轴对称,关于原点成中心对称,从几何上讲,曲线 $\Gamma$ 是到两定点 $\left( {1,0}\right)$ 和 $\left( {-1,0}\right)$ 的距离乘积为 $m$ 的点的轨迹,由 $f\left( {x,y}\right)  = 0$ 可得 $y =  \pm  \sqrt{\sqrt{4{x}^{2} + {m}^{2}} - {x}^{2} - 1}$ ,因此它在 $x$ 轴上方和下方分别是两个函数的图像,这两个函数图像在 $x$ 轴上有公共点 (方程 $y = 0$ 的解相同),由 $\sqrt{4{x}^{2} + {m}^{2}} - {x}^{2} - 1 \geq  0$ 得 $1 - m \leq  {x}^{2} \leq  1 + m,0 < m < 1$ 时, $- \sqrt{1 + m} \leq  x \leq   - \sqrt{1 - m}$ 或 $\sqrt{1 - m} \leq  x \leq  \sqrt{1 + m}$ ,所以曲线 $\Gamma$ 与 $y$ 轴无公共点,曲线 $\Gamma$ 是在 $y$ 轴两侧的两个曲线构成,是双轨道曲线; 当 $m > 1$ 时, $- \sqrt{1 + m} \leq  x \leq  \sqrt{1 + m}$ ,结合对称性知,曲线 $\Gamma$ 是一个封闭曲线,是单轨道曲线 (实际上上述过程中只要对 $m$ 取一个特定值讨论即可),命题 $p\text{ 、 }q$ 均正确.
        
        \ps 卡西尼卵形线.
    \end{jieda}

    \item (2025$\cdot$虹口区一模$\cdot$16)设数列 $\left\{  {a}_{n}\right\}$ 的前四项分别为 ${a}_{1}$ 、 ${a}_{2}$ 、 ${a}_{3}$ 、 ${a}_{4}$ ，对于以下两个命题，说法正确的是 \dotfill{(\kh{A})}

    \ding{172}  存在等比数列 $\left\{  {a}_{n}\right\}$ 以及锐角 $\alpha$ ，使 $\left\{  {\sin \alpha ,\cos \alpha ,\tan \alpha }\right\}   = \left\{  {{a}_{1},{a}_{2},{a}_{3}}\right\}$ 成立；

    \ding{173}  对任意等差数列 $\left\{  {a}_{n}\right\}$ 以及锐角 $\alpha$ ，均不能使 $\left\{  {\sin \alpha ,\cos \alpha ,\tan \alpha ,\cot \alpha }\right\}   = \left\{  {{a}_{1},{a}_{2},{a}_{3},{a}_{4}}\right\}$ 成立.
    
    \xx{\ding{172} 是真命题， \ding{173} 是真命题}{\ding{172} 是真命题， \ding{173} 是假命题}{\ding{172} 是假命题， \ding{173} 是真命题}{\ding{172} 是假命题， \ding{173} 是假命题}

    \begin{jieda}
        \ding{172}正确即$sin\alpha \cdot tan\alpha = cos^2\alpha$，即$sin^2\alpha = cos^3\alpha$
        
        设$f(\alpha) = sin^2\alpha - cos^3\alpha$在$\alpha \in \left(0, \dfrac{\pi}{2}\right)$严增，$f(0) = -1, f\left(\dfrac{\pi}{2}\right) = 1$，根据函数零点存在定理可知$f(\alpha)$在$\left(0, \dfrac{\pi}{2}\right)$有零点，故\ding{172}正确

        假设存在，则$sin\alpha + cot\alpha = cos\alpha + tan\alpha$，即$sin^2\alpha cos\alpha + cos^2\alpha = cos^2\alpha sin\alpha + sin^2\alpha$，即$sin\alpha cos\alpha (sin\alpha - cos\alpha) = sin^2\alpha - cos^2\alpha$，即$sin\alpha cos\alpha = sin\alpha + cos\alpha$

        前者值域为$\left(0, 1\right]$，后者值域为$\left(1, \dfrac{\sqrt{2}}{2}\right]$，方程无解，故\ding{173}正确.
    \end{jieda}
\end{enumerate}

\section*{第 2 章 不 等 式}

\begin{center}

\begin{tabular}{|p{4cm}|p{12cm}|}
\hline
考点 & 课标要求与命题趋势 \\
\cline{1-2}
1. 等式与不等式的性质 & 梳理并掌握等式的性质，在理解不等式的概念的基础上掌握其性质，并利用不等式的性质比较大小 \\
\cline{1-2}
2. 不等式的求解 & \ding{172}  从函数观点看一元二次方程. 会结合一元二次函数的图像，判断一元二次方程实根的存在性及实根的个数, 了解函数的零点与方程根的关系 \ding{173}  从函数观点看一元二次不等式. (1) 经历从实际情境中抽象出一元二次不等式的过程. 了解一元二次不等式的现实意义. 能借助一元二次函数求解一元二次不等式,并能用集合和区间表示一元二次不等式的解集. (2) 借助一元二次函数的图像，了解一元二次不等式与相应函数、方程的联系 \ding{174}  会解一元一次不等式组、简单分式不等式、绝对值不等式和简单的指数、对数不等式 \\
\cline{1-2}
3. 基本不等式及其应用 & \ding{172}  理解、掌握基本不等式及其推论，结合具体实例，能用基本不等式解决简单的最大值或最小值问题; 会使用应用条件: “一正,二定,三相等”,能正确处理常数“1”求最值,能用拼凑等思想合理使用基本不等式求最值 \ding{173}  能够从函数观点认识方程和不等式，感悟数学知识之间的关联，认识函数的重要性. 掌握等式与不等式的性质, 以及基本不等式及其应用 \ding{174}  本节内容是上海高考的常考内容，一般会结合条件等式考查拼凑思想来使用基本不等式求最值, 或者和其他版块关联, 难度中等偏上 \\
\cline{1-2}
\hline
\end{tabular}
\end{center}
\clearpage
\subsection*{考点1 等式与不等式的性质}

\subsubsection*{一、填空题}
\begin{enumerate}
    \item (2025$\cdot$杨浦区一模$\cdot$5)已知 $\left| {x + 3}\right|  + \left| {x - 5}\right|  = 8$ ，则实数 $x$ 的取值范围为\blkda[2]{$\left\lbrack  {-3,5}\right\rbrack$}.
    \item (2024$\cdot$虹口区一模$\cdot$11)设 $a \in  \mathbf{R}$ ，若关于 $x$ 的方程 ${2x}\left| x\right|  - \left( {a - 2}\right) x + \left| x\right|  - a + 1 = 0$ 有 3 个不同的实数解，则实数 $a$ 的取值范围为\blkda[2]{$\left( {9, + \infty }\right)$}.
    \begin{jieda}
        ${2x}\left| x\right|  - \left( {a - 2}\right) x + \left| x\right|  - a + 1 = 0 \Longleftrightarrow (2x+1)(|x|+1) = a(x+1) \\
        \Longleftrightarrow a = \left\{\begin{array}{l}
            2x+1, x \geq 0 \\
            -2\left(x+1 + \dfrac{1}{x+1}\right)+5
        \end{array}\right.$，数形结合可知$a \in \left( {9, + \infty }\right)$时满足要求.
    \end{jieda}
\end{enumerate}

\subsubsection*{二、选择题}
\begin{enumerate}
    \item (2024$\cdot$松江区一模$\cdot$13)英国数学家哈利奥特最先使用“<”和“>”符号，并逐渐被数学界接受，不等号的引用对不等式的发展影响深远，对任意实数 $a$ 、 $b$ 、 $c$ 、 $d$ ，下列命题是真命题是 \dotfill{(\kh{D})}
    \xx{若 ${a}^{2} < {b}^{2}$ ,则 $a < b$}{若 $a < b$ ,则 ${ac} < {bc}$}{若 $a < b,c < d$ ,则 ${ac} < {bd}$}{若 $a < b,c < d$ ,则 $a + c < b + d$}

    \item (2024$\cdot$上海春季卷$\cdot$13)已知 $a$ 、 $b$ 、 $c \in  \mathbf{R}$ ， $b > c$ ，则下列不等式恒成立的是 \dotfill{(\kh{B})}\xx{$a + {b}^{2} > a + {c}^{2}$}{${a}^{2} + b > {a}^{2} + c$}{$a{b}^{2} > a{c}^{2}$}{${a}^{2}b > {a}^{2}c$}

    \item (2025$\cdot$崇明区二模$\cdot$13)若 $a > b > 0,c > d$ ，则下列结论正确的是 \dotfill{(\kh{D})}
    \xx{$a - b < 0$}{${ac} > {bd}$}{$a{c}^{2} > b{c}^{2}$}{$\dfrac{a}{{c}^{2} + 1} > \dfrac{b}{{c}^{2} + 1}$}

    \item (2024$\cdot$崇明区二模$\cdot$13)若 $a > b$ ， $c < 0$ ，则下列不等式成立的是 \dotfill{(\kh{D})}
    \xx{$a > b - c$}{$\dfrac{a}{c} > \dfrac{b}{c}$}{$a + c < b + c$}{$a{c}^{2} > b{c}^{2}$}

    \item (2025$\cdot$嘉定区二模$\cdot$13)已知实数 $a$ 、 $b$ 满足 $a > b$ ，则下列不等式中，不恒成立的是 \dotfill{(\kh{B})}
    \xx{${a}^{3} > {b}^{3}$}{$\dfrac{a}{b} > 1$}{${a}^{2} + {b}^{2} > {2ab}$}{${2}^{a} > {2}^{b}$}

    \item (2025$\cdot$闵行区二模$\cdot$13)已知 $a \in  \mathbf{R}$ ，则下列结论不恒成立的是 \dotfill{(\kh{B})}
    \xx{$a + 1 > a$}{$a + \dfrac{1}{a} \geq  2$}{$\left| {1 - a}\right|  + \left| {a + 2}\right|  \geq  3$}{${a}^{2} + a + 1 > 0$}

    \item (2024$\cdot$闵行区一模$\cdot$13)已知 $a$ 、 $b \in  \mathbf{R}$ ， $a > b$ ，则下列不等式中不一定成立的是 \dotfill{(\kh{C})}
    \xx{$a + 2 > b + 2$}{${2a} > {2b}$}{${a}^{2} > {b}^{2}$}{${2}^{a} > {2}^{b}$}

    \item (2025$\cdot$浦东新区一模$\cdot$13)若实数 $a$ 、 $b$ 满足 ${a}^{2} > {b}^{2}$ ，下列不等式中恒成立的是 \dotfill{(\kh{C})}
    \xx{$a > b$}{$\dfrac{1}{a} < \dfrac{1}{b}$}{${a}^{2} + {b}^{2} > {2ab}$}{$a > \left| b\right|$}

    \item (2025$\cdot$青浦区二模$\cdot$14)若正数$m$、$n$、$a$均不为1，则下列不等式中与“$m>n$”等价的是\dotfill{(\kh{C})}
    \xx{${a}^{m} > {a}^{n}$}{${\log }_{a}m > {\log }_{a}n$}{${m}^{a} > {n}^{a}$}{${\log }_{m}a > {\log }_{n}a$}

    \item (2024$\cdot$杨浦区二模$\cdot$14)已知实数 $a$ 、 $b$ 、 $c$ 、 $d$ 满足 $a > b > 0 > c > d$ ，则下列不等式一定正确的是\dotfill{(\kh{C})}
    \xx{$a + d > b + c$}{${ad} > {bc}$}{$a + c > b + d$}{${ac} > {bd}$}

    \item (2024$\cdot$崇明区一模$\cdot$14)若 $x > y > 0$ ，则下列不等式正确的是 \dotfill{(\kh{C})}
    \xx{$\left| x\right|  < \left| y\right|$}{${x}^{2} < {y}^{2}$}{$\dfrac{1}{x} < \dfrac{1}{y}$}{$\dfrac{x + y}{2} \leq  \sqrt{xy}$}
\end{enumerate}

\subsection*{考点2 不等式的求解}

\subsubsection*{一、填空题}
\begin{enumerate}
    \item (2025$\cdot$黄浦区二模$\cdot$1)不等式 $\dfrac{x}{x - 2} < 0$ 的解集是\blkda[2]{$\left( {0,2}\right)$}.
    \item (2025$\cdot$浦东新区二模$\cdot$1)已知 $x \in  \mathbf{R}$ ，不等式 $\dfrac{x - 2}{x} < 0$ 的解为\blkda[2]{$\left( {0,2}\right)$}.
    \item (2025$\cdot$普陀区二模$\cdot$1)不等式 $\left| {1 + \dfrac{1}{x}}\right|  < 0$ 的解集是\blkda[2]{$\left( {-1,0}\right)$}.
    \item (2025$\cdot$崇明区二模$\cdot$1)不等式 $\left| {x - 1}\right|  < 2$ 的解为\blkda[2]{$\left( {-1,3}\right)$}.
    \item (2023$\cdot$上海秋季卷$\cdot$1)不等式 $\left| {x - 2}\right|  < 1$ 的解集为\blkda[2]{$\left( {1,3}\right)$}.
    \item (2025$\cdot$徐汇区一模$\cdot$1)不等式 ${x}^{2} - {4x} + 3 < 0$ 的解集为\blkda[2]{$\left( {1,3}\right)$}.
    \item (2024$\cdot$徐汇区一模$\cdot$2)不等式 $\dfrac{1}{x} > 1$ 的解集是\blkda[2]{$\left( {0,1}\right)$}.
    \item (2024$\cdot$嘉定区一模$\cdot$1)不等式 ${x}^{2} - x - 6 < 0$ 的解集为\blkda[2]{$\left( {-2,3}\right)$}.
    \item (2024$\cdot$崇明区二模$\cdot$2)不等式 $x\left( {x - 1}\right)  < 0$ 的解为\blkda[2]{$\left( {0,1}\right)$}.
    \item (2025$\cdot$静安区一模$\cdot$2)不等式 $\left| {{2x} - 1}\right|  < 3$ 的解集为\blkda[2]{$\left( {-1,2}\right)$}.
    \item (2025$\cdot$上海秋季卷$\cdot$2)不等式 $\dfrac{x - 1}{x - 3} < 0$ 的解集为\blkda[2]{$\left( {1,3}\right)$}.
    \item (2025$\cdot$杨浦区二模$\cdot$2)不等式 $\dfrac{x + 1}{x - 2} < 0$ 的解集为\blkda[2]{$\left( {-1,2}\right)$}.
    \item (2025$\cdot$闵行区二模$\cdot$2)设 $x \in  \mathbf{R}$ ，则不等式 $\left| {x - 2}\right|  \leq  5$ 的解集为\blkda[2]{$\left\lbrack  {-3,7}\right\rbrack$}.
    \item (2025$\cdot$嘉定区二模$\cdot$2)不等式 $\dfrac{x - 2}{x + 1} < 0$ 的解集为\blkda[2]{$\left( {-1,2}\right)$}.
    \item (2025$\cdot$静安区二模$\cdot$2)不等式 $\dfrac{{2x} - 1}{x + 2} < 0$ 的解集为\blkda[2]{$\left( {-2,\dfrac{1}{2}}\right)$}.
    \item (2025$\cdot$虹口区二模$\cdot$2)不等式 $\dfrac{x - 2}{x + 1} \leq  0$ 的解集是\blkda[2]{$\left( {-1,2}\right)$}.
    \item (2025$\cdot$上海春季卷$\cdot$2)不等式 $\dfrac{x}{x - 1} < 0$ 的解集为\blkda[2]{$\left( {0,1}\right)$}.
    \item (2024$\cdot$青浦区二模$\cdot$2)不等式 $\left| {x - 2}\right|  > 1$ 的解集为\blkda[3]{$\left( {-\infty ,1}\right)  \cup  \left( {3, + \infty }\right)$}.
    \item (2025$\cdot$闵行区一模$\cdot$2)不等式 $\dfrac{{2x} - 1}{x - 1} < 0$ 的解集为\blkda[2]{$\left( {\dfrac{1}{2},1}\right)$}.
    \item (2025$\cdot$黄浦区一模$\cdot$2)不等式 ${x}^{2} - {3x} + 2 < 0$ 的解集为\blkda[2]{$\left( {1,2}\right)$}.
    \item (2025$\cdot$崇明区一模$\cdot$2)不等式 $\dfrac{x - 1}{{2x} + 1} < 0$ 的解为\blkda[2]{$\left( {-\dfrac{1}{2},1}\right)$}.
    \item (2023$\cdot$上海春季卷$\cdot$3)若不等式 $\left| {x - 1}\right|  \leq  2$ ，则 $x$ 的取值范围是\blkda[2]{$[-1, 3]$}.
    \item (2024$\cdot$上海秋季卷$\cdot$3)已知 $x \in  \mathbf{R}$ ，则不等式 ${x}^{2} - {2x} - 3 < 0$ 的解集为\blkda[2]{$(-1, 3)$}.
    \item (2025$\cdot$奉贤区一模$\cdot$3)已知 $x \in  \mathbf{R}$ ，则不等式 ${x}^{2} - x + 2 > 0$ 的解集为\blkda[2]{$\mathbf{R}$}.
    \item (2025$\cdot$青浦区一模$\cdot$3)不等式 $\dfrac{x - 3}{x + 1} \geq  2$ 的解集为\blkda[2]{$\lbrack  - 5, - 1)$}.
    \item (2024$\cdot$金山区一模$\cdot$3)不等式 $\dfrac{x - 1}{x + 2} > 0$ 的解集是\blkda[3.5]{$(-\infty, -2) \cup (1, +\infty)$}.
    \item (2025$\cdot$杨浦区一模$\cdot$3)不等式 $\dfrac{x + 2}{x - 1} < 0$ 的解集为\blkda[2]{$\left( {-2,1}\right)$}.
    \item (2025$\cdot$上海春季卷$\cdot$8)关于 $x$ 的方程 $\left| {x - 1}\right|  + \left| {\pi  - x}\right|  = \pi  - 1$ 的解集为\blkda[2]{$\left\lbrack  {1,\pi }\right\rbrack$}.
    \item (2025$\cdot$浦东新区一模$\cdot$8)已知函数 $y = f\left( x\right)$ 的表达式为 $f\left( x\right)  = \left\{  \begin{array}{l} {\left( \dfrac{1}{2}\right) }^{x} - 3,x \leq  0, \\  {x}^{2},x > 0, \end{array}\right.$ 则不等式 $f\left( x\right)  \leq  1$ 的解集为\blkda[2]{$\left\lbrack  {-2,1}\right\rbrack$}.
    \item (2025$\cdot$杨浦区二模$\cdot$8)不等式 $\left| {x + 3}\right|  + \left| {a - x}\right|  > 6$ 对一切实数 $x$ 恒成立，则实数 $a$ 的取值范围为\blkda[3.5]{$\left( {-\infty , - 9}\right)  \cup  \left( {3, + \infty }\right)$}.
    \item (2024$\cdot$静安区一模$\cdot$8)若不等式 $\left| {x - 3}\right|  + \left| {x - 5}\right|  \geq  a$ 对所有实数 $x$ 恒成立，则实数 $a$ 的取值范围是\blkda[2]{$( - \infty ,2\rbrack$}.
    \item (2024$\cdot$静安区一模$\cdot$10)不等式 ${\log }_{2}x + \dfrac{x}{2} < 4$ 的解集为\blkda[2]{$\left( {0,4}\right)$}.
    \item (2024$\cdot$虹口区二模$\cdot$12)已知关于 $x$ 的不等式 $\left( {\ln x - {kx}}\right) \left\lbrack  {{x}^{2} - \left( {k + 3}\right) x + 4}\right\rbrack   \leq  0$ 对任意 $x \in  \left( {0, + \infty }\right)$ 均成立，则实数 $k$ 的取值范围为\blkda[2]{$\left\lbrack  {\dfrac{1}{{e}},1}\right\rbrack$}.
    \begin{jieda}
        $\left( {\ln x - {kx}}\right) \left\lbrack  {{x}^{2} - \left( {k + 3}\right) x + 4}\right\rbrack   \leq  0$ 对任意 $x \in  (0$ , $+ \infty )$ 均成立
        
        即$\left( {\dfrac{\ln x}{x} - k}\right) \left\lbrack  {x + \dfrac{4}{x} - \left( {k + 3}\right) }\right\rbrack   \leq  0$在$(0, +\infty)$恒成立
        
        即$ \left( {k - \dfrac{\ln x}{x}}\right)  \left\lbrack  {k - \left( {x + \dfrac{4}{x} - 3}\right) }\right\rbrack   \leq  0$在$(0, +\infty)$恒成立
        
        即$\dfrac{\ln x}{x} \leq  k \leq  x + \dfrac{4}{x} - 3$在$(0, +\infty)$恒成立
        
        ${y}_{1} = \dfrac{\ln x}{x} \leq  \dfrac{1}{{e}},{y}_{2} = x + \dfrac{4}{x} - 3 \geq  2\sqrt{x \times  \dfrac{4}{x}} - 3 = 1$ , 所以实数 $k$ 的取值范围为 $k \in  \left\lbrack  {\dfrac{1}{{e}},1}\right\rbrack$ .

        \begin{center}
        \includegraphics[width=0.2\textwidth]{bodyxb/src/12_423_1184_221_274_0.jpg}
        \end{center}
    \end{jieda}
\end{enumerate}
\subsection*{考点3 基本不等式及其应用}

\subsubsection*{一、填空题}
\begin{enumerate}
    \item (2024$\cdot$闵行区一模$\cdot$3)若 ${xy} = 1$ $(x, y \in  \mathbf{R})$，则 ${x}^{2} + 4{y}^{2}$ 的最小值为\blkda[2]{4}.
    \begin{jieda}
        ${x}^{2} + 4{y}^{2} \geq  {2x} \cdot  {2y} = {4xy} = 4$ ,当且仅当 $x = {2y} =  \pm  \sqrt{2}$ 时，上式取得最小值 4.
    \end{jieda}

    \item (2025$\cdot$闵行区一模$\cdot$4)已知正实数 $a$ 、 $b$ 满足 ${ab} = 1$ ，则 $\dfrac{1}{a} + \dfrac{1}{b}$ 的最小值为\blkda[2]{2}.
    \begin{jieda}
        $\dfrac{1}{a} + \dfrac{1}{b} = a+b \geq  2\sqrt{ab} = 2$ ,当且仅当 $a = b = 1$ 时取等号,故 $\dfrac{1}{a} + \dfrac{1}{b}$ 的最小值为2.
    \end{jieda}

    \item (2024$\cdot$徐汇区一模$\cdot$4)若实数 $x$ 、 $y$ 满足 $x + y = 2$ ，则 ${x}^{2} + {y}^{2}$ 的最小值为\blkda[2]{2}.
    \begin{jieda}
        ${x}^{2} + {y}^{2} = {\left( x + y\right) }^{2} - 2xy = 4 - 2xy \geq 4 - (x^2+y^2) \Longrightarrow  {x}^{2} + {y}^{2} \geq  2$ .
    \end{jieda}

    \item (2024$\cdot$徐汇区二模$\cdot$4)若正数 $a$ 、 $b$ 满足 $\dfrac{1}{a} + \dfrac{1}{b} = 1$ ，则 ${2a} + b$ 的最小值为\blkda[2]{$3 + 2\sqrt{2}$}.
    \begin{jieda}
        由已知 ${2a} + b = \left( {{2a} + b}\right) \left( {\dfrac{1}{a} + \dfrac{1}{b}}\right)  = 3 + \dfrac{2a}{b} + \dfrac{b}{a} \geq 3 + 2\sqrt{2}$ ,当且仅当 $\dfrac{2a}{b} = \dfrac{b}{a}$ ,即 $a = 1 + \dfrac{\sqrt{2}}{2},b = 1 + \sqrt{2}$ 时等号成立，故所求最小值是 $3 + 2\sqrt{2}$.
    \end{jieda}

    \item (2025$\cdot$杨浦区二模$\cdot$5)已知 $a > 0$ ， $b > 0$ ， $a + b = 2$ ，则 ${ab}$ 的最大值为\blkda[2]{1}.
    \begin{jieda}
        由基本不等式得 ${ab} \leq  {\left( \dfrac{a + b}{2}\right) }^{2} = 1$
    \end{jieda}

    \item (2024$\cdot$闵行区二模$\cdot$5)正实数 $a$ 、 $b$ 满足 $a + {2b} = 1$ ，则 ${ab}$ 的最大值为\blkda[2]{$\dfrac{1}{8}$}.
    \begin{jieda}
        ${2ab} \leq  {\left( \dfrac{a + {2b}}{2}\right) }^{2} \Longrightarrow  {ab} \leq  \dfrac{1}{8}$.
    \end{jieda}

    \item (2024$\cdot$普陀区二模$\cdot$5)若实数 $a$ 、 $b$ 满足 $a - {2b} \geq  0$ ，则 ${2}^{a} + \dfrac{1}{{4}^{b}}$ 的最小值为\blkda[2]{2}.
    \begin{jieda}
        因为 ${2}^{a} > 0,\dfrac{1}{{4}^{b}} > 0,a - {2b} \geq  0$ ,所以 ${2}^{a} + \dfrac{1}{{4}^{b}} = {2}^{a} +  \dfrac{1}{{2}^{2b}} \geq  2\sqrt{{2}^{a} \cdot  \dfrac{1}{{2}^{2b}}} = 2\sqrt{{2}^{a - {2b}}} \geq  2\sqrt{{2}^{0}} = 2$ ,当且仅当 ${2}^{a} = \dfrac{1}{{2}^{2b}}$ , 即 $a = b = 0$ 时等号成立，所以 ${2}^{a} + \dfrac{1}{{4}^{b}}$ 的最小值为 2.
    \end{jieda}

    \item (2023$\cdot$上海春季卷$\cdot$6)已知正实数 $a$ 、 $b$ 满足 $a + {4b} = 1$ ，则 ${ab}$ 的最大值为\blkda[2]{$\dfrac{1}{16}$}.
    \begin{jieda}
        $a + {4b} = 1 \geq  2\sqrt{4ab} \Longrightarrow  0 < {ab} \leq  \dfrac{1}{16}$.
    \end{jieda}

    \item (2025$\cdot$浦东新区一模$\cdot$6)已知实数 $a$ 、 $b$ 满足 $a + {2b} = 1$ ，则 ${3}^{a} + {9}^{b}$ 的最小值为\blkda[2]{$2\sqrt{3}$}.
    \begin{jieda}
        因为 $a + {2b} = 1$ ,所以 ${3}^{a} + {9}^{b} = {3}^{a} + {3}^{2b} \geq  2\sqrt{{3}^{a} \times  {3}^{2b}} =  2\sqrt{{3}^{a + {2b}}} = 2\sqrt{3}$ ,当且仅当 ${3}^{a} = {3}^{2b}$ ,即 $a = \dfrac{1}{2},b = \dfrac{1}{4}$ 时取等号，故 ${3}^{a} + {9}^{b}$ 的最小值为 $2\sqrt{3}$.
    \end{jieda}

    \item (2024$\cdot$松江区一模$\cdot$6)已知 $\lg a + \lg b = 1$ ，则 $a + {2b}$ 的最小值为\blkda[2]{$4\sqrt{5}$}.
    \begin{jieda}
        $\lg a + \lg b = 1 \Longrightarrow  {ab} = {10},a + {2b} \geq  2\sqrt{2ab} = 4\sqrt{5}$
    \end{jieda}

    \item (2025$\cdot$崇明区一模$\cdot$6)已知 $x$ 、 $y$ 为正实数，且满足 ${4x} + y = {40}$ ，则 ${xy}$ 的最大值是\blkda[2]{100}.
    \begin{jieda}
        因为 ${4x} + y = {40}$ ,所以 ${xy} = \dfrac{1}{4} \times  {4x} \times  y \leq  \dfrac{1}{4} \times  {\left( \dfrac{{4x} + y}{2}\right) }^{2} = {100}$ ,当且仅当 ${4x} = y$ ,即 $x = 5,y = {20}$ 时,等号成立. 即 ${xy}$ 的最大值为 100.
    \end{jieda}

    \item (2025$\cdot$黄浦区一模$\cdot$6)若正数 $x$ 、 $y$ 满足 $x + {4y} = 1$ ，则 ${xy}$ 的最大值为\blkda[2]{$\dfrac{1}{16}$}.
    \begin{jieda}
        \begin{dingautolist}{172}
            \item 因为正数 $x\text{ 、 }y$ 满足 $x + {4y} = 1$ ,所以 $x = 1 - {4y} \Longrightarrow  {xy} = y\left( {1 - {4y}}\right)  =  - 4{y}^{2} + y =  - 4{\left( y - \dfrac{1}{8}\right) }^{2} + \dfrac{1}{16}$ ,所以当 $y = \dfrac{1}{8}$ 时, ${xy}$ 最大值为 $\dfrac{1}{16}$.
            \item 因为 $x > 0,y > 0$ ,所以 ${xy} = \dfrac{1}{4} \cdot  x \cdot  {4y} \leq  \dfrac{1}{4} \cdot  {\left( \dfrac{x + {4y}}{2}\right) }^{2} = \dfrac{1}{4} \cdot  \dfrac{1}{4} = \dfrac{1}{16}$ ,所以 ${\left( xy\right) }_{\max } = \dfrac{1}{16}.$
        \end{dingautolist}
    \end{jieda}

    \item (2024$\cdot$嘉定区一模$\cdot$7)已知实数 $a$ 、 $b$ 满足 ${ab} =  - 6$ ，则 ${a}^{2} + {b}^{2}$ 的最小值为\blkda[2]{12}.
    \begin{jieda}
        ${a}^{2} + {b}^{2} \geq  2\left| {ab}\right|  = {12}$ .
    \end{jieda}

    \item (2024$\cdot$上海春季卷$\cdot$7)已知实数 ${ab} = 1$ ，则 $4{a}^{2} + 9{b}^{2}$ 的最小值为\blkda[2]{12}.
    \begin{jieda}
        由 $4{a}^{2} + 9{b}^{2} \geq  2\sqrt{{36}{a}^{2}{b}^{2}} = {12}$ ,当且仅当 ${2a} = {3b}$ 时,即 $a = \dfrac{\sqrt{6}}{2},b = \dfrac{\sqrt{6}}{3}$ 或 $a =  - \dfrac{\sqrt{6}}{2},b =  -\dfrac{\sqrt{6}}{3}$ 时 $4{a}^{2} +  9{b}^{2}$ 取最小值为 12.
    \end{jieda}

    \item (2024$\cdot$奉贤区二模$\cdot$7)某商品的成本 $C$ 与产量 $q$ 之间满足关系式 $C = C\left( q\right)$ ，定义平均成本 $\overline{C} = \dfrac{C\left( q\right) }{q}$ ，假设 $C\left( q\right)  = \dfrac{1}{4}{q}^{2} + {100}$ ，当产量等于\blkda[2]{20}时，平均成本最少.    
    \begin{jieda}
        由题知 $\overline{C} = \dfrac{C\left( q\right) }{q} = \dfrac{\dfrac{1}{4}{q}^{2} + {100}}{q} = \dfrac{q}{4} + \dfrac{100}{q} \geq  2\sqrt{\dfrac{q}{4} \cdot  \dfrac{100}{q}} = {10}$ ,当且仅当 $\dfrac{q}{4} = \dfrac{100}{q}$ ,即 $q = {20}$ 时取等号, 故答案为:20.
    \end{jieda}

    \item (2025$\cdot$静安区一模$\cdot$8)若用 $t$ 替换命题“对于任意实数 $d$ ，有 ${d}^{2} \geq  0$ ，且等号当且仅当 $d = 0$ 时成立”中的 $d$ ，即可推出平均值不等式“任意两个正数的算术平均值不小于它们的几何平均值，且等号当且仅当这两个正数相等时成立”. 则 $t =$ \blkda[2]{$\sqrt{a} - \sqrt{b}$}.

    \item (2025$\cdot$上海秋季卷$\cdot$8)设 $a$ 、 $b > 0$ ， $a + \dfrac{1}{b} = 1$ ，则 $b + \dfrac{1}{a}$ 的最小值为\blkda[2]{4}.
    \begin{jieda}
        易知 $b + \dfrac{1}{a} = \left( {b + \dfrac{1}{a}}\right) \left( {a + \dfrac{1}{b}}\right)  = {ab} + \dfrac{1}{ab} + 2 \geq  2\sqrt{{ab} \cdot  \dfrac{1}{ab}} + 2 = 4$ ,当且仅当 ${ab} = 1$ ,即 $a = \dfrac{1}{2},b = 2$ 时取得最小值.
    \end{jieda}

    \item (2024$\cdot$静安区二模$\cdot$8)在下列关于实数 $a$ 、 $b$ 的四个不等式中，恒成立的是\blkda[2]{\ding{173}  \ding{174}  \ding{175}}(请填入全部正确的序号).
    \ding{172}  $a + b \geq  2\sqrt{ab}$ ;  \ding{173}  ${\left( \dfrac{a + b}{2}\right) }^{2} \geq  {ab}$ ;  \ding{174}  $\left| a\right|  - \left| b\right|  \leq  \left| {a - b}\right|$ ； \ding{175}  ${a}^{2} + {b}^{2} \geq  {2b} - 1$.

    \item (2025$\cdot$长宁区一模$\cdot$10)若正实数 $a$ 、 $b$ 满足 ${ab} = {2a} + b$ ，则 $a + {2b}$ 的最小值是\blkda[2]{9}.
    \begin{jieda}
        $\dfrac{1}{a} + \dfrac{2}{b} = 1$，$a+2b = (a+2b)\left(\dfrac{1}{a} + \dfrac{2}{b}\right) = 5 + \dfrac{2b}{a} + \dfrac{2a}{b} \geq 9$，等号成立时 $a = 3,b = 3$.
    \end{jieda}

    \item (2024$\cdot$长宁区一模$\cdot$10) 若 “存在 $x > 0$ ，使得 ${x}^{2} + {ax} + 1 < 0$”是假命题，则实数 $a$ 的取值范围为\blkda[2]{$\lbrack  - 2, + \infty )$}.
    \begin{jieda}
        由题意可得: “任意 $x > 0$ ,使得 ${x}^{2} + {ax} + 1 \geq  0$ ” 是真命题，等价于 “任意 $x > 0$ ,使得 $x + \dfrac{1}{x} \geq   - a$ ” 是真命题,因为 $x + \dfrac{1}{x} \geq  2\sqrt{x \cdot  \dfrac{1}{x}} = 2$ ,当且仅当 $x = \dfrac{1}{x}$ ,即 $x = 1$ 时,等号成立,所以 $2 \geq   - a$，所以实数 $a$ 的取值范围 $\lbrack  - 2, +\infty )$.
    \end{jieda}

    \item (2025$\cdot$静安区二模$\cdot$10)已知 $a > b$ ，且 ${ab} = 1$ ，则 $\dfrac{{a}^{2} + {b}^{2}}{a - b}$ 的最小值是\blkda[2]{$2\sqrt{2}$}.
    \begin{jieda}
        $a > 1 > b = \dfrac{1}{a}$，$\dfrac{{a}^{2} + {b}^{2}}{a - b} = \dfrac{a^2+\dfrac{1}{a^2}}{a - \dfrac{1}{a}}$，令$t = a - \dfrac{1}{a}, t > 0$，原式$ = \dfrac{t^2+2}{t} = t + \dfrac{2}{t} \geq 2\sqrt{2}$.
    \end{jieda}

    \item (2025$\cdot$青浦区二模$\cdot$11)道路通行能力指单位时间(1 小时)内通过道路上指定断面的最大车辆数, 是度量道路疏导交通能力的指标. 同时为了行驶安全, 车辆之间必须保持一定的安全距离. 为了研究某城市道路通行能力，现给出如下假设:

    假设 1: 车身长度均为 ${4.8}{m}$ ;

    假设 2: 所有车辆以相同的速度 $v$ (单位: ${{km}}/{h}$ )匀速行驶;

    假设 3: 安全距离 $d$ (单位: ${m}$ )与车辆速度 $v$ 近似满足 $d = {3.2} + {0.6522v} + {0.01}{v}^{2}$.

    该城市道路通行能力的最大值约为\blkda[2]{821}. (结果保留整数)
    \begin{jieda}
        1小时 $= {3600}$ 秒,车辆速度 $v$ (千米/小时) 换算为米/秒是 $\dfrac{v}{3.6}$ 米/秒. 1 小时内通过的车辆数$n = \dfrac{{3600} \times  \dfrac{v}{3.6}}{{4.8} + {3.2} + {0.6522v} + {0.01}{v}^{2}} = \dfrac{1000v}{8 + {0.6522v} + {0.01}{v}^{2}} =  \dfrac{1000}{{0.01v} + \dfrac{8}{v} + {0.6522}}.$

        ${0.01v} + \dfrac{8}{v} \geq  2\sqrt{{0.01v} \times  \dfrac{8}{v}} = \dfrac{2\sqrt{2}}{5},$，当且仅当 ${0.01v} = \dfrac{8}{v}$ 时等号成立.

        所以 $n \leq  \dfrac{1000}{\dfrac{2\sqrt{2}}{5} + {0.6522}} \approx  {821}$
    \end{jieda}

    \item (2025$\cdot$金山区一模$\cdot$10)某海滨浴场平面图是如图所示的半圆，其中 $O$ 是圆心,直径 ${MN}$ 为 ${400}{m},P$ 是弧 ${MN}$ 的中点. 一个急救中心 $A$ 在栈桥 ${OP}$ 中点上，计划在弧 ${NP}$ 上设置一个瞭望台 $B$ ，并在 ${AB}$ 间修建浮桥. 已知 $\angle {ABO}$ 越大，瞭望台 $B$ 处的视线范围越大，则 $B$ 处的视线范围最大时， ${AB}$ 的长度为\blkda[2]{173}$m$. (结果精确到 $1{m}$ )

    \begin{center}
    \includegraphics[width=0.2\textwidth]{bodyxb/src/16_1151_687_312_219_0.jpg}
    \end{center}
    \begin{jieda}
        令 ${AB} = x,\angle {ABO} = \theta  \in  \left( {0,\dfrac{\pi }{2}}\right)$，且 ${OA} = {100},{OB} = {200}$ ,在 $\bigtriangleup {ABO}$ 中,

        $\cos \theta  = \dfrac{{x}^{2} + {200}^{2} - {100}^{2}}{{2x} \times  {200}} = \dfrac{1}{400} \cdot  \left( {x + \dfrac{30000}{x}}\right)  \geq \dfrac{1}{400} \cdot  2\sqrt{x \cdot  \dfrac{30000}{x}} = \dfrac{\sqrt{3}}{2},$

        当且仅当 $x = {100}\sqrt{3} \approx  {173}$ 米时, $\cos \theta$ 取最小值，此时 $\angle {ABO} = \theta$ 最大.
    \end{jieda}

    \item (2025$\cdot$闵行区一模$\cdot$11)如图，某小区内有一块矩形区域 ${ABCD}$ ，其中 ${AB} = {40}{m},{AD} = {20}{m}$ ，点 $E$ 、 $F$ 分别为 ${AB}$ 、 ${CD}$ 的中点，左右两个扇形区域为花坛 (两个扇形的圆心分别为 $A$ 、 $B$ ，半径均为 ${20}{m}$ )，其余区域为草坪,现规划在草坪上修建一个三角形的儿童游乐区,且三角形的一个顶点在线段 ${EF}$ 上，另外两个顶点在线段 ${CD}$ 上，则该游乐区面积的最大值为\blkda[2]{137} ${{m}}^{2}$. (结果保留整数)

    \begin{center}
    \includegraphics[width=0.2\textwidth]{bodyxb/src/16_1156_975_310_191_0.jpg}
    \end{center}
    \begin{jieda}
        设游乐区所在的三角形为 $\bigtriangleup {PQM},M$ 在线段 ${EF}$ 上， $P$ 、 $Q$ 在线段 ${DC}$ 上，如图所示，当 ${PM}$ 、 ${QM}$ 分别与圆弧相切时， ${S}_{PQM}$ 取得最大值

        \begin{center}
        \includegraphics[width=0.2\textwidth]{bodyxb/src/13_1081_1719_226_141_0.jpg}
        \end{center}
        
        设 ${PM}$ 与圆弧相切于点 $G$，$PG = PD, MG = ME$，故$PF+PM+MF = 40$

        设$PF = a, MF = b$，则$S_{\bigtriangleup {PQM}} = ab$，$a^2+b^2 = (40-a-b)^2$，即$(a-40)(b-40) = 800$，$a = \dfrac{800}{b-40} + 40$，$ab = \dfrac{800b}{b-40} + 40b = 800+\dfrac{32000}{b-40} + 40b = 2400 - \left(40(40-b)+\dfrac{32000}{40-b}\right) \leq 2400-2\sqrt{1280000} = 2400 - 1600\sqrt{2} \approx 137$.


    \end{jieda}

\end{enumerate}
\subsubsection*{二、选择题}
\begin{enumerate}
    \item (2024$\cdot$青浦区二模$\cdot$13)函数 $y = {3x} + \dfrac{1}{x}\left( {x > 0}\right)$ 的最小值是 \dotfill{(\kh{D})}
    \xx{4}{5}{$3\sqrt{2}$}{$2\sqrt{3}$}

    \item (2025$\cdot$松江区一模$\cdot$13)已知 $a > b > 0$ ，以下四个数中最大的是 \dotfill{(\kh{D})}

    \xx{$b$}{$\sqrt{ab}$}{$\dfrac{a + b}{2}$}{$\sqrt{\dfrac{{a}^{2} + {b}^{2}}{2}}$}

    \item (2024$\cdot$浦东新区二模$\cdot$14)已知 $a \in  \mathbf{R}$ ，则下列结论不恒成立的是 \dotfill{(\kh{B})}

    \xx{$a\left( {1 - a}\right)  \leq  \dfrac{1}{4}$}{$a + \dfrac{1}{a} \geq  2$}{$\left| {a - 1}\right|  + \left| {a + 2}\right|  \geq  3$}{$\sin a + \dfrac{1}{2 + \sin a} \geq  0$}

    \item (2024$\cdot$黄浦区一模$\cdot$15)若实数 $a$ 、 $b$ 满足 ${a}^{2} + {b}^{2} = 1 + \left| {ab}\right|$ ，则必有 \dotfill{(\kh{D})}
    \xx{${a}^{2} + {b}^{2} \geq  2$}{${a}^{2} - {b}^{2} \leq  1$}{$a - b \leq  1$}{$a + b \leq  2$}
    \begin{jieda}
        对于 ${A},{a}^{2} + {b}^{2} - 1 = \left| {ab}\right|  \leq  \dfrac{{a}^{2} + {b}^{2}}{2} \Longrightarrow  {a}^{2} + {b}^{2} \leq  2$ , A 错误; 
        
        对于 B,取特殊值, $a = {5b},{a}^{2} + {b}^{2} = 1 + \left| {ab}\right|  \Longrightarrow  {26}{b}^{2} =  1 + 5{b}^{2},{b}^{2} = \dfrac{1}{21},{a}^{2} - {b}^{2} = \dfrac{8}{7} > 1$ , B 错误; 
        
        对于 ${C}$ ,取特殊值 ${a}^{2} + {b}^{2} = 1 + \left| {ab}\right| ,a = 1,b =  - 1,a - b = 2,{C}$ 错误; 
        
        对于 ${D}$ , ${a}^{2} + {b}^{2} = 1 + \left| {ab}\right|  \Longrightarrow  {\left( a + b\right) }^{2} = 1 + \left| {ab}\right|  + {2ab} \leq  1 + 3\left| {ab}\right|  \leq  1 + 3 \times  {\left( \dfrac{a + b}{2}\right) }^{2} \Longrightarrow  \dfrac{1}{4}{\left( a + b\right) }^{2} \leq  1 \Longrightarrow  {\left( a + b\right) }^{2} \leq  4 \Longrightarrow   - 2 \leq  a +  b \leq  2,{D}$ 正确. 故选 D.
    \end{jieda}

    \item (2024$\cdot$松江区一模$\cdot$16)关于曲线 $M : {x}^{\frac{1}{2}} + {y}^{\frac{1}{2}} = 1$ ，有下述两个结论: \ding{172}  曲线 $M$ 上的点到坐标原点的距离最小值是 $\dfrac{\sqrt{2}}{2}$ ； \ding{173}  曲线 $M$ 与坐标轴围成的图形的面积不大于 $\dfrac{1}{2}$ ，则下列说法正确的是\dotfill{(\kh{C})}
    \xx{\ding{172}  \ding{173} 都正确}{\ding{172} 正确， \ding{173} 错误}{\ding{172} 错误， \ding{173} 正确}{\ding{172}  \ding{173} 错误}
    \begin{jieda}
        画图即可.
    \end{jieda}
\end{enumerate}

\section*{第 3 章 函 数}

\begin{center}

\begin{tabular}{|p{4cm}|p{12cm}|}
\hline
考点 & 课标要求与命题趋势 \\
\cline{1-2}
1. 函数的定义域与值域 & \ding{172}  在初中用变量之间的依赖关系描述函数的基础上，用集合语言和对应关系刻画函数， 建立完整的函数概念，体会集合语言和对应关系在刻画函数概念中的作用. 了解构成函数的要素，能求简单函数的定义域与值域 \ding{173}  在实际情境中，会根据不同的需要选择恰当的方法(如图像法、列表法、解析法)表示函数,理解函数图像的作用 \ding{174}  通过具体实例，了解简单的分段函数，并能简单应用 \\
\cline{1-2}
2. 函数单调性的判断及其应用 & 借助函数图像，会用符号语言表达函数的单调性、最大值、最小值，理解它们的作用和实际意义，会求简单函数的最值 \\
\cline{1-2}
3. 函数的奇偶性及其应用 & 结合具体函数，了解函数奇偶性的概念和几何意义. 会运用函数图像理解和研究函数的奇偶性及其应用 \\
\cline{1-2}
4. 函数图像 & 能够利用函数的单调性、奇偶性、对称性、周期性等性质描述、作出一些简单的基本初等函数的图像，会利用函数的图像分析函数的一些性质 \\
\cline{1-2}
5. 函数零点及其应用 & \ding{172}  结合学过的函数图像，了解函数零点与方程解的关系 \ding{173}  结合具体连续函数及其图像的特点，了解函数零点存在定理，探索用二分法求方程近似解的思路并会画程序框图，能借助计算工具用二分法求方程近似解，了解用二分法求方程近似解具有一般性规律 \ding{174}  掌握函数零点的定义，会用零点存在定理判断零点所在区间，会求解零点相关问题，这是高考命题的高频考点 \\
\cline{1-2}
6. 函数模型 & \ding{172}  理解函数模型是描述客观世界中变量关系和规律的重要数学语言和工具. 在实际情境中, 会选择合适的函数类型刻画现实问题的变化规律 \ding{173}  结合现实情境中的具体问题，利用计算工具，比较对数函数、一元一次函数、指数函数增长速度的差异，理解“对数增长”“直线上升”“指数爆炸”等术语的现实含义 \ding{174}  重点提升数学抽象、逻辑推理、数学运算和数学建模素养 \ding{175} 收集、阅读一些现实生活、生产实际或者经济领域中的数学模型，体会人们是如何借助函数刻画实际问题的，感悟数学模型中参数的现实意义 \\
\cline{1-2}
7. 函数综合问题 & \ding{172}  掌握处理函数综合问题的数学思想和方法 \ding{173}  函数的形成与发展:收集、阅读函数的形成与发展的历史资料，撰写小论文，论述函数发展的过程、重要结果、主要人物、关键事件及其对人类文明的贡献 \ding{174}该内容是上海高考的必考内容，一般会以抽象函数作为载体，考查函数的单调性、奇偶性、周期性及对称性, 是上海高考一轮复习的重点内容 \\
\cline{1-2}
\hline
\end{tabular}
\end{center}
\clearpage
\subsection*{考点1 函数的定义域与值域}

\subsubsection*{一、填空题}

1. (2024$\cdot$松江区二模$\cdot$1)函数 $y = \lg \left( {x - 2}\right)$ 的定义域是\blkda[2]{}.

2. (2024$\cdot$宝山区一模$\cdot$1)函数 $f\left( x\right)  = \lg \left( {x - 1}\right)$ 的定义域是\blkda[2]{}.

3. (2024$\cdot$虹口区一模$\cdot$2)函数 $y = \lg \left( {x - 2}\right)  + \dfrac{1}{\sqrt{5 - x}}$ 的定义域为\blkda[2]{}.

4.(2025$\cdot$松江区一模$\cdot$3)函数 $y = \lg \left( {{3x} + 1}\right)  + \sqrt{1 - x}$ 的定义域是\blkda[2]{}.

5.(2025$\cdot$金山区一模$\cdot$3)已知函数 $y = f\left( x\right)$ 的表达式为 $f\left( x\right)  = \left\{  \begin{array}{l} {2}^{x},x \leq  2, \\  {\log }_{2}x,x > 2, \end{array}\right.$ 则 $f\left( 4\right)$ 的值为\blkda[2]{}.

6.(2024$\cdot$金山区二模$\cdot$3)函数 $y = {\log }_{2}\dfrac{2 + x}{1 - x}$ 的定义域是\blkda[2]{}.

7.(2024$\cdot$静安区二模$\cdot$3)函数 $y = \ln \dfrac{1 - x}{2 + x}$ 的定义域为\blkda[2]{}.

8.(2025$\cdot$宝山区二模$\cdot$4)已知函数 $f\left( x\right)  = \left\{  \begin{array}{l} {x}^{2} + 2,x < 1, \\  f\left( {x - 2}\right) ,x \geq  1, \end{array}\right.$ 则 $f\left( 4\right)  =$ \blkda[2]{}.

9.(2023$\cdot$上海秋季卷$\cdot$5)已知函数 $f\left( x\right)  = \left\{  \begin{array}{l} {2}^{x},x > 0, \\  1,x \leq  0, \end{array}\right.$ 则 $f\left( x\right)$ 的值域为\blkda[2]{}.

10. (2024$\cdot$嘉定区二模$\cdot$6)函数 $y = \left| {x - 1}\right|  + \left| {x - 4}\right|$ 的值域为\blkda[2]{}.

11. (2025$\cdot$闵行区二模$\cdot$8)已知函数 $f\left( x\right)  = \left\{  \begin{array}{l} \left( {1 - {2m}}\right) x + {3m},x < 1, \\  {x}^{2},x \geq  1 \end{array}\right.$ 的值域为 $\mathbf{R}$ ，则实数 $m$ 的取值范围是\blkda[2]{}.

12.(2024$\cdot$上海春季卷$\cdot$9)已知 $f\left( x\right)  = {x}^{2},g\left( x\right)  = \left\{  \begin{array}{l} f\left( x\right) ,x \geq  0, \\   - f\left( {-x}\right) ,x < 0, \end{array}\right.$ 则 $g\left( x\right)  \leq  2 - x$ 的解集为\blkda[2]{}.

\subsubsection*{二、解答题}

13.(2025$\cdot$闵行区一模$\cdot$18)已知 $f\left( x\right)  = \left\{  \begin{array}{l} {x}^{2} - {ax},x \geq  0, \\  x + \dfrac{1}{x},x < 0. \end{array}\right.$

(1)若 $a = 1$ ，求函数 $y = f\left( x\right)$ 的值域；

(2)若存在 $\varphi  \in  \left( {0,\dfrac{\pi }{4}}\right)$ ，使得 $f\left( {\sin \varphi }\right)  = f\left( {\cos \varphi }\right)$ ，求实数 $a$ 的取值范围.

\subsection*{考点2 函数单调性的判断及其应用}

\subsubsection*{一、填空题}

1.(2024$\cdot$杨浦区一模$\cdot$4)函数 $y = \left| {x - 3}\right|  + \left| {5 - x}\right|$ 的最小值为\blkda[2]{}.

2.(2024$\cdot$虹口区一模$\cdot$9)已知 $y = f\left( x\right)$ 是定义在(-1,1)上的函数，若 $f\left( x\right)  = {3x} + \sin x$ ，且 $f\left( {1 - {a}^{2}}\right)  + f\left( {1 - a}\right)  < 0$ ，则实数 $a$ 的取值范围为\blkda[2]{}.

3.(2025$\cdot$崇明区二模$\cdot$10)已知 $f\left( x\right)  = \left\{  \begin{array}{l}  - {x}^{2} + {ax},x \leq  1, \\  {ax} - 1,x > 1. \end{array}\right.$ 若函数 $y = f\left( x\right)$ 有两个极值点，则实数 $a$ 的取值范围是\blkda[2]{}.

4.(2024$\cdot$青浦区一模$\cdot$10)已知函数 $y = \left\{  \begin{array}{l} {x}^{2} - {2x} + 2,x \geq  0, \\  x + \dfrac{a}{x} + {3a},x < 0 \end{array}\right.$ 的值域为 $\mathbf{R}$ ,则实数 $a$ 的取值范围为\blkda[2]{}.

5.(2025$\cdot$宝山区一模$\cdot$11)某物流公司为了扩大业务量，计划改造一间高为 $6{m}$ ，底面积为 ${24}{{m}}^{2}$ ，且背面靠墙的长方体形状的仓库. 因仓库的背面靠墙,无须建造费用,设仓库前面墙体的长为 $x{m}  \left( {4 \leq  x \leq  6}\right)$. 现有甲、乙两支工程队参加竞标，甲队的报价方案为:仓库前面新建墙体每平方米 400 元，左右两面新建墙体每平方米 300 元，屋顶和地面以及其他共计 28800 元；乙队给出的整体报价为 $\left( {1 + \dfrac{2}{x}}\right)  \times  {6k} \times  {10}^{4}$ 元 $\left( {k > 0}\right)$. 不考虑其他因素，若乙队要确保竞标成功，则实数 $k$ 的取值范围是\blkda[2]{}.

\subsubsection*{二、选择题}

6.(2025$\cdot$闵行区一模$\cdot$14)下列函数中，在区间 $\left( {0, + \infty }\right)$ 上是严格减函数的为 \dotfill{(\kh{})}

\xx{}{}{}{}$y = {x}^{\dfrac{1}{2}}$ B. $y = \dfrac{1}{{x}^{2} + 1}$

C. $y = {2}^{x}$ D. $y = \lg \left| x\right|$

7.(2025$\cdot$崇明区一模$\cdot$13)下列函数中，在其定义域上既是奇函数又是严格增函数的是 \dotfill{(\kh{})}

\xx{}{}{}{}$y = {x}^{3}$ B. $y = {{e}}^{x}$

C. $y = \lg x$ D. $y = \sin x$

8.(2024$\cdot$黄浦区二模$\cdot$15)设函数 $f\left( x\right)  = \left\{  \begin{array}{l}  - {x}^{2} + {ax} + {20}, - 4 \leq  x \leq  0, \\  a{x}^{2} - {2x} + 3,0 < x \leq  4. \end{array}\right.$ 若 $f\left( x\right)  > 0$ 恒成立，则实数 $a$ 的取值范围是 \dotfill{(\kh{})}

\xx{}{}{}{}$\left( {1, + \infty }\right)$ B. $\left( {0,\dfrac{1}{3}}\right)$

C. $\left( {\dfrac{5}{16},1}\right)$ D. $\left( {\dfrac{1}{3},1}\right)$

\subsubsection*{三、解答题}

9.(2025$\cdot$嘉定区二模$\cdot$21)已知函数 $y = f\left( x\right)$ ，其中 $f\left( x\right)  = {x}^{2} + {ax} + b\left( {a\text{ 、 }b \in  \mathbf{R}}\right)$ ，定义集合 $S\left( f\right)  =  \left\{  {\left( {x,y}\right)  \mid  y = f\left( x\right) ,x \in  \mathbf{R}}\right\}$. 对于点 $P\left( {p,q}\right)$ ,定义集合 $D\left( P\right)  = \{ \left( {x,y}\right)  \mid  \left| {x - p}\right|  \leq  1,\left( {x,y}\right)  \in  S\left( f\right) \}$. 若对任意 $\left( {x,y}\right)  \in  D\left( P\right)$ ，均有 $\left| {y - q}\right|  \leq  1$ ，则称点 $P$ 为平衡点.

(1)当 $a = b = 0$ 时,判断点 $P\left( {0,0}\right)$ 是否为平衡点;

(2)当 $a = 0$ 时,求实数 $b$ 的取值范围，使得点 $P\left( {0,0}\right)$ 是平衡点；

(3)求所有实数 $a$ 和 $b$ ，使得点 $P\left( {0,0}\right)$ 是平衡点.

考点3 > 函数的奇偶性及其应用

\subsubsection*{一、填空题}

1. (2024$\cdot$黄浦区一模$\cdot$2)若函数 $y = \left( {x + 1}\right) \left( {x - a}\right)$ 为偶函数，则实数 $a$ 的值为\blkda[2]{}.

2.(2024$\cdot$上海秋季卷$\cdot$4)已知 $f\left( x\right)  = {x}^{3} + a,x \in  \mathbf{R}$ ，且 $y = f\left( x\right)$ 是奇函数，则 $a =$ \blkda[2]{}.

3. (2025$\cdot$杨浦区一模$\cdot$4)已知函数 $y = {x}^{2} + {ax} + 1$ 是偶函数，则实数 $a$ 的值为\blkda[2]{}.

4.(2025$\cdot$宝山区二模$\cdot$5)设 $f\left( x\right)  = \left( {m - 1}\right) {x}^{2} + {3x} + \left( {2 - n}\right)$ ，若函数 $y = f\left( x\right)$ 是奇函数，则 $m + n =$ \blkda[2]{}.

5.(2025$\cdot$徐汇区一模$\cdot$5)设 $a$ 、 $b \in  \mathbf{R}$ ， $f\left( x\right)  = {x}^{3} + 3\sin x + b$ ，若函数 $y = f\left( x\right)$ 是定义在 $\left\lbrack  {-a,{2a} - 1}\right\rbrack$ 上的奇函数，则 $a + b =$ \blkda[2]{}.

6.(2025$\cdot$宝山区一模$\cdot$5)已知 $a$ 、 $b$ 为实数，且函数 $y = {x}^{2} + {ax} + 1$ ( $x \in  \left\lbrack  {b,4}\right\rbrack$ ) 是偶函数，则 $a - b =$ \blkda[2]{}.

7.(2025$\cdot$闵行区一模$\cdot$7)已知函数 $y = \left\{  \begin{array}{l} {\log }_{2}x,x > 0, \\  f\left( x\right) ,x < 0 \end{array}\right.$ 为奇函数，则 $f\left( {-8}\right)  =$ \blkda[2]{}.

8.(2025$\cdot$静安区一模$\cdot$11)记 $f\left( x\right)  = {x}^{2} + \left( {{a}^{2} + {b}^{2} - 1}\right) x + {a}^{2} + {2ab} - {b}^{2}$ ，若函数 $y = f\left( x\right)$ 是偶函数，则该函数图像与 $y$ 轴交点的纵坐标的最大值为\blkda[2]{}.

9.(2024$\cdot$松江区一模$\cdot$11)若函数 $y = f\left( x\right)$ 是定义在 $\mathbf{R}$ 上的不恒为零的偶函数，且对任意实数 $x$ 都有 $x \cdot  f\left( {x + 2}\right)  - \left( {x + 2}\right)  \cdot  f\left( x\right)  + 2$ ，则 $f\left( {2023}\right)  =$ \blkda[2]{}.

\subsubsection*{二、选择题}

10.(2023$\cdot$上海春季卷$\cdot$13)下列函数是偶函数的是 \dotfill{(\kh{})}

\xx{}{}{}{}$y = \sin x$ B. $y = \cos x$ C. $y = {x}^{3}$ D. $y = {2}^{x}$

11.(2024$\cdot$长宁区一模$\cdot$13)下列函数中既是奇函数又是增函数的是 \dotfill{(\kh{})}

\xx{}{}{}{}$f\left( x\right)  = {2x}$ B. $f\left( x\right)  = {x}^{2}$ C. $f\left( x\right)  = \ln x$ D. $f\left( x\right)  = {{e}}^{x}$

12.(2025$\cdot$虹口区二模$\cdot$14)下列函数中为奇函数的是 \dotfill{(\kh{})}

\xx{}{}{}{}$y = \sqrt{x}$ B. $y = {x}^{-2}$ C. $y = \sin \left( {x + \dfrac{\pi }{2}}\right)$ D. $y = \tan \left( {x + \pi }\right)$

13.(2024$\cdot$闵行区二模$\cdot$14)已知 $y = f\left( x\right) \left( {x \in  \mathbf{R}}\right)$ 为奇函数，当 $x > 0$ 时， $f\left( x\right)  = {\log }_{2}x - 1$ ，则集合 $\{ x \mid  f\left( {-x}\right)  - f\left( x\right)  < 0\}$ 可表示为 \dotfill{(\kh{})}

\xx{}{}{}{}$\left( {2, + \infty }\right)$ B. $\left( {-\infty , - 2}\right)$

C. $\left( {-\infty , - 2}\right)  \cup  \left( {2, + \infty }\right)$ D. $\left( {-2,0}\right)  \cup  \left( {2, + \infty }\right)$

14.(2024$\cdot$上海秋季卷$\cdot$16)已知函数 $y = f\left( x\right)$ 的定义域为 $\mathbf{R}$ ，定义集合 $M = \left\{  {{x}_{0} \mid  {x}_{0} \in  \mathbf{R},x \in  \left( {-\infty }\right. }\right.$ ， $\left. {{x}_{0}),f\left( x\right)  < f\left( {x}_{0}\right) }\right\}$ ,在使得 $M = \left\lbrack  {-1,1}\right\rbrack$ 的所有 $f\left( x\right)$ 中,下列成立的是 \dotfill{(\kh{})}

\xx{}{}{}{}存在 $f\left( x\right)$ 是偶函数 B. 存在 $f\left( x\right)$ 在 $x = 2$ 处取最大值

C. 存在 $f\left( x\right)$ 是严格增函数 D. 存在 $f\left( x\right)$ 在 $x =  - 1$ 处取到极小值

15.(2024$\cdot$崇明区二模$\cdot$16)已知函数 $y = f\left( x\right)$ 的定义域为 $D,{x}_{1}\text{ 、 }{x}_{2} \in  D$.

命题 $p$ : 若当 $f\left( {x}_{1}\right)  + f\left( {x}_{2}\right)  = 0$ 时,都有 ${x}_{1} + {x}_{2} = 0$ ,则函数 $y = f\left( x\right)$ 是 $D$ 上的奇函数.

命题 $q$ : 若当 $f\left( {x}_{1}\right)  < f\left( {x}_{2}\right)$ 时,都有 ${x}_{1} < {x}_{2}$ ,则函数 $y = f\left( x\right)$ 是 $D$ 上的增函数. 下列说法正确的是 \dotfill{(\kh{})}

\xx{}{}{}{}$p\text{ 、 }q$ 都是真命题 B. $p$ 是真命题, $q$ 是假命题

C. $p$ 是假命题， $q$ 是真命题 D. $p$ 、 $q$ 都是假命题

\subsubsection*{三、解答题}

16.(2025$\cdot$浦东新区二模$\cdot$17)已知函数 $y = f\left( x\right)$ 的表达式 $f\left( x\right)  = \dfrac{1}{{{e}}^{x} + 1} - a$.

(1)若函数 $y = f\left( x\right)$ 是奇函数，求实数 $a$ 的值；

(2)对任意实数 $x \in  \left\lbrack  {-1,1}\right\rbrack$ ，不等式 $f\left( x\right)  \leq  0$ 恒成立，求实数 $a$ 的取值范围.

17.(2023$\cdot$上海秋季卷$\cdot$18)已知 $f\left( x\right)  = \dfrac{{x}^{2} + \left( {{3a} + 1}\right) x + c}{x + a},a\text{ 、 }c \in  \mathbf{R}$.

(1)当 $a = 0$ 时,是否存在实数 $c$ ，使得函数 $y = f\left( x\right)$ 为奇函数；

(2)若函数 $y = f\left( x\right)$ 的图像过点 $\left( {1,3}\right)$ ，且与 $x$ 轴负半轴有两个不同交点，求实数 $c$ 的值及实数 a 的取值范围.

考点4 函数图像

\subsubsection*{一、填空题}

1. (2024$\cdot$嘉定区一模$\cdot$9)关于 $x$ 的方程 $\left| {{x}^{2} - {3x} + 2}\right|  = {mx}$ 有三个不同的实数解，则实数 $m$ 的值为\blkda[2]{}.

2. (2025$\cdot$虹口区一模$\cdot$8)已知 $f\left( x\right)  = \left\{  \begin{array}{l} {x}^{2} - x,x \geq  0, \\  f\left( {-x}\right) ,x < 0, \end{array}\right.$ 则 $f\left( x\right)  \leq  6$ 的解集是\blkda[2]{}.

3. (2025$\cdot$闵行区二模$\cdot$10)已知函数 $y = g\left( x\right)$ 的图像是折线段 ${ABC}$ ，且 $A\left( {0,0}\right)$ 、 $B\left( {\dfrac{1}{2},\dfrac{1}{2}}\right)$ 、 $C\left( {1,0}\right)$ ，则函数 $y = x \cdot  g\left( x\right) \left( {0 \leq  x \leq  1}\right)$ 的图像与 $x$ 轴围成的图形面积为\blkda[2]{}.

4.(2025$\cdot$普陀区一模$\cdot$11)设 $t \in  \mathbf{R}$ ，直线 $l : x + y - t = 0$ 与曲线 ${C}_{1} : y = \dfrac{1}{4}{x}^{2}\left( {0 \leq  x \leq  4}\right)$ 和曲线 ${C}_{2} : y = 2{x}^{\dfrac{1}{2}}$ 分别交于 $P\text{ 、 }Q$ 两点,则 $\left| {PQ}\right|$ 的最大值是\blkda[2]{}.

5.(2025$\cdot$普陀区一模$\cdot$12)设 $a > b > 0$ ，函数 $y = f\left( x\right)$ 的表达式为 $f\left( x\right)  = \left| {x - \dfrac{1}{x} + \ln x}\right|$ ，若 $f\left( a\right)  =  f\left( b\right)$ ，且关于 $x$ 的方程 $\left| {{x}^{2} + {ax} + {2ab}}\right|  + \left| {{x}^{2} - {ax} + {2ab}}\right|  = {2a}\left| x\right|$ 的整数解有且仅有 4 个，则 $a$ 的取值范围是\blkda[2]{}.

\subsubsection*{二、选择题}

\begin{center}
\includegraphics[width=0.2\textwidth]{bodyxb/src/22_1125_1827_332_191_0.jpg}
\end{center}


第 6 题图

6.(2024$\cdot$奉贤区二模$\cdot$14)已知函数 $y = f\left( x\right)$ ，其中 $f\left( x\right)  = {x}^{2} + 1,y =  g\left( x\right)$ ，其中 $g\left( x\right)  = 4\sin x$ ，则图像如图所示的函数可能是\blkda[2]{}\dotfill{(\kh{})}

\xx{}{}{}{}$y = \dfrac{g\left( x\right) }{f\left( x\right) }$ B. $y = \dfrac{f\left( x\right) }{g\left( x\right) }$

C. $y = f\left( x\right)  + g\left( x\right)  - 1$ D. $y = f\left( x\right)  - g\left( x\right)  - 1$

\subsection*{考点5 函数零点及其应用}

\subsubsection*{一、填空题}

1.(2024$\cdot$徐汇区一模$\cdot$6)函数 $y = \lg \left( {{2x} + 1}\right)  + \lg x$ 的零点是\blkda[2]{}.

2.(2024$\cdot$青浦区一模$\cdot$8)若函数 $y = \cos \left( {x + \varphi }\right)$ 是奇函数，则该函数的所有零点是\blkda[2]{}.

3. (2024$\cdot$青浦区二模$\cdot$9)对于函数 $y = f\left( x\right)$ ，其中 $f\left( x\right)  = \left\{  \begin{array}{l} {\left( x - 1\right) }^{3},0 \leq  x < 2, \\  \dfrac{2}{x},x \geq  2, \end{array}\right.$ 若关于 $x$ 的方程 $f\left( x\right)  = {kx}$ 有两个不同的根，则实数 $k$ 的取值范围是\blkda[2]{}.

4. (2025$\cdot$杨浦区一模$\cdot$10)已知 $f\left( x\right)  = \left\{  \begin{array}{l} {x}^{\dfrac{1}{3}},0 \leq  x \leq  a, \\  {\log }_{3}x,x > a, \end{array}\right.$ 其中实数 $a > 0$. 若函数 $y = f\left( x\right)  - 2$ 有且仅有 2 个零点,则 $a$ 的取值范围为\blkda[2]{}.

5.(2024$\cdot$金山区一模$\cdot$11)若函数 $y = \left| {\left( {1 - {x}^{2}}\right) \left( {{x}^{2} + {ax} + b}\right) }\right|  - c$ (c ≠ 0)的图像关于直线 $x =  - 2$ 对称，且该函数有且仅有 7 个零点，则 $a + b + c$ 的值为\blkda[2]{}.

6.(2025$\cdot$崇明区一模$\cdot$11)已知 $f\left( x\right)  = A\sin \left( {{\omega x} + \dfrac{\pi }{6}}\right) \left( {A > 0,\omega  > 0}\right)$ ，若函数 $y = f\left( x\right)$ 在区间 $\left\lbrack  {0,{2\pi }}\right\rbrack$ 上有且仅有 3 个零点和 1 个极小值点，则 $\omega$ 的取值范围是\blkda[2]{}.

7.(2025$\cdot$普陀区二模$\cdot$12)设 $a \in  \mathbf{R}$ ，函数 $y = f\left( x\right)$ 的表达式为 $f\left( x\right)  = \left\{  \begin{array}{l} {{e}}^{x},x > 0, \\  \dfrac{1}{x + a},x < 0. \end{array}\right.$ 若函数 $g\left( x\right)  =  f\left( x\right) f\left( {x + 1}\right)  - 1$ 恰有三个零点，则 $a$ 的取值范围是\blkda[2]{}.

\subsubsection*{二、选择题}

8.(2025$\cdot$宝山区一模$\cdot$14)下列函数中，在区间 $\left( {0, + \infty }\right)$ 上是严格增函数且存在零点的是\blkda[2]{} \dotfill{(\kh{})}

\xx{}{}{}{}$y = {{e}}^{x}$ B. $y = \sqrt{x} + 2$ C. $y =  - {\log }_{\dfrac{1}{2}}x$ D. $y = {\left( x - 2\right) }^{2}$

9.(2025$\cdot$青浦区一模$\cdot$15)已知函数 $y = f\left( x\right)$ 是定义在 $\mathbf{R}$ 上的奇函数，且当 $x > 0$ 时， $f\left( x\right)  =  \left( {x - 1}\right) \left( {x - 3}\right)  + {0.01}$ ,则关于函数 $y = f\left( x\right)$ 在 $\mathbf{R}$ 上的零点的说法正确的是 \dotfill{(\kh{})}

\xx{}{}{}{}有 4 个零点,其中只有一个零点在区间 $\left( {-3, - 1}\right)$ 上

B. 有 4 个零点,其中两个零点在区间 $\left( {-3, - 1}\right)$ 上,另外两个零点在区间 $\left( {1,3}\right)$ 上

C. 有 5 个零点,两个正零点中一个在区间 $\left( {0,1}\right)$ 上,一个在区间 $\left( {3, + \infty }\right)$ 上

D. 有 5 个零点,都不在 $\left( {0,1}\right)$ 上

10.(2024$\cdot$松江区二模$\cdot$15)已知某个三角形的三边长为 $a$ 、 $b$ 、 $c$ ，其中 $a < b$. 若 $a$ 、 $b$ 是函数 $y = a{x}^{2} -  {bx} + c$ 的两个零点,则 $a$ 的取值范围是 \dotfill{(\kh{})}

\xx{}{}{}{}$\left( {\dfrac{1}{2},1}\right)$ B. $\left( {\dfrac{1}{2},\dfrac{\sqrt{5} - 1}{2}}\right)$ C. $\left( {0,\dfrac{\sqrt{5} - 1}{2}}\right)$ D. $\left( {\dfrac{\sqrt{5} - 1}{2},1}\right)$

\subsubsection*{三、解答题}

11.(2024$\cdot$上海春季卷$\cdot$17)已知 $f\left( x\right)  = \sin \left( {{\omega x} + \dfrac{\pi }{3}}\right)$ ， $\omega  > 0$.

(1)若 $\omega  = 1$ ，当 $x \in  \left\lbrack  {0,\pi }\right\rbrack$ 时，求 $y = f\left( x\right)$ 的值域；

(2)已知 $a > \pi \left( {a \in  \mathbf{R}}\right)$ ，且 $y = f\left( x\right)$ 的最小正周期为 $\pi$ ，若 $y = f\left( x\right)$ 在 $x \in  \left\lbrack  {\pi ,a}\right\rbrack$ 上有三个零点， 求 $a$ 的取值范围.

12.(2025$\cdot$上海秋季卷$\cdot$21)已知函数 $y = f\left( x\right)$ 的定义域为 R. 对于正实数 $a$ ，定义集合 ${M}_{a} = \{ x \mid  f\left( {x + a}\right)  = f\left( x\right) \}$.

(1)若 $f\left( x\right)  = \sin x$ ，判断 $\dfrac{\pi }{3}$ 是否是 ${M}_{\pi }$ 中的元素，请说明理由；

(2)若 $f\left( x\right)  = \left\{  {\begin{array}{l} x + 2,x < 0, \\  \sqrt{x},x \geq  0, \end{array}{M}_{a} \neq  \varnothing }\right.$ ，求 $a$ 的取值范围；

(3)若 $y = f\left( x\right)$ 是偶函数，当 $x \in  (0,1\rbrack$ 时， $f\left( x\right)  = 1 - x$ ，且对任意 $a \in  \left( {0,2}\right)$ ，均有 ${M}_{a} \subseteq  {M}_{2}$. 写出 $y = f\left( x\right) ,x \in  \left( {1,2}\right)$ 的解析式,并证明: 对任意实数 $c$ ,函数 $y = f\left( x\right)  - c$ 在 $\left\lbrack  {-3,3}\right\rbrack$ 上至多有 9 个零点.

\subsection*{考点6 二 函数模型}

\subsubsection*{一、填空题}

1.(2025$\cdot$虹口区一模$\cdot$11)2024 年 10 月 30 日，“神舟十九号”载人飞船发射成功，标志着中国空间站建设进入新阶段. 在飞船竖直升空过程中，某位记者用照相机在同一位置以同一姿势连续拍照两次. 已知 “神舟十九号” 飞船船体实际长度为 $H$ ,且在照片上飞船船体长度为 $h$ ,比较两张照片,相对于照片中的同一固定参照物飞船上升了 $m$. 假设该记者连按拍照键间的反应时间为 $t$ ,并忽略相机曝光时长，若用平均速度估算瞬时速度，则拍照时飞船的瞬时速度为\blkda[2]{}. (用含有 $H\text{ 、 }h\text{ 、 }m\text{ 、 }t$ 的式子表示)

2.(2025$\cdot$松江区一模$\cdot$12)交通信号灯由红灯、绿灯、黄灯组成. 黄灯设置的时长与路口宽度、限定速度、停车距离有关. 根据路况不同，道路的限定速度一般在 ${30}{{km}}/{h}$ 至 ${70}{{km}}/{h}$ 之间. 由相关数据， 驾驶员反应距离 ${s}_{1}$ (单位: ${m}$ )关于车速 $v$ (单位: ${m}/{s}$ )的函数模型为: ${s}_{1} = {0.7584v}$ ; 刹车距离 ${s}_{2}$ (单位: ${m}$ ) 关于车速 $v$ (单位: ${m}/{s}$ ) 的函数模型为: ${s}_{2} = {0.072}{v}^{2}$ ,反应距离与刹车距离之和称为停车距离. 已知某个十字路口宽度为 ${30}{m}$ ,为保证通行安全,黄灯亮的时间是允许限速车辆离停车线距离小于停车距离的汽车通过十字路口，则该路口黄灯亮的时间最多为\blkda[2]{}s.(结果精确到 ${0.01}{s}$ )

\subsubsection*{二、选择题}

3.(2025$\cdot$松江区一模$\cdot$16)设函数 $y = f\left( x\right)$ 与 $y = g\left( x\right)$ 均是定义在 $\mathbf{R}$ 上的函数，有以下两个命题:  \ding{172}  若 $y = f\left( x\right)$ 是周期函数，且是 $\mathbf{R}$ 上的减函数，则函数 $y = f\left( x\right)$ 必为常值函数；

 \ding{173}  若对任意的 $a$ 、 $b \in  \mathbf{R}$ ，有 $\left| {f\left( a\right)  - f\left( b\right) }\right|  \leq  \left| {g\left( a\right)  - g\left( b\right) }\right|$ 成立，且 $y = g\left( x\right)$ 是 $\mathbf{R}$ 上的增函数，则 $y = f\left( x\right)  - g\left( x\right)$ 是 $\mathbf{R}$ 上的增函数. 则以下选项正确的是 \dotfill{(\kh{})}

\xx{}{}{}{} \ding{172} 是真命题， \ding{173} 是假命题

B. 两个都是真命题

C.  \ding{172} 是假命题， \ding{173} 是真命题

D. 两个都是假命题

4.(2025$\cdot$浦东新区一模$\cdot$16)设函数 $y = F\left( x\right)$ 、 $y = G\left( x\right)$ 的定义域均为 $\mathbf{R}$ ，值域分别为 $A$ 、 $B$ ，且 $A \cap  B = \varnothing$. 若集合 $S$ 满足以下两个条件:(1) $A \cup  B \subseteq  S$ ；(2) ${\complement }_{S}\left( {A \cup  B}\right)$ 是有限集，则称 $y = F\left( x\right)$ 和 $y = G\left( x\right)$ 是 $S$ -互补函数. 给出以下两个命题: \ding{172}  存在函数 $y = f\left( x\right)$ ，使得 $y = {2}^{f\left( x\right) }$ 和 $y = {\log }_{2}f\left( x\right)$ 是 $\left\lbrack  {0,{16}}\right\rbrack$ 一互补函数;  \ding{173}  存在函数 $y = g\left( x\right)$ ,使得 $y = \sin g\left( x\right)$ 和 $y = {tang}\left( x\right)$ 是 $\lbrack 0, + \infty )$ 一互补函数. 则 \dotfill{(\kh{})}

\xx{}{}{}{} \ding{172}  \ding{173} 都是真命题

B.  \ding{172} 是真命题， \ding{173} 是假命题

C.  \ding{172} 是假命题， \ding{173} 是真命题

D.  \ding{172}  \ding{173} 都是假命题

\subsubsection*{三、解答题}

5.(2023$\cdot$上海春季卷$\cdot$19)为了节能环保、节约材料，定义建筑物的“体形系数”为S= $\dfrac{{F}_{0}}{{V}_{0}}$ ，其中 ${F}_{0}$ 为建筑物暴露在空气中的面积 (单位: ${{m}}^{2}$ ), ${V}_{0}$ 为建筑物的体积 (单位: ${{m}}^{3}$ ).

(1)若有一圆柱形建筑物的底面半径为 $R$ ,高度为 $H$ ,求该建筑物的“体形系数” (结果用含 $R$ 、 $H$ 的代数式表示)；

(2)定义建筑物的“形状因子”为 $f = \dfrac{{L}^{2}}{A}$ ，其中 $A$ 为底面面积， $L$ 为建筑底面周长，又定义 $T$ 为总建筑面积,即为每层建筑面积总和 (每层建筑面积为每一层的底面面积). 现有一垂直于底面的宿舍,已知该宿舍的层高为 $3{m}$ ,有 $n$ 层,其“形状因子” $f = {18}$ ,总建筑面积为 $T = {10000}{{m}}^{2}$ , 试求当该宿舍的层数 $n$ 为多少时,“体形系数” $S$ 最小.

6.(2024$\cdot$松江区一模$\cdot$19)为了鼓励居民节约用气，某市对燃气收费实行阶梯计价，普通居民燃气收费标准如下:

第一档:年用气量在 $0 \sim  {310}$ (含) ${{m}}^{3}$ ，价格为 $a$ 元/ ${{m}}^{3}$ ；

第二档:年用气量在 310\textasciitilde 520(含) ${{m}}^{3}$ ，价格为 $b$ 元/ ${{m}}^{3}$ ；

第三档:年用气量在 ${520}{{m}}^{3}$ 以上，价格为 $c$ 元 $/{{m}}^{3}$.

(1)请写出普通居民的年度燃气费用(单位:元)关于年度的燃气用量(单位: ${{m}}^{3}$ )的函数解析式 (用含 $a$ 、 $b$ 、 $c$ 的式子表示)；

(2)已知某户居民 2023 年部分月份用气量与缴费情况如下表,求 $a$ 、 $b$ 、 $c$ 的值.

\begin{center}

\begin{tabular}{|c|c|c|c|c|c|c|c|c|}
\hline
月份 & 1 & 2 & 3 & 4 & 5 & 9 & 10 & 12 \\
\cline{1-9}
当月燃气用量 $\left( {{m}}^{3}\right)$ & 56 & 80 & 66 & 58 & 60 & 53 & 55 & 63 \\
\cline{1-9}
当月燃气费(元) & 168 & 240 & 198 & 174 & 183 & 174.9 & 186 & 264.6 \\
\cline{1-9}
\hline
\end{tabular}
\end{center}

7.(2025$\cdot$静安区一模$\cdot$21)如果函数 $y = f\left( x\right)$ 满足以下两个条件，我们就称函数 $y = f\left( x\right)$ 为 $U$ 型函数.

 \ding{172}  对任意的 $x \in  \left\lbrack  {0,1}\right\rbrack$ ，有 $f\left( x\right)  \geq  1$ ， $f\left( 1\right)  = 3$ ；

 \ding{173}  对于任意的 $x$ 、 $y \in  \left\lbrack  {0,1}\right\rbrack$ ，若 $x + y \leq  1$ ，则 $f\left( {x + y}\right)  \geq  f\left( x\right)  + f\left( y\right)  - 1$.

求证:(1) $y = {3}^{x}$ 是 $U$ 型函数；

(2) $U$ 型函数 $y = f\left( x\right)$ 在 $\left\lbrack  {0,1}\right\rbrack$ 上为增函数；

(3)对于 $U$ 型函数 $y = f\left( x\right)$ ,有 $f\left( \dfrac{1}{{3}^{n}}\right)  \leq  \dfrac{2}{{3}^{n}} + 1$ ( $n$ 为正整数).

8.(2025$\cdot$闵行区一模$\cdot$21)设函数 $y = f\left( x\right)$ 的定义域为 $\mathbf{R}$ ，集合 $M = \{ x \mid  f\left( x\right)  = a,x \in  \mathbf{R}\}$. 若 $M$ 中有且仅有一个元素,则称 $a$ 为函数 $y = f\left( x\right)$ 的一个“S值”.

(1)设 $f\left( x\right)  = {x}^{2} - {2x}$ ，求 $y = f\left( x\right)$ 的“S值”；

(2)设 $g\left( x\right)  = 3{x}^{4} - \left( {{4k} + 4}\right) {x}^{3} + {6k}{x}^{2} + 1$ ，且 $0 < k \leq  1$ ，若 $y = g\left( x\right)$ 的函数值中不存在 “S值”，求实数 $k$ 取值的集合;

(3)已知定义域为 $\mathbf{R}$ 的函数 $y = h\left( x\right)$ 的图像是一条连续曲线，且函数 $y = h\left( x\right)$ 的所有函数值均为“$S$ 值”,若 $m < n$ ,证明: $y = h\left( x\right)$ 在 $\left\lbrack  {m,n}\right\rbrack$ 上为严格增函数的一个充要条件是 $h\left( m\right)  < h\left( n\right)$.

9.(2024$\cdot$长宁区二模$\cdot$21)设函数 $y = f\left( x\right)$ 的定义域为 $D$ ，若存在实数 $k$ ，使得对于任意 $x \in  D$ ，都有 $f\left( x\right)  \leq  k$ ,则称函数 $y = f\left( x\right)$ 有上界,实数 $k$ 的最小值为函数 $y = f\left( x\right)$ 的上确界; 记集合 ${M}_{n} =  \left\{  {f\left( x\right) \left| {y = \dfrac{f\left( x\right) }{{x}^{n}}\text{ 在区间 }\left( {0, + \infty }\right) }\right. \text{ 上是严格增函数 }}\right\}.$

(1)求函数 $y = \dfrac{2}{x - 1}\left( {2 < x < 6}\right)$ 的上确界；

(2)若 $f\left( x\right)  = {x}^{3} - h{x}^{2} + {2x}\ln x \in  {M}_{1}$ ，求 $h$ 的最大值；

(3)设函数 $y = f\left( x\right)$ 一定义域为 $\left( {0, + \infty }\right)$. 若 $f\left( x\right)  \in  {M}_{2}$ ，且 $y = f\left( x\right)$ 有上界，求证: $f\left( x\right)  < 0$ ,且存在函数 $y = f\left( x\right)$ ,它的上确界为 0.

10.(2024$\cdot$虹口区二模$\cdot$21)若函数 $y = f\left( x\right)$ 满足:对任意 ${x}_{1}$ 、 ${x}_{2} \in  \mathbf{R}$ ， ${x}_{1} + {x}_{2} \neq  0$ ，都有 $\dfrac{f\left( {x}_{1}\right)  + f\left( {x}_{2}\right) }{{x}_{1} + {x}_{2}} > 0$ ， 则称函数 $y = f\left( x\right)$ 具有性质 $P$.

(1)设 $f\left( x\right)  = {{e}}^{x},g\left( x\right)  = {x}^{3} + x$ ，分别判断 $y = f\left( x\right)$ 与 $y = g\left( x\right)$ 是否具有性质 $P$ ，并说明理由；

(2)设 $f\left( x\right)  = x + a\sin {2x}$ ，若函数 $y = f\left( x\right)$ 具有性质 $P$ ，求实数 $a$ 的取值范围；

(3)已知函数 $y = f\left( x\right)$ 具有性质 $P$ ，且图像是一条连续曲线，若 $y = f\left( x\right)$ 在 $\mathbf{R}$ 上是严格增函数， 求证: $y = f\left( x\right)$ 是奇函数.

11. (2024$\cdot$杨浦区二模$\cdot$21)函数 $y = f\left( x\right)$ 、 $y = g\left( x\right)$ 的定义域均为 $\mathbf{R}$ ，若对任意两个不同的实数。 $b$ ,均有 $f\left( a\right)  + g\left( b\right)  > 0$ 或 $f\left( b\right)  + g\left( a\right)  > 0$ 成立,则称 $y = f\left( x\right)$ 与 $y = g\left( x\right)$ 为相关函数对.

(1)判断函数 $f\left( x\right)  = x + 1$ 与 $g\left( x\right)  =  - x + 1$ 是否为相关函数对，并说明理由;

(2)已知 $f\left( x\right)  = {{e}}^{x}$ 与 $g\left( x\right)  =  - x + k$ 为相关函数对，求实数 $k$ 的取值范围；

(3)已知函数 $y = f\left( x\right)$ 与 $y = g\left( x\right)$ 为相关函数对，且存在正实数 $M$ ，对任意实数 $x \in  \mathbf{R}$ 均有 $\left| {f\left( x\right) }\right|  \leq  M$. 求证: 存在实数 $m\text{ 、 }n\left( {m < n}\right)$ ,使得对任意 $x \in  \left( {m,n}\right)$ ,均有 $f\left( x\right)  + g\left( x\right)  \geq  - \dfrac{1}{2024}$

考点7 又 函数综合问题

\subsubsection*{一、填空题}

1.(2024$\cdot$松江区一模$\cdot$10)已知函数 $f\left( x\right)  =  - {x}^{2} + {6x} + m,g\left( x\right)  = 2\sin \left( {{2x} + \dfrac{\pi }{3}}\right)$ ,对任意 ${x}_{0} \in  \left\lbrack  {0,\dfrac{\pi }{4}}\right\rbrack$ ,存在 ${x}_{1}\text{ 、 }{x}_{2} \in  \left\lbrack  {-1,3}\right\rbrack$ ,使得 $f\left( {x}_{1}\right)  \leq  g\left( {x}_{0}\right)  \leq  f\left( {x}_{2}\right)$ ,则实数 $m$ 的取值范围是\blkda[2]{}.

2.(2024$\cdot$长宁区一模$\cdot$12)设 $f\left( x\right)  = \left| {{\log }_{2}x + {ax} + b}\right| \left( {a > 0}\right)$ ，记函数 $y = f\left( x\right)$ 在区间 $\left\lbrack  {t,t + 1}\right\rbrack  \left( {t > 0}\right)$ 上的最大值为 ${M}_{t}\left( {a,b}\right)$ ，若对任意 $b \in  \mathbf{R}$ ，都有 ${M}_{t}\left( {a,b}\right)  \geq  a + 1$ ，则实数 $t$ 的最大值为\blkda[2]{}.

3. (2024$\cdot$徐汇区一模$\cdot$12)已知函数 $y = f\left( x\right)$ ，其中 $f\left( x\right)  = \left| {\dfrac{{2}^{x + 1}}{{2}^{x} + {2}^{-x}} - 1 - a}\right|$ ，存在实数 ${x}_{1}$ 、 ${x}_{2}$ 、 $\cdots$ 、 ${x}_{n}$ ，使得 $\mathop{\sum }\limits_{{i = 1}}^{{n - 1}}f\left( {x}_{i}\right)  = f\left( {x}_{n}\right)$ 成立，若正整数 $n$ 的最大值为 8，则实数 $a$ 的取值范围是\blkda[2]{}.

\subsubsection*{二、选择题}

4.(2025$\cdot$闵行区二模$\cdot$15)已知定义在 $\mathbf{R}$ 上的函数 $y = f\left( x\right)$ ，满足以下两个条件:(1) $f\left( x\right)  > 0$ 对任意 $x \in  \mathbf{R}$ 恒成立，且 $f\left( 1\right)  = 1$ ；(2)对任意 ${x}_{1}$ 、 ${x}_{2} \in  \mathbf{R}$ 都有 $f\left( {{x}_{1} + 2{x}_{2}}\right)  = {4f}\left( {x}_{1}\right)  \cdot  {f}^{2}\left( {x}_{2}\right)$ ，则下列关于函数 $y = f\left( x\right)$ 的表述中正确的个数为 \dotfill{(\kh{})}

 \ding{172}  $f\left( 0\right)  = \dfrac{1}{2}$ ;  \ding{173}  $f\left( x\right)  \cdot  f\left( {-x}\right)  = \dfrac{1}{4}$ ;  \ding{174}  函数 $y = f\left( x\right)$ 有最小值.

\xx{}{}{}{}0 B. 1 C. 2 D. 3

5.(2025$\cdot$普陀区一模$\cdot$16)在平面直角坐标系中，将函数 $y = f\left( x\right)$ 的图像绕坐标原点 $O$ 逆时针旋转 $\dfrac{\pi }{4}$ 后,所得曲线仍然是某个函数的图像,则称函数 $y = f\left( x\right)$ 为“ $R$ 函数”. 对于命题:

 \ding{172}  设 $m \in  \mathbf{R}$ ，若函数 $g\left( x\right)  = \left( {m - 1}\right) x + \dfrac{1}{x}$ 为“ $R$ 函数”，则 $m > 1$ ；

 \ding{173}  设 $k \in  \mathbf{R}$ ，若函数 $h\left( x\right)  = \dfrac{k\left( {x + 1}\right) }{{{e}}^{x}}$ 为“ $R$ 函数”，则满足条件的 $k$ 的整数值至少有 4 个. 则下列结论中正确的是 \dotfill{(\kh{})}

\xx{}{}{}{} \ding{172} 为真， \ding{173} 为真 B.  \ding{172} 为真， \ding{173} 为假 C.  \ding{172} 为假， \ding{173} 为真 D.  \ding{172} 为假， \ding{173} 为假

6.(2024$\cdot$杨浦区一模$\cdot$16)函数 $y = f\left( x\right)$ 满足:对于任意 $x \in  \mathbf{R}$ 都有 $f\left( x\right)  = f\left( {a}^{x}\right)$ ，(常数 $a > 1$ ， $a \neq  1)$. 给出以下两个命题: \ding{172}  无论 $a$ 取何值，函数 $y = f\left( x\right)$ 不是 $\left( {0, + \infty }\right)$ 上的严格增函数； \ding{173}  当 $0 < a < 1$ 时，存在无穷多个开区间 ${I}_{1},{I}_{2},\cdots ,{I}_{n}$ ，使得 ${I}_{1} \supset  {I}_{2} \supset  \cdots  \supset  {I}_{n} \supset  \cdots$ ，且集合 $\{ y \mid  y = f\left( x\right)$ ， $\left. {x \in  {I}_{n}}\right\}   = \left\{  {y \mid  y = f\left( x\right) ,x \in  {I}_{n + 1}}\right\}$ 对任意正整数 $n$ 都成立，则 \dotfill{(\kh{})}

\xx{}{}{}{} \ding{172}  \ding{173} 都正确 B.  \ding{172} 正确， \ding{173} 不正确 C.  \ding{172} 不正确， \ding{173} 正确 D.  \ding{172}  \ding{173} 都不正确

\subsubsection*{三、解答题}

7.(2024$\cdot$普陀区二模$\cdot$21)对于函数 $y = f\left( x\right) ,x \in  {D}_{1}$ 和 $y = g\left( x\right) ,x \in  {D}_{2}$ ，设 ${D}_{1} \cap  {D}_{2} = D$ ，若 ${x}_{1}$ 、 ${x}_{2} \in  D$ ,且 ${x}_{1} \neq  {x}_{2}$ ,皆有 $\left| {f\left( {x}_{1}\right)  - f\left( {x}_{2}\right) }\right|  \leq  t\left| {g\left( {x}_{1}\right)  - g\left( {x}_{2}\right) }\right| \left( {t > 0}\right)$ 成立，则称函数 $y = f\left( x\right)$ 与 $y = g{\left( x\right) }^{u}$ 具有性质 $H{\left( t\right) }^{n}$.

(1)判断函数 $f\left( x\right)  = {x}^{2},x \in  \left\lbrack  {1,2}\right\rbrack$ 与 $g\left( x\right)  = {2x}$ 是否“具有性质 $H\left( 2\right)$ ”，并说明理由;

(2)若函数 $f\left( x\right)  = 2 + {x}^{2},x \in  (0,1\rbrack$ 与 $g\left( x\right)  = \dfrac{1}{x}$ “具有性质 $H\left( t\right)$ ”,求 $t$ 的取值范围；

(3)若函数 $f\left( x\right)  = \dfrac{1}{{x}^{2}} + {2\ln }x - 3$ 与 $y = g\left( x\right)$ “具有性质 $H\left( 1\right)$ ”，且函数 $y = g\left( x\right)$ 在区间 $\left( {0, + \infty }\right)$ 上存在两个零点 ${x}_{1}\text{ 、 }{x}_{2}$ ,求证 ${x}_{1}^{2} + {x}_{2}^{2} > 2$.

8.(2024$\cdot$静安区一模$\cdot$21)如果函数 $y = f\left( x\right)$ 满足以下两个条件，我们就称 $y = f\left( x\right)$ 为 $L$ 型函数.

 \ding{172}  对任意的 $x \in  \left( {0,1}\right)$ ，总有 $f\left( x\right)  > 0$ ；

 \ding{173}  当 ${x}_{1} > 0,{x}_{2} > 0,{x}_{1} + {x}_{2} < 1$ 时,总有 $f\left( {{x}_{1} + {x}_{2}}\right)  < f\left( {x}_{1}\right)  + f\left( {x}_{2}\right)$ 成立.

(1)记 $g\left( x\right)  = {x}^{2} + \dfrac{1}{2}$ ，求证: $y = g\left( x\right)$ 为 $L$ 型函数；

(2)设 $b \in  \mathbf{R}$ ，记 $p\left( x\right)  = \ln \left( {x + b}\right)$ ，若 $y = p\left( x\right)$ 是 $L$ 型函数，求 $b$ 的取值范围；

(3)是否存在 $L$ 型函数 $y = r\left( x\right)$ 满足:对于任意的 $m \in  \left( {0,4}\right)$ ，都存在 ${x}_{0} \in  \left( {0,1}\right)$ ，使得等式 $r\left( {x}_{0}\right)  = m$ 成立? 请说明理由.

9.(2024$\cdot$长宁区一模$\cdot$21)若函数 $y = f\left( x\right)$ 与 $y = g\left( x\right)$ 满足:对任意 ${x}_{1}$ 、 ${x}_{2} \in  \mathbf{R}$ ，都有 $\mid  f\left( {x}_{1}\right)  -  f\left( {x}_{2}\right) \left|  \geq  \right| g\left( {x}_{1}\right)  - g\left( {x}_{2}\right)  \mid$ ,则称函数 $y = f\left( x\right)$ 是函数 $y = g\left( x\right)$ 的“约束函数”. 已知函数 $y =  f\left( x\right)$ 是函数 $y = g\left( x\right)$ 的“约束函数”.

(1)若 $f\left( x\right)  = {x}^{2}$ ，判断函数 $y = g\left( x\right)$ 的奇偶性，并说明理由；

(2)若 $f\left( x\right)  = {ax} + {x}^{3}\left( {a > 0}\right) ,g\left( x\right)  = \sin x$ ，求实数 $a$ 的取值范围；

(3)若 $y = g\left( x\right)$ 为严格减函数， $f\left( 0\right)  < f\left( 1\right)$ ，且函数 $y = f\left( x\right)$ 的图像是连续曲线，求证: $y =  f\left( x\right)$ 是 $\left( {0,1}\right)$ 上的严格增函数.

\section*{第 4 章 幂函数、指数函数与对数函数}

\begin{center}

\begin{tabular}{|p{4cm}|p{12cm}|}
\hline
考点 & 课标要求与命题趋势 \\
\cline{1-2}
1. 幂函数 & 通过具体实例,结合 $y = x\text{ 、 }y = \dfrac{1}{x}\text{ 、 }y = {x}^{2}\text{ 、 }y = \sqrt{x}\text{ 、 }y = {x}^{3}$ 的图像,理解它们的变化规律, 了解幂函数的性质 \\
\cline{1-2}
2. 指数函数及其应用 & \ding{172}  通过对有理指数幂 ${a}^{\frac{m}{n}}$ ( $a > 0$ ,且 $a \neq  1,m\text{ 、 }n$ 为整数，且 $n > 0$ )、实数指数幂 ${a}^{x}$ ( $a > 0$ ，且 $a \neq  1,x \in  \mathbf{R}$ ) 含义的认识,了解指数幂的拓展过程,掌握指数幂的运算性质 \ding{173}  通过具体实例，了解指数函数的实际意义，理解指数函数的概念 \ding{174}  能用描点法或借助计算工具画出具体指数函数的图像，探索并理解指数函数的单调性与特殊点 \\
\cline{1-2}
3. 对数函数及其应用 & \ding{172}  理解对数的概念和运算性质，知道用换底公式能将一般对数转化成自然对数或常用对数 \ding{173}  通过具体实例，了解对数函数的概念. 能用描点法或借助计算工具画出具体对数函数的图像，探索并了解对数函数的单调性与特殊点 \ding{174}  知道对数函数 $y = {\log }_{a}x$ 与指数函数 $y = {a}^{x}$ 互为反函数 $\left( {a > 0\text{ ,且 }a \neq  1}\right)$ \ding{175} 收集、阅读对数概念的形成与发展的历史资料，撰写小论文，论述对数发明的过程以及对数对简化运算的作用 \\
\cline{1-2}
4. 幂、指数、对数函数值大小比较 & 掌握幂函数、指数函数、对数函数的图像与性质, 掌握幂、指数、对数的相关运算, 对指数、 对数、幂函数值的进行大小比较，都是上海高考命题的方向 \\
\cline{1-2}
\hline
\end{tabular}
\end{center}
\clearpage
\subsection*{考点1}

\subsubsection*{一、填空题}

1.(2024$\cdot$宝山区二模$\cdot$3)将 $\sqrt{{a}^{2}\sqrt{a}}$ (其中 $a > 0$ )化为有理数指数幂的形式为\blkda[2]{}.

2.(2025$\cdot$徐汇区二模$\cdot$4)已知幂函数 $y = f\left( x\right)$ 的图像过点 $\left( {3,\dfrac{\sqrt{3}}{3}}\right)$ ，则该幂函数的值域是\blkda[2]{}.

3.(2024$\cdot$崇明区二模$\cdot$6)已知幂函数 $y = f\left( x\right)$ 的图像经过点(2，4)，则 $f\left( 3\right)  =$ \blkda[2]{}.

4. (2024$\cdot$浦东新区二模$\cdot$6)已知 $y = f\left( x\right)$ 是奇函数，当 $x \geq  0$ 时， $f\left( x\right)  = {x}^{\dfrac{2}{3}}$ ，则 $f\left( {-\dfrac{8}{125}}\right)$ 的值是\blkda[2]{}.

5.(2024$\cdot$静安区一模$\cdot$7)下列幂函数在区间 $\left( {0, + \infty }\right)$ 上是严格增函数，且图像关于原点成中心对称的是\blkda[2]{}. (请填入全部正确的序号)

 \ding{172}  $y = {x}^{\dfrac{1}{2}}$ ;  \ding{173}  $y = {x}^{\dfrac{1}{3}}$ ;  \ding{174}  $y = {x}^{\dfrac{2}{3}}$ ； \ding{175}  $y = {x}^{-\dfrac{1}{3}}$.



\begin{center}
\includegraphics[width=0.2\textwidth]{bodyxb/src/33_1094_187_289_241_0.jpg}
\end{center}


第 6 题图

6.(2025$\cdot$长宁区一模$\cdot$7)已知 $a \in  \left\{  {-1, - \dfrac{2}{3}, - \dfrac{1}{3},\dfrac{1}{3},\dfrac{2}{3},1,2,3}\right\}$ ，函数 $y = {x}^{a}$ 的大致图像如图所示，则 $a =$ \blkda[2]{}.

7.(2025$\cdot$黄浦区一模$\cdot$9)若 $f\left( x\right)  = {x}^{3},g\left( x\right)  = \left\{  \begin{array}{l} f\left( x\right) ,x \geq  0, \\  f\left( {-x}\right) ,x < 0, \end{array}\right.$ 则不等式 $g\left( x\right)  <  - x$ 的解集为\blkda[2]{}.

\subsubsection*{二、选择题}

8.(2022$\cdot$上海春季卷$\cdot$13)以下函数的定义域为 $\mathbf{R}$ 的是 \dotfill{(\kh{})}

\xx{}{}{}{}$y = {x}^{-1}$ B. $y = {x}^{-\dfrac{1}{2}}$ C. $y = {x}^{\dfrac{1}{3}}$ D. $y = {x}^{\dfrac{1}{2}}$

9.(2025$\cdot$上海春季卷$\cdot$14)幂函数 $y = {x}^{a}$ 在(0，+∞)上是严格减函数，且经过(-1，-1)，则 $a$ 的值可能是 \dotfill{(\kh{})}

\xx{}{}{}{}$- \dfrac{2}{3}$ B. $- \dfrac{1}{3}$ C. $\dfrac{1}{3}$ D. 3

\subsubsection*{三、解答题}

10.(2025$\cdot$黄浦区二模$\cdot$17)已知 $f\left( x\right)  = {2}^{x}$.

(1)若 $f\left( x\right)  - f\left( {2 - x}\right)  = 3$ ，求 $x$ 的值；

(2)是否存在实数 $a$ ，使函数 $y = f\left( x\right)  + {af}\left( {-x}\right)$ 是奇函数？请说明理由.

11. (2025$\cdot$虹口区二模$\cdot$18)已知函数 $y = f\left( x\right)$ 的表达式为 $f\left( x\right)  = {2}^{x},x \in  \mathbf{R}$.

(1)解不等式: ${\log }_{2}\left\lbrack  {f\left( x\right)  - 1}\right\rbrack   + {\log }_{2}\left\lbrack  {f\left( x\right) }\right\rbrack   \leq  1$ ；

( 2 )若存在实数 ${x}_{0} \in  \left\lbrack  {0,\dfrac{\pi }{2}}\right\rbrack$ ，使得 $f\left( {m\sin {x}_{0}}\right)$ ， $f\left( {2 + \sin {2{x}_{0}}}\right)$ ， $f\left( {m\cos {x}_{0}}\right)$ 成等比数列，求实数 $m$ 的最小值.

\subsection*{考点2 分 指数函数及其应用}

\subsubsection*{一、填空题}

1.(2025$\cdot$奉贤区一模$\cdot$4)设 $f\left( x\right)  = \left\{  \begin{array}{l} \ln x + 1,x > 0, \\  {2}^{x} + 1,x \leq  0, \end{array}\right.$ 若 $f\left( {x}_{0}\right)  = 1$ ，则函数 $y = {a2} + b$ \blkda[2]{}.

2.(2024$\cdot$宝山区一模$\cdot$6)设 $a$ 、 $b$ 为常数，若 $a > 1$ ， $b <  - 1$ ，则函数 $y = {a}^{x} + b$ 的图像必定不经过第\blkda[2]{}象限.

3. (2024$\cdot$宝山区一模$\cdot$7)设函数 $f\left( x\right)  = \left\{  \begin{array}{l} \dfrac{1}{2}x - 1\left( {x \geq  0}\right) , \\  \dfrac{1}{x}\left( {x < 0}\right) , \end{array}\right.$ 若 $f\left( a\right)  = a$ ，则实数 $a$ 的值为\blkda[2]{}.

4.(2025$\cdot$崇明区一模$\cdot$8)已知 $f\left( x\right)  = \left\{  \begin{array}{l} {2}^{x} - 1,x > 1, \\  x - 1,x \leq  1, \end{array}\right.$ 关于 $x$ 的方程 $f\left( x\right)  = 2$ 的解 $x =$ \blkda[2]{}.

5.(2024$\cdot$奉贤区一模$\cdot$8)已知函数 $y = f\left( x\right)$ 是定义在 $\mathbf{R}$ 上的奇函数，当 $x \geq  0$ 时， $f\left( x\right)  = {{e}}^{x} + b$ ( $b$ 是常数),则 $f\left( {-\ln 2}\right)  =$ \blkda[2]{}.

6.(2024$\cdot$杨浦区二模$\cdot$8)若函数 $g\left( x\right)  = \left\{  \begin{array}{l} {2}^{x} - 1,x \leq  0, \\  f\left( x\right) ,x > 0 \end{array}\right.$ 为奇函数，则函数 $y = f\left( x\right) ,x \in  \left( {0, + \infty }\right)$ 的值域为\blkda[2]{}.

7. (2025$\cdot$松江区一模$\cdot$9)已知函数 $y = f\left( x\right)$ 的表达式为 $f\left( x\right)  = \left\{  \begin{array}{l} {3}^{x},x \geq  0, \\  \dfrac{1}{{3}^{x}},x < 0, \end{array}\right.$ 则满足 $f\left( m\right)  \geq  f\left( {m + 2}\right)$ 的实数 $m$ 的最大值为\blkda[2]{}.

8.(2023$\cdot$上海春季卷$\cdot$9)已知函数 $f\left( x\right)  = {2}^{-x} + 1$ ，且 $g\left( x\right)  = \left\{  \begin{array}{l} {\log }_{2}\left( {x + 1}\right) ,x \geq  0, \\  f\left( {-x}\right) ,x < 0, \end{array}\right.$ 则方程 $g\left( x\right)  = 2$ 的解为\blkda[2]{}.

9.(2024$\cdot$崇明区二模$\cdot$9)已知函数 $y = \left\{  \begin{array}{l} {3}^{x} - x,x \leq  0, \\  f\left( x\right) ,x > 0 \end{array}\right.$ 为奇函数，则 $f\left( 2\right)  =$ \blkda[2]{}.

10.(2024$\cdot$浦东新区二模$\cdot$9)已知 $f\left( x\right)  = {2}^{x} + x$ ，则不等式 $f\left( \left| {{2x} - 3}\right| \right)  < 3$ 的解集为\blkda[2]{}.

11.(2024$\cdot$松江区二模$\cdot$11)已知 $0 < a < 2$ ，函数 $y = \left\{  \begin{array}{l} \left( {a - 2}\right) x + {4a} + 1,x \leq  2, \\  2{a}^{x - 1},x > 2, \end{array}\right.$ 若该函数存在最小值， 则实数 $a$ 的取值范围是\blkda[2]{}.

\subsubsection*{二、选择题}

12.(2024$\cdot$徐汇区二模$\cdot$13)在下列函数中，值域为 $\mathbf{R}$ 的偶函数是 \dotfill{(\kh{})}

\xx{}{}{}{}$y = {x}^{\dfrac{1}{3}}$ B. $y = \lg \left| x\right|$ C. $y = {{e}}^{x} + {{e}}^{-x}$ D. $y = {x}^{3}\cos x$

13.(2025$\cdot$上海秋季卷$\cdot$14)设 $a > 0,s \in  \mathbf{R}$. 下列各项中，能推出 ${a}^{s} > a$ 的一项是 \dotfill{(\kh{})}

\xx{}{}{}{}$a > 1$ ，且 $s > 0$ B. $a > 1$ ，且 $s < 0$ C. $0 < a < 1$ ，且 $s > 0$ D. $0 < a < 1$. 且 $s < 0$

14.(2024$\cdot$奉贤区一模$\cdot$14)函数 $y = \dfrac{{2}^{x} - 1}{{2}^{x} + 1}$ 在定义域(一∞，+∞)上是 \dotfill{(\kh{})}

\xx{}{}{}{}严格增的奇函数 B. 严格增的偶函数 C. 严格减的奇函数 D. 严格减的偶函数

15.(2024$\cdot$金山区二模$\cdot$16)设 $f\left( x\right)  = {x}^{3} - {3x}$ ，有如下两个命题:

 \ding{172}  函数 $y = f\left( x\right)$ 的图像与圆 ${x}^{2} + {y}^{2} = 1$ 有且只有两个公共点；

 \ding{173}  存在唯一的正方形 ${ABCD}$ ，其四个顶点都在函数 $y = f\left( x\right)$ 的图像上.

则下列说法正确的是 \dotfill{(\kh{})}

\xx{}{}{}{} \ding{172} 正确， \ding{173} 正确 B.  \ding{172} 正确， \ding{173} 不正确 C.  \ding{172} 不正确， \ding{173} 正确 D.  \ding{172} 不正确， \ding{173} 不正确

\subsubsection*{三、解答题}

16.(2024$\cdot$黄浦区二模$\cdot$17)设 $a \in  \mathbf{R}$ ，函数 $f\left( x\right)  = \dfrac{{2}^{x} + a}{{2}^{x} - 1}$.

(1)求 $a$ 的值，使得 $y = f\left( x\right)$ 为奇函数；

(2)若 $f\left( 2\right)  = a$ ，求满足 $f\left( x\right)  > a$ 的实数 $x$ 的取值范围.

17.(2024$\cdot$浦东新区一模$\cdot$17)已知函数 $y = f\left( x\right)$ ，其中 $f\left( x\right)  = \dfrac{{4}^{x} + k}{{2}^{x}}\left( {k \in  \mathbf{R}}\right)$.

(1)是否存在实数 $k$ ，使函数 $y = f\left( x\right)$ 是奇函数？若存在，请写出证明；

(2)当 $k = 1$ 时,若关于 $x$ 的不等式 $f\left( x\right)  \geq  a$ 恒成立，求实数 $a$ 的取值范围.

18.(2025$\cdot$奉贤区一模$\cdot$17)已知函数 $y = f\left( x\right)$ ，其中 $f\left( x\right)  = {a}^{x}$ (常数 $a > 0$ 且 $a \neq  1$ ).

(1)若函数 $y = f\left( x\right)$ 的图像过点 $\left( {2,9}\right)$ ，求关于 $x$ 的不等式 $f\left( \left| {{2x} - 1}\right| \right)  > 3$ 的解集；

(2)若存在 $x \in  (0,1\rbrack$ ，使得数列 $f\left( 1\right) ,f\left( {tx}\right) ,f\left( {{x}^{2} + 2}\right)$ 是等比数列，求实数 $t$ 的取值范围.

\subsection*{考点3 对数函数及其应用}

\subsubsection*{一、填空题}

1.(2024$\cdot$上海春季卷$\cdot$1)函数 $f\left( x\right)  = {\log }_{2}x$ 的定义域为\blkda[2]{}.

2.(2025$\cdot$嘉定区一模$\cdot$1)函数 $y = {\log }_{2}\left( {{x}^{2} - 1}\right)$ 的定义域为\blkda[2]{}.

3.(2025$\cdot$浦东新区一模$\cdot$1)若对数函数 $y = {\log }_{a}x$ ( $a > 0$ 且 $a \neq  1$ )的图像经过点 $\left( {4,2}\right)$ ，则实数 $a =$ \blkda[2]{}.

4.(2025$\cdot$虹口区一模$\cdot$2)函数 $y = \ln \dfrac{x}{x - 1}$ 的定义域是\blkda[2]{}.

5.(2025$\cdot$徐汇区一模$\cdot$2)已知函数 $y = f\left( x\right)$ ，其中 $f\left( x\right)  = \left\{  \begin{array}{l} \ln x,x > 0, \\   - 1,x \leq  0, \end{array}\right.$ 则 $f\left( 1\right)  =$ \blkda[2]{}.

6.(2025$\cdot$虹口区一模$\cdot$5)设 $a > 0$ 且 $a \neq  1$ ，则函数 $y = 2 + {\log }_{a}x$ 的图像恒过的定点坐标为\blkda[2]{}.

7.(2025$\cdot$静安区二模$\cdot$5)已知 ${\log }_{3}2 = a$ ，则 ${\log }_{2}{48} =$ \blkda[2]{}.(请用含 $a$ 的代数式表达)

8.(2024$\cdot$长宁区二模$\cdot$5)已知 ${2}^{a} = {3}^{b} = 6$ ，则 $\dfrac{1}{a} + \dfrac{1}{b} =$ \blkda[2]{}.

9.(2025$\cdot$宝山区二模$\cdot$7)已知函数 $y = {a}^{x + 1} - {\log }_{a}\left( {x + 2}\right)  + 1$ ( $a > 0$ 且 $a \neq  1$ )的图像经过定点 $A$ ，则点 $A$ 的坐标为\blkda[2]{}.

10. (2024$\cdot$长宁区二模$\cdot$8)已知函数 $y = f\left( x\right)$ 是定义域为 $\mathbf{R}$ 的奇函数，当 $x > 0$ 时， $f\left( x\right)  = {\log }_{2}x$ ，若 $f\left( a\right)  > 1$ ，则实数 $a$ 的取值范围为\blkda[2]{}.

11. (2025$\cdot$宝山区一模$\cdot$8)若 ${9}^{a} = {4}^{b} = m$ ，且 $\dfrac{1}{a} + \dfrac{1}{b} = 2$ ，则 $m =$ \blkda[2]{}.

12.(2024$\cdot$青浦区二模$\cdot$8)已知 $f\left( x\right)  = \lg x - 1,g\left( x\right)  = \lg x - 3$ ，若 $\left| {f\left( x\right) }\right|  + \left| {g\left( x\right) }\right|  =  \mid  f\left( x\right)  +  g\left( x\right)  \mid$ ，则满足条件的 $x$ 的取值范围是\blkda[2]{}.

13.(2024$\cdot$长宁区一模$\cdot$8)在有声世界，声强级是表示声强度相对大小的指标，其值 $y$ (单位: ${{dB}}$ )定义为 $y = {10}\lg \dfrac{I}{{I}_{0}}$. 其中 $I$ 为声场中某点的声强度,单位为 ${W}/{{m}}^{2},{I}_{0} = {10}^{-{12}}{W}/{{m}}^{2}$ 为基准值. 若 $I = {10}{W}/{{m}}^{2}$ ，则其相应的声强级为\blkda[2]{} ${{dB}}$.

14.(2024$\cdot$松江区二模$\cdot$8)已知函数 $y = \left| {{\log }_{2}x}\right|$ ，记 $y = f\left( x\right)$ ，若 $f\left( {x}_{1}\right)  = f\left( {x}_{2}\right)$ ( ${x}_{1} \neq  {x}_{2}$ )，则 $4{x}_{1} + {x}_{2}$ 的最小值为\blkda[2]{}.

15.(2025$\cdot$嘉定区一模$\cdot$9)已知 $f\left( x\right)  = \ln \left( {x + 1}\right)$ ， $g\left( x\right)  = \left\{  \begin{array}{l} f\left( x\right) ,x \geq  0, \\  f\left( {-x}\right) ,x < 0, \end{array}\right.$ 则 $g\left( x\right)  > x + 2 - {e}$ 的解集为\blkda[2]{}.

16.(2025$\cdot$青浦区一模$\cdot$11)若函数 $y = {\log }_{\dfrac{1}{2}}\left( {a{x}^{2} - {8x} + {15}}\right)$ 在区间(1,2)上严格增，则实数 $a$ 的取值范围是\blkda[2]{}.

\subsubsection*{二、选择题}

17.(2025$\cdot$静安区一模$\cdot$14)污水处理厂通过清除污水中的污染物获得清洁用水并生产肥料. 该厂的污水处理装置每小时从处理池清除掉 12\% 的污染残留物. 要使处理池中的污染物水平降到最初的 10\%，大约需要的时间为(参考数据:lo 0.88\%=0.0555) \dotfill{(\kh{})}

\xx{}{}{}{}${14}{h}$ B. ${18}{h}$ C. ${20}{h}$ D. ${24}{h}$

\subsubsection*{三、解答题}

18.(2025$\cdot$闵行区二模$\cdot$17)已知函数 $y = f\left( x\right)$ 是定义在 $\mathbf{R}$ 上的奇函数，当 $x > 0$ 时， $f\left( x\right)  = {\log }_{2}x$. (1)求 $f\left( {-2}\right)  + f\left( 0\right)$ 的值；

(2)若 $g\left( x\right)  = f\left( x\right)  \cdot  f\left( \dfrac{x}{4}\right) ,x \in  \left\lbrack  {1,8}\right\rbrack$ ，求函数 $y = g\left( x\right)$ 的值域.

19.(2025$\cdot$金山区一模$\cdot$17)已知常数 $a > 1$ ，函数 $y = f\left( x\right)$ 的表达式为 $f\left( x\right)  = {\log }_{a}\left( {x + 2}\right)  - {\log }_{a}\left( {2 - x}\right)$.

(1)证明:函数 $y = f\left( x\right)$ 是奇函数；

(2)若函数 $y = f\left( x\right)$ 在区间 $\left\lbrack  {0,1}\right\rbrack$ 上的最大值为 2，求实数 $a$ 的值.

20. (2025$\cdot$青浦区二模$\cdot$17)对于函数 $y = f\left( x\right)$ ，其中 $f\left( x\right)  = {\log }_{a}x\left( {a > 0,a \neq  1}\right)$.

(1)若函数 $y = f\left( x\right)$ 的图像过点 $\left( {4,2}\right)$ ，求 $f\left( {{2x} - 2}\right)  < f\left( x\right)$ 的解集；

(2)求证:当 $a = \sqrt{2}$ 时，存在 $x$ 使得 $f\left( {x + 1}\right)$ 、 $f\left( {ax}\right)$ 、 $f\left( {x + 2}\right)$ 成等差数列.

21. (2024$\cdot$徐汇区二模$\cdot$17)已知函数 $y = f\left( x\right)$ ，其中 $f\left( x\right)  = {\log }_{\dfrac{1}{2}}\dfrac{2 + x}{x - 2}$.

(1)求证: $y = f\left( x\right)$ 是奇函数；

(2)若关于 $x$ 的方程 $f\left( x\right)  = {\log }_{\dfrac{1}{2}}\left( {x + k}\right)$ 在区间 $\left\lbrack  {3,4}\right\rbrack$ 上有解,求实数 $k$ 的取值范围.

22.(2024$\cdot$上海秋季卷$\cdot$18)若 $f\left( x\right)  = {\log }_{a}x\left( {a > 0,a \neq  1}\right)$.

(1)若函数 $y = f\left( x\right)$ 的图像过点 $\left( {4,2}\right)$ ，求 $f\left( {{2x} - 2}\right)  < f\left( x\right)$ 的解集；

(2)存在 $x$ 使得 $f\left( {x + 1}\right) ,f\left( {ax}\right) ,f\left( {x + 2}\right)$ 成等差数列,求 $a$ 的取值范围.

\subsection*{考点4 > 幂指对函数值大小比较}

\subsubsection*{一、填空题}

1. (2025$\cdot$静安区一模$\cdot$12)已知 $\lg {x}_{1}$ 、 $\lg {x}_{2}$ 、 $\lg {x}_{3}$ 、 $\lg {x}_{4}$ 、 $\lg {x}_{5}$ 是从大到小连续的正整数，且 ${\left( \lg {x}_{4}\right) }^{2} <  \lg {x}_{1} \cdot  \lg {x}_{5}$ ，则 ${x}_{1}$ 的最小值为\blkda[2]{}.

\subsubsection*{二、选择题}

2.(2024$\cdot$宝山区二模$\cdot$13)已知 $a > b > 0$ ，则 \dotfill{(\kh{})}

\xx{}{}{}{}${a}^{2} > {b}^{2}$ B. ${2}^{a} < {2}^{b}$ C. ${a}^{\dfrac{1}{2}} < {b}^{\dfrac{1}{2}}$ D. ${\log }_{\dfrac{1}{2}}a > {\log }_{\dfrac{1}{2}}b$

3.(2024$\cdot$杨浦区一模$\cdot$13)已知实数 $a$ 、 $b$ 满足 $a > b$ ，则下列不等式恒成立的是 \dotfill{(\kh{})}

\xx{}{}{}{}${a}^{2} > {b}^{2}$ B. ${a}^{3} > {b}^{3}$ C. $\left| a\right|  > \left| b\right|$ D. ${a}^{-1} > {b}^{-1}$

4.(2025$\cdot$青浦区一模$\cdot$13)已知 $x$ 、 $y \in  \mathbf{R}$ 且满足 $x > y$ ，则下列关系式恒成立的是 \dotfill{(\kh{})}

\xx{}{}{}{}$\dfrac{1}{{x}^{2} + 1} < \dfrac{1}{{y}^{2} + 1}$ B. $\ln \left( {{x}^{2} + 1}\right)  > \ln \left( {{y}^{2} + 1}\right)$

C. $\sin x > \sin y$ D. ${x}^{3} > {y}^{3}$

\subsubsection*{三、解答题}

5.(2025$\cdot$杨浦区二模$\cdot$17)已知函数 $y = f\left( x\right)$ 是定义在 $\mathbf{R}$ 上的偶函数.

(1)当 $x \in  \lbrack 0, + \infty )$ 时， $f\left( x\right)  = {\log }_{2}\left( {x + 1}\right)$ ，求 $x \in  \left( {-\infty ,0}\right)$ 时， $y = f\left( x\right)$ 的表达式；

(2)当 $x \in  \lbrack 0, + \infty )$ 时， $f\left( x\right)  = {x}^{3} + {2025}{}^{x}$ ，若实数 $t$ 满足 $f\left( {{3t} - 2}\right)  < f\left( {t - 1}\right)$ ，求 $t$ 的取值范围.

6.(2025$\cdot$宝山区二模$\cdot$18)已知函数 $y = f\left( x\right)$ ，其中 $f\left( x\right)  = {a}^{x}\left( {a > 0\text{ 且 }a \neq  1}\right)$ ，

(1)若 $f\left( 2\right)  = 4$ ，求方程 $f\left( x\right)  - f\left( {-x}\right)  = 2$ 的解；

(2)已知 $0 < a < 1$ ，若关于 $x$ 的不等式 ${\left\lbrack  f\left( mx\right) \right\rbrack  }^{2} \geq  f\left( {{x}^{2} + 1}\right)$. $f\left( {x + 3}\right)$ 在区间 $\left\lbrack  {1,2}\right\rbrack$ 上恒成立， 求实数 $m$ 的最大值.

7.(2025$\cdot$崇明区二模$\cdot$18)已知 $f\left( x\right)  = {\log }_{3}\left( {x + a}\right)  + {\log }_{3}\left( {6 - x}\right)$.

(1)是否存在实数 $a$ ，使得函数 $y = f\left( x\right)$ 是偶函数？若存在，求实数 $a$ 的值；若不存在，请说明理由;

(2)若 $a >  - 3$ 且 $a \neq  0$ ，解关于 $x$ 的不等式 $f\left( x\right)  \leq  f\left( {6 - x}\right)$.

8.(2025$\cdot$奉贤区二模$\cdot$18)函数 $y = f\left( x\right)$ ，其中 $f\left( x\right)  = {{e}}^{-\dfrac{{\left( x - \mu \right) }^{2}}{2}}$.

(1)若函数 $y = f\left( x\right)$ 是偶函数，当 $f\left( t\right)  = \dfrac{1}{\sqrt{{e}}}$ 时，求 $t$ 的值；

(2)求函数 $y = f\left( x\right)$ 的值域，并证明对任意的正实数 $k$ 和实数 $x$ ，不等式 $k + \dfrac{1}{k} \geq  {2f}\left( x\right)$ 恒成立.

9.(2025$\cdot$徐汇区二模$\cdot$18)已知函数 $y = f\left( x\right)$ ，其中 $f\left( x\right)  = {\log }_{2}x$.

(1)解关于 $x$ 的不等式 $f\left( {{3x} - 2}\right)  < f\left( {{2x} + 1}\right)$ ；

(2)若存在唯一的实数 ${x}_{0}$ ，使得 $f\left( {x}_{0}\right)$ 、 $f\left( {{x}_{0} - a}\right)$ 、 $f\left( 2\right)$ 依次成等差数列，求实数 $a$ 的取值范围.

10.(2025$\cdot$徐汇区二模$\cdot$21)对于函数 $y = h\left( x\right)$ ，记 ${h}^{\left( 0\right) }\left( x\right)  = h\left( x\right)$ ， ${h}^{\left( 1\right) }\left( x\right)  = {\left( h\left( x\right) \right) }^{\prime }$ ， $\cdots$ ， ${h}^{\left( n + 1\right) }\left( x\right)  = {\left( {h}^{\left( n\right) }\left( x\right) \right) }^{\prime }\left( {n \in  \mathbf{N}}\right)$. 如果 $n$ 是满足 ${h}^{\left( n\right) }\left( x\right)  = h\left( x\right)$ 的最小正整数，则称 $n$ 是函数 $y =  h\left( x\right)$ 的“最小导周期”.

(1)已知函数 $y = f\left( x\right)$ ，其中 $f\left( x\right)  = a\sin \left( {x + t}\right)  + b\cos \left( {x + t}\right)$ ，求证:对任意实数 $a$ 、 $b$ 、 $t$ ，都有 ${f}^{\left( 4\right) }\left( x\right)  = f\left( x\right) ;$

(2)设 $m\text{ 、 }n \in  \mathbf{R},g\left( x\right)  = {{e}}^{mx} + n\cos x$ ，若函数 $y = g\left( x\right)$ 的最小导周期为 2，记 $M\left( {a,b}\right)  =  \sqrt{{\left( a - b\right) }^{2} + {\left( a + 1 + g\left( b\right) \right) }^{2}}$ ，当实数 $a$ 、 $b$ 变化时，求 $M\left( {a,b}\right)$ 的最小值；

(3)设 $\omega  > 1,h\left( x\right)  = \cos {\omega x}$ ，若函数 $y = h\left( x\right)$ 满足 ${h}^{\left( 2\right) }\left( x\right)  \leq  x$ 对 $x \in  \left( {0, + \infty }\right)$ 恒成立，且存在 ${x}_{0} \in  \left( {0, + \infty }\right)$ 使得 ${h}^{\left( 2\right) }\left( {x}_{0}\right)  = {x}_{0}$ ,试用 $\omega$ 表示 ${x}_{0}$ ,并证明 $\dfrac{\pi }{2\omega } < {x}_{0} < \dfrac{\pi }{\omega }$.

11. (2025$\cdot$杨浦区二模$\cdot$21)已知函数 $y = f\left( x\right)$ 的导函数为 $y = {f}^{\prime }\left( x\right)$ ，若函数 $y = f\left( x\right)$ 的定义域为 $\mathbf{R}$ ，且不等式 $f\left( x\right)  > {f}^{\prime }\left( x\right)$ 对任意 $x \in  \mathbf{R}$ 成立，则称函数 $y = f\left( x\right)$ 是“超导函数”.

(1)判断 $f\left( x\right)  = {{e}}^{x} + 1$ 是否为“超导函数”，并说明理由;

(2)若函数 $y = g\left( x\right)$ 与 $y = h\left( x\right)$ 都是“超导函数”，且对任意 $x \in  \mathbf{R}$ ，都有 ${h}^{\prime }\left( x\right)  > 0,{g}^{\prime }\left( x\right)  < 0$ ，记 $F\left( x\right)  = g\left( x\right) h\left( x\right)$ ,求证: 函数 $y = F\left( x\right)$ 是 “超导函数”;

(3)已知函数 $y = \varphi \left( x\right)$ 是 “超导函数” 且 $\varphi \left( 1\right)  = {e}$ ，若有且仅有一个实数 $t$ 满足 $\varphi \left( {\ln t + 1 - {at}}\right)  =  {{e}}^{\ln t + 1 - {at}}$ ,求 $a$ 的取值范围.

\section*{第 5 章 三 角 函 数}

\begin{center}

\begin{tabular}{|p{4cm}|p{12cm}|}
\hline
考点 & 课标要求与命题趋势 \\
\cline{1-2}
1. 任意角的正弦、 余弦、正切、余切及同角三角比的基本关系 & \ding{172}  了解任意角的概念和弧度制，能进行弧度与角度的互化，体会引入弧度制的必要性 \ding{173}  借助单位圆理解任意角三角函数(正弦、余弦、正切)的定义，并能利用三角函数的定义解决相关问题，理解并掌握同角三角函数的基本关系式，能够利用公式化简求值； 能画出这些三角函数的图像，了解三角函数的周期性、奇偶性、最大(小)值\\
\cline{1-2}
2. 诱导公式及其应用 & 借助单位圆的对称性,利用定义推导出诱导公式 $\left( {\alpha  \pm  \dfrac{\pi }{2},\alpha  \pm  \pi }\right)$的正弦、余弦、正切，能够运用诱导公式解决相关问题，该内容是上海高考的必考内容，一般会考查三角函数化简求值或特殊角求三角函数值，需加强复习备考 \\
\cline{1-2}
3. 两角和与差的正弦、余弦、正切公式及二倍角公式 & \ding{172}  借助图像理解正弦函数在 $\left\lbrack  {-\dfrac{\pi }{2},\dfrac{\pi }{2}}\right\rbrack$ 上、余弦函数在 $\left\lbrack  {0,{2\pi }}\right\rbrack$ 上、正切函数在 $\left( {-\dfrac{\pi }{2},\dfrac{\pi }{2}}\right)$ 上的性质 \ding{173}  经历推导两角差余弦公式的过程，知道两角差余弦公式的意义 \ding{174}  能从两角差的余弦公式推导出两角和与差的正弦、余弦、正切公式，能推导二倍角的正弦、余弦、正切公式，了解它们的内在联系 \ding{175}  能运用这些公式解决相关的求值与化简问题，该内容是上海高考的必考内容，一般会考查两角和与差的正弦、余弦、正切公式及倍角公式变形应用和半角公式变形应用， 同时也需掌握升幂公式和降幂公式，掌握拼凑角的思路 \\
\cline{1-2}
4. 简单的三角恒等变换 & \ding{172}  能运用常用的三角公式进行简单的恒等变换(包括推导出积化和差、和差化积、半角公式，这三组公式不要求记忆) \ding{173} 结合具体实例,了解 $y = A\sin \left( {{\omega x} + \varphi }\right)$ 的实际意义; 能借助图像理解参数 $\omega$ 、 $\varphi$ 、 $A$ 的意义，了解参数的变化对函数图像的影响 \\
\cline{1-2}
5. 解三角形 & 掌握正弦定理、余弦定理及其相关变形应用，会用三角形的面积公式解决与面积有关的计算问题,会用正弦定理、余弦定理等知识和方法解决三角形中的综合问题,会利用基本不等式和相关函数性质解决三角形中的最值及范围问题，该内容是上海高考的常考内容 \\
\cline{1-2}
6. 解三角形之实际应用 & 一般考查正、余弦定理和三角形面积公式在解三角形中的应用，结合三角函数及三角恒等变换等知识点进行综合考查, 也常结合基本不等式和相关函数性质等知识点求解范围及最值，需重点复习 \\
\cline{1-2}
7. 三角函数的图像与性质 & 能用五点作图法作出正弦、余弦和正切函数图像，理解并掌握三角函数的图像与性质; 会求参数及函数解析式. 该内容是上海高考的必考内容,一般会综合考查三角函数的图像与性质的综合应用，需加强复习备考 \\
\cline{1-2}
8. 三角函数图像的伸缩平移变换 & 会先平移后伸缩或先伸缩后平移来综合解决三角函数的伸缩平移变换，会用三角函数解决简单的实际问题，体会可以利用三角函数构建刻画事物周期变化的数学模型 \\
\cline{1-2}
\hline
\end{tabular}
\end{center}
\clearpage
\subsection*{考点1 任意角的正弦、余弦、正切、余切及同角三角比的基本关系}

\subsubsection*{一、填空题}
\begin{enumerate}
    \item (2024$\cdot$松江区二模$\cdot$4)已知点 $A$ 的坐标为 $\left( {\dfrac{1}{2},\dfrac{\sqrt{3}}{2}}\right)$ ，将 ${OA}$ 绕坐标原点 $O$ 逆时针旋转 $\dfrac{\pi }{2}$ 至 ${OP}$ ，则点 $P$ 的坐标为\blkda[2]{$\left( {-\dfrac{\sqrt{3}}{2},\dfrac{1}{2}}\right)$}.
    
    \item (2025$\cdot$闵行区二模$\cdot$4)若 $\alpha$ 为第二象限角，且 $\sin \alpha  = \dfrac{4}{5}$ ，则 $\tan \alpha  =$ \blkda[2]{$- \dfrac{4}{3}$}.

\end{enumerate}
\subsubsection*{二、选择题}
\begin{enumerate}
    \item (2025$\cdot$虹口区一模$\cdot$13)已知 $\alpha  \in  \left( {0,\pi }\right)$ ，则“$\sin \left( {\pi  - \alpha }\right)  = \dfrac{1}{2}$ ”是“$\cos \alpha  = \dfrac{\sqrt{3}}{2}$”的\dotfill{(\kh{C})}
    \xx{充要条件}{充分非必要条件}{必要非充分条件}{既非充分又非必要条件} 
    
    \item (2024$\cdot$静安区一模$\cdot$14)设$\alpha$是第一象限的角，则 $\dfrac{\alpha }{2}$ 所在的象限为 \dotfill{(\kh{C})}
    \xx{第一象限}{第三象限}{第一象限或第三象限}{第二象限或第四象限}

    \item (2025$\cdot$奉贤区一模$\cdot$14)若函数 $y = {\log }_{2}\sin x + {\log }_{2}\cos x$ ，则下列命题正确的是 \dotfill{(\kh{C})}
    \xx{函数是偶函数}{函数定义域是 $\left( {0,\dfrac{\pi }{2}}\right)$}{函数最大值 $-1$}{函数的最小正周期为 $\pi$}
    \begin{jieda}
        由题意知 $\left\{  {\begin{array}{l} \sin x > 0 \\  \cos x > 0 \end{array} \Longrightarrow  x \in  \left( {{2k\pi },{2k\pi } + \dfrac{\pi }{2}}\right) ,k \in  \mathbf{Z}}\right.$，定义域不关于原点对称,则不具有奇偶性，最小正周期为$2\pi$.
        
        
        $y = {\log }_{2}\sin x +  {\log }_{2}\cos x = {\log }_{2}\sin x\cos x = {\log }_{2}\left( {\dfrac{1}{2}\sin {2x}}\right)$ ,因为 $\dfrac{1}{2}\sin {2x} \in  \left( {0,\dfrac{1}{2}}\right\rbrack$ ,所以 ${\log }_{2}\left( {\dfrac{1}{2}\sin {2x}}\right)  \in  ( - \infty , - 1\rbrack$ ,则函数的最大值为 $-1$
    \end{jieda}
\end{enumerate}
\subsection*{考点2 诱导公式及其应用}

\subsubsection*{一、填空题}
\begin{enumerate}
    \item (2024$\cdot$闵行区一模$\cdot$2)若 $\sin \alpha  = \dfrac{1}{3}$ ，则 $\sin \left( {\pi  - \alpha }\right)  =$ \blkda[2]{$\dfrac{1}{3}$}.

    \item (2024$\cdot$闵行区二模$\cdot$3)始边与 $x$ 轴的正半轴重合的角 $\alpha$ 的终边过点(3，-4)，则 $\sin \left( {\alpha  + \pi }\right)  =$ \blkda[2]{$\dfrac{4}{5}$}.

    \item (2024$\cdot$青浦区一模$\cdot$3)已知 $\alpha$ 满足 $\cos \alpha  = m$ ，则 $\sin \left( {\alpha  + \dfrac{\pi }{2}}\right)  =$ \blkda[2]{$m$}(结果用含有 $m$ 的式子表示).

    \item (2024$\cdot$普陀区二模$\cdot$3)若 $\cos \left( {\dfrac{\pi }{3} - \alpha }\right)  = \dfrac{3}{5}$ ，则 $\sin \left( {\dfrac{\pi }{6} + \alpha }\right)  =$ \blkda[2]{$\dfrac{3}{5}$}.

    \item (2024$\cdot$嘉定区一模$\cdot$4)已知 $\tan \alpha  = 2$ ，则 $\tan \left( {\alpha  + \dfrac{\pi }{2}}\right)  =$ \blkda[2]{$- \dfrac{1}{2}$}.

    \item (2024$\cdot$金山区一模$\cdot$5)已知角 $\alpha$ 、 $\beta$ 的终边关于原点 $O$ 对称，则 $\cos \left( {\alpha  - \beta }\right)  =$ \blkda[2]{$-1$}.

    \item (2025$\cdot$奉贤区二模$\cdot$6)已知 $\theta$ 是斜率为 $-1$ 的直线的倾斜角，则 $\sin \left( {\theta  - \dfrac{\pi }{2}}\right)  =$ \blkda[2]{$\dfrac{\sqrt{2}}{2}$}.
\end{enumerate}

\subsubsection*{二、选择题}
\begin{enumerate}
    \item (2025$\cdot$青浦区二模$\cdot$13)“函数 $y = \sin \left( {x + \varphi }\right)$ 为偶函数”是“ $\varphi  = \dfrac{\pi }{2}$ ”的 \dotfill{(\kh{B})}
    \xx{充分不必要条件}{必要不充分条件}{充要条件}{既不充分也不必要条件}

    \item (2025$\cdot$普陀区二模$\cdot$14)设 $m \in  \mathbf{R}$ ，在平面直角坐标系 ${xOy}$ 中，角 $\alpha$ 的顶点在坐标原点，始边与 $x$ 轴的正半轴重合,若角 $\alpha$ 的终边经过点 $P\left( {-3,m}\right)$ ,且 $\sin \left( {\pi  + {2\alpha }}\right)  > 0$ ,则角 $\alpha$ 属于 \dotfill{(\kh{B})}
    \xx{第一象限}{第二象限}{第三象限}{第四象限}
    \begin{jieda}
        因为 $\sin \left( {\pi  + {2\alpha }}\right)  =  - \sin {2\alpha } > 0$ ，所以 $\sin {2\alpha } < 0$ ，所以 $2\cos \alpha \sin \alpha  < 0$ ，所以 $\sin \alpha, \cos \alpha$ 异号，所以 $\alpha$ 在第二、四象限， 又 $P\left( {-3,m}\right)$ ,所以 $\alpha$ 在第二象限.
    \end{jieda}
\end{enumerate}
\subsection*{考点3 两角和与差的正弦、余弦、正切公式及二倍角公式}

\subsubsection*{一、填空题}
\begin{enumerate}
    \item (2024$\cdot$虹口区二模$\cdot$1)若 $\sin x =  - \dfrac{3}{5}$ ，则 $\cos {2x} =$ \blkda[2]{$\dfrac{7}{25}$}.

    \item (2025$\cdot$松江区一模$\cdot$2)若 $\sin \alpha  = \dfrac{4}{5}$ ，则 $\cos {2\alpha } =$ \blkda[2]{$- \dfrac{7}{25}$}.

    \item (2025$\cdot$青浦区二模$\cdot$2)函数 $y = 2\sin x - 3\cos x$ 的值域是\blkda[2]{$\left\lbrack  {-\sqrt{13},\sqrt{13}}\right\rbrack$}.

    \item (2024$\cdot$宝山区二模$\cdot$2)已知 $\tan \alpha  = 3$ ，则 $\tan \left( {\alpha  - \dfrac{\pi }{4}}\right)  =$ \blkda[2]{$\dfrac{1}{2}$}.

    \item (2025$\cdot$虹口区一模$\cdot$3)若 $\tan \alpha  = 5$ ，则 $\tan {2\alpha } =$ \blkda[2]{$- \dfrac{5}{12}$}.

    \item (2025$\cdot$虹口区二模$\cdot$3)若 $\sin \alpha  = \dfrac{2}{3}$ ，则 $\cos {2\alpha } =$ \blkda[2]{$\dfrac{1}{9}$}.

    \item (2024$\cdot$杨浦区一模$\cdot$3)若 $\sin \alpha  = \dfrac{3}{5}$ ，则 $\cos {2\alpha } =$ \blkda[2]{$\dfrac{7}{25}$}.

    \item (2023$\cdot$上海秋季卷$\cdot$4)已知 $\tan A = 3$ ，则 $\tan {2A} =$ \blkda[2]{$- \dfrac{3}{4}$}.

    \item (2025$\cdot$浦东新区二模$\cdot$4)若 $f\left( x\right)  = \cos {3x}\cos x + \sin {3x}\sin x$ ，则函数 $y = f\left( x\right)$ 的最小正周期为\blkda[2]{$\pi$}.

    \item (2024$\cdot$杨浦区二模$\cdot$4)已知 $\sin \alpha  = \dfrac{1}{3}$ ，则 $\cos {2\alpha } =$ \blkda[2]{$\dfrac{7}{9}$}.

    \item (2024$\cdot$松江区一模$\cdot$5)已知 $\sin \theta  = \dfrac{3}{5}$ ， $\theta  \in  \left( {0,\dfrac{\pi }{2}}\right)$ ，则 $\tan \left( {\theta  - \dfrac{\pi }{4}}\right)$ 的值为\blkda[2]{$- \dfrac{1}{7}$}.

    \item (2025$\cdot$上海春季卷$\cdot$5)已知 $\tan \alpha  = 1$ ，则 $\cos \left( {\alpha  + \dfrac{\pi }{4}}\right)  =$ \blkda[2]{0}.

    \item (2025$\cdot$嘉定区二模$\cdot$6)已知 $\theta  \in  \mathbf{R}$ ，若 $\tan \theta  + \cot \theta  = 5$ ，则 $\sin {2\theta } =$ \blkda[2]{$\dfrac{2}{5}$}.

    \item (2025$\cdot$黄浦区二模$\cdot$6)函数 $f\left( x\right)  = \sqrt{3}\sin x + \sin \left( {\dfrac{\pi }{2} + x}\right)$ 的最大值是\blkda[2]{2}.

    \item (2025$\cdot$静安区二模$\cdot$6)已知 $\sin \left( {\dfrac{\pi }{4} - x}\right)  = \dfrac{3}{5}$ ，则 $\sin {2x}$ 的值为\blkda[2]{$\dfrac{7}{25}$}.

    \item (2025$\cdot$徐汇区二模$\cdot$6)已知 $\cos \theta  =  - \dfrac{3}{5}$ ， $\theta  \in  \left( {0,\pi }\right)$ ，则 $\tan \left( {\theta  - \dfrac{\pi }{4}}\right)$ 的值为\blkda[2]{7}.

    \item (2024$\cdot$虹口区一模$\cdot$6)已知 $\cos x =  - \dfrac{1}{3}$ ，且 $x$ 为第三象限的角，则 $\tan {2x} =$ \blkda[2]{$- \dfrac{4\sqrt{2}}{7}$}.
\end{enumerate}

\subsubsection*{二、选择题}
\begin{enumerate}
    \item (2024$\cdot$上海秋季卷$\cdot$14)下列函数 $f\left( x\right)$ 的最小正周期是 ${2\pi }$ 的是 \dotfill{(\kh{A})}
    \xx{$\sin x + \cos x$}{$\sin x\cos x$}{${\sin }^{2}x + {\cos }^{2}x$}{${\sin }^{2}x - {\cos }^{2}x$}

    \item (2024$\cdot$黄浦区二模$\cdot$14)函数 $y = 1 - 2{\cos }^{2}\left( {x - \dfrac{\pi }{4}}\right)$ 是 \dotfill{(\kh{A})}
    \xx{最小正周期为 $\pi$ 的奇函数}{最小正周期为 $\pi$ 的偶函数}{最小正周期为 $\dfrac{\pi }{2}$ 的奇函数}{最小正周期为 $\dfrac{\pi }{2}$ 的偶函数}
    \begin{jieda}
        $y = 1 - 2{\cos }^{2}\left( {x - \dfrac{\pi }{4}}\right)  =   - \cos \left( {{2x} - \dfrac{\pi }{2}}\right)  =  - \sin {2x}$ ,因为 $f\left( {-x}\right)  =  - \sin \left( {-{2x}}\right)  =  \sin {2x} =  - f\left( x\right)$ ,所以为奇函数,周期 $T = \dfrac{2\pi }{2} = \pi$ ,所以此函数最小正周期为 $\pi$ 的奇函数
    \end{jieda}

    \item (2024$\cdot$长宁区一模$\cdot$15)设点 $P$ 是以原点为圆心的单位圆上的动点，它从初始位置 ${P}_{0}\left( {1,0}\right)$ 出发，沿单位圆按逆时针方向转动角 $\alpha \left( {0 < \alpha  < \dfrac{\pi }{2}}\right)$ 后到达点 ${P}_{1}$ ，然后继续沿单位圆按逆时针方向转动角 $\dfrac{\pi }{4}$ 到达 ${P}_{2}$. 若点 ${P}_{2}$ 的横坐标为 $- \dfrac{3}{5}$ ，则点 ${P}_{1}$ 的纵坐标 \dotfill{(\kh{D})}
    \xx{$\dfrac{\sqrt{2}}{10}$}{$\dfrac{\sqrt{2}}{5}$}{$\dfrac{3\sqrt{2}}{5}$}{$\dfrac{7\sqrt{2}}{10}$}
    \begin{jieda}
        由题可知 $\cos \left( {\alpha  + \dfrac{\pi }{4}}\right)  =  - \dfrac{3}{5}$ ,且 $0 < \alpha  < \dfrac{\pi }{2}$ ,因为 $\dfrac{\pi }{4} < \alpha  + \dfrac{\pi }{4} < \dfrac{3\pi }{4}$ ,可知 $\dfrac{\pi }{2} < \alpha  + \dfrac{\pi }{4} < \dfrac{3\pi }{4}$ ,则 $\sin \left( {\alpha  + \dfrac{\pi }{4}}\right)  =  \sqrt{1 - {\cos }^{2}\left( {\alpha  + \dfrac{\pi }{4}}\right) } = \dfrac{4}{5}$ ,所以 $\sin \alpha  = \sin \left\lbrack  {\left( {\alpha  + \dfrac{\pi }{4}}\right)  - \dfrac{\pi }{4}}\right\rbrack   =  \sin \left( {\alpha  + \dfrac{\pi }{4}}\right) \cos \dfrac{\pi }{4} - \cos \left( {\alpha  + \dfrac{\pi }{4}}\right) \sin \dfrac{\pi }{4} = \dfrac{\sqrt{2}}{2}\left( {\dfrac{4}{5} + \dfrac{3}{5}}\right)  =  \dfrac{7\sqrt{2}}{10}$.
    \end{jieda}
\end{enumerate}
\subsection*{考点4 简单的三角恒等变换}

\subsubsection*{一、填空题}
\begin{enumerate}
    \item (2024$\cdot$奉贤区二模$\cdot$6)已知 $\alpha  \in  \left\lbrack  {0,\pi }\right\rbrack$ ，且 ${2\cos {2\alpha } - 3\cos \alpha } = 5$ ，则$\alpha = $\blkda[2]{$\pi$}.
    \begin{jieda}
        已知 $2\cos {2\alpha } - 3\cos \alpha  = 5$ ,由倍角公式得 $4{\cos }^{2}\alpha  -  3\cos \alpha  - 7 = \left( {4\cos \alpha  - 7}\right) \left( {\cos \alpha  + 1}\right)  = 0$ ,由 $\alpha  \in  \left\lbrack  {0,\pi }\right\rbrack  ,\cos \alpha  \in  \left\lbrack  {-1,1}\right\rbrack$ ,解得 $\cos \alpha  =  - 1$ ,则 $\alpha  = \pi$.
    \end{jieda}
\end{enumerate}

\subsubsection*{二、选择题}
\begin{enumerate}
    \item (2025$\cdot$黄浦区一模$\cdot$15)设$0 \leq  x < {2\pi }$ ，满足 $\sin \left( {x + \dfrac{\pi }{6}}\right)  = \sin x + \sin \dfrac{\pi }{6}$ 的 $x$ 的个数为\dotfill{(\kh{C})}
    \xx{0 个}{1 个}{2 个}{无数个}
    \begin{jieda}
        \begin{dingautolist}{172}
            \item $\sin \left( {x + \dfrac{\pi }{6}}\right)  = \sin x + \sin \dfrac{\pi }{6} \Leftrightarrow 2sin\left(\dfrac{x}{2}+\dfrac{\pi}{12}\right)cos\left(\dfrac{x}{2}+\dfrac{\pi}{12}\right) = 2sin\left(\dfrac{x}{2}+\dfrac{\pi}{12}\right)cos\left(\dfrac{x}{2}-\dfrac{\pi}{12}\right)$

            $sin\left(\dfrac{x}{2}+\dfrac{\pi}{12}\right) = 0$或$cos\left(\dfrac{x}{2}+\dfrac{\pi}{12}\right) = cos\left(\dfrac{x}{2}-\dfrac{\pi}{12}\right)$，$x = \dfrac{11\pi}{6}$或$0$.
            \item 设$f(x) = \sin \left( {x + \dfrac{\pi }{6}}\right) - \sin x - \sin \dfrac{\pi }{6}$
            
            $f'(x) = cos\left(x+\dfrac{\pi}{6}\right) - cosx$，当$x \in \left[0, \dfrac{11\pi}{12}\right), f(x)$严减，当$x \in \left[\dfrac{11\pi}{12}, \dfrac{23\pi}{12}\right), f(x)$严增，当$x \in \left[\dfrac{23\pi}{12}, 2\pi\right), f(x)$严减

            $f(0) = 0, f\left(\dfrac{11\pi}{12}\right) < 0, f\left(\dfrac{23\pi}{12}\right) > 0, f(2\pi) = 0$
        \end{dingautolist}
    \end{jieda}
    \item (2025$\cdot$闵行区一模$\cdot$15)设 $f\left( x\right)  = \left( {\sin x - \cos x}\right) \left( {\cos x - \tan x}\right) \left( {\tan x - \sin x}\right)$ ，若 $\alpha$ 、 $\beta$ 为同一象限的角,且不存在 $\alpha \text{ 、 }\beta$ ,使得 $f\left( \alpha \right) f\left( \beta \right)  < 0$ ,则 $\alpha \text{ 、 }\beta$ 所在的象限为 \dotfill{(\kh{D})}
    \xx{第一象限}{第二象限}{第三象限}{第四象限}
    \begin{jieda}
        第四象限内三角比的大小关系不改变.
    \end{jieda}
\end{enumerate}

\subsection*{考点5 解三角形}

\subsubsection*{一、填空题}
\begin{enumerate}
    \item (2024$\cdot$徐汇区二模$\cdot$3)在 $\bigtriangleup  {ABC}$ 中， ${AC} = 1$ ， $\angle C = \dfrac{2\pi }{3}$ ， $\angle A = \dfrac{\pi }{6}$ ，则 $\bigtriangleup  {ABC}$ 的外接圆半径为\blkda[2]{1}.
    \begin{jieda}
        由已知 $\angle B = \dfrac{\pi }{6}$ ,设三角形外接圆半径为 $R$ ,则 ${2R} =  \dfrac{AC}{\sin B} = \dfrac{1}{\sin \dfrac{\pi }{6}} = 2$ ,所以 $R = 1$.
    \end{jieda}

    \item (2025$\cdot$黄浦区二模$\cdot$4)在 $\bigtriangleup  {ABC}$ 中，若 $A = {45}^{ \circ  }$ ， $B = {30}^{ \circ  }$ ， $a = 2\sqrt{6}$ ，则 $b =$ \blkda[2]{$2\sqrt{3}$}.
    \begin{jieda}
        由正弦定理得 $\dfrac{a}{\sin A} = \dfrac{b}{\sin B}$ ,故 $\dfrac{2\sqrt{6}}{\sin {45}^{ \circ  }} = \dfrac{b}{\sin {30}^{ \circ  }}$ ,即 $\dfrac{2\sqrt{6}}{\dfrac{\sqrt{2}}{2}} = \dfrac{b}{\dfrac{1}{2}}$ ,解得 $b = 2\sqrt{3}$.
    \end{jieda}

    \item (2025$\cdot$嘉定区一模$\cdot$4)在 $\bigtriangleup  {ABC}$ 中，若 ${AB} = 5,{BC} = \sqrt{21},{CA} = 4$ ，则 $\angle A =$ \blkda[2]{$\dfrac{\pi }{3}$}.
    \begin{jieda}
        由余弦定理知 $\cos A = \dfrac{{b}^{2} + {c}^{2} - {a}^{2}}{2bc} = \dfrac{{16} + {25} - {21}}{2 \times  4 \times  5} =  \dfrac{1}{2}$ ,因为 $A \in  \left( {0,\pi }\right)$ ,所以 $A = \dfrac{\pi }{3}$.
    \end{jieda}

    \item (2024$\cdot$长宁区二模$\cdot$4)在 $\bigtriangleup  {ABC}$ 中，内角 $A$ 、 $B$ 、 $C$ 所对的边分别为 $a$ 、 $b$ 、 $c$ ，若 ${a}^{2} = {b}^{2} + {bc} + {c}^{2}$ ，则 $A =$ \blkda[2]{${120}^{ \circ }$}.
    \begin{jieda}
        $\cos A = \dfrac{{b}^{2} + {c}^{2} - {a}^{2}}{2bc} = \dfrac{-{bc}}{2bc} =  - \dfrac{1}{2}$ ,所以 $A = {120}^{ \circ  }$.
    \end{jieda}

    \item (2025$\cdot$松江区一模$\cdot$4)在 $\bigtriangleup  {ABC}$ 中，设角 $A$ 、 $B$ 及 $C$ 所对边的边长分别为 $a$ 、 $b$ 及 $c$ ，若 $a = 4,b =  \sqrt{3},C = \dfrac{5}{6}\pi$ ,则边长 $c =$ \blkda[2]{$\sqrt{31}$}.
    \begin{jieda}
        ${c}^{2} = {a}^{2} + {b}^{2} - {2ab}\cos C = {16} + 3 - 2 \times  4 \times  \sqrt{3} \times  \left( {-\dfrac{\sqrt{3}}{2}}\right)  = {31} \Longrightarrow  c = \sqrt{31}.$
    \end{jieda}

    \item (2025$\cdot$浦东新区一模$\cdot$5)在 $\bigtriangleup  {ABC}$ 中， ${BC} = 5$ ， $\angle B = {45}^{ \circ  }$ ， $\angle C = {105}^{ \circ  }$ ，则 ${AC} =$ \blkda[2]{$5\sqrt{2}$}.
    \begin{jieda}
        $\angle A =  {30}^{ \circ  }$ ,由正弦定理得, ${AC} = \dfrac{{BC}\sin B}{\sin A} = \dfrac{5 \times  \dfrac{\sqrt{2}}{2}}{\dfrac{1}{2}} = 5\sqrt{2}$.
    \end{jieda}

    \item (2024$\cdot$上海春季卷$\cdot$5)在 $\bigtriangleup  {ABC}$ 中，若 ${BC} = 2,\angle A = {45}^{ \circ  },\angle B = {30}^{ \circ  }$ ，则 ${AC} =$ \blkda[2]{$\sqrt{2}$}.
    \begin{jieda}
        在 $\bigtriangleup {ABC}$ 中,由正弦定理 $\dfrac{AC}{\sin B} = \dfrac{BC}{\sin A} \Longrightarrow  {AC} =  \dfrac{{BC} \times  \sin B}{\sin A} = \dfrac{2 \times  \dfrac{1}{2}}{\dfrac{\sqrt{2}}{2}} = \sqrt{2}.$
    \end{jieda}

    \item (2025$\cdot$青浦区二模$\cdot$6)已知 $\bigtriangleup  {ABC}$ 的角 $A$ 、 $B$ 、 $C$ 对应边长分别为 $a = 4,b = 5,c = 6$ ，则 $A =$ \blkda[2]{$\arccos \dfrac{3}{4}$}.
    \begin{jieda}
        根据余弦定理得 $\cos A = \dfrac{{b}^{2} + {c}^{2} - {a}^{2}}{2bc}$ ,把 $a = 4,b =$ 5, $c = 6$ 代入可得 $\cos A = \dfrac{{5}^{2} + {6}^{2} - {4}^{2}}{2 \times  5 \times  6} = \dfrac{3}{4}$ ,因为 $A \in  \left( {0,\pi }\right)$ , 所以 $A = \arccos \dfrac{3}{4}$.
    \end{jieda}

    \item (2025$\cdot$普陀区一模$\cdot$6)设 $\bigtriangleup  {ABC}$ 的内角 $A$ 、 $B$ 、 $C$ 的对边分别为 $a$ 、 $b$ 、 $c$ ，若 $b = 4$ ， $\sin \left( {A + \dfrac{\pi }{3}}\right)  = 0$ ， $\bigtriangleup {ABC}$ 的面积为 $\sqrt{3}$ ，则 $a$ 的值为\blkda[2]{$\sqrt{21}$}.
    \begin{jieda}
        因为 $\sin \left( {A + \dfrac{\pi }{3}}\right)  = 0$ ,且 $A \in  \left( {0,\pi }\right)$ ,所以 $A = \dfrac{2\pi }{3}$. 
        
        由 ${S}_{\bigtriangleup {ABC}} = \dfrac{1}{2}{bc}\sin A = \sqrt{3}$ ,得 $\dfrac{1}{2} \times  {4c} \times  \dfrac{\sqrt{3}}{2} = \sqrt{3}$ ,解得 $c = 1$. 
        
        根据余弦定理,可得 ${a}^{2} = {b}^{2} + {c}^{2} - {2bc}\cos \dfrac{2\pi }{3} = {16} + 1 - 2 \times  4 \times  1 \times  \left( {-\dfrac{1}{2}}\right)  = {21}$ ,所以 $a = \sqrt{21}$.
    \end{jieda}

    \item (2025$\cdot$静安区一模$\cdot$6)在 $\bigtriangleup {ABC}$ 中，已知 ${BC} = 5$ ， ${AC} = 4$ ， $A = {2B}$ ，则 $\cos B$ 的值为\blkda[2]{$\dfrac{5}{8}$}.
    \begin{jieda}
        在 $\bigtriangleup {ABC}$ 中,由正弦定理得 $\dfrac{BC}{\sin A} = \dfrac{AC}{\sin B}$ ,而 ${BC} =  5,{AC} = 4,A = {2B}$ ,因此 $\dfrac{5}{\sin {2B}} = \dfrac{4}{\sin B}$ ,即 $\dfrac{5}{2\sin B\cos B} = \dfrac{4}{\sin B}$ , 所以 $\cos B = \dfrac{5}{8}$.
    \end{jieda}

    \item (2024$\cdot$黄浦区二模$\cdot$6)在 $\bigtriangleup  {ABC}$ 中， $\cos A =  - \dfrac{3}{5}$ ， ${AB} = 1$ ， ${AC} = 5$ ，则 ${BC} =$ \blkda[2]{$4\sqrt{2}$}.
    \begin{jieda}
        在 $\bigtriangleup {ABC}$ 中,根据余弦定理可得: $\cos A =  \dfrac{A{B}^{2} + A{C}^{2} - B{C}^{2}}{2 \cdot  {AB} \cdot  {AC}}$ ，设 ${BC} = x\left( {x > 0}\right)$ ，则 $- \dfrac{3}{5} = \dfrac{1 + {25} - {x}^{2}}{2 \times  1 \times  5}$ ，整理可得 ${x}^{2} = {32}$ ,解得 $x = 4\sqrt{2}$ ,故 ${BC} = 4\sqrt{2}$.
    \end{jieda}

    \item (2025$\cdot$奉贤区二模$\cdot$7)已知 $\bigtriangleup  {ABC}$ ，“ $A$ 、 $B$ 、 $C$ 成等差数列且 $\sin A$ 、 $\sin B$ 、 $\sin C$ 成等比数列”是“$\bigtriangleup  {ABC}$ 是正三角形”的\blkda[2]{充要}条件.
    \begin{jieda}
        在 $\bigtriangleup  {ABC}$ 中，由 $A$ 、 $B$ 、 $C$ 成等差数列，得 ${2B} = A +  C$ ,而 $A + B + C = \pi$ ,则 $B = \dfrac{\pi }{3}$ ,由 $\sin A\text{ 、 }\sin B\text{ 、 }\sin C$ 成等比数列,得 ${\sin }^{2}B = \sin A\sin C$ ,由正弦定理得 ${b}^{2} = {ac}$ ,由余弦定理得 ${b}^{2} = {a}^{2} + {c}^{2} - {2ac}\cos B$ ,即 ${ac} = {a}^{2} + {c}^{2} - {ac}$ ,解得 $a = c$ , 因此 $\bigtriangleup {ABC}$ 是正三角形
        
        若 $\bigtriangleup {ABC}$ 是正三角形,则 $A = B =  C = \dfrac{\pi }{3},\sin B = \sin A = \sin C = \dfrac{\sqrt{3}}{2}$ ,因此 $A\text{ 、 }B\text{ 、 }C$ 成等差数列且 $\sin A\text{ 、 }\sin B\text{ 、 }\sin C$ 成等比数列,所以 “ $A\text{ 、 }B\text{ 、 }C$ 成等差数列且 $\sin A\text{ 、 }\sin B\text{ 、 }\sin C$ 成等比数列” 是 “ $\bigtriangleup  {ABC}$ 是正三角形” 的充要条件.
    \end{jieda}

    \item (2025$\cdot$青浦区一模$\cdot$7)在 $\bigtriangleup  {ABC}$ 中，已知 $\angle {ACB} = {120}^{ \circ  }$ ， ${AB} = {2\sqrt{7}}$ ，若 ${BC} = {2AC}$ ，则 $\bigtriangleup  {ABC}$ 的面积为\blkda[2]{$2\sqrt{3}$}.
    \begin{jieda}
        在 $\bigtriangleup {ABC}$ 中, $\angle {ACB} = {120}^{ \circ  },{AB} = 2\sqrt{7},{BC} =  {2AC}$ ,由余弦定理得 ${28} = A{B}^{2} = A{C}^{2} + B{C}^{2} - {2AC} \cdot {BC}\cos {120}^{ \circ  } = {7A}{C}^{2}$ ，解得 $A{C}^{2} = 4$ ，所以 $\bigtriangleup  {ABC}$ 的面积为 $\dfrac{1}{2}{AC} \times  {BC} \times  \sin {120}^{ \circ  } = A{C}^{2} \cdot  \dfrac{\sqrt{3}}{2} = 2\sqrt{3}$.
    \end{jieda}

    \item (2025$\cdot$虹口区二模$\cdot$7)若 $\bigtriangleup  {ABC}$ 的三条边的长分别为4、5、6，则 $\bigtriangleup  {ABC}$ 的外接圆面积为\blkda[2]{$\dfrac{64\pi }{7}$}. (结果保留 $\pi$ )
    \begin{jieda}
        不妨设 $a = 4,b = 5,c = 6$ ,由余弦定理可得 $\cos C =  \dfrac{{a}^{2} + {b}^{2} - {c}^{2}}{2ab} = \dfrac{{16} + {25} - {36}}{2 \times  4 \times  5} = \dfrac{1}{8}$，故 $\sin C =  \sqrt{1 - {\cos }^{2}C} = \sqrt{1 - {\left( \dfrac{1}{8}\right) }^{2}} = \dfrac{3\sqrt{7}}{8}$ ,设 $\bigtriangleup {ABC}$ 的外接圆半径为 $R$ ,则 ${2R} = \dfrac{c}{\sin C} = \dfrac{6}{\dfrac{3\sqrt{7}}{8}} = \dfrac{16}{\sqrt{7}}$ ,所以, $R = \dfrac{8}{\sqrt{7}}$ ,因此, $\bigtriangleup {ABC}$ 的外接圆的面积为 $\pi R^{2} = \pi  \times  {\left( \dfrac{8}{\sqrt{7}}\right) }^{2} = \dfrac{64\pi }{7}$.
    \end{jieda}

    \item (2024$\cdot$虹口区二模$\cdot$7)已知一个三角形的三边长分别为2、3、4，则这个三角形外接圆的直径为\blkda[2]{$\dfrac{16}{15}\sqrt{15}$}.
    \begin{jieda}
        不妨设 $a = 2,b = 3,c = 4$ ,由余弦定理可得 $\cos C =  \dfrac{{a}^{2} + {b}^{2} - {c}^{2}}{2ab} = \dfrac{4 + 9 - 16}{2 \times  2 \times  3} = -\dfrac{1}{4}$，故 $\sin C =  \sqrt{1 - {\cos }^{2}C} = \sqrt{1 - {\left(-\dfrac{1}{4}\right) }^{2}} = \dfrac{\sqrt{15}}{4}$ ,设 $\bigtriangleup {ABC}$ 的外接圆半径为 $R$ ,则 ${2R} = \dfrac{c}{\sin C} = \dfrac{4}{\dfrac{\sqrt{15}}{4}} = \dfrac{16}{15}\sqrt{15}$.
    \end{jieda}

    \item (2025$\cdot$崇明区二模$\cdot$8)在 $\bigtriangleup  {ABC}$ 中，若 $c = 3$ ， $C = \dfrac{\pi }{3}$ ，其面积为 $\sqrt{3}$ ，则 $a + b =$ \blkda[2]{$\sqrt{21}$}.
    \begin{jieda}
        已知 $C = \dfrac{\pi }{3},{S}_{\bigtriangleup {ABC}} = \sqrt{3}$ ,代入面积公式可得: $\sqrt{3} =  \dfrac{1}{2}{ab}\sin \dfrac{\pi }{3}$ ，则 $\sqrt{3} = \dfrac{1}{2}{ab} \times  \dfrac{\sqrt{3}}{2}$ ，可得: ${ab} = 4$ . 根据余弦定理 ${c}^{2} = {a}^{2} + {b}^{2} - {2ab}\cos C$ ,可得 ${3}^{2} = {a}^{2} + {b}^{2} - {2ab}\cos \dfrac{\pi }{3}$ ,则 $9 =  {a}^{2} + {b}^{2} - {ab}$ . 即 $9 = {\left( a + b\right) }^{2} - {2ab} - {ab} = {\left( a + b\right) }^{2} - {3ab}$ ,把 ${ab} = 4$ 代入可得: $9 = {\left( a + b\right) }^{2} - 3 \times  4$ ,即 ${\left( a + b\right) }^{2} = 9 + {12} =$ 21. 由于 $a\text{ 、 }b$ 为边长,可得 $a + b = \sqrt{21}$.
    \end{jieda}

    \item (2023$\cdot$上海秋季卷$\cdot$8)已知 $\bigtriangleup  {ABC}$ 的角 $A$ 、 $B$ 、 $C$ 对应边长分别为 $a$ 、 $b$ 、 $c$ ， $a = 4$ ， $b = 5$ ， $c = 6$ ，则 $\sin A =$ \blkda[2]{$\dfrac{\sqrt{7}}{4}$}.
    \begin{jieda}
        $\cos A = \dfrac{{b}^{2} + {c}^{2} - {a}^{2}}{2bc} = \dfrac{{25} + {36} - {16}}{2 \times  5 \times  6} = \dfrac{3}{4}$ ,所以 $\sin A  = \dfrac{\sqrt{7}}{4}$.
    \end{jieda}

    \item (2024$\cdot$黄浦区一模$\cdot$8)在 $\bigtriangleup  {ABC}$ 中，三个内角 $A$ 、 $B$ 、 $C$ 的对边长分别为 $a$ 、 $b$ 、 $c$ ，若 $5{a}^{2} - 5{b}^{2} +  {6bc} - 5{c}^{2} = 0$ ，则 $\sin {2A}$ 的值为\blkda[2]{$\dfrac{24}{25}$}.
    \begin{jieda}
        $\left\{  {\begin{array}{l} 5{a}^{2} - 5{b}^{2} + {6bc} - 5{c}^{2} = 0 \Longrightarrow  {b}^{2} + {c}^{2} - {a}^{2} = \dfrac{6}{5}{bc} \\  \cos A = \dfrac{{b}^{2} + {c}^{2} - {a}^{2}}{2bc} \end{array} \Longrightarrow  }\right.  \cos A = \dfrac{3}{5},\sin A = \dfrac{4}{5},\sin {2A} = 2\sin A\cos A = \dfrac{24}{25}.$
    \end{jieda}

    \item (2024$\cdot$松江区一模$\cdot$9)在 $\bigtriangleup  {ABC}$ 中，设角 $A$ 、 $B$ 及 $C$ 所对边的边长分别为 $a$ 、 $b$ 及 $c$ ，若 $a = 3$ ， $c =$ 5， $B = {2A}$ ，则边长 $b =$ \blkda[2]{$2\sqrt{6}$}.
    \begin{jieda}
        $\dfrac{a}{\sin A} = \dfrac{b}{\sin B} = {2R} \Longrightarrow  \dfrac{a}{\sin A} = \dfrac{b}{\sin B} = \dfrac{b}{\sin {2A}} \Longrightarrow  b =  6\cos A;\cos B = \cos {2A} \Longrightarrow  \dfrac{{a}^{2} + {c}^{2} - {b}^{2}}{2ac} = 2{\cos }^{2}A - 1$ 代入得 $\dfrac{9 + {25} - {36}{\cos }^{2}A}{30} = 2{\cos }^{2}A - 1 \Longrightarrow  \cos A = \dfrac{\sqrt{6}}{3} \Longrightarrow  b = 2\sqrt{6}.$
    \end{jieda}

    \item (2024$\cdot$上海秋季卷$\cdot$11)已知点 $B$ 在点 $C$ 正北方向，点 $D$ 在点 $C$ 的正东方向， ${BC} = {CD}$ ，存在点$A$满足 $\angle {BAC} = {16.5}^{ \circ  }$ ， $\angle {DAC} = {37}^{ \circ  }$ ，则 $\angle {BCA} =$ \blkda[2]{${7.8}^{ \circ  }$}.(精确到 ${0.1}^{ \circ  }$)
    \begin{jieda}
        设 $\angle {BCA} = \theta$ ，则 $\angle {ACD} = {90}^{\circ} - \theta$，$\angle ADC = \ang{53} + \theta$，$\angle ABC = \ang{163.5} - \theta$
        
        由正弦定理，在 $\bigtriangleup  {DCA}$ 中，$\dfrac{CD}{sin\ang{37}} = \dfrac{AC}{sin(\ang{53} + \theta)}$，在 $\bigtriangleup {BCA}$ 中，$\dfrac{BC}{sin\ang{16.5}} = \dfrac{AC}{sin(\ang{163.5} - \theta)}$

        故$\dfrac{sin(\ang{53} + \theta)}{cos\ang{53}} = \dfrac{sin(\ang{16.5} + \theta)}{sin\ang{16.5}}$，即$tan\ang{53}cos\theta + sin\theta = cos\theta + cot\ang{16.5}sin\theta$
        
        故$tan\theta = \dfrac{tan\ang{53}-1}{cot\ang{16.5}-1}$，$\theta = arctan\left(\dfrac{tan\ang{53}-1}{cot\ang{16.5}-1}\right) \approx \ang{7.8}$

        \ps 题目条件集中在角上，故先标出角，利用正弦定理建立等量关系
        \begin{center}
            \includegraphics[width=0.15\textwidth]{bodyxb/src/35_1339_1865_192_231_0.jpg}
        \end{center}
    \end{jieda}

    \item (2025$\cdot$静安区二模$\cdot$12)在边长为 1 的正三角形 ${ABC}$ 的边 ${AB}$ 、 ${AC}$ 上分别取 $D$ 、 $E$ 两点，若沿线段 ${DE}$ 折叠该三角形时,顶点 $A$ 恰好落在边 ${BC}$ 上. 则线段 ${AD}$ 的长度的最小值为\blkda[2]{$2\sqrt{3} -3$}.
    \begin{jieda}
        以$BC$中点为原点建系，$A\left(0, \dfrac{\sqrt{3}}{2}\right)$，设$A$翻折之后为$A'(t, 0), t \in \left[-\dfrac{1}{2}, \dfrac{1}{2}\right]$，根据几何关系易知$AA'$中垂线必然与线段${AB}$ 、 ${AC}$相交，即$DE$为$AA'$中垂线.
        
        $\vv{AA'} = \left(-t, \dfrac{\sqrt{3}}{2}\right)$，故设$AA': -tx+\dfrac{\sqrt{3}}{2}y+c = 0$，代入$AA'$中点坐标$\left(\dfrac{t}{2}, \dfrac{\sqrt{3}}{4}\right)$，可得$c = \dfrac{4t^2-3}{8}$，因此$DE: -tx+\dfrac{\sqrt{3}}{2}y+\dfrac{4t^2-3}{8} = 0$，直线$AB: y = \sqrt{3}x + \dfrac{\sqrt{3}}{2}$，两者交点为$D$.

        $\dfrac{3-2t}{2}x_D+\dfrac{4t^2+3}{8} = 0$，即$x_D = \dfrac{4t^2+3}{4(2t-3)} \in \left[-\dfrac{1}{2}, \dfrac{3}{2} - \sqrt{3}\right]$，故$BD \leq 4 - 2\sqrt{3}$，$AD \geq 2\sqrt{3} -3$.
    \end{jieda}
\end{enumerate}

\subsubsection*{二、选择题}
\begin{enumerate}
    \item (2025$\cdot$杨浦区二模$\cdot$13) $\bigtriangleup  {ABC}$ 中，“$sinA = sinB$”是 “$A = B$” 的 \dotfill{(\kh{C})}
    \xx{充分非必要条件}{必要非充分条件}{充要条件}{既非充分也非必要条件}
    \item (2024$\cdot$普陀区一模$\cdot$14)在 $\bigtriangleup  {ABC}$ 中，角 $A$ 、 $B$ 、 $C$ 所对的边分别为 $a$ 、 $b$ 、 $c$ ，若 $a = \sqrt{3}$ ，且 $c - {2b} +  2\sqrt{3}\cos C = 0$ ,则该三角形外接圆的半径为 \dotfill{(\kh{A})}
    \xx{1}{$\sqrt{3}$}{2}{$2\sqrt{3}$}
    \begin{jieda}
        因为 $a = \sqrt{3},c - {2b} + {2a}\cos C = 0$ ,所以 $\sin C - 2\sin B +  2\sin A\cos C = 0$ . 所以 $\sin C - 2\sin \left( {A + C}\right)  + 2\sin A\cos C = 0$ ,所以 $\sin C - 2\sin A\cos C - 2\sin C\cos A + 2\sin A\cos C = 0$ ,所以 $\sin C - 2\sin C\cos A = 0$ ,因为 $\sin C > 0$ ,所以 $\cos A = \dfrac{1}{2},A \in  \left( {0,\pi }\right)$ ,所以 $A = \dfrac{\pi }{3}$ ,设该三角形外接圆的半径为 $r$ ,由正弦定理得 $\dfrac{a}{\sin A} = \dfrac{\sqrt{3}}{\dfrac{\sqrt{3}}{2}} = 2 = {2r}$ ,所以 $r = 1$
    \end{jieda}
\end{enumerate}


\subsubsection*{三、解答题}
\begin{enumerate}
    \item (2024$\cdot$宝山区二模$\cdot$17)在 $\bigtriangleup  {ABC}$ 中，角 $A$ 、 $B$ 、 $C$ 的对边分别为 $a$ 、 $b$ 、 $c$ ，已知 ${\sin }^{2}A + {\sin }^{2}C =  {\sin }^{2}B + \sin A\sin C$.
    \begin{enumerate}
        \item 求角 $B$ 的大小
        \item 若 $\bigtriangleup  {ABC}$ 的面积为 $\sqrt{3}$ ，求 $a + c$ 的最小值，并判断此时 $\bigtriangleup  {ABC}$ 的形状.
    \end{enumerate}
    \begin{jieda}
        \begin{enumerate}
            \item 由正弦定理得 ${a}^{2} + {c}^{2} = {b}^{2} + {ac}$ ，又由余弦定理得 $\cos B  = \dfrac{{a}^{2} + {c}^{2} - {b}^{2}}{2ac} = \dfrac{ac}{2ac} = \dfrac{1}{2}$ ,因为 $B$ 是三角形内角,所以 $B = \dfrac{\pi }{3}$ 
            \item 由三角形面积公式 ${S}_{\bigtriangleup {ABC}} = \dfrac{1}{2}{ac}\sin B = \dfrac{1}{2}{ac}\sin \dfrac{\pi }{3} =  \dfrac{\sqrt{3}}{4}{ac} = \sqrt{3}$ ，得 ${ac} = 4$ ，因为 $a + c \geq  2\sqrt{ac} = 4$ ，当且仅当 $a =  c = 2$ 时取等号,所以 $a + c$ 的最小值为 4,此时 $\bigtriangleup {ABC}$ 为等边三角形.
        \end{enumerate}
    \end{jieda}
    \item (2024$\cdot$静安区二模$\cdot$17)在 $\bigtriangleup  {ABC}$ 中，角 $A$ 、 $B$ 、 $C$ 的对边分别为 $a$ 、 $b$ 、 $c$ ，已知 $a = 3$ ， $b = 5$ ， $c = 7$.
    \begin{enumerate}
        \item 求角 $C$ 的大小；
        \item 求 $\sin \left( {A + C}\right)$ 的值.
    \end{enumerate}
    \begin{jieda}
        \begin{enumerate}
            \item 由余弦定理，有 $\cos C = \dfrac{{a}^{2} + {b}^{2} - {c}^{2}}{2ab} =  - \dfrac{1}{2}$ ，所以 $C = \dfrac{2\pi }{3}$
            \item 由正弦定理，有 $\dfrac{b}{\sin B} = \dfrac{c}{\sin C}$ ，即 $\sin B = \dfrac{b\sin C}{c} =  \dfrac{5\sqrt{3}}{14}$ . 所以 $\sin \left( {A + C}\right)  = \sin \left( {\pi  - B}\right)  = \sin B = \dfrac{5\sqrt{3}}{14}$.
        \end{enumerate}
    \end{jieda}

    \item (2024$\cdot$奉贤区一模$\cdot$17)在 $\bigtriangleup  {ABC}$ 中，设角 $A$ 、 $B$ 及 $C$ 所对边的边长分别为 $a$ 、 $b$ 及 $c$. 已知、 $\sqrt{3}c =  \sqrt{3}b\cos A + a\sin B$.
    \begin{enumerate}
        \item 求角 $B$ 的大小
        \item 当 $a = 2\sqrt{2}$ 、 $b = 2\sqrt{3}$ 时，求边长 $c$ 和 $\bigtriangleup {ABC}$ 的面积 $S$.
    \end{enumerate}
    \begin{jieda}
        \begin{enumerate}
            \item 由正弦定理得 $\sqrt{3}\sin C = \sqrt{3}\sin B \cdot  \cos A + \sin A$ - $\sin B$ ，由于 $C = \pi  - \left( {A + B}\right)$ ，得 $\sqrt{3}\sin \left( {A + B}\right)  = \sqrt{3}\sin B$ . $\cos A + \sin A \cdot  \sin B$ ,展开得 $\sqrt{3}\sin A \cdot  \cos B + \sqrt{3}\cos A \cdot  \sin B =  \sqrt{3}\sin B \cdot  \cos A + \sin A \cdot  \sin B$ ,化简得 $\sqrt{3}\cos B = \sin B$ ,则 $\tan B = \sqrt{3}$ ，所以 $B = \dfrac{\pi }{3}$ 
            
            \item 由正弦定理，得 $\dfrac{2\sqrt{3}}{\sin \dfrac{\pi }{3}} = \dfrac{2\sqrt{2}}{\sin A} = \dfrac{c}{\sin C}$ , $\dfrac{2\sqrt{3}}{\sin {60}^{ \circ  }} = \dfrac{2\sqrt{2}}{\sin A} =  \dfrac{c}{\sin C},\sin A = \dfrac{\sqrt{2}}{2}$ ,因为 $a < b$ ,所以 $A$ 是锐角,即 $A = \dfrac{\pi }{4}$ ,因为 $A + C = \dfrac{2\pi }{3}$ ,所以 $C = \dfrac{5\pi }{12},c = \dfrac{2\sqrt{3}}{\sin \dfrac{\pi }{3}} \times  \sin C = \sqrt{6} + \sqrt{2}$ ,所以 ${S}_{\bigtriangleup {ABC}} = \dfrac{1}{2}{ab}\sin C = \dfrac{1}{2} \times  2\sqrt{3} \times  2\sqrt{2} \times  \sin \dfrac{5\pi }{12} = 3 + \sqrt{3}.$
        \end{enumerate}
    \end{jieda}

    \item (2024$\cdot$松江区二模$\cdot$17)设 $f\left( x\right)  = {\sin }^{2}\dfrac{\omega }{2}x + \sqrt{3}\cos \dfrac{\omega }{2}x\sin \dfrac{\omega }{2}x\left( {\omega  > 0}\right)$ ，函数 $y = f\left( x\right)$ 图像的两条相邻对称轴之间的距离为 $\pi$.
    \begin{enumerate}
        \item 求函数 $y = f\left( x\right)$ 的解析式
        \item 在 $\bigtriangleup  {ABC}$ 中，设角 $A$ 、 $B$ 及 $C$ 所对边的边长分别为 $a$ 、 $b$ 及 $c$ ，若 $a = \sqrt{3},b = \sqrt{2},f\left( A\right)  = \dfrac{3}{2}$ ，求角 $C$.
    \end{enumerate}
    \begin{jieda}
        \begin{enumerate}
            \item $f\left( x\right)  = {\sin }^{2}\dfrac{\omega }{2}x + \sqrt{3}\cos \dfrac{\omega }{2}x\sin \dfrac{\omega }{2}x = \dfrac{1 - \cos {\omega x}}{2} +  \dfrac{\sqrt{3}}{2}\sin {\omega x} = \sin \left( {{\omega x} - \dfrac{\pi }{6}}\right)  + \dfrac{1}{2}$ . 因为函数 $y = f\left( x\right)$ 图像的相邻两条对称轴之间的距离为 $\pi$ ,所以 $T = {2\pi }$ ,即 $T = \dfrac{2\pi }{\omega } = {2\pi }$ . 所以 $\omega  = 1$ . 所以 $f\left( x\right)  = \sin \left( {x - \dfrac{\pi }{6}}\right)  + \dfrac{1}{2}$
            \item 由 $f\left( A\right)  = \dfrac{3}{2}$ ，得 $\sin \left( {A - \dfrac{\pi }{6}}\right)  + \dfrac{1}{2} = \dfrac{3}{2}$ ，所以 $\sin (A -  \left. \dfrac{\pi }{6}\right)  = 1$ . 因为 $A \in  \left( {0,\pi }\right)$ ,所以 $A = \dfrac{2\pi }{3}$ ,由 $\dfrac{a}{\sin A} = \dfrac{b}{\sin B}$ ,得 $\dfrac{\sqrt{3}}{\dfrac{\sqrt{3}}{2}} =  \dfrac{\sqrt{2}}{\sin B}$ ,化简得 $\sin B = \dfrac{\sqrt{2}}{2}$ ,所以角 $B = \dfrac{\pi }{4}$ . 所以角 $C = \pi  - \dfrac{2\pi }{3}  - \dfrac{\pi }{4} = \dfrac{\pi }{12}.$
        \end{enumerate}
    \end{jieda}
    
    \item (2024$\cdot$虹口区一模$\cdot$17)设 $\bigtriangleup  {ABC}$ 的内角 $A$ 、 $B$ 、 $C$ 所对的边分别为 $a$ 、 $b$ 、 $c$ ，若 $\vv{m} = (\sin A + \sin B - \sin C,\sin A),\vv{n} = \left( {c,b + c - a}\right)$ ,且 $\vv{m}//\vv{n}$.
    \begin{enumerate}
        \item 求角 $B$ 的大小
        \item 若 $\bigtriangleup  {ABC}$ 为锐角三角形，求 $y = \sin A + \sin C$ 的取值范围.
    \end{enumerate}
    \begin{jieda}
        \begin{enumerate}
            \item 因为 $\vv{m}//\vv{n}$ ，所以(sinA + sinB - sinC)·(b + c - a)= $c \cdot  \sin A$ ,由正弦定理,可得 $\left( {a + b - c}\right)  \cdot  \left( {b + c - a}\right)  = {ac}$ ,即 ${ac} = {a}^{2} + {c}^{2} - {b}^{2}$ ,于是,由余弦定理,得 $\cos B = \dfrac{{a}^{2} + {c}^{2} - {b}^{2}}{2ac} =  \dfrac{1}{2}$ ,又 $B \in  \left( {0,\pi }\right)$ ,所以 $B = \dfrac{\pi }{3}$
            \item 由(1)可知 $A + C = \dfrac{2\pi }{3} \Longrightarrow  y = \sin A + \sin C = \sin A + \sin \left( {\dfrac{2\pi }{3} - }\right.  A) = \dfrac{3}{2}\sin A + \dfrac{\sqrt{3}}{2}\cos A = \sqrt{3}\sin \left( {A + \dfrac{\pi }{6}}\right)$ ,由 $\bigtriangleup {ABC}$ 为锐角三角形,得 $0 < A < \dfrac{\pi }{2}$ ,且 $0 < \dfrac{2\pi }{3} - A < \dfrac{\pi }{2} \Longrightarrow  \dfrac{\pi }{6} < A < \dfrac{\pi }{2}$ ,从而 $\dfrac{\pi }{3} < A + \dfrac{\pi }{6} < \dfrac{2\pi }{3} \Longrightarrow  y = \sin A + \sin C = \sqrt{3}\sin \left( {A + \dfrac{\pi }{6}}\right)$ 的取值范围为 $\left( {\dfrac{3}{2},\sqrt{3}}\right\rbrack$.
        \end{enumerate}
    \end{jieda}

    \item (2025$\cdot$长宁区一模$\cdot$17)在 $\bigtriangleup  {ABC}$ 中，角 $A$ 、 $B$ 、 $C$ 所对的边分别为 $a$ 、 $b$ 、 $c$ ，且 $b\sin A - \sqrt{3}a\cos B = 0$.
    \begin{enumerate}
        \item 求角 $B$ 的大小
        \item 若 $b = 2,{\bigtriangleup {ABC}}$ 的面积为 $\sqrt{3}$ ，请判断 ${\bigtriangleup {ABC}}$ 的形状，并说明理由.
    \end{enumerate}
    \begin{jieda}
        \begin{enumerate}
            \item 由正弦定理可得 $\sin B\sin A - \sqrt{3}\sin A\cos B = 0$ ，因为 $\sin A \neq  0$ ，所以 $\tan B = \sqrt{3}$ ， $B \in  \left( {0,\pi }\right)$ ，所以 $B = \dfrac{\pi }{3}$
            
            \item ${S}_{\bigtriangleup {ABC}} = \dfrac{1}{2}{ac}\sin B = \sqrt{3}$ ，所以 ${ac} = 4$ ，由余弦定理 ${b}^{2} =  {a}^{2} + {c}^{2} - {2ac}\cos B = 4$ ,得 ${a}^{2} + {c}^{2} = 8$ ,即 $\dfrac{16}{{c}^{2}} + {c}^{2} = 8$ ,解得 $c =  2,a = 2$ ,所以 $\bigtriangleup {ABC}$ 是等边三角形.
        \end{enumerate}
    \end{jieda}

    \item (2024$\cdot$闵行区二模$\cdot$17)在锐角 $\bigtriangleup  {ABC}$ 中，角 $A$ 、 $B$ 、 $C$ 所对边的边长分别为 $a$ 、 $b$ 、 $c$ ，且 ${2b}\sin A -  \sqrt{3}a = 0.$
    \begin{enumerate}
        \item 求角 $B$
        \item 求 $\sin A + \sin C$ 的取值范围.
    \end{enumerate}
    \begin{jieda}
        \begin{enumerate}
            \item 因为 ${2b}\sin A - \sqrt{3}a = 0$ ，所以 $2\sin A\sin B - \sqrt{3}\sin A =$ 0,又因为 $\sin A \neq  0$，所以 $\sin B = \dfrac{\sqrt{3}}{2}$，所以 $B = \dfrac{\pi }{3}$.

            \item 因为 $\bigtriangleup  {ABC}$ 为锐角三角形，所以 $A \in  \left( {\dfrac{\pi }{6},\dfrac{\pi }{2}}\right)$ ，则 $\sin A +  \sin C = \sin A + \sin \left( {\dfrac{2\pi }{3} - A}\right)  = \dfrac{3}{2}\sin A + \dfrac{\sqrt{3}}{2}\cos A = \sqrt{3}\sin (A +  \left. \dfrac{\pi }{6}\right)$ ,所以 $\sin A + \sin C \in  \left( {\dfrac{3}{2},\sqrt{3}}\right\rbrack$.
        \end{enumerate}
    \end{jieda}

    \item (2024$\cdot$嘉定区一模$\cdot$17)已知 $\bigtriangleup  {ABC}$ 中， $\vv{CA} \cdot  \vv{CB} = 1$ ， $\bigtriangleup  {ABC}$ 的面积 $S = \dfrac{1}{2}$ 
    \begin{enumerate}
        \item 求角 $C$ 的值
        \item 若 $\sin A\cos A = \dfrac{\sqrt{3}}{4},a = 2$ ，求 $c$.
    \end{enumerate}
    \begin{jieda}
        \begin{enumerate}
            \item $\vv{CA} \cdot  \vv{CB} = 1 \Longrightarrow  {ab}\cos C = 1,S = \dfrac{1}{2} \Longrightarrow  {ab}\sin C = 1$ ，两式相除得: $\tan C = 1$ ,所以 $C = \dfrac{\pi }{4}$ 
            \item $\sin A\cos A = \dfrac{\sqrt{3}}{2} \Longrightarrow  \sin {2A} = \dfrac{\sqrt{3}}{2}$ ，所以 $A = \dfrac{\pi }{6}$ 或 $\dfrac{\pi }{3},B = \dfrac{7\pi }{12}$ 或 $\dfrac{5\pi }{12},B > A$ ,由 (1) 得 ${ab} = \sqrt{2}$ ,因为 $a = 2$ ,所以 $b = \dfrac{\sqrt{2}}{2}$ ,则 $b < a$ ,由大边对大角得: $B < A$ ,由此矛盾,所以不存在三角形,于是 $c$ 无解.
        \end{enumerate}
    \end{jieda}

    \item (2025$\cdot$宝山区一模$\cdot$18)在 $\bigtriangleup  {ABC}$ 中，已知 ${b}^{2} + {c}^{2} = {a}^{2} + {bc}$.
    \begin{enumerate}
        \item 若 $\sin C = {2\sin B}$ ，且 $b = 2$ ，求 $\bigtriangleup  {ABC}$ 的面积
        \item 若 $b + c = 1$ ，求 $a$ 的取值范围.
    \end{enumerate}
    \begin{jieda}
        \begin{enumerate}
            \item 由正弦定理得 $\dfrac{c}{b} = \dfrac{\sin C}{\sin B} = 2$ ，又 $b = 2$ ，从而 $c = 4$ ，由 ${b}^{2} + {c}^{2} = {a}^{2} + {bc}$ 得 $\cos A = \dfrac{{b}^{2} + {c}^{2} - {a}^{2}}{2bc} = \dfrac{bc}{2bc} = \dfrac{1}{2}$ ，从而 $A = \dfrac{\pi }{3}$ ， 所以 $\bigtriangleup {ABC}$ 的面积 $S = \dfrac{1}{2}{bc}\sin A = \dfrac{1}{2} \times  2 \times  4 \times  \sin \dfrac{\pi }{3} = 2\sqrt{3}$
            \item $1 = b + c = \dfrac{a}{sinA} \cdot (sinB + sinC) = \dfrac{a}{sinA} \cdot \left(sinB + sin\left(\dfrac{2\pi}{3} - B\right)\right) = \dfrac{a}{sinA} \cdot \sqrt{3}sin\left(B + \dfrac{\pi}{6}\right)$
            
            $a = \dfrac{sinA}{\sqrt{3}sin\left(B + \dfrac{\pi}{6}\right)} = \dfrac{1}{2sin\left(B + \dfrac{\pi}{6}\right)}$，$B+\dfrac{\pi}{6} \in \left(\dfrac{\pi}{6}, \dfrac{5\pi}{6}\right)$，故$a \in \left[\dfrac{1}{2}, 1\right)$
        \end{enumerate}
    \end{jieda}

    \item (2024$\cdot$崇明区二模$\cdot$18)在锐角三角形 $\bigtriangleup  {ABC}$ 中，角 $A$ 、 $B$ 、 $C$ 的对边分别为 $a$ 、 $b$ 、 $c$. 若 $\vv{CD}$ 为 $\vv{CA}$ 在 $\vv{CB}$ 方向上的投影向量，且满足 ${2c}\sin B = \sqrt{5}\left| \vv{CD}\right|$.
    \begin{enumerate}
        \item 求 $\cos C$ 的值
        \item 若 $b = \sqrt{3},a = {3c}\cos B$ ，求 $\bigtriangleup  {ABC}$ 的周长
    \end{enumerate}
    \begin{jieda}
        \begin{enumerate}
            \item 由题意,得 $\left| \vv{CD}\right|  = b\cos C$ ,又 ${2c}\sin B = \sqrt{5}\left| \vv{CD}\right|$ ,所以 ${2c}\sin B = \sqrt{5}b\cos C$ ,由正弦定理 $\dfrac{b}{\sin B} = \dfrac{c}{\sin C}$ ,得 $2\sin C\sin B =  \sqrt{5}\sin B\cos C$ ,又 $\sin B \neq  0$ ,所以 $2\sin C = \sqrt{5}\cos C$ ,因为 $C$ 为锐角,所以 $\cos C = \sqrt{1 - {\sin }^{2}C} = \dfrac{2}{3}$
            \item 由 $a = {3c}\cos B$ ，得 $\cos B = \dfrac{a}{3c}$ ，由余弦定理得 $\cos B =  \dfrac{{a}^{2} + {c}^{2} - {b}^{2}}{2ac}$ ,所以 $\dfrac{{a}^{2} + {c}^{2} - {b}^{2}}{2ac} = \dfrac{a}{3c}$ ,得 ${a}^{2} + 3{c}^{2} = 3{b}^{2}$
            
            由余弦定理得 $\cos C = \dfrac{{a}^{2} + {b}^{2} - {c}^{2}}{2ab} = \dfrac{2}{3}$ ,得 $3{a}^{2} + 3{b}^{2} - 3{c}^{2} =  4\sqrt{3}a$
            
            联立两式，解得 $a = \sqrt{3},c = \sqrt{2}$ ，故 $\bigtriangleup  {ABC}$ 的周长为 $a + b + c = 2\sqrt{3} + \sqrt{2}$.
        \end{enumerate}
    \end{jieda}

    \item (2024$\cdot$崇明区一模$\cdot$18)在 $\bigtriangleup  {ABC}$ 中， $a = 5$ ， $b = 6$.
    \begin{enumerate}
        \item 若 $\cos B =  - \dfrac{4}{5}$ ，求 $A$ 和 $\bigtriangleup  {ABC}$ 外接圆半径 $R$ 的值
        \item 若 $\bigtriangleup  {ABC}$ 的面积 $S = \dfrac{{15}\sqrt{7}}{4}$ ，求 $c$ 的值.
    \end{enumerate}
    \begin{jieda}
        \begin{enumerate}
            \item 因为 $\cos B =  - \dfrac{4}{5},B \in  \left( {0,\pi }\right)$ ,所以 $\sin B =  \sqrt{1 - {\cos }^{2}B} = \dfrac{3}{5}$ ,由正弦定理,得 $\dfrac{a}{\sin A} = \dfrac{b}{\sin B} = {2R}$ ,即 $\dfrac{5}{\sin A} =  \dfrac{6}{\dfrac{3}{5}} = {2R}$ ,所以 $\sin A = \dfrac{1}{2},R = 5$ ,因为 $a < b$ ,所以 $A \in  (0$ , $\left. \dfrac{\pi }{2}\right)$ ,因此 $A = \dfrac{\pi }{6},R = 5$
            \item 由 ${S}_{\bigtriangleup {ABC}} = \dfrac{1}{2}{ab}\sin C$ 得 $\sin C = \dfrac{2{S}_{\bigtriangleup {ABC}}}{ab} = \dfrac{2 \times  \dfrac{{15}\sqrt{7}}{4}}{5 \times  6} =  \dfrac{\sqrt{7}}{4}$ ,于是 $\cos C =  \pm  \sqrt{1 - {\sin }^{2}C} =  \pm  \dfrac{3}{4}$
            \begin{dingautolist}{172}
                \item 当 $\cos C = \dfrac{3}{4}$ 时,由余弦定理,得 ${c}^{2} = {5}^{2} + {6}^{2} - 2 \times  5 \times  6 \times  \dfrac{3}{4} = {16}$
                \item 当 $\cos C =  - \dfrac{3}{4}$ 时,由余弦定理,得 ${c}^{2} = {5}^{2} + {6}^{2} - 2 \times  5 \times  6 \times  \left( {-\dfrac{3}{4}}\right)  = {106}$ . 所以, $c = 4$ 或 $c = \sqrt{106}$.
            \end{dingautolist}
        \end{enumerate}
    \end{jieda}

    \item (2025$\cdot$上海春季卷$\cdot$18)在 $\bigtriangleup  {ABC}$ 中，角 $A$ 、 $B$ 、 $C$ 所对的边分别为 $a$ 、 $b$ 、 $c$ ，且 $c = 5$.
    \begin{enumerate}
        \item 若 $\dfrac{a}{4b} = \dfrac{\sin B}{\sin A},C = \dfrac{\pi }{2}$ ,求 $a$
        \item 若 ${ab} = {20}$ ，求 $\bigtriangleup  {ABC}$ 的面积的最大值.
    \end{enumerate}
    \begin{jieda}
        \begin{enumerate}
            \item 由 $\dfrac{a}{4b} = \dfrac{\sin B}{\sin A}$ ，根据正弦定理得 $\dfrac{a}{4b} = \dfrac{b}{a}$ ，化简得 $a =  {2b}$ . 因为 $C = \dfrac{\pi }{2},c = 5$ ,所以 ${a}^{2} + {b}^{2} = {c}^{2} = {25}$ ,即 $4{b}^{2} + {b}^{2} =$ 25，解得 $b = \sqrt{5}$ ，所以 $a = {2b} = {2\sqrt{5}}$
            \item 由余弦定理可得 $\cos C = \dfrac{{a}^{2} + {b}^{2} - {c}^{2}}{2ab} = \dfrac{1}{40}\left( {{a}^{2} + {b}^{2} - {25}}\right)  \geq  \dfrac{1}{40}\left( {{2ab} - {25}}\right)  = \dfrac{3}{8}$ ,当且仅当 $a = b = 2\sqrt{5}$ 时,等号成立. 所以 ${\cos }^{2}C \geq  \dfrac{9}{64}$ ,即 $1 - {\sin }^{2}C \geq  \dfrac{9}{64}$ ,所以 ${\sin }^{2}C \leq  \dfrac{55}{64}$ ,又因为 $C \in  \left( {0,\pi }\right)$ ,所以 $\sin C > 0$ ,故 $\sin C \leq  \dfrac{\sqrt{55}}{8}$ ,由 ${S}_{\bigtriangleup {ABC}} =  \dfrac{1}{2}{ab}\sin C$ 得 ${S}_{\bigtriangleup {ABC}} = \dfrac{1}{2} \times  {20} \times  \sin C = {10}\sin C \leq  {10} \times  \dfrac{\sqrt{55}}{8} =  \dfrac{5\sqrt{55}}{4}$ ,故当 $a = b = 2\sqrt{5}$ 时, ${\left( {S}_{\bigtriangleup {ABC}}\right) }_{\max } = \dfrac{5\sqrt{55}}{4}$.
        \end{enumerate}
    \end{jieda}
    
    \item (2023$\cdot$上海春季卷$\cdot$18)在 $\bigtriangleup  {ABC}$ 中，角 $A$ 、 $B$ 、 $C$ 对应边为 $a$ 、 $b$ 、 $c$ ，其中 $b = 2$.
    \begin{enumerate}
        \item 若 $A + C = {120}^{ \circ  }$ ，且 $a = {2c}$ ，求边长 $c$ 的值
        \item 若 ${A - C} = {15}^{ \circ  },{a = \sqrt{2}c\sin A}$ ，求 ${S}_{\bigtriangleup {ABC}}$
    \end{enumerate}
    \begin{jieda}
        \begin{enumerate}
            \item $B = \dfrac{\pi }{3}$ ,由余弦定理, $\cos B = \dfrac{4{c}^{2} + {c}^{2} - {2}^{2}}{2 \cdot  {2c} \cdot  c} = \dfrac{1}{2} \Longrightarrow  {c}^{2} = \dfrac{4}{3} \Longrightarrow  c = \dfrac{2\sqrt{3}}{3};$
            \item 由正弦定理， $a = \sqrt{2}c\sin A \Longrightarrow  \sin A = \sqrt{2}\sin C\sin A \Longrightarrow  \sin C =  \dfrac{\sqrt{2}}{2} \Longrightarrow  C = \dfrac{\pi }{4}$ 或 $\dfrac{3\pi }{4}$ ,所以 $A = C + \dfrac{\pi }{12} = \dfrac{\pi }{3}$ 或 $\dfrac{5\pi }{6}$ (出现两个钝角,舍去),即 $B = \dfrac{5\pi }{12}$ ,由正弦定理, $\dfrac{b}{\sin B} = \dfrac{a}{\sin A} \Longrightarrow  a = 3\sqrt{2} -  \sqrt{6}$ ,所以 ${S}_{\bigtriangleup {ABC}} = \dfrac{1}{2}{ab}\sin C = 3 - \sqrt{3}$ .
        \end{enumerate}
    \end{jieda}

    \item (2024$\cdot$嘉定区二模$\cdot$18)在 $\bigtriangleup  {ABC}$ 中，角 $A$ 、 $B$ 、 $C$ 的对边分别为 $a$ 、 $b$ 、 $c$ ， ${\cos }^{2}B - {\sin }^{2}B =  - \dfrac{1}{2}$.
    \begin{enumerate}
        \item 求角 $B$ ，并计算 $\sin \left( {B + \dfrac{\pi }{6}}\right)$ 的值
        \item 若 $b = \sqrt{3}$ ，且 $\bigtriangleup  {ABC}$ 是锐角三角形，求 $a + {2c}$ 的最大值
    \end{enumerate}
    \begin{jieda}
        \begin{enumerate}
            \item 因为 ${\cos }^{2}B - {\sin }^{2}B =  - \dfrac{1}{2}$ ，所以 $\cos {2B} =  - \dfrac{1}{2}$ 又 $B \in  \left( {0,\pi }\right)$ ,所以 $B = \dfrac{\pi }{3}$ 或 $B = \dfrac{2\pi }{3}$ . 所以 $\sin \left( {B + \dfrac{\pi }{6}}\right)  =  \sin \dfrac{\pi }{2} = 1$ 或 $\sin \left( {B + \dfrac{\pi }{6}}\right)  = \sin \dfrac{5\pi }{6} = \dfrac{1}{2}$
            \item 由正弦定理， $a + {2c} = \dfrac{b}{\sin B}\left( {\sin A + 2\sin C}\right)  = 2(\sin A +  2\sin C)$ ,由 $\bigtriangleup {ABC}$ 是锐角三角形,故 $\dfrac{\pi }{6} < A < \dfrac{\pi }{2},f\left( A\right)  =  2\left\lbrack  {\sin A + 2\sin \left( {\dfrac{2\pi }{3} - A}\right) }\right\rbrack   = 2\left( {2\sin A + \sqrt{3}\cos A}\right)  \leq  2\sqrt{7}$ ,当等号成立时, $\tan A = \dfrac{2\sqrt{3}}{3} > \dfrac{\sqrt{3}}{3} = \tan \dfrac{\pi }{6}$ 满足题意,所以 $a + {2c}$ 的最大值为 $2\sqrt{7}$.
        \end{enumerate}
    \end{jieda}

    \item (2024$\cdot$闵行区一模$\cdot$18)在 $\bigtriangleup  {ABC}$ 中，角 $A$ 、 $B$ 、 $C$ 所对边的边长分别为 $a$ 、 $b$ 、 $c$ ，且 $a - {2c}\cos B = c$.
    \begin{enumerate}
        \item 若 $\cos B = \dfrac{1}{3},c = 3$ ，求 $b$ 的值
        \item 若 $\bigtriangleup  {ABC}$ 为锐角三角形，求 $\sin C$ 的取值范围.
    \end{enumerate}
    \begin{jieda}
        \begin{enumerate}
            \item 将 $\cos B = \dfrac{1}{3},c = 3$ 带人条件中可得 $a = 5$ ，由余弦定理 ${b}^{2} = {a}^{2} + {c}^{2} - {2ac}\cos B$ 可得 $b = 2\sqrt{6}$
            \item 因为 $a - {2c}\cos B = c$ ，由正弦定理可得 $\sin A - 2\sin C\cos B =  \sin C$ ,因为 $\sin A = \sin \left( {B + C}\right)  = \sin B\cos C + \cos B\sin C$ ,所以 $\sin B\cos C - \sin C\cos B = \sin C,\sin \left( {B - C}\right)  = \sin C$ ,因为 $B -  C \in  \left( {-\dfrac{\pi }{2},\dfrac{\pi }{2}}\right) ,C \in  \left( {0,\dfrac{\pi }{2}}\right)$ ,所以 $B - C = C$ ,即 $B = {2C}$ ,又因为 $\bigtriangleup {ABC}$ 为锐角三角形,所以 $C \in  \left( {\dfrac{\pi }{6},\dfrac{\pi }{4}}\right) ,\sin C \in  \left( {\dfrac{1}{2},\dfrac{\sqrt{2}}{2}}\right)$
        \end{enumerate}
    \end{jieda}

    \item (2024$\cdot$杨浦区二模$\cdot$18)已知 $f\left( x\right)  = \sin {\omega x}\left( {\omega  > 0}\right)$.
    \begin{enumerate}
        \item 若 $y = f\left( x\right)$ 的最小正周期为 ${2\pi }$ ,判断函数 $F\left( x\right)  = f\left( x\right)  + f\left( {x + \dfrac{\pi }{2}}\right)$ 的奇偶性,并说明理由
        \item 已知 $\omega  = 2,{\bigtriangleup {ABC}}$ 中， $a$ 、 $b$ 、 $c$ 分别是角 $A$ 、 $B$ 、 $C$ 所对的边，若 $f\left( {A + \dfrac{\pi }{3}}\right)  = 0,a = 2,b = 3$ ，求 $c$ 的值.
    \end{enumerate}
    \begin{jieda}
        \begin{enumerate}
            \item 由 $T = \dfrac{2\pi }{\omega }$ 知: $\omega  = 1$ ，此时 $F\left( x\right)  = f\left( x\right)  + f(x +  \left. \dfrac{\pi }{2}\right)  = \sin x + \sin \left( {x + \dfrac{\pi }{2}}\right)  = \sin x + \cos x,F\left( {-x}\right)  + F\left( x\right)  =  \sin \left( {-x}\right)  + \cos \left( {-x}\right)  + \sin x + \cos x = 2\cos x \neq  0;F\left( {-x}\right)  -  F\left( x\right)  = \sin \left( {-x}\right)  + \cos \left( {-x}\right)  - \sin x - \cos x =  - 2\sin x \neq  0$ 所以 $F\left( x\right)$ 非奇非偶.
            \item $\omega  = 2$ ，则 $f\left( x\right)  = \sin {2x}$ ， $f\left( {A + \dfrac{\pi }{3}}\right)  = \sin \left( {{2A} + \dfrac{2\pi }{3}}\right)  = 0$ ， 因为 $0 < A < \pi$ ,所以 $A = \dfrac{\pi }{6}$ 
            
            根据余弦定理有$a^2 = b^2 + c^2 - 2bccosA$，即$4 = 9 + c^2 - 3\sqrt{3}c$，解得$c = \dfrac{3\sqrt{3} \pm \sqrt{7}}{2}$
        \end{enumerate}
    \end{jieda}

    \item (2024$\cdot$宝山区一模$\cdot$18)在 $\bigtriangleup  {ABC}$ 中，角 $A$ 、 $B$ 、 $C$ 的对边分别为 $a$ 、 $b$ 、 $c$.
    \begin{enumerate}
        \item 若 ${2a}\sin B = \sqrt{3}b$ 求角 $A$ 的大小
        \item 若 ${BC}$ 边上的高等于 $\dfrac{a}{2}$ ，求 $\dfrac{c}{b} + \dfrac{b}{c}$ 的最大值.
    \end{enumerate}
    \begin{jieda}
        \begin{enumerate}
            \item 根据正弦定理得 $2\sin A\sin B = \sqrt{3}\sin B$ ,所以 $\sin A =  \dfrac{\sqrt{3}}{2}$ ，所以 $A = \dfrac{\pi }{3}$ 或 $\dfrac{2\pi }{3}$
            \item $\dfrac{c}{b} + \dfrac{b}{c} = \dfrac{b^2+c^2}{bc} = \dfrac{a^2+2bccosA}{bc} = \dfrac{a^2}{bc}+2cosA = \dfrac{sin^2A}{sinBsinC}+2cosA$
            
            $\dfrac{a}{2} = bsinC$，故$sinA = 2sinBsinC$，因此$\dfrac{sin^2A}{sinBsinC} = 2sinA$

            $\dfrac{c}{b} + \dfrac{b}{c} = 2\left( {\sin A + \cos A}\right)  = 2\sqrt{2}\sin \left( {A + \dfrac{\pi }{4}}\right)$ . 当$A = \dfrac{\pi }{4}$ 时 $\dfrac{{b}^{2} + {c}^{2}}{bc}$ 有最大值 $2\sqrt{2}$.

            \ps 高的条件用面积处理亦可.
        \end{enumerate}
    \end{jieda}

    \item (2025$\cdot$崇明区一模$\cdot$18)在 $\bigtriangleup  {ABC}$ 中，已知点 $D$ 是 ${BC}$ 边上一点，且 ${BD} = 1$ ， ${CD} = 3$.
    \begin{enumerate}
        \item 若 ${AD}\bot {BC}$ ，且 $\angle {ABD} = 2\angle {ACD}$ ，求 ${AD}$ 的长
        \item 若 $\angle {ABD} = {55}^{ \circ  },\angle {ACD} = {32}^{ \circ  }$ ，求 ${AD}$ 的长. (结果精确到 0.01)
    \end{enumerate}
    \begin{jieda}
        \begin{enumerate}
            \item 因为 ${AD} \bot  {BC}$ ,所以 $\tan \angle {ABD} = \dfrac{AD}{1},\tan \angle {ACD} =  \dfrac{AD}{3}$ ,又 $\angle {ABD} = 2\angle {ACD}$ ,所以 $\tan \angle {ABD} = \dfrac{2\tan \angle {ACD}}{1 - {\tan }^{2}\angle {ACD}}$ 即 ${AD} = \dfrac{2 \times  \dfrac{AD}{3}}{1 - {\left( \dfrac{AD}{3}\right) }^{2}}$ ,解得 ${AD} = \sqrt{3}$
            \item 在 $\bigtriangleup {ABC}$ 中， $\angle {BAC} = {93}^{ \circ  },{BC} = 4$ ，由正弦定理得 $\dfrac{BC}{\sin \angle {BAC}} = \dfrac{AC}{\sin \angle {ABC}}$ ,所以 ${AC} = \dfrac{{BC} \cdot  \sin \angle {ABC}}{\sin \angle {BAC}} \approx 3.281$，在 $\bigtriangleup {ACD}$ 中,由余弦定理得${AD} = \sqrt{A{C}^{2} + D{C}^{2} - {2AC} \cdot  {DC} \cdot  \cos \angle {ACD}} \approx  {1.75}.$
        \end{enumerate}
    \end{jieda}

    \item (2024$\cdot$青浦区一模$\cdot$18)在 $\bigtriangleup {ABC}$ 中，角 $A$ 、 $B$ 、 $C$ 所对的边分别为 $a$ 、 $b$ 、 $c$ ，且满足 ${a}^{2} + {c}^{2} - {b}^{2} +  {ac} = 0.$
    \begin{enumerate}
        \item 求角 $B$ 的大小
        \item 若 $b = 2\sqrt{3}$ ，求 $\bigtriangleup  {ABC}$ 的周长的最大值.
    \end{enumerate}
    \begin{jieda}
        \begin{enumerate}
            \item 因为 ${a}^{2} + {c}^{2} - {b}^{2} =  - {ac}$ ，由余弦定理得 $\cos \angle B =  \dfrac{{a}^{2} + {c}^{2} - {b}^{2}}{2ac} =  - \dfrac{1}{2},\angle B = {120}^{ \circ  }$
            
            \item 由正弦定理得， $a = 4\sin A$ ， $c = 4\sin \left( {{60}^{ \circ  } - A}\right)$ ，所以 $\bigtriangleup {ABC}$ 的周长为 $a + b + c = 4\sin A + 4\sin \left( {{60}^{ \circ  } - A}\right)  + 2\sqrt{3} =  4\sin \left( {A + {60}^{ \circ  }}\right)  + 2\sqrt{3},{0}^{ \circ  } < A < {60}^{ \circ  }$ ,当 $A = {30}^{ \circ  }$ 时, $\bigtriangleup {ABC}$ 的周长的最大值为 $4 + 2\sqrt{3}$.
        \end{enumerate}
    \end{jieda}

    \item (2025$\cdot$杨浦区 一模$\cdot$18)已知 $\bigtriangleup  {ABC}$ 的内角 $A$ 、 $B$ 、 $C$ 所对边的长度分别为 $a$ 、 $b$ 、 $c$.
    \begin{enumerate}
        \item 若 ${\left( a + b\right) }^{2} - {c}^{2} = 4,C = {60}^{ \circ  }$ ，求 $\bigtriangleup  {ABC}$ 的面积
        \item 若 $\dfrac{\cos C}{c} = \dfrac{\cos A}{{3b} - a}$ ，求 $\sin C$ 的值.
    \end{enumerate}
    \begin{jieda}
        \begin{enumerate}
            \item 由 ${\left( a + b\right) }^{2} - {c}^{2} = 4$ ，得 ${a}^{2} + {b}^{2} - {c}^{2} = 4 - {2ab}$ ，由余弦定理得 ${c}^{2} = {a}^{2} + {b}^{2} - {2ab}\cos C$ ，即 ${a}^{2} + {b}^{2} - {c}^{2} = {ab}$ ，所以 4- ${2ab} = {ab}$ ,即 ${ab} = \dfrac{4}{3}$ ,所以 $\bigtriangleup {ABC}$ 的面积为 $\dfrac{1}{2}{ab}\sin C = \dfrac{1}{2} \times  \dfrac{4}{3} \times  \dfrac{\sqrt{3}}{2} = \dfrac{\sqrt{3}}{3}$
            
            \item 因为 $\dfrac{\cos C}{c} = \dfrac{\cos A}{{3b} - a}$ ，由正弦定理得 $\dfrac{\cos C}{\sin C} = \dfrac{\cos A}{3\sin B - \sin A}$ ， 可得 $3\sin B\cos C - \sin A\cos C = \cos A\sin C$ ,则 $3\sin B\cos C =  \cos A\sin C + \sin A\cos C = \sin \left( {A + C}\right)  = \sin B$ ,因为 $B \in  \left( {0,\pi }\right)$ , 所以 $\sin B > 0$ ,则 $\cos C = \dfrac{1}{3}$ ,又 $C \in  \left( {0,\pi }\right)$，所以 $\sin C = \sqrt{1 - {\cos }^{2}C} = \dfrac{2\sqrt{2}}{3}$.
        \end{enumerate}
    \end{jieda}
\end{enumerate}

\subsection*{考点6 解三角形之实际应用}

\subsubsection*{一、填空题}
\begin{enumerate}
    \item (2025$\cdot$静安区一模$\cdot$10)如图所示，小明和小宁家都住在东方明珠塔附近的同一幢楼上，小明家在 $A$ 层,小宁家位于小明家正上方的 $B$ 层,已知 ${AB} = a$. 小明在家测得东方明珠塔尖的仰角为 $\alpha$ ,小宁在家测得东方明珠塔尖的仰角为 $\beta$ ,则他俩所住的这幢楼与东方明珠塔之间的距离 $d =$ \blkda[2]{$\dfrac{a}{\tan \alpha  - \tan \beta }$}.
    \begin{center}
        \includegraphics[width=0.2\textwidth]{bodyxb/src/54_214_1748_332_248_0.jpg}
    \end{center}
    \begin{jieda}
        分别过点 $A\text{ 、 }B$ 作 ${CD}$ 的垂线,垂足分别为 $M\text{ 、 }N$ ,则根据正切函数的定义得 ${CM} = d\tan \alpha ,{CN} =  d\tan \beta$ ,则 ${MN} = {AB} = d\tan \alpha  -  d\tan \beta  = a$ ,解得 $d = \dfrac{a}{\tan \alpha  - \tan \beta }$. 

        \begin{center}
        \includegraphics[width=0.2\textwidth]{bodyxb/src/38_1317_700_262_209_0.jpg}
        \end{center}
    \end{jieda}

    \item (2024$\cdot$奉贤区一模$\cdot$10)某林场为了及时发现火情，设立了两个观测点 $A$ 和 $B$. 某日两个观测点的林场人员都观测到 $C$ 处出现火情. 在 $A$ 处观测到火情发生在北偏西 ${40}^{ \circ  }$ 方向，而在 $B$ 处观测到火情在北偏西 ${60}^{ \circ  }$ 方向. 已知 $B$ 在 $A$ 的正东方向 ${10}{{km}}$ 处 (如图所示),则 $\left| {{BC} - {AC}}\right|  =$ \blkda[2]{$7.8$} ${{km}}$. (精确到 ${0.1}{{km}}$ )

    \begin{center}
    \includegraphics[width=0.2\textwidth]{bodyxb/src/54_650_1745_305_248_0.jpg}
    \end{center}
    \begin{jieda}
        $\dfrac{BC}{\sin {130}^{ \circ  }} = \dfrac{AC}{\sin {30}^{ \circ  }} = \dfrac{AB}{\sin {20}^{ \circ  }} = \dfrac{10}{\sin {20}^{ \circ  }} \Longrightarrow  \left| {{BC} - {AC}}\right|  =  \dfrac{10}{\sin {20}^{ \circ  }}\left( {\sin {130}^{ \circ  } - \sin {30}^{ \circ  }}\right)  \approx  {7.8}\left( {{km}}\right).$
    \end{jieda}

    \item (2024$\cdot$徐汇区一模$\cdot$10)某建筑物内一个水平直角型过道(如图所示)，两过道的宽度均为 $3{m}$ ，有一个水平截面为矩形的设备需要水平通过直角型过道. 若该设备水平截面矩形的宽 ${BC}$ 为 $1{m}$ ，则该设备能水平通过直角型过道的长 ${AB}$ 不超过\blkda[2]{$6\sqrt{2} - 2$} ${m}$.

    \begin{center}
    \includegraphics[width=0.2\textwidth]{bodyxb/src/54_1058_1743_292_241_0.jpg}
    \end{center}
    \begin{jieda}
        \begin{dingautolist}{172}
            \item 设$\angle{ABO} = \alpha$，$O$到$CD$距离为$|AB|sin\alpha cos\alpha + 1 \leq 3\sqrt{2}$恒成立，故$|AB| \leq 6\sqrt{2} - 2$. 经检验$|AB| = 6\sqrt{2} - 2$可行.
            
            \ps 由几何直观可知$\alpha = \dfrac{\pi}{4}$时可以通过即可.
            \item 如图所示, 当该设备恰好卡住时,过点 $C$ 作垂直于过道墙面的直线, 垂足分别为 $G\text{ 、 }H$ . 记 $\angle {CBG} = \theta  \in  \left( {0,\dfrac{\pi }{2}}\right)$ ,则 ${CG} = \sin \theta ,{CH} = 3 - \sin \theta$ , ${CM} = \dfrac{3 - \sin \theta }{\cos \theta }$ ,同理可得 ${DM} = \dfrac{3 - \cos \theta }{\sin \theta }$ ,因此 ${AB} = {CM} +  {DM} = \dfrac{3\left( {\sin \theta  + \cos \theta }\right)  - 1}{\sin \theta \cos \theta }$ ,令 $t = \sin \theta  + \cos \theta  = \sqrt{2}\sin \left( {\theta  + \dfrac{\pi }{4}}\right)  \in  (1,\sqrt{2}\rbrack$ ,则 ${AB} = \dfrac{{3t} - 1}{\dfrac{{t}^{2} - 1}{2}} = \dfrac{2\left( {{3t} - 1}\right) }{{t}^{2} - 1}$ ,令 $s = {3t} - 1 \in  (2,3\sqrt{2} -$ 1],则 $t = \dfrac{s + 1}{3}$ . 故 ${AB} = \dfrac{2s}{\dfrac{{\left( s + 1\right) }^{2}}{9} - 1} = \dfrac{18}{s - \dfrac{8}{s} + 2}$ ,因为 $y =  s = \dfrac{8}{s} + 2$ 在 $\left( {2,3\sqrt{2} - 1}\right\rbrack$ 上为严格增函数,所以 $s = 3\sqrt{2} - 1$ , 即 $t = \sqrt{2}$ 即 $\theta  = \dfrac{\pi }{4}$ 时, $A{B}_{\min } = \dfrac{3 - \dfrac{\sqrt{2}}{2}}{\dfrac{\sqrt{2}}{2}} + \dfrac{3 - \dfrac{\sqrt{2}}{2}}{\dfrac{\sqrt{2}}{2}} = 6\sqrt{2} - 2$ ,即设备能水平通过直角型过道的长 ${AB}$ 不超过 $6\sqrt{2} - 2$ 米.
            \item 设$\angle{ABO} = \alpha$，$CD: \dfrac{x}{|AB|cos\alpha + \dfrac{1}{sin\alpha}} + \dfrac{y}{|AB|sin\alpha + \dfrac{1}{cos\alpha}} = 1$，由$O,M$在$CD$异侧或$M$在$CD$上，则$\dfrac{3}{|AB|cos\alpha + \dfrac{1}{sin\alpha}} + \dfrac{3}{|AB|sin\alpha + \dfrac{1}{cos\alpha}} \geq 1$，即$\dfrac{3(sin\alpha+cos\alpha)-1}{sin\alpha cos\alpha}\geq |AB|$，同上，得$|AB| \leq 6\sqrt{2} - 2$
        \end{dingautolist}
    \end{jieda}

    \item (2025$\cdot$奉贤区一模$\cdot$10)申辉中学高一(8)班设计了一个“水滴状”班徽的平面图 (如图所示),徽章由等腰三角形 ${ABC}$ 及以弦 ${BC}$ 和劣弧 ${BC}$ 所围成的弓形所组成，其中 ${AB} = {AC}$ ，劣弧 ${BC}$ 所在的圆为三角形的外接圆， 圆心为 $O$. 已知 $\angle {BAC} = \theta$ ， $\theta  \in  \left( {0\text{ ， }\dfrac{\pi }{2}}\right)$ ，外接圆的半径是 2，则该图形的面积为\blkda[2]{${4\theta } + 4\sin \theta$}. (用含 $\theta$ 的表达式表示)
    \begin{jieda}
        由正弦定理得 $\dfrac{\left| AB\right| }{\sin \left( \dfrac{\pi  - \theta }{2}\right) } = \dfrac{\left| AB\right| }{\cos \dfrac{\theta }{2}} = 4 \Longrightarrow  \left| {AB}\right|  =  4\cos \dfrac{\theta }{2}$ ,且可知 $\angle {BOC} = {2\theta }$ ,则该图形的面积为 $S = \dfrac{1}{2}\left| {AB}\right| \left| {AC}\right| \sin \theta  + {2\theta } \times  2 - \dfrac{1}{2} \times  2 \times  2 \times  \sin {2\theta },$

        $S = \dfrac{1}{2}{\left( 4\cos \dfrac{\theta }{2}\right) }^{2}\sin \theta  + {4\theta } - 2\sin {2\theta } = 8{\left( \cos \dfrac{\theta }{2}\right) }^{2}\sin \theta  + {4\theta } - 2\sin {2\theta }  = 4\left( {\cos \theta  + 1}\right) \sin \theta  + {4\theta } - 2\sin {2\theta } = 2\sin {2\theta } + 4\sin \theta  + {4\theta } - 2\sin {2\theta }  = 4\sin \theta  + {4\theta }$.
    \end{jieda}

    \begin{center}
    \includegraphics[width=0.2\textwidth]{bodyxb/src/55_1138_170_262_274_0.jpg}
    \end{center}

    \item (2025$\cdot$浦东新区二模$\cdot$10)如图所示，某建筑物 ${OP}$ 垂直于地面，从地面点 $A$ 处测得建筑物顶部 $P$ 的仰角为 ${30}^{ \circ  }$ ,从地面点 $B$ 处测得建筑物顶部 $P$ 的仰角为 ${45}^{ \circ  }$ ,已知 $A\text{ 、 }B$ 相距 ${100}{m},\angle {AOB} = {60}^{ \circ  }$ ,则该建筑物 ${OP}$ 高度约为\blkda[2]{$66.4$}m. (保留一位小数)
    \begin{center}
    \includegraphics[width=0.2\textwidth]{bodyxb/src/55_120_1349_275_261_0.jpg}
    \end{center}
    \begin{jieda}
        ${OA} = \sqrt{3}{OP}$，${OB} = {OP}$ . 在 $\bigtriangleup  {AOB}$ 中，已知 ${AB} =$ 100 米, $\angle {AOB} = {60}^{ \circ  }$ . 根据余弦定理 $A{B}^{2} = O{A}^{2} + O{B}^{2} - 2 {OA} \cdot  {OB} \cdot  \cos \angle {AOB}$
        
        ${100}^{2} = {\left( \sqrt{3}OP\right) }^{2} + O{P}^{2} - 2 \cdot  \sqrt{3}{OP} \cdot  {OP} \cdot  \cos {60}^{ \circ  }$，则 ${OP} = \sqrt{\dfrac{{10000}\left( {4 + \sqrt{3}}\right) }{13}} \approx  {66.4}$.
    \end{jieda}

    \item (2025$\cdot$闵行区二模$\cdot$11)如图，现对某景区一长 ${AB} = {600}{m}$ ，宽 ${AD} = {360}{m}$ 的矩形空地进行建设. 规划在边 ${AB}$ 、 ${AD}$ 上分别取点 $M$ 、 $N$ 修建人行步道(不考虑宽度)，且满足点 $A$ 关于步道 ${MN}$ 的对称点 $E$ 在边 ${DC}$ 上. 在 $\bigtriangleup  {AMN}$ 内种植花卉，在 $\bigtriangleup  {EMN}$ 内搭建娱乐设施，其余区域规划为露营区，则人行步道 ${MN}$ 的最短距离为\blkda[2]{$468$}m. (结果精确到 $1{m}$ )
    \begin{center}
    \includegraphics[width=0.2\textwidth]{bodyxb/src/55_421_1387_307_222_0.jpg}
    \end{center}
    \begin{jieda}
        \begin{dingautolist}{172}
            \item 以$A$为原点建系，设$E(t, 360), t \in [0, 600]$，则可设$MN: tx+360y+c=0$，代入$AE$中点坐标$\left(\dfrac{t}{2}, 180\right)$，可得$c = -\dfrac{t^2+360^2}{2}$，即$MN: tx+360y-\dfrac{t^2+360^2}{2}=0$，故$M\left(\dfrac{t^2+360^2}{2t},0\right), N\left(0, \dfrac{t^2+360^2}{720}\right)$，故$|MN| = \dfrac{1}{2}\sqrt{\left(t+\dfrac{360^2}{t}\right)^2+\left(360+\dfrac{t}{360}\right)^2}$，计算器打表可得最小值为$270\sqrt{3}$
            \item 以$A$为原点建系，设$\angle AMN = \alpha$，$MN: \dfrac{x}{|MN|cos\alpha} + \dfrac{y}{|MN|sin\alpha} = 1$，因为$E$在$y = 360$上，因此$AE$中点纵坐标应为$180$，$MN$上纵坐标为$180$的点为$\left(|MN|cos\alpha-180cot\alpha, 180\right)$，故$\vv{n} = \left(|MN|cos\alpha-180cot\alpha, 180\right)$为$MN$法向量，因此$\dfrac{180}{|MN|cos\alpha} = \dfrac{|MN|cos\alpha-180cot\alpha}{|MN|sin\alpha}$，整理得$|MN| = \dfrac{180}{sin\alpha - sin^3\alpha}$，根据几何关系易知$\dfrac{\sqrt{10}}{10} \leq  \sin \theta  \leq  \dfrac{\sqrt{2}}{2}$ , 令 $\sin \theta  = t,t \in  \left\lbrack  {\dfrac{\sqrt{10}}{10},\dfrac{\sqrt{2}}{2}}\right\rbrack$ ,则 ${MN} = \dfrac{180}{t - {t}^{3}}$ . 令 $f\left( t\right)  = t - {t}^{3}$ , 则 ${f}^{\prime }\left( t\right)  = 1 - 3{t}^{2}$ ,令 ${f}^{\prime }\left( t\right)  > 0$ ,则 $t \in  \left\lbrack  {\dfrac{\sqrt{10}}{10},\dfrac{\sqrt{3}}{3}}\right)$ ,函数 $y =  f\left( t\right)$ 严格增; 令 ${f}^{\prime }\left( t\right)  < 0$ ,则 $t \in  \left( {\dfrac{\sqrt{3}}{3},\dfrac{\sqrt{2}}{2}}\right\rbrack$ ,函数 $y = f\left( t\right)$ 严格减; 所以 $f{\left( t\right) }_{\max } = f\left( \dfrac{\sqrt{3}}{3}\right)  = \dfrac{2\sqrt{3}}{9}$ ,所以 $M{N}_{\min } = {270}\sqrt{3} \approx$ 468,即人行步道 ${MN}$ 的最短距离为 ${468}{m}$ .
            \item 由题意, ${Rt}\bigtriangleup {MAN} \cong  {Rt}\bigtriangleup {MEN}$ ,设 $\angle {AMN} = \theta  \left( {0 < \theta  < \dfrac{\pi }{2}}\right)$ ,则 $\angle {DNE} = \pi  - 2\left( {\dfrac{\pi }{2} - \theta }\right)  = {2\theta }$ ,在 ${Rt}\bigtriangleup {DEN}$ 中, $\cos {2\theta } = \dfrac{DN}{NE} = \dfrac{{360} - {AN}}{AN}$ ,得 ${AN} = \dfrac{360}{\cos {2\theta } + 1} = \dfrac{180}{{\cos }^{2}\theta }$ ,则 ${MN} = \dfrac{AN}{\sin \theta } = \dfrac{180}{\sin \theta {\cos }^{2}\theta } = \dfrac{180}{\sin \theta \left( {1 - {\sin }^{2}\theta }\right) }$ ,由于 $\left\{  {\begin{array}{l} {AN} = \dfrac{180}{{\cos }^{2}\theta } = {180}\left( {1 + {\tan }^{2}\theta }\right)  \leq  {360} \\  {AM} = \dfrac{AN}{\tan \theta } = {180}\left( {\tan \theta  + \dfrac{1}{\tan \theta }}\right)  \leq  {600} \end{array}\text{ ,解得 }\dfrac{1}{3} \leq  \tan \theta  \leq  1}\right.$ ,则 ${\tan }^{2}\theta  = \dfrac{{\sin }^{2}\theta }{{\cos }^{2}\theta } = \dfrac{{\sin }^{2}\theta }{1 - {\sin }^{2}\theta } \in  \left\lbrack  {\dfrac{1}{9},1}\right\rbrack$ ,解得 $\dfrac{\sqrt{10}}{10} \leq  \sin \theta  \leq  \dfrac{\sqrt{2}}{2}$ , 令 $\sin \theta  = t,t \in  \left\lbrack  {\dfrac{\sqrt{10}}{10},\dfrac{\sqrt{2}}{2}}\right\rbrack$ ,则 ${MN} = \dfrac{180}{t - {t}^{3}}$ . 令 $f\left( t\right)  = t - {t}^{3}$ , 则 ${f}^{\prime }\left( t\right)  = 1 - 3{t}^{2}$ ,令 ${f}^{\prime }\left( t\right)  > 0$ ,则 $t \in  \left\lbrack  {\dfrac{\sqrt{10}}{10},\dfrac{\sqrt{3}}{3}}\right)$ ,函数 $y =  f\left( t\right)$ 严格增; 令 ${f}^{\prime }\left( t\right)  < 0$ ,则 $t \in  \left( {\dfrac{\sqrt{3}}{3},\dfrac{\sqrt{2}}{2}}\right\rbrack$ ,函数 $y = f\left( t\right)$ 严格减; 所以 $f{\left( t\right) }_{\max } = f\left( \dfrac{\sqrt{3}}{3}\right)  = \dfrac{2\sqrt{3}}{9}$ ,所以 $M{N}_{\min } = {270}\sqrt{3} \approx$ 468,即人行步道 ${MN}$ 的最短距离为 ${468}{m}$ .
        \end{dingautolist}
    \end{jieda}

    \item (2025$\cdot$徐汇区二模$\cdot$11)如图所示，某处有一块圆心角为 $\dfrac{2}{3}\pi$ 的扇形绿地 ${AOB},$ 扇形的半径为 ${20}{m},{AB}$ 是一条原有的人行直路,由于工程建设需要,现要在绿地中建一条直路 ${OC}$ ,以便在图中阴影部分区域分类堆放物料. 为了尽量减少对绿地的破坏(不计路宽)，则原直路 ${AB}$ 与新直路 ${OC}$ 的交叉点 $D$ 到 $O$ 的距离为\blkda[2]{${10}\sqrt{2}$} ${m}$.
    \begin{center}
    \includegraphics[width=0.2\textwidth]{bodyxb/src/55_752_1405_316_200_0.jpg}
    \end{center}
    \begin{jieda}
        过点 $O$ 作 ${OH} \bot  {AB}$ 垂足为 $H$ ,可得 ${OH} = {10}$ , ${AB} = {20}\sqrt{3}$ ，设 $\angle {BOD} = \alpha ,{BD} = m$ ，在 $\bigtriangleup  {BOD}$ 中，由正弦定理得 $\dfrac{m}{\sin \alpha } = \dfrac{OB}{\sin \angle {ODB}}$ ，因为 $\sin \angle {ODB} = \sin \left( {\alpha  + \angle {OBD}}\right)  =  \sin \left( {\alpha  + \dfrac{\pi }{6}}\right)$ ,所以 $m = \dfrac{{20}\sin \alpha }{\sin \left( {\alpha  + \dfrac{\pi }{6}}\right) }$ ,又由阴影部分的面积: $S = {S}_{\text{ 扇形 }{OBC}} - {S}_{\bigtriangleup {OBD}} + {S}_{\bigtriangleup {OAD}} = \dfrac{1}{2}\alpha  \times  {\left( {20}\right) }^{2} - \dfrac{1}{2}m \times  {10} +  \dfrac{1}{2}\left( {{20}\sqrt{3} - m}\right)  \times  {10} = {200\alpha } + {100}\sqrt{3} - {10m} = {200\alpha } + {100}\sqrt{3} -  \dfrac{{200}\sin \alpha }{\sin \left( {\alpha  + \dfrac{\pi }{6}}\right) } = {200}\left( {\alpha  - \dfrac{\sin \alpha }{\sin \left( {\alpha  + \dfrac{\pi }{6}}\right) }}\right)  + {100}\sqrt{3}$ ,其中 $0 < \alpha  \leq  \dfrac{\pi }{2}$ ,令 $f\left( \alpha \right)  = \alpha  - \dfrac{\sin \alpha }{\sin \left( {\alpha  + \dfrac{\pi }{6}}\right) },0 < \alpha  \leq  \dfrac{\pi }{2}$ ,可得 ${f}^{\prime }\left( \alpha \right)  = 1 -  \dfrac{\cos \alpha \sin \left( {\alpha  + \dfrac{\pi }{6}}\right)  - \sin x\cos \left( {\alpha  + \dfrac{\pi }{6}}\right) }{{\sin }^{2}\left( {\alpha  + \dfrac{\pi }{6}}\right) } = 1 - \dfrac{\dfrac{1}{2}}{{\sin }^{2}\left( {\alpha  + \dfrac{\pi }{6}}\right) }$ ,令 ${f}^{\prime }\left( \alpha \right)  = 0$ ,可得 $\sin \left( {\alpha  + \dfrac{\pi }{6}}\right)  = \dfrac{\sqrt{2}}{2}$ ,解得 $\alpha  = \dfrac{\pi }{12}$ ,当 $0 < \alpha  < \dfrac{\pi }{12}$ 时, ${f}^{\prime }\left( \alpha \right)  < 0$ ; 当 $\dfrac{\pi }{12} < \alpha  < \dfrac{\pi }{2}$ 时, ${f}^{\prime }\left( \alpha \right)  > 0$ ,所以 $f\left( \alpha \right)$ 在 $\left( {0,\dfrac{\pi }{12}}\right)$ 上严格减,在 $\left( {\dfrac{\pi }{12},\dfrac{\pi }{2}}\right)$ 上严格增,当 $\alpha  = \dfrac{\pi }{12}$ 时,函数 $y = f\left( \alpha \right)$ 取得最小值,则 $\angle {DOH} =  \dfrac{\pi }{4}$ ,所以 $\bigtriangleup {ODH}$ 为等腰直角三角形, 因为 $\left| {OH}\right|  = {10}$ ,所以 $\left| {OD}\right|  = {10}\sqrt{2}$.
        \begin{center}
            \includegraphics[width=0.2\textwidth]{bodyxb/src/39_1298_991_232_141_0.jpg}
        \end{center}
    \end{jieda}

    \item (2025$\cdot$上海秋季卷$\cdot$11)小申同学观察发现，生活中有些时候影子可以完全投射在斜面上. 某斜面上有两根长为 $1{m}$ 的垂直于水平面放置的杆子，与斜面的接触点分别为 $A$ 、 $B$ ，它们在阳光的照射下呈现出影子. 阳光可视为平行光,其中一根杆子的影子在水平面上,长度为 ${0.4}{m}$ ,另一根杆子的影子完全在斜面上，长度为 ${0.45}{m}$. 则斜面的底角 $\theta  =$ \blkda[2]{${12.58}^{ \circ  }$}.(结果用角度制表示，精确到 0.01 \textdegree )
    \begin{center}
    \includegraphics[width=0.2\textwidth]{bodyxb/src/55_1095_1384_335_219_0.jpg}
    \end{center}
    \begin{jieda}
        如图,在 $A$ 处, $\tan x = \dfrac{1}{0.4} = {2.5}$ ,在 $B$ 处满足 $\tan \angle {CED} = {2.5}$ ,(其中 ${ED}//$ 水平面, ${CE}$ 是射过 $B$ 处杆子最高点的光线,光线交斜面于 $E$ ),故设 ${BD} = y$ ,则 ${ED} =  \dfrac{1 + y}{2.5}$ ,由勾股定理, ${y}^{2} + {\left( \dfrac{1 + y}{2.5}\right) }^{2} = {0.45}^{2}$ ,解得 $y \approx  {0.098}$ , 于是 $\theta  = \arcsin \dfrac{0.098}{0.45} \approx  {12.58}^{ \circ  }$.

        \begin{center}
        \includegraphics[width=0.3\textwidth]{bodyxb/src/39_989_1509_409_236_0.jpg}
        \end{center}
    \end{jieda}


    \item (2025$\cdot$闵行区二模$\cdot$11)已知某星球的球心为 $F$ ，半径为 $R$ ，该星球的卫星的运行轨道是以 $F$ 为一个焦点的椭圆,该椭圆的离心率为 $\dfrac{3}{5}$ ，卫星运行过程中离该星球表面最近的距离为 $R$ ，若当卫星处于某位置时, 用卫星上的光学仪器观测该星球, 把光学仪器的镜头与星球表面被观测点的连线称为视线，任意两条视线所成的最大夹角称为张角，则卫星运行过程中张角的最小值为\blkda[2]{${14.4}^{ \circ}$}. (精确到 $\left. {0.{1}^{ \circ  }}\right)$
    \begin{jieda}
        设椭圆轨道的半长轴为 $a$ ,焦距为 ${2c}$ ,由题意可得, 卫星运行过程中离该星球表面最近的距离为 $R$ ,在近日点即 $a - c = {2R}$ ,又椭圆的离心率为 $\dfrac{3}{5}$ ,即 $\dfrac{c}{a} = \dfrac{3}{5}$ ,由以上可得 $a = {5R},c = {3R}$，易知在远日点时张角最小, 设此时张角为 $\theta$ ,则 $\sin \left( \dfrac{\theta }{2}\right)  = \dfrac{R}{8R} = \dfrac{1}{8}$ ,即 $\theta  = 2\arcsin \left( \dfrac{1}{8}\right)  \approx {14.4}^{ \circ}$. 
    \end{jieda}

    \item (2024$\cdot$金山区二模$\cdot$11)某临海地区为保障游客安全修建了海上救生栈道，如图所示，线段 ${BC}$ 、 ${CD}$ 是救生栈道的一部分，其中 ${BC} = {300}{m}$ ， ${CD} = {800}{m}$ ， $B$ 在 $A$ 的北偏东 ${30}^{ \circ  }$ 方向， $C$ 在 $A$ 的正北方向， $D$ 在 $A$ 的北偏西 ${80}^{ \circ  }$ 方向，且 $\angle B = {90}^{ \circ  }$. 若救生艇在 $A$ 处载上遇险游客需要尽快抵达救生栈道 $B - C - D$ ，则最短距离为\blkda[2]{$475$}m. (结果精确到 $1{m}$ )
    \begin{center}
    \includegraphics[width=0.2\textwidth]{bodyxb/src/56_275_950_325_222_0.jpg}
    \end{center}
    \begin{jieda}
        作 ${AE} \bot  {CD}$ 交与 $E$ ,由题意得 $\angle B = {90}^{ \circ  },\angle {CAB} = {30}^{ \circ  },{BC} = {300}{m}$ ，所以 ${AB} = {BC}\cot {30}^{ \circ  } = {300}\sqrt{3}{m},{AC} =  {600}{m}$ ,在 $\bigtriangleup {ADC},\cos {80}^{ \circ  } = \dfrac{{600}^{2} + A{D}^{2} - {800}^{2}}{2 \times  {600} \times  {AD}}$ ,解得 ${AD} \approx  {643.50}{m},{S}_{\bigtriangleup {ADC}} = \dfrac{1}{2}{AD} \times  {AC} \times  \sin {80}^{ \circ  } \approx  {190117.1367}{{m}}^{2},{AE} =  \dfrac{2{S}_{\bigtriangleup {ADC}}}{CD} \approx  {475.29} < {300}\sqrt{3} = {AB}$ ,所以最短距离为 ${475}{m}$ .

        \begin{center}
        \includegraphics[width=0.2\textwidth]{bodyxb/src/40_634_285_233_164_0.jpg}
        \end{center}
    \end{jieda}

    \item (2025$\cdot$奉贤区二模$\cdot$11)中企联合大厦是目前奉贤区的第一高楼，是奉贤“美奉贤强”的一个缩影. 某数学建模兴趣小组的同学们去实地进行测量, 经过多次的测量, 最终在平行于地面的同一水平面上选取三个点: 点 $B$ 、点 $C$ 、点 $O$ 作为测量基点. 设大厦的最高点为 $A$ ,在点 $B$ 处测得点 $A$ 的仰角为 $\angle {ABO} = \theta  = {71.5}^{ \circ  }$ ,在点 $C$ 处测得点 $A$ 的仰角为 $\angle {ACO} = {44}^{ \circ  }$ ,又测得 ${BC} = {221}{m}$ , $\angle {BOC} = {119}^{ \circ  }$ ,如图所示. 现作出以下几个假设:
    
    \begin{center}
    \begin{minipage}{0.45\textwidth}
    \begin{dingautolist}{172}
        \item 直线 ${AO}$ 垂直于平面 ${OBC}$
        \item 平面 ${OBC}$ 到地面的距离等于测角仪高度，在计算过程中测角仪高度忽略不计
        \item 其他次要因素等忽略不计
    \end{dingautolist}
    \end{minipage}
    \begin{minipage}{0.45\textwidth}
        \includegraphics[width=0.5\textwidth]{bodyxb/src/56_671_908_300_262_0.jpg}
    \end{minipage}
    \end{center}
    根据以上信息估算奉贤第一高楼的高度约\blkda[2]{180}m. (结果保留整数)
    \begin{jieda}
        设高度为 $h$ ,在 $\bigtriangleup {ABO}$ 中,根据正弦定理有: $\dfrac{h}{\sin {71.5}^{ \circ  }} = \dfrac{BO}{\sin \left( {{90}^{ \circ  } - {71.5}^{ \circ  }}\right) }$ ,即 ${BO} = \tan {18.5}^{ \circ  }h \approx  {0.33h}$ ,在 $\bigtriangleup {ACO}$ 中,根据正弦定理有: $\dfrac{h}{\sin {44}^{ \circ  }} = \dfrac{CO}{\sin \left( {{90}^{ \circ  } - {44}^{ \circ  }}\right) }$ ,即 ${CO} =  \tan {46}^{ \circ  }h \approx  {1.04h}$ ,由余弦定理可知: $B{C}^{2} = B{O}^{2} + C{O}^{2} -  {2BO} \cdot  {CO}\cos {119}^{ \circ  } = {0.33}^{2}{h}^{2} + {1.04}^{2}{h}^{2} - 2 \times  {0.33h} \times  {1.04h}  \times  \left( {-{0.48}}\right)  \approx  {0.11}{h}^{2} + {1.08}{h}^{2} + {0.33}{h}^{2} \approx  {1.5}{h}^{2} = {221}^{2}$ ,解得: $h \approx  {180}$.
    \end{jieda}

    \item (2025$\cdot$黄浦区一模$\cdot$11)一个机器零件的形状是有缺口的圆形铁片，如图中实线部分为裁剪后的形状. 已知这个圆的半径是 ${13}{{cm}},{AB} = 8{{cm}},{BC} = 6{{cm}}$ ，且 ${AB}\bot {BC}$ ，则圆心到点 $B$ 的距离约为\blkda[2]{$7.3$}cm. (结果精确到 ${0.1}{{cm}}$ )
    \begin{center}
    \includegraphics[width=0.2\textwidth]{bodyxb/src/56_1042_919_248_253_0.jpg}
    \end{center}
    \begin{jieda}
        如图所示,设圆心为 $D,{AC}$ 的中点为 $E$ ,则 ${AD} = {13}$ ,在 $\operatorname{Rt}\bigtriangleup {ABC}$ 中, ${AC} =  \sqrt{{8}^{2} + {6}^{2}} = {10} = {2AE} \Longrightarrow  {AE} = 5,\cos \angle {BAC}  = \dfrac{AB}{AC} = \dfrac{8}{10} = \dfrac{4}{5} \Longrightarrow  \sin \angle {BAC} = \dfrac{3}{5}$ ,在 ${{Rt}}\bigtriangleup {DEA}$ 中, $\cos \angle {DAE} = \dfrac{AE}{AD} = \dfrac{5}{13} \Longrightarrow  \sin \angle {DAE} = \dfrac{12}{13}$ ,在 $\bigtriangleup {DAB}$ 中, $\cos \angle {BAD} = \cos \left( {\angle {DAE} - \angle {BAC}}\right)  = \dfrac{5}{13} \times  \dfrac{4}{5} +  \dfrac{12}{13} \times  \dfrac{3}{5} = \dfrac{56}{65} \Longrightarrow  B{D}^{2} = {13}^{2} + {8}^{2} - 2 \cdot  {13} \cdot  8 \cdot  \dfrac{56}{65} = {53.8} \Longrightarrow  {BD} \approx  {7.3}{{cm}}.$

        \begin{center}
        \includegraphics[width=0.2\textwidth]{bodyxb/src/40_691_903_177_185_0.jpg}
        \end{center}
    \end{jieda}
\end{enumerate}
\subsubsection*{二、选择题}
\begin{enumerate}
    \item (2024$\cdot$嘉定区二模$\cdot$15)嘉定区某校学习小组开展测量太阳高度角的数学活动. 太阳高度角是指某时刻太阳光线和地平面所成的角. 测量时, 假设太阳光线均为平行的直线, 地面为水平平面. 如图所示，两竖直墙面所成的二面角为 ${120}^{ \circ  }$ ，墙的高度均为 $3{m}$. 在时刻 $t$ ，实地测量得在太阳光线照射下的两面墙在地面的阴影宽度分别为 $1{m}$ 、 ${1.5}{m}$. 在线查阅嘉定区的气象资料，当天的太阳高度角和对应时间的部分数据如表所示，则时刻 $t$ 最可能为 \dotfill{(\kh{B})}

    \begin{center}
    \begin{minipage}{0.45\textwidth}
        \includegraphics[width=0.5\linewidth]{bodyxb/src/56_990_1633_315_281_0.jpg}
    \end{minipage}
    \begin{minipage}{0.45\textwidth}
    \begin{tabular}{|c|c|c|c|}
    \hline
    太阳高度角 & 时间 & 太阳高度角 & 时间 \\
    \cline{1-4}
    43.13 \textdegree  & 08:30 & 68.53 \textdegree  & 10:30 \\
    \cline{1-4}
    49.53 \textdegree  & 09:00 & 74.49 \textdegree  & 11:00 \\
    \cline{1-4}
    55.93 \textdegree  & 09:30 & 79.60 \textdegree  & 11:30 \\
    \cline{1-4}
    62.29 \textdegree  & 10:00 & 82.00 \textdegree  & 12:00 \\
    \cline{1-4}
    \hline
    \end{tabular} 
    \end{minipage}
    \end{center}
    \xx{$09:00$}{$10:00$}{$11:00$}{$12:00$}
    \begin{jieda}
        如图所示: $\bigtriangleup {BCD}$ 中,由余弦定理可得, $B{D}^{2} = B{C}^{2} + C{D}^{2} - 2 \times  {BC} \times  {CD} \times  \cos \alpha  = 1.{5}^{2} + {1}^{2} - 2 \times$ 1. $5 \times  1 \times  \cos {60}^{ \circ  } = \dfrac{7}{4},{BD} = \dfrac{\sqrt{7}}{2},A\text{ 、 }B$ 、 $C\text{ 、 }D$ 四点共圆,由正弦定理可得,

        \begin{center}
        \includegraphics[width=0.2\textwidth]{bodyxb/src/40_644_1428_225_222_0.jpg}
        \end{center}
        

        ${2r} = \dfrac{BD}{\sin {60}^{ \circ  }} = \dfrac{\dfrac{\sqrt{7}}{2}}{\dfrac{\sqrt{3}}{2}} = \dfrac{\sqrt{21}}{3}$ ,太阳高度角为 $\theta ,\tan \theta  = \dfrac{OA}{AC} = \dfrac{3}{\dfrac{\sqrt{21}}{3}} =  \dfrac{3}{7}\sqrt{21},\theta  \approx  {63}^{ \circ  }$ ,比对太阳高度角和对应时间的部分数据,可得此时为 10 点,故选 B.
    \end{jieda}
\end{enumerate}
\subsubsection*{三、解答题}
\begin{enumerate}
    \item (2024$\cdot$长宁区一模$\cdot$19)汽车转弯时遵循阿克曼转向几何原理，即转向时所有车轮中垂线交于一点,该点称为转向中心: 如图 1,某汽车四轮中心分别为 $A\text{ 、 }B\text{ 、 }C\text{ 、 }D$ ,向左转向,左前轮转向角为 $\alpha$ ,右前轮转向角为 $\beta$ ,转向中心为 $O$. 设该汽车左右轮距 ${AB}$ 为 $w{m}$ ,前后轴距 ${AD}$ 为 $l{m}$.
    \begin{center}
    \includegraphics[width=0.7\textwidth]{bodyxb/src/57_302_446_933_324_0.jpg}
    \end{center}
    \begin{enumerate}
        \item 试用 $w$ 、 $l$ 和 $\alpha$ 表示 $\tan \beta$
        \item 如图2，有一直角弯道， $M$ 为内直角顶点， ${EF}$ 为上路边,路宽均为3.5m，汽车行驶其中，左轮 $A$ 、 $D$ 与路边 ${FS}$ 相距 $2{m}$. 试依据如下假设，对问题 * 做出判断，并说明理由.

        假设:  \ding{172}  转向过程中,左前轮转向角 $\alpha$ 的值始终为 ${30}^{ \circ  }$ ;  \ding{173}  设转向中心 $O$ 到路边 ${EF}$ 的距离为 $d$ ,若 ${OB} < d$ 且 ${OM} < {OD}$ ,则汽车可以通过,否则不能通过;  \ding{174}  $w = {4.570},l = {2.680}$. 
        
        问题 * : 可否选择恰当转向位置, 使得汽车通过这一弯道?
    \end{enumerate}
    \begin{jieda}
        \begin{enumerate}
            \item ${OD} = \dfrac{AD}{\tan \alpha } =  \dfrac{l}{\tan \alpha },{OC} = {OD} + {DC} = \dfrac{l}{\tan \alpha } + w$ ,所以 $\tan \beta  = \dfrac{BC}{OC} = \dfrac{l}{\dfrac{l}{\tan \alpha } + w}$ 
            
            \item 以 ${EF}$ 和 ${FS}$ 分别为 $x$ 轴和 $y$ 轴建立坐标系,则 $M(-{3.5}, -{3.5}),{OD} = \dfrac{l}{\tan \alpha } = \sqrt{3}l = {4.642},{OB} = \sqrt{{l}^{2} + {\left( \sqrt{3}l + w\right) }^{2}} =$ 6.766,设 $O\left( {a,b}\right) \left( {a < 0,b < 0}\right)$ ,可得 $a =  - \sqrt{3}l - 2 =  - {6.642},d =  - b,{OM} = \sqrt{{\left( {3.5} + a\right) }^{2} + {\left( {3.5} + b\right) }^{2}} = \sqrt{{9.875} + {\left( {3.5} + b\right) }^{2}}$ ,由 ${OM} < {OD}$ ,可得 ${9.872} + {\left( {3.5} + b\right) }^{2} < {21.548}$ ,解得 $-6.917 < b <  - {0.83}$ ,由 ${OB} < d$ ,得 $b <  - {6.766}$ ,所以当$-6.917 < b <  - {6.766}$ 时, ${BO} < d$ 且 ${OM} < {OD}$ ,此时汽车可以通过弯道. 可知选择恰当转向位置, 汽车可以通过弯道.

            \begin{center}
            \includegraphics[width=0.3\textwidth]{bodyxb/src/40_1229_354_348_266_0.jpg}
            \end{center}
        \end{enumerate}
    \end{jieda}

    \item (2024$\cdot$浦东新区一模$\cdot$19)某街道规划建一座口袋公园. 如图所示，公园由扇形 ${AOC}$ 区域和 $\bigtriangleup  {COD}$ 区域组成，其中 $A$ 、 $O$ 、 $D$ 三点共线，扇形半径 ${OA}$ 为 ${30}{m}$. 规划口袋公园建成后，扇形 ${AOC}$ 区域将作为花草展示区, $\bigtriangleup {COD}$ 区域作为亲水平台区,两个区域的所有边界修建休闲步道.
    \begin{enumerate}
        \item 若 $\angle {AOC} = \dfrac{\pi }{3},{OD} = {2OA}$ ，求休闲步道总长；(精确到 $m$)
        \item 若 $\angle {ODC} = \dfrac{\pi }{6}$ ，在前期民意调查时发现，绝大部分居民对亲水平台区更感兴趣. 请你根据民意调查情况，从该区域面积最大或周长最长的视角出发，选择其中一个方案，设计 $\bigtriangleup  {COD}$ 的形状.
    \end{enumerate}

    \begin{center}
    \includegraphics[width=0.3\textwidth]{bodyxb/src/58_1033_486_369_175_0.jpg}
    \end{center}
    \begin{jieda}
        \begin{enumerate}
            \item 休闲步道总长为 $\overset{\frown}{AC} + {2OA} + {OD} + {CD} = \dfrac{\pi }{3} \times  {30} +  {120} + \sqrt{{30}^{2} + {60}^{2} - 2 \times  {30} \times  {60} \times  \cos \dfrac{2\pi }{3}} = {10\pi } + {120} + {30}\sqrt{7} \approx$ 231 米. 所以休闲步道总长为 231 米
            
            \item 方案一:设 $\angle {COD} = \theta$ ， $\theta  \in  \left( {0,\dfrac{5\pi }{6}}\right)$ ， $\bigtriangleup  {COD}$ 中，由正弦定理得 $\dfrac{OC}{\sin \dfrac{\pi }{6}} = \dfrac{OD}{\sin \left( {\dfrac{5\pi }{6} - \theta }\right) } = \dfrac{CD}{\sin \theta }$ ,得 ${OD} = {60}\sin \left( {\dfrac{5\pi }{6} - \theta }\right)$ , ${CD} = {60}\sin \theta ,\theta  \in  \left( {0,\dfrac{5\pi }{6}}\right)$ ,故 $\bigtriangleup {COD}$ 的面积 $S = \dfrac{1}{2} \times  {30} \times  {60}\sin \left( {\dfrac{5\pi }{6} - \theta }\right)  \cdot  \sin \theta  = {900}\sin \left( {\dfrac{5\pi }{6} - \theta }\right)  \cdot  \sin \theta  = {450}\sin ({2\theta } -  \left. \dfrac{\pi }{3}\right)  + {225}\sqrt{3}$ ,因为 $\theta  \in  \left( {0,\dfrac{5\pi }{6}}\right)$ ,所以 ${2\theta } - \dfrac{\pi }{3} \in  \left( {-\dfrac{\pi }{3},\dfrac{4\pi }{3}}\right)$ , 当 ${2\theta } - \dfrac{\pi }{3} = \dfrac{\pi }{2}$ ,即 $\theta  = \dfrac{5\pi }{12}$ 时有 ${S}_{\max } = {450} + {225}\sqrt{3}$ 平方米,因此,当亲水平台区的面积最大时, $\bigtriangleup {COD}$ 是以 ${OC}$ 为底边的等腰三角形.

            方案二: 设 $\angle {COD} = \theta ,\theta  \in  \left( {0,\dfrac{5\pi }{6}}\right) ,\bigtriangleup {COD}$ 中,由正弦定理得 $\dfrac{OC}{\sin \dfrac{\pi }{6}} = \dfrac{OD}{\sin \left( {\dfrac{5\pi }{6} - \theta }\right) } = \dfrac{CD}{\sin \theta }$ ,得 ${OD} = {60}\sin \left( {\dfrac{5\pi }{6} - \theta }\right) ,{CD} =  {60}\sin \theta ,\theta  \in  \left( {0,\dfrac{5\pi }{6}}\right)$ ,故 $\bigtriangleup  {COD}$ 的周长 $L = {60}\sin \theta  + {60}\sin \left( {\dfrac{5\pi }{6} - }\right.  \theta ) + {30} = \left( {{60} + {30}\sqrt{3}}\right) \sin \theta  + {30}\cos \theta  + {30} = {60}\sqrt{2 + \sqrt{3}}\sin (\theta  +  \left. \dfrac{\pi }{12}\right)  + {30}$ ,因为 $\theta  \in  \left( {0,\dfrac{5\pi }{6}}\right)$ ,所以 $\theta  + \dfrac{\pi }{12} \in  \left( {\dfrac{\pi }{12},\dfrac{11\pi }{12}}\right)$ ,当 $\theta  +  \dfrac{\pi }{12} = \dfrac{\pi }{2}$ ,即 $\theta  = \dfrac{5\pi }{12}$ 时有 ${L}_{\max } = {60}\sqrt{2 + \sqrt{3}} + {30} = {30}\sqrt{6} +  {30}\sqrt{2} + {30}$ 米,因此,当亲水平台区的周长最长时, $\bigtriangleup {COD}$ 是以 ${OC}$ 为底边的等腰三角形.
        \end{enumerate}
    \end{jieda}


    \item (2025$\cdot$松江区一模$\cdot$19)为了打造美丽社区，某小区准备将一块由一个半圆和长方形组成的空地进行美化，如图所示，长方形的边 ${AB}$ 为半圆的直径， $O$ 为半圆的圆心， ${AB} = {2AD} = {200}{m}$. 现要将此空地规划出一个等腰三角形区域 ${PMN}$ (底边 ${MN} \bot  {CD}$ ) 种植观赏树木,其余区域种植花卉. 设 $\angle {MOB} = \theta ,\theta  \in  \left( {0,\dfrac{\pi }{2}}\right\rbrack$.

    \begin{center}
    \includegraphics[width=0.2\textwidth]{bodyxb/src/58_1082_1395_326_287_0.jpg}
    \end{center}
    \begin{enumerate}
        \item 当 $\theta  = \dfrac{\pi }{3}$ 时，求 $\bigtriangleup  {PMN}$ 的面积
        \item 求三角形区域 ${PMN}$ 面积的最大值
    \end{enumerate}
    \begin{jieda}
        \begin{enumerate}
            \item 设 ${MN}$ 与 ${AB}$ 相交于点 $E$ ，则 $\left| {MN}\right|  = \left| {ME}\right|  + \left| {EN}\right|  =  {100}\sin \dfrac{\pi }{3} + {100} = {50}\sqrt{3} + {100},\left| {AE}\right|  = {100} + {100}\cos \dfrac{\pi }{3} =$ 150,易知 $\left| {AE}\right|$ 等于 $P$ 到 ${MN}$ 的距离,故 ${S}_{\bigtriangleup {PMN}} = \dfrac{1}{2}\left| {MN}\right|$ . $\left| {AE}\right|  = \dfrac{1}{2} \times  \left( {{50}\sqrt{3} + {100}}\right)  \times  {150} = {3750}\sqrt{3} + {7500}\left( {{m}}^{2}\right)$
            \item 过点 $P$ 作 ${PF} \bot  {MN}$ 于点 $F$ ，则 ${PF} = {AE} = {100} + {100}\cos \theta$ ,而 $\left| {MN}\right|  = \left| {ME}\right|  + \left| {EN}\right|  =  {100} + {100}\sin \theta$ ,则 ${S}_{\bigtriangleup {PMN}} = \dfrac{1}{2}\left| {MN}\right|$ . $\left| {PF}\right|  = {5000}\left( {1 + \sin \theta }\right) \left( {1 + \cos \theta }\right)  =  {5000}\left( {1 + \sin \theta  + \cos \theta  + \sin \theta \cos \theta }\right)$ ,设 $\sin \theta  + \cos \theta  = t$ ,因为 $\theta  \in  \left( {0,\dfrac{\pi }{2}}\right\rbrack$ ,所以 $\theta  + \dfrac{\pi }{4} \in  \left( {\dfrac{\pi }{4},\dfrac{3\pi }{4}}\right\rbrack$ ,故 $t = \sin \theta  + \cos \theta  = \sqrt{2}\sin \left( {\theta  + \dfrac{\pi }{4}}\right)  \in  \left\lbrack  {1,\sqrt{2}}\right\rbrack$ ,而 $\sin \theta \cos \theta  = \dfrac{{t}^{2} - 1}{2}$ ,得 ${S}_{\bigtriangleup {PMN}} = {5000}(1 + t +  \left. \dfrac{{t}^{2} - 1}{2}\right)  = {2500}{\left( t + 1\right) }^{2}$ ,当 $t = \sqrt{2}$ 时, ${\left( {S}_{\bigtriangleup {PMN}}\right) }_{\max } = {2500}(\sqrt{2} +  1{)}^{2} = {7500} + {5000}\sqrt{2}\left( {{m}}^{2}\right)$ ,故三角形区域 ${PMN}$ 面积的最大值为 $\left( {{7500} + {5000}\sqrt{2}}\right) {{m}}^{2}$ .

            \begin{center}
            \includegraphics[width=0.2\textwidth]{bodyxb/src/41_573_151_247_215_0.jpg}
            \end{center}
        \end{enumerate}
    \end{jieda}

    
    \item (2024$\cdot$黄浦区一模$\cdot$19)某公园的一个角形区域 ${AOB}$ 如图所示，其中 ${\angle {AOB}} = \dfrac{2\pi }{3}$. 现拟用长度为 ${100}{m}$ 的隔离档板 (折线 ${DCE}$ ) 与部分围墙 (折线 ${DOE}$ ) 围成一个花卉育苗区 ${ODCE}$ ,要求满足 ${OD} = {OC} = {OE}$.

    \begin{center}
    \includegraphics[width=0.2\textwidth]{bodyxb/src/59_1083_382_311_215_0.jpg}
    \end{center}
    \begin{enumerate}
        \item 设 $\angle {DOC} = \dfrac{\pi }{3} + \alpha \left( {-\dfrac{\pi }{3} < \alpha  < \dfrac{\pi }{3}}\right)$ ，试用 $\alpha$ 表示 ${OD}$
        \item 为使花卉育苗区的面积最大，应如何设计？请说明理由.
    \end{enumerate}
    \begin{jieda}
        \begin{enumerate}
            \item 由 $\angle {DOC} = \dfrac{\pi }{3} + \alpha \left( {-\dfrac{\pi }{3} < \alpha  < \dfrac{\pi }{3}}\right) ,\angle {AOB} = \dfrac{2\pi }{3}$ , 可知 $\angle {COE} = \dfrac{\pi }{3} - \alpha$ ,作 ${OF} \bot  {CD}$ ,垂足为 $F$ ,由 ${OD} = {OC}$ , 可知 ${CF} = {DF}$ 且 $\angle {DOF} = \dfrac{1}{2}\angle {DOC} = \dfrac{\pi }{6} + \dfrac{\alpha }{2}$ ,在直角 $\bigtriangleup {DOF}$ 中, ${DF} = {OD}\sin \left( {\dfrac{\pi }{6} + \dfrac{\alpha }{2}}\right)$ ,故 ${CD} = {2OD}\sin \left( {\dfrac{\pi }{6} + }\right.  \left. \dfrac{\alpha }{2}\right)$ ,同理可得 ${EC} = {2OC}\sin \left( {\dfrac{\pi }{6} - \dfrac{\alpha }{2}}\right)  = {2OD}\sin \left( {\dfrac{\pi }{6} - }\right.  \left. \dfrac{\alpha }{2}\right)$ ,所以 ${2OD}\sin \left( {\dfrac{\pi }{6} + \dfrac{\alpha }{2}}\right)  + {2OD}\sin \left( {\dfrac{\pi }{6} - \dfrac{\alpha }{2}}\right)  = {100}$ ,可得 ${OD} = \dfrac{50}{\sin \left( {\dfrac{\pi }{6} + \dfrac{\alpha }{2}}\right)  + \sin \left( {\dfrac{\pi }{6} - \dfrac{\alpha }{2}}\right) }  = \dfrac{50}{\cos \dfrac{\alpha }{2}}$ (米);

            \begin{center}
            \includegraphics[width=0.2\textwidth]{bodyxb/src/41_595_1140_224_158_0.jpg}
            \end{center}
            \item 设花卉育苗区的面积为 $S$ 平方米,则$S = \dfrac{1}{2}O{D}^{2}\sin \left( {\dfrac{\pi }{3} + \alpha }\right)  + \dfrac{1}{2}O{D}^{2}\sin \left( {\dfrac{\pi }{3} - \alpha }\right)  = \dfrac{1}{2} \times  \dfrac{{50}^{2}}{{\cos }^{2}\dfrac{\alpha }{2}}\left\lbrack  {\sin \left( {\dfrac{\pi }{3} + \alpha }\right)  + \sin \left( {\dfrac{\pi }{3} - \alpha }\right) }\right\rbrack  .  S = \dfrac{1}{2} \times  \dfrac{{50}^{2}}{{\cos }^{2}\dfrac{\alpha }{2}} \times  \sqrt{3}\cos \alpha  = \dfrac{{50}^{2}\sqrt{3}\cos \alpha }{1 + \cos \alpha } = {2500}\sqrt{3}\lbrack 1 -  \left. \dfrac{1}{1 + \cos \alpha }\right\rbrack$ . 当且仅当 $\cos \alpha  = 1$ 且 $- \dfrac{\pi }{3} < \alpha  < \dfrac{\pi }{3}$ ,即 $\alpha  = 0$ 时, $S$ 取最大值,此时 ${OD} = {50}$ 米. 故使 $\angle {DOC} = \dfrac{\pi }{3}$ ,且 ${OD} = {50}$ 米, 可使花卉育苗区的面积最大.
        \end{enumerate}
    \end{jieda}

    \item (2024$\cdot$徐汇区一模$\cdot$19)2023 年，杭州亚运会首次启用机器狗搬运赛场上的运动装备. 如图所示， 在某项运动赛事扇形场地 ${OAB}$ 中, $\angle {AOB} = \dfrac{\pi }{2},{OA} = {500}{m}$ ，点 $Q$ 是弧 ${AB}$ 的中点， $P$ 为线段 ${OQ}$ 上一点(不与点 $O$ 、 $Q$ 重合). 为方便机器狗运输装备，现需在场地中铺设三条轨道 ${PO}$ 、 ${PA}$ 、${PB}$. 记 $\angle {APQ} = \theta$ ,三条轨道的总长度为 $y{m}$.

    \begin{center}
    \includegraphics[width=0.2\textwidth]{bodyxb/src/59_1117_1319_283_275_0.jpg}
    \end{center}

    \begin{enumerate}
        \item 将 $y$ 表示成 $\theta$ 的函数,并写出 $\theta$ 的取值范围
        \item 当三条轨道的总长度最小时,求轨道 ${PO}$ 的长
    \end{enumerate}
    \begin{jieda}
        \begin{enumerate}
            \item 因为点 $Q$ 是弧 ${AB}$ 的中点，由对称性，知 ${PA} = {PB}$ ， $\angle {AOP} = \angle {BOP} = \dfrac{\pi }{4}$ ,又 $\angle {APO} = \pi  - \theta ,\angle {OAP} = \theta  - \dfrac{\pi }{4}$ , ${OA} = {500}$ ,由正弦定理,得 $\dfrac{AP}{\sin \dfrac{\pi }{4}} = \dfrac{OA}{\sin \left( {\pi  - \theta }\right) } = \dfrac{OP}{\sin \left( {\theta  - \dfrac{\pi }{4}}\right) }$ ,所以 ${AP} = \dfrac{{250}\sqrt{2}}{\sin \theta },{OP} = \dfrac{{500}\sin \left( {\theta  - \dfrac{\pi }{4}}\right) }{\sin \theta }$ . 所以 $y = {AP} + {BP} +  {OP} = {2AP} + {OP} = \dfrac{{500}\sqrt{2} + {500}\sin \left( {\theta  - \dfrac{\pi }{4}}\right) }{\sin \theta } = {250}\sqrt{2} \cdot \dfrac{2 + \sin \theta  - \cos \theta }{\sin \theta }$ ,因为 $\angle {APQ} > \angle {AOP}$ ,所以 $\theta  > \dfrac{\pi }{4},\angle {AQO} =  \angle {OAQ} = \dfrac{1}{2}\left( {\pi  - \dfrac{\pi }{4}}\right)  = \dfrac{3}{8}\pi$ ,所以 $\theta  \in  \left( {\dfrac{\pi }{4},\dfrac{5\pi }{8}}\right)$
            \item 
            \begin{dingautolist}{172}
                \item 由(1)得: $y = {250}\sqrt{2} + {250}\sqrt{2} \cdot \dfrac{2 - \cos \theta }{\sin \theta },\theta  \in  \left( {\dfrac{\pi }{4},\dfrac{5\pi }{8}}\right)$ . 记 $t = \dfrac{2 - \cos \theta }{\sin \theta }$ ,则 $t\sin \theta  + \cos \theta  = 2$ ,由辅助角公式可得: $\sqrt{{t}^{2} + 1}\sin \left( {\theta  + \varphi }\right)  = 2 \Longrightarrow  \sin \left( {\theta  + \varphi }\right)  = \dfrac{2}{\sqrt{{t}^{2} + 1}} \leq  1$ ,解得 $t \geq  \sqrt{3}$ ,当 $t = \sqrt{3}$ 时,可有 $\sin \left( {\theta  + \dfrac{\pi }{6}}\right)  = 1 \Longrightarrow  \theta  = \dfrac{\pi }{3} \in  \left( \dfrac{\pi }{4}\right.$ , $\left. \dfrac{5\pi }{8}\right)$ ,等号可以取得. 故当 $\theta  = \dfrac{\pi }{3}$ 时,三条轨道的总长度最小,此时 ${OP} = \dfrac{250}{3}\left( {3\sqrt{2} - \sqrt{6}}\right)$.
                
                \ps 可以先通过斜率法判断可以在相切的位置取到最值，再采用辅助角公式解决.
                \item 由 (1) 得: $y = {250}\sqrt{2} + {250}\sqrt{2} \cdot  \dfrac{2 - \cos \theta }{\sin \theta },\theta  \in  \left( \dfrac{\pi }{4}\right.$ , $\left. \dfrac{5\pi }{8}\right)$ . 记 $t = \dfrac{2 - \cos \theta }{\sin \theta },x = \tan \dfrac{\theta }{2} \in  \left( {\tan \dfrac{\pi }{8},\tan \dfrac{5\pi }{16}}\right)$ ,则由万能置换公式可得: $t = \dfrac{2 - \dfrac{1 - {x}^{2}}{1 + {x}^{2}}}{\dfrac{2x}{1 + {x}^{2}}} = \dfrac{3{x}^{2} + 1}{2x} = \dfrac{1}{2}\left( {{3x} + \dfrac{1}{x}}\right)  \geq  \sqrt{{3x} \cdot  \dfrac{1}{x}} = \sqrt{3}$ ,当且仅当 $x = \dfrac{\sqrt{3}}{3}$ 即 $\theta  = \dfrac{\pi }{3}$ 时等号成立. 故当 $\theta  = \dfrac{\pi }{3}$ ,三条轨道的总长度最小,此时 ${OP} = \dfrac{250}{3}\left( {3\sqrt{2} - \sqrt{6}}\right)$.
                \item 令 $f\left( \theta \right)  = \dfrac{2 + \sin \theta  - \cos \theta }{\sin \theta },\theta  \in  \left( {\dfrac{\pi }{4},\dfrac{5\pi }{8}}\right)$ . 由 ${f}^{\prime }\left( \theta \right)  =  \dfrac{1 - 2\cos \theta }{{\sin }^{2}\theta } = 0$ ,解得 $\theta  = \dfrac{\pi }{3}$ ,则有

                \begin{center}
                \begin{tabular}{|c|c|c|c|}
                \hline
                $\theta$ & $\dfrac{\pi }{4} < \theta  < \dfrac{\pi }{3}$ & $\theta  = \dfrac{\pi }{3}$ & $\dfrac{\pi }{3} < \theta  < \dfrac{5\pi }{8}$ \\
                \cline{1-4}
                ${f}^{\prime }\left( \theta \right)$ & $< 0$ & $= 0$ & >0 \\
                \cline{1-4}
                $f\left( \theta \right)$ & 严格减 & 极小值 & 严格增 \\
                \cline{1-4}
                \hline
                \end{tabular}
                \end{center}

                所以当 $\theta  = \dfrac{\pi }{3}$ ,即 ${OP} = \dfrac{250}{3}\left( {3\sqrt{2} - \sqrt{6}}\right)$ 米时, $f\left( \theta \right)$ 有唯一的极小值,即是最小值,则 $f{\left( \theta \right) }_{\min } = f\left( \dfrac{\pi }{3}\right)$ ,三条轨道总长度的最小值为 ${250}\sqrt{6} + {250}\sqrt{2}$ . 故当 $\theta  = \dfrac{\pi }{3}$ 时,三条轨道的总长度最小,此时 ${OP} = \dfrac{250}{3}\left( {3\sqrt{2} - \sqrt{6}}\right)$ .
            \end{dingautolist}
        \end{enumerate}
    \end{jieda}
    
    \item (2024$\cdot$杨浦区一模$\cdot$19)某数学建模小组研究挡雨棚(图 1)，将它抽象为柱体(图 2)，平面 ${ABC}$ : 与 ${A}_{1}{B}_{1}{C}_{1}$ 全等且所在平面平行, $\bigtriangleup {ABC}$ 与 $\bigtriangleup {A}_{1}{B}_{1}{C}_{1}$ 各边表示挡雨棚支架,支架 $A{A}_{1}\text{ 、 }B{B}_{1}$ 、 $C{C}_{1}$ 垂直于平面 ${ABC}$. 雨滴下落方向与外墙 (所在平面) 所成角为 $\dfrac{\pi }{6}\left({\text{即}\angle {AOB} = \dfrac{\pi }{6}}\right)$ ,挡雨棚有效遮挡的区域为矩形 $A{A}_{1}{O}_{1}O.\left( {O\text{ 、 }{O}_{1}}\right.$ 分别在 ${CA}\text{ 、 }{C}_{1}{A}_{1}$ 延长线上

    \begin{center}
    \includegraphics[width=0.6\textwidth]{bodyxb/src/60_299_432_952_426_0.jpg}
    \end{center}
    \begin{enumerate}
        \item 挡雨板(曲面 $B{B}_{1}{C}_{1}C$ )的面积可以视为曲线段 ${BC}$ 与线段 $B{B}_{1}$ 长的乘积. 已知 ${OA} = {1.5}{m}$ ， ${AC} = {0.3}{m},A{A}_{1} = 2{m}$ ,小组成员对曲线段 ${BC}$ 有两种假设,分别为:  \ding{172}  其为直线段且 $\angle {ACB} = \dfrac{\pi }{3}$ ;  \ding{173}  其为以 $O$ 为圆心的圆弧. 请分别计算这两种假设下挡雨板的面积 (精确到 0.1 ${{m}}^{2}$ )
        \item 小组拟自制 $\bigtriangleup {ABC}$ 部分的支架用于测试(图 3)，其中 ${AC} = {0.6}{m}$ ， $\angle {ABC} = \dfrac{\pi }{2}$ ， $\angle {CAB} = \theta$ ， 其中 $\dfrac{\pi }{6} < \theta  < \dfrac{\pi }{2}$ ，求有效遮挡区域高 ${OA}$ 的最大值.
    \end{enumerate}
    \begin{jieda}
        \begin{enumerate}
            \item \ding{172}  其为直线段且 $\angle {ACB} = \dfrac{\pi }{3}$ 时， ${AC} = {0.3}$ 米，所以在 Rt $\bigtriangleup {ACB}$ 中, ${AC} = {BC}\cos \dfrac{\pi }{3}$ ,即 ${BC} = {2AC} = {0.6}$ (米). 所以 $S = {BC} \cdot  B{B}_{1} = {BC} \cdot  A{A}_{1} = {0.6} \times  2 = {1.2}$ (平方米)

            \ding{173}  其为以 $O$ 为圆心的圆弧时,此时圆的半径为 ${OA} + {AC} =$ 1.8(米),圆心角 $\angle {AOB} = \dfrac{\pi }{6}$ ,所以圆弧 ${BC}$ 的长 $l = {1.8} \times  \dfrac{\pi }{6} = \dfrac{3\pi }{10}$ ,所以 $S = l \cdot  B{B}_{1} = l \cdot  A{A}_{1} = 2 \times  \dfrac{3\pi }{10} \approx  {1.9}$ (平方米)

            \item 由题意， ${AB} = {AC}\cos \theta  = 0.6\cos \theta$ ， $\angle {ABO} = \theta  - \dfrac{\pi }{6}$ ，由正弦定理可得: $\dfrac{OA}{\sin \left( {\theta  - \dfrac{\pi }{6}}\right) } = \dfrac{AB}{\sin \dfrac{\pi }{6}}$ ，即 ${OA} = {{1.2}\cos }\theta \sin \left( {\theta  - \dfrac{\pi }{6}}\right)  = 1.2\cos \theta \left( {\dfrac{\sqrt{3}}{2}\sin \theta  - \dfrac{1}{2}\cos \theta }\right)  = 1.2\left( {\dfrac{\sqrt{3}}{2}\sin \theta \cos \theta  - \dfrac{1}{2}{\cos }^{2}\theta }\right)  =$ 0.6 $\left( {\dfrac{\sqrt{3}}{2}\sin {2\theta } - \dfrac{1}{2}\cos {2\theta } - \dfrac{1}{2}}\right)  = {0.6}\sin \left( {{2\theta } - \dfrac{\pi }{6}}\right)  - {0.3}$ ,其中 $\dfrac{\pi }{6} < \theta  < \dfrac{\pi }{2}$ ,当 ${2\theta } - \dfrac{\pi }{6} = \dfrac{\pi }{2}$ ,即 $\theta  = \dfrac{\pi }{3}$ 时, $O{A}_{\max } = 0.6 - 0.3 =$ 0.3 (米). 即有效遮挡区域高 ${OA}$ 的最大值为 0.3 米.
        \end{enumerate}
    \end{jieda}

    \item (2024$\cdot$静安区一模$\cdot$20)如图所示，某公园东北角处有一座小山，山顶有一根垂直于水平地平面的钢制笔直旗杆 ${AB}$ ，公园内的小山下是一个水平广场(虚线部分). 某高三班级数学老师留给同学们的周末作业是:进入该公园，提出与测量有关的问题，在广场上实施测量，并运用数学知识解决问题. 老师提供给同学们的条件是:已知 ${AB} = {10}{m}$ ，规定使用的测量工具只有一只小小的手持激光测距仪. (如下图，该测距仪能准确测量它到它发出的激光投射在物体表面上的光点之间的距离)

    \begin{center}
    \includegraphics[width=0.7\textwidth]{bodyxb/src/61_305_444_923_330_0.jpg}
    \end{center}
    
    \begin{enumerate}
        \item 甲同学来到通往山脚下的笔直小路 $l$ 上，他提出的问题是:如何测量小山的高度？干是，他站在点 $C$ 处,独立地实施了测量,并运用数学知识解决了问题. 请写出甲同学的解决问题方案,并用假设的测量数据(字母表示)表示出小山的高度 $H$ 
        \begin{jieda}
            \begin{dingautolist}{172}
                \item 设点 $A$ 在水平面的投影点为 $O$ . 用测距仪测得 ${CA} = m,{CB} = n$ . 在 $\bigtriangleup {ABC}$ 中， $\cos \angle {BAC} =  \dfrac{{100} + {m}^{2} - {n}^{2}}{20m}$ ,在 $\bigtriangleup {AOC}$ 中, $\cos \angle {OAC} =  - \dfrac{{100} + {m}^{2} - {n}^{2}}{20m}$ , 所以 $H = m\cos \angle {OAC} = \dfrac{{n}^{2} - {m}^{2} - {100}}{20}$ .
                \item 在平面 ${ABC}$ 上,以点 $C$ 为原点,向量 $\vv{CO}$ 为 $x$ 轴,建立平面直角坐标系 ${xCy}$ ,设点 $A\left( {x,H}\right)$ ,则 $B(x,H +$ 10),用测距仪测得 ${CA} = m,{CB} = n$ ，则 $\left\{  \begin{array}{l} {x}^{2} + {H}^{2} = {m}^{2} \\  {x}^{2} + {\left( H + {10}\right) }^{2} = {n}^{2} \end{array}\right.$ ， 解得 $H = \dfrac{{n}^{2} - {m}^{2} - {100}}{20}$
            \end{dingautolist}
        \end{jieda}
        \item 乙同学是在一阵大风过后进入公园的. 广场上的人纷纷议论:旗杆 ${AB}$ 似乎是由于在根部 $A$ 处松动产生了倾斜. 她提出的问题是: 如何检验旗杆 ${AB}$ 是否还垂直于地面? 并且设计了一个不用计算就能解决问题的独立测量方案. 请你写出她的方案, 并说明理由.
        \begin{jieda}
            用测距仪测得 ${CA} = m,{CB} = n$ ，在广场上从点 $C$ 移动至点 $D$ ，使得 ${DB} = n$ ，再移至点 $E$ ，使得 ${EB} = n$ ，此时再测量 ${DA}$ 、 ${EA}$ ，若 ${CA} = {DA} = {EA}$ ，则可知旗杆 ${AB}$ 垂直于地面，否则就是倾斜了. 理由如下:已知 ${CB} = {DB}$ ， ${CA} = {DA}$ ， 设点 $M$ 是 ${CD}$ 的中点，则在等腰 $\bigtriangleup  {CBD}$ 中， ${BM}\bot {CD}$ . 同理 ${AM} \bot  {CD}$ ,又 ${AB}\text{ 、 }{BM} \subset$ 平面 ${ABM}$ ,所以 ${CD} \bot$ 平面 ${ABM}$ ; 又因为 ${AB} \subset$ 平面 ${ABM}$ ,故 ${AB} \bot  {CD}$ . 同理可证 ${AB} \bot  {DE}$ . 综上所述,旗杆 ${AB}$ 垂直于地面.
            \begin{center}
            \includegraphics[width=0.2\textwidth]{bodyxb/src/42_649_1431_205_197_0.jpg}
            \end{center}
        \end{jieda}
        %\item 已知(1)中的小路 $l$ 是东西方向，且与点 $A$ 所确定的平面垂直于地平面. 又已知在(2)中的乙同学已经断定旗杆 ${AB}$ 大致向广场方向倾斜. 如果你是该班级的同学，你会提出怎样的有实际意义的问题? 请写出实施测量与解决问题的方案, 并说明理由. (如果需要, 可通过假设的测量数据或运算结果列式说明，不必计算)
    \end{enumerate}
\end{enumerate}
\subsection*{考点7 三角函数的图像与性质}

\subsubsection*{一、填空题}
\begin{enumerate}
    \item (2025$\cdot$宝山区一模$\cdot$2)函数 $y = \cos {2x} + 1$ 的最小正周期为\blkda[2]{$\pi$}.

    \item (2025$\cdot$杨浦区一模$\cdot$2)函数 $y = \sin {2x}$ 的最小正周期是\blkda[2]{$\pi$}.

    \item (2024$\cdot$嘉定区一模$\cdot$3)函数 $y = \sin {\pi x}$ 的最小正周期为\blkda[2]{2}.

    \item (2025$\cdot$上海秋季卷$\cdot$5)函数 $y = \cos x$ 在 $\left\lbrack  {-\dfrac{\pi }{2},\dfrac{\pi }{4}}\right\rbrack$ 上的值域为\blkda[2]{$\left\lbrack  {0,1}\right\rbrack$}.

    \item (2024$\cdot$普陀区一模$\cdot$6)若函数 $y = \tan {3x}$ 在区间 $\left( {m,\dfrac{\pi }{6}}\right)$ 上是严格增函数，则实数 $m$ 的取值范围为\blkda[2]{$\left\lbrack  {-\dfrac{\pi }{6},\dfrac{\pi }{6}}\right)$}.

    \item (2025$\cdot$崇明区二模$\cdot$6)函数 $y = 2\sin \left( {{\omega x} - \dfrac{\pi }{16}}\right) \left( {\omega  > 0}\right)$ 的最小正周期是 $\pi$ ，则 $\omega  =$ \blkda[2]{2}.

    \item (2025$\cdot$静安区二模$\cdot$8)设某港口水的深度 $y\left( {m}\right)$ 关于时间 $t\left( {h}\right)$ 的函数近似满足 $y = h + A\sin \left( {{\omega t} + \varphi }\right)$. 根据某一天的测量，港口水的深度在早上 3 点达到最大值 ${18}{m}$ ，之后持续减少，并在上午 9 点达到最小值 ${14}{m}$. 则该港口水的深度 $y\left( {m}\right)$ 关于时间 $t\left( {h}\right)$ 的函数的近似表达式为\blkda[2]{$y = 2\sin \dfrac{\pi }{6}t + {16}$}.
    \begin{jieda}
        由题意可知 $\left\{  \begin{array}{l} h + A = {18} \\  h - A = {14} \end{array}\right.$ ,解得 $\left\{  \begin{array}{l} h = {16} \\  A = 2 \end{array}\right. ,\dfrac{T}{2} = 6$ ,所以 $T = \dfrac{2\pi }{\omega } = {12}$ ,所以 $\omega  = \dfrac{\pi }{6}$ ,所以 $y = {16} + 2\sin \left( {\dfrac{\pi }{6}t + \varphi }\right)$ ,又当 $t = 3$ 时,函数取得最大值 18,所以 $\sin \left( {\dfrac{\pi }{2} + \varphi }\right)  = 1$ ,所以 $\dfrac{\pi }{2} +  \varphi  = \dfrac{\pi }{2} + {2k\pi },k \in  \mathbf{Z}$ ,所以 $\varphi  = {2k\pi },k \in  \mathbf{Z}$ ,所以 $y = {16} + 2\sin \left( {\dfrac{\pi }{6}t}\right.  + {2k\pi }) = 2\sin \dfrac{\pi }{6}t + {16}$.
    \end{jieda}

    \item (2025$\cdot$金山区一模$\cdot$9)在$(0,2\pi)$内，使 $\sin x \geq  \cos x$ 成立的 $x$ 的取值范围是\blkda[2]{$\left\lbrack  {\dfrac{\pi }{4},\dfrac{5\pi }{4}}\right\rbrack$}.

    \item (2024$\cdot$金山区一模$\cdot$9)已知函数 $y = \sin \left( {\omega x}\right) \left( {\omega  > 0}\right)$ 在区间 $\left\lbrack  {0,\pi }\right\rbrack$ 上是严格增函数，且其图像关于点 $\left( {{4\pi },0}\right)$ 对称，则 $\omega$ 的值为\blkda[2]{$\dfrac{1}{4}$ 或 $\dfrac{1}{2}$}.
    \begin{jieda}
        因为 $x \in  \left\lbrack  {0,\pi }\right\rbrack$ ,则 ${\omega x} \in  \left\lbrack  {0,{\omega \pi }}\right\rbrack$ ,函数 $y = \sin \left( {\omega x}\right)  \left( {\omega  > 0}\right)$ 在区间 $\left\lbrack  {0,\pi }\right\rbrack$ 上是严格增函数,所以 $0 < {\omega \pi } \leq  \dfrac{\pi }{2}$ ,即 $0 < \omega  \leq  \dfrac{1}{2}$ ; 又因为 $y = \sin \left( {\omega x}\right)$ 的图像关于点 $\left( {{4\pi },0}\right)$ 对称,则 ${\omega x} = {k\pi }\left( {k \in  \mathbf{Z}}\right)$ ,则 $x = \dfrac{k\pi }{\omega }\left( {k \in  \mathbf{Z}}\right)$ ,所以 ${4\pi } = \dfrac{k\pi }{\omega }\left( {k \in  \mathbf{Z}}\right)$ ,解得 $\omega  = \dfrac{k}{4}\left( {k \in  \mathbf{Z}}\right)$ ,结合 $0 < \omega  \leq  \dfrac{1}{2}$ ,所以 $\omega  = \dfrac{1}{4}$ 或 $\dfrac{1}{2}$.
    \end{jieda}

    \item (2025$\cdot$普陀区二模$\cdot$9)设 $k \geq  1,k \in  \mathbf{N},0 < \varphi  < \dfrac{\pi }{2}$ ，函数 $y = f\left( x\right)$ 的表达式为 $f\left( x\right)  = 2\sin \left( {{kx} + \varphi }\right)$ ， 则对任意的实数 $a$ 皆有 $\{ f\left( x\right)  \mid  a < x < a + 2\}  = \{ f\left( x\right)  \mid  x \in  \mathbf{R}\}$ 成立的一个充分条件是\blkda[2]{$k = 4$}.
    \begin{jieda}
        因为函数 $f\left( x\right)  = 2\sin \left( {{kx} + \varphi }\right)$ ,要使 $\{ f\left( x\right)  \mid  a < x <  a + 2\}  = \{ f\left( x\right)  \mid  x \in  \mathbf{R}\}$ ,则周期 $T = \dfrac{2\pi }{k} \leq  2$ ,即 $k \geq  \pi  \approx  {3.14}$ , 因为 $k \in  \mathbf{N}$ ,所以一个充分条件是 $k = 4$
    \end{jieda}

    \item (2024$\cdot$黄浦区一模$\cdot$10)若 $\varphi$ 是一个三角形的内角，且函数 $y = 3\sin \left( {{2x} + \varphi }\right)$ 在区间 $\left\lbrack  {-\dfrac{\pi }{4},\dfrac{\pi }{6}}\right\rbrack$ 上是单调函数，则 $\varphi$ 的取值范围是\blkda[2]{$0 < \varphi  \leq  \dfrac{\pi }{6}$}.
    \begin{jieda}
        $x \in  \left\lbrack  {-\dfrac{\pi }{4},\dfrac{\pi }{6}}\right\rbrack  ,{2x} + \varphi  \in  \left\lbrack  {\varphi  - \dfrac{\pi }{2},\varphi  + \dfrac{\pi }{3}}\right\rbrack  ,\varphi  \in  \left( {0,\pi }\right)  \Longrightarrow  \varphi  - \dfrac{\pi }{2} \in  \left( {-\dfrac{\pi }{2},\dfrac{\pi }{2}}\right) ,\varphi  + \dfrac{\pi }{3} \in  \left( {\dfrac{\pi }{3},\dfrac{4\pi }{3}}\right)$ ,函数 $y =  3\sin \left( {{2x} + \varphi }\right)$ 在区间 $\left\lbrack  {-\dfrac{\pi }{4},\dfrac{\pi }{6}}\right\rbrack$ 上是单调函数 $\Longrightarrow  \varphi  + \dfrac{\pi }{3} \leq  \dfrac{\pi }{2},0 < \varphi  \leq  \dfrac{\pi }{6}.$
    \end{jieda}

    \item (2024$\cdot$杨浦区一模$\cdot$10)函数 $f\left( x\right)  = \cos \left( {{\omega x} + \varphi }\right) \varphi  \in  \left( {0,{2\pi }}\right)$ 在 $x \in  \mathbf{R}$ 上是单调增函数，且图像关于原点对称，则满足条件的数对 $\left( {\omega ,\varphi }\right)$ 为\blkda[2]{$\left( {0,\dfrac{\pi }{2}}\right) ,\left( {0,\dfrac{3\pi }{2}}\right)$}.
    \begin{jieda}
        当 $\omega  \neq  0$ 时, $f\left( x\right)  = \cos \left( {{\omega x} + \varphi }\right) \varphi  \in  \left( {0,{2\pi }}\right)$ 在 $x \in  \mathbf{R}$ 上必有增有减,不合题意,故 $\omega  = 0$ ,此时 $f\left( x\right)  = \cos \varphi ,\varphi  \in  \left( {0,{2\pi }}\right)$ 为常值函数,由其图像关于原点对称,所以 $f\left( x\right)  =  \cos \varphi  = 0$ ,所以 $\varphi  = \dfrac{\pi }{2}$ 或 $\dfrac{3\pi }{2}$ ,故满足条件的数对为 $\left( {0,\dfrac{\pi }{2}}\right) ,\left(0, \dfrac{3\pi }{2}\right)$
    \end{jieda}

    \item (2024$\cdot$嘉定区二模$\cdot$10)已知 $f\left( x\right)  = \dfrac{2}{\sin x} + \dfrac{2}{\cos x},x \in  \left( {0,\dfrac{\pi }{2}}\right)$ ，则函数 $y = f\left( x\right)$ 的最小值为\blkda[2]{$4\sqrt{2}$}.
    \begin{jieda}
        $f\left( x\right)  = \dfrac{2}{\sin x} + \dfrac{2}{\cos x} = \dfrac{2\left( {\sin x + \cos x}\right) }{\sin x\cos x}$ ,令 $t = \sin x +  \cos x,t \in  (1,\sqrt{2}\rbrack ,\dfrac{2t}{\dfrac{{t}^{2} - 1}{2}} = \dfrac{4t}{{t}^{2} - 1} = \dfrac{4}{t - \dfrac{1}{t}}$ ,当 $t = \sqrt{2}$ 时 $f\left( x\right)$ 取得最小值 $4\sqrt{2}$.
    \end{jieda}

    \item (2024$\cdot$奉贤区二模$\cdot$12)函数 $y = \sin \left( {{wx} + \varphi }\right) \left( {w > 0,\left| \varphi \right|  < \dfrac{\pi }{2}}\right)$ 的图像记为曲线 $F$. 如图所示, $A\text{ 、 }B\text{ 、 }C$ 是曲线 $F$ 与坐标轴相交的三个点,直线 ${BC}$ 与曲线 $F$ 的图像交于点 $M$ ,若直线 ${AM}$ 的斜率为 ${k}_{1}$ ,直线 ${BM}$ 的斜率为 ${k}_{2},{k}_{2} \neq  2{k}_{1}$ ，则直线 ${AB}$ 的斜率为\blkda[2]{$\dfrac{{k}_{1}{k}_{2}}{2{k}_{1} - {k}_{2}}$}. (${k}_{1}$ 、 ${k}_{2}$ 表示).
    \begin{center}
    \includegraphics[width=0.2\textwidth]{bodyxb/src/62_1078_1724_331_250_0.jpg}
    \end{center}
    \begin{jieda}
        设$\vv{MC} = (m, k_2m), \vv{MA} = (n, k_1n)$，则$\vv{MB} = (2m, 2k_2m)$，$\vv{AB} = (2m - n, k_2m)$且有$k_1n = k_2m$，故$k_{AB} = \dfrac{k_2m}{2m - n} = \dfrac{k_2m}{2m - \dfrac{k_2m}{k_1}} = \dfrac{{k}_{1}{k}_{2}}{2{k}_{1} - {k}_{2}}$
    \end{jieda}
\end{enumerate}

\subsubsection*{二、选择题}
\begin{enumerate}
    \item (2025$\cdot$金山区一模$\cdot$13)函数 $y = 1 - 2{\sin }^{2}x$ 是 \dotfill{(\kh{B})}
    \xx{最小正周期为 $\pi$ 的奇函数}{最小正周期为 $\pi$ 的偶函数}{最小正周期为 $\dfrac{\pi }{2}$ 的奇函数}{最小正周期为 $\dfrac{\pi }{2}$ 的偶函数}
    \item (2024$\cdot$静安区二模$\cdot$13)函数 $y = 2\sin x - \cos x\left( {x \in  \mathbf{R}}\right)$ 的最小正周期为 \dotfill{(\kh{A})}
    \xx{${2\pi }$}{$\pi$}{$\dfrac{3\pi }{2}$}{$\dfrac{\pi }{2}$}

    \item (2025$\cdot$长宁区一模$\cdot$15)已知函数 $y = \sin \left( {{\omega x} + \dfrac{\pi }{6}}\right) \left( {\omega  > 0}\right)$ 在区间 $\left( {-\dfrac{\pi }{2},\dfrac{\pi }{3}}\right)$ 上单调递增，则 $\omega$ 的取值范围是 \dotfill{(\kh{A})}
    \xx{$(0,1\rbrack$}{$\left( {0,1}\right)$}{$\left( {1,\dfrac{4}{3}}\right)$}{$\left( {0,\dfrac{6}{5}}\right\rbrack$}
    \begin{jieda}
        $x \in  \left( {-\dfrac{\pi }{2},\dfrac{\pi }{3}}\right)  \Longrightarrow  {\omega x} + \dfrac{\pi }{6} \in  \left( {-\dfrac{\pi }{2}\omega  + \dfrac{\pi }{6},\dfrac{\pi }{3}\omega  + }\right.  \left. \dfrac{\pi }{6}\right)$ ,因 $y = \sin x$ 在 $\left( {-\dfrac{\pi }{2} + {2k\pi },\dfrac{\pi }{2} + {2k\pi }}\right)$ (其中 $k \in  \mathbf{Z}$ ) 上单调递增,则 $\left\{  {\begin{array}{l} \dfrac{\pi }{3}\omega  + \dfrac{\pi }{6} \leq  \dfrac{\pi }{2} + {2k\pi } \\   - \dfrac{\pi }{2}\omega  + \dfrac{\pi }{6} \geq   - \dfrac{\pi }{2} + {2k\pi } \end{array} \Longrightarrow  \left\{  {\begin{array}{l} \omega  \leq  1 + {6k} \\  \omega  \leq  \dfrac{4}{3} - {4k} \end{array},k \in \mathbf{Z}}\right. }\right.$. 又因 $\omega  > 0$ ,则取 $k = 0 \Longrightarrow  \omega  \leq  1$ ,则 $0 < \omega  \leq  1$
    \end{jieda}

    \item (2023$\cdot$上海秋季卷$\cdot$15)已知 $a > 0$ ，函数 $y = \sin x$ 在区间 $\left\lbrack  {a,{2a}}\right\rbrack$ 上的最小值为 $s$ ，在 $\left\lbrack  {{2a},{3a}}\right\rbrack$ 上的最小值为 $t$ ,当 $a$ 变化时,下列不可能的是 \dotfill{(\kh{D})}
    \xx{$s > 0,t > 0$}{$s > 0,t < 0$}{$s < 0,t < 0$}{$s < 0,t > 0$}
    \begin{jieda}
        $a = \dfrac{\pi }{4}$ 时, ${A}$ 可能; $a = \dfrac{5\pi }{12}$ 时, ${B}$ 可能; $a = \dfrac{7\pi }{12}$ 时, ${C}$ 可能；D选项，若 $s < 0$ ，则 ${2a} > \pi$ ，若 $t > 0$ ，则 $\left\lbrack  {{2a},{3a}}\right\rbrack$ 的区间长度 $a < \pi$ ，同时 $\sin {2a} > 0$ 且 $\sin {3a} > 0$ ，所以 ${2a} \in  \left( {0,\pi }\right)$ 且 ${3a} \in  \left( {0,\pi }\right)$ ,与前面的 ${2a} > \pi$ 矛盾,故 ${D}$ 不可能.
    \end{jieda}

    \item (2025$\cdot$上海春季卷$\cdot$16)已知 $a \in  \mathbf{R}$ ，不等式 $\left\lbrack  {\tan \left( {\dfrac{\pi }{6}x}\right)  - a}\right\rbrack   \cdot  \left\lbrack  {\tan \left( {\dfrac{\pi }{6}x}\right)  - a - 1}\right\rbrack   < 0$ 在$(0,2025)$中的整数解有 $m$ 个. 关于 $m$ 的个数,以下不可能的是 \dotfill{(\kh{D})}
    \xx{0}{338}{674}{1012}
    \begin{jieda}
        由题意,得 $a < \tan \left( {\dfrac{\pi }{6}x}\right)  < a + 1$ ,设 $f\left( x\right)  =  \tan \left( {\dfrac{\pi }{6}x}\right)$ ,则 $T = \dfrac{\pi }{\dfrac{\pi }{6}} = 6$ ,因为 ${2025} \div  6 = {337}\cdots 3$ ,所以有 337 个完整周期及多出来 $x \in  \left( {0,3}\right)$ 上的图像,在一个周期内整数解的函数值为: $f\left( 1\right)  = \dfrac{\sqrt{3}}{3},f\left( 2\right)  = \sqrt{3},f\left( 4\right)  =  - \sqrt{3},f\left( 5\right)  =  - \dfrac{\sqrt{3}}{3},f\left( 6\right)  = 0$ ,当 $a \geq  \sqrt{3}$ 时,因为 $\{ x \in  \mathbf{Z} \mid  f\left( x\right) \}  \notin  \left( {a,a + 1}\right)$ , 所以不等式无解,即 $m = 0$ ; 当 $\dfrac{\sqrt{3}}{3} < a < \sqrt{3} < a + 1$ 时, $m = {338}$ ; 当 $a <  - \dfrac{\sqrt{3}}{3} < 0 < a + 1$ 时, $m = 2 \times  {337} = {674}$ ; 由排除法,选 D.

        \begin{center}
        \includegraphics[width=0.9\textwidth]{bodyxb/src/44_200_604_668_183_0.jpg}
        \end{center}
    \end{jieda}

    \item (2024$\cdot$闵行区二模$\cdot$16)已知 $f\left( x\right)  = \sin x$ ，集合 $D = \left\lbrack  {-\dfrac{\pi }{2},\dfrac{\pi }{2}}\right\rbrack  ,\Gamma  = \left\{  {\left( {x,y}\right) \left| {{2f}\left( x\right)  + f\left( y\right)  = 0}\right. }\right.$ ， $y \in  D\} ,\Omega  = \{ \left( {x,y}\right)  \mid  {{2f}\left( x\right)  + f\left( y\right)  \geq  0},x,y \in  D\}$. 关于下列两个命题的判断，说法正确的是\dotfill{(\kh{A})} 命题 \ding{172} : 集合 $\Gamma$ 表示的平面图形是中心对称图形; 命题 \ding{173} : 集合 $\Omega$ 表示的平面图形的面积不大于 $\dfrac{5{\pi }^{2}}{12}$.
    \xx{\ding{172} 真命题， \ding{173} 假命题}{\ding{172} 假命题， \ding{173} 真命题}{\ding{172}  \ding{173} 均为真命题}{\ding{172}  \ding{173} 均为假命题}
    \begin{jieda}
        对于命题 \ding{172}  $\Gamma  = \{ \left( {x,y}\right)  {2f}\left( x\right)  + f\left( y\right)  = 0,x\text{ 、 }y \in  D\}$ ,以 $\left( {-x, - y}\right)$ 代入曲线仍成立，所以集合Γ表示的图形关于原点中心对称，故 \ding{172} 是真命题；

        对于命题 \ding{173} ，若$(x_0, y_0) \in \Omega$，则$(-x_0, -y_0) \notin \Omega$，若$(x_0, y_0) \notin \Omega$，则$(-x_0, -y_0) \in \Omega$，其图形面积应为$\dfrac{\pi^2}{2}$
    \end{jieda}

    \item (2024$\cdot$嘉定区二模$\cdot$16)已知函数 $y = f\left( x\right) \left( {x \in  \mathbf{R}}\right)$ 的最小正周期是 ${T}_{1}$ ，函数 $y = g\left( x\right) \left( {x \in  \mathbf{R}}\right)$ 的最小正周期是 ${T}_{2}$ ,且 ${T}_{1} = k{T}_{2}\left( {k > 1}\right)$ ,对于命题甲: 函数 $y = f\left( x\right)  + g\left( x\right) \left( {x \in  \mathbf{R}}\right)$ 可能不是周期函数; 命题乙: 若函数 $y = f\left( x\right)  + g\left( x\right) \left( {x \in  \mathbf{R}}\right)$ 的最小正周期是 ${T}_{3}$ ,则 ${T}_{3} \geq  {T}_{1}$. 下列选项正确的是 \dotfill{(\kh{C})}
    \xx{甲和乙均为真命题}{甲和乙均为假命题}{甲为真命题且乙为假命题}{甲为假命题且乙为真命题}
    \begin{jieda}
        对于命题甲: 如果两个周期函数相加,一般来说,它们的和函数的周期会是这两个周期的公倍数(如果存在公倍数的话)。对于周期函数,如果 ${T}_{1}$ 和 ${T}_{2}$ 的比值是有理数,那么它们存在公共周期；如果比值是无理数,通常和函数不是周期函数。
        
        $f\left( x\right)  = \cos \sqrt{2}{\pi x},g\left( x\right)  = \cos x,{T}_{1} =  \sqrt{2},{T}_{2} = {2\pi },y = \cos \sqrt{2}{\pi x} + \cos x$ 不是周期函数; 所以命题甲为真命题
        
        对于命题乙: $f\left( x\right)  = \cos {5x} + \cos x,g\left( x\right)  =  - \cos x,{T}_{1} = {2\pi },{T}_{2} = {2\pi },f\left( x\right)  + g\left( x\right)  = \cos {5x},{T}_{3} = \dfrac{2\pi }{5}$ ,所以命题乙为假命题;

        \ps 经典反例$y = f(x), y = -f(x)$，周期相同，相加后没有最小正周期.
    \end{jieda}
\end{enumerate}
\subsubsection*{三、解答题}
\begin{enumerate}
    \item (2025$\cdot$青浦区一模$\cdot$17)已知 $f\left( x\right)  = \left( {2{\cos }^{2}x - 1}\right) \sin {2x} + \dfrac{1}{2}\cos {4x}$.
    \begin{enumerate}
        \item 求函数 $y = f\left( x\right)$ 的最小正周期及最大值
        \item 若 $\alpha  \in  \left( {\dfrac{\pi }{2},\pi }\right)$ ，且 $f\left( \alpha \right)  = \dfrac{\sqrt{2}}{2}$ ，求 $\alpha$ 的值.
    \end{enumerate}
    \begin{jieda}
        \begin{enumerate}
            \item 因为 $f\left( x\right)  = \left( {2{\cos }^{2}x - 1}\right) \sin {2x} + \dfrac{1}{2}\cos {4x} =  \cos {2x}\sin {2x} + \dfrac{1}{2}\cos {4x} = \dfrac{1}{2}\left( {\sin {4x} + \cos {4x}}\right)  = \dfrac{\sqrt{2}}{2}\sin ({4x} +  \left. \dfrac{\pi }{4}\right)$ ,所以 $y = f\left( x\right)$ 的最小正周期为 $\dfrac{\pi }{2}$ ,最大值为 $\dfrac{\sqrt{2}}{2}$
            
            \item 因为 $\alpha  \in  \left( {\dfrac{\pi }{2},\pi }\right)$ ，所以 ${4\alpha } + \dfrac{\pi }{4} \in  \left( {\dfrac{9\pi }{4},\dfrac{17\pi }{4}}\right)$ . 因为 $f\left( \alpha \right)  =  \dfrac{\sqrt{2}}{2}$ ,所以 $f\left( \alpha \right)  = \dfrac{\sqrt{2}}{2}\sin \left( {{4\alpha } + \dfrac{\pi }{4}}\right)  = \dfrac{\sqrt{2}}{2}$ ,即 $\sin \left( {{4\alpha } + \dfrac{\pi }{4}}\right)  = 1$ . 所以 ${4\alpha } + \dfrac{\pi }{4} = \dfrac{5\pi }{2}$ ,故 $\alpha  = \dfrac{9\pi }{16}$ .
        \end{enumerate}
    \end{jieda}
    \item (2025$\cdot$浦东新区 一模$\cdot$17)已知函数 $y = f\left( x\right)$ 的表达式为 $f\left( x\right)  = \sin {\omega x},\omega  > 0$.
    \begin{enumerate}
        \item 若函数 $y = f\left( x\right)$ 的最小正周期为 $\dfrac{\pi }{2}$ ，求 $\omega$ 的值及 $y = f\left( x\right)$ 的单调增区间
        \item 若 $\omega  = 2$ ，设函数 $y = g\left( x\right)$ 的表达式为 $g\left( x\right)  = f\left( x\right)  + \sqrt{3}\cos {2x}$. 求当 $x \in  \left\lbrack  {0,\dfrac{\pi }{2}}\right\rbrack$ 时， $y =  g\left( x\right)$ 的值域.
    \end{enumerate}
    \begin{jieda}
        \begin{enumerate}
            \item 因为 $\omega  > 0$ ，所以 $\dfrac{2\pi }{\omega } = \dfrac{\pi }{2}$ ，解得 $\omega  = 4$ ， $f\left( x\right)  = \sin {4x}$ ，令 $- \dfrac{\pi }{2} + {2k\pi } \leq  {4x} \leq  \dfrac{\pi }{2} + {2k\pi },k \in  \mathbf{Z}$ ,解得 $- \dfrac{\pi }{8} + \dfrac{k\pi }{2} \leq  x \leq  \dfrac{\pi }{8} +  \dfrac{k\pi }{2},k \in  \mathbf{Z}$ ,故单调递增区间为 $\left\lbrack  {-\dfrac{\pi }{8} + \dfrac{k\pi }{2},\dfrac{\pi }{8} + \dfrac{k\pi }{2}}\right\rbrack  ,k \in  \mathbf{Z}$
            \item $\omega  = 2,g\left( x\right)  = \sin {2x} + \sqrt{3}\cos {2x} = 2\sin \left( {{2x} + \dfrac{\pi }{3}}\right) ,x \in  \left\lbrack  {0,\dfrac{\pi }{2}}\right\rbrack$ 时, ${2x} + \dfrac{\pi }{3} \in  \left\lbrack  {\dfrac{\pi }{3},\dfrac{4\pi }{3}}\right\rbrack$ ,故 $\sin \left( {{2x} + \dfrac{\pi }{3}}\right)  \in  \left\lbrack  {-\dfrac{\sqrt{3}}{2},1}\right\rbrack$ , 所以 $g\left( x\right)  = 2\sin \left( {{2x} + \dfrac{\pi }{3}}\right)  \in  \left\lbrack  {-\sqrt{3},2}\right\rbrack$ .
        \end{enumerate}
    \end{jieda}
    \item (2025$\cdot$徐汇区一模$\cdot$17)已知 $f\left( x\right)  = a\sin {\omega x} + b\cos {\omega x}\left( {\omega  > 0}\right)$ ，若定义在 $\mathbf{R}$ 上的函数 $y = f\left( x\right)$ 的最小正周期为 $\pi$ ,且对任意的 $x \in  \mathbf{R}$ ,都有 $f\left( x\right)  \leq  f\left( \dfrac{\pi }{12}\right)  = 4$.
    \begin{enumerate}
        \item 求实数 $a$ 、 $b$ 的值
        \item 设 ${x}_{1}\text{ 、 }{x}_{2} \in  \left( {0,\pi }\right)$ ，当 ${x}_{1} \neq  {x}_{2}$ 时， $f\left( {x}_{1}\right)  = f\left( {x}_{2}\right)  =  - 2$ ，求 ${x}_{1} + {x}_{2}$ 的值
    \end{enumerate}
    \begin{jieda}
        \begin{enumerate}
            \item $f\left( x\right)  = a\sin {\omega x} + b\cos {\omega x} = \sqrt{{a}^{2} + {b}^{2}}\sin \left( {{\omega x} + \varphi }\right)$ ,由 $y = f\left( x\right)$ 的最小正周期为 $\pi$ ,知 $\omega  = 2$ ,因为 $f\left( x\right)  \leq  f\left( \dfrac{\pi }{12}\right)  =  a\sin \dfrac{\pi }{6} + b\cos \dfrac{\pi }{6} = 4$ ,所以 $\left\{  {\begin{array}{l} \dfrac{1}{2}a + \dfrac{\sqrt{3}}{2}b = 4 \\  {a}^{2} + {b}^{2} = {16} \end{array} \Longrightarrow  \left\{  \begin{array}{l} a = 2 \\  b = 2\sqrt{3} \end{array}\right. }\right.$
            \item 由(1)可得: $f\left( x\right)  = 2\sin {2x} + 2\sqrt{3}\cos {2x} = 4\sin \left( {{2x} + }\right.  \left. \dfrac{\pi }{3}\right)$ ,因为 $4\sin \left( {{2x} + \dfrac{\pi }{3}}\right)  =  - 2 \Longrightarrow  \sin \left( {{2x} + \dfrac{\pi }{3}}\right)  =  - \dfrac{1}{2}$ ,所以 ${2x} + \dfrac{\pi }{3} = {2k\pi } - \dfrac{\pi }{6}$ 或 ${2x} + \dfrac{\pi }{3} = {2k\pi } + \dfrac{7}{6}\pi$ ,即 $x = {k\pi } -  \dfrac{\pi }{4}$ 或 $x = {k\pi } + \dfrac{5}{12}\pi ,k \in  \mathbf{Z}$ ,又 ${x}_{1},{x}_{2} \in  \left( {0,\pi }\right)$ ,则不妨令 ${x}_{1} =  \dfrac{3}{4}\pi ,{x}_{2} = \dfrac{5}{12}\pi$ ,故 ${x}_{1} + {x}_{2} = \dfrac{7}{6}\pi$ .
        \end{enumerate}
    \end{jieda}

    \item (2024$\cdot$静安区一模$\cdot$17)记 $f\left( x\right)  = {\sin }^{2}x - {\cos }^{2}x + 2\sqrt{3}\sin x\cos x + \lambda \left( {x \in  \mathbf{R}}\right)$ ，其中 $\lambda$ 为实常数.
    \begin{enumerate}
        \item 求函数 $y = f\left( x\right)$ 的最小正周期
        \item 若函数 $y = f\left( x\right)$ 的图像经过点 $\left( {\dfrac{\pi }{2},0}\right)$ ，求该函数在区间 $\left\lbrack  {0,\dfrac{2}{3}\pi }\right\rbrack$ 上的最大值和最小值.
    \end{enumerate}
    \begin{jieda}
        \begin{enumerate}
            \item $f\left( x\right)  =  - \cos {2x} + \sqrt{3}\sin {2x} + \lambda  = 2\sin \left( {{2x} - \dfrac{\pi }{6}}\right)  +  \lambda$ . 所以函数 $y = f\left( x\right)$ 的最小正周期为 $\pi$
            
            \item 因为 $f\left( \dfrac{\pi }{2}\right)  = 1 + \lambda  = 0$ ，所以 $\lambda  =  - 1$ ，则 $f\left( x\right)  =  2\sin \left( {{2x} - \dfrac{\pi }{6}}\right)  - 1$ . 令 ${2x} - \dfrac{\pi }{6} = t$ ,则 $t \in  \left\lbrack  {-\dfrac{\pi }{6},\dfrac{7\pi }{6}}\right\rbrack$ . 当 ${2x} - \dfrac{\pi }{6} =  - \dfrac{\pi }{6}$ 或 $\dfrac{7\pi }{6}$ ,即 $x = 0$ 或 $\dfrac{2\pi }{3}$ 时, $f{\left( x\right) }_{\min } =  - 2$ . 当 ${2x} - \dfrac{\pi }{6} = \dfrac{\pi }{2}$ ,即 $x = \dfrac{\pi }{3}$ 时, $f{\left( x\right) }_{\max } = 1$ .
        \end{enumerate}
    \end{jieda}

    \item (2025$\cdot$虹口区一模$\cdot$17)设 $f\left( x\right)  = \sin {\omega x}\left( {\omega  > 0}\right)$.
    \begin{enumerate}
        \item 当函数 $y = f\left( x\right)$ 的最小正周期为 ${2\pi }$ 时,求 $y = f\left( x\right)  + \cos x$ 在 $\left\lbrack  {0,\dfrac{\pi }{2}}\right\rbrack$ 上的最大值
        \item 若 $\omega  = 2$ ，且在 $\bigtriangleup  {ABC}$ 中，角 $A$ 、 $B$ 、 $C$ 所对的边长为 $a$ 、 $b$ 、 $c$ ，锐角 $A$ 满足 $f\left( {A + \dfrac{\pi }{6}}\right)  = 0$ ， $\vv{AB} \cdot  \vv{AC} = 4$ ，求 $a$ 的最小值.
    \end{enumerate}
    \begin{jieda}
        \begin{enumerate}
            \item 因为函数 $y = f\left( x\right)$ 的最小正周期为 ${2\pi }$ ，所以 $\omega  = 1$ ， 故 $y = f\left( x\right)  + \cos x = \sin x + \cos x = \sqrt{2}\sin \left( {x + \dfrac{\pi }{4}}\right)$ ,因为 $x \in  \left\lbrack  {0,\dfrac{\pi }{2}}\right\rbrack$ ,所以 $x + \dfrac{\pi }{4} \in  \left\lbrack  {\dfrac{\pi }{4},\dfrac{3\pi }{4}}\right\rbrack$ ,故当 $x + \dfrac{\pi }{4} = \dfrac{\pi }{2}$ ,即 $x = \dfrac{\pi }{4}$ 时,取到最大值 $\sqrt{2}$
            \item $f\left( x\right)  = \sin {2x}$ ，当 $f\left( {A + \dfrac{\pi }{6}}\right)  = 0$ 时， $\sin \left( {{2A} + \dfrac{\pi }{3}}\right)  = 0$ ， 即 ${2A} + \dfrac{\pi }{3} = {k\pi },k \in  \mathbf{Z}$ ，由于 $A$ 为锐角，解得 $A = \dfrac{\pi }{3}$ ，由 $\vv{AB}$ . $\vv{AC} = 4 \Longrightarrow  {cb}\cos \dfrac{\pi }{3} = 4 \Longrightarrow  {bc} = 8$ ，所以 ${a}^{2} = {b}^{2} + {c}^{2} - {2bc}\cos \dfrac{\pi }{3} \geq  {2bc} - {bc} = {bc} = 8$ ,等号当且仅当 $b = c = 2\sqrt{2}$ 时成立,此时 $a$ 的最小值为 $2\sqrt{2}$ .
        \end{enumerate}
    \end{jieda}

    \item (2024$\cdot$青浦区二模$\cdot$17)对于函数 $y = f\left( x\right)$ ，其中 $f\left( x\right)  = 2\sin x\cos x + 2\sqrt{3}{\cos }^{2}x - \sqrt{3},x \in  \mathbf{R}$.
    \begin{enumerate}
        \item 求函数 $y = f\left( x\right)$ 的单调增区间
        \item 在锐角三角形 ${ABC}$ 中，若 $f\left( A\right)  = 1,\vv{AB} \cdot  \vv{AC} = \sqrt{2}$ ，求 $\bigtriangleup  {ABC}$ 的面积.
    \end{enumerate}
    \begin{jieda}
        \begin{enumerate}
            \item $f\left( x\right)  = 2\sin x\cos x + 2\sqrt{3}{\cos }^{2}x - \sqrt{3} = 2\sin x\cos x +  \sqrt{3}\left( {2{\cos }^{2}x - 1}\right)  = \sin {2x} + \sqrt{3}\cos {2x} = 2\sin \left( {{2x} + \dfrac{\pi }{3}}\right)$ ,令 $u =  {2x} + \dfrac{\pi }{3}$ ,则 $f\left( u\right)  = 2\sin u$ ,函数 $u = {2x} + \dfrac{\pi }{3}$ 为严格增函数, 当 $u \in  \left\lbrack  {{2k\pi } - \dfrac{\pi }{2},{2k\pi } + \dfrac{\pi }{2}}\right\rbrack  ,k \in  \mathbf{Z}$ 时函数 $f\left( u\right)  = 2\sin u$ 为严格增函数,即 ${2k\pi } - \dfrac{\pi }{2} \leq  {2x} + \dfrac{\pi }{3} \leq  {2k\pi } + \dfrac{\pi }{2},k \in  \mathbf{Z}$ ,得 ${k\pi } - \dfrac{5\pi }{12} \leq  x \leq  {k\pi } + \dfrac{\pi }{12},k \in  \mathbf{Z}$ ,所以函数 $y = f\left( x\right)$ 的单调增区间是 $\left\lbrack  {{k\pi } - \dfrac{5\pi }{12},{k\pi } + \dfrac{\pi }{12}}\right\rbrack  ,\left( {k \in  \mathbf{Z}}\right)$
            
            \item 由已知 $f\left( A\right)  = 2\sin \left( {{2A} + \dfrac{\pi }{3}}\right)  = 1$ ，所以 $\sin \left( {{2A} + \dfrac{\pi }{3}}\right)  =  \dfrac{1}{2}$ ,因为 $0 < A < \dfrac{\pi }{2}$ ,所以 $\dfrac{\pi }{3} < {2A} + \dfrac{\pi }{3} < \dfrac{4\pi }{3}$ ,即 ${2A} + \dfrac{\pi }{3} =  \dfrac{5\pi }{6}$ ,所以 $A = \dfrac{\pi }{4}$ ,又 $\vv{AB} \cdot  \vv{AC} = \left| \vv{AB}\right|  \cdot  \left| \vv{AC}\right| \cos A = \sqrt{2}$ ,所以 $\left| \vv{AB}\right|  \cdot  \left| \vv{AC}\right|  = 2$ ,所以 $\bigtriangleup {ABC}$ 的面积 $S = \dfrac{1}{2}\left| \vv{AB}\right|$ . $\left| \vv{AC}\right| \sin A = \dfrac{1}{2} \times  2 \times  \dfrac{\sqrt{2}}{2} = \dfrac{\sqrt{2}}{2}$.
        \end{enumerate}
    \end{jieda}

    \item (2024$\cdot$金山区二模$\cdot$17)已知函数 $y = f\left( x\right)$ ，记 $f\left( x\right)  = \sin \left( {{\omega x} + \varphi }\right) ,\omega  > 0,0 < \varphi  < \pi ,x \in  \mathbf{R}$.
    \begin{enumerate}
        \item 若函数 $y = f\left( x\right)$ 的最小正周期为 $\pi$ ,当 $f\left( \dfrac{\pi }{6}\right)  = 1$ 时,求 $\omega$ 和 $\varphi$ 的值
        \item 若 $\omega  = 1$ ， $\varphi  = \dfrac{\pi }{6}$ ，函数 $y = {f}^{2}\left( x\right)  - {2f}\left( x\right)  - a$ 有零点，求实数 $a$ 的取值范围.
    \end{enumerate}
    \begin{jieda}
        \begin{enumerate}
            \item 因为函数 $y = f\left( x\right)$ 的最小正周期 $\dfrac{2\pi }{\omega } = \pi$ ，所以 $\omega  = 2$ . 当 $x = \dfrac{\pi }{6}$ 时, $\sin \left( {\dfrac{\pi }{3} + \varphi }\right)  = 1$ ,所以 $\dfrac{\pi }{3} + \varphi  = {2k\pi } + \dfrac{\pi }{2}\left( {k \in  \mathbf{Z}}\right)$ ,得 $\varphi  = {2k\pi } + \dfrac{\pi }{6}\left( {k \in  \mathbf{Z}}\right)$ ,因为 $0 < \varphi  < \pi$ ,所以取 $k = 0$ 得 $\varphi  = \dfrac{\pi }{6}$
            
            \item 当 $\omega  = 1,\varphi  = \dfrac{\pi }{6}$ 时， $f\left( x\right)  = \sin \left( {x + \dfrac{\pi }{6}}\right) ,x \in  \mathbf{R}$ ,设 $t =  f\left( x\right)  = \sin \left( {x + \dfrac{\pi }{6}}\right)  \in  \left\lbrack  {-1,1}\right\rbrack$ . 由题意得， ${t}^{2} - {2t} - a = 0$ 在 $x \in  \left\lbrack  {-1,1}\right\rbrack$ 有解.
            \begin{dingautolist}{172}
                \item 化简得 $a = {t}^{2} - {2t}$ . 又因为 $g\left( t\right)  = {t}^{2} - {2t}$ 在 $t \in  \left\lbrack  {-1,1}\right\rbrack$ 上严格减，所以 $a \in  \left\lbrack  {-1,3}\right\rbrack$.
                \item 记 $g\left( t\right)  = {t}^{2} - {2t} - a$ ,由根的分布可得 $\left\{  \begin{array}{l} \Delta  \geq  0, \\  g\left( {-1}\right)  \geq  0 \end{array}\right.$ , 所以 $a \in  \left\lbrack  {-1,3}\right\rbrack$.
            \end{dingautolist}
        \end{enumerate}
    \end{jieda}

    \item (2025$\cdot$闵行区二模$\cdot$18)已知 $f\left( x\right)  = \sin \left( {{\omega x} + \varphi }\right) \left( {\omega  > 0,0 < \varphi  < \pi }\right)$ ，函数 $y = f\left( x\right)$ 的部分图像如图所示，图中最高点 $S\left( {\dfrac{\pi }{3},1}\right)$ ，最低点 $T\left( {\dfrac{4\pi }{3}, - 1}\right)$.

    \begin{center}
    \includegraphics[width=0.2\textwidth]{bodyxb/src/66_1066_931_332_263_0.jpg}
    \end{center}

    \begin{enumerate}
        \item 求函数 $y = f\left( x\right)$ 的解析式
        \item 若 $\bigtriangleup  {ABC}$ 的内角 $A$ 、 $B$ 、 $C$ 所对的边分别为 $a$ 、 $b$ 、 $c$ ，若 $A \neq  B$ ， $f\left( A\right)  = f\left( B\right) ,c = 2$ ,求 $\bigtriangleup {ABC}$ 面积的取值范围.
    \end{enumerate}
    \begin{jieda}
        \begin{enumerate}
            \item 因为 $f\left( x\right)  = \sin \left( {{\omega x} + \varphi }\right)$ 图像经过 $S\left( {\dfrac{\pi }{3},1}\right)$ , $T\left( {\dfrac{4\pi }{3}, - 1}\right)$ ,所以 $\dfrac{4\pi }{3} - \dfrac{\pi }{3} = \dfrac{T}{2}$ 得周期 $T = {2\pi }$ ,由 $\dfrac{2\pi }{\omega } = {2\pi }$ 得, $\omega  = 1$ . 又 $f\left( \dfrac{\pi }{3}\right)  = 1$ 得 $\dfrac{\pi }{3} + \varphi  = {2k\pi } + \dfrac{\pi }{2},k \in  \mathbf{Z}$ ,又因为 $0 <  \varphi  < \pi$ ,所以 $\varphi  = \dfrac{\pi }{6}$ ,所以 $f\left( x\right)  = \sin \left( {x + \dfrac{\pi }{6}}\right)$
            
            \item 因为 $f\left( A\right)  = f\left( B\right)$ ，又 $0 < A + B < \pi$ ，结合图像对称性可知: $A + B = \dfrac{2\pi }{3}$ ,则 $C = \dfrac{\pi }{3}$ ,又 $c = 2$ ,由正弦定理得: $\dfrac{a}{\sin A} =  \dfrac{b}{\sin B} = \dfrac{c}{\sin C} = \dfrac{4\sqrt{3}}{3}$ ,则 $a = \dfrac{4\sqrt{3}}{3}\sin A,b = \dfrac{4\sqrt{3}}{3}\sin B$ ,所以 ${ab} = \dfrac{16}{3}\sin A\sin B = \dfrac{16}{3}\sin A\sin \left( {\dfrac{2\pi }{3} - A}\right)  = \dfrac{16}{3}\sin A\left( {\dfrac{\sqrt{3}}{2}\cos A + }\right.  \left. {\dfrac{1}{2}\sin A}\right)  = \dfrac{8\sqrt{3}}{3}\sin A\cos A + \dfrac{8}{3}{\sin }^{2}A = \dfrac{4\sqrt{3}}{3}\sin {2A} + \dfrac{4}{3} -  \dfrac{4}{3}\cos {2A} = \dfrac{8}{3}\sin \left( {{2A} - \dfrac{\pi }{6}}\right)  + \dfrac{4}{3}$ ,由 $A + B = \dfrac{2\pi }{3},A \neq  B$ ,可得 $A \in  \left( {0,\dfrac{\pi }{3}}\right)  \cup  \left( {\dfrac{\pi }{3},\dfrac{2\pi }{3}}\right)$ ,所以 ${2A} - \dfrac{\pi }{6} \in  \left( {-\dfrac{\pi }{6},\dfrac{\pi }{2}}\right)  \cup  \left( \dfrac{\pi }{2}, \dfrac{7\pi }{6}\right)$ ,则 $\sin \left( {{2A} - \dfrac{\pi }{6}}\right)  \in  \left( {-\dfrac{1}{2},1}\right)$ ,故 ${ab} \in  \left( {0,4}\right)$ ,于是可得 $\bigtriangleup {ABC}$ 的面积为 $\dfrac{1}{2}{ab}\sin C = \dfrac{1}{2}{ab} \cdot  \dfrac{\sqrt{3}}{2} = \dfrac{\sqrt{3}}{4}{ab} \in  \left( {0,\sqrt{3}}\right)$ , 故 $\bigtriangleup {ABC}$ 面积的取值范围为 $\left( {0,\sqrt{3}}\right)$.
        \end{enumerate}
    \end{jieda}

    \item (2025$\cdot$普陀区一模$\cdot$18)设函数 $y = f\left( x\right)$ 的表达式为 $f\left( x\right)  = \sin \left( {\omega x}\right)$ ，其中 $\omega  > 0$.
    \begin{enumerate}
        \item 设 $\omega  = 1,m \in  \mathbf{R}$ ，若有且只有一个 ${x}_{0} \in  \left( {0,m}\right)$ ，使得函数 $y = f\left( {x + \dfrac{\pi }{4}}\right)$ 取得最小值，求 $m$ 的取值范围
        \item 若对任意的 $x \in  \mathbf{R}$ ，皆有 $f\left( x\right)  + f\left( {\dfrac{2\pi }{3} - x}\right)  = 0$ 成立，且函数 $y = f\left( x\right)$ 在区间 $\left( {-\dfrac{\pi }{8},0}\right)$ 上是严格增函数,求函数 $y = f\left( x\right)$ 的最小正周期.
    \end{enumerate}
    \begin{jieda}
        \begin{enumerate}
            \item 当 $\omega  = 1$ 时， $f\left( x\right)  = \sin x$ ，所以函数 $y = f\left( {x + \dfrac{\pi }{4}}\right)  =  \sin \left( {x + \dfrac{\pi }{4}}\right)$ 在区间 $\left( {0,m}\right)$ 上只有一个最小值点，又因为 $\dfrac{\pi }{4} <  x + \dfrac{\pi }{4} < \dfrac{\pi }{4} + m$ ,由正弦函数的图像可知: $\dfrac{3\pi }{2} < m + \dfrac{\pi }{4} \leq  \dfrac{7\pi }{2}$ , 解得 $\dfrac{5\pi }{4} < m \leq  \dfrac{13\pi }{4}$ ,所以 $m$ 的取值范围为 $\left( {\dfrac{5\pi }{4},\dfrac{13\pi }{4}}\right\rbrack$ 
            
            \item 由 $f\left( x\right)  + f\left( {\dfrac{2\pi }{3} - x}\right)  = 0$ ，可知函数 $y = f\left( x\right)$ 关于点 $\left( {\dfrac{\pi }{3},0}\right)$ 对称. 因此 $\sin \left( \dfrac{\omega \pi }{3}\right)  = 0$ ,解得 $\omega  = {3k}$ ,其中 $k$ 为整数. 由于函数在区间 $\left( {-\dfrac{\pi }{8},0}\right)$ 上是严格增函数,所以 $\dfrac{T}{2} \geq  0 - \left( {-\dfrac{\pi }{8}}\right)  = \dfrac{\pi }{8}$ ,所以 $T = \dfrac{2\pi }{\omega } \geq  \dfrac{\pi }{4}$ ,结合 $\omega  = {3k}$ ,其中 $k$ 为整数,所以 $0 < \omega  \leq  8$ ,又 $\omega  = {3k}$ ,其中 $k$ 为整数,所以 $\omega  = 3$ 或 $\omega  = 6$ ,当 $\omega  = 6$ 时, $f\left( x\right)  = \sin \left( {6x}\right)$ ,函数在区间 $\left( {-\dfrac{\pi }{8},0}\right)$ 上不是严格增函数; 当 $\omega  = 3$ 时, $f\left( x\right)  = \sin \left( {3x}\right)$ ,函数在区间 $\left( {-\dfrac{\pi }{8},0}\right)$ 上是严格增函数,且关于点 $\left( {\dfrac{\pi }{3},0}\right)$ 对称. 所以 $\omega  = 3$ . 因此函数 $y = f\left( x\right)$ 的最小正周期为 $T = \dfrac{2\pi }{3}$ .
        \end{enumerate}
    \end{jieda}

    \item (2025$\cdot$普陀区二模$\cdot$18)设 $\omega  > 0$ ，函数 $y = f\left( x\right)$ 的表达式为 $f\left( x\right)  = \cos \left( {{\omega x} - \dfrac{\pi }{6}}\right)$.
    \begin{enumerate}
        \item 若 $\omega  = 2$ ，设 $\bigtriangleup  {ABC}$ 的内角 $A$ 、 $B$ 、 $C$ 的对边分别为 $a$ 、 $b$ 、 $c$ ， $,f\left( C\right)  =  - \dfrac{\sqrt{3}}{2}$ ，且 ${c}^{2} - {a}^{2} - {b}^{2} = 4$ ，求 $\bigtriangleup  {ABC}$ 的面积
        \item 对任意的 ${x}_{0} \in  \mathbf{R}$ ，皆有 $f\left( {x}_{0}\right)  \leq  f\left( \dfrac{\pi }{3}\right)$ 成立，且该函数在区间 $\left( {\dfrac{\pi }{6},\dfrac{\pi }{2}}\right)$ 上不存在最小值，求函数 $y = f\left( x\right)$ 在 $x \in  \left( {\dfrac{4\pi }{3},\dfrac{17\pi }{6}}\right)$ 的单调区间.
    \end{enumerate}
    \begin{jieda}
        \begin{enumerate}
            \item 因为 $\omega  = 2$ ，所以 $f\left( x\right)  = \cos \left( {{2x} - \dfrac{\pi }{6}}\right)$ ，所以 $f\left( C\right)  =  \cos \left( {{2C} - \dfrac{\pi }{6}}\right)  =  - \dfrac{\sqrt{3}}{2}$ ,解得 ${2C} - \dfrac{\pi }{6} = \dfrac{7\pi }{6}$ 或 $\dfrac{5\pi }{6}$ ,所以 $C = \dfrac{\pi }{2}$ 或 $\dfrac{2\pi }{3}$ ,若 $C = \dfrac{\pi }{2}$ ,则 ${c}^{2} = {a}^{2} + {b}^{2}$ ,不符; 所以 $C = \dfrac{2\pi }{3}$ ,所以 $\cos C = \dfrac{{a}^{2} + {b}^{2} - {c}^{2}}{2ab} =  - \dfrac{1}{2}$ ,所以 ${c}^{2} = {a}^{2} + {b}^{2} + {ab}$ ,由 ${c}^{2} -  {a}^{2} - {b}^{2} = 4$ ,得 $4 + {a}^{2} + {b}^{2} = {a}^{2} + {b}^{2} + {ab}$ ,所以 ${ab} = 4,{S}_{\bigtriangleup {ABC}} =  \dfrac{1}{2}{ab}\sin C = \dfrac{1}{2} \times  4 \times  \dfrac{\sqrt{3}}{2} = \sqrt{3}$
            
            \item 由 $f\left( \dfrac{\pi }{3}\right)  = 1$ ，得 $\dfrac{\omega \pi }{3} - \dfrac{\pi }{6} = {2k\pi }$ ，所以 $\omega  = {6k} + \dfrac{1}{2}$ ， $k \in  \mathbf{Z}$ ，令 $t = {\omega x} - \dfrac{\pi }{6}$ ,因为 $x \in  \left( {\dfrac{\pi }{6},\dfrac{\pi }{2}}\right)$ ,所以 $t \in  \left( {\dfrac{\pi }{6}\omega  - \dfrac{\pi }{6},\dfrac{\pi }{2}\omega  - }\right.  \left. \dfrac{\pi }{6}\right)$ ,又函数在 $\left( {\dfrac{\pi }{6},\dfrac{\pi }{2}}\right)$ 无最小值,所以 $\left\{  \begin{array}{l} \dfrac{\pi }{6}\omega  - \dfrac{\pi }{6} >  - \dfrac{\pi }{2} \\  \dfrac{\pi }{2}\omega  - \dfrac{\pi }{6} < \dfrac{3\pi }{2} \end{array}\right.$ ,解得 $\omega  \in  \left( {0,\dfrac{10}{3}}\right)$ ,所以 $k = 0$ ,则 $\omega  = \dfrac{1}{2}$ ,所以 $f\left( x\right)  = \cos \left( {\dfrac{1}{2}x}\right.  \left. {-\dfrac{\pi }{6}}\right)$ ,令 ${2k\pi } \leq  \dfrac{1}{2}x - \dfrac{\pi }{6} \leq  {2k\pi } + \pi$ ,所以 ${4k\pi } + \dfrac{\pi }{3} \leq  x \leq  {4k\pi }  + \dfrac{7\pi }{3},k \in  \mathbf{Z}$ ,当 $k = 0$ 时, $\dfrac{\pi }{3} \leq  x \leq  \dfrac{7\pi }{3}$ ,因为 $x \in  \left( {\dfrac{4\pi }{3},\dfrac{17\pi }{6}}\right)$ ,所以 $y = f\left( x\right)$ 在 $\left( {\dfrac{4\pi }{3},\dfrac{7\pi }{3}}\right)$ 单调递减,在 $\left( {\dfrac{7\pi }{3},\dfrac{17\pi }{6}}\right)$ 单调递增.
        \end{enumerate}
    \end{jieda}

    \item (2025$\cdot$黄浦区一模$\cdot$18)已知 $f\left( x\right)  = \sin x$.
    \begin{enumerate}
        \item 求函数 $y = f\left( x\right)  \cdot  f\left( {\dfrac{\pi }{2} - x}\right)$ 的最小正周期
        \item 求函数 $y = f\left( {{2x} + \dfrac{\pi }{3}}\right) ,x \in  \left\lbrack  {0,\dfrac{\pi }{2}}\right\rbrack$ 的单调减区间.
    \end{enumerate}
    \begin{jieda}
        \begin{enumerate}
            \item 由 $f\left( x\right)  = \sin x$ ,得 $f\left( {\dfrac{\pi }{2} - x}\right)  = \sin \left( {\dfrac{\pi }{2} - x}\right)  =  \cos x$ ,则函数 $y = f\left( x\right)  \cdot  f\left( {\dfrac{\pi }{2} - x}\right)  = \sin x \cdot  \cos x = \dfrac{1}{2}\sin {2x}$ , 故最小正周期为 $\dfrac{2\pi }{2} = \pi$
            
            \item 由 $f\left( x\right)  = \sin x$ ，得 $y = f\left( {{2x} + \dfrac{\pi }{3}}\right)  = \sin \left( {{2x} + \dfrac{\pi }{3}}\right) ,x \in  \left\lbrack  {0,\dfrac{\pi }{2}}\right\rbrack$ ; 由 $0 \leq  x \leq  \dfrac{\pi }{2}$ ,得 $\dfrac{\pi }{3} \leq  {2x} + \dfrac{\pi }{3} \leq  \dfrac{4\pi }{3}$ ,令 $\dfrac{\pi }{2} \leq  {2x} +  \dfrac{\pi }{3} \leq  \dfrac{4\pi }{3}$ ,解得 $\dfrac{\pi }{12} \leq  x \leq  \dfrac{\pi }{2}$ ,故单调减区间为 $\left\lbrack  {\dfrac{\pi }{12},\dfrac{\pi }{2}}\right\rbrack$ .
        \end{enumerate}
    \end{jieda}
\end{enumerate}
\subsection*{考点8 三角函数图像的伸缩平移变换}

\subsubsection*{一、填空题}

\begin{enumerate}
    \item (2024$\cdot$闵行区一模$\cdot$7)若将函数 $y = \sin \left( {{2x} + \varphi }\right) \left( {0 < \varphi  < \pi }\right)$ 的图像向右平移 $\dfrac{\pi }{3}$ 个单位,得到的图像所对应的函数为奇函数,则 $\varphi  =$ \blkda[2]{$\dfrac{2\pi }{3}$}.
    \begin{jieda}
        原函数有对称中心$\left(-\dfrac{\pi}{3}, 0\right)$，故$2 \times \left(-\dfrac{\pi}{3}\right) + \varphi = k\pi, k \in \mathbf{Z}$，故$\varphi = \dfrac{2\pi}{3} + k\pi, k \in \mathbf{Z}$，又因为$0 < \varphi  < \pi $，所以$\varphi = \dfrac{2\pi}{3}$
    \end{jieda}

    \item (2024$\cdot$虹口区一模$\cdot$8)已知函数 $f\left( x\right)  = \cos \left( {{\omega x} + \varphi }\right) (\omega  > 0$ ， $\left| \varphi \right|  < \dfrac{\pi }{2})$ 的部分图像如图所示，则 $f\left( x\right)  =$ \blkda[2]{$\cos \left( {{2x} - \dfrac{\pi }{6}}\right)$}.
    \begin{center}
    \includegraphics[width=0.2\textwidth]{bodyxb/src/68_1078_908_343_267_0.jpg}
    \end{center}

    \item (2024 $\cdot$奉贤区一模$\cdot$9)设函数 $y = \sin {\omega x}\left( {\omega  > 0}\right)$ 在区间 $\left( {0,{2\pi }}\right)$ 上恰有三个极值点，则 $\omega$ 的取值范围为\blkda[2]{$\left( {\dfrac{5}{4},\dfrac{7}{4}}\right\rbrack$}.
    \begin{jieda}
        $x \in  \left( {0,{2\pi }}\right) ,{\omega x}  \in  \left( {0,{2\pi \omega }}\right)  \Longrightarrow  \dfrac{5\pi }{2} < {2\pi \omega } \leq  \dfrac{7\pi }{2} \Longrightarrow  \omega  \in  \left( {\dfrac{5}{4},\dfrac{7}{4}}\right\rbrack$.

        \begin{center}
        \includegraphics[width=0.4\textwidth]{bodyxb/src/46_1219_1003_360_168_0.jpg}
        \end{center}
    \end{jieda}

    \item (2024$\cdot$浦东新区一模$\cdot$10)如图所示，已知函数 $y = A\sin \left( {{\omega x} + \varphi }\right)  \left( {A > 0,\omega  > 0,0 < \varphi  < \dfrac{\pi }{2}}\right)$ 的图像与 $y$ 轴的交点为 $\left( {0,1}\right)$ ,并已知其在 $y$ 轴右侧的第一个最高点和第一个最低点的坐标分别为 $\left( {{x}_{0},2}\right)$ 和 $\left( {{x}_{0} + {2\pi }, - 2}\right)$. 记 $y = f\left( x\right)$ ,则 $f\left( \dfrac{\pi }{3}\right)  =$ \blkda[2]{$\sqrt{3}$}.

    \begin{center}
    \includegraphics[width=0.3\textwidth]{bodyxb/src/68_1024_1329_399_220_0.jpg}
    \end{center}
    \begin{jieda}
        $\dfrac{T}{2} = {2\pi },T = {4\pi } = \dfrac{2\pi }{\omega } \Longrightarrow  \omega  = \dfrac{1}{2},A = 2 \Longrightarrow  f\left( x\right)  =  2\sin \left( {\dfrac{1}{2}x + \varphi }\right) ,f\left( 0\right)  = 1 \Longrightarrow  \varphi  = \dfrac{\pi }{6},f\left( x\right)  = 2\sin \left( {\dfrac{1}{2}x + \dfrac{\pi }{6}}\right)  \Longrightarrow  f\left( \dfrac{\pi }{3}\right)  = \sqrt{3}.$
    \end{jieda}

    \item (2024$\cdot$普陀区一模$\cdot$10)设函数 $y = \sin \left( {{2x} + \varphi }\right) \left( {0 < \varphi  < \dfrac{\pi }{2}}\right)$ 的图像与直线 $y = t$ 相交的连续的三个公共点从左到右依次记为 $A\text{ 、 }B\text{ 、 }C$ ，若 $\left| {BC}\right|  = 2\left| {AB}\right|$ ，则正实数 $t$ 的值为\blkda[2]{$\dfrac{1}{2}$}.
    \begin{jieda}
        作出函数 $y = \sin \left( {{2x} + \varphi }\right)$ ,的大致图象,如图,令 ${2x} + \varphi  = {2k\pi } + \dfrac{\pi }{2},k \in  {Z}$ ,

        解得 $x = {k\pi } + \dfrac{\pi }{4} - \dfrac{\varphi }{2},k \in  {Z}$ ,

        \begin{center}
        \includegraphics[width=0.2\textwidth]{bodyxb/src/1_87_1144_381_303_0.jpg}
        \end{center}


        则函数 $y = \sin \left( {{2x} + \varphi }\right)$ 的图象与直线 $y = t\left( {0 < t < 1}\right)$ 连续的三个公共点 $A,B,C$ ,(可以同时往左或往右移动正整数倍周期长度)

        即 $A,B$ 关于直线 $x = {k\pi } + \dfrac{\pi }{4} - \dfrac{\varphi }{2},k \in  {Z}$ 对称, $\left| {AC}\right|  = T = \dfrac{2\pi }{2} = \pi$ ,

        由于 $\left| {BC}\right|  = 2\left| {AB}\right|$ ,故 $\left| {AB}\right|  = \dfrac{1}{3}\left| {AC}\right|  = \dfrac{\pi }{3}$ ,

        而 $A,B$ 关于直线 $x = {k\pi } + \dfrac{\pi }{4} - \dfrac{\varphi }{2},k \in  {Z}$ 对称,

        故 $A$ 点横坐标为 ${k\pi } + \dfrac{\pi }{4} - \dfrac{\varphi }{2} - \dfrac{\pi }{6} = {k\pi } + \dfrac{\pi }{12} - \dfrac{\varphi }{2}$ ,

        将 $A$ 点横坐标代入 $y = \sin \left( {{2x} + \varphi }\right)$ ，得 $\sin \left( {\dfrac{\pi }{6} + {2k\pi }}\right)  = t = \dfrac{1}{2}$
    \end{jieda}

    \item (2025$\cdot$黄浦区二模$\cdot$12)设 $a$ 、 $b$ 为常数， $f\left( x\right)  = \left| {a + \sin x}\right|  + \left| {a - \sin x}\right|$ ，若对任意的 $b \in  \left( {1,2}\right)$ ，函数 $y = f\left( x\right)  - b$ 在区间 $\left\lbrack  {0,{2\pi }}\right\rbrack$ 上恰有 4 个零点，则 $a$ 的取值范围是\blkda[2]{$\left[-\dfrac{1}{2}, \dfrac{1}{2}\right]$}.
    \begin{jieda}
        由 $x \in  \left\lbrack  {0,{2\pi }}\right\rbrack$ ,则 $\sin x \in  \left\lbrack  {-1,1}\right\rbrack$ ,又 $b \in  \left( {1,2}\right)$
        
        当 $a \geq  1,f\left( x\right)  = a + \sin x + a - \sin x = {2a} \geq  2$ ,此时 $y = {2a} - b$ 无零点
        
        当 $a \leq   - 1,f\left( x\right)  =  - a - \sin x - a + \sin x =  - {2a} \geq  2$ ,此时 $y =  - {2a} - b$ 无零点
        
        当 $0 \leq  \left| a\right|  < 1,f\left( x\right)  = \left\{  \begin{array}{l} 2\left| a\right| ,0 \leq  \left| {\sin x}\right|  \leq  \left| a\right| \\  2\left| {\sin x}\right| ,\left| a\right|  < \left| {\sin x}\right|  \leq  1 \end{array}\right.$ 如图,此时 $f\left( x\right)  \in  \left\lbrack  {2\left| a\right| ,2}\right\rbrack$ ,而 $b \in  \left( {1,2}\right)$ ,要使 $f\left( x\right)  = b$ 在区间 $\left\lbrack  {0,{2\pi }}\right\rbrack$ 上恰有 4 个根,则 $2\left| a\right|  \leq  1$ ,则 $- \dfrac{1}{2} \leq  a \leq  \dfrac{1}{2}$.
        \begin{center}
            \includegraphics[width=0.4\textwidth]{bodyxb/src/47_552_466_266_169_0.jpg}
        \end{center}
    \end{jieda}

    \item (2025$\cdot$徐汇区二模$\cdot$12)设实数 $\omega  > 0$ ，若 $f\left( x\right)  = \sin {\omega x}$ 满足对任意 ${x}_{1} \in  \left\lbrack  {0,\pi }\right\rbrack$ ，都存在 ${x}_{2} \in  \left\lbrack  {\pi ,{2\pi }}\right\rbrack$ ，使得 $f\left( {x}_{1}\right)  + f\left( {x}_{2}\right)  = 0$ 成立，则 $\omega$ 的最小值是\blkda[2]{$\dfrac{3}{4}$}.
    \begin{jieda}
        换元$t = \omega x$，故$\forall t_1 \in [0, \omega\pi]$，$\exists t_2 \in [\omega\pi, 2\omega\pi]$使得$sint_1 + sint_2 = 0$

        假设$\omega < \dfrac{\pi}{2}$，则$sint_1, sint_2 \geq 0$，不满足要求.

        故$sint_1$可以取到$1$，故$sint_2$应可以取到$-1$，故$\omega \geq \dfrac{3}{4}$，经检验$\omega = \dfrac{3}{4}$满足要求.
    \end{jieda}

    \item (2025$\cdot$闵行区一模$\cdot$12)已知 $f\left( x\right)  = \left| {\sin {\omega x}}\right|$ ，若存在 ${x}_{1}$ 、 ${x}_{2} \in  \left\lbrack  {{\omega \pi },{2\omega \pi }}\right\rbrack$ ，且 ${x}_{1} \neq  {x}_{2}$ ，使得 $\dfrac{1}{f\left( {x}_{1}\right)  + 1} + \dfrac{1}{f\left( {x}_{2}\right)  + 1} = 1$ 成立，则 $\omega$ 的取值范围是\blkda[4]{$\left\lbrack  {\dfrac{\sqrt{5}}{2},\dfrac{\sqrt{6}}{2}}\right\rbrack   \cup  \left\lbrack  {\dfrac{\sqrt{7}}{2}, + \infty }\right)$}.
    \begin{jieda}
        原问题转化为$y = sinx$在$[\omega^2\pi, 2\omega^2\pi]$上存在至少两个对称轴.
    \end{jieda}
\end{enumerate}

\subsubsection*{二、选择题}
\begin{enumerate}
    \item (2024$\cdot$虹口区二模$\cdot$14)设 $f\left( x\right)  = \sin {2x} + \sqrt{3}\cos {2x}$ ，将函数 $y = f\left( x\right)$ 的图像沿 $x$ 轴向右平移 $\dfrac{\pi }{6}$ 个单位,得到函数 $y = g\left( x\right)$ 的图像,则 \dotfill{(\kh{D})}
    \xx{函数 $y = g\left( x\right)$ 是偶函数}{函数 $y = g\left( x\right)$ 的图像关于直线 $x = \dfrac{\pi }{2}$ 对称}{函数 $y = g\left( x\right)$ 在 $\left\lbrack  {\dfrac{\pi }{4},\dfrac{\pi }{2}}\right\rbrack$ 上是严格增函数}{函数 $y = g\left( x\right)$ 在 $\left\lbrack  {\dfrac{\pi }{6},\dfrac{2\pi }{3}}\right\rbrack$ 上的值域为 $\left\lbrack  {-\sqrt{3},2}\right\rbrack$}
    \item (2025$\cdot$嘉定区二模$\cdot$15)已知关于 $x$ 的不等式 $\left| {\sin {2x}}\right|  > \cos x$ 在区间 $\left\lbrack  {0,{2\pi }}\right\rbrack$ 内有 $k$ 个整数解，则 $k$ 的值为 \dotfill{(\kh{C})}
    \xx{3}{4}{5}{6}
    \begin{jieda}
        由题设 $2\left| {\sin x\cos x}\right|  > \cos x$ ,显然 $x \notin  \left\{  {0,\dfrac{\pi }{2},\dfrac{3\pi }{2},{2\pi }}\right\}$ , 当 $x \in  \left( {0,\dfrac{\pi }{2}}\right)  \cup  \left( {\dfrac{3\pi }{2},{2\pi }}\right)$ ,则 $\left| {\sin x}\right|  > \dfrac{1}{2}$ ,此时 $x \in  \left( \dfrac{\pi }{6}\right.$ , $\left. \dfrac{\pi }{2}\right)  \cup  \left( {\dfrac{3\pi }{2},\dfrac{11\pi }{6}}\right)$ ,当 $x \in  \left( {\dfrac{\pi }{2},\dfrac{3\pi }{2}}\right)$ ,则 $2\left| {\sin x\cos x}\right|  \geq  0 >  \cos x$ ,此时 $x \in  \left( {\dfrac{\pi }{2},\dfrac{3\pi }{2}}\right)$ ,所以,整数解有1,2,3,4,5,共 5 个整数解.

        \ps 绘制振动图像即可.
    \end{jieda}
    \item (2024$\cdot$崇明区二模$\cdot$15)设函数 $f\left( x\right)  = \sin \left( {x - \dfrac{\pi }{6}}\right)$ ，若对于任意 $\alpha  \in  \left\lbrack  {-\dfrac{5\pi }{6}, - \dfrac{\pi }{2}}\right\rbrack$ ，在区间 $\left\lbrack  {0,m}\right\rbrack$ 上总存在唯一确定的 $\beta$ ,使得 $f\left( \alpha \right)  + f\left( \beta \right)  = 0$ ,则 $m$ 的最小值为 \dotfill{(\kh{B})}
    \xx{$\dfrac{\pi }{6}$}{$\dfrac{\pi }{2}$}{$\dfrac{7\pi }{6}$}{$\pi$}
    \begin{jieda}
        $f(\beta)$的值域包含$\left[0, \dfrac{\sqrt{3}}{2}\right]$，并且$\left[0, \dfrac{\sqrt{3}}{2}\right]$内的函数值和自变量一一对应.
    \end{jieda}
    \item (2025$\cdot$闵行区二模$\cdot$15)已知函数 $y = \cos \left( {\dfrac{\pi }{6}x + \dfrac{\pi }{3}}\right)$ 在区间 $\left\lbrack  {a,a + 9}\right\rbrack$ 上既有最大值又有最小值， 则关于实数 $a$ 的取值，以下不可能的是 \dotfill{(\kh{D})}
    \xx{2024}{2025}{2026}{2027}
    \begin{jieda}
        由题意可得函数的周期为 $T = \dfrac{2\pi }{\dfrac{\pi }{6}} = {12}$ ,最大值点满足 $\dfrac{\pi }{6}x + \dfrac{\pi }{3} = {2k\pi },k \in  \mathbf{Z}$ ,解得 $x = {12k} - 2,k \in  \mathbf{Z}$ ,最小值点满足 $\dfrac{\pi }{6}x + \dfrac{\pi }{3} = {2k\pi } + \pi ,k \in  \mathbf{Z}$ ,解得 $x = {12k} + 4,k \in  \mathbf{Z}$ ,因为函数 $y = \cos \left( {\dfrac{\pi }{6}x + \dfrac{\pi }{3}}\right)$ 在区间 $\left\lbrack  {a,a + 9}\right\rbrack$ 上既有最大值又有最小值,区间的长度为 9,对于 ${A}$ ,若 $a = {2024}$ ,当 $k = {169}$ 时，最大值点为 2026，最小值点为 2032，此时位于区间 [2024,2033]内，故 A 正确；对于 B，若 $a = {2025}$ ，当 $k = {169}$ 时，最大值点为 2026，最小值点为 2032，此时位于区间 [2025,2034]内，故 B 正确；对于 C，若 $a = {2026}$ ，当 $k = {169}$ 时，最大值点为 2026，最小值点为 2032，此时位于区间 [2026,2035]内，故 C 正确；对于 D，若 $a = {2027}$ ，当 $k = {169}$ 时，最大值点为 2026，当 $k = {170}$ 时，最大值点为 2038，此时不位于区间 $\left\lbrack  {{2027},{2036}}\right\rbrack$ 内，故 D 错误. 故选:D.
    \end{jieda}
    \item (2025$\cdot$奉贤区二模$\cdot$16)若 ${5\pi }$ 是函数 $y = \cos {nx}\sin \dfrac{2000}{{n}^{2}}x$ 的一个周期，则正整数 $n$ 所有可能取值个数是 \dotfill{(\kh{B})}
    \xx{2}{3}{4}{9}
    \begin{jieda}
        由题意可得 $\cos n\left( {x + {5\pi }}\right)  \cdot  \sin \dfrac{2000}{{n}^{2}}\left( {x + {5\pi }}\right)  =  \cos \left( {{nx} + {5n\pi }}\right)  \cdot  \sin \left( {\dfrac{2000}{{n}^{2}}x + \dfrac{10000}{{n}^{2}}\pi }\right)  = \cos {nx}\sin \dfrac{2000}{{n}^{2}}x$
        \begin{dingautolist}{172}
            \item ${5n\pi } = 2{k}_{1}\pi ,{k}_{1} \in  \mathbf{Z},\dfrac{10000}{{n}^{2}}\pi  = 2{k}_{2}\pi ,{k}_{2} \in  \mathbf{Z}$，解得 $n = \dfrac{2}{5}{k}_{1},{k}_{1} \in  \mathbf{Z},n =  {50}\sqrt{\dfrac{2}{{k}_{2}}},{k}_{2} \in  \mathbf{Z}$，由 $n$ 为正整数，且 50 的因数为1,2,5,10,25,50，则 ${k}_{2}$ 的取值可能有2,8,50,200,1250,5000,此时 $n$ 的可能取值有 50,25,10,5,2,1,由 ${k}_{1} \in  \mathbf{Z}$ ,则 $n$ 为 2 的倍数,故 $n$ 的可能取值有 50,10,2.
            \item ${5n\pi } = 2{k}_{1}\pi  + \pi$ , ${k}_{1} \in  \mathbf{Z},\dfrac{10000}{{n}^{2}}\pi  = 2{k}_{2}\pi  + \pi ,{k}_{2} \in  \mathbf{Z}$，解得$n = \dfrac{2{k}_{1} + 1}{5},{k}_{1} \in  \mathbf{Z},n = {100}\sqrt{\dfrac{1}{2{k}_{2} + 1}}$ , ${k}_{2} \in  \mathbf{Z}$，由 $n$ 为正整数，且 100 的因数为1,2,4,5,10,20,25,50, 100,则奇数 $\sqrt{2{k}_{2} + 1}$ 的取值只可能有 5,25,此时 $n$ 的可能取值有 20,4,由 ${k}_{1} \in  \mathbf{Z}$ ,则 $n$ 奇数,所以此时 $n$ 无取值.
        \end{dingautolist}
    \end{jieda}
\end{enumerate}
\subsubsection*{三、解答题}
\begin{enumerate}
    \item (2024$\cdot$长宁区二模$\cdot$17)某同学用“五点法”画函数 $f\left( x\right)  = \sin \left( {{\omega x} + \varphi }\right) \left( {\omega  > 0}\right)$ 在某一个周期内的图像时，列表并填入了部分数据，如下表:
    \begin{center}
    \begin{tabular}{|c|c|c|c|c|c|}
    \hline
    ${\omega x} + \varphi$ & 0 & $\dfrac{\pi }{2}$ & $\pi$ & $\dfrac{3\pi }{2}$ & ${2\pi }$ \\
    \cline{1-6}
    $x$ & $\Delta$ & $\dfrac{\pi }{6}$ & $\dfrac{5\pi }{12}$ & $\dfrac{2\pi }{3}$ & $\dfrac{11\pi }{12}$ \\
    \cline{1-6}
    $\sin \left( {{\omega x} + \varphi }\right)$ & 0 & 1 & $\Delta$ & -1 & 0 \\
    \cline{1-6}
    \hline
    \end{tabular}
    \end{center}
    \begin{enumerate}
        \item 请将上表 $\Delta$ 处的数据补充完整，并直接写出函数 $y = f\left( x\right)$ 的解析式
        \item 设 $\omega  = 1,\varphi  = 0,g\left( x\right)  = {f}^{2}\left( x\right)  + f\left( x\right) f\left( {\dfrac{\pi }{2} - x}\right) \left( {x \in  \left\lbrack  {0,\dfrac{\pi }{2}}\right\rbrack  }\right)$ ，求函数 $y = g\left( x\right)$ 的值域.
    \end{enumerate}
    \begin{jieda}
        \begin{enumerate}
            \item 由题意 $\left\{  \begin{array}{l} \omega  \cdot  \dfrac{\pi }{6} + \varphi  = \dfrac{\pi }{2} \\  \omega  \cdot  \dfrac{2\pi }{3} + \varphi  = \dfrac{3\pi }{2} \end{array}\right.$ ，解得 $\omega  = 2$ ， $\varphi  = \dfrac{\pi }{6}$ ，所以函数 $y = f\left( x\right)$ 的解析式为 $f\left( x\right)  = \sin \left( {{2x} + \dfrac{\pi }{6}}\right) ,{2x} + \dfrac{\pi }{6} = 0$ 时,解得 $x =  - \dfrac{\pi }{12}$ ,当 $x = \dfrac{5\pi }{12}$ 时, ${2x} + \dfrac{\pi }{6} = \pi ,\sin \left( {{2x} + \dfrac{\pi }{6}}\right)  = 0$ ,将表中 $\Delta$ 处的数据补充完整如下表:

            \begin{center}
            \begin{tabular}{|c|c|c|c|c|c|}
            \hline
            ${\omega x} + \varphi$ & 0 & $\dfrac{\pi }{2}$ & $\pi$ & $\dfrac{3\pi }{2}$ & ${2\pi }$ \\
            \cline{1-6}
            $x$ & $- \dfrac{\pi }{12}$ & $\dfrac{\pi }{6}$ & $\dfrac{5\pi }{12}$ & $\dfrac{2\pi }{3}$ & $\dfrac{11\pi }{12}$ \\
            \cline{1-6}
            $\sin \left( {{\omega x} + \varphi }\right)$ & 0 & 1 & 0 & -1 & 0 \\
            \cline{1-6}
            \hline
            \end{tabular}
            \end{center}

            \item 若 $\omega  = 1,\varphi  = 0$ ，则 $g\left( x\right)  = {\sin }^{2}x + \sin x\sin \left( {\dfrac{\pi }{2} - x}\right)  =  {\sin }^{2}x + \sin x\cos x = \dfrac{1 - \cos {2x}}{2} + \dfrac{1}{2}\sin {2x} = \dfrac{\sqrt{2}}{2}\sin \left( {{2x} - \dfrac{\pi }{4}}\right)  +  \dfrac{1}{2}\left( {x \in  \left\lbrack  {0,\dfrac{\pi }{2}}\right\rbrack  }\right)$ ,因为 $x \in  \left\lbrack  {0,\dfrac{\pi }{2}}\right\rbrack$ ,所以 ${2x} - \dfrac{\pi }{4} \in  \left\lbrack  {-\dfrac{\pi }{4}}\right.$ , $\left. \dfrac{3\pi }{4}\right\rbrack$ ,进而 $\sin \left( {{2x} - \dfrac{\pi }{4}}\right)  \in  \left\lbrack  {-\dfrac{\sqrt{2}}{2},1}\right\rbrack$ ,所以函数 $y = g\left( x\right)$ 的值域为 $\left\lbrack  {0,\dfrac{\sqrt{2} + 1}{2}}\right\rbrack$.
        \end{enumerate}
    \end{jieda}
    \item (2024$\cdot$浦东新区二模$\cdot$17)已知函数 $y = f\left( x\right)$ ，其中 $f\left( x\right)  = \sin x$.
    \begin{enumerate}
        \item 求 $f\left( {x - \dfrac{\pi }{4}}\right)  = \dfrac{\sqrt{3}}{2}$ 在 $x \in  \left\lbrack  {0,\pi }\right\rbrack$ 上的解
        \item 已知 $g\left( x\right)  = \sqrt{3}f\left( x\right) f\left( {x + \dfrac{\pi }{2}}\right)  - f\left( x\right) f\left( {x + \pi }\right)$ ，若关于 $x$ 的方程 $g\left( x\right)  - m = \dfrac{1}{2}$ 在 $x \in  \left\lbrack  {0,\dfrac{\pi }{2}}\right\rbrack$ 时有解,求实数 $m$ 的取值范围.
    \end{enumerate}
    \begin{jieda}
        \begin{enumerate}
            \item 由题,原式等价于求 $\sin \left( {x - \dfrac{\pi }{4}}\right)  = \dfrac{\sqrt{3}}{2}$ 在 $x \in  \left\lbrack  {0,\pi }\right\rbrack$ 上的解. 从而有 $x - \dfrac{\pi }{4} = {2k\pi } + \dfrac{\pi }{3}$ 或 $x - \dfrac{\pi }{4} = {2k\pi } + \dfrac{2\pi }{3},k \in$ Z,解得 $x = {2k\pi } + \dfrac{7\pi }{12}$ 或 $x = {2k\pi } + \dfrac{11\pi }{12},k \in  \mathbf{Z}$ ,又 $x \in  \left\lbrack  {0,\pi }\right\rbrack$ , 所以 $x = \dfrac{7\pi }{12}$ 或 $x = \dfrac{11\pi }{12}$ . 因此 $f\left( {x - \dfrac{\pi }{4}}\right)  = \dfrac{\sqrt{3}}{2}$ 在 $x \in  \left\lbrack  {0,\pi }\right\rbrack$ 上的解为 $\dfrac{7\pi }{12},\dfrac{11\pi }{12}$ ;

            \begin{center}
            \includegraphics[width=0.4\textwidth]{bodyxb/src/48_1278_373_301_212_0.jpg}
            \end{center}
            

            \item 由题, $g\left( x\right)  = \sqrt{3}\sin x\sin \left( {x + \dfrac{\pi }{2}}\right)  - \sin x\sin \left( {x + \pi }\right)  =  \sqrt{3}\sin x\cos x + {\sin }^{2}x = \dfrac{\sqrt{3}}{2}\sin {2x} + \dfrac{1 - \cos {2x}}{2} = \sin \left( {{2x} - \dfrac{\pi }{6}}\right)  +  \dfrac{1}{2}$ ,故 $g\left( x\right)  - m = \dfrac{1}{2}$ 在 $x \in  \left\lbrack  {0,\dfrac{\pi }{2}}\right\rbrack$ 时有解,等价于 $m =  \sin \left( {{2x} - \dfrac{\pi }{6}}\right)$ 在 $x \in  \left\lbrack  {0,\dfrac{\pi }{2}}\right\rbrack$ 时有解. 可知 ${2x} - \dfrac{\pi }{6} \in  \left\lbrack  {-\dfrac{\pi }{6}}\right.$ , $\left. \dfrac{5\pi }{6}\right\rbrack$ ,因而 $\sin \left( {{2x} - \dfrac{\pi }{6}}\right)  \in  \left\lbrack  {-\dfrac{1}{2},1}\right\rbrack$ ,所以实数 $m$ 的取值范围是 $\left\lbrack  {-\dfrac{1}{2},1}\right\rbrack$.
        \end{enumerate}
    \end{jieda}
    \item (2024$\cdot$普陀区二模$\cdot$18)设函数 $f\left( x\right)  = \sin \left( {{\omega x} + \varphi }\right) ,\omega  > 0,0 < \varphi  < \pi$ ，它的最小正周期为 $\pi$.
    \begin{enumerate}
        \item 若函数 $y = f\left( {x - \dfrac{\pi }{12}}\right)$ 是偶函数，求 $\varphi$ 的值
        \item 在 $\bigtriangleup  {ABC}$ 中，角 $A$ 、 $B$ 、 $C$ 的对边分别为 $a$ 、 $b$ 、 $c$ ，若 $a = 2$ ， $A = \dfrac{\pi }{6}$ ， $f\left( \dfrac{B - \varphi }{2}\right)  = \dfrac{\sqrt{3}}{4}c$ ，求 $b$ 的值.
    \end{enumerate}
    \begin{jieda}
        \begin{enumerate}
            \item 因为函数 $f\left( x\right)  = \sin \left( {{\omega x} + \varphi }\right)$ 的最小正周期为 $\pi$ ，且 $\omega  > 0$ ,所以 $\dfrac{2\pi }{\omega } = \pi$ ,即 $\omega  = 2$ ,则 $y = f\left( {x - \dfrac{\pi }{12}}\right)  = \sin \left({2x} + {\varphi  - \dfrac{\pi }{6}}\right)$ ,又函数 $y = f\left( {x - \dfrac{\pi }{12}}\right)$ 是偶函数,则 $\varphi  - \dfrac{\pi }{6} = {k\pi } +  \dfrac{\pi }{2},k \in  \mathbf{Z}$ ,即 $\varphi  = {k\pi } + \dfrac{2\pi }{3}$ ,又 $0 < \varphi  < \pi$ ,则 $\varphi  = \dfrac{2\pi }{3}$ ;
            
            \item 由 $f\left( \dfrac{B - \varphi }{2}\right)  = \dfrac{\sqrt{3}}{4}c$ ，得 $\sin B = \dfrac{\sqrt{3}}{4}c$ ，又 $a = 2$ ， $A = \dfrac{\pi }{6}$ ，则 $\sin B = \dfrac{b\sin A}{a} = \dfrac{b}{4} = \dfrac{\sqrt{3}}{4}c$ ,即 $b = \sqrt{3}c$ ,由余弦定理,得 ${a}^{2} = {b}^{2} +  {c}^{2} - {2bc}\cos A = 3{c}^{2} + {c}^{2} - 2\sqrt{3}c \cdot  c \cdot  \dfrac{\sqrt{3}}{2}$ ,解得 $c = 2$ ,则 $b =  2\sqrt{3}$.
        \end{enumerate}
    \end{jieda}
    \item (2025$\cdot$嘉定区一模$\cdot$18)已知 $f\left( x\right)  = 2\cos \left( {{\omega x} + \dfrac{3\pi }{4}}\right)$ ，其中 $\omega  > 0$.
    \begin{enumerate}
        \item 若 $\omega  = 2$ ，求函数 $y = f\left( x\right) ,x \in  \left\lbrack  {-\dfrac{\pi }{4},\dfrac{\pi }{4}}\right\rbrack$ 的值域
        \item 若 $f\left( \dfrac{\pi }{4}\right)  = 0$ ，且函数 $y = f\left( x\right)$ 在 $\left( {\dfrac{\pi }{4},\dfrac{\pi }{3}}\right)$ 内有极小值，但无极大值，求 $\omega$ 的值.
    \end{enumerate}
    \begin{jieda}
        \begin{enumerate}
            \item $\omega  = 2,f\left( x\right)  = 2\cos \left( {{2x} + \dfrac{3\pi }{4}}\right)$ ，因为 $x \in  \left\lbrack  {-\dfrac{\pi }{4},\dfrac{\pi }{4}}\right\rbrack$ ， 所以 ${2x} + \dfrac{3\pi }{4} \in  \left\lbrack  {\dfrac{\pi }{4},\dfrac{5\pi }{4}}\right\rbrack$ ,所以 $\cos \left( {{2x} + \dfrac{3\pi }{4}}\right)  \in  \left\lbrack  {-1,\dfrac{\sqrt{2}}{2}}\right\rbrack$ , 因此函数 $y = f\left( x\right) ,x \in  \left\lbrack  {-\dfrac{\pi }{4},\dfrac{\pi }{4}}\right\rbrack$ 的值域为 $\left\lbrack  {-2,\sqrt{2}}\right\rbrack$ ;
            
            \item $\omega \cdot \dfrac{\pi }{4} + \dfrac{3\pi }{4} = {2k\pi } + \dfrac{\pi }{2},k \in  \mathbf{Z}$ ,可得 $\omega  = {8k} - 1,k \in  \mathbf{Z}$
            
            $\left\{  \begin{array}{l} \dfrac{\pi }{4} + \dfrac{1}{4}T < \dfrac{\pi }{3} \\  \dfrac{\pi }{4} + \dfrac{3}{4}T \geq  \dfrac{\pi }{3} \end{array}\right.$ ，所以 $\left\{  \begin{array}{l} \dfrac{\pi }{4} + \dfrac{1}{4} \cdot  \dfrac{2\pi }{\omega } < \dfrac{\pi }{3} \\  \dfrac{\pi }{4} + \dfrac{3}{4} \cdot  \dfrac{2\pi }{\omega } \geq  \dfrac{\pi }{3} \end{array}\right.$ ，即 $6 < \omega  \leq  {18}$  
            
            因此$\omega  = 7$ 或 $15$.
        \end{enumerate}
    \end{jieda}
\end{enumerate}

\section*{第 6 章 数 列}

\begin{center}

\begin{tabular}{|p{4cm}|p{12cm}|}
\hline
考点 & 课标要求与命题趋势 \\
\cline{1-2}
1. 数列的增减性 & 通过日常生活和数学中的实例, 了解数列的概念和表示方法 (列表、图像、通项公式), 理解数列是一种特殊的函数; 能利用数列的周期性、单调性解决简单的问题. 该内容是上海高考的必考内容，常考查利用与关系求通项或项及通项公式构造的相关应用，需综合复习 \\
\cline{1-2}
2. 数列的递推公式 & 能利用等差数列、等比数列的定义以及递推关系求数列的通项公式 \\
\cline{1-2}
3. 等差数列及其前 $n$ 项和 & \ding{172}  通过生活中的实例，理解等差数列的概念和通项公式的意义 \ding{173}  探索并掌握等差数列的变化规律，建立并熟练掌握通项公式和前 $n$ 项和公式，该内容是上海高考的必考内容，一般给出数列为等差数列，或通过构造为等差数列，求通项公式及前 $n$ 项和. 能运用等差数列解决简单的实际问题和数学问题,感受数学模型的现实意义与应用，需综合复习 \ding{174}  能在具体的问题情境中发现数列的等差关系，并能用等差数列的有关知识解决相应的问题 \ding{175}  体会等差数列与一元一次函数的关系，感受数列与函数的共性与差异，体会数学的整体性 \\
\cline{1-2}
4. 等比数列及其前 $n$ 项和 & \ding{172}  通过生活中的实例，理解并掌握等比数列的概念和通项公式的意义 \ding{173}  探索并掌握等比数列的前 $n$ 项和公式. 该内容是上海高考的必考内容,一般给出的数列为等比数列,或通过构造为等比数列,求通项公式及前 $n$ 项和,需综合复习 \ding{174}  能在具体的问题情境中，发现数列的等比关系，并解决相应的实际问题和数学问题， 感受数学模型的现实意义与应用 \ding{175}  体会等比数列与指数函数的关系，感受数列与函数的共性与差异，体会数学的整体性 \\
\cline{1-2}
5. 数列中的数学文化 & 数列是一类特殊的函数, 是数学重要的研究对象, 是研究其他类型函数的基本工具, 在日常生活中也有着广泛的应用 \\
\cline{1-2}
6. 数列求和 & 熟练掌握裂项相消求和和错位相减求和. 该内容是上海高考的常考内容, 常考查裂项相消求和、错位相减求和、奇偶并项求和, 需重点综合复习 \\
\cline{1-2}
7. 数列中的不等式. 最值及范围问题 & 在了解等差数列与一元一次函数、等比数列与指数函数的联系的基础上,综合利用函数、不等式的思想方法, 解决数列中的不等式、最值及范围问题 \\
\cline{1-2}
8. 数列与其他知识点的关联问题 &  研究数列通项、求和的方法灵活多样, 可以与向量、函数、三角、概率统计以及新定义知识等结合, 往往难度较大. 如与函数结合, 解决问题的主要思维方法有: \ding{172}  基本方法求通项:定义法、累乘法； \ding{173}  不等式的证明，借助构造函数利用导数分析单调性，求最值； \ding{174}  新定义考查，如结合导数的最值分析和不等式的放缩思维 \\
\cline{1-2}
\hline
\end{tabular}
\end{center}
\clearpage
\subsection*{考点1 数列的增减性}

\subsubsection*{一、填空题}

1. (2025$\cdot$嘉定区一模$\cdot$8)已知数列 $\left\{  {a}_{n}\right\}$ 的通项公式为 ${a}_{n} =  - n + c$ ，其中 $c$ 为常数，设数列 $\left\{  {a}_{n}\right\}$ 的前 $n$ 项和为 ${S}_{n}$ ，若 ${S}_{6} > {S}_{5}$ 且 ${S}_{6} > {S}_{7}$ ，则 $c$ 的取值范围为\blkda[2]{}.

2.(2024$\cdot$虹口区二模$\cdot$8)已知等比数列 $\left\{  {a}_{n}\right\}$ 是严格减数列，其前 $n$ 项和为 ${S}_{n}$ ， ${a}_{1} = 2$ ，若 ${a}_{1},2{a}_{2},3{a}_{3}$ 成等差数列,则 $\mathop{\lim }\limits_{{n \rightarrow  \infty }}{S}_{n} =$ \blkda[2]{}.

3.(2024$\cdot$普陀区一模$\cdot$9)若数列 $\left\{  {a}_{n}\right\}$ 满足 ${a}_{1} = {12},{a}_{n + 1} = {a}_{n} + {2n}$ ( $n \geq  1$ ， $n \in  \mathbf{N}$ ，则 $\dfrac{{a}_{n}}{n}$ 的最小值是\blkda[2]{}.

\subsubsection*{二、选择题}

4.(2025$\cdot$宝山区一模$\cdot$16)设 $\bigtriangleup {A}_{n}{B}_{n}{C}_{n}$ 的三边长分别为 ${a}_{n}$ 、 ${b}_{n}$ 、 ${c}_{n}$ ，面积为 ${S}_{n}$ ( $n$ 为正整数). 若 ${b}_{1}$ 一 ${c}_{1} = \dfrac{1}{2}{a}_{1}$ ,其中 ${c}_{1} > \dfrac{1}{4}{a}_{1},{a}_{n + 1} = {a}_{n},{b}_{n + 1} = {c}_{n} + \dfrac{1}{4}{a}_{n},{c}_{n + 1} = {b}_{n} + \dfrac{1}{4}{a}_{n}$ ,则 \dotfill{(\kh{})}

\xx{}{}{}{}$\left\{  {S}_{n}\right\}$ 为严格减数列

B. $\left\{  {S}_{n}\right\}$ 为严格增数列

C. $\left\{  {S}_{{2n} - 1}\right\}$ 为严格增数列, $\left\{  {S}_{2n}\right\}$ 为严格减数列

D. $\left\{  {S}_{{2n} - 1}\right\}$ 为严格减数列, $\left\{  {S}_{2n}\right\}$ 为严格增数列

5.(2024$\cdot$徐汇区一模$\cdot$16)已知数列 $\left\{  {a}_{n}\right\}$ 为无穷数列. 若存在正整数 $l$ ，使得对任意的正整数 $n$ ，均有 ${a}_{n + l} \leq  {a}_{n}$ ,则称数列 $\left\{  {a}_{n}\right\}$ 为“$l$ 阶弱减数列”. 有以下两个命题: \ding{172}  数列 $\left\{  {b}_{n}\right\}$ 为无穷数列且 ${b}_{n} =  \cos n - \dfrac{n}{2}\left( {n\text{ 为正整数 }}\right)$ ,则数列 $\left\{  {b}_{n}\right\}$ 是“$l$ 阶弱减数列”的充要条件是 $l \geq  4$ ;  \ding{173}  数列 $\left\{  {c}_{n}\right\}$ 为无穷数列且 ${c}_{n} = {an} + \dfrac{1 - {q}^{n}}{1 - q}\left( {n\text{ 为正整数 }}\right)$ ,若存在 $a \in  \mathbf{R}$ ,使得数列 $\left\{  {c}_{n}\right\}$ 是“2 阶弱减数列”,则 $- 1 \leq  q < 1$. 那么 \dotfill{(\kh{})}

\xx{}{}{}{} \ding{172} 是真命题， \ding{173} 是假命题 B.  \ding{172} 是假命题， \ding{173} 是真命题

C.  \ding{172}  \ding{173} 都是真命题 D.  \ding{172}  \ding{173} 都是假命题

6.(2025$\cdot$闵行区一模$\cdot$16)已知数列 $\left\{  {a}_{n}\right\}$ 满足 ${a}_{n + 1} = \left| {{a}_{n} + 1}\right|  + \lambda \left| {{a}_{n} - 1}\right|$ ，其中 $\lambda$ 为常数. 对于下述两个命题:

 \ding{172}  对于任意的 $\lambda  > 0$ ，任意的 ${a}_{1} \in  \mathbf{R}$ ，都有 $\left\{  {a}_{n}\right\}$ 是严格增数列；

 \ding{173}  对于任意的 $\lambda  < 0$ ，存在 ${a}_{1} \in  \mathbf{R}$ ，使得 $\left\{  {a}_{n}\right\}$ 是严格减数列.

以下说法正确的为 \dotfill{(\kh{})}

\xx{}{}{}{} \ding{172} 真命题， \ding{173} 假命题 B.  \ding{172} 假命题， \ding{173} 真命题

C.  \ding{172}  \ding{173} 均为真命题 D.  \ding{172}  \ding{173} 均为假命题



7.(2025$\cdot$嘉定区一模$\cdot$16)已知数列 $\left\{  {a}_{n}\right\}$ 满足 ${a}_{n + 1} = r{a}_{n}\left( {1 - {a}_{n}}\right) \left( {n = 1,2,3,\cdots }\right)$ ， ${a}_{1} \in  \left( {0,1}\right)$ ，给出以下四个结论:

 \ding{172}  当 $r = 2$ 时，只存在有限个 ${a}_{1}$ ，使得对任意正整数 $n$ ，都有 ${a}_{n + 1} > {a}_{n}$ ；

 \ding{173}  当 $r = 2$ 时，存在 ${a}_{1}$ 和正整数 $P$ ，当 $n > P$ 时， ${a}_{n + 1} - {a}_{n} < \dfrac{1}{2025}$ ；

 \ding{174}  当 $r = 3$ 时，存在 ${a}_{1}$ 和正整数 $P$ ，当 $n > P$ 时， ${a}_{n + 1} = {a}_{n}$ ；

 \ding{175}  当 $r =  - 3$ 时,不存在 ${a}_{1}$ ,使得对任意正整数 $n$ ,且 $n \geq  3$ ,都有 ${a}_{n} > 0$. 其中正确结论是 \dotfill{(\kh{})}

\xx{}{}{}{} \ding{172}  \ding{173}  B.  \ding{173}  \ding{174}  C.  \ding{174}  \ding{175}  D.  \ding{173}  \ding{175} 

\subsubsection*{三、解答题}

8.(2024$\cdot$金山区一模$\cdot$18)已知数列 $\left\{  {a}_{n}\right\}$ 满足 ${\log }_{2}{a}_{n + 1} = 1 + {\log }_{2}{a}_{n}$ ，且 ${a}_{1} = 2$.

(1)求 ${a}_{10}$ 的值;

(2)若数列 $\left\{  {{a}_{n} + \dfrac{\lambda }{{a}_{n}}}\right\}$ 为严格增数列,其中 $\lambda$ 是常数,求 $\lambda$ 的取值范围.

考点2 数 数列的递推公式

\subsubsection*{一、填空题}

1.(2024$\cdot$静安区一模$\cdot$5)计算 $\mathop{\sum }\limits_{{i = 1}}^{{+\infty }}{\left( \dfrac{2}{3}\right) }^{i} =$ \blkda[2]{}.

2.(2024$\cdot$徐汇区二模$\cdot$5)已知数列 $\left\{  {a}_{n}\right\}$ 的前 $n$ 项和为 ${S}_{n}$ ，若 ${S}_{n} = \dfrac{3}{2}{a}_{n} - \dfrac{1}{2}$ ( $n$ 是正整数)，则 ${a}_{5} =$ \blkda[2]{}.

3. (2024$\cdot$宝山区二模$\cdot$8)在数列 $\left\{  {a}_{n}\right\}$ 中， ${a}_{1} = 2$ ，且 ${a}_{n} = {a}_{n - 1} + \lg \dfrac{n}{n - 1}\left( {n \geq  2}\right)$ ，则 ${a}_{100} =$ \blkda[2]{}.

4.(2024$\cdot$青浦区一模$\cdot$11)已知数列 $\left\{  {a}_{n}\right\}$ 的通项公式为 ${a}_{n} = \left| {n - {18}}\right|$ ，记 ${S}_{n} = \mathop{\sum }\limits_{{i = 1}}^{n}{a}_{i}$ ，若 ${S}_{n + {30}} - {S}_{n} =$ 225，则正整数 $n$ 的值为\blkda[2]{}.

5.(2024$\cdot$普陀区二模$\cdot$11)设 $k$ 、 $m$ 、 $n$ 是正整数， ${S}_{n}$ 是数列 $\left\{  {a}_{n}\right\}$ 的前 $n$ 项和， ${a}_{1} = 2$ ， ${S}_{n} = {a}_{n + 1} + 1$ ，若 $m = \mathop{\sum }\limits_{{i = 1}}^{k}{t}_{i}\left( {{S}_{i} - 1}\right)$ ,且 ${t}_{i} \in  \{ 0,1\}$ ,记 $f\left( m\right)  = {t}_{1} + {t}_{2} + \cdots  + {t}_{k}$ ,则 $f\left( {2024}\right)  =$ \blkda[2]{}.

6.(2025$\cdot$黄浦区一模$\cdot$12)设常数 $b$ 为整数，数列 $\left\{  {a}_{n}\right\}$ 的通项公式为 ${a}_{n} = {\left( n + b\right) }^{2} + \dfrac{b}{2}$ ，若 ${a}_{m} + {a}_{m + 1} +  {a}_{m + 2}\left( {m \geq  1,m \in  \mathbf{Z}}\right)$ 的最小值为 -7,则 $b =$ \blkda[2]{}.

7.(2024$\cdot$浦东新区一模$\cdot$12)已知数列 $\left\{  {a}_{n}\right\}$ 满足 ${a}_{1} = a$ ，且对任意正整数 $n$ ，关于 $x$ 的实系数方程 ${x}^{2} + 2\sqrt{{a}_{n + 1}}x + 2{a}_{n}\left( {2 - {a}_{n}}\right)  = 0$ 都有两个相等的实根. 若 ${a}_{2024} = 0$ ,则满足条件的不同实数 $a$ 的个数为\blkda[2]{}个.

\subsubsection*{二、选择题}

8.(2025$\cdot$上海秋季卷$\cdot$16)已知数列 $\left\{  {a}_{n}\right\}$ 、 $\left\{  {b}_{n}\right\}$ 、 $\left\{  {c}_{n}\right\}$ 的通项公式分别为 ${a}_{n} = {10n} - 9,{b}_{n} = {2}^{n},{c}_{n} = \lambda {a}_{n} +  \left( {1 - \lambda }\right) {b}_{n}$. 若对任意的 $\lambda  \in  \left\lbrack  {0,1}\right\rbrack$ ， ${a}_{n}$ 、 ${b}_{n}$ 、 ${c}_{n}$ 的值均能构成三角形，则满足条件的正整数 $n$ 有\blkda[2]{}匈椒

\xx{}{}{}{}4 个 B. 3 个 C. 1 个 D. 无数个

9.(2025$\cdot$嘉定区二模$\cdot$16)设数列 $\left\{  {a}_{n}\right\}$ 满足 ${a}_{n} = \sin \left( \dfrac{n\pi }{3}\right)  + {\left( -1\right) }^{n}\cos \left( \dfrac{n\pi }{4}\right)$ ，记其前 $n$ 项和为 ${S}_{n}$ ，前 $n$ 项积为 ${T}_{n}$ ,则下列结论正确的是 \dotfill{(\kh{})}

\xx{}{}{}{}数列 $\left\{  {S}_{n}\right\}$ 和数列 $\left\{  {T}_{n}\right\}$ 均不是周期数列

B. 数列 $\left\{  {S}_{n}\right\}$ 是周期数列,数列 $\left\{  {T}_{n}\right\}$ 不是周期数列

C. 数列 $\left\{  {S}_{n}\right\}$ 不是周期数列,数列 $\left\{  {T}_{n}\right\}$ 是周期数列

D. 数列 $\left\{  {S}_{n}\right\}$ 和数列 $\left\{  {T}_{n}\right\}$ 均为周期数列

10.(2024$\cdot$宝山区二模$\cdot$16)数列 $\left\{  {a}_{n}\right\}$ 中， ${S}_{n}$ 是其前 $n$ 项的和，若对任意正整数 $n$ ，总存在正整数 $m$ ， 使得 ${S}_{n} = {a}_{m}$ ,则称数列 $\left\{  {a}_{n}\right\}$ 为 “某数列”. 现有如下两个命题:  \ding{172}  等比数列 $\left\{  {2}^{n}\right\}$ 为 “某数列”;  \ding{173}  对任意的等差数列 $\left\{  {a}_{n}\right\}$ ，总存在两个“某数列” $\left\{  {b}_{n}\right\}$ 和 $\left\{  {c}_{n}\right\}$ ，使得 ${a}_{n} = {b}_{n} + {c}_{n}$. 下列选项中正确的是 \dotfill{(\kh{})}

\xx{}{}{}{} \ding{172}  \ding{173} 均为真命题 B.  \ding{172} 为真命题， \ding{173} 为假命题

C.  \ding{172} 为假命题， \ding{173} 为真命题 D.  \ding{172}  \ding{173} 均为假命题

11.(2025$\cdot$徐汇区一模$\cdot$16)已知数列 $\left\{  {a}_{n}\right\}$ 的前 $n$ 项和为 ${S}_{n}$ ，设 ${t}_{n} = \dfrac{{S}_{n}}{n}$ ( $n$ 为正整数). 若存在常数 $c$ ， 使得任意两两不相等的正整数 $i\text{ 、 }j\text{ 、 }k$ ,都有 $\left( {i - j}\right) {t}_{k} + \left( {j - k}\right) {t}_{i} + \left( {k - i}\right) {t}_{j} = c$ ,则称数列 $\left\{  {a}_{n}\right\}$ 为 “轮换均值数列”. 现有下列两个命题: \ding{172}  任意等差数列 $\left\{  {a}_{n}\right\}$ 都是“轮换均值数列”， \ding{173}  存在公比不为 1 的等比数列 $\left\{  {b}_{n}\right\}$ 是“轮换均值数列”. 则下列说法正确的是 \dotfill{(\kh{})}

\xx{}{}{}{} \ding{172} 是真命题， \ding{173} 是假命题 B.  \ding{172} 是假命题， \ding{173} 是真命题

C.  \ding{172}  \ding{173} 都是真命题 D.  \ding{172}  \ding{173} 都是假命题

12.(2025$\cdot$杨浦区一模$\cdot$16)设无穷数列 $\left\{  {a}_{n}\right\}$ 的前 $n$ 项和为 ${S}_{n}$ ，且对任意的正整数 $n$ ， ${a}_{n + 1} = \dfrac{{S}_{n}}{{a}_{n}}$ ，则 $\mathop{\sum }\limits_{{i = 1}}^{5}{a}_{2i} - \mathop{\sum }\limits_{{i = 1}}^{6}{a}_{{2i} - 1}$ 的值可能为 \dotfill{(\kh{})}

\xx{}{}{}{}-6 B. 0

C. 6 D. 12

13.(2025$\cdot$长宁区一模$\cdot$16)数列 $\left\{  {a}_{n}\right\}$ 为严格增数列，且对任意的正整数 $n$ ，都有 $\dfrac{{a}_{n + 1}}{n + 1} \geq  \dfrac{{a}_{n}}{n}$ ，则称数列 $\left\{  {a}_{n}\right\}$ 满足“性质 $\Omega$ ”. 现给出两个命题:

 \ding{172}  存在等差数列 $\left\{  {a}_{n}\right\}$ 满足“性质 $\Omega$ ”；

 \ding{173}  任意等比数列 $\left\{  {a}_{n}\right\}$ ，若首项 ${a}_{1} > 0$ ，则 $\left\{  {a}_{n}\right\}$ 满足“性质 $\Omega$ ”.

下列判断正确的是 \dotfill{(\kh{})}

\xx{}{}{}{} \ding{172}  \ding{173} 都是真命题 B.  \ding{172} 是真命题， \ding{173} 是假命题

C.  \ding{172} 是假命题， \ding{173} 是真命题 D.  \ding{172}  \ding{173} 都是假命题

14.(2024$\cdot$长宁区二模$\cdot$16)设数列 $\left\{  {a}_{n}\right\}$ 的前 $n$ 项和为 ${S}_{n}$ ，若存在非零常数 $c$ ，使得对任意正整数 $n$ ， 都有 $2\sqrt{{S}_{n}} = {a}_{n} + c$ ,则称数列 $\left\{  {a}_{n}\right\}$ 具有性质 $p$. 现有如下两个命题: \ding{172}  存在等差数列 $\left\{  {a}_{n}\right\}$ 具有性质 $p$ ;  \ding{173}  不存在等比数列 $\left\{  {a}_{n}\right\}$ 具有性质 $p$. 下列判断正确的是 \dotfill{(\kh{})}

\xx{}{}{}{} \ding{172}  \ding{173} 均是真命题 B.  \ding{172} 是真命题， \ding{173} 是假命题

C.  \ding{172} 是假命题， \ding{173} 是真命题 D.  \ding{172}  \ding{173} 均是假命题

15.(2025$\cdot$奉贤区一模$\cdot$16)已知数列 $\left\{  {a}_{n}\right\}$ 不是常数列，前 $n$ 项和为 ${S}_{n}$ ， ${a}_{n} > 0$. 若对任意正整数列，存在正整数 $m$ ,使得 $\left| {{S}_{n} - {a}_{m}}\right|  < {a}_{1}$ ,则称 $\left\{  {a}_{n}\right\}$ 是 “可控数列”. 现给出两个命题:

 \ding{172}  若各项均为正整数的等差数列 $\left\{  {a}_{n}\right\}$ 满足: $d = 3$ ，则 $\left\{  {a}_{n}\right\}$ 是“可控数列”；

 \ding{173}  若等比数列 $\left\{  {a}_{n}\right\}$ 是“可控数列”，则其公比为 $q \in  \left( {0,1}\right)$.

下列判断正确的是 \dotfill{(\kh{})}

\xx{}{}{}{} \ding{172} 与 \ding{173} 均为真命题 B.  \ding{172} 与 \ding{173} 均为假命题

C.  \ding{172} 为假命题， \ding{173} 为真命题 D.  \ding{172} 为真命题， \ding{173} 为假命题

16.(2024$\cdot$黄浦区二模$\cdot$16)设数列 $\left\{  {a}_{n}\right\}$ 的前 $n$ 项和为 ${S}_{n}$ ，若对任意的正整数 $n$ ， ${S}_{n}$ 都是数列 $\left\{  {a}_{n}\right\}$ 中的项，则称数列 $\left\{  {a}_{n}\right\}$ 为“ $T$ 数列”. 对于命题: \ding{172}  存在“ $T$ 数列” $\left\{  {a}_{n}\right\}$ ，使得数列 $\left\{  {S}_{n}\right\}$ 为公比不为 1 的等比数列;  \ding{173}  对于任意的实数 ${a}_{1}$ ,都存在实数 $d$ ,使得以 ${a}_{1}$ 为首项、 $d$ 为公差的等差数列 $\left\{  {a}_{n}\right\}$ 为“ $T$ 数列”. 下列判断正确的是 \dotfill{(\kh{})}

\xx{}{}{}{} \ding{172} 和 \ding{173} 均为真命题 B.  \ding{172} 和 \ding{173} 均为假命题

C.  \ding{172} 是真命题， \ding{173} 是假命题 D.  \ding{172} 是假命题， \ding{173} 是真命题

\subsection*{考点 3 等差数列及其前 $n$ 项和}

\subsubsection*{一、填空题}

1.(2025$\cdot$奉贤区二模$\cdot$1)等差数列首项为 1 ，公差是 3 ，则第 5 项等于\blkda[2]{}.

2.(2025$\cdot$上海秋季卷$\cdot$3)已知等差数列 $\left\{  {a}_{n}\right\}$ 的首项 ${a}_{1} =  - 3$ ，公差 $d = 2$ ，则该数列的前 6 项和为\blkda[2]{}.

3.(2024$\cdot$浦东新区二模$\cdot$3)已知等差数列 $\left\{  {a}_{n}\right\}$ 满足 ${a}_{1} + {a}_{6} = {12},{a}_{4} = 7$ ，则\blkda[2]{}。

4.(2024$\cdot$宝山区一模$\cdot$3)已知等差数列 $\left\{  {a}_{n}\right\}$ 的前 $n$ 项和为 ${S}_{n}$ ，若 ${a}_{4} + {a}_{13} = 1$ 则 ${S}_{16} =$ \blkda[2]{}.

5.(2025$\cdot$普陀区二模$\cdot$4)设 $n \geq  1,n \in  \mathbf{N}$ ， ${S}_{n}$ 是等差数列 $\left\{  {a}_{n}\right\}$ 的前 $n$ 项和，若 ${a}_{3n} = 3{a}_{n} \neq  0$ ，则 $\dfrac{{S}_{5}}{{a}_{10}}$ 的值为\blkda[2]{}.

6.(2025$\cdot$静安区一模$\cdot$4)设 $\left\{  {a}_{n}\right\}$ 是等差数列， ${a}_{1} =  - 6,{a}_{3} = 0$ ，则该数列的前 8 项的和 ${S}_{8}$ 的值为\blkda[2]{}.

7.(2024$\cdot$崇明区二模$\cdot$5)若等差数列 $\left\{  {a}_{n}\right\}$ 的首项 ${a}_{1} = 1$ ，前 5 项和 ${S}_{5} = {25}$ ，则 ${a}_{5} =$ \blkda[2]{}.

8.(2025$\cdot$普陀区一模$\cdot$5)设 $n \geq  1$ ， $m \geq  1$ ， $m$ 、 $n \in  \mathbf{N}$ ，等差数列 $\left\{  {a}_{n}\right\}$ 的首项 ${a}_{1} = 0$ ，公差 $d \neq  0$ ，若 ${a}_{m} =  \mathop{\sum }\limits_{{i = 1}}^{{11}}{a}_{i}$ ，则 $m$ 的值为\blkda[2]{}.

9.(2025$\cdot$浦东新区二模$\cdot$6)设数列 $\left\{  {a}_{n}\right\}$ 为等差数列，其前 $n$ 项和为 ${S}_{n}$ ，已知 ${a}_{3} + {a}_{15} = 2$ ，则 ${S}_{17} =$ \blkda[2]{}.

10.(2024$\cdot$上海春季卷$\cdot$6)在等差数列 $\left\{  {a}_{n}\right\}$ 中， ${a}_{n} = n + c$ ， ${S}_{n}$ 是数列 $\left\{  {a}_{n}\right\}$ 的前 $n$ 项和，若 ${S}_{7} < 0$ ，则 $c$ 的取值范围是\blkda[2]{}.

11.(2025$\cdot$青浦区一模$\cdot$6)已知数列 $\left\{  {a}_{n}\right\}$ 满足 ${a}_{1} + 2{a}_{2} + 3{a}_{3} + \cdots  + n{a}_{n} = n\left( {n + 2}\right)$ ，则 ${a}_{66} =$ \blkda[2]{}.

12.(2024$\cdot$松江区二模$\cdot$7)已知等差数列 $\left\{  {a}_{n}\right\}$ 的公差为 2，前 $n$ 项和为 ${S}_{n}$ ，若 ${a}_{3} = {S}_{5}$ ，则使得 ${S}_{n} < {a}_{n}$ 成立的 $n$ 的最大值为\blkda[2]{}.

13.(2025$\cdot$浦东新区一模$\cdot$7)若等差数列 $\left\{  {a}_{n}\right\}$ 满足 ${a}_{7} + {a}_{8} + {a}_{9} = 0,{a}_{7} + {a}_{10} = 1$ ，则 ${a}_{1} =$ \blkda[2]{}.

14.(2025$\cdot$闵行区二模$\cdot$8)已知 $\left\{  {a}_{n}\right\}$ 是等差数列，若 ${a}_{3}$ 、 ${a}_{7}$ 分别是函数 $y = {x}^{2} - {4x} + 2$ 的两个零点， 则 ${a}_{5} =$ \blkda[2]{}.

15.(2025$\cdot$金山区一模$\cdot$8)已知 ${S}_{n}$ 是等差数列 $\left\{  {a}_{n}\right\}$ 的前 $n$ 项和，若 $2{a}_{1} + 3{a}_{11} = {20}$ ，则 ${S}_{13}$ 的值为\blkda[2]{}.

16.(2024$\cdot$黄浦区二模$\cdot$10)已知数列 $\left\{  {a}_{n}\right\}$ 是给定的等差数列，其前 $n$ 项和为 ${S}_{n}$ ，若 ${a}_{9}{a}_{10} < 0$ ，且当 $m =  {m}_{0}$ 与 $n = {n}_{0}$ 时， $\left| {{S}_{m} - {S}_{n}}\right|$ ( $m,n \in  \left\{  {x \mid  x \leq  {30},x \in  {\mathbf{N}}^{ * }}\right\}$ )取得最大值，则 $\left| {{m}_{0} - {n}_{0}}\right|$ 的值为\blkda[2]{}.

17.(2025$\cdot$黄浦区二模$\cdot$11)设 $\left\{  {a}_{n}\right\}$ 为等差数列，其前 $n$ 项和为 ${S}_{n}$ ，若 $\left( {{S}_{8} - {S}_{7}}\right)$ ( ${S}_{9} - {S}_{7}$ ) $< 0$ ，则满足 ${S}_{m}{S}_{m + 1} < 0$ 的正整数 $m =$ \blkda[2]{}.

18.(2024$\cdot$杨浦区二模$\cdot$12)已知实数 $a$ 满足: \ding{172}  $a \in  \lbrack 0,{2\pi }$ )； \ding{173}  存在实数 $b$ 、 $c\left( {a < b < c < {2\pi }}\right)$ ，使得 $a,b,c$ 是等差数列， $\cos b,\cos a,\cos c$ 也是等差数列，则实数 $a$ 的取值范围是\blkda[2]{}.

\subsubsection*{二、选择题}

19.(2025$\cdot$松江区一模$\cdot$14)渐进式延迟退休方案是指采取较缓而稳妥的方式逐步延长退休年龄. 对于男职工，新方案将延迟法定退休年龄每 4 个月延迟 1 个月，逐步将男职工的法定退休年龄从原 60 周岁延迟至 63 周岁. 如果男职工延迟法定退休年龄部分对照表如下表所示:

\begin{center}

\begin{tabular}{|p{4cm}|p{12cm}|}
\hline
出生时间 & 改革后法定退休年龄 \\
\cline{1-2}
1965 年 1 月— 4 月 & 60 岁+1 个月 \\
\cline{1-2}
1965 年 5 月—8 月 & 60 岁+2 个月 \\
\cline{1-2}
1965 年 9 月—12 月 & 60 岁+3 个月 \\
\cline{1-2}
1966 年 1 月—4 月 & 60 岁+4 个月 \\
\cline{1-2}
...... &...... \\
\cline{1-2}
\hline
\end{tabular}
\end{center}

那么 1974 年 5 月出生的男职工退休年龄为 \dotfill{(\kh{})}

\xx{}{}{}{}62 岁 3 个月 B. 62 岁 4 个月 C. 62 岁 5 个月 D. 63 岁

20. (2024$\cdot$奉贤区一模$\cdot$16)已知等差数列 $\left\{  {a}_{n}\right\}$ 的前 $n$ 项和为 ${S}_{n}$ ，且关于正整数 $n$ 的不等式(S，m- ${2022})\left( {{S}_{n + 1} - {2022}}\right)  < 0$ 与不等式 $\left( {{S}_{n} - {2023}}\right) \left( {{S}_{n + 1} - {2023}}\right)  < 0$ 的解集均为 $M$.

命题 $\alpha$ : 集合 $M$ 中元素的个数一定是偶数个;

命题 $\beta$ : 若数列 $\left\{  {a}_{n}\right\}$ 的公差 $d > 0$ ,且 ${n}_{0} \in  M$ ,则 ${a}_{{n}_{0} + 1} > 1$.

下列说法中正确的是 \dotfill{(\kh{})}

\xx{}{}{}{}命题 $\alpha$ 是真命题，命题 $\beta$ 是假命题 B. 命题 $\alpha$ 是假命题，命题 $\beta$ 是真命题

C. 命题 $\alpha$ 、命题 $\beta$ 都是假命题 D. 命题 $\alpha$ 、命题 $\beta$ 都是真命题

21.(2025$\cdot$杨浦区二模$\cdot$16)设 $A$ 是由 $k$ 个二次函数组成的集合，对于连续的正整数 $1,2,3,\cdots$ ， 2025，存在二次函数 $y = {f}_{i}\left( x\right)  \in  A$ (1 $\leq  i \leq  {2025}$ ， $i \in  \mathbf{Z}$ ， $y = {f}_{i}\left( x\right)$ 可重复)，使得 ${f}_{1}\left( 1\right)$ ， ${f}_{2}\left( 2\right)$ ， ${f}_{3}\left( 3\right) ,\cdots ,{f}_{2025}\left( {2025}\right)$ 是等差数列，则 $k$ 的最小可能值是 \dotfill{(\kh{})}

\xx{}{}{}{}507 B. 1013 C. 1519 D. 2025

22.(2025$\cdot$宝山区二模$\cdot$16)若对任意正整数 $n$ ，数列 $\left\{  {a}_{n}\right\}$ 的前 $n$ 项和 ${S}_{n}$ 都是完全平方数，则称数列 $\left\{  {a}_{n}\right\}$ 为“完全平方数列”. 有如下两个命题:  \ding{172}  若数列 $\left\{  {b}_{n}\right\}$ 的前 $n$ 项和 ${T}_{n} = {\left( n - t\right) }^{2}$ ( $t$ 为正整数),则使得数列 $\left\{  \left| {b}_{n}\right| \right\}$ 为“完全平方数列”的 $t$ 值有且仅有一个;  \ding{173}  存在无穷多个“完全平方数列”的等差数列. 对于以上两个命题, 下列选项中正确的是 \dotfill{(\kh{})}

\xx{}{}{}{} \ding{172}  \ding{173} 都是真命题 B.  \ding{172} 是真命题， \ding{173} 是假命题

C.  \ding{172} 是假命题， \ding{173} 是真命题 D.  \ding{172}  \ding{173} 都是假命题

\subsubsection*{三、解答题}

23.(2024$\cdot$虹口区二模$\cdot$17)已知等差数列 $\left\{  {a}_{n}\right\}$ 满足 ${a}_{2} = 5,{a}_{9} + 7 = 2{a}_{6}$.

(1)求 $\left\{  {a}_{n}\right\}$ 的通项公式；

(2)设数列 $\left\{  {b}_{n}\right\}$ 前 $n$ 项和为 ${S}_{n}$ ，且 ${b}_{n} = {a}_{n + 1}^{2} - {a}_{n}^{2}$ ，若 ${S}_{m} > {432}$ ，求正整数 $m$ 的最小值.

\subsection*{考点4 又 等比数列及其前 $n$ 项和}

\subsubsection*{一、填空题}

1.(2023$\cdot$上海秋季卷$\cdot$3)若数列 $\left\{  {a}_{n}\right\}$ 是首项为 3、公比为 2 的等比数列，则 ${S}_{6} =$ \blkda[2]{}.

2.(2024$\cdot$虹口区一模$\cdot$3)设等比数列 $\left\{  {a}_{n}\right\}$ 的前 $n$ 项和为 ${S}_{n}$ ，若 ${a}_{2} = 1,{S}_{2} = 4$ ，则 $\mathop{\lim }\limits_{{n \rightarrow  \infty }}{S}_{n} =$ \blkda[2]{}.

3. (2025$\cdot$嘉定区二模$\cdot$4)已知等比数列 $\left\{  {a}_{n}\right\}$ 的首项为 1，公比为 $q$ ，其前 $n$ 项和为 ${S}_{n}$. 若 ${S}_{3} > {S}_{4}$ ，则 $q$ 的取值范围为\blkda[2]{}.

4.(2024$\cdot$崇明区一模$\cdot$4)已知等比数列 $\left\{  {a}_{n}\right\}$ 首项 ${a}_{1} = 1$ ，公比 $q = 2$ ，则 ${S}_{5} =$ \blkda[2]{}.

5.(2024$\cdot$金山区二模$\cdot$5)设公比为 2 的等比数列 $\left\{  {a}_{n}\right\}$ 的前 $n$ 项和为 ${S}_{n}$ ，若 ${S}_{2024} - {S}_{2022} = 6$ ，则 ${a}_{2024} =$ \blkda[2]{}.

6.(2024$\cdot$闵行区二模$\cdot$6)等比数列 $\left\{  {a}_{n}\right\}$ 的前 $n$ 项和为 ${S}_{n}$ ，若 ${a}_{2} = 3,{a}_{5} = {81}$ ，则 ${S}_{5} =$ \blkda[2]{}.

7.(2024$\cdot$杨浦区二模$\cdot$6)各项为正的等比数列 $\left\{  {a}_{n}\right\}$ 满足 ${a}_{1} = 2,{a}_{2} + {a}_{3} = {12}$ ，则通项公式 ${a}_{n} =$ \blkda[2]{}.

8.(2024$\cdot$奉贤区一模$\cdot$6)已知数列 $\left\{  {a}_{n}\right\}$ 是各项为正的等比数列， ${a}_{1} = 1,{a}_{5} = 1$ ，则其前 10 项和 ${S}_{10} =$ \blkda[2]{}.

9.(2025$\cdot$黄浦区二模$\cdot$7)已知等比数列 $\left\{  {a}_{n}\right\}$ 为严格增数列，其前 $n$ 项和为 ${S}_{n}$ ，若 ${a}_{1}{a}_{10} = {a}_{8}$ ， ${S}_{4} - {S}_{1} = \dfrac{21}{4}$ ，则该数列的公比为\blkda[2]{}.

10.(2025$\cdot$上海春季卷$\cdot$7)已知 $\left\{  {a}_{n}\right\}$ 是首项为 1、公差为 1 的等差数列， $\left\{  {b}_{n}\right\}$ 是首项为 1、公比为 $q  \left( {q > 0}\right)$ 的等比数列. 若数列 $\left\{  {{a}_{n} \cdot  {b}_{n}}\right\}$ 的前三项和为 2，则 $q =$ \blkda[2]{}.

11.(2024$\cdot$静安区二模$\cdot$7)已知等比数列的前 $n$ 项和为 ${S}_{n} = {\left( \dfrac{1}{2}\right) }^{n} + a$ ，则 $a$ 的值为\blkda[2]{}.

12.(2025$\cdot$青浦区二模$\cdot$7)数列 $\left\{  {a}_{n}\right\}$ 中， ${a}_{1} = 1$ ， ${a}_{n} = 2{a}_{n + 1}$ ，则 $\mathop{\sum }\limits_{{i = 1}}^{{+\infty }}{a}_{i} =$ \blkda[2]{}.

13.(2024$\cdot$普陀区二模$\cdot$8)设等比数列 $\left\{  {a}_{n}\right\}$ 的公比为 $q\left( {n \geq  1,n \in  \mathbf{N}}\right)$ ，则 “12 ${a}_{2},{a}_{4},2{a}_{3}$ 成等差数列” 的一个充分非必要条件是\blkda[2]{}.

14.(2025$\cdot$松江区一模$\cdot$8)已知数列 $\left\{  {a}_{n}\right\}$ 是等比数列，若 ${\log }_{2}{a}_{1} + {\log }_{2}{a}_{4} = 3$ ， ${{2}}^{{a}_{2}} \cdot  {2}^{{a}_{3}} = {64}$ ，则 ${a}_{10} =$ \blkda[2]{}.

15.(2024$\cdot$闵行区一模$\cdot$11)已知数列 $\left\{  {a}_{n}\right\}$ 为无穷等比数列，若 $\mathop{\sum }\limits_{{i = 1}}^{{+\infty }}{a}_{i} =  - 2$ ，则 $\mathop{\sum }\limits_{{i = 1}}^{{+\infty }}\left| {a}_{i}\right|$ 的取值范围为\blkda[2]{}.

16.(2024$\cdot$上海秋季卷$\cdot$12)无穷等比数列 $\left\{  {a}_{n}\right\}$ 满足首项 ${a}_{1} > 0,q > 1$ ，记 ${I}_{n} = \left\{  {x - y \mid  x,y \in  \left\lbrack  {{a}_{1},{a}_{2}}\right\rbrack   \cup  }\right.  \left. \left\lbrack  {{a}_{n},{a}_{n + 1}}\right\rbrack  \right\}$ ,若对任意正整数 $n$ 集合 ${I}_{n}$ 是闭区间,则 $q$ 的取值范围是\blkda[2]{}.

\subsubsection*{二、选择题}

17.(2024$\cdot$青浦区二模$\cdot$15)设 ${S}_{n}$ 是首项为 ${a}_{1}$ 、公比为 $q$ 的等比数列 $\left\{  {a}_{n}\right\}$ 的前 $n$ 项和，且 ${S}_{2023} <  {S}_{2025} < {S}_{2024}$ ,则 \dotfill{(\kh{})}

\xx{}{}{}{}${a}_{1} > 0$ B. $q > 0$ C. $\left| {S}_{n}\right|  \leq  \left| {a}_{1}\right|$ D. $\left| {S}_{n}\right|  < \left| q\right|$

\subsubsection*{三、解答题}

18.(2024$\cdot$黄浦区一模$\cdot$17)已知等比数列 $\left\{  {a}_{n}\right\}$ 是严格增数列，其第 3、4、5 项的乘积为 1000，并且这三项分别乘以 4、3、2 后，所得的三个数依次成等差数列.

(1)求数列 $\left\{  {a}_{n}\right\}$ 的通项公式；

(2)若对任意的正整数 $n$ ，数列 $\left\{  {b}_{n}\right\}$ 的前 $n$ 项和 ${S}_{n} = 3\left( {1 - {2}^{n}}\right)$ ， 的量 $\left( {{a}_{n},{b}_{n}}\right)$ 的模为 ${t}_{n}$ ，求数列 $\left\{  {t}_{n}\right\}$ 的前 $n$ 项和.

19.(2024$\cdot$奉贤区二模$\cdot$17)已知 $\left\{  {a}_{n}\right\}$ 是公差 $d = 2$ 的等差数列，其前 5 项和为 15， $\left\{  {b}_{n}\right\}$ 是公比 $q$ 为实数的等比数列， ${b}_{1} = 1$ ， ${b}_{4} - {b}_{2} = 6$.

(1)求 $\left\{  {a}_{n}\right\}$ 和 $\left\{  {b}_{n}\right\}$ 的通项公式；

(2)设 ${c}_{n} = {2}^{{a}_{n}} + {b}_{2n}\left( {n \geq  1,n \in  \mathbf{N}}\right)$ ，计算 $\mathop{\sum }\limits_{{i = 1}}^{n}{c}_{i}$.

\subsection*{考点5 > 数列中的数学文化}

\subsubsection*{一、填空题}

1.(2024$\cdot$静安区一模$\cdot$11)在国家开发西部的号召下，某企业得到了一笔 400 万元的无息贷款用做设备更新. 据预测，该企业设备更新后，第 1 个月收入为 20 万元，在接下来的 5 个月中，每月收入都比上个月增长 20\%，从第 7 个月开始，每个月的收入都比前一个月增加 2 万元. 则从新设备使用开始计算，该企业用所得收入偿还 400 万无息贷款只需\blkda[2]{}个月. (结果取整)

2.(2024$\cdot$宝山区二模$\cdot$11)某区域的地形大致如下左图，某部门负责该区域的安全警戒，在哨位 $O$ 的正上方安装探照灯对警戒区域进行探查扫描.

假设 1: 警戒区域为空旷的扇环形平地 ${A}_{1}{A}_{n}{B}_{n}{B}_{1}$ ;

假设 2: 视探照灯为点 $M$ ,且距离地面 ${20}{m}$ ;

假设 3: 探照灯 $M$ 照射在地面上的光斑是椭圆.

当探照灯 $M$ 以某一俯角从 ${A}_{k}{A}_{k + 1}$ 侧扫描到 ${B}_{k}{B}_{k + 1}$ 侧时,记为一次扫描,此过程中照射在地面上的光斑形成一个扇环 ${S}_{k}\left( {k = 1,2,3,\cdots }\right)$. 由此,通过调整 $M$ 的俯角,逐次扫描形成扇环 ${S}_{1}\text{ 、 }{S}_{2}\text{ 、 }{S}_{3}\cdots$. 第一次扫描时,光斑的长轴为 ${EF},\left| {OE}\right|  = {30}{m}$ ,此时在探照灯 $M$ 处测得点 $F$ 的俯角为 ${30}^{ \circ  }$ (如下右图). 记 $\left| {{A}_{k}{A}_{k + 1}}\right|  = {d}_{k}$ ,经测量知 $\left| {{A}_{1}{A}_{n}}\right|  = {80}{m}$ ,且 $\left\{  {d}_{k}\right\}$ 是公差约为 ${0.1}{m}$ 的等差数列,则至少需要经过\blkda[2]{}次扫描，才能将整个警戒区域扫描完毕.

\begin{center}
\includegraphics[width=0.6\textwidth]{bodyxb/src/80_412_1248_765_374_0.jpg}
\end{center}


第 2 题图

3.(2024$\cdot$杨浦区二模$\cdot$11)某钢材公司积压了部分圆钢，经清理知共有 2024 根，每根圆钢的直径为 ${10}{{cm}}$. 现将它们堆放在一起. 若堆成纵断面为等腰梯形 (如图所示,每一层的根数比上一层根数多 1 根),且为杜绝安全隐患,堆放高度不得高于 $\dfrac{3}{2}{m}$ ，若堆放占用场地面积最小，则最下层圆钢根数为\blkda[2]{}.

\begin{center}
\includegraphics[width=0.2\textwidth]{bodyxb/src/80_1186_1735_237_135_0.jpg}
\end{center}


第 3 题图

\subsubsection*{二、解答题}

4.(2024$\cdot$虹口区一模$\cdot$19)2022 年 12 月底,某厂的废水池已储存废水 800 t，以后每月新产生的 2 t 废水也存入废水池. 该厂 2023 年开始对处理后的废水进行排放，1 月底排放 10 t 处理后的废水， 计划以后每月月底排放一次,每月排放处理后的废水比上月增加 $2{t}$.

(1)若按计划排放，该厂在哪一年的几月份排放后，可第一次将废水池中的废水排放完毕？

(2)该厂加强科研攻关，提升废水处理技术，经过深度净化的废水可以再次利用. 该厂从 2023 年 7 月开始对该月计划排放的废水进行深度净化，首次净化废水 $5{t}$ ，以后每月比上月提高 20\% 的净化能力. 试问: 哪一年的几月份开始,该厂当月排放的废水能被全部净化?

\subsubsection*{一、填空题}

1. (2024$\cdot$普陀区一模$\cdot$4)设 ${S}_{n}$ 是等差数列 $\left\{  {a}_{n}\right\}$ 的前 $n$ 项和( $n \geq  1$ ， $n \in  \mathbf{N}$ )，若 ${a}_{2} + {a}_{4} = 9 - {a}_{6}$ ，则 ${S}_{7} =$ \blkda[2]{}.

2.(2024$\cdot$杨浦区一模$\cdot$5)等差数列 $\left\{  {a}_{n}\right\}$ 中，若 ${a}_{1} + {a}_{2} = 6,{a}_{2} + {a}_{3} = {10}$ ，则 $\left\{  {a}_{n}\right\}$ 的前 10 项和为\blkda[2]{}.

3. (2024$\cdot$浦东新区一模$\cdot$5)已知 ${S}_{n}$ 是等差数列 $\left\{  {a}_{n}\right\}$ 的前 $n$ 项和，若 ${a}_{n} = {2n} - 3$ ，则满足 ${S}_{m} = {24}$ 的正整数 $m$ 的值为\blkda[2]{}.

4.(2025$\cdot$浦东新区二模$\cdot$12)已知数列 $\left\{  {a}_{n}\right\}$ ， ${a}_{1} = 1$ ， ${a}_{n} \in  \{ 1, - 1\}$ ， $\left( {n \geq  2}\right)$ ，并且前 $n$ 项的和 ${S}_{n}$ 满足:  \ding{172}  存在小于 1013 的正整数 $t$ ，使得 ${S}_{{2t} + 1} =  - 1$ ； \ding{173}  对任意的正整数 $k$ 和 $m$ ，都有 $\left| {{S}_{2m} - {S}_{{2k} - 1}}\right|  \leq$ 1. 满足以上条件的数列 $\left\{  {S}_{n}\right\}  \left( {1 \leq  n \leq  {2025}}\right)$ 共有\blkda[2]{}个.

\subsubsection*{二、选择题}

5.(2024$\cdot$松江区二模$\cdot$16)设 ${S}_{n}$ 为数列 $\left\{  {a}_{n}\right\}$ 的前 $n$ 项和，有以下两个命题: \ding{172}  若 $\left\{  {a}_{n}\right\}$ 是公差不为零的等差数列且 $k \in  \mathbf{N},k \geq  2$ ,则 ${S}_{1} \cdot  {S}_{2} \cdot  \cdots  \cdot  {S}_{{2k} - 1} = 0$ 是 ${a}_{1} \cdot  {a}_{2} \cdot  \cdots  \cdot  {a}_{k} = 0$ 的必要非充分条件;  \ding{173}  若 $\left\{  {a}_{n}\right\}$ 是等比数列且 $k \in  \mathbf{N},k \geq  2$ ,则 ${S}_{1} \cdot  {S}_{2} \cdot  \cdots  \cdot  {S}_{k} = 0$ 的充要条件是 ${a}_{k} + {a}_{k + 1} = 0$. 则下列判断中正确的是 \dotfill{(\kh{})}

\xx{}{}{}{} \ding{172} 是真命题， \ding{173} 是假命题 B.  \ding{172} 是假命题， \ding{173} 是真命题

C.  \ding{172}  \ding{173} 都是真命题 D.  \ding{172}  \ding{173} 都是假命题

6.(2023$\cdot$上海春季卷$\cdot$16)数列 $\left\{  {a}_{n}\right\}$ 的各项均为实数， ${S}_{n}$ 为其前 $n$ 项和，对任意 $k > 2$ 022 都有 $\left| {S}_{k}\right|  > \left| {S}_{k + 1}\right|$ ,则下列说法可能正确的是 \dotfill{(\kh{})}

\xx{}{}{}{}${a}_{1},{a}_{3},{a}_{5},\cdots ,{a}_{{2n} - 1}$ 为等差数到, ${a}_{2},{a}_{4},{a}_{6},\cdots ,{a}_{2n}$ 为等比数列

B. ${a}_{1},{a}_{3},{a}_{5},\cdots ,{a}_{{2n} - 1}$ 为等比数列, ${a}_{2},{a}_{4},{a}_{6},\cdots ,{a}_{2n}$ 为等差数列

C. ${a}_{1},{a}_{2},{a}_{3},\cdots ,{a}_{2022}$ 为等差数列, ${a}_{2022},{a}_{2023},{a}_{2024},\cdots ,{a}_{n}$ 为等比数列

D. ${a}_{1},{a}_{2},{a}_{3},\cdots ,{a}_{2022}$ 为等比数列, ${a}_{2022},{a}_{2023},{a}_{2024},\cdots ,{a}_{n}$ 为等差数列

\subsubsection*{三、解答题}

7.(2024$\cdot$徐汇区一模$\cdot$17)已知等差数列 $\left\{  {a}_{n}\right\}$ 的前 $n$ 项和为 ${S}_{n}$ ， ${a}_{1} = 2$ ， ${S}_{5} = {20}$.

(1)求数列 $\left\{  {a}_{n}\right\}$ 的通项公式；

(2)若等比数列 $\left\{  {b}_{n}\right\}$ 的公比为 $q = \dfrac{1}{2}$ ，且满足 ${a}_{4} + {b}_{4} = 9$ ，求数列 $\left\{  {{a}_{n} - {b}_{n}}\right\}$ 的前 $n$ 项和 ${T}_{n}$.

8.(2024$\cdot$嘉定区一模$\cdot$18)已知数列 $\left\{  {a}_{n}\right\}$ 的前 $n$ 项和为 ${S}_{n}$ ， ${S}_{n} = {n}^{2} + n$ ，其中 $n \in  \mathbf{N}$ ， $n \geq  1$.

(1)求 $\left\{  {a}_{n}\right\}$ 的通项公式；

(2)求数列 $\left\{  \dfrac{1}{{a}_{n}{a}_{n + 1}}\right\}$ 的前 $n$ 项和 ${H}_{n}$.

\subsection*{考点7 》数列与不等式、最值及范围问题}

\subsubsection*{一、填空题}

1. (2025$\cdot$闵行区二模$\cdot$10)已知数列 $\left\{  {a}_{n}\right\}$ 为等差数列，数列 $\left\{  {b}_{n}\right\}$ 为等比数列，且 ${a}_{1} = {b}_{1} = 1,{a}_{2} = {b}_{2} =  t$ ，若 ${a}_{6} + {b}_{6} > {38}$ ，则实数 $t$ 的取值范围为\blkda[2]{}.

2.(2024$\cdot$黄浦区一模$\cdot$11)设 ${a}_{1},{a}_{2},{a}_{3},\cdots ,{a}_{n}$ 是首项为 3 且公比为 $3\sqrt{3}$ 的等比数列，则满足不等式 ${\log }_{3}{a}_{1} - {\log }_{3}{a}_{2} + {\log }_{3}{a}_{3} - {\log }_{3}{a}_{4} + \cdots  + {\left( -1\right) }^{n + 1}{\log }_{3}{a}_{n} > {18}$ 的最小正整数 * $n$ 的值为\blkda[2]{}.

\subsubsection*{二、选择题}

3.(2024$\cdot$杨浦区一模$\cdot$15)等比数列 $\left\{  {a}_{n}\right\}$ 的首项 ${a}_{1} = \dfrac{1}{64}$ ，公比为 $q$ ，数列 $\left\{  {b}_{n}\right\}$ 满足 ${b}_{n} = {\log }_{0.5}{a}_{n}$ (n 是正整数)，若当且仅当 $n = 4$ 时， $\left\{  {b}_{n}\right\}$ 的前 $n$ 项和 ${B}_{n}$ 取得最大值，则 $q$ 的取值范围是\blkda[2]{}. \dotfill{(\kh{})}

\xx{}{}{}{}$\left( {3,2\sqrt{3}}\right)$ B. $\left( {3,4}\right)$

C. $\left( {2\sqrt{2},4}\right)$ D. $\left( {2\sqrt{2},3\sqrt{2}}\right)$

4.(2025$\cdot$普陀区一模$\cdot$15)设(a)>0且 $a \neq  1$ ， $k$ 、 $m$ 、 $n$ 都是正整数，数列 $\left\{  {a}_{n}\right\}$ 的通项公式为\blkda[2]{}； $\left\{  \begin{array}{l} \left( {a - 6}\right) n + {21}\left( {1 \leq  n \leq  m}\right) , \\  {a}^{n - 3}\left( {n > m}\right) , \end{array}\right.$ 记数列 $\left\{  {a}_{n}\right\}$ 中前 $k$ 项的最小值为 ${h}_{k}$ ，由所有 ${h}_{k}$ 的值所组成的集合记为 $A$. 若集合 $A$ 中仅有四个元素,则下列说法中错误的是 \dotfill{(\kh{})}

\xx{}{}{}{}当 $m = 3$ 时, $a$ 的取值范围是 $\left( {1,6}\right)$

B. 不存在 $a$ 和 $m$ 的值,使得 ${a}_{4} \notin  A$

C. 当 $m = 4$ 时， $a$ 的取值范围是(3,6)

D. 存在 $a$ 和 $m$ 的值，使得 ${a}_{5} \in  A$

5.(2025$\cdot$崇明区一模$\cdot$16)已知数列 $\left\{  {a}_{n}\right\}$ ，若存在数列 $\left\{  {b}_{n}\right\}$ 满足对任意正整数 $n$ ，都有 $\left( {{a}_{n} - {b}_{n}}\right) \left( {{a}_{n + 1} - }\right.  \left. {b}_{n + 1}\right)  < 0$ ,则称数列 $\left\{  {b}_{n}\right\}$ 是 $\left\{  {a}_{n}\right\}$ 的交错数列. 有下列两个命题:  \ding{172}  对任意给定的等差数列 $\left\{  {a}_{n}\right\}$ ,不存在等差数列 $\left\{  {b}_{n}\right\}$ ,使得 $\left\{  {b}_{n}\right\}$ 是 $\left\{  {a}_{n}\right\}$ 的交错数列;  \ding{173}  对任意给定的等比数列 $\left\{  {a}_{n}\right\}$ ,都存在等比数列 $\left\{  {b}_{n}\right\}$ ,使得 $\left\{  {b}_{n}\right\}$ 是 $\left\{  {a}_{n}\right\}$ 的交错数列. 下列结论正确的是 \dotfill{(\kh{})}

\xx{}{}{}{} \ding{172} 与 \ding{173} 都是真命题 B.  \ding{172} 为真命题， \ding{173} 为假命题

C.  \ding{172} 为假命题， \ding{173} 为真命题 D.  \ding{172} 与 \ding{173} 都是假命题

6.(2024 普陀区二模$\cdot$16)设 ${S}_{n}$ 是数列 $\left\{  {a}_{n}\right\}$ 的前 $n$ 项和 $\left( {n \geq  1,n \in  \mathbf{N}}\right)$ ，若数列 $\left\{  {a}_{n}\right\}$ 满足:对任意的 $n \geq  2$ ,存在大于 1 的整数 $m$ ,使得 $\left( {{S}_{m} - {a}_{n}}\right) \left( {{S}_{m} - {a}_{n + 1}}\right)  < 0$ 成立,则称数列 $\left\{  {a}_{n}\right\}$ 是“ $G$ 数列”. 现给出如下两个结论: \ding{172}  存在等差数列 $\left\{  {a}_{n}\right\}$ 是“ $G$ 数列”;  \ding{173}  任意等比数列 $\left\{  {a}_{n}\right\}$ 都不是“ $G$ 数列”. 下列结论正确的是 \dotfill{(\kh{})}

\xx{}{}{}{} \ding{172}  \ding{173} 都成立 B.  \ding{172} 成立， \ding{173} 不成立

C.  \ding{172} 不成立， \ding{173} 成立 D.  \ding{172}  \ding{173} 都不成立

\subsubsection*{三、解答题}

7.(2025$\cdot$普陀区一模$\cdot$21)设 $t > 1,n \geq  1,n \in  \mathbf{N}$ ，若正项数列 $\left\{  {a}_{n}\right\}$ 满足 $\dfrac{1}{t}{a}_{n} < {a}_{n + 1} < {a}_{n}$ ，则称数列 $\left\{  {a}_{n}\right\}$ 具有性质“ $P\left( t\right)$ ”.

(1)设 $m \geq  1,m \in  \mathbf{N}$ ，若数列 ${10},7,m,4,3$ 具有性质“ $P\left( 2\right)$ ”，求满足条件的 $m$ 的值；

(2)设数列 $\left\{  {a}_{n}\right\}$ 的通项公式为 ${a}_{n} = \left( {n + 1}\right) {\left( \dfrac{t}{9}\right) }^{n}$ ，问是否存在 $t$ 使得数列 $\left\{  {a}_{n}\right\}$ 具有性质“ $P\left( t\right)$ ”？ 若存在,求出满足条件的 $t$ 的取值范围; 若不存在,请说明理由;

(3)设函数 $y = f\left( x\right)$ 的表达式为 $f\left( x\right)  = \ln \left( {{{e}}^{x} - 1}\right)  - \ln x$ ，数列 $\left\{  {a}_{n}\right\}$ 的前 $n$ 项和为 ${S}_{n}$ ，且满足 ${a}_{1} = \dfrac{2}{3},{a}_{n + 1} = f\left( {a}_{n}\right)$ ,证明: 数列 $\left\{  {a}_{n}\right\}$ 具有性质“ $P\left( 3\right)$ ”,并比较 ${S}_{n}$ 与 $1 - \dfrac{1}{{3}^{n}}$ 的大小.

\subsection*{考点8 数列与其他知识点的关联问题}

\subsubsection*{一、填空题}

1.(2024$\cdot$奉贤区一模$\cdot$7)一工厂生产了某种产品 16800 件，它们来自甲、乙、丙三条生产线，为检查这批产品的质量, 决定采用分层抽样的方法进行抽样. 已知甲、乙、丙三条生产线抽取的个体数组成一个等差数列，则乙生产线生产了\blkda[2]{}件产品.

2. (2024$\cdot$闵行区一模$\cdot$8)已知 $f\left( x\right)  = {x}^{2} - {8x} + {10},x \in  \mathbf{R}$ ，数列 $\left\{  {a}_{n}\right\}$ 是公差为 1 的等差数列，若 $f\left( {a}_{1}\right)  + f\left( {a}_{2}\right)  + f\left( {a}_{3}\right)$ 的值最小，则 ${a}_{1} =$ \blkda[2]{}.

3.(2024$\cdot$宝山区一模$\cdot$11)已知函数 $f\left( x\right)  = {\left( x + 1\right) }^{3} + 1$ ，正项等比数列 $\left\{  {a}_{n}\right\}$ 满足 ${a}_{1012} = \dfrac{1}{10}$ ，则 $\mathop{\sum }\limits_{{k = 1}}^{{2023}}f\left( {\lg {a}_{k}}\right)  =$

\subsubsection*{二、选择题}

4.(2024$\cdot$嘉定区一模$\cdot$15)已知等差数列 $\left\{  {a}_{n}\right\}$ ，公差为 $d$ ， $f\left( x\right)  = \left| {x - {a}_{1}}\right|  + \left| {x - {a}_{2}}\right|$ ，则下列命题正确的是 \dotfill{(\kh{})}

\xx{}{}{}{}函数 $y = f\left( x\right) \left( {x \in  \mathbf{R}}\right)$ 可能是奇函数

B. 若函数 $y = f\left( x\right) \left( {x \in  \mathbf{R}}\right)$ 是偶函数，则 $d = 0$

C. 若 $d = 0$ ，则函数 $y = f\left( x\right) \left( {x \in  \mathbf{R}}\right)$ 是偶函数

D. 若 $d \neq  0$ ，则函数 $y = f\left( x\right) \left( {x \in  \mathbf{R}}\right)$ 的图像是轴对称图形

5.(2025$\cdot$静安区二模$\cdot$16)设函数 $y = f\left( x\right)$ 的定义域为 $\mathbf{R}$ ，若 $f\left( 0\right)  = {2025}$ ，且对任意 $x \in  \mathbf{R}$ ，满足 $f\left( {x + 1}\right)  - f\left( x\right)  \leq  {2}^{x},f\left( {x + 2}\right)  - f\left( x\right)  \geq  3 \times  {2}^{x}$ ,则 $f\left( {2025}\right)$ 的值为 \dotfill{(\kh{})}

\xx{}{}{}{}${2}^{2025} + {2024}$ B. ${2}^{2024} + {2025}$ C. ${2}^{2025} + {2025}$ D. 以上答案均不对

6.(2025$\cdot$崇明区二模$\cdot$16)数列 $\left\{  {a}_{n}\right\}$ 是等差数列，周期数列 $\left\{  {b}_{n}\right\}$ 满足 ${b}_{n} = \cos \left( {a}_{n}\right)$ ，若集合 $X =  \left\{  {x \mid  x = {b}_{n},n\text{ 是正整数 }}\right\}$ 中恰有三个元素,则数列 $\left\{  {b}_{n}\right\}$ 的周期 $T$ 的取值不可能是 \dotfill{(\kh{})}

\xx{}{}{}{}4 B. 5 C. 6 D. 7

\subsubsection*{三、解答题}

7.(2024$\cdot$长宁区一模$\cdot$17)已知等差数列 $\left\{  {a}_{n}\right\}$ 的前 $n$ 项和为 ${S}_{n}$ ，公差 $d = 2$.

(1)若 ${S}_{10} = {100}$ ，求 $\left\{  {a}_{n}\right\}$ 的通项公式；

(2)从集合 $\left\{  {{a}_{1},{a}_{2},{a}_{3},{a}_{4},{a}_{5},{a}_{6}}\right\}$ 中任取 3 个元素，记这 3 个元素能成等差数列为事件 $A$ ，求事件 $A$ 发生的概率 $P\left( A\right)$.

8.(2024$\cdot$松江区一模$\cdot$18)已知数列 $\left\{  {a}_{n}\right\}$ 为等差数列， $\left\{  {b}_{n}\right\}$ 是公比为 2 的等比数列，且 ${a}_{2} - {b}_{2} = {a}_{3} -  {b}_{3} = {b}_{4} - {a}_{4}$

(1)证明: ${a}_{1} = {b}_{1}$ ；

(2)若集合 $M = \left\{  {k \mid  {b}_{k} = {a}_{m} + {a}_{1},1 \leq  m \leq  {50}}\right\}$ ，求集合 $M$ 中的元素个数.

9.(2024$\cdot$崇明区一模$\cdot$21)已知 $f\left( x\right)  = {mx} + \sin x$ ( $m \in  \mathbf{R}$ 且 $m \neq  0$ ).

(1)若函数 $y = f\left( x\right)$ 是实数集 $\mathbf{R}$ 上的严格增函数，求实数 $m$ 的取值范围；

(2)已知数列 $\left\{  {a}_{n}\right\}$ 是等差数列(公差 $d \neq  0$ )， ${b}_{n} = f\left( {a}_{n}\right)$. 是否存在数列 $\left\{  {a}_{n}\right\}$ 使得数列 $\left\{  {b}_{n}\right\}$ 是等差数列? 若存在,请写出一个满足条件的数列 $\left\{  {a}_{n}\right\}$ ,并证明此时的数列 $\left\{  {b}_{n}\right\}$ 是等差数列; 若不存在,请说明理由;

(3) 若 $m = 1$ ,是否存在直线 $y = {kx} + b$ 满足:  \ding{172}  对任意的 $x \in  \mathbf{R}$ 都有 $f\left( x\right)  \geq  {kx} + b$ 成立;  \ding{173}  存在 ${x}_{0} \in  \mathbf{R}$ 使得 $f\left( {x}_{0}\right)  = k{x}_{0} + b$ ? 若存在,请求出满足条件的直线方程; 若不存在,请说明理由.

10.(2024$\cdot$杨浦区一模$\cdot$21)设函数 $f\left( x\right)  = x + A\sin \dfrac{\pi x}{2},x \in  \mathbf{R}$ (其中常数 $A \in  \mathbf{R},A > 0$ )，无穷数列 $\left\{  {a}_{n}\right\}$ 满足: 首项 ${a}_{1} > 0,{a}_{n + 1} = f\left( {a}_{n}\right)$.

(1)判断函数 $y = f\left( x\right)$ 的奇偶性,并说明理由;

(2)若数列 $\left\{  {a}_{n}\right\}$ 是严格增数列，求证:当 $A < 4$ 时，数列 $\left\{  {a}_{n}\right\}$ 不是等差数列；

(3)当 $A = 8$ 时，数列 $\left\{  {a}_{n}\right\}$ 是否可能为公比小于 0 的等比数列? 若可能，求出所有公比的值；若不可能, 请说明理由.

11.(2023$\cdot$上海秋季卷$\cdot$21)已知函数 $f\left( x\right)  = \ln x$ ，过函数上的点 $\left( {{a}_{1},f\left( {a}_{1}\right) }\right)$ 作切线交 $y$ 轴于(0， $\left. {a}_{2}\right) ,{a}_{2} > 0$ ,过函数上的点 $\left( {{a}_{2},f\left( {a}_{2}\right) }\right)$ 作切线交 $y$ 轴于 $\left( {0,{a}_{3}}\right)$ ,以此类推,直至 ${a}_{m} \leq  0$ 时则停止操作,得到数列 $\left\{  {a}_{n}\right\}  ,{m}_{n} \in  \mathbf{N},1 < n \leq  m$.

(1)证明: ${a}_{n + 1} = \ln {a}_{n} - 1$ ；

(2)试比较 ${a}_{n + 1}$ 和 ${a}_{n} - 2$ 的大小；

(3)若正整数 $k \geq  3$ ，是否存在 $k$ 使得 ${a}_{1},{a}_{2},{a}_{3},\cdots ,{a}_{k}$ 依次成等差数列? 若存在，求出 $k$ 的所有取值; 若不存在, 请说明理由.

12.(2024$\cdot$浦东新区二模$\cdot$21)已知函数 $y = f\left( x\right)$ 及其导函数 $y = {f}^{\prime }\left( x\right)$ 的定义域均为 $D$ ，设 ${x}_{0} \in  D$ ,曲线 $y = f\left( x\right)$ 在点 $\left( {{x}_{0},f\left( {x}_{0}\right) }\right)$ 处的切线交 $x$ 轴于点 $\left( {{x}_{1},0}\right)$. 当 $n \geq  1$ 时,设曲线 $y = f\left( x\right)$ 在点 $\left( {{x}_{n},f\left( {x}_{n}\right) }\right)$ 处的切线交 $x$ 轴于点 $\left( {{x}_{n + 1},0}\right)$. 依此类推,称得到的数列 $\left\{  {x}_{n}\right\}$ 为函数 $y = f\left( x\right)$ 关于 ${x}_{0}$ 的“ $N$ 数列”.

(1)若 $f\left( x\right)  = \ln x,\left\{  {x}_{n}\right\}$ 是函数 $y = f\left( x\right)$ 关于 ${x}_{0} = \dfrac{1}{{e}}$ 的“ $N$ 数列”,求 ${x}_{1}$ 的值;

(2)若 $f\left( x\right)  = {x}^{2} - 4,\left\{  {x}_{n}\right\}$ 是函数 $y = f\left( x\right)$ 关于 ${x}_{0} = 3$ 的“ $N$ 数列”，记 ${a}_{n} = {\log }_{3}\dfrac{{x}_{n} + 2}{{x}_{n} - 2}$ ，证明: $\left\{  {a}_{n}\right\}$ 是等比数列,并求出其公比;

(3)若 $f\left( x\right)  = \dfrac{x}{a + {x}^{2}}$ ，则对任意给定的非零实数 $a$ ，是否存在 ${x}_{0} \neq  0$ ，使得函数 $y = f\left( x\right)$ 关于 ${x}_{0}$ 的“$N$ 数列” $\left\{  {x}_{n}\right\}$ 为周期数列? 若存在，求出所有满足条件的 ${x}_{0}$ ；若不存在，请说明理由.

13.(2024$\cdot$青浦区二模$\cdot$21)若无穷数列 $\left\{  {a}_{n}\right\}$ 满足:存在正整数 $T$ ，使得 ${a}_{n + T} = {a}_{n}$ 对一切正整数 $n$ 成立,则称 $\left\{  {a}_{n}\right\}$ 是周期为 $T$ 的周期数列.

(1)若 ${a}_{n} = \sin \left( {\dfrac{\pi n}{m} + \dfrac{\pi }{3}}\right)$ (其中正整数 $m$ 为常数， $n \in  \mathbf{N},n \geq  1$ )，判断数列 $\left\{  {a}_{n}\right\}$ 是否为周期数列，并说明理由;

(2)若 ${a}_{n + 1} = {a}_{n} + \sin {a}_{n}\left( {n \in  \mathbf{N},n \geq  1}\right)$ ，判断数列 $\left\{  {a}_{n}\right\}$ 是否为周期数列，并说明理由；

(3)设 $\left\{  {b}_{n}\right\}$ 是无穷数列，已知 ${a}_{n + 1} = {b}_{n} + \sin {a}_{n}\left( {n \in  \mathbf{N},n \geq  1}\right)$. 求证:“存在 ${a}_{1}$ ，使得 $\left\{  {a}_{n}\right\}$ 是周期数列”的充要条件是“$\left\{  {b}_{n}\right\}$ 是周期数列”.

14.(2024$\cdot$青浦区一模$\cdot$21)已知有穷等差数列 $\left\{  {a}_{n}\right\}$ : ${a}_{1},{a}_{2},\cdots ,{a}_{m}\left( {m \geq  3,m \in  \mathbf{N}}\right)$ 的公差 $d$ 大于零.

(1)证明: $\left\{  {a}_{n}\right\}$ 不是等比数列；

(2)是否存在指数函数 $y = f\left( x\right)$ 满足: $y = f\left( x\right)$ 在 $x = {a}_{1}$ 处的切线交 $x$ 轴于 $\left( {{a}_{2},0}\right) ,y = f\left( x\right)$ 在 $x = {a}_{2}$ 处的切线交 $x$ 轴于 $\left( {{a}_{3},0}\right) ,\cdots ,y = f\left( x\right)$ 在 $x = {a}_{m - 1}$ 处的切线交 $x$ 轴于 $\left( {{a}_{m},0}\right)$ ? 若存在,请写出函数 $y = f\left( x\right)$ 的表达式,并说明理由; 若不存在,也请说明理由;

(3)若数列 $\left\{  {a}_{n}\right\}$ 中所有项按照某种顺序排列后可以构成等比数列 $\left\{  {b}_{n}\right\}$ ，求出所有可能的 $m$ 的取值.

15.(2024$\cdot$徐汇区二模$\cdot$21)已知各项均不为 0 的数列 $\left\{  {a}_{n}\right\}$ 满足 ${a}_{n + 2}{a}_{n} = {a}_{n + 1}{a}_{n} + {a}_{n + 1}^{2}$ (n后置). 数), ${a}_{1} = {a}_{2} = 1$ ,定义函数 $y = {f}_{n}\left( x\right)  = 1 + \mathop{\sum }\limits_{{k = 1}}^{n}\dfrac{1}{k!}{x}^{k}\left( {x \geq  0}\right)$ , e 是自然对数的底数.

(1)求证:数列 $\left\{  \dfrac{{a}_{n + 1}}{{a}_{n}}\right\}$ 是等差数列，并求数列 $\left\{  {a}_{n}\right\}$ 的通项公式；

(2)记函数 $y = {g}_{n}\left( x\right)$ ，其中 ${g}_{n}\left( x\right)  = 1 - {{e}}^{-x}{f}_{n}\left( x\right)$.

 \ding{172}  证明:对任意 $x \geq  0,0 \leq  {g}_{3}\left( x\right)  \leq  {f}_{4}\left( x\right)  - {f}_{3}\left( x\right)$ ;

 \ding{173}  数列 $\left\{  {b}_{n}\right\}$ 满足 ${b}_{n} = \dfrac{{2}^{n - 1}}{{a}_{n}}$ ，设 ${T}_{n}$ 为数列 $\left\{  {b}_{n}\right\}$ 的前 $n$ 项和. 数列 $\left\{  {T}_{n}\right\}$ 的极限的严格定义为: 若存在一个常数 $T$ ，使得对任意给定的正实数 $u$ (不论它多么小)，总存在正整数 $m$ 满足: 当 $n \geq  m$ 时，恒有 $\left| {{T}_{n} - T}\right|  < u$ 成立，则称 $T$ 为数列 $\left\{  {T}_{n}\right\}$ 的极限. 试根据以上定义求出数列 $\left\{  {T}_{n}\right\}$ 的极限 $T$.

\section*{第 7 章 平面解析几何}

\begin{center}

\begin{tabular}{|p{4cm}|p{12cm}|}
\hline
考点 & 课标要求与命题趋势 \\
\cline{1-2}
1. 直线方程与圆的方程 & \ding{172}  在平面直角坐标系中，结合具体图形，探索确定直线位置的几何要素 \ding{173}  理解直线的倾斜角和斜率的概念，经历用代数方法刻画直线斜率的过程，掌握过两点的直线斜率的计算公式; 能根据斜率判定两条直线平行或垂直 \ding{174}  根据确定直线位置的几何要素，探索并掌握直线方程的几种形式(点斜式、两点式及一般式);能用解方程组的方法求两条直线的交点坐标 \ding{175}  探索并掌握平面上两点间的距离公式、点到直线的距离公式,会求两条平行直线间的距离. 熟练掌握距离计算及其参数求解, 该内容是上海高考常考内容, 通常和圆结合在一起考查,需重点练习 \\
\cline{1-1}
2. 直线与圆的位置关系及其应用 &  \\
\cline{1-2}
3. 椭圆方程及其性质 & \ding{172}  了解圆锥曲线的实际背景，感受圆锥曲线在刻画现实世界和解决实际问题中的作用. 了解椭圆、抛物线的简单应用 \ding{173}  经历从具体情境中抽象出椭圆的过程，掌握椭圆的定义、标准方程及简单几何性质，关注离心率的求解及应用，是上海高考高频考点 \ding{174}  通过圆锥曲线与方程的学习，进一步体会数形结合的思想 \\
\cline{1-2}
4. 双曲线方程及其性质 & \ding{172}  了解双曲线的定义、几何图形和标准方程，以及它们的简单几何性质 \ding{173}  熟练掌握双曲线的方程及其性质应用，关注离心率的求解及应用，这是上海高考高频考点 \\
\cline{1-2}
5. 抛物线方程及其性质 & 熟练掌握抛物线的方程及其性质应用, 这是上海高考高频考点 \\
\cline{1-2}
6. 直线与圆锥曲线的位置关系及其应用 & \ding{172}  熟练掌握直线与圆锥曲线的位置关系，并会求解最值及范围，该内容也是上海高考命题执点:掌握求曲线方程及轨迹方程的基本方法和数学思想 \ding{173} 收集、阅读平面解析几何的形成与发展的历史资料，撰写小论文、论述平面解析几何发展的过程、重要结果、主要人物、关键事件及其对人类文明的贡献 \\
\cline{1-2}
\hline
\end{tabular}
\end{center}
\clearpage
\subsection*{考点1 直线方程与圆的方程}

\subsubsection*{一、填空题}

1.(2024$\cdot$上海春季卷$\cdot$2)直线 $x - y + 1 = 0$ 的倾斜角为\blkda[2]{}.

2.(2025$\cdot$嘉定区一模$\cdot$2)直线: ${3x} - y + 1 = 0$ 的倾斜角为\blkda[2]{}.(用反三角函数表示)

3. (2025$\cdot$奉贤区一模$\cdot$2)若直线 ${l}_{1} : x + {ay} - 2 = 0$ 与直线 ${l}_{2} : {ax} + y - 2 = 0$ 互相垂直，则 $a =$ \blkda[2]{}.

4.(2025$\cdot$闵行区一模$\cdot$3)直线、 $\sqrt{3}x - y + 1 = 0$ 的倾斜角为\blkda[2]{}.

5.(2024$\cdot$徐汇区一模$\cdot$3)已知直线 $l : y = {kx} + 2$ 经过点(1，1)，则直线 $l$ 倾斜角的大小为\blkda[2]{}.

6.(2025$\cdot$浦东新区二模$\cdot$3)设圆 $C$ 方程为 ${x}^{2} + {y}^{2} + {4x} - {6y} + {10} = 0$ ，则圆 $C$ 的半径为\blkda[2]{}.

7.(2023$\cdot$上海春季卷$\cdot$4)已知圆 $C$ 的一般方程为 ${x}^{2} + {2x} + {y}^{2} = 0$ ，则圆 $C$ 的半径为\blkda[2]{}.

\subsubsection*{二、选择题}

8.(2025$\cdot$崇明区二模$\cdot$4)直线 $x =  - 2$ 与直线 $\sqrt{3}x - y + 1 = 0$ 的夹角为\blkda[2]{}.

9.(2025$\cdot$长宁区一模$\cdot$4)以(3,4)为圆心， \ding{51} 另为半径的圆的标准方程是\blkda[2]{}.

10.(2025$\cdot$青浦区一模$\cdot$4)已知直线 ${l}_{1} : x + \left( {1 + m}\right) y + m - 2 = 0$ 与直线 ${l}_{2} : {mx} + {2y} + 8 = 0$ 平行， 则 $m =$ \blkda[2]{}.

11.(2025$\cdot$虹口区二模$\cdot$5)若直线 $l$ 与直线 $y = x$ 平行，且经过圆 ${x}^{2} - {2x} + {y}^{2} = 0$ 的圆心，则 $l$ 的方程为\blkda[2]{}.

12.(2025$\cdot$奉贤区二模$\cdot$5)直线 3x+4y-5=0 上的动点 P 和直线 3x+4y+10=0 上的动点 Q，则点 $P$ 与点 $Q$ 之间距离的最小值是\blkda[2]{}.

13.(2025$\cdot$金山区一模$\cdot$6)以((3，4)为圆心且过点(1，-3)的圆的标准方程是\blkda[2]{}.

14.(2024$\cdot$长宁区二模$\cdot$6)直线: ${2x} - y - 3 = 0$ 与直线 $x - {3y} - 5 = 0$ 的夹角大小为\blkda[2]{}.

15.(2023$\cdot$上海秋季卷$\cdot$7)已知圆 ${x}^{2} + {y}^{2} - {4y} - m = 0$ 的面积为 $\pi$ ，则 $m =$ \blkda[2]{}.

16.(2024$\cdot$青浦区二模$\cdot$7)已知直线 ${l}_{1}$ 的倾斜角比直线 ${l}_{2} : y = x\tan {80}^{ \circ  }$ 的倾斜角小 ${20}^{ \circ  }$ ，则 ${l}_{1}$ 的斜率为\blkda[2]{}.

17.(2024$\cdot$奉贤区一模$\cdot$11)已知直线 ${l}_{1} : y - 2 = 0$ 和直线 ${l}_{2} : x + 1 = 0$ ，则曲线 ${\left( x - 1\right) }^{2} + {y}^{2} = 1$ 上一动点 $P$ 到直线 ${l}_{1}$ 和直线 ${l}_{2}$ 的距离之和的最小值是\blkda[2]{}.

18.(2024$\cdot$宝山区一模$\cdot$12)设点 $P$ 在直线 $l : {2x} - y - 5 = 0$ 上，点 $Q$ 在曲线 $\Gamma  : y = x + \ln x$ 上，线段 ${PQ}$ 的中点为 $M,O$ 为坐标原点，则 $\left| \vv{OM}\right|$ 的最小值为\blkda[2]{}.

\subsubsection*{二、选择题}

19.(2024$\cdot$嘉定区一模$\cdot$13)直线倾斜角的取值范围为 \dotfill{(\kh{})}

\xx{}{}{}{}$\left\lbrack  {0,\dfrac{\pi }{2}}\right)$ B. $\left\lbrack  {0,\dfrac{\pi }{2}}\right\rbrack$ C. $\lbrack 0,\pi )$ D. $\left\lbrack  {0,\pi }\right\rbrack$

20.(2025$\cdot$闵行区二模$\cdot$16)已知点 $A\left( {{x}_{1},{y}_{1}}\right)$ 在圆 ${x}^{2} + {y}^{2} = 9$ 上，点 $B\left( {{x}_{2},{y}_{2}}\right)$ 在圆 ${x}^{2} + {y}^{2} = {12}$ 上,且 ${x}_{1}{x}_{2} + {y}_{1}{y}_{2} = {x}_{1} + {x}_{2} - 1,O$ 为坐标原点. 对于以下两个命题,判断正确的是

 \ding{172}  在坐标平面内存在点 $P$ ，使得 ${AP}\bot {BP}$ 恒成立； \ding{173}  $\bigtriangleup  {OAB}$ 面积的最小值为 $\sqrt{22}$.

\xx{}{}{}{} \ding{172}  \ding{173} 都是真命题 B.  \ding{172} 是假命题， \ding{173} 是真命题

C.  \ding{172} 是真命题， \ding{173} 是假命题 D.  \ding{172}  \ding{173} 都是假命题

\subsection*{考点2}

\subsubsection*{一、填空题}

1.(2024$\cdot$徐汇区二模$\cdot$6)若圆 ${x}^{2} + {y}^{2} = 4$ 与圆 ${x}^{2} + {y}^{2} - {2ax} + {a}^{2} - 1 = 0$ 内切，则 $a$ 等于\blkda[2]{}.

2. (2025$\cdot$嘉定区二模$\cdot$7)直线 $l : y = x + 1$ 与圆 $C : {x}^{2} + {y}^{2} - {4x} - {2y} = 0$ 相交所得的弦长为\blkda[2]{}.

3.(2024$\cdot$浦东新区二模$\cdot$8)已知圆 ${C}_{1} : {x}^{2} + {y}^{2} - {2ax} + {a}^{2} - 1 = 0\left( {a > 0}\right)$ ，圆 ${C}_{2} : {x}^{2} + {y}^{2} - {4y} - 5 =$ 0 ，若两圆相交，则实数 $a$ 的取值范围为\blkda[2]{}.

4.(2025$\cdot$黄浦区二模$\cdot$8)已知 $a$ 为常数，圆 ${\left( x - a\right) }^{2} + {\left( y + a - 2\right) }^{2} = {r}^{2}\left( {r > 0}\right)$ 与圆 ${x}^{2} + {y}^{2} = 1$ 有公共点,当 $r$ 取到最小值时, $a$ 的值为\blkda[2]{}.

5.(2025$\cdot$浦东新区一模$\cdot$10)某地要建造一个市民休闲公园长方形 ABCD， 如图所示,边 ${AB} = 2{{km}}$ ,边 ${AD} = 1{{km}}$ ,其中区域 ${ADE}$ 开挖成一个人工湖, 其他区域为绿化风景区. 经测算, 人工湖在公园内的边界是一段圆弧, 且 $A\text{ 、 }D$ 位于圆心 $O$ 的正北方向, $E$ 位于圆心 $O$ 的北偏东 ${60}^{ \circ  }$ 方向. 拟定在圆弧 $P$ 处修建一座渔人码头,供游客湖中泛舟,并在公园的边 ${DC}\text{ 、 }{CB}$ 开设两扇门 $M$ 、 $N$ ，修建步行道 ${PM}$ 、 ${PN}$ 通往渔人码头，且 ${PM}\bot {CD}$ 、 ${PN} \bot  {CB}$ ，则步行道 ${PM}$ 、 ${PN}$ 长度之和的最小值是\blkda[2]{} ${{km}}$. (精确到 0.001)

\begin{center}
\includegraphics[width=0.2\textwidth]{bodyxb/src/92_1109_245_289_285_0.jpg}
\end{center}


第 5 题图

6.(2024$\cdot$上海春季卷$\cdot$11)已知正方形展区 ${ABCD}$ 边长为 ${1.2}{{km}},E$ 距 ${AB}\text{ 、 }{AD}$ 的距离都为 ${0.2}{{km}},F$ 距 ${BC}\text{ 、 }{CD}$ 的距离都为 ${0.4}{{km}}$ ,若有一个圆形跑道经过 $E\text{ 、 }F$ 两点,且与 ${AD}$ 只有一个交点,则圆形跑道的周长为\blkda[2]{}.

\begin{center}
\includegraphics[width=0.2\textwidth]{bodyxb/src/92_1110_587_288_292_0.jpg}
\end{center}


第 6 题图

7.(2024$\cdot$崇明区二模$\cdot$12)已知实数 ${x}_{1}$ 、 ${x}_{2}$ 、 ${y}_{1}$ 、 ${y}_{2}$ 满足: ${x}_{1}^{2} + {y}_{1}^{2} = 1,{x}_{2}^{2} +  {y}_{2}^{2} = 1,{x}_{1}{y}_{2} - {y}_{1}{x}_{2} = 1$ ,则 $\left| {{x}_{1} + {y}_{1} - 2}\right|  + \left| {{x}_{2} + {y}_{2} - 2}\right|$ 的最大值是\blkda[2]{}.

8.(2024$\cdot$静安区二模$\cdot$12)我们称如图的曲线为“爱心线”，其上的任意一点 $P  \left( {x,y}\right)$ 都满足方程 ${x}^{2} - 2\left| x\right| y + {y}^{2} - \sqrt{2}\left| x\right|  + 2\sqrt{2}y = 0$ ,现将一边在 $x$ 轴上,另外两个顶点在爱心线上的矩形称为 “心吧”. 若已知点 $M\left( {\dfrac{\sqrt{2}}{2}, - \sqrt{2}}\right)$ 到 “爱心线” 上任意一点的最小距离为 $d$ ，则用 $d$ 表示 “心吧” 面积的最大值为\blkda[2]{}.

\begin{center}
\includegraphics[width=0.2\textwidth]{bodyxb/src/92_1139_956_287_270_0.jpg}
\end{center}


第 8 题图

9.(2025$\cdot$嘉定区一模$\cdot$12)已知实数 ${x}_{1}$ 、 ${x}_{2}$ 、 ${y}_{1}$ 、 ${y}_{2}$ 满足: ${x}_{1}^{2} + {y}_{1}^{2} = 1$ ， ${x}_{2}^{2} + {y}_{2}^{2} =$ 1， ${x}_{1}{x}_{2} + {y}_{1}{y}_{2} = \dfrac{1}{2}$ ，则 $\dfrac{\left| {x}_{1} + {y}_{1} - 1\right| }{\sqrt{2}} + \dfrac{\left| {x}_{2} + {y}_{2} - 1\right| }{\sqrt{2}}$ 的最小值为\blkda[2]{}.

\subsubsection*{二、选择题}

10.(2024$\cdot$普陀区二模$\cdot$15)直线 $l$ 经过定点 $P\left( {2,1}\right)$ ，且与 $x$ 轴正半轴、 $y$ 轴正半轴分别相交于 $A$ 、 $B$ 两点, $O$ 为坐标原点,动圆 $M$ 在 $\bigtriangleup {OAB}$ 的外部,且与直线 $l$ 及两坐标轴的正半轴均相切,则 $\bigtriangleup {OAB}$ 周长的最小值是 \dotfill{(\kh{})}

\xx{}{}{}{}3 B. 5 C. 10 D. 12

11.(2024$\cdot$黄浦区一模$\cdot$16)在平面直角坐标系 ${xOy}$ 中，对于定点 $P\left( {a,b}\right)$ ，记点集 $\{ \left( {x,y}\right)  \mid  \left| {x - a}\right|  \leq  1,\left| {y - b}\right|  \leq  1\}$ 中距离原点 $O$ 最近的点为点 ${Q}_{P}$ ，此最近距离为 $f\left( P\right)$. 当点 $P$ 在曲线 ${x}^{2} + {y}^{2} -  {8x} - {4y} + {16} = 0$ 上运动时，关于下列结论: \ding{172} \xx{}{}{}{}${Q}_{P}$ 的轨迹是一个圆； \ding{173}  $f\left( P\right)$ 的取值范围是 $\left\lbrack  {\sqrt{10} - 2,\sqrt{10} + 2}\right\rbrack$.

下列判断正确的是 \dotfill{(\kh{})}

\xx{}{}{}{} \ding{172}  \ding{173} 都成立 B.  \ding{172} 成立， \ding{173} 不成立 C.  \ding{172} 不成立， \ding{173} 成立 D.  \ding{172}  \ding{173} 都不成立

\subsection*{考点3 椭圆方程及其性质}

\subsubsection*{一、填空题}

1. (2025$\cdot$黄浦区一模$\cdot$3)椭圆 $\dfrac{{x}^{2}}{4} + \dfrac{{y}^{2}}{3} = 1$ 的焦距为\blkda[2]{}.

2.(2025$\cdot$静安区二模$\cdot$3)椭圆 $\dfrac{{x}^{2}}{4} + \dfrac{{y}^{2}}{3} = 1$ 的离心率是\blkda[2]{}.

3.(2025$\cdot$静安区一模$\cdot$5)到点 ${F}_{1}\left( {-3,0}\right)$ 、 ${F}_{2}\left( {3,0}\right)$ 距离之和为 10 的动点 $P$ 的轨迹方程为\blkda[2]{}.

4.(2024$\cdot$青浦区二模$\cdot$6)椭圆 $\dfrac{{x}^{2}}{{a}^{2}} + {y}^{2} = 1\left( {a > 1}\right)$ 的离心率为 $\dfrac{\sqrt{3}}{2}$ ，则 $a =$ \blkda[2]{}.

5.(2024$\cdot$杨浦区一模$\cdot$6)若椭圆 $\dfrac{{x}^{2}}{{a}^{2}} + {y}^{2} = 1\left( {a > 1}\right)$ 长轴长为 4，则其离心率为\blkda[2]{}.

6.(2025$\cdot$普陀区一模$\cdot$7)设椭圆 $C : \dfrac{{x}^{2}}{{a}^{2}} + \dfrac{{y}^{2}}{{b}^{2}} = 1\left( {a > b > 0}\right)$ 的左、右焦点分别为 ${F}_{1}$ 、 ${F}_{2}$ ，左顶点为 $A$ ， 若椭圆 $C$ 的离心率为 $\dfrac{1}{3}$ ，则 $\dfrac{\left| {F}_{2}A\right| }{\left| A{F}_{1}\right| }$ 的值为\blkda[2]{}.

7.(2025$\cdot$松江区一模$\cdot$10)已知点 $P$ 为椭圆 $\dfrac{{x}^{2}}{4} + \dfrac{{y}^{2}}{3} = 1$ 上任意一点， ${EF}$ 为圆 $N$ :( $x - 1{\text{ ) }}^{2} + {y}^{2} = 4$ 的任意一条直径，则 $\vv{PE} \cdot  \vv{PF}$ 的取值范围是\blkda[2]{}.

8.(2025$\cdot$徐汇区一模$\cdot$10)已知椭圆 $\dfrac{{x}^{2}}{{a}^{2}} + \dfrac{{y}^{2}}{{b}^{2}} = 1\left( {a > b > 0}\right)$ 的左、右焦点分别为 ${F}_{1}$ 、 ${F}_{2}$ ， $P$ 为椭圆上一点，且 $\angle {P{F}_{2}{F}_{1}} = \dfrac{\pi }{3}$ ，若此椭圆的离心率为 $\sqrt{3} - 1$ ，则 $\angle {P{F}_{1}{F}_{2}}$ 的大小为\blkda[2]{}.

9.(2025$\cdot$杨浦区一模$\cdot$11)中国探月工程又称“嫦娥工程”，是中国航天活动的第三个里程碑. 在探月过程中，月球探测器需要进行变轨，即从一条椭圆轨道变到另一条不同的椭圆轨道上. 若变轨前后的两条椭圆轨道均以月球中心为一个焦点, 变轨后椭圆轨道上的点与月球中心的距离最小值保持不变，而距离最大值扩大为变轨前的 4 倍，椭圆轨道的离心率扩大为变轨前的 2.5 倍，则变轨前的椭圆轨道的离心率为\blkda[2]{}. (精确到 0.01)

\subsubsection*{二、选择题}

\begin{center}
\includegraphics[width=0.2\textwidth]{bodyxb/src/93_1180_1772_286_245_0.jpg}
\end{center}


第 10 题图

10.(2024$\cdot$普陀区一模$\cdot$13)若椭圆 $\Gamma$ 的两个顶点和焦点都在圆 $O : {x}^{2} + {y}^{2} = 4$ 上，如图所示，则下列结论正确的是 \dotfill{(\kh{})}

\xx{}{}{}{}椭圆 $\Gamma$ 的方程是 $\dfrac{{x}^{2}}{4} + \dfrac{{y}^{2}}{2} = 1$

B. 过椭圆 $\Gamma$ 上的点作圆 $O$ 的切线,一定有两条

C. 圆 $O$ 上的点与椭圆 $\Gamma$ 上的点的距离的最大值是 $2\left( {\sqrt{2} + 1}\right)$

D. 直线 $x + y + 2\sqrt{2} = 0$ 与椭圆 $\Gamma$ 有交点,与圆 $O$ 无交点

11. (2024$\cdot$虹口区一模$\cdot$16)已知曲线 $\Gamma$ 的对称中心为 $O$ ，若对于 $\Gamma$ 上的任意一点 $A$ ，都存在 $\Gamma$ 上两点 $B$ 、 $C$ ，使得 $O$ 为 $\bigtriangleup {ABC}$ 的重心，则称曲线 $\Gamma$ 为“自稳定曲线”. 现有如下两个命题:

 \ding{172}  任意椭圆都是“自稳定曲线”； \ding{173} 存在双曲线是“自稳定曲线”. 则 \dotfill{(\kh{})}

\xx{}{}{}{} \ding{172} 是假命题， \ding{173} 是真命题 B.  \ding{172} 是真命题， \ding{173} 是假命题

C.  \ding{172}  \ding{173} 都是假命题 D.  \ding{172}  \ding{173} 都是真命题

\subsubsection*{三、解答题}

12.(2025$\cdot$青浦区二模$\cdot$19)如图所示，椭圆 ${C}_{1} : \dfrac{{x}^{2}}{8} + \dfrac{{y}^{2}}{{b}^{2}} = 1\left( {0 < b < 2\sqrt{2}}\right)$ 与双曲线 ${C}_{2} : \dfrac{{x}^{2}}{{b}^{2}} - {y}^{2} = 1$ 在第一象限的公共点为 $A\left( {{x}_{A},{y}_{A}}\right) \left( {{x}_{A} > 0}\right)$. 曲线 $\Gamma$ 由两段曲线组成: 当 $x \leq  {x}_{A}$ 时,曲线 $\Gamma$ 与椭圆 ${C}_{1}$ 重合,当 $x > {x}_{A}$ 时,曲线 $\Gamma$ 与双曲线 ${C}_{2}$ 重合.

(1)当 ${x}_{A} = 2$ 时,求 $b$ 的值;

(2)已知 $b = \sqrt{2}$ ，直线 $l$ 过点 $D\left( {2,0}\right)$ 与曲线 $\Gamma$ 交于 $E$ 、 $F$ 两点，若 $\vv{AD} \cdot  \vv{EF} = 2$ ，求直线 $l$ 的方程；

(3)已知 $A\left( {2,1}\right)$ ，斜率为 $k\left( {k \geq  1}\right)$ 的直线 $m$ 过点 $P\left( {0,1}\right)$ 与曲线 $\Gamma$ 交于 $M$ ， $N$ 两点，若 ${S}_{\bigtriangleup {AMN}} =  \lambda \tan \angle {MAN}$ ,求实数 $\lambda$ 的最大值.

\begin{center}
\includegraphics[width=0.5\textwidth]{bodyxb/src/94_778_881_692_249_0.jpg}
\end{center}


第 12 题图

13.(2025$\cdot$上海春季卷$\cdot$20)在平面直角坐标系中，已知曲线 $\Gamma  : \dfrac{{x}^{2}}{4} + {y}^{2} = 1\left( {y \geq  0}\right)$ ，点 $P$ 、 $Q$ 分别为 $\Gamma$ 上不同的两点,设点 $T\left( {t,0}\right)$.

(1)求 $\Gamma$ 所在椭圆的离心率；

(2)若 $t = 1$ ，点 $Q$ 在 $y$ 轴上，点 $T$ 到直线 ${PQ}$ 的距离为 $\dfrac{\sqrt{5}}{5}$ ，求点 $P$ 的坐标；

(3)是否存在 $t$ ，使得 $\bigtriangleup  {TPQ}$ 是以点 $T$ 为直角顶点的等腰直角三角形？若存在，求 $t$ 的取值范围；若不存在，请说明理由.

\begin{center}
\includegraphics[width=0.2\textwidth]{bodyxb/src/94_1138_1672_325_208_0.jpg}
\end{center}


第 13 题图

14.(2025$\cdot$嘉定区二模$\cdot$20)已知椭圆 $C : \dfrac{{x}^{2}}{9} + {y}^{2} = 1$ ， $F$ 为椭圆的右焦点，过椭圆上一点 $P\left( {3,0}\right)$ 的直线 ${l}_{1}$ 交椭圆于另一点 $Q$ ，点 $M$ 为椭圆上任意一点.

(1)求 $\left| {MF}\right|$ 的最小值；

(2)当直线 ${l}_{1}$ 的斜率为 1 时,求 $\bigtriangleup  {PQM}$ 面积的最大值及此时点 $M$ 的坐标；

(3)若直线 ${PQ}$ 与直线 ${l}_{2} : x =  - 3$ 交于点 $D$ ，点 $D$ 不在 $x$ 轴上， $Q$ 关于原点的对称点为点 $R$ ，直线 ${PR}$ 与 ${l}_{2}$ 交于点 $E$ ，求线段 $\left| {DE}\right|$ 的取值范围.

\begin{center}
\includegraphics[width=0.2\textwidth]{bodyxb/src/95_1042_444_342_192_0.jpg}
\end{center}


第 14 题图

15.(2024$\cdot$上海春季卷$\cdot$20)在平面直角坐标系 ${xOy}$ 中，已知点 $A$ 为椭圆 $\Gamma  : \dfrac{{x}^{2}}{6} + \dfrac{{y}^{2}}{2} = 1$ 上一点， ${F}_{1}$ 、 ${F}_{2}$ 分别为椭圆的左、右焦点.

(1)若点 $A$ 的横坐标为 2，求 $\left| {A{F}_{1}}\right|$ 的长；

(2)设 $\Gamma$ 的上、下顶点分别为 ${M}_{1}$ 、 ${M}_{2}$ ，记 $\bigtriangleup  A{F}_{1}{F}_{2}$ 的面积为 ${S}_{1}$ ， $\bigtriangleup  A{M}_{1}{M}_{2}$ 的面积为 ${S}_{2}$ ，若 ${S}_{1} \geq  {S}_{2}$ ，求 $\left| {OA}\right|$ 的取值范围；

(3)在 $x$ 轴上方，设直线 $A{F}_{2}$ 与 $\Gamma$ 交于点 $B$ ，与 $y$ 轴交于点 $K$ ， $K{F}_{1}$ 延长线与 $\Gamma$ 交于点 $C$ ,在 $x$ 轴上方是否存在点 $C$ ,使得 $\vv{{F}_{1}A} +  \vv{{F}_{1}B} + \vv{{F}_{1}C} = \lambda \left( {\vv{{F}_{2}A} + \vv{{F}_{2}B} + \vv{{F}_{2}C}}\right) \left( {\lambda  \in  \mathbf{R}}\right)$ 成立？若存在，请求出点 $C$ 的坐标; 若不存在,请说明理由.

\begin{center}
\includegraphics[width=0.2\textwidth]{bodyxb/src/95_1068_1331_329_260_0.jpg}
\end{center}


第 15 题图

16. (2024$\cdot$金山区二模$\cdot$20)已知椭圆 $\Gamma  : \dfrac{{x}^{2}}{4} + \dfrac{{y}^{2}}{3} = 1$ 的右焦点为 $F$ ，直线 $l$ 与椭圆 $\Gamma$ 交于不同的两点 $M\left( {{x}_{1},{y}_{1}}\right)$ 、 $N\left( {{x}_{2},{y}_{2}}\right)$.

(1)证明:点 $M$ 到右焦点 $F$ 的距离为 $2 - \dfrac{{x}_{1}}{2}$ ；

(2)设点 $Q\left( {0,\dfrac{1}{2}}\right)$ ，当直线 $l$ 的斜率为 $\dfrac{1}{2}$ ，且 $\vv{QF}$ 与 $\vv{QM} + \vv{QN}$ 平行时，求直线 $l$ 的方程；

(3)当直线 $l$ 与 $x$ 轴不垂直，且 $\bigtriangleup  {MNF}$ 的周长为 4 时，试判断直线 $l$ 与圆 $C$ : ${x}^{2} + {y}^{2} = 3$ 的位置关系,并证明你的结论.

17.(2025$\cdot$宝山区一模$\cdot$20)已知椭圆 $\Gamma  : \dfrac{{x}^{2}}{9} + \dfrac{{y}^{2}}{3} = 1$ ，直线 $l$ 经过椭圆 $\Gamma$ 的右顶点 $P$ 且与椭圆交于另一点 $A$ ,设线段 ${AP}$ 的中点为 $M$.

(1)求椭圆 $\Gamma$ 的焦距和离心率；

(2)若 ${k}_{OM} =  - \dfrac{1}{3}$ ，求直线 ${AP}$ 的方程；

(3)过点 $P$ 再作一条直线与椭圆 $\Gamma$ 交于点 $B$ ，线段 ${BP}$ 的中点为 N. 若 ${OM}\bot {ON}$ ，则直线 ${AB}$ 是否经过定点？若经过定点，求出定点坐标; 若不经过定点, 请说明理由.

\begin{center}
\includegraphics[width=0.2\textwidth]{bodyxb/src/96_1139_1402_331_251_0.jpg}
\end{center}


第 17 题图

18.(2025$\cdot$嘉定区一模$\cdot$20)在平面直角坐标系 ${xOy}$ 中，已知椭圆 $\Gamma  : \dfrac{{x}^{2}}{5} + \dfrac{{y}^{2}}{4} = 1,{F}_{1}\text{ 、 }{F}_{2}$ 是其左、右焦点,过椭圆 $\Gamma$ 右焦点 ${F}_{2}$ 的直线 ${PQ}$ 交椭圆于 $P\text{ 、 }Q$ 两点.

(1)若 $\vv{P{F}_{1}} \cdot  \vv{P{F}_{2}} = 3$ ，求点 $P$ 的坐标；

(2)若 $\bigtriangleup {F}_{1}{PQ}$ 的面积为 $\dfrac{40}{21}$ ，求直线 ${PQ}$ 的方程；

(3)设直线 $l$ 与椭圆 $\Gamma$ 交于 $A$ 、 $B$ 两点， $M$ 为线段 ${AB}$ 的中点. 当 ${k}_{OM} \cdot  {k}_{AB} = {k}_{OA} \cdot  {k}_{OB}$ 时， $\bigtriangleup {OAB}$ 的面积是否为定值？如果是,请求出这个定值; 如果不是,请说明理由.

19.(2024$\cdot$黄浦区二模$\cdot$20)如图所示，已知 ${\Gamma }_{1}$ 是中心在坐标原点、焦点在 $x$ 轴上的椭圆， ${\Gamma }_{2}$ 是以 ${\Gamma }_{1}$ 的焦点 ${F}_{1}$ 、 ${F}_{2}$ 为顶点的等轴双曲线，点 $M\left( {\dfrac{5}{3},\dfrac{4}{3}}\right)$ 是 ${\Gamma }_{1}$ 与 ${\Gamma }_{2}$ 的一个交点，动点 $P$ 在 ${\Gamma }_{2}$ 的右支上且异于顶点.

(1)求 ${\Gamma }_{1}$ 与 ${\Gamma }_{2}$ 的方程；

(2)若直线 $P{F}_{2}$ 的倾斜角是直线 $P{F}_{1}$ 的倾斜角的 2 倍,求点 $P$ 的坐标；

(3)设直线 $P{F}_{1}\text{ 、 }P{F}_{2}$ 的斜率分别为 ${k}_{1}\text{ 、 }{k}_{2}$ ，直线 $P{F}_{1}$ 与 ${\Gamma }_{1}$ 相交于点 $A\text{ 、 }B$ ,直线 $P{F}_{2}$ 与 ${\Gamma }_{1}$ 相交于点 $C\text{ 、 }D,\left| {A{F}_{1}}\right|  \cdot  \left| {B{F}_{1}}\right|  = m,\left| {C{F}_{2}}\right|  \cdot  \left| {D{F}_{2}}\right|  = n$ ,求证: ${k}_{1}{k}_{2} = 1$ 且存在常数 $s$ 使得 $m + n = {smn}$.

\begin{center}
\includegraphics[width=0.2\textwidth]{bodyxb/src/97_1054_1394_337_329_0.jpg}
\end{center}


第 19 题图

20.(2024$\cdot$嘉定区二模$\cdot$20)如图所示，已知三点 $A$ 、 $B$ 、 $P$ 都在椭圆 $\dfrac{{x}^{2}}{4} + \dfrac{{y}^{2}}{2} = 1$ 上.

(1)若点 $A$ 、 $B$ 、 $P$ 都是椭圆的顶点，求 $\bigtriangleup  {ABP}$ 的面积；

\begin{center}
\includegraphics[width=0.3\textwidth]{bodyxb/src/98_1093_294_352_278_0.jpg}
\end{center}


第 20 题图

(2)若直线 ${AB}$ 的斜率为 1，求弦 ${AB}$ 中点 $M$ 的轨迹方程；

(3)若直线 ${AB}$ 的斜率为 2，设直线 ${PA}$ 的斜率为 ${k}_{PA}$ ，直线 ${PB}$ 的斜率为 ${k}_{PB}$ ,是否存在定点 $P$ ,使得 ${k}_{PA} + {k}_{PB} = 0$ 恒成立? 若存在,求出所有满足条件的点 $P$ ; 若不存在,请说明理由.

21. (2024$\cdot$浦东新区二模$\cdot$20)已知椭圆 $C : \dfrac{{x}^{2}}{2} + {y}^{2} = 1$ ，点 ${F}_{1}$ 、 ${F}_{2}$ 分别为椭圆的左、右焦点.

(1)若椭圆上点 $P$ 满足 $P{F}_{2} \bot  {F}_{1}{F}_{2}$ ，求 $\left| {P{F}_{1}}\right|$ 的值；

(2)点 $A$ 为椭圆的右顶点，定点 $T\left( {t,0}\right)$ 在 $x$ 轴上，若点 $S$ 为椭圆上一动点，当 $\left| {ST}\right|$ 取得最小值时,点 $S$ 恰与点 $A$ 重合,求实数 $t$ 的取值范围;

(3)已知 $m$ 为常数，过点 ${F}_{2}$ 且法向量为 $\left( {1, - m}\right)$ 的直线 $l$ 交椭圆于 $M$ ， $N$ 两点，若椭圆 $C$ 上存在点 $R$ 满足 $\vv{OR} = \lambda \vv{OM} + \mu \vv{ON}\left( {\lambda ,\mu  \in  \mathbf{R}}\right)$ ,求 ${\lambda \mu }$ 的最大值.

22.(2024$\cdot$普陀区二模$\cdot$20)设椭圆 $\Gamma  : \dfrac{{x}^{2}}{{a}^{2}} + {y}^{2} = 1\left( {a > 1}\right)$ ， $\Gamma$ 的离心率是短轴长的 $\dfrac{\sqrt{2}}{4}$ 倍，直线 $l$ 交 $\Gamma$ 于 $A$ 、 $B$ 两点， $C$ 是 $\Gamma$ 上异于 $A$ 、 $B$ 的一点， $O$ 是坐标原点.

(1)求椭圆 $\Gamma$ 的方程；

(2)若直线 $l$ 过 $\Gamma$ 的右焦点 $F$ ，且 $\vv{CO} = \vv{OB}$ ， $\vv{CF} \cdot  \vv{AB} = 0$ ，求 ${S}_{\bigtriangleup {CBF}}$ 的值；

(3)设直线 $l$ 的方程为 $y = {kx} + m\left( {k\text{ 、 }m \in  \mathbf{R}}\right)$ ，且 $\vv{OA} + \vv{OB} = \vv{CO}$ ，求 $\left| \vv{AB}\right|$ 的取值范围.

23.(2024$\cdot$徐汇区二模$\cdot$20)已知椭圆 $C : \dfrac{{x}^{2}}{4} + \dfrac{{y}^{2}}{3} = 1,{A}_{1}$ 、 ${A}_{2}$ 分别为椭圆 $C$ 的左、右顶点， ${F}_{1}$ 、 ${F}_{2}$ 分别为左、右焦点,直线 $l$ 交椭圆 $C$ 于 $M$ 、 $N$ 两点 $\left( l\right.$ 不过点 $\left. {A}_{2}\right)$.

(1)若 $Q$ 为椭圆 $C$ 上(除 ${A}_{1}$ 、 ${A}_{2}$ 外)任意一点，求直线 $Q{A}_{1}$ 和 $Q{A}_{2}$ 的斜率之积；

(2)若 $\vv{N{F}_{1}} = 2\vv{{F}_{1}M}$ ，求直线 $l$ 的方程；

(3)若直线 $M{A}_{2}$ 与直线 $N{A}_{2}$ 的斜率分别是 ${k}_{1}$ 、 ${k}_{2}$ ，且 ${k}_{1}{k}_{2} =  - \dfrac{9}{4}$ ，求证:直线 $l$ 过定点.

24.(2024$\cdot$杨浦区二模$\cdot$20)已知椭圆 $\Gamma  : \dfrac{{x}^{2}}{{a}^{2}} + \dfrac{{y}^{2}}{{b}^{2}} = 1\left( {a > b > 0}\right)$ 的上顶点为 $A\left( {0,1}\right)$ ，离心率 $e = \dfrac{\sqrt{3}}{2}$ ， 过点 $P\left( {-2,1}\right)$ 的直线 $l$ 与椭圆 $\Gamma$ 交于 $B$ 、 $C$ 两点,直线 ${AB}$ 、 ${AC}$ 分别与 $x$ 轴交于点 $M$ 、 $N$.

(1)求椭圆 $\Gamma$ 的方程；

(2)已知命题“对任意直线 $l$ ，线段 ${MN}$ 的中点为定点”为真命题，求 $\bigtriangleup  {AMN}$ 重心的坐标；

(3)是否存在直线 $l$ ，使得 ${S}_{\bigtriangleup {AMN}} = 2{S}_{\bigtriangleup {ABC}}$ ？若存在，求出所有满足条件的直线 $l$ 的方程；若不存在，请说明理由. (其中 ${S}_{\bigtriangleup {AMN}}$ 、 ${S}_{\bigtriangleup {ABC}}$ 分别表示 $\bigtriangleup  {AMN}$ 、 $\bigtriangleup  {ABC}$ 的面积)

25.(2024$\cdot$宝山区一模$\cdot$20)以坐标原点为对称中心，焦点在 $x$ 轴上的椭圆 $\Gamma$ 过点 $A$ (-2,0)，且离心率为 $\dfrac{\sqrt{3}}{2}$.

(1)求椭圆 $\Gamma$ 的方程；

(2)若点 $B\left( {1,0}\right)$ ，动点 $M$ 满足 $\left| {MA}\right|  = 2\left| {MB}\right|$ ，求动点 $M$ 的轨迹所围成的图形的面积；

(3)过圆 ${x}^{2} + {y}^{2} = 4$ 上一点 $P$ (不在坐标轴上)作椭圆 $\Gamma$ 的两条切线 ${l}_{1}$ 、 ${l}_{2}$. 记 ${OP}$ 、 ${l}_{1}$ 、 ${l}_{2}$ 的斜率分别为 ${k}_{0}$ 、 ${k}_{1}$ 、 ${k}_{2}$ ，求证: ${k}_{0}\left( {{k}_{1} + {k}_{2}}\right)  =  - 2$.

26.(2024$\cdot$长宁区二模$\cdot$20)已知椭圆 $\Gamma  : \dfrac{{x}^{2}}{6} + \dfrac{{y}^{2}}{3} = 1$ ， $O$ 为坐标原点.

(1)求 $\Gamma$ 的离心率 $e$ ；

(2)设点 $N\left( {1,0}\right)$ ，点 $M$ 在 $\Gamma$ 上，求 $\left| {MN}\right|$ 的最大值和最小值；

(3)点 $T\left( {2,1}\right)$ ，点 $P$ 在直线 $x + y = 3$ 上，过点 $P$ 且与 ${OT}$ 平行的直线 $l$ 与 $\Gamma$ 交于 $A$ 、 $B$ 两点； 试探究: 是否存在常数 $\lambda$ ,使得 $\left| {\vv{PA} \cdot  \vv{PB}}\right|  = \lambda {\left| \vv{PT}\right| }^{2}$ 恒成立? 若存在,求出该常数的值; 若不存在, 说明理由.

27.(2024$\cdot$奉贤区一模$\cdot$20)已知椭圆 $\dfrac{{x}^{2}}{{a}^{2}} + \dfrac{{y}^{2}}{{b}^{2}} = 1\left( {a > b > 0}\right)$ 的焦距为 $2\sqrt{3}$ 、离心率为 $\dfrac{\sqrt{3}}{2}$ ，椭圆的左右焦点分别为 ${F}_{1}$ 、 ${F}_{2}$ ，直角坐标原点记为 $O$. 设点 $P\left( {0,t}\right)$ ，过点 $P$ 作倾斜角为锐角的直线 $l$ 与椭圆交于不同的两点 $B$ 、 $C$.

(1)求椭圆的方程；

(2)设椭圆上有一动点 $T$ ，求 $\vv{PT} \cdot  \left( {\vv{T{F}_{1}} - \vv{T{F}_{2}}}\right)$ 的取值范围；

(3)设线段 ${BC}$ 的中点为 $M$ ，当 $t \geq  \sqrt{2}$ 时，判别椭圆上是否存在点 $Q$ ，使得非零向量 $\vv{OM}$ 与向量 $\vv{PQ}$ 平行，请说明理由.

28.(2024$\cdot$长宁区一模$\cdot$20)已知椭圆 $\Gamma  : \dfrac{{x}^{2}}{4} + \dfrac{{y}^{2}}{2} = 1,{F}_{1}$ 、 ${F}_{2}$ 为 $\Gamma$ 的左、右焦点，点 $A$ 在 $\Gamma$ 上，直线 $l$ 与圆 $C : {x}^{2} + {y}^{2} = 2$ 相切.

(1)求 $\bigtriangleup  A{F}_{1}{F}_{2}$ 的周长;

(2)若直线 $l$ 经过 $\Gamma$ 的右顶点,求直线 $l$ 的方程；

(3)设点 $D$ 在直线 $y = 2$ 上， $O$ 为原点，若 ${OA}\bot {OD}$ ，求证:直线 ${AD}$ 与圆 $C$ 相切.

\begin{center}
\includegraphics[width=0.2\textwidth]{bodyxb/src/102_1132_452_318_272_0.jpg}
\end{center}


第 28 题图

29.(2025$\cdot$浦东新区一模$\cdot$20)已知椭圆 $\dfrac{{x}^{2}}{8} + \dfrac{{y}^{2}}{4} = 1$ 的左、右焦点分别为 ${F}_{1}$ 、 ${F}_{2}$ ，过坐标原点的直线交椭圆于 $A$ 、 $B$ 两点,点 $A$ 在第一象限.

(1)若 $\left| {OA}\right|  = \sqrt{6}$ ，求点 $A$ 的坐标；

(2)求 $\left| {\vv{A{F}_{1}} + 3\vv{A{F}_{2}}}\right|$ 的取值范围；

(3)若 ${AE}\bot x$ 轴，垂足为 $E$ ，连接 ${BE}$ 并延长交椭圆于点 $C$ ，求 $\bigtriangleup {ABC}$ 面积的最大值.

\begin{center}
\includegraphics[width=0.2\textwidth]{bodyxb/src/102_1118_1397_331_256_0.jpg}
\end{center}


第 29 题图



30. (2025$\cdot$普陀区二模$\cdot$20)设 $a > b > 0,m > 0$ ，点 $A$ 、 $F$ 分别是椭圆 $\Gamma  : \dfrac{{x}^{2}}{{a}^{2}} + \dfrac{{y}^{2}}{{b}^{2}} = 1$ 的上顶点与右焦点,且 $\left| {AF}\right|  = 2$ ,直线 $l : x - {my} - 1 = 0$ 经过点 $F$ 且与 $\Gamma$ 交于 $P\text{ 、 }Q$ 两点, $O$ 是坐标原点.

(1)求椭圆 $\Gamma$ 的方程；

(2)若 $m = \sqrt{3}$ ，点 $M$ 是 $x$ 轴上的一点，且 $\bigtriangleup  {MPQ}$ 的面积为 $\dfrac{6}{13}$ ，求点 $M$ 的坐标；

(3)若点 $G$ 在直线 $x = 5$ 上，向量 $\vv{PG}$ 在直线 $l$ 上的投影为向量 $\vv{PF}$ ，求证: $\angle {PGQ} < \dfrac{\pi }{4}$.

31. (2025$\cdot$浦东新区二模$\cdot$20)已知椭圆 ${C}_{1}$ 的方程为 $\dfrac{{x}^{2}}{3} + {y}^{2} = 1$ ，右顶点为 $A$ ，上顶点为 $B$ ，椭圆 ${C}_{2}$ 的中心位于坐标原点,两个椭圆的离心率相等.

(1)若椭圆 ${C}_{2}$ 的方程是 $\dfrac{{x}^{2}}{{a}^{2}} + \dfrac{{y}^{2}}{2a} = 1\left( {a > 0}\right)$ ，焦点在 $x$ 轴上，求 $a$ 的值；

(2)设椭圆 ${C}_{2}$ 的焦点在 $x$ 轴上，直线 ${AB}$ 与 ${C}_{2}$ 相交于点 $C$ 、 $D$ ，若 $\left| {CD}\right|  = 3\left| {AB}\right|$ ，求 ${C}_{2}$ 的标准方程；

(3)设椭圆 ${C}_{2}$ 的焦点在 $y$ 轴上，点 $P$ 在 ${C}_{1}$ 上，点 $Q$ 在 ${C}_{2}$ 上. 若存在 $\bigtriangleup {APQ}$ 是等腰直角三角形,且 $\left| {AP}\right|  = \left| {AQ}\right|$ ,求 ${C}_{2}$ 的长轴的取值范围.

32.(2025$\cdot$长宁区一模$\cdot$20)已知椭圆的左、右焦点分别为 ${F}_{1}\left( {-1,0}\right)$ 、 ${F}_{2}\left( {1,0}\right)$ ，且经过点 $P\left( {-1,\dfrac{3}{2}}\right)$.

(1)求该椭圆的离心率；

(2)点 $Q$ 为椭圆上一点，且位于第三象限，若 $\bigtriangleup  {PQ}{F}_{2}$ 的面积为 3，求点 $Q$ 的坐标；

(3) $A$ 、 $B$ 、 $C$ 、 $D$ 是椭圆上不重合的四个点， ${AB}$ 与 ${CD}$ 相交于点 ${F}_{1}$ ， 且 $\vv{AB} \cdot  \vv{CD} = 0$ ，求 $\left| {AB}\right|  + \left| {CD}\right|$ 的取值范围.

\begin{center}
\includegraphics[width=0.2\textwidth]{bodyxb/src/104_1120_446_334_318_0.jpg}
\end{center}


第 32 题图

33. ( 2024 年闭行区二模$\cdot$20 )如图所示，已知椭圆 ${C}_{1} : \dfrac{{x}^{2}}{4} + {y}^{2} = 1$ 和抛物线 ${C}_{2} : {x}^{2} = {2py}\left( {p > 0}\right)$ ， ${C}_{2}$ 的焦点 $F$ 是 ${C}_{1}$ 的上顶点,过 $F$ 的直线交 ${C}_{2}$ 于 $M\text{ 、 }N$ 两点,连接 ${NO}\text{ 、 }{MO}$ 并延长之,分别交 ${C}_{1}$ 于 $A\text{ 、 }B$ 两点,连接 ${AB}$ ,设 $\bigtriangleup {OMN}\text{ 、 }\bigtriangleup {OAB}$ 的面积分别为 ${S}_{\bigtriangleup {OMN}}\text{ 、 }{S}_{\bigtriangleup {OAB}}$.

(1)求 $p$ 的值；

\begin{center}
\includegraphics[width=0.2\textwidth]{bodyxb/src/104_1110_1331_340_247_0.jpg}
\end{center}


第 33 题图

(2)求 $\vv{OM} \cdot  \vv{ON}$ 的值；

(3)求 $\dfrac{{S}_{\bigtriangleup {OMN}}}{{S}_{\bigtriangleup {OAB}}}$ 的取值范围.

34. (2024$\cdot$青浦区一模$\cdot$20)已知椭圆 $\Gamma$ 的离心率是 $\dfrac{1}{2}$ ，长轴长为 4，椭圆的中心是坐标原点，焦点在 $x$ 轴上.

(1)求椭圆 $\Gamma$ 的标准方程；

(2)已知 $A$ 、 $B$ 、 $C$ 是椭圆 $\Gamma$ 上三个不同的点， $F$ 是椭圆 $\Gamma$ 的右焦点，若原点 $O$ 是 $\bigtriangleup  {ABC}$ 的重心， 求 $\left| {FA}\right|  + \left| {FB}\right|  + \left| {FC}\right|$ 的值;

(3)已知 $T\left( {1,1}\right)$ ，椭圆 $\Gamma$ 上四个动点 $M$ 、 $N$ 、 $P$ 、 $Q$ 满足 $\vv{MT} = 3\vv{TQ}$ ， $\vv{NT} = 3\vv{TP}$ ，求直线 ${MN}$ 的方程.

\begin{center}
\includegraphics[width=0.2\textwidth]{bodyxb/src/105_1109_485_305_208_0.jpg}
\end{center}


第 34 题图

35.(2025$\cdot$闵行区二模$\cdot$20)已知椭圆 $\Gamma  : \dfrac{{x}^{2}}{4} + \dfrac{{y}^{2}}{3} = 1$ ，左、右焦点分别为 ${F}_{1}$ 、 ${F}_{2}$ ，上、下顶点分别为 $A$ 、 $B$ ，左、右顶点分别为 $C$ 、 $D$ ， $P$ 、 $Q$ 是 $\Gamma$ 上异于椭圆顶点的两点.

(1)求 $\bigtriangleup A{F}_{1}{F}_{2}$ 的周长；

(2)若点 $Q$ 在第一象限且满足 $\bigtriangleup {ABQ}$ 的面积比 $\bigtriangleup {F}_{1}{F}_{2}Q$ 的面积大，求点 $Q$ 的横坐标的取值范围;

(3)记点 $A$ 在直线 ${PQ}$ 上的投影为 $H$ ，且直线 ${CP}$ 的斜率是直线 ${DQ}$ 的斜率的 3 倍，试判断:过点 $A$ 、 $H$ 、 $O$ ( $O$ 为坐标原点)三点的圆是否为定圆？若是，求出该圆的方程；若不是，请说明理由.

36.(2024$\cdot$松江区一模$\cdot$20)已知椭圆 $\Gamma  : \dfrac{{y}^{2}}{{a}^{2}} + \dfrac{{x}^{2}}{{b}^{2}} = 1\left( {a > b > 0}\right)$ 的离心率为 $\dfrac{\sqrt{2}}{2}$ ，其上焦点 $F$ 与抛物线 $K : {x}^{2} = {4y}$ 的焦点重合.

(1)椭圆 $\Gamma$ 的方程；

(2)若过点 $F$ 的直线交椭圆 $\Gamma$ 于点 $A$ 、 $B$ ，同时交抛物线 $K$ 于点 $C$ 、 $D$ (如图 1 所示，点 $C$ 在椭圆与抛物线第一象限交点上方),试比较线段 ${AC}$ 与 ${BD}$ 长度的大小,并说明理由;

(3)若过点 $F$ 的直线交椭圆 $\Gamma$ 于点 $A\text{ 、 }B$ ，过点 $F$ 与直线 ${AB}$ 垂直的直线 ${EG}$ 交抛物线 $K$ 于点 $E$ 、 $G$ (如图 2 所示)，试求四边形 ${AEBG}$ 面积的最小值.

\begin{center}
\includegraphics[width=0.4\textwidth]{bodyxb/src/106_974_593_501_300_0.jpg}
\end{center}


第 36 题图

37.(2025$\cdot$青浦区一模$\cdot$20)已知椭圆 $C : \dfrac{{x}^{2}}{4} + \dfrac{{y}^{2}}{3} = 1$ ， $F$ 为椭圆 $C$ 的右焦点，过点 $F$ 的直线 $l$ 交椭圆 $C$ 于 $A$ 、 $B$ 两点.

(1)若直线 $l$ 垂直于 $x$ 轴,求椭圆 $C$ 的弦 ${AB}$ 的长度;

\begin{center}
\includegraphics[width=0.2\textwidth]{bodyxb/src/106_1137_1284_331_245_0.jpg}
\end{center}


第 37 题图

(2)设点 $P\left( {-3,0}\right)$ ，当 $\angle {PAB} = {90}^{ \circ  }$ 时，求点 $A$ 的坐标；

(3)设点 $M\left( {3,0}\right)$ ，记 ${MA}$ 、 ${MB}$ 的斜率分别为 ${k}_{1}$ 和 ${k}_{2}$ ，求 ${k}_{1} + {k}_{2}$ 的取值范围.

38.(2024$\cdot$崇明区二模$\cdot$20)已知椭圆 $\Gamma  : \dfrac{{x}^{2}}{2} + {y}^{2} = 1$ ， $A$ 为 $\Gamma$ 的上顶点， $P$ 、 $Q$ 是 $\Gamma$ 上不同于点 $A$ 的两点.

(1)求椭圆 $\Gamma$ 的离心率；

(2)若 $F$ 是椭圆 $\Gamma$ 的右焦点， $B$ 是椭圆 下顶点， $R$ 是直线 ${AF}$ 上一点. 若 $\bigtriangleup  {ABR}$ 有一个内角为 $\dfrac{\pi }{3}$ ,求点 $R$ 的坐标;

(3)作 ${AH} \bot  {PQ}$ ，垂足为 $H$. 若直线 ${AP}$ 与直线 ${AQ}$ 的斜率之和为 2，是否存在 $x$ 轴上的点 $M$ ，使得 $\left| \vv{MH}\right|$ 为定值? 若存在,请求出点 $M$ 的坐标; 若不存在,请说明理由.

39. (2025$\cdot$黄浦区二模$\cdot$20)椭圆 $\Gamma  : \dfrac{{x}^{2}}{{a}^{2}} + \dfrac{{y}^{2}}{{b}^{2}} = 1\left( {a > b > 0}\right)$ 的左、右焦点分别为 ${F}_{1}\left( {-c,0}\right)$ 、 ${F}_{2}\left( {c,0}\right)  \left( {c > 0}\right)$ ,过点 ${F}_{1}$ 的直线 $l$ 与 $\Gamma$ 交于点 $P$.

(1)若 $c = 2$ ，点 $P$ 的坐标为 $\left( {2,\sqrt{2}}\right)$ ，求点 ${F}_{2}$ 到直线 $l$ 的距离；

(2)当 $b \leq  c$ 时,求满足 $P{F}_{1} \bot  P{F}_{2}$ 的点 $P$ 的个数；

(3)设直线 $l$ 与 $\Gamma$ 的另一个交点为 $Q,\vv{{F}_{1}Q} = \lambda \vv{QP}\left( {\lambda  \in  \mathbf{R}}\right)$ ，点 $P$ 的横坐标为 $\dfrac{c}{2}$ ，若 $\Gamma$ 的离心率 $e >  \dfrac{1}{2}$ ，求 $\lambda$ 的取值范围.

40. (2024$\cdot$奉贤区二模$\cdot$20)已知曲线 $C : \dfrac{{x}^{2}}{4} + \dfrac{{y}^{2}}{2} = 1,O$ 是坐标原点，过点 $T\left( {1,0}\right)$ 的直线 ${l}_{1}$ 与曲线 $C$ 交于 $P\text{ 、 }Q$ 两点.

(1)当 ${l}_{1}$ 与 $x$ 轴垂直时,求 $\bigtriangleup {OPQ}$ 的面积;

(2)过圆 ${x}^{2} + {y}^{2} = 6$ 上任意一点 $M$ 作直线 ${MA}$ 、 ${MB}$ ，分别与曲线 $C$ 切于 $A$ 、 $B$ 两点，求证: ${MA} \bot  {MB}$ ;

(3)过点 $N\left( {n,0}\right) \left( {n > 2}\right)$ 的直线 ${l}_{2}$ 与双曲线 $\dfrac{{x}^{2}}{4} - {y}^{2} = 1$ 交于 $R$ 、 $S$ 两点 $\left( {{l}_{1}\text{ 、 }{l}_{2}}\right.$ 不与 $x$ 轴重合). 记直线 ${TR}$ 的斜率为 ${k}_{TR}$ ,直线 ${TS}$ 斜率为 ${k}_{TS}$ ,当 $\angle {ONP} = \angle {ONQ}$ 时,求证: $n$ 与 ${k}_{TR} + {k}_{TS}$ 都是定值.

\begin{center}
\includegraphics[width=0.2\textwidth]{bodyxb/src/108_639_684_330_315_0.jpg}
\end{center}


第 40 题(2)图

\begin{center}
\includegraphics[width=0.3\textwidth]{bodyxb/src/108_1035_721_415_278_0.jpg}
\end{center}


第 40 题(3)图

41.(2025$\cdot$奉贤区一模$\cdot$20)椭圆 $\Gamma  : \dfrac{{x}^{2}}{{a}^{2}} + {y}^{2} = 1\left( {a > 1}\right)$ 的左、右焦点分别为 ${F}_{1}$ 、 ${F}_{2}$ ，设 $P\left( {{x}_{0},{y}_{0}}\right)$ 是第一象限内椭圆上的一点, $P{F}_{1}$ 的延长线交椭圆于点 $Q\left( {{x}_{1},{y}_{1}}\right)$.

(1)若椭圆的离心率 $\dfrac{\sqrt{2}}{2}$ ，求 $a$ 的值；

(2)若 $a = \sqrt{2},\overline{PQ} \cdot  \vv{O{F}_{1}} = \dfrac{12}{5}$ ，求 ${x}_{0}$ ；

(3)若 $a = 2$ ，过点 $T\left( {0,t}\right)$ 的直线 $l$ 与椭圆 $\Gamma$ 交于 $M$ 、 $N$ 两点，且 $\left| {MN}\right|  = 2$ ，则当 $t \geq  0$ 时，判断符合要求的直线有几条, 并说明理由.

\begin{center}
\includegraphics[width=0.2\textwidth]{bodyxb/src/108_763_1638_307_231_0.jpg}
\end{center}


第 41 题(2)图

\begin{center}
\includegraphics[width=0.2\textwidth]{bodyxb/src/108_1141_1636_309_230_0.jpg}
\end{center}


第 41 题(3)图

42.(2025$\cdot$上海秋季卷$\cdot$20)已知椭圆 $\Gamma  : \dfrac{{x}^{2}}{{a}^{2}} + \dfrac{{y}^{2}}{5} = 1\left( {a > \sqrt{5}}\right) ,M\left( {0,m}\right) \left( {m > 0}\right)$ ， $A$ 是 $\Gamma$ 的右顶点.

(1)若 $\Gamma$ 的焦点 $\left( {2,0}\right)$ ，求离心率 $e$ ；

(2)若 $a = 4$ ，且 $\Gamma$ 上存在一点 $P$ ，满足 $\vv{PA} = 2\vv{MP}$ ，求 $m$ ；

(3)已知 ${AM}$ 的中垂线 $l$ 的斜率为 $2,l$ 与 $\Gamma$ 交于 $C$ 、 $D$ 两点， $\angle {CMD}$ 为钝角，求 $a$ 的取值范围.

\subsection*{考点 4 双曲线方程及其性质}

\subsubsection*{一、填空题}

1. (2024$\cdot$松江区一模$\cdot$2)双曲线 $\dfrac{{x}^{2}}{3} - {y}^{2} = 1$ 的右焦点坐标是\blkda[2]{}.

2. (2024$\cdot$崇明区一模$\cdot$2)双曲线 ${x}^{2} - \dfrac{{y}^{2}}{4} = 1$ 的焦距是\blkda[2]{}.

3. (2024$\cdot$奉贤区一模$\cdot$3)曲线 ${x}^{2} - 2{y}^{2} = 1$ 的渐近线方程为\blkda[2]{}.

4.(2024$\cdot$黄浦区一模$\cdot$4)若双曲线 $\dfrac{{x}^{2}}{16} - \dfrac{{y}^{2}}{m} = 1$ 经过点 $\left( {{4\sqrt{2}},3}\right)$ ，则此双曲线的离心率为\blkda[2]{}.

5.(2024$\cdot$金山区一模$\cdot$4)双曲线 ${x}^{2} - \dfrac{{y}^{2}}{2} = 1$ 的离心率是\blkda[2]{}.

6.(2025$\cdot$青浦区一模$\cdot$5)两条渐近线互相垂直的双曲线的离心率为\blkda[2]{}.

7.(2025$\cdot$嘉定区一模$\cdot$5)已知双曲线 $C : \dfrac{{x}^{2}}{3} - \dfrac{{y}^{2}}{2} = 1$ ，则双曲线 $C$ 的离心率为\blkda[2]{}.

8.(2025$\cdot$崇明区一模$\cdot$5)双曲线 ${x}^{2} - \dfrac{{y}^{2}}{4} = 1$ 的渐近线方程是\blkda[2]{}.

9.(2024$\cdot$嘉定区一模$\cdot$5)双曲线 $\dfrac{{x}^{2}}{4} - \dfrac{{y}^{2}}{5} = 1$ 的离心率为\blkda[2]{}.

10.(2024$\cdot$闵行区一模$\cdot$6)若双曲线 $\dfrac{{x}^{2}}{{a}^{2}} - \dfrac{{y}^{2}}{{b}^{2}} = 1\left( {a > 0,b > 0}\right)$ 的离心率为 $\sqrt{2}$ ，则该双曲线的渐近线方程为\blkda[2]{}.

11. (2024$\cdot$虹口区一模$\cdot$7)双曲线 ${x}^{2} - \dfrac{{y}^{2}}{4} = 1$ 的两条渐近线夹角的余弦值为\blkda[2]{}.

12.(2024$\cdot$上海春季卷$\cdot$8)在 $\bigtriangleup  {ABC}$ 中， ${AB} = 5,{AC} = 7,{BC} = 6$ ，若以 $B$ 、 $C$ 为焦点，且过点 $A$ 的双曲线的离心率 $e =$ \blkda[2]{}.

13.(2024$\cdot$金山区二模$\cdot$8)已知双曲线 $C : \dfrac{{x}^{2}}{{a}^{2}} - \dfrac{{y}^{2}}{{b}^{2}} = 1\left( {a > 0,b > 0}\right)$ ，给定的四点 ${P}_{1}\left( {4, - 3}\right)$ 、 ${P}_{2}\left( {3,4}\right)$ 、 ${P}_{3}\left( {-4,3}\right) \text{ 、 }{P}_{4}\left( {-2,0}\right)$ 中恰有三个点在双曲线 $C$ 上，则该双曲线 $C$ 的离心率是\blkda[2]{}.

14.(2025$\cdot$静安区一模$\cdot$9)以双曲线 $\dfrac{{x}^{2}}{4} - \dfrac{{y}^{2}}{m} = 1$ 的离心率为半径，以右焦点为圆心的圆与双曲线的渐近线相切，则 $m$ 的值为\blkda[2]{}.

\begin{center}
\includegraphics[width=0.2\textwidth]{bodyxb/src/110_409_617_341_279_0.jpg}
\end{center}


第 13 题图

\begin{center}
\includegraphics[width=0.2\textwidth]{bodyxb/src/110_910_588_329_305_0.jpg}
\end{center}


第 14 题图

15.(2024$\cdot$松江区二模$\cdot$9)已知 ${F}_{1}$ 、 ${F}_{2}$ 分别为双曲线 $C : \dfrac{{x}^{2}}{{a}^{2}} - \dfrac{{y}^{2}}{{b}^{2}} = 1\left( {a > 0,b > 0}\right)$ 的左、右焦点，过 ${F}_{1}$ 的直线 $l$ 与双曲线 $C$ 的左、右两支分别交于 $A$ 、 $B$ 两点,若 $\left| {AB}\right|  : \left| {B{F}_{2}}\right|  : \left| {A{F}_{2}}\right|  = 3 : 4 :$ 5，则该双曲线的离心率为\blkda[2]{}.

16.(2025$\cdot$浦东新区一模$\cdot$9)已知双曲线 ${x}^{2} - \dfrac{{y}^{2}}{3} = 1$ 的左、右焦点分别为 ${F}_{1}$ 、 ${F}_{2}$ ，双曲线上的点 $P$ 在第一象限，且 $P{F}_{2}$ 与双曲线的一条渐近线平行，则 $\bigtriangleup  P{F}_{1}{F}_{2}$ 的面积为\blkda[2]{}.

17.(2024$\cdot$闵行区二模$\cdot$10)双曲线 $\Gamma  : {x}^{2} - \dfrac{{y}^{2}}{6} = 1$ 的左右焦点分别为 ${F}_{1}$ 、 ${F}_{2}$ ，过坐标原点的直线与 $\Gamma$ 相交于 $A$ 、 $B$ 两点，若 $\left| {{F}_{1}B}\right|  = 2\left| {{F}_{1}A}\right|$ ，则 $\vv{{F}_{2}A} \cdot  \vv{{F}_{2}B} =$ \blkda[2]{}.

18.(2025$\cdot$上海春季卷$\cdot$10)已知双曲线 $\dfrac{{x}^{2}}{{a}^{2}} - \dfrac{{y}^{2}}{6 - {a}^{2}} = 1\left( {a > 0}\right)$ 的左、右焦点分别为 ${F}_{1}$ 、 ${F}_{2}$. 通过 ${F}_{2}$ 且倾斜角为 $\dfrac{\pi }{3}$ 的直线与双曲线交于第一象限的点 $A$ ,延长 $A{F}_{2}$ 至 $B$ 使得 ${AB} = A{F}_{1}$. 若 $\bigtriangleup B{F}_{1}{F}_{2}$ 的面积为 $3\sqrt{6}$ ，则 $a$ 的值为\blkda[2]{}.

19.(2025$\cdot$徐汇区二模$\cdot$10)已知双曲线 $\dfrac{{x}^{2}}{{a}^{2}} - \dfrac{{y}^{2}}{{b}^{2}} = 1\left( {a > 0,b > 0}\right)$ 的左焦点为 ${F}_{1}$ 、右焦点为 ${F}_{2}$. 若双曲线的右支上存在一点 $P$ ，使得直线 $P{F}_{1}$ 与以双曲线的实轴为直径的圆相切，切点为线段 $P{F}_{1}$ 的中点，则该双曲线的离心率为\blkda[2]{}.

20.(2025$\cdot$虹口区一模$\cdot$10)双曲线 ${C}_{1} : \dfrac{{x}^{2}}{{a}^{2}} - \dfrac{{y}^{2}}{{b}^{2}} = 1$ 的左、右焦点分别为 ${F}_{1}$ 和 ${F}_{2}$ ，若以点 ${F}_{2}$ 为焦点的抛物线 ${C}_{2} : {y}^{2} = {2px}\left( {p > 0}\right)$ 与 ${C}_{1}$ 在第一象限交于点 $P$ ，且 $\angle {P{F}_{1}{F}_{2}} = \dfrac{\pi }{4}$ ，则 ${C}_{1}$ 的离心率为\blkda[2]{}.

21. (2024$\cdot$虹口区二模$\cdot$10)从某个角度观察篮球(如图 1) 可以得到一个对称的平面图形 (如图 2), 篮球的外轮廓为圆 $O$ ,将篮球的表面粘合线视为坐标轴和双曲线,若坐标轴和双曲线与圆 $O$ 的交点将圆的周长 8 等分，且 $\left| {AB}\right|  = \left| {BC}\right|  = \left| {CD}\right|$ ,则该双曲线的离心率为\blkda[2]{}.

\begin{center}
\includegraphics[width=0.2\textwidth]{bodyxb/src/111_912_223_200_214_0.jpg}
\end{center}


\begin{center}
\includegraphics[width=0.2\textwidth]{bodyxb/src/111_1115_176_277_270_0.jpg}
\end{center}


图1 图2

第 21 题图

22. (2024$\cdot$浦东新区二模$\cdot$11)已知双曲线 $\dfrac{{x}^{2}}{{a}^{2}} - \dfrac{{y}^{2}}{{b}^{2}} = 1$ ( $a >$

$0,b > 0)$ 的左、右焦点分别为 ${F}_{1}\text{ 、 }{F}_{2},M$ 为双曲线上一点，若 $\angle {F}_{1}M{F}_{2} = \dfrac{2\pi }{3},{OM} = \dfrac{\sqrt{21}}{3}b$ ，则双曲线的离心率为\blkda[2]{}.

23.(2025$\cdot$崇明区二模$\cdot$11)已知双曲线 ${x}^{2} - \dfrac{{y}^{2}}{{b}^{2}} = 1\left( {b > 0}\right)$ 的左、右焦点为 ${F}_{1}$ 、 ${F}_{2}$ ，以 $O$ 为顶点， ${F}_{2}$ 为焦点作抛物线交双曲线于 $P$ ，且 $\angle {P{F}_{1}{F}_{2}} = {45}^{ \circ  }$ ，则 ${b}^{2} =$ \blkda[2]{}.

\subsubsection*{二、选择题}

24.(2024$\cdot$嘉定区二模$\cdot$13)双曲线 ${\Gamma }_{1} : \dfrac{{x}^{2}}{4} - \dfrac{{y}^{2}}{2} = 1$ 和双曲线 ${\Gamma }_{2} : \dfrac{{y}^{2}}{4} - \dfrac{{x}^{2}}{2} = 1$ 具有相同的\dotfill{(\kh{})}

\xx{}{}{}{}焦点 B. 顶点 C. 渐近线 D. 离心率

25.(2024$\cdot$静安区二模$\cdot$15)设 $a > 1$ ，则双曲线 $\dfrac{{x}^{2}}{{a}^{2}} - \dfrac{{y}^{2}}{{\left( a + 1\right) }^{2}} = 1$ 的离心率， $e$ 的取值范围是\dotfill{(\kh{})}

\xx{}{}{}{}$\left( {\sqrt{2},2}\right)$ B. $\left( {\sqrt{2},\sqrt{5}}\right)$ C. $\left( {2,5}\right)$ D. $\left( {2,\sqrt{5}}\right)$

26.(2025$\cdot$上海秋季卷$\cdot$15)已知 $A\left( {0,1}\right)$ 、 $B\left( {1,2}\right)$ ， $C$ 在 $\Gamma  : {x}^{2} - {y}^{2} = 1$ ( $x \geq  1$ ， $y \geq  0$ )上，则 $\bigtriangleup  {ABC}$ 的面积 \dotfill{(\kh{})}

\xx{}{}{}{}有最大值，但没有最小值 B. 没有最大值, 但有最小值

C. 既有最大值, 也有最小值 D. 既没有最大值，也没有最小值

27.(2025$\cdot$普陀区二模$\cdot$15)设 $a > 0$ ， $b > 0$ ，点 $M\left( {{x}_{0},{y}_{0}}\right)$ ， $O$ 是坐标原点， $\left| {OM}\right|  = a$ ， $F$ 是双曲线， $\dfrac{{x}^{2}}{{a}^{2}} - \dfrac{{y}^{2}}{{b}^{2}} = 1$ 的左焦点，若直线 $l : {x}_{0}x + {y}_{0}y = {a}^{2}$ 经过点 $F$ ，且与双曲线 $\Gamma$ 的右支在第一象限内交于 $P$ 点,则双曲线 $\Gamma$ 的离心率的一个可能的值是 \dotfill{(\kh{})}

\xx{}{}{}{}$\dfrac{\sqrt{5}}{2}$ B. $\dfrac{\sqrt{6}}{2}$ C. $\dfrac{\sqrt{7}}{2}$ D. $\dfrac{\sqrt{11}}{2}$

28.(2025$\cdot$金山区一模$\cdot$15)古希腊数学家阿波罗尼奥斯用不同的平面截同一圆锥，得到了圆锥曲线,其中的一种如图所示. 用过点 $M$ 且垂直于圆锥底面的平面截两个全等的对顶圆锥得到双曲线的一部分,已知高 ${PO} = 2$ ,底面圆的半径为 4, $M$ 为母线 ${PB}$ 的中点,平面与底面的交线 ${EF} \bot  {AB}$ ,则双曲线的两条渐近线所成角的余弦值为 \dotfill{(\kh{})}

\begin{center}
\includegraphics[width=0.2\textwidth]{bodyxb/src/111_1147_1734_258_287_0.jpg}
\end{center}


第 28 题图

\xx{}{}{}{}$\dfrac{5}{6}$ B. $\dfrac{4}{5}$

C. $\dfrac{3}{4}$ D. $\dfrac{3}{5}$

\subsubsection*{三、解答题}

29. (2024$\cdot$静安区一模$\cdot$19)已知双曲线 $C : \dfrac{{x}^{2}}{2} - {y}^{2} = 1$ ，点 $M$ 的坐标为 $\left( {0,1}\right)$.

(1)设直线 $l$ 过点 $M$ ，斜率为 $\dfrac{1}{2}$ ，它与双曲线 $C$ 交于 $A$ 、 $B$ 两点，求线段 ${AB}$ 的长；

(2)设点 $P$ 在双曲线 $C$ 上， $Q$ 是点 $P$ 关于 $y$ 轴的对称点. 记 $k = \vv{MP} \cdot  \vv{MQ}$ ,求 $k$ 的取值范围.

\begin{center}
\includegraphics[width=0.3\textwidth]{bodyxb/src/112_1044_472_416_369_0.jpg}
\end{center}


第 29 题图

30. (2025$\cdot$松江区一模$\cdot$20)如果一条双曲线的实轴和虚轴分别是一个椭圆的长轴和短轴，则称它们为“共轴”曲线. 若双曲线 ${C}_{1}$ 与椭圆 ${C}_{2}$ 是 “共轴” 曲线，且椭圆 ${C}_{2} : \dfrac{{x}^{2}}{9} + \dfrac{{y}^{2}}{{b}^{2}} = 1\left( {0 < b < 3}\right) ,{e}_{1}{e}_{2} =  \dfrac{4\sqrt{5}}{9}\left( {{e}_{1}\text{ 、 }{e}_{2}}\right.$ 分别为曲线 ${C}_{1}\text{ 、 }{C}_{2}$ 的离心率). 已知点 $M\left( {1,0}\right)$ ,点 $P$ 为双曲线 ${C}_{1}$ 上任意一点.

(1)求双曲线 ${C}_{1}$ 的方程；

(2)延长线段 ${PM}$ 到点 $Q$ ，且 $\left| {PM}\right|  = 2\left| {MQ}\right|$ ，若点 $Q$ 在椭圆 ${C}_{2}$ 上，试求点 $P$ 的坐标；

(3)若点 $P$ 在双曲线 ${C}_{1}$ 的右支上，点 $A$ 、 $B$ 分别为双曲线 ${C}_{1}$ 的左、右顶点，直线 ${PM}$ 交双曲线的左支于点 $R$ ，直线 ${AP}$ 、 ${BR}$ 的斜率分别为 ${k}_{AP}$ 、 ${k}_{BR}$ ，是否存在实数 $\lambda$ ，使得 ${k}_{AP} = \lambda {k}_{BR}$ ？若存在，求出 $\lambda$ 的值；若不存在，请说明理由.

31.(2025$\cdot$黄浦区一模$\cdot$20)双曲线 $\Gamma  : \dfrac{{x}^{2}}{{a}^{2}} - \dfrac{{y}^{2}}{{b}^{2}} = 1\left( {a > 0,b > 0}\right)$ 的左、右焦点分别为 ${F}_{1}\left( {-c,0}\right)$ 、 ${F}_{2}\left( {c,0}\right) \left( {c > 0}\right)$ ,过点 ${F}_{1}$ 的直线 $l$ 与 $\Gamma$ 右支在 $x$ 轴上方交于点 $A$.

(1)若 $a = \sqrt{5}$ ，点 $A$ 的坐标为 $\left( {3,4}\right)$ ，求 $c$ 的值；

(2)若 $A{F}_{2} \bot  {F}_{1}{F}_{2}$ ，且 $a$ ， $b$ ， $c$ 是等比数列，求证:直线 $l$ 的斜率为定值；

(3)设直线 $l$ 与 $\Gamma$ 左支的交点为 $B,c = 3$ ，求当且仅当 $a$ 满足什么条件时，存在直线 $l$ ，使得 $\left| {AB}\right|  = \left| {A{F}_{2}}\right|$ 成立?

32. (2025$\cdot$虹口区二模$\cdot$20)已知点 ${F}_{1}$ 、 ${F}_{2}$ 是双曲线 $C$ : $\dfrac{{x}^{2}}{{a}^{2}} - {y}^{2} = 1\left( {a > 0}\right)$ 的左、右焦点.

(1)若 $y = x$ 是双曲线 $C$ 的一条渐近线,求 $C$ 的离心率；

(2)当 $a = \sqrt{2}$ 时,若双曲线 $C$ 上存在一点 $P$ 满足 $\left| {P{F}_{1}}\right|  + \left| {P{F}_{2}}\right|  = 4$ ,求 $\bigtriangleup P{F}_{1}{F}_{2}$ 的面积;

(3)若在双曲线 $C$ 上分别存在两点 $A$ 和 $B$ ，点 $A$ 在第一象限，点 $B$ 在第二象限，使得四边形 ${F}_{1}{BA}{F}_{2}$ 的面积为 $2\sqrt{{a}^{2} + 1}$ ，且存在实数 $\lambda$ 使 $\overline{{F}_{2}A} = \lambda \overline{{F}_{1}B}$ ，求实数 $a$ 的取值范围.

33. (2025$\cdot$徐汇区一模$\cdot$20)已知过点 $P\left( {3,\sqrt{2}}\right)$ 的双曲线 $C$ 的渐近线方程为 $x \pm  \sqrt{3}y = 0$. 如图所示， 过双曲线 $C$ 的右焦点 $F$ 作与坐标轴都不垂直的直线 $l$ 交 $C$ 的右支于 $A$ 、 $B$ 两点.

(1)求双曲线 $C$ 的标准方程；

(2)已知点 $Q\left( {\dfrac{3}{2},0}\right)$ ，求证: $\angle {AQF} = \angle {BQF}$ ；

(3)若以 ${AB}$ 为直径的圆被直线 $x = \dfrac{3}{2}$ 截得的劣弧为 $\overset{\frown}{MN}$ ,则 $\overset{\frown}{MN}$ 所对圆心角的大小是否为定值？若是,求出该定值; 若不是,请说明理由.

\begin{center}
\includegraphics[width=0.3\textwidth]{bodyxb/src/114_1090_443_374_246_0.jpg}
\end{center}


第 33 题图

34.(2024$\cdot$徐汇区一模$\cdot$20)已知双曲线 $E : \dfrac{{x}^{2}}{{a}^{2}} - \dfrac{{y}^{2}}{{b}^{2}} = 1\left( {a > 0,b > 0}\right)$ 的离心率为 $e$.

(1)若 $e = \sqrt{2}$ ，且双曲线 $E$ 经过点 $\left( {\sqrt{2},1}\right)$ ，求双曲线 $E$ 的方程；

(2)若 $a = 2$ ，双曲线 $E$ 的左、右焦点分别为 ${F}_{1}$ 、 ${F}_{2}$ ，焦点到双曲线 $E$ 的渐近线的距离为 $\sqrt{3}$ ，点 $M$ 在第一象限且在双曲线 $E$ 上,若 $\left| {M{F}_{1}}\right|  = 8$ ,求 $\cos \angle {F}_{1}M{F}_{2}$ 的值;

(3)设圆 $O : {x}^{2} + {y}^{2} = 4,k\text{ 、 }m \in  \mathbf{R}$. 若动直线 $l : y = {kx} + m$ 与圆 $O$ 相切，且 $l$ 与双曲线 $E$ 交于 $A$ 、 $B$ 时,总有 $\angle {AOB} = \dfrac{\pi }{2}$ ,求双曲线 $E$ 离心率 $e$ 的取值范围.

35.(2024$\cdot$杨浦区一模$\cdot$20)已知双曲线 $\Gamma  : \dfrac{{x}^{2}}{3} - \dfrac{{y}^{2}}{12} = 1,A\left( {2,2}\right)$ 是双曲线 $\Gamma$ 上一点.

(1)若椭圆 $C$ 以双曲线 $\Gamma$ 的顶点为焦点，长轴长为 $4\sqrt{3}$ ，求椭圆 $C$ 的标准方程；

(2)设 $P$ 是第一象限中双曲线 $\Gamma$ 渐近线上一点， $Q$ 是双曲线 $\Gamma$ 上一点，且 $\vv{PA} = \vv{AQ}$ ，求 $\bigtriangleup  {POQ}$ 的面积 $S\left( {O\text{ 为坐标原点 }}\right)$ ；

(3)当直线 $l : y =  - {4x} + m$ (常数 $m \in  \mathbf{R}$ )与双曲线 $\Gamma$ 的左支交于 $M$ 、 $N$ 两点时，分别记直线 ${AM}\text{ 、 }{AN}$ 的斜率为 ${k}_{1}\text{ 、 }{k}_{2}$ ,求证: ${k}_{1} + {k}_{2}$ 为定值.

36.(2025$\cdot$普陀区一模$\cdot$20)设 $a > 0,m > 0,{F}_{1}$ 、 ${F}_{2}$ 分别是双曲线 $\Gamma  : \dfrac{{x}^{2}}{{a}^{2}} - {y}^{2} = 1$ 的左、右焦点,直线 $l : x - {my} - 2 = 0$ 经过点 ${F}_{2}$ 与 $\Gamma$ 的右支交于 $A\text{ 、 }B$ 两点,点 $O$ 是坐标原点.

(1)若点 $M$ 是 $\Gamma$ 上的一点, $\left| {M{F}_{1}}\right|  = 2$ ，求 $\left| {M{F}_{2}}\right|$ 的值；

(2)设 $\lambda$ 、 $\mu  \in  \mathbf{R}$ ，点 $P$ 在直线 $x = 6$ 上，若点 $O$ 、 $A$ 、 $P$ 、 $B$ 满足: $\vv{OA} = \lambda \vv{BP}$ ， $\vv{OB} = \mu \vv{AP}$ ，求点 $P$ 的坐标;

(3)设 ${AO}$ 的延长线与 $\Gamma$ 交于 $G$ 点,若向量 $\vv{OA}$ 与 $\vv{OB}$ 满足: $\vv{OA} \cdot  \vv{OB} \geq  {17}$ ，求 $\bigtriangleup  {GAB}$ 的面积 $S$ 的取值范围.

37.(2025$\cdot$闵行区二模$\cdot$20)已知双曲线 $\Gamma  : {x}^{2} - \dfrac{{y}^{2}}{3} = 1$ 的右焦点为 $F$ ，过点 $F$ 的直线 $l$ 交双曲线 $\Gamma$ 右支于 $A$ 、 $B$ 两点(点 $A$ 在 $x$ 轴上方)，点 $C$ 在双曲线 $\Gamma$ 上，直线 ${AC}$ 交 $x$ 轴于点 $Q$ (点 $Q$ 在点 $F$ 的右侧).

(1)求双曲线 $\Gamma$ 的渐近线方程；

(2)若点 $A\left( {2,3}\right)$ ，且 $\tan \angle {BAC} = \dfrac{1}{2}$ ，求点 $C$ 的坐标；

(3)若 $\bigtriangleup  {ABC}$ 的重心 $G$ 在 $x$ 轴上，记 $\bigtriangleup  {AFG}$ 、 $\bigtriangleup  {CQG}$ 的面积分别为 ${S}_{1}$ 、 ${S}_{2}$ ，求 $\dfrac{{S}_{1}}{{S}_{2}}$ 的最小值.

38.(2024$\cdot$浦东新区一模$\cdot$20)已知双曲线 $C : {x}^{2} - \dfrac{{y}^{2}}{3} = 1$ 的左、右焦点分别为 ${F}_{1}$ 、 ${F}_{2}$ ， $P$ 为双曲线右支上一点.

(1)求双曲线 $C$ 的离心率；

(2)设过点 $P$ 和 ${F}_{2}$ 的直线 $l$ 与双曲线 $C$ 的右支有另一交点为 $Q$ ，求 $\vv{OP} \cdot  \vv{OQ}$ 的取值范围；

(3)过点 $P$ 分别作双曲线 $C$ 两条渐近线的垂线,垂足分别为 $M$ 、 $N$ 两点,是否存在点 $P$ ，使得 $\left| {PM}\right|  + \left| {PN}\right|  = \sqrt{2}$ ? 若存在,求出点 $P$ 的坐标; 若不存在,请说明理由.

39.(2024$\cdot$普陀区一模$\cdot$20)设双曲线 $\Gamma  : \dfrac{{x}^{2}}{{t}^{2}} - {y}^{2} = 1\left( {t > 0}\right)$ ，点 ${F}_{1}$ 是 $\Gamma$ 的左焦点，点 $O$ 为坐标原点.

(1)若 $\Gamma$ 的离心率为 $\dfrac{\sqrt{10}}{3}$ ，求双曲线 $\Gamma$ 的焦距；

(2)过点 ${F}_{1}$ 且一个法向量为 $\vv{n} = \left( {t, - 1}\right)$ 的直线与 $\Gamma$ 的一条渐近线相交于点 $M$ ，若 ${S}_{\bigtriangleup {MO}{F}_{1}} = \dfrac{1}{2}$ ， 求双曲线 $\Gamma$ 的方程；

(3)若 $t = \sqrt{2}$ ，直线 $l : {kx} - y + m = 0\left( {k > 0,m \in  \mathbf{R}}\right)$ 与 $\Gamma$ 交于 $P$ 、 $Q$ 两点， $\left| {\vv{OP} + \vv{OQ}}\right|  = 4$ ，求直线 $l$ 的斜率 $k$ 的取值范围.

40. (2024$\cdot$宝山区二模$\cdot$20)已知双曲线 ${x}^{2} - \dfrac{{y}^{2}}{2} = 1$ 的左、右顶点分别为 $A$ 、 $B$ ，设点 $P$ 在第一象限且在双曲线上， $O$ 为坐标原点.

(1)求双曲线的两条渐近线夹角的余弦值；

\begin{center}
\includegraphics[width=0.2\textwidth]{bodyxb/src/117_1097_1276_299_299_0.jpg}
\end{center}


第 40 题图

(2)若 $\vv{PA} \cdot  \vv{PB} \leq  9$ ，求 $\left| \vv{OP}\right|$ 的取值范围；

(3)椭圆 $C$ 的长轴长为 $2\sqrt{2}$ ，且短轴的端点恰好是 $A$ 、 $B$ 两点，直线 ${AP}$ 与椭圆的另一个交点为 $Q$. 记 $\bigtriangleup {POA}$ 、 $\bigtriangleup {QAB}$ 的面积分别为 ${S}_{1}$ 、 ${S}_{2}$. 求 ${S}_{1}^{2} - {S}_{2}^{2}$ 的最小值，并写出取最小值时点 $P$ 的坐标. 41.(2024$\cdot$青浦区二模$\cdot$20)已知双曲线 $\Gamma  : \dfrac{{x}^{2}}{4} - \dfrac{{y}^{2}}{5} = 1,{F}_{1}\text{ 、 }{F}_{2}$ 分别为其左、右焦点.

\begin{center}
\includegraphics[width=0.2\textwidth]{bodyxb/src/118_1130_302_330_319_0.jpg}
\end{center}


第 41 题图

(1)求 ${F}_{1}$ 、 ${F}_{2}$ 的坐标和双曲线 $\Gamma$ 的渐近线方程；

(2)如图， $P$ 是双曲线 $\Gamma$ 右支在第一象限内一点，圆 $C$ 是 $\bigtriangleup  P{F}_{1}{F}_{2}$ 的内切圆,设圆与 $P{F}_{1}$ 、 $P{F}_{2}$ 、 ${F}_{1}{F}_{2}$ 分别切于点 $D$ 、 $E$ 、 $F$ ，当圆 $C$ 的面积为 ${4\pi }$ 时，求直线 $P{F}_{2}$ 的斜率；

(3)是否存在过点 ${F}_{2}$ 的直线 $l$ 与双曲线 $\Gamma$ 的左、右两支分别交于 $A\text{ 、 }B$ 两点,且使得 $\angle {F}_{1}{AB} = \angle {F}_{1}{BA}$ ? 若存在,求出直线 $l$ 的方程；若不存在，请说明理由.

42. (2025$\cdot$宝山区二模$\cdot$20)已知双曲线 $C : {x}^{2} - \dfrac{{y}^{2}}{3} = 1,{F}_{1}$ 、 ${F}_{2}$ 分别是其左、右焦点，直线 $l$ 与双曲线 $C$ 的右支交于 $A$ 、 $B$ 两点.

(1)当直线 $l$ 过点 ${F}_{2}$ ，且 $\left| {AB}\right|  = 6$ 时，求 $\bigtriangleup  {AB}{F}_{1}$ 的周长；

(2)已知点 $N\left( {-2,3}\right)$ ，若直线 ${AN}$ 、 ${BN}$ 的斜率之和为\blkda[2]{}。且: ${aQ} = {aH} > {AN}$ ，且两 $N$ 阶\blkda[2]{}为\blkda[2]{}，与 $y$ 轴交于点 $R$ 、 $S$ 时，求 $\bigtriangleup  {RSN}$ 的面积；

(3)已知直线 $l$ 过点 ${F}_{2},P$ 是双曲线 $C$ 上一点且位于第一象限，且满足 $\vv{OQ} = 2\vv{OP}$ 的点 $Q$ 在线段 ${AB}$ 上,若 $\vv{AB} = 2\vv{Q{F}_{2}}$ ,求点 $P$ 的坐标.

\begin{center}
\includegraphics[width=0.2\textwidth]{bodyxb/src/118_1177_1367_287_357_0.jpg}
\end{center}


第 42 题图

43.(2024$\cdot$上海秋季卷$\cdot$20)已知双曲线 $\Gamma  : {x}^{2} - \dfrac{{y}^{2}}{{b}^{2}} = 1\left( {b > 0}\right)$ 左右顶点分别为 ${A}_{1}$ 、 ${A}_{2}$ ，过点 $M\left( {-2,0}\right)$ 的直线 $l$ 交双曲线 $\Gamma$ 于 $P$ 、 $Q$ 两点.

(1)若离心率 $e = 2$ 时，求 $b$ 的值；

(2)若 $b = \dfrac{2\sqrt{6}}{3},{\bigtriangleup M}{A}_{2}P$ 为等腰三角形时，且点 $P$ 在第一象限，求点 $P$ 的坐标；

(3)连接 ${QO}$ 并延长,交双曲线 $\Gamma$ 于点 $R$ ,若 $\vv{{A}_{1}R} \cdot  \vv{{A}_{2}P} = 1$ ，求 $b$ 的取值范围.

\subsection*{考点5 抛物线方程及其性质}

\subsubsection*{一、填空题}

1.(2024$\cdot$宝山区二模$\cdot$1)抛物线 ${x}^{2} = {4y}$ 的焦点坐标为\blkda[2]{}.

2.(2024$\cdot$静安区一模$\cdot$1)准线方程为 $x + 1 = 0$ 的抛物线标准方程为\blkda[2]{}.

3.(2024$\cdot$普陀区一模$\cdot$1)若抛物线 ${x}^{2} = {my}$ 的顶点到它的准线距离为 $\dfrac{1}{2}$ ，则正实数 $m =$ \blkda[2]{}.

4.(2025$\cdot$普陀区一模$\cdot$2)已知某抛物线的准线方程为 $y = 1$ ，则该抛物线的标准方程为\blkda[2]{}.

5.(2024$\cdot$杨浦区二模$\cdot$2)抛物线 ${y}^{2} = {8x}$ 的准线方程为\blkda[2]{}.

6.(2024$\cdot$嘉定区二模$\cdot$2)抛物线 ${y}^{2} = {4x}$ 的准线方程为\blkda[2]{}.

7.(2024$\cdot$黄浦区二模$\cdot$2)抛物线 ${y}^{2} = {4x}$ 的焦点到准线的距离是\blkda[2]{}.

8.(2024$\cdot$虹口区二模$\cdot$3)过抛物线 ${y}^{2} = {4x}$ 焦点的弦 ${AB}$ 的中点横坐标为 2，则弦 ${AB}$ 的长度为\blkda[2]{}.

9.(2025$\cdot$黄浦区二模$\cdot$3)抛物线 ${y}^{2} = x$ 的焦点到其顶点的距离为\blkda[2]{}.

10.(2025$\cdot$宝山区二模$\cdot$3)抛物线 ${x}^{2} = {4y}$ 的准线方程是\blkda[2]{}.

11.(2024$\cdot$奉贤区二模$\cdot$3)抛物线 ${y}^{2} = {4x}$ 上一点到点(1,0)的距离最小值为\blkda[2]{}.

12.(2025$\cdot$普陀区二模$\cdot$5)设 $m > 2$ ，抛物线 $C : {x}^{2} = {2my}$ 上的点 $P$ 到 $C$ 的焦点的距离为 5，点 $P$ 到 $y$ 轴的距离为 3，则 $m$ 的值为\blkda[2]{}.

13.(2025$\cdot$奉贤区一模$\cdot$7)已知抛物线 ${x}^{2} = {ay}\left( {a > 0}\right)$ 上有一点 $P$ 到准线的距离为 6，点 $P$ 到 $x$ 轴的距离为 4 ，则抛物线的焦点坐标为\blkda[2]{}.

14.(2024$\cdot$上海秋季卷$\cdot$7)已知抛物线 ${y}^{2} = {4x}$ 上有一点 $P$ 到准线的距离为 9，那么点 $P$ 到 $x$ 轴的距离为\blkda[2]{}.

15.(2025$\cdot$浦东新区二模$\cdot$8)设 $M\left( {x,y}\right)$ 为抛物线 ${y}^{2} = {4x}$ 上任意一点，若 $x + {2y} + m$ 的最小值为 1,则 $m$ 的值为\blkda[2]{}.

16.(2024$\cdot$普陀区二模$\cdot$10)已知抛物线 ${y}^{2} = 4\sqrt{3}x$ 的焦点 $F$ 是双曲线 $\Gamma$ 的右焦点，过点 $F$ 的直线 $l$ 的法向量 $\vv{n} = \left( {1, - \sqrt{3}}\right) ,l$ 与 $y$ 轴以及 $\Gamma$ 的左支分别相交 $A$ 、 $B$ 两点,若 $\vv{BF} = 2\vv{BA}$ ,则双曲线 $\Gamma$ 的实轴长为\blkda[2]{}.

17.(2025$\cdot$青浦区二模$\cdot$10)已知点 $P$ 是抛物线 ${y}^{2} = {8x}$ 上一动点，点 $Q$ 在圆 ${\left( x - 5\right) }^{2} + {y}^{2} = 1$ 上运动，则 $P$ 与 $Q$ 两点间最短距离为\blkda[2]{}.

18.(2024$\cdot$长宁区二模$\cdot$10)已知抛物线 $\Gamma  : {y}^{2} = {4x}$ 的焦点为 $F$ ，准线为 $l$ ，点 $M$ 在 $\Gamma$ 上， ${MN} \bot  l$ ， $\angle {NFM} = {30}^{ \circ  }$ ，则点 $M$ 的横坐标为\blkda[2]{}.

19.(2024$\cdot$杨浦区$\cdot$一模$\cdot$11)已知抛物线 ${y}^{2} = {2px}$ ( $p > 0$ )的焦点为 $F$ ，第一象限的 $A$ 、 $B$ 两点在抛物线上,且满足 $\left| {BF}\right|  - \left| {AF}\right|  = 4,\left| {AB}\right|  = 4\sqrt{2}$. 若线段 ${AB}$ 中点的纵坐标为 4,则抛物线的方程为\blkda[2]{}.

20.(2024$\cdot$浦东新区 一模$\cdot$11)已知曲线 ${C}_{1} : {y}^{2} = {4x}\left( {x \leq  1}\right)$ ，曲线 ${C}_{2} : {x}^{2} - {2x} + {y}^{2} - {28} = 0\left( {x \geq  6}\right)$ ， 若 $\bigtriangleup {ABC}$ 顶点 $A$ 的坐标为 $\left( {1,0}\right)$ ，顶点 $B$ 、 $C$ 分别在曲线 ${C}_{1}$ 和 ${C}_{2}$ 上运动，则 $\bigtriangleup {ABC}$ 周长的最小值为\blkda[2]{}.

\subsubsection*{二、选择题}

21.(2024$\cdot$金山区二模$\cdot$13)若抛物线 ${y}^{2} = {2px}\left( {p > 0}\right)$ 的焦点是椭圆 $\dfrac{{x}^{2}}{16} + \dfrac{{y}^{2}}{4} = 1$ 的一个顶点，则 $p$ 的值为 \dotfill{(\kh{})}

\xx{}{}{}{}2 B. 3

C. 4 D. 8

22.(2025$\cdot$徐汇区一模$\cdot$14)下列抛物线中,焦点坐标为 $\left( {0,\dfrac{1}{8}}\right)$ 的是 \dotfill{(\kh{})}

\xx{}{}{}{}${y}^{2} = \dfrac{1}{2}x$ B. ${y}^{2} = \dfrac{1}{4}x$

C. ${x}^{2} = \dfrac{1}{2}y$ D. ${x}^{2} = \dfrac{1}{4}y$

23.(2024$\cdot$青浦区二模$\cdot$14)已知点 $P\left( {2,2\sqrt{2}}\right)$ 是抛物线 $C : {y}^{2} = {2px}\left( {p > 0}\right)$ 上一点，到抛物线 $C$ 的准线的距离为 $d,M$ 是 $x$ 轴上一点，则 “点 $M$ 的坐标为 $\left( {1,0}\right)$ ”是“$d = \left| {PM}\right|$”的 \dotfill{(\kh{})}

\xx{}{}{}{}充分不必要条件

B. 必要不充分条件

C. 充要条件

D. 既不充分也不必要条件

\subsubsection*{三、解答题}

24.(2025$\cdot$金山区一模$\cdot$20)已知椭圆 ${\Gamma }_{1} : \dfrac{{x}^{2}}{4} + \dfrac{{y}^{2}}{3} = 1$ ，抛物线 ${\Gamma }_{2} : {y}^{2} = {2px}\left( {p > 0}\right)$ 与 ${\Gamma }_{1}$ 有一个相同的焦点 $F$. 过点 $F$ 作互相垂直的两条直线 $l$ 与 ${l}^{\prime }$ ，直线 $l$ 与 ${\Gamma }_{1}$ 交于点 $A$ 、 $B$ ，直线 ${l}^{\prime }$ 与 ${\Gamma }_{2}$ 交于点 $C\text{ 、 }D$.

(1)求椭圆 ${\Gamma }_{1}$ 的离心率及抛物线 ${\Gamma }_{2}$ 的方程；

(2)若直线 $l$ 的倾斜角为 $\dfrac{3\pi }{4}$ ，求 ${AB}$ 中点 $M$ 的坐标；

(3)四边形 ${ACBD}$ 的面积是否存在最小值？若存在，求出最小值；若不存在，请说明理由.

25.(2025$\cdot$崇明区二模$\cdot$20)已知抛物线 $\Gamma  : {x}^{2} = {4y}$ ，过点 $P\left( {a,b}\right)$ 的直线 $l$ 与抛物线 $\Gamma$ 交于点 $A$ 、 $B$ ， 与 $y$ 轴交于点 $C$.

(1)若点 $A$ 位于第一象限，且点 $A$ 到抛物线 $\Gamma$ 的焦点的距离等于 3，求点 $A$ 的坐标；

(2)若点 $A$ 坐标为 $\left( {4,4}\right)$ ，且点 $B$ 恰为线段 ${AC}$ 的中点，求原点 $O$ 到直线 $l$ 的距离；

\begin{center}
\includegraphics[width=0.2\textwidth]{bodyxb/src/121_1076_1392_338_287_0.jpg}
\end{center}


第 25 题图

(3)若抛物线 $\Gamma$ 上存在定点 $D$ 使得满足题意的点 $A\text{ 、 }B$ 都有 ${DA} \bot$

${DB}$ ,求 $a\text{ 、 }b$ 满足的关系式.

26.(2024$\cdot$虹口区一模$\cdot$20)已知点 $M\left( {m,4}\right)$ 在抛物线 $\Gamma  : {x}^{2} = {2py}\left( {p > 0}\right)$ 上，点 $F$ 为 $\Gamma$ 的焦点，且 $\left| {MF}\right|  = 5$ ,过点 $F$ 的直线 $l$ 与 $\Gamma$ 及圆 ${x}^{2} + {\left( y - 1\right) }^{2} = 1$ 依次相交于点 $A\text{ 、 }B\text{ 、 }C\text{ 、 }D$ ,如图所示.

(1)求抛物线 $\Gamma$ 的方程及点 $M$ 的坐标；

\begin{center}
\includegraphics[width=0.3\textwidth]{bodyxb/src/122_1016_299_391_308_0.jpg}
\end{center}


第 26 题图

(2)求证: $\left| {AC}\right|$. $\left| {BD}\right|$ 为定值；

(3)过 $A$ 、 $B$ 两点分别作 $\Gamma$ 的切线 ${l}_{1}$ 、 ${l}_{2}$ ，且 ${l}_{1}$ 与 ${l}_{2}$ 相交于点 $P$ ,求 $\bigtriangleup {ACP}$ 与 $\bigtriangleup {BDP}$ 的面积之和的最小值.

27.(2024$\cdot$嘉定区一模$\cdot$20)抛物线 ${y}^{2} = {4x}$ 上有一动点 $P\left( {s,t}\right)$ ， $t > 0$. 过点 $P$ 作抛物线的切线 $l$ ，再过点 $P$ 作直线 $m$ ,使得 $m \bot  l$ ,直线 $m$ 和抛物线的另一个交点为 $Q$.

(1)当 $s = 1$ 时,求切线 $l$ 的直线方程;

(2)当直线 $l$ 与抛物线准线的交点在 $x$ 轴上时,求三角形 ${OPQ}$ 的面积(点 $O$ 是坐标原点);

(3)求出线段 $\left| {PQ}\right|$ 关于 $s$ 的表达式，并求 $\left| {PQ}\right|$ 的最小值.

28.(2024$\cdot$闵行区一模$\cdot$20)已知 $0 < p < 4$ ，曲线 ${\Gamma }_{1}$ 、 ${\Gamma }_{2}$ 的方程分别为 ${y}^{2} = {2px}\left( {0 \leq  x \leq  8,y \geq  0}\right)$ 和 ${x}^{2} = {2py}\left( {0 \leq  y \leq  8,x \geq  0}\right) ,{\Gamma }_{1}$ 与 ${\Gamma }_{2}$ 在第一象限内相交于点 $K\left( {{x}_{K},{y}_{K}}\right)$.

(1)若 $\left| {OK}\right|  = 4\sqrt{2}$ ，求 $p$ 的值；

(2)若 $p = 2$ ，定点 $T$ 的坐标为 $\left( {4,0}\right)$ ，动点 $M$ 在直线 $y = x$ 上，动点 $N\left( {{x}_{N},{y}_{N}}\right)$ ( $0 \leq  {x}_{N} \leq  4$ )在曲线 ${\Gamma }_{2}$ 上,求 $\left| {MN}\right|  + \left| {MT}\right|$ 的最小值;

(3)已知点 $A\left( {{x}_{1},{y}_{1}}\right) \left( {0 \leq  {x}_{1} \leq  {x}_{K}}\right)$ 、 $B\left( {{x}_{2},{y}_{2}}\right) \left( {{x}_{K} < {x}_{2} \leq  8}\right)$ 在曲线 ${\Gamma }_{1}$ 上,点 $A\text{ 、 }B$ 关于直线 $y = x$ 的对称点分别为 $C\text{ 、 }D$ ,设 $\left| {AC}\right|$ 的最大值为 $m,\left| {BD}\right|$ 的最大值为 $t$ ,若 $\dfrac{m}{t} \in  \left\lbrack  {\dfrac{1}{2},2}\right\rbrack$ ,求实数 $p$ 的取值范围.

\begin{center}
\includegraphics[width=0.2\textwidth]{bodyxb/src/123_1120_467_285_292_0.jpg}
\end{center}


第 28 题图

29. (2025$\cdot$徐汇区二模$\cdot$20)已知抛物线 $C : {y}^{2} = {4x}$ ，点 $F$ 是抛物线 $C$ 的焦点.

(1)求点 $F$ 的坐标及点 $F$ 到准线 $l$ 的距离；

(2)过点 $F$ 作相互垂直的两条直线 ${l}_{1}\text{ 、 }{l}_{2}\text{ ， }{l}_{1}$ 交抛物线 $C$ 于点 ${P}_{1}\text{ 、 }{P}_{2}\text{ ， }{l}_{2}$ 交抛物线 $C$ 于点 ${Q}_{1}$ 、 ${Q}_{2}$ ,求证: $\dfrac{1}{\left| {P}_{1}{P}_{2}\right| } + \dfrac{1}{\left| {Q}_{1}{Q}_{2}\right| }$ 为定值,并求出该定值;

(3)过点 $F$ 且斜率为 $\sqrt{3}$ 的直线交抛物线 $C$ 于 $A$ 、 $B$ 两点，设点 $P$ 不在直线 ${AB}$ 上且 ${PF}$ 为 $\bigtriangleup {PAB}$ 的内角平分线,求 $\bigtriangleup {PAB}$ 面积的最大值.

30.(2024$\cdot$崇明区一模$\cdot$20)已知抛物线 ${\Gamma }_{1} : {y}^{2} = {4x},{\Gamma }_{2} : {y}^{2} = {2x}$ ，直线 $l$ 交抛物线，\blkda[2]{}， 抛物线 ${\Gamma }_{2}$ 于点 $B$ 、 $C$ ，其中点 $A$ 、 $B$ 位于第一象限.

(1)若点 $A$ 到抛物线 ${\Gamma }_{1}$ 焦点的距离为 2，求点 $A$ 的坐标；

(2)若点 $A$ 的坐标为 $\left( {4,4}\right)$ ，且线段 ${AC}$ 的中点在 $x$ 轴上，求原点 $O$ 到直线 $l$ 的距离；

(3)若 $\vv{AB} = 2\vv{CD}$ ，求 $\bigtriangleup  {AOD}$ 与 $\bigtriangleup  {BOC}$ 的面积之比.

31. (2023$\cdot$上海秋季卷$\cdot$20)已知抛物线 $\Gamma  : {y}^{2} = {4x},A$ 为第一象限内 $\Gamma$ 上的一点，设 $A$ 的纵坐标为 $a\left( {a > 0}\right)$.

(1)若点 $A$ 到 $\Gamma$ 的准线的距离为 3，求 $a$ 的值；

(2)若 $a = 4,B$ 为 $x$ 轴上的一点，且线段 ${AB}$ 的中点在 $\Gamma$ 上，求点 $B$ 的坐标及原点 $O$ 到直线 ${AB}$ 的距离;

(3)设直线 $l : x =  - 3,P$ 是第一象限内， $\Gamma$ 上异于 $A$ 的动点，直线 ${AP}$ 与直线 $l$ 交于点 $Q$ ，点 $H$ 为点 $P$ 在直线 $l$ 上的投影，若点 $A$ 满足性质 “当点 $P$ 变化时， $\left| {HQ}\right|  > 4$ 恒成立”，求 $a$ 的取值范围.

32. (2024$\cdot$金山区一模$\cdot$20)已知三条直线 ${l}_{i} : y = {kx} + {m}_{i}\left( {i = 1,2,3}\right)$ 分别与抛物线 $\Gamma  : {y}^{2} = {8x}$ 交于点 ${A}_{i}\text{ 、 }{B}_{i},T\left( {t,0}\right)$ 为 $x$ 轴上一定点,且 ${m}_{1} < {m}_{2} < {m}_{3} <  - t$ ,记点 $T$ 到直线 ${l}_{i}$ 的距离为 ${d}_{i}$ , $\bigtriangleup T{A}_{i}{B}_{i}$ 的面积为 ${S}_{i}$.

(1)若直线 ${l}_{3}$ 的倾斜角为 ${45}^{ \circ  }$ ，且过抛物线 $\Gamma$ 的焦点 $F$ ，求直线 ${l}_{3}$ 的方程；

(2)若 $\vv{O{A}_{1}} \cdot  \vv{O{B}_{1}} = 0$ ，且 $k{m}_{1} \neq  0$ ，证明:直线 ${l}_{1}$ 过定点；

(3)当 $k = 1$ 时,是否存在点 $T$ ，使得 ${S}_{1}$ ， ${S}_{2}$ ， ${S}_{3}$ 成等比数列， ${d}_{1}$ ， ${d}_{2}$ ， ${d}_{3}$ 也成等比数列? 若存在， 请求出点 $T$ 的坐标；若不存在，请说明理由.

33.(2025$\cdot$杨浦区一模$\cdot$20)如图所示，已知抛物线 $\Gamma  : {y}^{2} = x$ ，点 $A$ 、 $B$ 、 $C$ 、 $D$ 是抛物线上的四个点， 其中 $A$ 、 $D$ 在第一象限， $B$ 、 $C$ 在第四象限，满足 ${AB}//{CD}$ ，线段 ${AC}$ 与 ${BD}$ 交于点 $H$. 记线段 ${AB}$ 与 ${CD}$ 的中点分别为 $M\text{ 、 }N$.

\begin{center}
\includegraphics[width=0.2\textwidth]{bodyxb/src/125_1128_1279_332_309_0.jpg}
\end{center}


第 33 题图

(1)求抛物线 $\Gamma$ 的焦点坐标；

(2)求证:点 $M$ 、 $H$ 、 $N$ 三点共线；

(3)若 $2\left| {HM}\right|  = \left| {HN}\right|  = 2$ ，求四边形 ${ABCD}$ 的面积.

\subsection*{考点6 直线与圆锥曲线的位置关系及其应用}

\subsubsection*{一、填空题}

1. (2025$\cdot$宝山区一模$\cdot$9)过双曲线 $\dfrac{{x}^{2}}{9} - \dfrac{{y}^{2}}{16} = 1$ 的左焦点 $F$ 作圆 ${x}^{2} + {y}^{2} = 9$ 的切线，切点为 $M$. 延长切线交双曲线的右支于点 $P$ ， $O$ 为坐标原点，点 $T$ 为线段 ${FP}$ 的中点，则 $\left| {OT}\right|  =$ \blkda[2]{}.

2.(2024$\cdot$崇明区一模$\cdot$12)已知正实数 $a$ 、 $b$ 、 $c$ 、 $d$ 满足 ${a}^{2} - {ab} + 1 = 0$ ， ${c}^{2} + {d}^{2} = 1$ ，则当 ${\left( a - c\right) }^{2} +  {\left( b - d\right) }^{2}$ 取得最小值时， ${ab} =$ \blkda[2]{}.

\subsubsection*{二、选择题}

3.(2024$\cdot$静安区一模$\cdot$16)记点 $P$ 到图形 $C$ 上每一个点的距离的最小值称为点 $P$ 到图形 $C$ 的距离， 那么平面内到定圆 $C$ 的距离与到定点 $A$ 的距离相等的点的轨迹不可能是 \dotfill{(\kh{})}

\xx{}{}{}{}圆 B. 椭圆

C. 双曲线的一支 D. 直线

\subsubsection*{三、解答题}

4.(2023$\cdot$上海春季卷$\cdot$20)椭圆 $\Gamma  : \dfrac{{x}^{2}}{{m}^{2}} + \dfrac{{y}^{2}}{3} = 1\left( {m > 0,m \neq  \sqrt{3}}\right)$.

(1)若 $m = 2$ ，求椭圆 $\Gamma$ 的离心率；

(2)设 ${A}_{1}$ 、 ${A}_{2}$ 为椭圆 $\Gamma$ 的左右顶点，椭圆 $\Gamma$ 上一点 $E$ 的纵坐标为 1，且 $\vv{E{A}_{1}} \cdot  \vv{E{A}_{2}} =  - 2$ ，求 $m$ 的值；

(3)过椭圆 $\Gamma$ 上一点 $P$ 作斜率为 $\sqrt{3}$ 的直线，与双曲线 $\dfrac{{y}^{2}}{5{m}^{2}} - \dfrac{{x}^{2}}{5} = 1$ 有一个公共点,求 $m$ 的取值范围.

\begin{center}
\includegraphics[width=0.2\textwidth]{bodyxb/src/126_1100_1331_329_309_0.jpg}
\end{center}


第 4 题图

5.(2025$\cdot$崇明区一模$\cdot$20)已知椭圆 $\Gamma  : \dfrac{{y}^{2}}{4} + \dfrac{{x}^{2}}{3} = 1$ ，点 ${F}_{1}$ 、 ${F}_{2}$ 分别是椭圆的下焦点和上焦点，过点 ${F}_{2}$ 的直线 $l$ 与椭圆交于 $A$ 、 $B$ 两点.

(1)若直线 $l$ 平行于 $x$ 轴,求线段 ${AB}$ 的长；

(2)若点 $A$ 在 $y$ 轴左侧，且 $\vv{{F}_{1}A} \cdot  \vv{{F}_{2}A} = \dfrac{9}{4}$ ，求直线 $l$ 的方程；

(3)已知椭圆上的点 $C$ 满足 $\left| {CA}\right|  = \left| {CB}\right|$ ，是否存在直线 $l$ 使得 $\bigtriangleup {ABC}$ 的重心在 $x$ 轴上？若存在,请求出直线 $l$ 的方程; 若不存在,请说明理由.

6.(2024$\cdot$虹口区二模$\cdot$20)已知椭圆 $\Gamma  : \dfrac{{x}^{2}}{{a}^{2}} + \dfrac{{y}^{2}}{{b}^{2}} = 1\left( {a > b > 0}\right)$ 的焦距为 $2\sqrt{3}$ ，点 $P\left( {0,1}\right)$ 在椭圆 $\Gamma$ 上，动直线 $l$ 与椭圆 $\Gamma$ 相交于不同的两点 $A$ 、 $B$ ，且直线 ${PA}$ 、 ${PB}$ 的斜率之积为 1.

(1)求椭圆 $\Gamma$ 的标准方程；

(2)若直线 ${PA}$ 的法向量为 $\vv{n} = \left( {1, - 2}\right)$ ，求直线 $l$ 的方程；

(3)是否存在直线 $l$ ，使得 $\bigtriangleup  {PAB}$ 为直角三角形？若存在，求出直线 $l$ 的斜率；若不存在，请说明理由.

7.(2025$\cdot$闵行区一模$\cdot$20)已知圆 $O : {x}^{2} + {y}^{2} = 1$ ，双曲线 $\Gamma  : {x}^{2} - \dfrac{{y}^{2}}{{b}^{2}} = 1$ ，直线 $l : y = {kx} + b$ ，其中 $k \in  \mathbf{R},b > 0.$

(1)当 $b = 2$ 时,求双曲线 $\Gamma$ 的离心率;

(2)若 $l$ 与圆 $O$ 相切，证明: $l$ 与双曲线 $\Gamma$ 的左右两支各有一个公共点；

(3)设 $l$ 与 $y$ 轴交于点 $P$ ，与圆 $O$ 交于点 $A$ 、 $B$ ，与双曲线 $\Gamma$ 的左右两支分别交于点 $C$ 、 $D$ ，四个点从左至右依次为 $C$ 、 $A$ 、 $B$ 、 $D$. 当 $k = \dfrac{\sqrt{2}}{2}$ 时，是否存在实数 $b$ ，使得 $\vv{PA} \cdot  \vv{PC} = \vv{PB} \cdot  \vv{PD}$ 成立？ 若存在,求出 $b$ 的值; 若不存在,说明理由.

8.(2025$\cdot$虹口区一模$\cdot$20)已知椭圆 $\Gamma  : \dfrac{{x}^{2}}{4} + {y}^{2} = 1$ 的左、右焦点分别为 ${F}_{1}$ 、 ${F}_{2}$ ，右顶点为 $A$ ，上顶点为 $B$ ,设 $P$ 为 $\Gamma$ 上的一点.

(1)当 $P{F}_{1} \bot  {F}_{1}{F}_{2}$ 时,求 $\left| {P{F}_{2}}\right|$ 的值;

( 2 )若点 $P$ 坐标为 $\left( {1,\dfrac{\sqrt{3}}{2}}\right)$ ，则在 $\Gamma$ 上是否存在点 $Q$ ，使 $\bigtriangleup  {APQ}$ 的面积为 $\dfrac{\sqrt{3} + 1}{2}$ ？若存在，请求出所有满足条件的点 $Q$ 的坐标；若不存在，请说明理由；

(3)已知点 $D$ 坐标为 $\left( {0,m}\right)$ ，过点 $P$ 和点 $D$ 的直线 $l$ 与椭圆 $\Gamma$ 交于另一点 $T$ ，当直线 $l$ 与 $x$ 轴和 $y$ 轴均不平行时,有 $\vv{PT} \cdot  \left( {\vv{BP} + \vv{BT}}\right)  = 0$ ,求实数 $m$ 的取值范围.



9.(2025$\cdot$奉贤区二模$\cdot$20)如图1所示，曲线 $\Gamma$ 是 $\dfrac{{x}^{2}}{4} + \dfrac{{y}^{2}}{{b}^{2}} = 1\left( {y \geq  0}\right)$ 与 $\dfrac{{y}^{2}}{{b}^{2}} - \dfrac{{x}^{2}}{4} = 1\left( {y \geq  0}\right)$ 组合的. (1) $\dfrac{{x}^{2}}{4} + \dfrac{{y}^{2}}{{b}^{2}} = 1\left( {y \geq  0}\right)$ 过点 $\left( {-1,\dfrac{3}{2}}\right)$ ，求 $\dfrac{{y}^{2}}{{b}^{2}} - \dfrac{{x}^{2}}{4} = 1\left( {y \geq  0}\right)$ 的渐近线方程；

(2) $b = \sqrt{2}$ ，设 $T\left( {0,t}\right) ,t \in  \left\lbrack  {-\dfrac{\sqrt{2}}{2},3\sqrt{2}}\right\rbrack$ ，曲线 $\Gamma$ 上找一个点 $Q$ ，使得 $\left| \vv{TQ}\right|$ 达到最小；

(3)若 $b = 1$ ，如图 2 所示，存在过点 $P\left( {0,1}\right)$ 的两条直线 ${l}_{1}$ 、 ${l}_{2}$ 与曲线 $\Gamma$ 的交点分别是 $A$ 、 $M$ 、 $B$ 、 $N$ ,点 $A$ 在第二象限,点 $M$ 在第一象限. 是否存在非零实数 $\lambda$ ,使得 $\vv{MN} = \lambda \vv{AB}$ 成立? 请说明理由.

\begin{center}
\includegraphics[width=0.6\textwidth]{bodyxb/src/129_715_646_711_221_0.jpg}
\end{center}


第 9 题图

10.(2025$\cdot$静安区二模$\cdot$20)如图所示，在直角坐标平面 ${xOy}$ 中， $\bigtriangleup {A}_{i}{B}_{i}{A}_{i + 1}$ 中( $i = 1,2,\cdots ,n,\cdots$ )为正三角形,且满足 $\vv{O{A}_{1}} = \left( {-\dfrac{1}{4},0}\right) ,\vv{{A}_{i}{A}_{i + 1}} = \left( {{2i} - 1,0}\right)$.

(1)求点 ${A}_{n}$ 的横坐标 ${X}_{n}$ 关于正整数 $n$ 的表达式；

(2)求证:点 ${B}_{1},{B}_{2},\cdots ,{B}_{n}\cdots$ 在抛物线 $\Gamma  : {y}^{2} = {3x}$ 上；

(3)过(2)中抛物线 $\Gamma  : {y}^{2} = {3x}$ 的焦点 $F$ 作两条互相垂直的弦 ${AC}$ 和 ${BD}$ ,求四边形 ${ABCD}$ 面积的最小值.

\begin{center}
\includegraphics[width=0.2\textwidth]{bodyxb/src/129_1097_1445_331_230_0.jpg}
\end{center}


第 10 题图

11.(2025$\cdot$杨浦区二模$\cdot$20)已知双曲线 $\Gamma$ 的标准方程为 ${x}^{2} - \dfrac{{y}^{2}}{2} = 1$ ，点 $P$ 是双曲线， 个动点.

(1)求双曲线 $\Gamma$ 的焦点坐标和渐近线方程；

(2)过点 $P$ 分别向两条渐近线作垂线，垂足为点 ${P}_{1}\text{ 、 }{P}_{2}$ ，求 $\vv{P{P}_{1}}\text{ 、 }\vv{P{P}_{2}}$ 的值；

(3)若 $\left| {OP}\right|  > \sqrt{2}$ ，如图所示，过 $P$ 作圆 $O : {x}^{2} + {y}^{2} = 2$ 的切线 $l$ ，切点为 $M$ ,交双曲线 $\Gamma$ 的左支于点 $Q$ ,分别交两条渐近线于点 $A$ 、 B. 设 $\left| {PQ}\right|  = \lambda \left| {AB}\right|$ ,求实数 $\lambda$ 的取值范围.

\begin{center}
\includegraphics[width=0.2\textwidth]{bodyxb/src/130_1144_442_331_371_0.jpg}
\end{center}


第 11 题图

\section*{第 8 章 平 面 向 量}

\begin{center}

\begin{tabular}{|p{4cm}|p{12cm}|}
\hline
考点 & 课标要求与命题趋势 \\
\cline{1-2}
1. 向量的概念及其线性运算 & \ding{172}  通过对力、速度、位移等的分析，了解平面向量的实际背景，理解平面向量的意义和两个向量相等的含义; 理解平面向量的几何表示和基本要素 \ding{173} 借助实例和平面向量的几何表示，掌握平面向量加、减运算及运算规则，理解其几何意义 \ding{174}  掌握平面向量的基本概念、线性运算及坐标运算，已知平面向量的关系，会求参数 \\
\cline{1-2}
2. 向量的基本定理及其应用 & \ding{172}  理解平面向量基本定理及其意义; 借助平面直角坐标系, 掌握平面向量的正交分解及坐标表示 \ding{173}  掌握基本定理的基底表示向量、能在平面几何图形中的应用 \\
\cline{1-2}
3. 向量的数量积与向量的坐标表示 & \ding{172}  通过实例分析，掌握平面向量数乘运算及运算规则，理解其几何意义；理解两个平面向量共线的含义;了解平面向量的线性运算性质及其几何意义 \ding{173}  通过物理中功等实例，理解平面向量数量积的概念及其物理意义，会计算平面向量的数量积;会用数量积判断两个平面向量的垂直关系 \ding{174}  通过几何直观，了解平面向量投影的概念以及投影向量的意义; 会求平面几何图形中的范围及最值等问题 \\
\cline{1-2}
4. 向量的夹角与向量的坐标表示 & \ding{172} 会用向量方法解决简单的平面几何问题、力学问题以及其他实际问题，体会向量在解决数学和实际问题中的作用 \ding{173} 借助向量的运算，探索三角形边长与角度的关系，掌握余弦定理、正弦定理；能用余弦定理、正弦定理解决简单的实际问题 \ding{174}  会用坐标表示平面向量的加、减运算与数乘运算; 能用坐标表示平面向量的数量积, 会表示两个平面向量的夹角; 能用坐标表示平面向量共线、垂直的条件 \\
\cline{1-2}
\hline
\end{tabular}
\end{center}
\clearpage
\subsection*{考点1 平面量的概念及其线性运算}

\subsubsection*{一、填空题}
\begin{enumerate}
    \item (2023$\cdot$上海春季卷$\cdot$2)已知向量 $\vv{a} = \left( {3,4}\right) ,\vv{b} = \left( {1,2}\right)$ ，则 $\vv{a} - 2\vv{b} =$ \blkda[2]{$\left( {1,0}\right)$}.

    \item (2024$\cdot$浦东新区一模$\cdot$4)已知直线 $l$ 的倾斜角为 $\dfrac{\pi }{3}$ ，请写出直线 $l$ 的一个法向量\blkda[2]{$\left( {\sqrt{3}, - 1}\right)$}.
    \begin{jieda}
        $\alpha  = \dfrac{\pi }{3},k = \sqrt{3},\vv{d} = \left( {1,\sqrt{3}}\right)  \Longrightarrow  \vv{n} = \left( {\sqrt{3}, - 1}\right)$.
    \end{jieda}

    \item (2025$\cdot$闵行区二模$\cdot$6)已知向量 $\vv{a} = \left( {3,4}\right) ,\vv{b} = \left( {\cos \theta ,\sin \theta }\right)$ ，若 $\vv{a}//\vv{b}$ ，则 $\tan \theta  =$ \blkda[2]{$\dfrac{4}{3}$}.
    \begin{jieda}
        因为 $\vv{a}//\vv{b}$ ,且 $\vv{a} = \left( {3,4}\right) ,\vv{b} = \left( {\cos \theta ,\sin \theta }\right)$ ,所以 $4\cos \theta  = 3\sin \theta$，所以 $\tan \theta  = \dfrac{4}{3}$
    \end{jieda}

    \item (2025$\cdot$上海春季卷$\cdot$12)在平面中， ${\vv{e}}_{1}$ 和 ${\vv{e}}_{2}$ 是互相垂直的单位向量，向量 $\vv{a}$ 满足 $\left|\vv{a} - 4\vv{e_1}\right|  = 2$ ，向量 $\vv{b}$ 满足 $\left|\vv{b} - 6\vv{e_2}\right|  = 1$ ，则 $\vv{b}$ 在 $\vv{a}$ 方向上的数量投影的最大值\blkda[2]{4}.
    \begin{jieda}
        设 $\vv{a} = \vv{OA},\vv{b} = \vv{OB}$ ,在坐标系中,不妨设 $4{\vv{e}}_{1} = \left( {0,4}\right)  = \vv{O{O}_{1}},6{\vv{e}}_{2} = \left( {6,0}\right)  = \vv{O{O}_{2}}$ ,则 $\left|\vv{OA} -  \vv{O{O}_{1}}\right|  = \left| \vv{{O}_{1}A}\right| = 2, \left|\vv{OB} - \vv{O{O}_{2}}\right| = \left|\vv{{O}_{2}B} \right| = 1$
        
        如图,显然 ${OA}$ 与圆 ${O}_{1}$ 相切在右侧时, $\vv{b}$ 在 $\vv{a}$ 方向上的数量投影最大, 当 ${BH} \bot  {OA}$ ,且 ${BH}$ 与圆 ${O}_{2}$ 相切时, $\vv{b}$ 在 $\vv{a}$ 方向上的数量投影 $\left| {OH}\right|$ 最大,由题意,得 $\angle {O}_{1}{OA} = {30}^{ \circ  } = \angle {HTO} \Longrightarrow  \left| {{O}_{2}T}\right|  = 2$ ,因为 $\bigtriangleup {HOT} \backsim  \bigtriangleup B{O}_{2}T$ ,所以 $\left| {OH}\right|  = 4$，故 $\vv{b}$ 在 $\vv{a}$ 方向上的数量投影的最大值为4.
        \begin{center}
            \includegraphics[width=0.4\textwidth]{bodyxb/src/95_516_1093_304_208_0.jpg}
        \end{center}
    \end{jieda}

    \item (2025$\cdot$奉贤区二模$\cdot$12) $\bigtriangleup  {ABC}$ 内一点 $F$ (如图 1)，式子 ${FA} + {FC} + {FB}$ 可以写成 $1 \times  {FA} + 1 \times  {FC} + 1 \times  {FB}$ ,这个式子中 ${FA}\text{ 、 }{FC}\text{ 、 }{FB}$ 的系数均为 1,以三个系数 1 作为边长可构造一个等边三角形，因此我们尝试把 $\bigtriangleup  {AFC}$ 绕点 $C$ 顺时针旋转 $\dfrac{\pi }{3}$ ，得到 $\bigtriangleup  {A}^{\prime }{F}^{\prime }C$ (如图 2)，所以 ${FA} + {FC} + {FB} =  {F}^{\prime }{A}^{\prime } + {F}^{\prime }F + {FB}$ ，显然 ${F}^{\prime }{A}^{\prime } + {F}^{\prime }F + {FB} \geq  {A}^{\prime }B$ ，当 ${A}^{\prime }$ 、 ${F}^{\prime }$ 、 $F$ 、 $B$ 四点共线时 (如图 3)， ${FA} +  {FC} + {FB}$ 最小.

    \begin{center}
    \includegraphics[width=0.8\textwidth]{bodyxb/src/132_276_480_308_217_0.jpg}
    \end{center}

    试用类似的方法解答下面这道题目: 已知 $\vv{a}$ 是平面内的任意一个向量,向量 $\vv{b}\text{ 、 }\vv{c}$ 满足 $\vv{b} \cdot  \vv{c} = 0$ , 且 $\left| \vv{b}\right|  = 4,\left| \vv{c}\right|  = 4$ ，则 $\sqrt{2}\left| {\vv{a} - \vv{b}}\right|  + \left| {\vv{a} - \vv{c}}\right|  + \left| {\vv{a} + \vv{c}}\right|$ 的最小值为\blkda[2]{$8\sqrt{2}$}.
    \begin{jieda}
        在如图所示的平面直角坐标系 ${xOy}$ 中,设 $B\left( {4,0}\right) \text{ 、 }{C}_{1}\left( {0,4}\right)$ 、 ${C}_{2}\left( {0, - 4}\right)$ ,不妨设 $\vv{a} = \vv{OA},\vv{b} = \vv{OB}$ , $\vv{c} = \vv{O{C}_{1}}$ ,由题意可得 $\sqrt{2}\left| {\vv{a} - \vv{b}}\right|  + \left| {\vv{a} - }\right|  \vv{c}\left| +\right| \vv{a} + \vv{c}\left| { = \sqrt{2}}\right| {AB}\left| +\right| A{C}_{1} \mid   +  \left| {A{C}_{2}}\right|$ ,将 $\bigtriangleup {AB}{C}_{2}$ 绕点 $B$ 逆时针旋转 $\dfrac{\pi }{2}$ 得到 $\bigtriangleup {A}^{\prime }B{C}^{\prime }$ ,则 $\left| {A{A}^{\prime }}\right|  = \sqrt{2}\left| {AB}\right| ,\left| {{A}^{\prime }{C}^{\prime }}\right|  = \left| {A{C}_{2}}\right|$ , 其中点 ${C}^{\prime }\left( {8, - 4}\right)$ ,故 $\sqrt{2}\left| {\vv{a} - \vv{b}}\right|  + \left| {\vv{a} - \vv{c}}\right|  + \left| {\vv{a} + \vv{c}}\right|  = \sqrt{2}\left| {AB}\right|  +  \left| {A{C}_{1}}\right|  + \left| {A{C}_{2}}\right|  = \left| {A{A}^{\prime }}\right|  + \left| {A{C}_{1}}\right|  + \left| {{A}^{\prime }{C}^{\prime }}\right|  \geq  \left| {{C}_{1}{C}^{\prime }}\right|  =  \sqrt{{\left( 8 - 0\right) }^{2} + {\left( -4 - 4\right) }^{2}} = 8\sqrt{2}$ ,当且仅当点 $A$ 与点 $B$ 重合时,此时,点 ${A}^{\prime }$ 也与点 $B$ 重合,等号成立,故 $\sqrt{2}\left| {\vv{a} - \vv{b}}\right|  + \left| \vv{a}  -  \vv{c}\right| +\left| \vv{a} + \vv{c} \right|$ 的最小值为 $8\sqrt{2}$.
        \begin{center}
        \includegraphics[width=0.2\textwidth]{bodyxb/src/95_578_1402_242_227_0.jpg}
        \end{center}
    \end{jieda}

    \item (2025$\cdot$嘉定区二模$\cdot$12)在平面直角坐标系中，一质点 $P$ 从原点 $O$ 出发，第一次从点 $O$ 移动到点 ${P}_{1}$ ,第二次从点 ${P}_{1}$ 移动到点 ${P}_{2},\cdots$ ,第 $k$ 次从点 ${P}_{k - 1}$ (规定 ${P}_{0} = O$ )移动到点 ${P}_{k}$. 记向量 ${\vv{v}}_{k} =  \vv{{P}_{k - 1}{P}_{k}}$ ，其模长为 $k$ ，方向与 $x$ 轴正方向成 ${\left( {90}k\right) }^{ \circ  }$ 角，设 $\vv{{S}_{n}}$ 为经过 $n$ 次移动的位移向量，即 $\vv{{S}_{n}} =  \vv{O{P}_{n}}$ ，则当 $\left| \vv{{S}_{n}}\right|  = \sqrt{85}$ 时， $n$ 的值为\blkda[2]{$13$}.
    \begin{jieda}
        根据题意可知 ${\vv{v}}_{k} = \vv{{P}_{k - 1}{P}_{k}}$ 的模长为 $k$ ,方向与 $x$ 轴正方向成 ${\left( {90}k\right) }^{ \circ  }$ 角,所以 ${\vv{v}}_{k} = \vv{{P}_{k - 1}{P}_{k}} = \left( {k\cos {\left( {90}k\right) }^{ \circ  },k\sin }\right.  \left. {\left( {90}k\right) }^{ \circ  }\right)$ ,所以 $\vv{{S}_{n}} = \vv{O{P}_{n}} = \left( {1 \times  \cos {90}^{ \circ  } + 2\cos {180}^{ \circ  } + 3\cos {270}^{ \circ  }}\right.  + \cdots  + n\cos {\left( {90}n\right) }^{ \circ  },1 \times  \sin {90}^{ \circ  } + 2\sin {180}^{ \circ  } + 3\sin {270}^{ \circ  } + \cdots  +  \left. {n\sin {\left( {90}n\right) }^{ \circ  }}\right)$
        
        ${\vv{S}}_{4} = \vv{O{P}_{4}} = \left( {-2 + 4,1 - 3}\right)  = \left( {2, - 2}\right),\left| {\vv{S}}_{4}\right|  = 2\sqrt{2}\\
        {\vv{S}}_{8} = \vv{O{P}_{8}} = \left( {-2 + 4 - 6 + 8,1 - 3 + 5 - 7}\right)  = (4$ , $- 4),\left| \vv{{S}_{8}}\right|  = 4\sqrt{2}\\
        \vv{{S}_{12}} = \vv{O{P}_{12}} = ( - 2 + 4 - 6 + 8 - {10} + {12}$ , $1 - 3 + 5 - 7 + 9 - {11}) = \left( {6, - 6}\right) ,\left| \vv{{S}_{12}}\right|  = 6\sqrt{2}\\
        \vv{{S}_{13}} = \vv{O{P}_{13}} =  \left( {-2 + 4 - 6 + 8 - {10} + {12},1 - 3 + 5 - 7 + 9 - {11} + {13}}\right)  = (6$ , 7), $\left| \vv{{S}_{13}}\right|  = \sqrt{85}$.
        
        \ps 找规律即可.
    \end{jieda}

    \item (2024$\cdot$虹口区一模$\cdot$12)设 $\vv{a_1}$ 、 $\vv{a_2}$ 、 $\vv{a_3}$ 、 $\vv{b_1}$ 、 $\vv{b_2}$ 、 $\vv{b_3}$ 是平面上两两不相等的向量，若 $\left|\vv{a_1} - \vv{a_2}\right|  = \left|\vv{a_2} - \vv{a_3}\right|  = \left|\vv{a_3} - \vv{a_1}\right| = 2$ ，且对任意的 $i$ 、 $j \in  \{ 1,2,3\}$ ，均有 $\left| {{\vv{a}}_{i} - {\vv{b}}_{j}}\right|  \in  \{ 1,\sqrt{3}\}$ ，则 $\left|\vv{b_1} - \vv{b_2}\right|  + \left| \vv{b_2} - \vv{b_3}\right|  +  \left|\vv{b_3} - \vv{b_1}\right|  =$\blkda[2]{$3$}.
    \begin{jieda}
        向量 $\vv{O{A}_{i}} = {\vv{a}}_{i},\vv{O{B}_{i}} = {\vv{b}}_{i},i\text{ 、 }j \in  \{ 1,2,3\}$ ,则 ${A}_{1}{A}_{2} =  {A}_{2}{A}_{3} = {A}_{3}{A}_{1} = 2$ ,因为 $\left| {{\vv{a}}_{i} - {\vv{b}}_{j}}\right|  \in  \{ 1,\sqrt{3}\}$ ,所以 $\left| {\vv{{A}_{i}{B}_{j}} \mid   = 1\text{ 或 }}\right|  \vv{{A}_{i}{B}_{j}} \mid   = \sqrt{3}$ ，又等边 $\bigtriangleup {A}_{1}{A}_{2}{A}_{3}$ 边长为 2，高线 ${A}_{1}{B}_{2} = {A}_{2}{B}_{3} =  {A}_{3}{B}_{1} = \sqrt{3}$ ,故 ${B}_{1}\text{ 、 }{B}_{2}\text{ 、 }{B}_{3}$ 分别为 ${A}_{1}{A}_{2}\text{ 、 }{A}_{2}{A}_{3}\text{ 、 }{A}_{3}{A}_{1}$ 中点 $\Longrightarrow  {B}_{1}{B}_{2} = {B}_{2}{B}_{3} = {B}_{3}{B}_{1} = 1$ ,即 $\left| {{\vv{b}}_{1} - {\vv{b}}_{2}}\right|  + \left| {{\vv{b}}_{2} - {\vv{b}}_{3}}\right|  +  \left| {{\vv{b}}_{3} - {\vv{b}}_{1}}\right|  = 1 + 1 + 1 = 3.$

        \begin{center}
        \includegraphics[width=0.3\textwidth]{bodyxb/src/95_1056_787_275_280_0.jpg}
        \end{center}
    \end{jieda}
\end{enumerate}

\subsubsection*{二、选择题}
\begin{enumerate}
    \item (2024$\cdot$杨浦区二模$\cdot$16)平面上的向量 $\vv{a}$ 、 $\vv{b}$ 满足: $\left| \vv{a}\right|  = 3$ ， $\left| \vv{b}\right|  = 4$ ， $\vv{a}\bot \vv{b}$ ，定义该平面上的向量集合 $A = \left\{ \vv{x}\left| \left|  \vv{x} + \vv{a}\right|  < \left|\vv{x} + \vv{b} \right| ,\vv{x} \cdot  \vv{a} > \vv{x} \cdot  \vv{b}\right. \right\}$. 给出如下两个结论:

    \ding{172}  对任意 $\vv{c} \in  A$ ，存在该平面的向量 $\vv{d} \in  A$ ，满足 $\left| {\vv{c} - \vv{d}}\right|  = 0.5$;

    \ding{173}  对任意 $\vv{c} \in  A$ ，存在该平面的向量 $\vv{d} \notin  A$ ，满足 $\left| {\vv{c} - \vv{d}}\right|  = 0.5$.

    则下面判断正确的是 \dotfill{(\kh{C})}
    \xx{\ding{172} 正确， \ding{173} 错误}{\ding{172} 错误， \ding{173} 正确}{\ding{172}  \ding{173} 均正确}{\ding{172}  \ding{173} 均错误}
    \begin{jieda}
        设 $\vv{x} = \left( {m,n}\right) ,\vv{a} = \left( {3,0}\right) ,\vv{b} = \left( {0,4}\right) ,\left| {\vv{x} + \vv{a}}\right|  <  \mid  \vv{x} +  \vv{b}\left|  \Longrightarrow  \right| \vv{x} + \vv{a}{\left. \right| }^{2} < {\left| \vv{x} + \vv{b}\right| }^{2}$ ,而 $\vv{x} + \vv{a} = \left( {m + 3,n}\right) ,\vv{x} + \vv{b} = (m, n + 4),{\left( m + 3\right) }^{2} + {n}^{2} < {m}^{2} + {\left( n + 4\right) }^{2} \Longrightarrow  {6m} - {8n} - 7 < 0$  
        
        $\vv{x} \cdot  \vv{a} > \vv{x} \cdot  \vv{b} \Longrightarrow  {3m} - {4n} > 0$
        
        $A$对应直线 ${6m} - {8n} - 7 = 0$ 和 ${3m} -  {4n} = 0$ 间的区域（不含边界），无论 $\vv{d}$ 是否属于 $A$ ,都有 $\left| {\vv{c} - \vv{d}}\right|  = 0.5$.
    \end{jieda}
\end{enumerate}
\subsection*{考点2 向量的基本定理及其应用}

\subsubsection*{一、填空题}
\begin{enumerate}
    \item (2025$\cdot$金山区一模$\cdot$2)已知向量 $\vv{a} = \left( {1,2}\right)$ ， $\vv{b} = \left( {3,k}\right)$ ，若 $\vv{a}//\vv{b}$ ，则实数 $k =$ \blkda[2]{6}.

    \item (2025$\cdot$闵行区二模$\cdot$3)已知向量 $\vv{a} = \left( {x,1}\right)$ ， $\vv{b} = \left( {1,2 - x}\right)$ ，若 $\vv{a}//\vv{b}$ ，则实数 $x =$ \blkda[2]{1}.

    \item (2025$\cdot$徐汇区二模$\cdot$3)在空间直角坐标系中，向量 $\vv{a} = \left( {-m,6,3}\right)$ 、 $\vv{b} = \left( {2,n,1}\right)$ ，若 $\vv{a}//\vv{b}$ ，则 $m + n =$ \blkda[2]{$-4$}.

    \item (2025$\cdot$上海春季卷$\cdot$4)已知 $\vv{a} = \left( {2,1}\right)$ ， $\vv{b} = \left( {1,x}\right)$ ，若 $\vv{a}//\vv{b}$ ，则 $x =$ \blkda[2]{$\dfrac{1}{2}$}.

    \item (2024$\cdot$长宁区一模$\cdot$4)设向量 $\vv{a} = \left( {1, - 2}\right)$ ， $\vv{b} = \left( {-1,m}\right)$ ，若 $\vv{a}//\vv{b}$ ，则 $m =$ \blkda[2]{$2$}.

    \item (2024$\cdot$上海秋季卷$\cdot$5)已知 $k \in  \mathbf{R},\vv{a} = \left( {2,5}\right) ,\vv{b} = \left( {6,k}\right)$ ，且 $\vv{a}//\vv{b}$ ，则 $k$ 的值为\blkda[2]{$15$}.

    \item (2025$\cdot$崇明区一模$\cdot$7)已知 $\vv{a} = \left( {1,2}\right)$ ， $\vv{b} = \left( {2,k}\right)$ ，如果 $\vv{a}//\vv{b}$ ，那么实数 $k$ 的值为\blkda[2]{$4$}.

    \item (2024$\cdot$青浦区一模$\cdot$9)已知向量 $\vv{d} = \left( {1, - 1}\right)$ 垂直于直线 $l$ 的法向量，过 $A\left( {1,1}\right)$ 、 $B\left( {-1,8}\right)$ 分别作直线 $l$ 的垂线，对应垂足为 ${A}_{1}$ 和 ${B}_{1}$ ，若 $\vv{{A_1}{B_1}} = \lambda \vv{d}$ ，则实数 $\lambda$ 的值为\blkda[2]{$- \dfrac{9}{2}$}.
    \begin{jieda}
        $\vv{{A_1}{B_1}}$是$\vv{AB} = (-2, 7)$在直线的方向向量$\vv{d}$上的投影向量，$\vv{{A_1}{B_1}} = \dfrac{\vv{AB} \cdot \vv{d}}{\left|\vv{d}\right|} \cdot \dfrac{\vv{d}}{\left|\vv{d}\right|} = -\dfrac{9}{2}\vv{d}$
    \end{jieda}

    \item (2024$\cdot$松江区二模$\cdot$10)已知正三角形 ${ABC}$ 的边长为 2，点 $D$ 满足 $\vv{CD} = m\vv{CA} + n\vv{CB}$ ，且 $m > 0,n > 0,{2m} + n = 1$ ，则 $\left| \vv{CD}\right|$ 的取值范围是\blkda[2]{$\left( {1,2}\right)$}.
    \begin{jieda}
        设$E$为$CA$中点，因此$\vv{CD} = 2m\vv{CE} + n\vv{CB}$，因为$m > 0,n > 0,{2m} + n = 1$，所以$D$在线段$BE$上（不含端点），因此$\left| \vv{CD}\right|$ 的取值范围是$\left( {1,2}\right)$.

        \begin{center}
            \includegraphics[width=0.2\textwidth]{bodyxb/src/96_701_1591_168_141_0.jpg}
        \end{center}
    \end{jieda}
\end{enumerate}

\subsubsection*{二、选择题}
\begin{enumerate}
    \item (2025$\cdot$宝山区二模$\cdot$13)已知向量 $\vv{a} = \left( {1,x}\right)$ ， $\vv{b} = \left( {3,1}\right)$ ，若 $\vv{a}//\vv{b}$ ，则 $x$ 的值为 \dotfill{(\kh{D})}
    \xx{$-3$}{$3$}{$- \dfrac{1}{3}$}{$\dfrac{1}{3}$}
\end{enumerate}

\subsection*{考点3 向量的数量积与向量的坐标表示}

\subsubsection*{一、填空题}
\begin{enumerate}
    \item (2023$\cdot$上海秋季卷$\cdot$2)若 $\vv{a} = \left( {-2,3}\right)$ ， $\vv{b} = \left( {1,2}\right)$ ，则 $\vv{a} \cdot  \vv{b} =$ \blkda[2]{$4$}.

    \item (2025$\cdot$浦东新区二模$\cdot$2)已知向量 $\vv{a} = \left( {1,2}\right)$ ， $\vv{b} = \left( {m,1}\right)$ ，若 $\vv{a}\bot \vv{b}$ ，则 $m =$ \blkda[2]{$-2$}.

    \item (2025$\cdot$嘉定区二模$\cdot$3)已知向量 $\vv{m} = \left( {1, - 2}\right)$ ， $\vv{n} = \left( {k,4}\right)$ ，若 $\vv{m}\bot \vv{n}$ ，则 $k =$ \blkda[2]{$8$}.

    \item (2024$\cdot$黄浦区二模$\cdot$3)若 $\vv{a} = \left( {3\cos \theta ,\sin \theta }\right)$ ， $\vv{b} = \left( {\cos \theta ,3\sin \theta }\right)$ ，其中 $\theta  \in  \mathbf{R}$ ，则 $\vv{a} \cdot  \vv{b} =$ \blkda[2]{$3$}.

    \item (2024$\cdot$宝山区二模$\cdot$4)已知向量 $\vv{a} = \left( {{2m},2}\right) ,\vv{b} = \left( {1,m + 1}\right)$ ，若 $\vv{a} \cdot  \vv{b} = {10}$ ，则实数 $m =$ \blkda[2]{$2$}.

    \item (2025$\cdot$崇明区二模$\cdot$5)已知 $\vv{a} = \left( {1,0}\right)$ ， $\vv{b} = \left( {2,1}\right)$ ，则 $\left| {\vv{a} + 2\vv{b}}\right|  =$ \blkda[2]{$\sqrt{29}$}.

    \item (2025$\cdot$青清区二模$\cdot$5)向量 $\vv{a} = (310, 118)$在向量 $\vv{b} = \left( {0,{2025}}\right)$ 方向上的数量投影是\blkda[2]{$118$}.
    \begin{jieda}
        向量 $\vv{a}$ 在向量 $\vv{b}$ 方向上的数量投影为 $\dfrac{\vv{a} \cdot  \vv{b}}{\left| \vv{b}\right| } = \dfrac{{310} \times  0 + {118} \times  {2025}}{2025} = {118}$.
    \end{jieda}

    \item (2024$\cdot$浦东新区一模$\cdot$6)已知向量 $\vv{a} = \left( {3,4}\right)$ ，向量 $\vv{b} = \left( {1,0}\right)$ ，则向量 $\vv{a}$ 在向量 $\vv{b}$ 上的投影为\blkda[2]{$\left( {3,0}\right)$}.
    \begin{jieda}
        $\left( \dfrac{\vv{a} \cdot  \vv{b}}{{\left| \vv{b}\right| }^{2}}\right) \vv{b} = \dfrac{3}{1} \times  \left( {1,0}\right)  = \left( {3,0}\right)$.
    \end{jieda}

    \item (2025$\cdot$杨浦区一模$\cdot$6)已知 $\left| \vv{a}\right|  = \left| \vv{b}\right|  = 1$ ，若$(2\vv{a} - \vv{b})\bot \vv{b}$ ，则向量 $\vv{a}$ 与 $\vv{b}$ 的夹角的余弦值为\blkda[2]{$\dfrac{1}{2}$}.
    \begin{jieda}
        设向量 $\vv{a}$ 与 $\vv{b}$ 的夹角为 $\theta \left( {\theta  \in  \left\lbrack  {0,\pi }\right\rbrack  }\right)$ ,若 $\left( {2\vv{a} - \vv{b}}\right)  \bot  \vv{b}$ ,则 $\left( {2\vv{a} - \vv{b}}\right)  \cdot  \vv{b} = 0$ ,所以 $2\vv{a} \cdot  \vv{b} - {\vv{b}}^{2} = 2\left| \vv{a}\right| \left| \vv{b}\right| \cos \theta  - {\left| \vv{b}\right| }^{2}  = 2\cos \theta  - 1 = 0$ ,可得 $\cos \theta  = \dfrac{1}{2}$.
    \end{jieda}

    \item (2024$\cdot$杨浦区一模$\cdot$7)已知向量 $\vv{a} = \left( {3,0}\right)$ 、 $\vv{b} = \left( {-2,2\sqrt{3}}\right)$ ，则 $\vv{b}$ 在 $\vv{a}$ 方向上的投影为\blkda[2]{$\left( {-2,0}\right)$}.
    \begin{jieda}
        向量 $\vv{a} = \left( {3,0}\right) ,\vv{b} = \left( {-2,2\sqrt{3}}\right)$ ,则 $\vv{b}$ 在 $\vv{a}$ 方向上的投影为 $\dfrac{\vv{a} \cdot  \vv{b}}{{\left| \vv{a}\right| }^{2}} \cdot  \vv{a} = \dfrac{-6}{9} \times  \left( {3,0}\right)  = \left( {-2,0}\right)$.
    \end{jieda}

    \item (2025$\cdot$长宁区一模$\cdot$8)已知向量 $\vv{a} = \left( {1,2}\right)$ 、 $\vv{b} = \left( {3, - 1}\right)$ ，则向量 $\vv{b}$ 在 $\vv{a}$ 方向上的投影的坐标是\blkda[2]{$\left( {\dfrac{1}{5},\dfrac{2}{5}}\right)$}.

    \item (2025$\cdot$宝山区二模$\cdot$9)已知 $\bigtriangleup  {ABC}$ 中， ${AB} = {AC} = 4$ ， $\angle {BAC} = \dfrac{2}{3}\pi$ ，点 $D$ 在线段 ${BC}$ 上，且 ${S}_{\bigtriangleup {ACD}} = 2{S}_{\bigtriangleup {ABD}}$ ，则 $\vv{AB} \cdot  \vv{AD}$ 的值为\blkda[2]{$8$}.
    \begin{jieda}
        ${S}_{\bigtriangleup {ACD}} = 2{S}_{\bigtriangleup {ABD}} \Longrightarrow {CD} = {2BD}$ ,所以 $\vv{AD} = \dfrac{2}{3}\vv{AB} + \dfrac{1}{3}\vv{AC}$ ,

        \begin{center}
        \includegraphics[width=0.2\textwidth]{bodyxb/src/96_1367_1919_210_112_0.jpg}
        \end{center}
        

        所以 $\vv{AB} \cdot  \vv{AD} = \vv{AB} \cdot  \left( {\dfrac{2}{3}\vv{AB} + \dfrac{1}{3}\vv{AC}}\right)  = \dfrac{2}{3}\vv{AB}{}^{2} + \dfrac{1}{3}\vv{AB}$ . $\vv{AC} = \dfrac{2}{3}{\left| \vv{AB}\right| }^{2} + \dfrac{1}{3}\left| \vv{AB}\right|  \cdot  \left| \vv{AC}\right| \cos \angle {BAC} = \dfrac{2}{3} \times  {4}^{2} +  \dfrac{1}{3} \times  4 \times  4 \times  \left( {-\dfrac{1}{2}}\right)  = 8$.
    \end{jieda}

    \item (2024$\cdot$虹口区二模$\cdot$9)已知平面向量 $\vv{a}$ 、 $\vv{b}$ 满足 $\left| \vv{a}\right|  = 3,\left| \vv{b}\right|  = 4,\vv{a} \cdot  \vv{b} = 4$ ，若平面向量 $\vv{c}$ 满足 $\left| {\vv{c} - \vv{b}}\right|  = 1$ ，则 $\left| {\vv{c} - \vv{a}}\right|$ 的最大值为\blkda[2]{$\sqrt{17} + 1$}.
    \begin{jieda}
        以$\vv{b}$为$x$轴正方向建系，$A(1, 2\sqrt{2}), B(4, 0)$，$|AB| = \sqrt{17}$，$C$在以$B$为圆心，$1$为半径的圆上，$|AC| \leq \sqrt{17} + 1$
    \end{jieda}

    \item (2024$\cdot$徐汇区一模$\cdot$9)在 $\bigtriangleup  {ABC}$ 中， ${AC} = {BC},{P}_{1}$ 、 ${P}_{2}$ 、 ${P}_{3}$ 为边 ${AB}$ 上的点，且 $8\left| \vv{{P}_{1}B}\right|  = 4\left| \vv{{P}_{2}B}\right|  =  2\left| \vv{{P}_{3}B}\right|  = \left| \vv{AB}\right|  = 8$ ,设 ${I}_{k} = \vv{{P}_{k}B} \cdot  \vv{{P}_{k}C}\left( {k = 1,2,3}\right)$ ,则 ${I}_{1} - {I}_{2} + {I}_{3} =$ \blkda[2]{1}.
    \begin{jieda}
        如图所示,由题意得 ${P}_{3}B = 4$ , ${P}_{2}B = 2,{P}_{1}B = 1$ ,且 ${P}_{3}C \bot  {P}_{3}B$ . 从而 ${I}_{1} = \vv{{P}_{1}B} \cdot  \vv{{P}_{1}C} =  - {P}_{1}B \cdot  {P}_{1}{P}_{3} =  - 3;  {I}_{2} = \vv{{P}_{2}B} \cdot  \vv{{P}_{2}C} =  - {P}_{2}B \cdot  {P}_{2}{P}_{3} =  - 4;  {I}_{3} = \vv{{P}_{3}B} \cdot  \vv{{P}_{3}C} = 0$ ; 因此 ${I}_{1} - {I}_{2} + {I}_{3} = 1$.

        \begin{center}
        \includegraphics[width=0.2\textwidth]{bodyxb/src/97_608_776_212_175_0.jpg}
        \end{center}
    \end{jieda}

    \item (2025$\cdot$普陀区一模$\cdot$10)设平面上四点 $P$ 、 $Q$ 、 $M$ 、 $N$ 满足: $\left| {PM}\right|  = \left| {PN}\right|  = 4$ ， $\left| {PQ}\right|  = 2$ ，若 $\vv{QM} \cdot  \vv{QN} = 0$ ，则 $\left| {MN}\right|$ 的最小值为\blkda[2]{$2\sqrt{7} - 2$}.
    \begin{jieda}
        以$P$为原点，$\vv{PQ}$为$x$轴正方向建系，$Q(2, 0)$，$M,N$在以$P$为圆心，$4$为半径的圆上，根据垂径定理可知$P$到$MN$距离越大，$|MN|$越小.
        
        取 ${MN}$ 的中点 $F$，则$|PF|$为$P$到$MN$距离，根据几何关系应满足$|PF| \leq |PQ| + |QF|$，设$|MN| = 2t, t > 0$，则$|PF| = \sqrt{16-t^2}$，$|QF| = \dfrac{|MN|}{2} = t$，因此$\sqrt{16-t^2} \leq 2+t$，解得$t \geq \sqrt{7}-1$，故$\left| {MN}\right|$ 的最小值为 $2\sqrt{7}-2$.
        \begin{center}
        \includegraphics[width=0.2\textwidth]{bodyxb/src/97_576_1182_244_234_0.jpg}
        \end{center}

        \ps $|PF|$变大$\Longleftrightarrow |QF|$变小，在三点共线时取得最值.
    \end{jieda}

    \item (2024$\cdot$闵行区一模$\cdot$10)若平面上的三个单位向量 $\vv{a}$ 、 $\vv{b}$ 、 $\vv{c}$ 满足 $\left| {\vv{a} \cdot  \vv{b}}\right|  = \dfrac{1}{2}$ ， $\left| {\vv{a} \cdot  \vv{c}}\right|  = \dfrac{\sqrt{3}}{2}$ ，则 $\vv{b} \cdot  \vv{c}$ 的所有可能的值组成的集合为\blkda[2]{$\left\{  {0,\dfrac{\sqrt{3}}{2}, - \dfrac{\sqrt{3}}{2}}\right\}$}.
    \begin{jieda}
        不妨设$\vv{a} = (1, 0)$，则$\vv{b} = \left(\pm \dfrac{1}{2}, \pm \dfrac{\sqrt{3}}{2}\right), \vv{c} = \left(\pm \dfrac{\sqrt{3}}{2}, \pm \dfrac{1}{2}\right)$，故$\vv{b} \cdot  \vv{c} \in \left\{  {0,\dfrac{\sqrt{3}}{2}, - \dfrac{\sqrt{3}}{2}}\right\}$
    \end{jieda}

    \item (2025$\cdot$闵行区一模$\cdot$10)已知 ${F}_{1}$ 、 ${F}_{2}$ 分别为椭圆 $\dfrac{{x}^{2}}{4} + \dfrac{{y}^{2}}{2} = 1$ 的左、右焦点，过 ${F}_{1}$ 的直线交椭圆于 $A$ 、 $B$ 两点. 若 $\vv{A{F}_{1}} \cdot  \vv{A{F}_{2}} = 0$ ，则 $\vv{A{F}_{2}} \cdot  \vv{B{F}_{2}} =$ \blkda[2]{4}.
    \begin{jieda}
        由题意,得 $c = \sqrt{4 - 2} =  \sqrt{2} = b$，故$A$应为短轴顶点，不妨取 $A(0,\sqrt{2})$，则 ${AB}$ 的方程为 $y = x + \sqrt{2}$ ,代入 ${x}^{2} + 2{y}^{2} = 4$ ,得 ${x}^{2} +  2\left( {{x}^{2} + 2\sqrt{2}x + 2}\right)  = 4 \Longrightarrow  3{x}^{2} + 4\sqrt{2}x = 0 \Longrightarrow  B\left( {-\dfrac{4\sqrt{2}}{3}, - \dfrac{\sqrt{2}}{3}}\right)  \Longrightarrow  \vv{A{F}_{2}} = \left( {\sqrt{2}, - \sqrt{2}}\right) ,\vv{B{F}_{2}} = \left( {\dfrac{7\sqrt{2}}{3},\dfrac{\sqrt{2}}{3}}\right)  \Longrightarrow  \vv{A{F}_{2}} \cdot  \vv{B{F}_{2}} = 4$.
        \begin{center}
        \includegraphics[width=0.3\textwidth]{bodyxb/src/97_1292_157_236_202_0.jpg}
        \end{center}
    \end{jieda}

    \item (2024$\cdot$崇明区二模$\cdot$11)已知 $A$ 、 $B$ 、 $C$ 是半径为 1 的圆上的三个不同的点，且 $\left| \vv{AB}\right|  = \sqrt{3}$ ，则 $\vv{AB} \cdot  \vv{AC}$ 的最小值是\blkda[2]{$\dfrac{3}{2} - \sqrt{3}$}.
    \begin{jieda}
        $\vv{AC}$在$\vv{AB}$上数量投影最小即可.
    \end{jieda}

    \item (2024$\cdot$金山区二模$\cdot$12)已知平面向量 $\vv{a}$ 、 $\vv{b}$ 、 $\vv{c}$ 满足: $\left| \vv{a}\right|  = \left| \vv{b}\right|  = 1$ ， $\vv{a} \cdot  \vv{c} = \vv{b} \cdot  \vv{c} = 1$ ，则 $\vv{a} \cdot  \vv{b} + {\vv{c}}^{2}$ 的最小值为\blkda[2]{$2\sqrt{2} - 1$}.
    \begin{jieda}
        设 $\vv{a} = \vv{OA},\vv{b} = \vv{OB},\vv{c} = \vv{OC},\vv{a} \cdot  \vv{c} = \vv{b} \cdot  \vv{c} = 1$，可得$\vv{OC}$ 在 $\vv{OA},\vv{OB}$ 上的数量投影都为 1,即 ${CA}\bot {OA},{CB}\bot {OB}$ ,不妨设 $\angle {AOC} =  \angle {BOC} = \alpha$，所以 $\left| \vv{OC}\right|  = \dfrac{1}{\cos \alpha }$，$\vv{a} \cdot  \vv{b} + {\vv{c}}^{2} = \left| \vv{a}\right|  \cdot  \left| \vv{b}\right| \cos {2\alpha } +  {\left| \vv{c}\right| }^{2} = 2{\cos }^{2}\alpha  - 1 + \dfrac{1}{{\cos }^{2}\alpha } \geq  2\sqrt{2{\cos }^{2}\alpha  \cdot  \dfrac{1}{{\cos }^{2}\alpha }} - 1 = 2\sqrt{2} - 1.$

        \begin{center}
        \includegraphics[width=0.2\textwidth]{bodyxb/src/97_1291_1447_239_199_0.jpg}
        \end{center}
    \end{jieda}

    \item (2025$\cdot$金山区一模$\cdot$12)已知 $O$ 为坐标原点，向量 $\vv{OA}$ 、 $\vv{OB}$ 满足 $\left| \vv{OA}\right|  + \left| \vv{OB}\right|  = 8$ ，将 $\vv{OA}$ 绕点 $O$ 按逆时针方向旋转 ${90}^{ \circ  }$ ，得到向量 $\vv{OC}$. 若 $\vv{OB} + \vv{OC} = \left( {-3,3}\right)$ ， $\vv{i} = \left( {1,0}\right)$ ，则$(\vv{OA} + \vv{OB}) \cdot  \vv{i}$ 的最大值为\blkda[2]{$4\sqrt{2}$}.
    \begin{jieda}
        设 $A\left( {x,y}\right)$ ,则 $C\left( {-y,x}\right)$ ,故 $\vv{OB} + \vv{OC} = \left( {{x}_{B} - y,{y}_{B}}\right.  + x) = \left( {-3,3}\right)$ ,所以 $\left\{  \begin{array}{l} {x}_{B} - y =  - 3 \\  {y}_{B} + x = 3 \end{array}\right.$ ,则 $\left\{  \begin{array}{l} {x}_{B} = y - 3 \\  {y}_{B} = 3 - x \end{array}\right.$ ,故 $B(y - 3$ , $3 - x)$ ,由 $\left| \vv{OA}\right|  + \left| \vv{OB}\right|  = \sqrt{{x}^{2} + {y}^{2}} + \sqrt{{\left( x - 3\right) }^{2} + {\left( y - 3\right) }^{2}} =$ 8,即 $A\left( {x,y}\right)$ 到原点距离与 $A\left( {x,y}\right)$ 到 $\left( {3,3}\right)$ 的距离之和为 8，易知 $A$ 的轨迹为一个椭圆,且 ${2a} = 8 \Longrightarrow  a = 4$ , ${2c} = 2\sqrt{3} \Longrightarrow  c = \sqrt{3}$
        
        如图,若 $P$ 为椭圆的右顶点,则 $\left| {OP}\right|  = \dfrac{8 - 3\sqrt{2}}{2} + 3\sqrt{2}  = \dfrac{8 + 3\sqrt{2}}{2}$ ,且直线 ${OP} : y = x$ ,易知 ${x}_{P} = {y}_{P} = \dfrac{\left| OP\right| }{\sqrt{2}} =  \dfrac{4\sqrt{2} + 3}{2}$ ,即 $P\left( {\dfrac{4\sqrt{2} + 3}{2},\dfrac{4\sqrt{2} + 3}{2}}\right)$ ,由 $\left( {\vv{OA} + \vv{OB}}\right)  \cdot  \vv{i} = (x +  y - 3,y - x + 3) \cdot  \left( {1,0}\right)  = x + y - 3$ ,令 $m = x + y - 3$
        
        要使 $\left( {\vv{OA} + \vv{OB}}\right)  \cdot  \vv{i}$ 最大,只需 $y =  - x + m + 3$ 的截距最大,即 $m$ 最大,结合图知,该直线过 $P$ 点即可，此时有 $m = \dfrac{4\sqrt{2} + 3}{2} + \dfrac{4\sqrt{2} + 3}{2} - 3 = 4\sqrt{2}$.

        \begin{center}
        \includegraphics[width=0.3\textwidth]{bodyxb/src/97_1292_1930_239_219_0.jpg}
        \end{center}
        

        \ps 首先得到 $A$ 的轨迹为一个椭圆及其右顶点坐标,再应用向量数量积的坐标表示化为求 $m = x + y - 3$ 最大为关键.
    \end{jieda}

    \item (2025$\cdot$宝山区一模$\cdot$12)已知平面向量 $\vv{a}$ 、 $\vv{b}$ 满足: $\left| \vv{a}\right|  = 1$ ， $\left| \vv{b}\right|  = m \in  \left( {1,2}\right)$ ，且对任意的单位向量 $\vv{c}$ 满足 $\left| {\vv{a} \cdot  \vv{c}}\right|  + \left| {\vv{b} \cdot  \vv{c}}\right|  \leq  \sqrt{6}$ ，则 $\vv{a} \cdot  \vv{b}$ 的最大值为\blkda[5]{$\left\{  \begin{array}{ll} m,&1 < m \leq \sqrt{6} - 1 \\  \dfrac{5 - {m}^{2}}{2},&\sqrt{6} - 1 <  m < 2 \end{array}\right.$}. (用含 $m$ 的式子表示)
    \begin{jieda}
        因为要求$\vv{a} \cdot  \vv{b}$ 的最大值，对于特定的$m$，$<\vv{a}, \vv{b}>$越小$\Longleftrightarrow \left|\vv{a}+\vv{b}\right|$越大，因此不妨先考虑$\langle\vv{a}, \vv{b}\rangle< \dfrac{\pi}{2}$的情况.

        根据三角不等式有$\left|\left(\vv{a}+\vv{b}\right) \cdot \vv{c} \right| \leq \left| {\vv{a} \cdot  \vv{c}}\right|  + \left| {\vv{b} \cdot  \vv{c}}\right|$
        
        根据数量积的定义有$\left|\left(\vv{a}+\vv{b}\right) \cdot \vv{c} \right| \leq \left|\vv{a}+\vv{b}\right| \cdot |\vv{c}| = \left|\vv{a}+\vv{b}\right|$

        因为$\langle\vv{a}, \vv{b}\rangle < \dfrac{\pi}{2}$，当$\vv{c}$与$\vv{a} + \vv{b}$共线时，两个不等式均取等号，所以$\left|\vv{a}+\vv{b}\right| \leq \sqrt{6}$. 下面尽可能使得$\left|\vv{a}+\vv{b}\right|$更大.
        \begin{dingautolist}{172}
            \item 当$m \leq \sqrt{6} - 1$
            
            $\langle\vv{a}, \vv{b}\rangle= 0$时，$\left|\vv{a}+\vv{b}\right|$取得最大值$\left|\vv{a}\right|+\left|\vv{b}\right| \leq \sqrt{6}$，故$\left(\vv{a} \cdot  \vv{b}\right)_{max} = m > 0$
            \item 当$m \in \left(\sqrt{6}-1, 2\right)$
            
            令$\left|\vv{a}+\vv{b}\right| = \sqrt{6}$，$6 = \left|\vv{a} + \vv{b}\right|^2 = 1+m^2+2\vv{a} \cdot  \vv{b}$，故$\left(\vv{a} \cdot  \vv{b}\right)_{max} = \dfrac{5-m^2}{2} > 0$
            
            \ps 此处根据三角形法则解三角形亦可.
        \end{dingautolist}
        $\langle\vv{a}, \vv{b}\rangle < \dfrac{\pi}{2}$的情况下已经找到$\left(\vv{a} \cdot  \vv{b}\right)_{max} = \left\{  \begin{array}{ll} m,&1 < m \leq \sqrt{6} - 1 \\  \dfrac{5 - {m}^{2}}{2},&\sqrt{6} - 1 <  m < 2 \end{array}\right.$，故无需考虑$\langle\vv{a}, \vv{b}\rangle\geq \dfrac{\pi}{2}$的情况.
    \end{jieda}

    \item (2024$\cdot$长宁区二模$\cdot$12)已知平面向量 $\vv{a}$ 、 $\vv{b}$ 、 $\vv{c}$ 满足: $\left| \vv{a}\right|  = \left| \vv{b}\right|  = \sqrt{10}$ ， $\left| \vv{c}\right|  = 2$ ，若 $\left( {\vv{c} - \vv{a}}\right) \cdot \left(\vv{c} - \vv{b}\right) = 0$ ,则 $\left| {\vv{a} - \vv{b}}\right|$ 的最小值为\blkda[2]{2}.
    \begin{jieda}
        方法同 (2025$\cdot$普陀区一模$\cdot$10).
    \end{jieda}

    \item (2025$\cdot$青浦区一模$\cdot$12)已知 $A$ 、 $B$ 、 $C$ 是单位圆上任意不同三点，则 $\dfrac{\vv{AB} \cdot  \vv{AC}}{\left| \vv{AB}\right| }$ 的取值范围是\blkda[2]{$\left( {-1,2}\right)$}.
    \begin{jieda}
        $\dfrac{\vv{AB} \cdot  \vv{AC}}{\left| \vv{AB}\right| }$即$\vv{AC}$在$\vv{AB}$上的数量投影，以$\vv{AB}$为$x$轴正方向，设$A(x_0, y_0), x_0 \in (-1, 0)$
        
        固定$x_0$，则$\vv{AC}$在$x$轴上的数量投影取值范围是$(-1-x_0, 1-x_0)$，因此$\dfrac{\vv{AB} \cdot  \vv{AC}}{\left| \vv{AB}\right| } \in \left( {-1,2}\right)$

        \twotu{\includegraphics[width=0.4\textwidth]{bodyxb/src/98_1022_1632_180_169_0.jpg}}{\includegraphics[width=0.4\textwidth]{bodyxb/src/98_1275_1632_189_169_0.jpg}}
    \end{jieda}

    \item (2025$\cdot$上海秋季卷$\cdot$12)已知 $f\left( x\right)  = \left\{  \begin{array}{ll} 1,&x > 0, \\  0,&x = 0, \\   - 1,&x < 0. \end{array}\right. \vv{a}\text{ 、 }\vv{b}\text{ 、 }\vv{c}$ 是平面内三个不同的单位向量. 若 $f\left( {\vv{a} \cdot  \vv{b}}\right)  + f\left( {\vv{b} \cdot  \vv{c}}\right)  + f\left( {\vv{c} \cdot  \vv{a}}\right)  = 0$ ，则 $\left| {\vv{a} + \vv{b} + \vv{c}}\right|$ 的取值范围是\blkda[2]{$\left( {1,\sqrt{5}}\right)$}.
        \begin{jieda}
            若 $f\left( {\vv{a} \cdot  \vv{b}}\right)  = f\left( {\vv{b} \cdot  \vv{c}}\right)  = f\left( {\vv{c} \cdot  \vv{a}}\right)  = 0$ ,则 $\vv{a} \cdot  \vv{b} = \vv{b} \cdot  \vv{c} = \vv{c} \cdot  \vv{a} = 0$ ,
            
            又三个向量均为平面内的单位向量,故向量 $\vv{a},\vv{b},\vv{c}$ 两两垂直,显然不成立; 故 $\left\{  {f\left( {\vv{a} \cdot  \vv{b}}\right) ,f\left( {\vv{b} \cdot  \vv{c}}\right) ,f\left( {\vv{c} \cdot  \vv{a}}\right) }\right\}   = \{  - 1,0,1\}$ .
            
            不妨设 $\left\{  \begin{array}{l} f\left( {\vv{a} \cdot  \vv{b}}\right)  = 1 \\  f\left( {\vv{b} \cdot  \vv{c}}\right)  = 0 \\  f\left( {\vv{c} \cdot  \vv{a}}\right)  =  - 1 \end{array}\right.$ ,则 $\vv{a} \cdot  \vv{b} > 0,\vv{b} \cdot  \vv{c} = 0,\vv{c} \cdot  \vv{a} < 0$ ,
            
            不妨设 $\vv{b} = \left( {1,0}\right) ,\vv{c} = \left( {0,1}\right) ,\vv{a} = \left( {\cos \theta ,\sin \theta }\right) ,\theta  \in  \lbrack 0,{2\pi })$ ,
            
            则 $\left\{  \begin{array}{l} \vv{a} \cdot  \vv{b} = \cos \theta  > 0 \\  \vv{c} \cdot  \vv{a} = \sin \theta  < 0 \end{array}\right.$ ,则 $\theta  \in  \left( {\dfrac{3}{2}\pi ,{2\pi }}\right)$ ,
            
            则 $\left| {\vv{a} + \vv{b} + \vv{c}}\right|  = \left| \left( {1 + \cos \theta ,1 + \sin \theta }\right) \right|  = \sqrt{{\left( 1 + \cos \theta \right) }^{2} + {\left( 1 + \sin \theta \right) }^{2}} = \sqrt{3 + 2\cos \theta  + 2\sin \theta }  = \sqrt{3 + 2\sqrt{2}\sin \left( {\theta  + \dfrac{\pi }{4}}\right) }$ ,
            
            由 $\theta  \in  \left( {\dfrac{3}{2}\pi ,{2\pi }}\right) ,\theta  + \dfrac{\pi }{4} \in  \left( {\dfrac{7}{4}\pi ,\dfrac{9}{4}\pi }\right)$ ,
            
            则 $\sin \left( {\theta  + \dfrac{\pi }{4}}\right)  \in  \left( {-\dfrac{\sqrt{2}}{2},\dfrac{\sqrt{2}}{2}}\right) ,2\sqrt{2}\sin \left( {\theta  + \dfrac{\pi }{4}}\right)  \in  \left( {-2,2}\right)$
            
            故 $\left| {\vv{a} + \vv{b} + \vv{c}}\right|  \in  \left( {1,\sqrt{5}}\right)$ .
        \end{jieda}

    \item (2024$\cdot$徐汇区二模$\cdot$12)如图所示，已知 $\bigtriangleup  {ABC}$ 满足 ${BC} = 8,{AC} =  {3AB},P$ 为 $\bigtriangleup {ABC}$ 所在平面内一点. 定义点集 $D = \left\{  {P\left| {\vv{AP} = {3\lambda }\vv{AB} + }\right. }\right.  \left. {\dfrac{1 - \lambda }{3}\vv{AC},\lambda  \in  \mathbf{R}}\right\}$. 若存在点 ${P}_{0} \in  D$ ,使得对任意 $P \in  D$ ,满足 $\left| \vv{AP}\right|  \geq  \left| \vv{A{P}_{0}}\right|$ 恒成立，则 $\left| \vv{A{P}_{0}}\right|$ 的最大值为\blkda[2]{$3$}.

    \begin{center}
    \includegraphics[width=0.2\textwidth]{bodyxb/src/135_1120_548_330_144_0.jpg}
    \end{center}
    \begin{jieda}
        取$M$满足$\vv{AM} = 3\vv{AB}$，$N$满足$\vv{AN} = \dfrac{1}{3}\vv{AC}$，故集合$D$为直线$MN$，$\left| \vv{AP}\right|  \geq  \left| \vv{A{P}_{0}}\right|$ 恒成立即$AP_0 \bot MN$，根据对称性可知$A$到$BC$的距离和$A$到$MN$距离相等，又因为$A$在半径为$3$的阿氏圆上，因此$\left| \vv{A{P}_{0}}\right|$ 的最大值为$3$.
        \begin{center}
        \includegraphics[width=0.2\textwidth]{bodyxb/src/99_578_509_242_178_0.jpg}
        \end{center}
    \end{jieda}
\end{enumerate}

\subsubsection*{二、选择题}
\begin{enumerate}
    \item (2025$\cdot$杨浦区二模$\cdot$15)已知 $A$ 、 $B$ 、 $C$ 是单位圆上的三个点，若 $\left| \vv{AB}\right|  = \sqrt{2}$ ，则 $\vv{AB} \cdot \vv{BC}$ 的最大值为 \dotfill{(\kh{D})}
    \xx{$\sqrt{2}$}{$1 + \dfrac{\sqrt{2}}{2}$}{$\sqrt{2} + 1$}{$\sqrt{2} - 1$}

    \item (2024$\cdot$静安区一模$\cdot$15)教材在推导向量的数量积的坐标表示公式“$\vv{a} \cdot  \vv{b} = {x}_{1}{x}_{2} + {y}_{1}{y}_{2}$ (其中 $\vv{a} = \left( {{x}_{1},{y}_{1}}\right) ,\vv{b} = \left( {{x}_{2},{y}_{2}}\right)$ )”的过程中,运用了以下哪些结论作为推理的依据 \dotfill{(\kh{D})}
    
    \ding{172}  向量坐标的定义；  \ding{173} 向量数量积的定义； \ding{174}  向量数量积的交换律； \ding{175}  向量数量积对数乘的结合律;  \ding{176}  向量数量积对加法的分配律.
    \xx{\ding{172}  \ding{174}  \ding{175}}{\ding{173}  \ding{175}  \ding{176}}{ \ding{172}  \ding{173}  \ding{174}  \ding{176}}{\ding{172}  \ding{173}  \ding{174}  \ding{175}  \ding{176}}
\end{enumerate}
\subsubsection*{三、解答题}

\begin{enumerate}
    \item (2025$\cdot$静安区二模$\cdot$17)已知向量 $\vv{a} = \left( {\cos \left( {x + \dfrac{\pi }{6}}\right) ,\sin \left( {x + \dfrac{\pi }{6}}\right) }\right) ,\vv{b} = (\cos \left( {x + \dfrac{\pi }{6}}\right) , - \sin (x +  \left. \dfrac{\pi }{6}\right)$ ,记 $f\left( x\right)  = \vv{a} \cdot  \vv{b}$.
    \begin{enumerate}
        \item 求函数 $y = f\left( x\right)$ 的最小正周期
        \item 若函数 $y = f\left( {x + \theta }\right)$ (其中常数 $\theta  \in  \left( {0,\dfrac{\pi }{2}}\right)$ )为奇函数，求 $\theta$ 的值.
    \end{enumerate} 
    \begin{jieda}
        \begin{enumerate}
            \item $f\left( x\right)  = \vv{a} \cdot  \vv{b} = {\cos }^{2}\left( {x + \dfrac{\pi }{6}}\right)  - {\sin }^{2}\left( {x + \dfrac{\pi }{6}}\right)  = \cos ({2x}  \left. {+\dfrac{\pi }{3}}\right)$ ,所以函数 $y = f\left( x\right)$ 的最小正周期 $T = \dfrac{2\pi }{2} = \pi$
            \item $y = f\left( {x + \theta }\right)  = \cos \left( {{2x} + {2\theta } + \dfrac{\pi }{3}}\right)$ ，因为函数 $y = f\left( {x + \theta }\right)$ 为奇函数，所以 ${2\theta } + \dfrac{\pi }{3} = \dfrac{\pi }{2} + {k\pi }$ ，解得 $\theta  = \dfrac{\pi }{12} + \dfrac{1}{2}{k\pi },k \in  \mathbf{Z}$ ,又因为 $\theta  \in  \left( {0,\dfrac{\pi }{2}}\right)$ ,所以 $\theta  = \dfrac{\pi }{12}$
        \end{enumerate}
    \end{jieda}
    
    \item (2025$\cdot$静安区一模$\cdot$18)已知向量 $\vv{a} = \left( {\cos \dfrac{3x}{2},\sin \dfrac{3x}{2}}\right) ,\vv{b} = \left( {\cos \dfrac{x}{2}, - \sin \dfrac{x}{2}}\right)$ ，且 $x \in  \left\lbrack 0,\dfrac{\pi }{2}\right\rbrack$.
    \begin{enumerate}
        \item 求 $\vv{a} \cdot  \vv{b}$ 及 $\left| {\vv{a} + \vv{b}}\right|$
        \item 记 $f\left( x\right)  = \vv{a} \cdot  \vv{b} - \left| {\vv{a} + \vv{b}}\right|$ ，求函数 $y = f\left( x\right)$ 的最小值.
    \end{enumerate}
    \begin{jieda}
        \begin{enumerate}
            \item 由题意得 $\vv{a} \cdot  \vv{b} = \cos \dfrac{3x}{2}\cos \dfrac{x}{2} - \sin \dfrac{3x}{2}\sin \dfrac{x}{2} =  \cos {2x}$ ,由于 $\left| \vv{a}\right|  = \sqrt{{\cos }^{2}\dfrac{3x}{2} + {\sin }^{2}\dfrac{3x}{2}} = 1,\left| \vv{b}\right|  = \sqrt{{\cos }^{2}\dfrac{x}{2} + {\left( -\sin \dfrac{x}{2}\right) }^{2}} = 1$ ,则 ${\left| \vv{a} + \vv{b}\right| }^{2} = {\vv{a}}^{2} + {\vv{b}}^{2} + 2\vv{a}$ . $\vv{b} = 2 + 2\vv{a} \cdot  \vv{b} = 2 + 2\cos {2x} = 4{\cos }^{2}x$ ,因为 $x \in  \left\lbrack  {0,\dfrac{\pi }{2}}\right\rbrack$ ,所以 $\left| {\vv{a} + \vv{b}}\right|  = 2\cos x$
            
            \item $f\left( x\right)  = \vv{a} \cdot  \vv{b} - \left| {\vv{a} + \vv{b}}\right|  = \cos {2x} - 2\cos x = 2{\cos }^{2}x - 2\cos x - 1  = 2{\left( \cos x - \dfrac{1}{2}\right) }^{2} - \dfrac{3}{2}$ ,因为 $x \in  \left\lbrack  {0,\dfrac{\pi }{2}}\right\rbrack$ ,则 $\cos x \in  \left\lbrack  {0,1}\right\rbrack$ , 则当 $\cos x = \dfrac{1}{2}$ ,即 $x = \dfrac{\pi }{3}$ 时,该函数取得最小值 $- \dfrac{3}{2}$ .
        \end{enumerate}
    \end{jieda}
\end{enumerate}
\subsection*{考点4 向量的夹角与向量的坐标表示}

\subsubsection*{一、填空题}
\begin{enumerate}
    \item (2024$\cdot$青浦区二模$\cdot$2)已知向量 $\vv{a} = \left( {-1,1}\right) ,\vv{b} = \left( {3,4}\right)$ ，则 $\langle \vv{a},\vv{b}\rangle  =$ \blkda[2]{$\arccos \dfrac{\sqrt{2}}{10}$}.

    \item (2024$\cdot$宝山区一模$\cdot$2)已知向量 $\vv{a} = \left( {{2m},1}\right) ,\vv{b} = \left( {{1,m} - 3}\right)$ ，若 $\vv{a}\bot \vv{b}$ ，则实数 $m =$ \blkda[2]{1}.

    \item (2024$\cdot$嘉定区一模$\cdot$2)已知 $\vv{a} = \left( {2,1}\right)$ ， $\vv{b} = \left( {-1,2}\right)$ ，则 $2\vv{a} + 3\vv{b} =$ \blkda[2]{$\left( {1,8}\right)$}.

    \item (2024$\cdot$金山区二模$\cdot$2)已知向量 $\vv{a} = \left( {1, - 3}\right)$ ， $\vv{b} = \left( {m,1}\right)$ ，若 $\vv{a}\bot \vv{b}$ ，则实数 $m$ 的值为\blkda[2]{3}.

    \item (2024$\cdot$崇明区二模$\cdot$3)已知向量 $\vv{a} = \left( {2,\lambda , - 1}\right)$ ， $\vv{b} = \left( {2,1, - 4}\right)$ ，若 $\vv{a}\bot \vv{b}$ ，则 $\lambda  =$ \blkda[2]{$-8$}.

    \item (2024$\cdot$松江区一模$\cdot$4)已知向量 $\vv{a} = \left( {1,2}\right)$ ， $\vv{b} = \left( {4,3}\right)$ ，则 $\vv{a} \cdot  \left( {2\vv{a} - \vv{b}}\right)  =$ \blkda[2]{0}.

    \item (2024$\cdot$静安区二模$\cdot$4)若单位向量 $\vv{a}$ 、 $\vv{b}$ 满足 $\vv{a}\bot \vv{b}$ ，则 $\left| {\vv{a} - \sqrt{3}\vv{b}}\right|  =$ \blkda[2]{2}.

    \item (2024$\cdot$黄浦区一模$\cdot$5)已知向量 $\vv{a} = \left( {0,2}\right) ,\vv{b} = \left( {\sqrt{3},1}\right)$ ，则向量 $\vv{a}$ 与 $\vv{b}$ 夹角的余弦值为\blkda[2]{$\dfrac{1}{2}$}.

    \item (2024$\cdot$奉贤区二模$\cdot$8)已知向量 $\vv{a} = \left( {1,1}\right) ,\vv{b} = \left( {2, - 1}\right)$ ，则 $\vv{b}$ 在 $\vv{a}$ 方向上的投影向量为\blkda[2]{$\left( {\dfrac{1}{2},\dfrac{1}{2}}\right)$}.

    \item (2025$\cdot$徐汇区二模$\cdot$8)已知 ${ABCD}$ 是正方形，点 $M$ 是 ${AB}$ 的中点，点 $E$ 在对角线 ${AC}$ 上，且 $\vv{AE} =  3\vv{EC}$ ，则 $\angle {MED}$ 的大小为\blkda[2]{$\dfrac{\pi }{2}$}.
    \begin{jieda}
        以点 $A$ 为坐标原点, $\vv{AB}$ 正方向为 $x$ 轴, $\vv{AD}$ 正方向为 $y$ 轴建立平面直角坐标系,设 ${AB} = 4$ ,则有 $A\left( {0,0}\right) \text{ 、 }M(2$ , $0)\text{ 、 }D\left( {0,4}\right) \text{ 、 }C\left( {4,4}\right)$ ,由 $\vv{AE} = 3\vv{EC}$ 有 $\vv{AE} = \dfrac{3}{4}\vv{AC} = \left( {3,3}\right)$ ,所以 $E\left( {3,3}\right)$ ,所以 $\vv{EM} = \left( {-1, - 3}\right) ,\vv{ED} = \left( {-3,1}\right)$ ,所以 $\vv{EM}$ . $\vv{ED} = 0$，所以 $\angle {MED} = \dfrac{\pi }{2}$.

        \begin{center}
        \includegraphics[width=0.2\textwidth]{bodyxb/src/100_413_1787_238_230_0.jpg}
        \end{center}
    \end{jieda}

    \item (2024$\cdot$普陀区二模$\cdot$9)若向量 $\vv{a}$ 在向量 $\vv{b}$ 上的投影为 $\dfrac{1}{3}\vv{b}$ ，且 $\left| {3\vv{a} - \vv{b}}\right|  = \left| {\vv{a} + \vv{b}}\right|$ ，则 $\cos \langle \vv{a},\vv{b}\rangle  =$ \blkda[2]{$\dfrac{\sqrt{3}}{3}$}.
    \begin{jieda}
        因为 $\vv{a}$ 在 $\vv{b}$ 上的投影为 $\dfrac{1}{3}\vv{b}$ ,所以 $\dfrac{\vv{a} \cdot  \vv{b}}{\left| \vv{b}\right| } \cdot  \dfrac{\vv{b}}{\left| \vv{b}\right| } = \dfrac{1}{3}\vv{b}$ , 则 $\dfrac{\vv{a} \cdot  \vv{b}}{{\left| \vv{b}\right| }^{2}} = \dfrac{1}{3}$ ,即 $\left| \vv{b}\right|  = \sqrt{3\vv{a} \cdot  \vv{b}}$ ,又 $\left| {3\vv{a} - \vv{b}}\right|  = \left| {\vv{a} + \vv{b}}\right|$ ,平方得 $8{\vv{a}}^{2} = 8\vv{a} \cdot  \vv{b}$ ,则 $\left| \vv{a}\right|  = \sqrt{\vv{a} \cdot  \vv{b}}$ ,即 $\cos \langle \vv{a},\vv{b}\rangle  = \dfrac{\vv{a} \cdot  \vv{b}}{\left| \vv{a}\right| \left| \vv{b}\right| } =  \dfrac{\vv{a} \cdot  \vv{b}}{\sqrt{3\vv{a} \cdot  \vv{b}}\sqrt{\vv{a} \cdot  \vv{b}}} = \dfrac{\sqrt{3}}{3}$.
    \end{jieda}

    \item (2024$\cdot$嘉定区二模$\cdot$11)在平面直角坐标系 ${xOy}$ 中，点 $P$ 在圆 ${x}^{2} + {y}^{2} = 1$ 上运动，定点 $A$ 、 $B$ 满足 $\vv{OA} \cdot  \vv{OB} = \dfrac{1}{2}$ 且 $\left| \vv{OA}\right|  = \left| \vv{OB}\right|  = 1$ ,若 $k \geq  \left| {\vv{OA} \cdot  \vv{OP}}\right|  + \left| {\vv{OB} \cdot  \vv{OP}}\right|$ 恒成立，则实数 $k$ 的取值范围为\blkda[2]{$\lbrack \sqrt{3}, + \infty)$}.
    \begin{jieda}
        以 $O$ 为原点建立直角坐标系， $A\left( {1,0}\right)$ ， $B\left( {\dfrac{1}{2},\dfrac{\sqrt{3}}{2}}\right)$ ， $P\left( {\cos \theta ,\sin \theta }\right)$ ， 则 $\left| {\vv{OA} \cdot  \vv{OP}}\right|  + \left| {\vv{OB} \cdot  \vv{OP}}\right|  \geq   \mid  (\vv{OA} +  \vv{OB}) \cdot  \vv{OP}\left|  = \right| \dfrac{3}{2}\cos \theta  + \dfrac{\sqrt{3}}{2}\sin \theta  \mid   =  \left| {\sqrt{3}\sin \left( {\theta  + \dfrac{\pi }{3}}\right) }\right|$ ,故 $k \geq  \sqrt{3}$ .

        \begin{center}
        \includegraphics[width=0.2\textwidth]{bodyxb/src/100_1341_1406_236_239_0.jpg}
        \end{center}
    \end{jieda}

    \item (2024$\cdot$崇明区一模$\cdot$11)已知不平行的两个向量 $\vv{a}$ 、 $\vv{b}$ 满足 $\left| \vv{a}\right|  = 1$ ， $\vv{a} \cdot  \vv{b} = \sqrt{3}$. 若对任意的 $t \in  \mathbf{R}$ ， 都有 $\left| {\vv{b} - t\vv{a}}\right|  \geq  2$ 成立，则 $\left| \vv{b}\right|$ 的最小值等于\blkda[2]{$\sqrt{7}$}.
    \begin{jieda}
        向量 $\vv{a}\text{ 、 }\vv{b}$ 满足 $\left| \vv{a}\right|  = 1$，$\vv{a} \cdot  \vv{b} = \sqrt{3} \Longrightarrow  \left| \vv{a}\right| \left| \vv{b}\right| \cos \theta  = \sqrt{3} \Longrightarrow  \left| \vv{b}\right| \cos \theta  = \sqrt{3}$ ,

        对任意的 $t \in  \mathbf{R}$ ,都有 $\left| {\vv{b} - t\vv{a}}\right|  \geq  2 \Longrightarrow  \left| {\vv{b} - t\vv{a}}\right|  = \left| \vv{b}\right| \sin \theta  \geq  2$ 成立,当 ${BA} \bot  {OA}$ 时 $m = \left| {\vv{b} - t\vv{a}}\right|$ 取得最小值,如图所示, $\left\{  {\begin{array}{l} \left| \vv{b}\right| \cos \theta  = \sqrt{3} \\  \left| \vv{b}\right| \sin \theta  \geq  2 \end{array} \Longrightarrow  \left| +\right| }\right.  {\left| \vv{b}\right| }^{2} \geq  7,\left| \vv{b}\right|  \geq  \sqrt{7}.$

        \begin{center}
        \includegraphics[width=0.2\textwidth]{bodyxb/src/100_1341_2020_236_140_0.jpg}
        \end{center}
    \end{jieda}

    \item (2025$\cdot$松江区一模$\cdot$11)已知平面向量 $\vv{a}$ 与 $\vv{b}$ 的夹角为 $\theta ,\vv{b} - \vv{a}$ 与 $\vv{a}$ 的夹角为 ${3\theta },\left| \vv{a}\right|  = 1,\vv{a}$ 和 $\vv{b} - \vv{a}$ 在 $\vv{b}$ 上的数量投影分别为 $x\text{ 、 }y$ ,则 $x\left( {y + \sin \theta }\right)$ 的取值范围是\blkda[2]{$\left\lbrack  {\dfrac{\sqrt{3} - 1}{4},\dfrac{\sqrt{2}}{2}}\right\rbrack$}.
    \begin{jieda}
        因为$\langle\vv{a}, \vv{b}\rangle = \theta, \langle\vv{b} - \vv{a}, \vv{a}\rangle = 3\theta$ ,所以 $\langle\vv{b} - \vv{a}, \vv{b}\rangle = 2\theta$，根据向量夹角定义可知$0 \leq  \theta  \leq  \dfrac{\pi }{3}$

        不妨设$\vv{b} = (t, 0)$，$\vv{a} = (cos\theta, sin\theta)$，因此$\vv{b} - \vv{a} = (t-cos\theta, -sin\theta)$，故$\left\{\begin{array}{l}
            x = cos\theta \\
            y = t - cos\theta
        \end{array}\right.$

        因为$\vv{b} - \vv{a}$的终点在第三象限或者第四象限，数形结合可知$\langle\vv{b} - \vv{a}, \vv{b}\rangle = 2\theta \Longleftrightarrow \dfrac{sin\theta}{t-cos\theta} = tan2\theta$，故$y = \dfrac{sin\theta}{tan2\theta}$，因此$x\left( {y + \sin \theta }\right) = cos\theta\left(\dfrac{sin\theta}{tan2\theta} + sin\theta\right) = \dfrac{sin2\theta}{2}\left(1 + \dfrac{cos2\theta}{sin2\theta}\right) = \dfrac{\sqrt{2}}{2}sin\left(2\theta + \dfrac{\pi}{4}\right)$

        因为 $\theta  \in  \left\lbrack  {0,\dfrac{\pi }{3}}\right\rbrack$ ,所以 ${2\theta } + \dfrac{\pi }{4}  \in  \left\lbrack  {\dfrac{\pi }{4},\dfrac{11}{12}\pi }\right\rbrack$ ,故当 ${2\theta } + \dfrac{\pi }{4} = \dfrac{11}{12}\pi$ 时, $\sin \left( {{2\theta } + \dfrac{\pi }{4}}\right)$ 取得最小值$\dfrac{\sqrt{3} - 1}{4}$，当 ${2\theta } + \dfrac{\pi }{4} = \dfrac{\pi }{2}$ 时， $\sin \left( {{2\theta } + \dfrac{\pi }{4}}\right)$ 取得最大值$\dfrac{\sqrt{2}}{2}$.
        
        故 $x\left( {y + \sin \theta }\right)$ 的取值范围为 $\left( {\dfrac{\sqrt{3} - 1}{4},\dfrac{\sqrt{2}}{2}}\right)$.

        \begin{center}
        \includegraphics[width=0.2\textwidth]{bodyxb/src/101_585_228_234_116_0.jpg}
        \end{center}
    \end{jieda}

    \item (2024$\cdot$金山区一模$\cdot$12)已知平面向量 $\vv{a}$ 、 $\vv{b}$ 、 $\vv{c}$ 满足 $\left| \vv{a}\right|  = 4\left| {\vv{b} - \vv{c}}\right|  = 2$ ， $\left| {\vv{a} + \vv{b}}\right|  = \left| {\vv{a} - \vv{b}}\right|  + \left| \vv{a}\right|$ ，且 $\langle \vv{a},\vv{c}\rangle  = \dfrac{\pi }{3}$ ，则 $\vv{a} \cdot  \vv{c}$ 的取值范围是\blkda[2]{}.
    \begin{jieda}
        根据题意不妨设 $\vv{a} = \left( {2,0}\right)  = \vv{OA},\vv{b} = \left( {x,y}\right)  = \vv{OB}$ , $\vv{c} = \vv{OC},O$ 为坐标原点,则 $\left| {\vv{a} + \vv{b}}\right|  = \left| {\vv{a} - \vv{b}}\right|  + \left| \vv{a}\right|  \Longrightarrow  \sqrt{{\left( x + 2\right) }^{2} + {y}^{2}} = \sqrt{{\left( x - 2\right) }^{2} + {y}^{2}} + 2$ ,即点 $B$ 到 $\left( {-2,0}\right)$ 的距离比到点 $A$ 的距离大 2,根据双曲线的定义可知 $B$ 的轨迹为双曲线的一支,以 2 为长轴,4 为焦距,则 ${x}^{2} - \dfrac{{y}^{2}}{3} = 1(x >$ 1),又 $\langle \vv{a},\vv{c}\rangle  = \dfrac{\pi }{3}$ ,易知 $C$ 点轨迹为 $l : y =  \pm  \sqrt{3}x\left( {x > 0}\right)$ ,显然 $C$ 点轨迹为 $B$ 点轨迹双曲线的渐近线
        
        \begin{center}
        \includegraphics[width=0.3\textwidth]{bodyxb/src/101_1254_224_276_222_0.jpg}
        \end{center}
        
        如图所示，由图形的对称性不妨设 $C\left( {m,\sqrt{3}m}\right)$ ,则 $\vv{a} \cdot  \vv{c} =  {2m}$ ,由题意 $\left| {\vv{b} - \vv{c}}\right|  = \dfrac{1}{2} = \left| \vv{CB}\right|$
        设$B(x_0, y_0)$，则$m = \dfrac{{x} + \sqrt{3} \cdot  \sqrt{3({x}^{2} - 1)} - \sqrt{1 - {\left( \sqrt{3}{x} - \sqrt{3({x}^{2} - 1)}\right) }^{2}}}{4}$，打表可得其最小值约为$0.692$，故$\vv{a} \cdot  \vv{c} \in [1.384, +\infty)$

        下为错误答案：

        当 ${BC} \bot  l$ 时,此时 $C$ 点横坐标 $m$ 最小，由点到直线的距离公式可知 $\left| {BC}\right|  = \dfrac{\left| \sqrt{3}x - y\right| }{2} = \dfrac{1}{2} \Longrightarrow  \left| {\sqrt{3}x - y}\right|  = 1$ ,而双曲线在渐近线 $y = \sqrt{3}x$ 下方,则 $\sqrt{3}x -  y = 1$ ,与双曲线方程联立 $\left\{  {\begin{array}{l} y = \sqrt{3}x - 1 \\  3{x}^{2} - {y}^{2} = 3 \end{array} \Longrightarrow  2\sqrt{3}x = 4 \Longrightarrow  }\right.  \left\{  {\begin{array}{l} x = \dfrac{2\sqrt{3}}{3} \\  y = 1 \end{array}\text{ ,即 }B\left( {\dfrac{2\sqrt{3}}{3},1}\right) \text{ ,则 }{l}_{BC} : y =  - \dfrac{1}{\sqrt{3}}\left( {x - \dfrac{2\sqrt{3}}{3}}\right)  + 1,}\right.$ 联立 $\left\{  {\begin{array}{l} y =  - \dfrac{1}{\sqrt{3}}\left( {x - \dfrac{2\sqrt{3}}{3}}\right)  + 1 \\  y = \sqrt{3}x \end{array} \Longrightarrow  x = \dfrac{5}{4\sqrt{3}}}\right.$ ,即 ${m}_{\min } = \dfrac{5}{4\sqrt{3}} =  \dfrac{5\sqrt{3}}{12} \Longrightarrow  \vv{a} \cdot  \vv{c} = {2m} \geq  \dfrac{5\sqrt{3}}{6}$ ,由双曲线的性质可知满足 $\left| {BC}\right|  =  \dfrac{1}{2}$ 的点 $C$ 横坐标无上限,故 $\vv{a} \cdot  \vv{c}$ 的取值范围是 $\left[ \dfrac{5\sqrt{3}}{6}, + \infty \right)$.
    \end{jieda}


    \item (2024$\cdot$金山区一模$\cdot$12改)已知平面向量 $\vv{a}$ 、 $\vv{b}$ 、 $\vv{c}$ 满足 $\left| \vv{a}\right| = |\vv{c}| = 2$ ， $\left| {\vv{a} + \vv{b}}\right|  = \left| {\vv{a} - \vv{b}}\right|  + \left| \vv{a}\right|$ ，且 $\langle \vv{a},\vv{c}\rangle  = \dfrac{\pi }{3}$，若有$\vv{b} = \lambda \vv{a} + (1-\lambda)\vv{c}$，则 $\lambda = $\blkda[2]{$\dfrac{1}{4}$}.
    \begin{jieda}
        根据题意不妨设 $\vv{a} = \left( {2,0}\right)  = \vv{OA},\vv{b} = \left( {x,y}\right)  = \vv{OB}$ , $\vv{c} = \vv{OC},O$ 为坐标原点,则 $\left| {\vv{a} + \vv{b}}\right|  = \left| {\vv{a} - \vv{b}}\right|  + \left| \vv{a}\right|  \Longrightarrow  \sqrt{{\left( x + 2\right) }^{2} + {y}^{2}} = \sqrt{{\left( x - 2\right) }^{2} + {y}^{2}} + 2$ ,即点 $B$ 到 $\left( {-2,0}\right)$ 的距离比到点 $A$ 的距离大 2,根据双曲线的定义可知 $B$ 的轨迹为双曲线的一支,以 2 为长轴,4 为焦距,则 ${x}^{2} - \dfrac{{y}^{2}}{3} = 1(x >$ 1),又 $\langle \vv{a},\vv{c}\rangle  = \dfrac{\pi }{3}$ ,易知 $C$ 点轨迹为 $l : y =  \pm  \sqrt{3}x\left( {x > 0}\right)$ ,显然 $C$ 点轨迹为 $B$ 点轨迹双曲线的渐近线，又因为$|\vv{c}| = 2$，所以不妨设$C(1, \sqrt{3})$，根据$\vv{b} = \lambda \vv{a} + (1-\lambda)\vv{c}$可知$A, B, C$共线，联立求得$B\left(\dfrac{5}{4}, \dfrac{3\sqrt{3}}{4}\right)$，故$\lambda = \dfrac{1}{4}$.
    \end{jieda}
\end{enumerate}

\subsubsection*{二、选择题}
\begin{enumerate}
    \item (2024$\cdot$嘉定区二模$\cdot$14)已知 $\vv{OA} = \left( {{x}_{1},{y}_{1}}\right) ,\vv{OB} = \left( {{x}_{2},{y}_{2}}\right)$ ，且 $\vv{OA}$ 、 $\vv{OB}$ 不共线，则 $\bigtriangleup {OAB}$ 的面积为 \dotfill{(\kh{B})}
    \xx{$\dfrac{1}{2}\left| {{x}_{1}{x}_{2} - {y}_{1}{y}_{2}}\right|$}{$\dfrac{1}{2}\left| {{x}_{1}{y}_{2} - {x}_{2}{y}_{1}}\right|$}{$\dfrac{1}{2}\left| {{x}_{1}{x}_{2} + {y}_{1}{y}_{2}}\right|$}{$\dfrac{1}{2}\left| {{x}_{1}{y}_{2} + {x}_{2}{y}_{1}}\right|$}
\end{enumerate}

\section*{第 9 章 复 数}

\begin{center}

\begin{tabular}{|p{4cm}|p{12cm}|}
\hline
考点 & 课标要求与命题趋势 \\
\cline{1-2}
1. 复数的概念与运算 & \ding{172} 通过方程的解，认识复数；掌握复数的分类依据、复数的实部、虚部及纯虚数等概念的数学意义 \ding{173}  理解、掌握复数的代数表示及其几何意义，理解并掌握共轭复数和两个复数相等的含义 \ding{174}  掌握复数代数表示的四则运算，了解复数加、减运算的几何意义 \\
\cline{1-2}
2. 复数的几何意义 &  \ding{172} 熟练掌握复数的几何意义; 该内容是上海高考必考内容,一般考查复数的四则运算、共轭复数、模长运算、几何意义，问题较为简单 \ding{173}  复数是一类重要的运算对象，有广泛的应用. 通过方程求解，理解引入复数的必要性. 了解数系的扩充，掌握复数的表示、运算及其几何意义 \\
\cline{1-2}
\hline
\end{tabular}
\end{center}
\clearpage
\subsection*{考点1 复数的概念与运算}

\subsubsection*{一、填空题}
\begin{enumerate}
    \item (2025$\cdot$宝山区二模$\cdot$1) $i$ 是虚数单位，则 $|1 + i| =$ \blkda[2]{$\sqrt{2}$}.

    \item (2024$\cdot$奉贤区二模$\cdot$1)已知复数 $z = \left( {3 - {4i}}\right)$ ，.i(i为虚数单位)，则 $z =$ \blkda[2]{$4 + 3{i}$}.

    \item (2024$\cdot$奉贤区一模$\cdot$1)若 $2 + a{i} = \left( {b{i} - 1}\right) {i}\left( {a\text{ 、 }b \in  \mathbf{R}}\right)$ ，其中 ${i}$ 是虚数单位，则 $a + b{i} =$ \blkda[2]{$- 1 - 2{i}$}.

    \item (2024$\cdot$徐汇区二模$\cdot$2)已知复数 $z = \dfrac{1 - {i}}{{i}}$ (i 为虚数单位)，则 $z \cdot  \overline{z} =$ \blkda[2]{$2$}.

    \item (2025$\cdot$普陀区二模$\cdot$2)已知复数 $z = \dfrac{3 - 2{i}}{{i}}$ ，其中 ${i}$ 为虚数单位，则 $z + \overline{z} =$ \blkda[2]{$-4$}.

    \item (2025$\cdot$徐汇区二模$\cdot$2)已知 ${i}$ 为虚数单位，复数 $z$ 满足 $z{i} - {i} = 1$ ，则 $\left| z\right|  =$ \blkda[2]{$\dfrac{1}{2}$}.

    \item (2025$\cdot$奉贤区二模$\cdot$2)已知 i 为虚数单位，复数 $z$ 满足 $z{i} - {i} = 1$ ，则 $\left| z\right|  =$ \blkda[2]{$\sqrt{2}$}.

    \item (2025$\cdot$崇明区二模$\cdot$2)已知复数 $\dfrac{z}{{i}} = 1 - 2{i}$ ( ${i}$ 为虚数单位)，则 $z =$ \blkda[2]{$2 + {i}$}.

    \item (2024$\cdot$静安区二模$\cdot$2)已知 ${i}$ 是虚数单位，复数 $z = \dfrac{m + {i}}{2 + {i}}$ 是纯虚数，则实数 $m$ 的值为\blkda[2]{$- \dfrac{1}{2}$}.

    \item (2024$\cdot$闵行区二模$\cdot$2)复数 $z$ 满足(2+i) $z = 3 + {4i}$ (i) 为虚数单位)，则 $\left| z\right|  =$ \blkda[2]{$\sqrt{5}$}.

    \item (2024$\cdot$浦东新区二模$\cdot$2)若复数 $z = 1 + 2{i}$ ( ${i}$ 是虚数单位)，则 $z \cdot  \overline{z} - z =$ \blkda[2]{$4 - 2{i}$}.

    \item (2024$\cdot$青浦区一模$\cdot$2)若复数 $z$ 满足 ${i}z = 3 + {i}$ ，则 $\left| z\right|  =$ \blkda[2]{$\sqrt{10}$}.

    \item (2024$\cdot$普陀区一模$\cdot$2)设 ${i}$ 为虚数单位，若复数 $z$ 满足 ${i}z = 1 + 2{i}$ ，则 $\left| {z - 1}\right|  =$ \blkda[2]{$\sqrt{2}$}.

    \item (2024$\cdot$杨浦区一模$\cdot$2)若复数 $z$ 满足 ${i}z =  - 2 + {i}$ (其中 ${i}$ 为虚数单位)，则 $\left| z\right|  =$ \blkda[2]{$\sqrt{5}$}.

    \item (2024$\cdot$长宁区一模$\cdot$2)复数 $z$ 满足 $z = \dfrac{1}{1 - {i}}$ ( ${i}$ 为虚数单位)，则 $\left| \overline{z}\right|  =$ \blkda[2]{$\dfrac{\sqrt{2}}{2}$}.

    \item (2025$\cdot$闵行区二模$\cdot$2)已知复数 $z$ 满足$(1+i) z = 1 - {i}$ ($i$为虚数单位)，则 $z =$ \blkda[2]{$-i$}.

    \item (2024$\cdot$松江区一模$\cdot$3)已知复数 $z = 2 + {i}$ (其中 ${i}$ 是虚数单位)，则 $\left| z\right|  =$ \blkda[2]{$\sqrt{5}$}.

    \item (2024$\cdot$杨浦区二模$\cdot$3)计算 $\dfrac{2 - {i}}{3 + {i}} =$ \blkda[2]{$\dfrac{1}{2} - \dfrac{1}{2}{i}$}.(其中 ${i}$ 为虚数单位)

    \item (2024 青浦区二模$\cdot$3)已知复数 $z = \dfrac{5{i}}{1 - {i}}$ ，则 ${Im}z =$ \blkda[2]{$\dfrac{5}{2}$}.

    \item (2024$\cdot$崇明区一模$\cdot$3)若复数 $z = {m}^{2} - 4 + \left( {m + 2}\right) {i}$ ( ${i}$ 为虚数单位)是纯虚数，则实数 $m$ 的值为\blkda[2]{$2$}.

    \item (2025$\cdot$浦东新区一模$\cdot$3)已知复数 ${z}_{1} = 3 + {i},{z}_{2} = a + 4{i},a \in  \mathbf{R}$ ，若 ${z}_{1} - {z}_{2}$ 为纯虚数，则 $\left| {z}_{2}\right|  =$ \blkda[2]{$5$}.

    \item (2025$\cdot$普陀区一模$\cdot$3)设 ${i}$ 为虚数单位，若复数 $z$ 满足 $z \cdot  \overline{z} + z - \overline{z} = 4 + 2{i}$ ，则 $\left| z\right|  =$ \blkda[2]{$2$}.

    \item (2025$\cdot$崇明区一模$\cdot$3)若复数 $z$ 满足 ${2z} + \overline{z} = 1 + {i}$ ，其中 ${i}$ 为虚数单位，则 $z =$ \blkda[2]{$\dfrac{1}{3} + {i}$}.

    \item (2025$\cdot$闵行区二模$\cdot$3)已知 ${i}$ 是虚数单位，则 $\left| \dfrac{1 + {i}}{{i}}\right|  =$ \blkda[2]{$\sqrt{2}$}.

    \item (2025$\cdot$上海春季卷$\cdot$3)已知复数 $z = \dfrac{2 + {i}}{{i}}$ ，其中 ${i}$ 为虚数单位，则 $\left| z\right|  =$ \blkda[2]{$\sqrt{5}$}.

    \item (2024$\cdot$上海春季卷$\cdot$3)若复数 $z$ 满足 $\dfrac{z}{1 + {i}} = {i}$ ( ${i}$ 为虚数单位)，则 $\overline{z} =$ \blkda[2]{$- 1 - {i}$}.

    \item (2025$\cdot$宝山区一模$\cdot$3)已知 ${i}$ 为虚数单位，若(a-2)+(2a-1)i为纯虚数，则实数 $a =$ \blkda[2]{$2$}.

    \item (2025$\cdot$嘉定区一模$\cdot$3) 如果复数 $z$ 满足 ${i} \cdot  \overline{z} = 1 + 2{i}$ ( ${i}$ 为虚数单位)，则 $z =$ \blkda[2]{$2 + {i}$}.

    \item (2025$\cdot$静安区一模$\cdot$3)已知 ${i}$ 是虚数单位， $\left( {m + {i}}\right)$ (1-2i)是纯虚数，则实数 $m$ 的值为\blkda[2]{$-2$}.

    \item (2024$\cdot$黄浦区一模$\cdot$3)已知复数 $z = 1 - {i}$ (i 为虚数单位)，则满足 $\overline{z} \cdot  w = z$ 的复数 $w$ 为\blkda[2]{$- {i}$}.

    \item (2024$\cdot$静安区一模$\cdot$4)已知 $a \in  \mathbf{R}$ ， ${i}$ 是虚数单位， $\dfrac{1}{a - {i}}$ 的虚部为\blkda[2]{$\dfrac{1}{{a}^{2} + 1}$}.

    \item (2025$\cdot$金山区一模$\cdot$4)已知复数 $z = 2 + {i}$ ，其中 ${i}$ 为虚数单位，则 $\left| {{i}z}\right|$ 的值为\blkda[2]{$\sqrt{5}$}.

    \item (2024$\cdot$崇明区二模$\cdot$4)若复数 $z$ 满足 ${i}z = 1 + {i}$ ( ${i}$ 为虚数单位)，则 $z =$ \blkda[2]{$1-i$}.

    \item (2024$\cdot$金山区二模$\cdot$4)已知复数 $z$ 满足 $z + {2\overline{z}} = 3 - {i}$ ，则 $\overline{z}$ 的模为\blkda[2]{$\sqrt{2}$}.

    \item (2024$\cdot$嘉定区二模$\cdot$5)已知 ${i}$ 是虚数单位，则 $\left| \dfrac{4 + {3{i}}}{{\left( 1 - 2{i}\right) }^{2}}\right|  =$ \blkda[2]{$1$}.

    \item (2024$\cdot$宝山区二模$\cdot$5)设实数 $x$ 、 $y$ 满足( $x + y$ i) $i - 2 + 4{i} = \left( {x - y{i}}\right) \left( {1 + {i}}\right)$ ( ${i}$ 为虚数单位)，则 $x + y =$ \blkda[2]{$2$}.

    \item (2025$\cdot$松江区一模$\cdot$5)若复数 $z$ 满足 ${i} \cdot  z = 2 + 3{i}$ (其中 ${i}$ 是虚数单位)，则复数 $z$ 的共轭复数 $\overline{z} =$ \blkda[2]{$3 + 2{i}$}.

    \item (2025$\cdot$黄浦区二模$\cdot$5)$i$为虚数单位，若复数 $z$ 满足 $z - \overline{z} = 2{i}$ 且 $\overline{z} = {i}z$ ，则 ${Re}z =$ \blkda[2]{$-1$}.

    \item (2023$\cdot$上海秋季卷$\cdot$6)已知复数 $z = 1 + {i}$ ，则 $\left| {1 - {i}z}\right|  =$ \blkda[2]{$\sqrt{5}$}.

    \item (2024$\cdot$杨浦区二模$\cdot$9)设复数 ${z}_{1}$ 与 ${z}_{2}$ 所对应的点为 ${Z}_{1}$ 与 ${Z}_{2}$ ，若 ${z}_{1} = 1 + {i},{z}_{2} = {i} \cdot  {z}_{1}$ ，则 $\left| \vv{{Z}_{1}{Z}_{2}}\right|  =$ \blkda[2]{$2$}.
    \begin{jieda}
        $Z_1(1, 1), Z_2(-1, 1)$，故$ \vv{{Z}_{1}{Z}_{2}} = (-2, 0)$
    \end{jieda}

    \item (2024$\cdot$黄浦区二模$\cdot$8)若实系数一元二次方程 ${x}^{2} + {ax} + b = 0$ 有一个虚数根的模为 4，则 $a$ 的取值范围是\blkda[2]{$\left( {-8,8}\right)$}.
    \begin{jieda}
        $16 = \left|z\right|^2 = z_1z_2 = b$，$\Delta = a^2 - 4b = a^2 - 64 < 0$，故$a \in (-8, 8)$
    \end{jieda}

    \item (2024$\cdot$上海秋季卷$\cdot$9)已知虚数 $z$ ，其实部为 1，且 $z + \dfrac{2}{z} = m\left( {m \in  \mathbf{R}}\right)$ ，则实数 $m$ 为\blkda[2]{$2$}.
    \begin{jieda}
        $m = 2Rez = 2$(当且仅当$\dfrac{2}{z} = \overline{z}$，$m \in \mathbf{R}$).
    \end{jieda}

    \item (2025$\cdot$浦东新区二模$\cdot$6)若关于 $x$ 的方程 ${x}^{2} - x + m = 0$ 的一个虚根的模为 2，则实数 $m$ 的值为\blkda[2]{$4$}.
    \begin{jieda}
        $4 = \left|z\right|^2 = z_1z_2 = m$.
    \end{jieda}

    \item (2025$\cdot$闵行区一模$\cdot$9)已知 $f\left( n\right)  = {{i}}^{n + 1} + {{i}}^{n + 2} + {{i}}^{n + 3} + {{i}}^{n + 4} + {{i}}^{n + 5}$ ( ${i}$ 为虚数单位， $n$ 为正整数)，当 ${n}_{1}$ 、 ${n}_{2}$ 取遍所有正整数时， $f\left( {n}_{1}\right)  + f\left( {n}_{2}\right)$ 的值中不同虚数的个数为\blkda[2]{$6$}.
    \begin{jieda}
        $f(n) \in \left\{1,-1, i, -i\right\}$，故$f\left( {n}_{1}\right)  + f\left( {n}_{2}\right) \in \left\{2, -2, 2i, -2i,1+i, 1-i, -1+i, -1-i\right\}$
    \end{jieda}

    \item (2025$\cdot$徐汇区一模$\cdot$8)已知复数 ${z}_{1}$ 和复数 ${z}_{2}$ 满足 ${z}_{1} + {z}_{2} = 3 + 4{i},\overline{{z}_{1}} - \overline{{z}_{2}} =  -2 + {i}$ ($i$ 为虚数单位),则 $\left| {{z}_{1}^{2} - {z}_{2}^{2}}\right|  =$\blkda[2]{$5\sqrt{5}$}.
    \begin{jieda}
        $z_1 - z_2 = -2-i$，$\left| {{z}_{1}^{2} - {z}_{2}^{2}}\right|  = \left|(z_1+z_2) \cdot (z_1-z_2)\right| = \left|z_1+z_2\right| \cdot \left|z_1-z_2\right| = 5\sqrt{5}$.
        
        \ps $\left\{\begin{array}{l}
            z_1 = \dfrac{1}{2} + \dfrac{3}{2}i \\
            z_2 = \dfrac{5}{2} + \dfrac{5}{2}i
        \end{array}\right.$也可.
    \end{jieda}
\end{enumerate}

\subsubsection*{二、选择题}
\begin{enumerate}
    \item (2024$\cdot$虹口区二模$\cdot$13)欧拉公式 ${{e}}^{\theta } = \cos \theta  + {i}\sin \theta$ 把自然对数的底数 ${e}$ ，虚数单位 i，三角函数 $\cos \theta$ 和 $\sin \theta$ 联系在一起，被称为 “数学的天桥”. 若复数 $z$ 满足 $\left( {z + {{e}}^{{i} \cdot  \frac{3\pi }{2}}}\right)  \cdot  {i} = 2 - {i}$ ，则 $\left| z\right|  =$ \dotfill{(\kh{A})}
    \xx{$\sqrt{2}$}{$2\sqrt{2}$}{$\sqrt{5}$}{$\sqrt{10}$}
    \begin{jieda}
        $\cos \dfrac{3\pi }{2} + {i}\sin \dfrac{3\pi }{2} =  - {i},\left( {z - {i}}\right) {i} = 2 - {i} \Longrightarrow  z = \dfrac{2 - {i}}{{i}} +  {i} =  - 1 - {i} \Longrightarrow  \left| z\right|  = \sqrt{2}$ 
    \end{jieda}
    \item (2024$\cdot$虹口区一模$\cdot$13)设 ${i}$ 为虚数单位，若 $z = \dfrac{2 - {i}}{1 + {{i}}^{2} - {{i}}^{5}}$ ，则 $\overline{z} =$ \dotfill{(\kh{A})}
    \xx{$1 - 2{i}$}{$1 + 2{i}$}{$2 - {i}$}{$2 + {i}$}
    \begin{jieda}
        $z = \dfrac{2 - {i}}{1 + {{i}}^{2} - {{i}}^{5}} = \dfrac{2 - {i}}{1 - 1 - {i}} = 1 + 2{i} \Longrightarrow  \bar{z} = 1 - 2{i}$
    \end{jieda}
    \item (2025$\cdot$长宁区一模$\cdot$13)已知复数 $z$ 和 $\overline{z}$ ，则下列说法正确的是 \dotfill{(\kh{A})}
    \xx{$z + \overline{z}$ 一定是实数}{$z - \overline{z}$ 一定是虚数}{若 $z + \overline{z} = 0$ ,则 $z$ 是纯虚数}{若 $z - \overline{z} = 0$ ,则 $z$ 是纯虚数} 
\end{enumerate}
\subsection*{考点2 复数的几何意义}

\subsubsection*{一、填空题}
\begin{enumerate}
    \item (2024$\cdot$普陀区二模$\cdot$1)已知复数 $z = 1 + {i}$ ，其中 ${i}$ 为虚数单位，则 $\overline{z}$ 在复平面内所对应的点的坐标为\blkda[2]{$\left( {1, - 1}\right)$}.

    \item (2025$\cdot$青浦区一模$\cdot$1)在复平面内，复数 $z = 1 + \dfrac{1}{2}{i}$ (其中 ${i}$ 是虚数单位)的共轭复数对应的点位于第\blkda[2]{四}象限.

    \item (2024$\cdot$松江区二模$\cdot$2)在复平面内，复数 $z$ 对应点的坐标是$(1,2)$，则 ${i} \cdot  z =$ \blkda[2]{$- 2 + {i}$}.

    \item (2024$\cdot$金山区一模$\cdot$2)在复平面内，若复数 $z$ 对应的点的坐标是$(-1，\sqrt{3})$，则 $z$ 的共轭复数 $\overline{z} =$ \blkda[2]{$- 1 - \sqrt{3}{i}$}.

    \item (2025$\cdot$虹口区一模$\cdot$7)已知非零复数 $z$ 满足 $\left| {z - 1}\right|  = 1,\left| {\overline{z} - {i}}\right|  = 1$ ，则 $z$ 的虚部为\blkda[2]{$-1$}.
    \begin{jieda}
        由图知,两圆交于 $\left( {0,0}\right)$ 、 $\left( {1, - 1}\right)$ ,因为 $z \neq  0$ ,所以 $z = 1 - {i} \Longrightarrow  Imz =  - 1$.

        \begin{center}
            \includegraphics[width=0.2\textwidth]{bodyxb/src/104_632_949_237_241_0.jpg}
        \end{center}
    \end{jieda}

    \item (2025$\cdot$杨浦区二模$\cdot$7)已知复数 $z$ 满足 $\left| {z - 1 + {i}}\right|  = 1$ ，其中 ${i}$ 为虚数单位，则 $\left| z\right|$ 的最小值为\blkda[2]{$\sqrt{2} - 1$}.
    \begin{jieda}
        在复平面内, $\left| {z - 1 + {i}}\right|  = 1$ 表示复数 $z$ 对应的点 $Z$ 与复数 $1 - {i}$ 对应点 $A\left( {1, - 1}\right)$ 的距离为 1,因此点 $Z$ 的轨迹是以 $A\left( {1, - 1}\right)$ 为圆心,1 为半径的圆, $\left| z\right|$ 表示点 $Z$ 到原点 $O$ 的距离，所以 $\left| z\right|$ 的最小值为 $\left| {OA}\right|  - 1 = \sqrt{2} - 1$.
    \end{jieda}

    \item (2025$\cdot$嘉定区二模$\cdot$8)已知复数 ${z}_{1}$ 、 ${z}_{2}$ 满足 $\left| {z}_{1}\right|  = 1$ ， $\left| {z}_{2}\right|  = 2$ ， $\left| {{z}_{1} - {z}_{2}}\right|  = \sqrt{7}$ ，则 $\left| {{z}_{1} + {z}_{2}}\right|$ 的值为\blkda[2]{$\sqrt{3}$}.
    \begin{jieda}
        设 $\vv{OA}$ 对应的复数为 ${z}_{1},\vv{OB}$ 对应的复数为 ${z}_{2}$ ,则 $\vv{OA} + \vv{OB}$ 对应的复数为 ${z}_{1} + {z}_{2},\vv{OA} - \vv{OB}$ 对应的复数为 ${z}_{1} -  {z}_{2}$ ,因为 $\left| {z}_{1}\right|  = 1,\left| {z}_{2}\right|  = 2,\left| {{z}_{1} - {z}_{2}}\right|  = \sqrt{7}$ ,由平行四边形的性质可得: ${\left| {z}_{1} + {z}_{2}\right| }^{2} + {\left| {z}_{1} - {z}_{2}\right| }^{2} = 2\left( {{\left| {z}_{1}\right| }^{2} + {\left| {z}_{2}\right| }^{2}}\right)$ ,所以 $\left| {{z}_{1} + {z}_{2}}\right|  = \sqrt{2{\left| {z}_{1}\right| }^{2} + 2{\left| {z}_{2}\right| }^{2} - {\left| {z}_{1} - {z}_{2}\right| }^{2}} = \sqrt{2 + 8 - 7} =  \sqrt{3}$
    \end{jieda}

    \item (2025$\cdot$奉贤区一模$\cdot$8)在复平面内， $O$ 为坐标原点，复数 ${z}_{1} = {i}\left( {-4 - {3{i}}}\right)$ ， ${z}_{2} = {12} + {5{i}}$ 对应的点分别为 ${Z}_{1}\text{ 、 }{Z}_{2}$ ,其中 ${i}$ 为虚数单位,则 $\left\langle  {\vv{O{Z}_{1}},\vv{O{Z}_{2}}}\right\rangle$ 的大小为\blkda[2]{$\arccos \dfrac{16}{65}$}.
    \begin{jieda}
        ${z}_{1} = {i}\left( {-4 - 3{i}}\right)  = 3 - 4{i},{z}_{2} = {12} + 5{i} \Longrightarrow  {Z}_{1}\left( {3, - 4}\right)$ , ${Z}_{2}\left( {{12},5}\right) ,\vv{O{Z}_{1}} = \left( {3, - 4}\right) ,\vv{O{Z}_{2}} = \left( {{12},5}\right) ,\cos \left\langle  {\vv{O{Z}_{1}},\vv{O{Z}_{2}}}\right\rangle   =  \dfrac{\vv{O{Z}_{1}} \cdot  \vv{O{Z}_{2}}}{\left| \vv{O{Z}_{1}}\right|  \cdot  \left| \vv{O{Z}_{2}}\right| } = \dfrac{{36} - {20}}{\sqrt{{3}^{2} + {\left( -4\right) }^{2}} \cdot  \sqrt{{12}^{2} + {5}^{2}}} = \dfrac{16}{65}$ ,则 $\left\langle  {\vv{O{Z}_{1}},\vv{O{Z}_{2}}}\right\rangle$ 的大小为 $\arccos \dfrac{16}{65}$.
    \end{jieda}

    \item (2025$\cdot$虹口区二模$\cdot$8)已知 $z$ 是实系数一元二次方程 ${x}^{2} + {px} + q = 0$ 的一个虚根，且 $\left| {z - 1}\right|  = 2$ ， 若 $z$ 在复平面上所对应的点在抛物线 ${y}^{2} = {4x}$ 上,则 $p =$ \blkda[2]{$-2$}.
    \begin{jieda}
        设 $z = a + b{i}\left( {a\text{ 、 }b \in  \mathbf{R}}\right)$ ,因为 $z$ 在复平面上所对应的点在抛物线 ${y}^{2} = {4x}$ 上,所以 ${b}^{2} = {4a},a > 0$ ，因为 $\left| {z - 1}\right|  =  \left| {a - 1 + b{i}}\right|  = 2$ ,故 ${\left( a - 1\right) }^{2} + {b}^{2} = 4$ ,即 ${\left( a - 1\right) }^{2} + {4a} = 4$ ,解得: $a = 1$ 或 $a =  - 3$ (舍去),故 ${b}^{2} = {4a} = 4,b =  \pm  2$ . 当 $z =  1 + 2{i}$ 时,代入方程 ${x}^{2} + {px} + q = 0$ 得, ${\left( 1 + 2{i}\right) }^{2} + p\left( {1 + 2{i}}\right)  + q = 0$ ,即 $p + q - 3 + \left( {{2p} + 4}\right) {i} = 0$ ,故 ${2p} + 4 = 0,p =  - 2$ ; 当 $z = 1 - 2{i}$ 时,代入方程 ${x}^{2} + {px} + q = 0$ 得, ${\left( 1 - 2{i}\right) }^{2} +  p\left( {1 - 2{i}}\right)  + q = 0$ ,即 $p + q - 3 - \left( {{2p} + 4}\right) {i} = 0$ ,故 ${2p} + 4 = 0$ , $p =  - 2$ . 综上可得, $p =  - 2$.
    \end{jieda}

    \item (2025$\cdot$青浦区二模$\cdot$9)已知复数 $z$ 、 $w$ 满足 $\left| z\right|  = 2,z = \left( {1 + {i}}\right) w$ (i 是虚数单位)，则 $\left| {w + 2}\right|$ 的最大值是\blkda[2]{$2 + \sqrt{2}$}.
    \begin{jieda}
        已知 $\left| z\right|  = 2,z = \left( {1 + {i}}\right) w$ ,则 $\left| z\right|  = \left| {\left( {1 + {i}}\right) w}\right|  =  \mid  1 +  {i}\left| \right| w\left| \text{ . 因为 }\right| 1 + {i} \mid   = \sqrt{{1}^{2} + {1}^{2}} = \sqrt{2}$ ,所以 $\left| w\right|  = \dfrac{\left| z\right| }{\left| 1 + {i}\right| } = \dfrac{2}{\sqrt{2}} =  \sqrt{2},\left| {w + 2}\right|$ 表示复数 $w$ 所对应的点到 $\left( {-2,0}\right)$ 的距离, $\left| w\right|  =  \sqrt{2}$ 说明 $w$ 对应的点在以原点为圆心, $\sqrt{2}$ 为半径的圆上,所以 $\left| {w + 2}\right|$ 的最大值为圆心 $\left( {0,0}\right)$ 到点 $\left( {-2,0}\right)$ 的距离加上半径, 即 $2 + \sqrt{2}$.
    \end{jieda}

    \item (2024 嘉定区一模$\cdot$11)已知复平面上一个动点 $Z$ 对应复数 $z$ ，若 $\left| {z - 4{i}}\right|  \leq  2$ ，其中 ${i}$ 是虚数单位，则向量 $\vv{OZ}$ 扫过的面积为\blkda[2]{$\dfrac{8}{3}\pi  + 4\sqrt{3}$}.
    \begin{jieda}
        如图所示: 向量 $\vv{OZ}$ 扫过的面积为 $S,S = \dfrac{2\sqrt{3} \times  2}{2} \times  2 - \dfrac{1}{3}\pi  \times  {2}^{2} + {4\pi } = \dfrac{8\pi }{3} + 4\sqrt{3}.$

        \begin{center}
            \includegraphics[width=0.2\textwidth]{bodyxb/src/104_1339_748_238_280_0.jpg}
        \end{center}
    \end{jieda}

    \item (2025$\cdot$上海秋季卷$\cdot$10)已知复数 $z$ 满足 ${z}^{2} = {\left( \overline{z}\right) }^{2},\left| z\right|  \leq  1$ ，则 $\left| {z - 2 - {3i}}\right|$ 的最小值是\blkda[2]{$2\sqrt{2}$}.
    \begin{jieda}
        ${z}^{2} = {\left( \overline{z}\right) }^{2} = \overline{z^2}$，故$z^2 \in \mathbf{R}$，即$z$是实数或纯虚数.
        
        由复数的几何意义知 $z$ 在复平面内对应的点 $Z\left( {a,b}\right)$ 在单位圆内部(含边界)的坐标轴上运动， 如图所示即线段 ${AB}\text{ 、 }{CD}$ 上运动,设 $E\left( {2,3}\right)$ ,则 $\left| {z - 2 - 3{i}}\right|  = \left| {ZE}\right|$ ,由图像可知 $\left| {BE}\right|  = \sqrt{10} > \left| {CE}\right|  = 2\sqrt{2}$ ,所以 ${\left| ZE\right| }_{\min } = 2\sqrt{2}$.
        \begin{center}
            \includegraphics[width=0.2\textwidth]{bodyxb/src/104_1340_1138_236_243_0.jpg}
        \end{center}
    \end{jieda}

    \item (2023$\cdot$上海春季卷$\cdot$11)设 ${z}_{1}\text{ 、 }{z}_{2} \in  \mathbf{C}$ 且 ${z}_{1} = {i}\overline{{z}_{2}}$ ，满足 $\left| {{z}_{1} - 1}\right|  = 1$ ，则 $\left| {{z}_{1} - {z}_{2}}\right|$ 的取值范围为\blkda[2]{$\left\lbrack  {0,2 + \sqrt{2}}\right\rbrack$}.
    \begin{jieda}
        满足 $\left| {{z}_{1} - 1}\right|  = 1,{z}_{1}$ 以 $G(1,0)$ 为圆心,半径为 1 的圆
        
        ${z}_{1} = {i}\overline{{z}_{2}} \Longleftrightarrow z_1 = \overline{\left(\dfrac{z_1}{i}\right)}$，$z_2$是$z_1$逆时针旋转$\dfrac{\pi}{2}$再关于实轴对称，故${z}_{2}$ 以 $M\left( {0,1}\right)$ 为圆心,半径为 1 的圆,所以 $\left| {{z}_{1} - {z}_{2}}\right|$ 的取值范围 $\lbrack 0$ , $2 + \sqrt{2}\rbrack$ .

        \begin{center}
        \includegraphics[width=0.2\textwidth]{bodyxb/src/104_1343_1638_236_236_0.jpg}
        \end{center}
    \end{jieda}

    \item (2025$\cdot$黄浦区一模$\cdot$10)$i$ 为虚数单位，若复数 ${z}_{1}$ 满足 $\left| {{z}_{1} - 1 + {i}}\right|  \leq  \sqrt{2}$ ，复数 ${z}_{2}$ 满足 $\left| {z}_{2}\right|  =  \left|{z}_{2} + 1-i\right|$,则 $\left| {{z}_{1} - {z}_{2}}\right|$ 的最小值为\blkda[2]{$\dfrac{\sqrt{2}}{2}$}.
    \begin{jieda}
        设 ${z}_{1}$ 对应点为 ${Z}_{1}$ ,则 ${Z}_{1}$ 的轨迹是以 $C(1, -1)$为圆心、 $\sqrt{2}$ 为半径的圆面(包括边界)内，设 $A\left( {-1,1}\right)$ ， ${z}_{2}$ 对应点为 ${Z}_{2}$ ,因为 $\left| {z}_{2}\right|  = \left| {{z}_{2} + 1 - {i}}\right|$ ,所以 $\left| {{Z}_{2}O}\right|  = \left| {{Z}_{2}A}\right|$ ,得 ${Z}_{2}$ 的轨迹是线段 ${OA}$ 的中垂线 $l : y = x + 1$ ,故 ${\left| {z}_{1} - {z}_{2}\right| }_{\min } =  {\left| {Z}_{1}{Z}_{2}\right| }_{\min } = {d}_{C \rightarrow  l} - \sqrt{2} = \dfrac{\left| 1 - \left( -1\right)  + 1\right| }{\sqrt{2}} - \sqrt{2} = \dfrac{3\sqrt{2}}{2} - \sqrt{2} = \dfrac{\sqrt{2}}{2}.$
    \end{jieda}

    \item (2025$\cdot$浦东新区一模$\cdot$12)已知在复数集中，等式 ${x}^{4} + {a}_{3}{x}^{3} + {a}_{2}{x}^{2} + {a}_{1}x + {a}_{0} = \left( {x - {z}_{1}}\right) (x -  \left. {z}_{2}\right) \left( {x - {z}_{3}}\right) \left( {x - {z}_{4}}\right)$ 对任意复数 $x$ 恒成立,复数 ${z}_{1}\text{ 、 }{z}_{2}\text{ 、 }{z}_{3}\text{ 、 }{z}_{4}$ 在复平面上对应的 4 个点为某个单位圆内接正方形的 4 个顶点, $\left\{  {{a}_{0},{a}_{1},{a}_{2},{a}_{3}}\right\}   \subset  \{ n \mid  1 \leq  n \leq  {2024},n \in  \mathbf{Z}\}$ ,则满足条件的不同集合 $\left\{  {{a}_{0},{a}_{1},{a}_{2},{a}_{3}}\right\}$ 个数为\blkda[2]{$10$}.
    \begin{jieda}
        根据共轭虚根定理可知，$z_1, z_2, z_3, z_4$以共轭虚根的形式存在.

        ${x}^{4} + {a}_{3}{x}^{3} + {a}_{2}{x}^{2} + {a}_{1}x + {a}_{0} = \left( {x - {z}_{1}}\right) \left( {x - {z}_{2}}\right) (x -  \left. {z}_{3}\right) \left( {x - {z}_{4}}\right)  = {x}^{4} - \left( {{z}_{1} + {z}_{2} + {z}_{3} + {z}_{4}}\right) {x}^{3} + \left( {{z}_{1}{z}_{2} + {z}_{1}{z}_{3} + {z}_{1}{z}_{4} + }\right.  \left. {{z}_{2}{z}_{3} + {z}_{2}{z}_{4} + {z}_{3}{z}_{4}}\right) {x}^{2} - \left( {{z}_{1}{z}_{2}{z}_{3} + {z}_{1}{z}_{2}{z}_{4} + {z}_{2}{z}_{3}{z}_{4}}\right) x +  {z}_{1}{z}_{2}{z}_{3}{z}_{4}$ ,对比系数可得 ${a}_{0} = {z}_{1}{z}_{2}{z}_{3}{z}_{4},{a}_{1} =  - \left( {{z}_{1}{z}_{2}{z}_{3} + }\right.  \left. {{z}_{1}{z}_{2}{z}_{4} + {z}_{2}{z}_{3}{z}_{4}}\right) ,{a}_{2} = {z}_{1}{z}_{2} + {z}_{1}{z}_{3} + {z}_{1}{z}_{4} + {z}_{2}{z}_{3} + {z}_{2}{z}_{4} + {z}_{3}{z}_{4}$ , ${a}_{3} =  - \left( {{z}_{1} + {z}_{2} + {z}_{3} + {z}_{4}}\right)$
        
        \begin{dingautolist}{172}
            \item 四个根两两互为共轭复数,故圆心在 $x$ 轴上,设单位圆的圆心为 $\left( {a,0}\right)$ ,不妨设 ${z}_{1} = a + \dfrac{\sqrt{2}}{2} + \dfrac{\sqrt{2}}{2}{i},{z}_{2} = a - \dfrac{\sqrt{2}}{2} + \dfrac{\sqrt{2}}{2}{i},{z}_{3} = a  - \dfrac{\sqrt{2}}{2} - \dfrac{\sqrt{2}}{2}{i},{z}_{4} = a + \dfrac{\sqrt{2}}{2} - \dfrac{\sqrt{2}}{2}{i},{a}_{0} = {z}_{1}{z}_{2}{z}_{3}{z}_{4} = \left( {a + \dfrac{\sqrt{2}}{2} + \dfrac{\sqrt{2}}{2}{i}}\right)  \left( {a - \dfrac{\sqrt{2}}{2} + \dfrac{\sqrt{2}}{2}{i}}\right) \left( {a - \dfrac{\sqrt{2}}{2} - \dfrac{\sqrt{2}}{2}{i}}\right) \left( {a + \dfrac{\sqrt{2}}{2} - \dfrac{\sqrt{2}}{2}{i}}\right)  = \left\lbrack  {\left( a + \dfrac{\sqrt{2}}{2}{i}\right) }^{2}-\dfrac{1}{2}\right\rbrack  \left\lbrack  {{\left( a - \dfrac{\sqrt{2}}{2}{i}\right) }^{2} - \dfrac{1}{2}}\right\rbrack   = \left\lbrack  {{a}^{2} + \sqrt{2}a{i} - 1}\right\rbrack  \left\lbrack  {{a}^{2} - \sqrt{2}a{i} - 1}\right\rbrack  \left\lbrack  {\left( {{a}^{2} - 1}\right)  + \sqrt{2}a{i}}\right\rbrack  \left\lbrack  {\left( {{a}^{2} - 1}\right)  - \sqrt{2}a{i}}\right\rbrack   = {\left( {a}^{2} - 1\right) }^{2} + 2{a}^{2} = {a}^{4} +1$
            
            类似计算可得 ${a}_{1} =  - \left( {{z}_{1}{z}_{2}{z}_{3} + {z}_{1}{z}_{2}{z}_{4} + {z}_{2}{z}_{3}{z}_{4}}\right)  =  - 4{a}^{3},{a}_{2} = {z}_{1}{z}_{2} + {z}_{1}{z}_{3} + {z}_{1}{z}_{4} + {z}_{2}{z}_{3} + {z}_{2}{z}_{4} + {z}_{3}{z}_{4} = 6{a}^{2},  {a}_{3} =  - \left( {{z}_{1} + {z}_{2} + {z}_{3} + {z}_{4}}\right)  =  - {4a}$ ,因为 $\left\{  {{a}_{0},{a}_{1},{a}_{2},{a}_{3}}\right\}   \subset  \{ n \mid  1 \leq  n \leq  {2024},n \in  \mathbf{Z}\}$ ,所以 $a$ 只能为负整数,又集合元素的互异性,从而可得 $a =  - 2, - 3, - 4, - 5, - 6$ ,此时集合 $\left\{  {{a}_{0},{a}_{1},{a}_{2},{a}_{3}}\right\}$ 的个数为 5 个.
            \item 四个根有 2 个为实数, 另外 2 个为共轭复数,故圆心在 $x$ 轴上,设单位圆的圆心为 $\left( {a,0}\right)$ ,不妨设 ${z}_{1} = a - 1,{z}_{2} = a + {i},{z}_{3} = a + 1,{z}_{4} = a - {i}$ ,计算可得 ${a}_{0} = {z}_{1}{z}_{2}{z}_{3}{z}_{4} = {a}^{4} - 1,{a}_{1} =  - \left( {{z}_{1}{z}_{2}{z}_{3} + {z}_{1}{z}_{2}{z}_{4} + }\right.  \left. {{z}_{2}{z}_{3}{z}_{4}}\right)  =  - 4{a}^{3},{a}_{2} = {z}_{1}{z}_{2} + {z}_{1}{z}_{3} + {z}_{1}{z}_{4} + {z}_{2}{z}_{3} + {z}_{2}{z}_{4} +  {z}_{3}{z}_{4} = 6{a}^{2},{a}_{3} =  - \left( {{z}_{1} + {z}_{2} + {z}_{3} + {z}_{4}}\right)  =  - {4a}$ ,因为 $\left\{  {{a}_{0},{a}_{1}}\right.$ , $\left. {{a}_{2},{a}_{3}}\right\}   \subset  \left\{  {n \mid  1 \leq  n \leq  {2024},n \in  \mathbf{Z}}\right\}$ ,所以 $a$ 只能为负整数,又集合元素的互异性,从而可得 $a =  - 2, - 3, - 4, - 5, - 6$ ,此时集合 $\left\{  {{a}_{0},{a}_{1},{a}_{2},{a}_{3}}\right\}$ 的个数为 5 个
        \end{dingautolist}
        综上，满足条件的不同 $\left\{  {{a}_{0},{a}_{1},{a}_{2},{a}_{3}}\right\}$ 的个数为 10 .
    \end{jieda}

    \item (2025$\cdot$杨浦区一模$\cdot$12)已知实数 $a > 0,{i}$ 是虚数单位，设集合 $A = \left\{  {z\left| {z = w + \dfrac{1}{w}}\right. }\right. ,\left| w\right|  > 1,w \in \mathbf{C},z \in  \mathbf{C}\left. \right\}$ ，集合 $B = \{ z \mid  \left| {z - 1 + {i}}\right|  = a,z \in  \mathbf{C}\}$ ，如果 $B \subseteq  A$ ，则 $a$ 的取值范围为\blkda[4]{$\left( {0,1}\right)  \cup  \left( {\sqrt{10}, + \infty }\right)$}.
    \begin{jieda}
        由于 $\left| w\right|  > 1$ ,设 $w = r\left( {\cos \theta  + {i}\sin \theta }\right) ,r > 1$ ,则 $w +  \dfrac{1}{w} = r\left( {\cos \theta  + {i}\sin \theta }\right)  + \dfrac{1}{r\left( {\cos \theta  + {i}\sin \theta }\right) } = r\left( {\cos \theta  + {i}\sin \theta }\right)  +  \dfrac{\cos \theta  - i\sin \theta }{r} = \left( {r + \dfrac{1}{r}}\right) \cos \theta  + {i}\left( {r - \dfrac{1}{r}}\right) \sin \theta ,$ 设 $z$ 对应点 $(x,  y)$ ,则 $x = \left( {r + \dfrac{1}{r}}\right) \cos \theta ,y = \left( {r - \dfrac{1}{r}}\right) \sin \theta$ ,所以 $\dfrac{{x}^{2}}{{\left( r + \dfrac{1}{r}\right) }^{2}}  + \dfrac{{y}^{2}}{{\left( r - \dfrac{1}{r}\right) }^{2}} = 1$ ,其中 $r + \dfrac{1}{r} > 2,r - \dfrac{1}{r} > 0$ ,当 $r > 1$ 时,该方程的几何意义为表示所有椭圆的并集,即平面上除去线段 ${A}_{1}{A}_{2}$ 的点的集合,其中 ${A}_{1}\left( {-2,0}\right) \text{ 、 }{A}_{2}\left( {2,0}\right)$ ,集合 $B = \{ z \mid  \left| {z - 1 + {i}}\right|  = a,z \in  \mathbf{C}\}$ 表示复平面上的圆,圆心为 $\left( {1, - 1}\right)$ ,半径为 $a$ ,如果 $B \subseteq  A$ ,则该圆与线段 ${A}_{1}{A}_{2}$ 无公共点,结合图形可知 $a$ 的取值范围为 $\left( {0,1}\right)  \cup  \left( {\sqrt{10}, + \infty }\right)$.

        \begin{center}
        \includegraphics[width=0.3\textwidth]{bodyxb/src/105_1289_581_239_191_0.jpg}
        \end{center}
    \end{jieda}

\end{enumerate}
\subsubsection*{二、选择题}
\begin{enumerate}
    \item (2025$\cdot$静安区二模$\cdot$14)若复数 $z = \dfrac{a{i}}{b + {i}}$ ( $a$ 、 $b \in  \mathbf{R}$ ， ${i}$ 是虚数单位)在复平面上对应的点位于第二象限, 则 \dotfill{(\kh{D})}
    \xx{$a > 0$ 且 $b > 0$}{$a > 0$ 且 $b < 0$}{$a < 0$ 且 $b > 0$}{$a < 0$ 且 $b < 0$}
    \begin{jieda}
        $z = \dfrac{a{i}}{b + {i}} = \dfrac{a{i}\left( {b - {i}}\right) }{\left( {b + {i}}\right) \left( {b - {i}}\right) } = \dfrac{a}{{b}^{2} + 1} + \dfrac{ab}{{b}^{2} + 1}{i}$ ,因为复数 $z = \dfrac{a{i}}{b + {i}}\left( {a\text{ 、 }b \in  \mathbf{R},{i}}\right.$ 是虚数单位) 在复平面上对应的点位于第二象限，所以 $\left\{  \begin{array}{l} \dfrac{a}{{b}^{2} + 1} < 0 \\  \dfrac{ab}{{b}^{2} + 1} > 0 \end{array}\right.$ ，解得 $\left\{  \begin{array}{l} a < 0 \\  b < 0 \end{array}\right.$
    \end{jieda}

    \item (2024$\cdot$宝山区一模$\cdot$15)已知 $z$ 是复数， $\overline{z}$ 是其共轭复数，则下列命题中正确的是 \dotfill{(\kh{B})}
    \xx{${z}^{2} = {\left| z\right| }^{2}$}{若 $\left| z\right|  = 1$ ,则 $\left| {z - 1 - {i}}\right|$ 的最大值为 $\sqrt{2} + 1$}{若 $z = {\left( 1 - 2{i}\right) }^{2}$ ,则复平面内 $\overline{z}$ 对应的点位于第一象限}{若 $1 - 3{i}$ 是关于 $x$ 的方程 ${x}^{2} + {px} + q = 0\left( {p\text{ 、 }q \in  \mathbf{R}}\right)$ 的一个根,则 $q =  - 8$}

    \item (2024$\cdot$闵行区一模$\cdot$15)已知复数 ${z}_{1}$ 、 ${z}_{2}$ 在复平面内对应的点分别为 $P$ 、 $Q$ ， $\left| {OP}\right|  = 5$ ( $O$ 为坐标原点),且 ${z}_{1}^{2} - {z}_{1}{z}_{2} \cdot  \sin \theta  + {z}_{2}^{2} = 0$ ,则对任意 $\theta  \in  \mathbf{R}$ ,下列选项中为定值的是 \dotfill{(\kh{A})}
    \xx{$\left| {OQ}\right|$}{ $\left| {PQ}\right|$}{$\bigtriangleup {OPQ}$ 的周长}{$\bigtriangleup {OPQ}$ 的面积}
    \begin{jieda}
        $sin\theta = \dfrac{z_1}{z_2} + \dfrac{z_2}{z_1}$，因为$sin\theta \in \mathbf{R}$，所以$\dfrac{z_1}{z_2} = \overline{\left(\dfrac{z_2}{z_1}\right)}$，故$\left|\dfrac{z_1}{z_2}\right| = 1$，故$|z_2| = |z_1| = 5$.
    \end{jieda}
\end{enumerate}
\section*{第 10 章 空间向量与立体几何}

\begin{center}

\begin{tabular}{|p{4cm}|p{12cm}|}
\hline
考点 & 课标要求与命题趋势 \\
\cline{1-2}
1. 空间几何体的结构特征、表面积和体积 & \ding{172}  利用实物、计算机软件等观察空间图形，认识柱、锥、台、球及简单组合体的结构特征， 能运用这些特征描述现实生活中简单物体的结构 \ding{173}  知道球、棱柱、棱锥、棱台的表面积和体积的计算公式, 能用公式解决简单的实际问题 \ding{174}  能用斜二测法画出简单空间图形(长方体、棱柱及其简单组合)的直观图 \ding{175} 借助长方体，在直观认识空间点、直线、平面的位置关系的基础上，抽象出空间点、直线、平面的位置关系的定义，了解基本事实(基本事实 1-4 也称公理)和定理 \\
\cline{1-2}
2. 空间向量及其应用 & \ding{172}  在平面直角坐标系的基础上，了解空间直角坐标系，感受建立空间直角坐标系的必要性，会用空间直角坐标系刻画点的位置；借助特殊长方体(所有被分别与坐标轴平行) 顶点的坐标; 探索并得出空间两点间的距离公式 \ding{173}  经历由平面向量推广到空间向量的过程，了解空间向量的概念；经历由平面向量的运算及其法则推广到空间向量的过程 \ding{174}  了解空间向量基本定理及其意义，掌握空间向量的正交分解及其坐标表示. 掌握空间向量的线性运算及其坐标表示; 掌握空间向量的数量积及其坐标表示; 了解空间向量投影的概念以及投影向量的意义 \ding{175}  能用向量语言表述直线和平面，理解直线的方向向量与平面的法向量；能用向量语言表述直线与直线、直线与平面、平面与平面的夹角以及垂直与平行关系; 能用向量方法证明必修内容中有关直线、平面位置关系的判定定理; 能用向量方法解决点到直线、点到平面、相互平行的直线、相互平行的平面的距离问题和简单夹角问题，并能描述解决这一类问题的程序，体会向量方法在研究几何问题中的作用 \\
\cline{1-2}
3. 直线、平面平行的判定与性质 4. 直线、平面垂直的判定与性质& \ding{172}  从定义和基本事实出发，借助长方体，通过直观感知，了解空间中直线与直线、直线与平面、平面与平面的平行和垂直的关系，归纳出相应的判定定理，并加以证明 \ding{173}  能用已获得的结论，证明空间基本图形位置关系的简单命题 \ding{174}  理解、掌握空间中点线面的位置关系及相关的图形和符号语言，熟练掌握平行关系的判定定理和性质定理及其应用，熟练掌握垂直关系的判定定理和性质定理及其应用， 该内容是上海高考的常考内容，需强化巩固复习 \\
\cline{1-2}
5. 空间角与距离 & \ding{172}  熟练掌握几何法和向量法求解空间角与空间距离，该内容是上海高考的常考内容，要熟练掌握方程思想求值，需强化巩固复习 \ding{173}  立体几何研究现实世界中物体的形状、大小与位置关系. 以长方体为载体, 认识和理解空间点、直线、平面的位置关系; 用数学语言表述有关平行、垂直的性质与判定, 并对某些结论进行论证; 了解一些简单几何体的表面积与体积的计算方法; 运用直观感知、操作确认、推理论证、度量计算等认识和探索空间图形的性质, 建立空间观念 \\
\cline{1-2}
6. 空间几何体的截面、交线及动态问题 & \ding{172} 了解柱、锥、台体及简单组合体的结构特征及其相关性质，会运用柱体、锥体、台体等组合体的表面积和体积的计算公式求解相关问题, 该内容是上海高考的常考内容, 一般给定柱、锥、台体及简单组合体, 求对应的表面积与体积, 需强化复习 \ding{173} 收集、阅读几何发展的历史资料，撰写小论文，论述几何发展的过程、重要结果、主要人物、关键事件及其对人类文明的贡献 \\
\cline{1-2}
\hline
\end{tabular}
\end{center}
\clearpage
\subsection*{考点1 空间几何体的结构特征、表面积和体积}

\subsubsection*{一、填空题}
\begin{enumerate}
    \item (2024$\cdot$虹口区二模$\cdot$2)已知一个球的表面积为 ${36\pi }$ ，则该球的体积为\blkda[2]{${36\pi }$}.

    \item (2025$\cdot$长宁区一模$\cdot$2)已知圆锥的底面半径为 1，母线长为 2，则该圆锥的体积是\blkda[2]{$\dfrac{\sqrt{3}}{3}\pi$}.(结果保留 $\pi$ )

    \item (2024$\cdot$嘉定区二模$\cdot$3)已知圆锥的母线长为 2，高为 1，则其体积为\blkda[2]{$\pi$}.

    \item (2024$\cdot$静安区一模$\cdot$3)若一个圆柱的底面半径和母线长都是 1，则这个圆柱的体积是\blkda[2]{$\pi$}.

    \item (2024$\cdot$虹口区一模$\cdot$4)已知一个圆锥的底面半径为 3，其侧面积为 15π，则该圆锥的体积为\blkda[2]{${12\pi}$}.

    \item (2025$\cdot$虹口区二模$\cdot$4)已知圆柱的底面半径为 1，母线长为 2，则该圆锥的侧面积为\blkda[2]{${4\pi}$}.

    \item (2024$\cdot$奉贤区二模$\cdot$4)已知圆锥的底面半径为 $\sqrt{2}$ ，母线长为 2，则该圆锥的侧面积为\blkda[2]{$2\sqrt{2}\pi$}.

    \item (2025$\cdot$闵行区二模$\cdot$4)已知圆柱的底面半径为 $\sqrt{3}$ ，高为 3，则圆柱的体积为\blkda[2]{${9\pi }$}.

    \item (2024$\cdot$黄浦区二模$\cdot$4)若一个圆柱的底面半径为 2，母线长为 3，则此圆柱的侧面积为\blkda[2]{${12\pi }$}.

    \item (2025$\cdot$黄浦区一模$\cdot$4)若圆柱的底面半径与高均为 1，则其侧面积为\blkda[2]{${2\pi }$}.

    \item (2025$\cdot$闵行区一模$\cdot$5)已知圆锥的高为 8，底面半径为 6，则该圆锥的侧面积为\blkda[2]{${60\pi }$}.

    \item (2024$\cdot$闵行区一模$\cdot$5)已知圆锥的底面周长为 ${4\pi }$ ，母线长为 3，则该圆锥的侧面积为\blkda[2]{${6\pi }$}.

    \item (2024$\cdot$松江区二模$\cdot$6)若一个圆锥的侧面展开图是面积为 ${2\pi }$ 的半圆面，则此圆锥的体积为\blkda[2]{$\dfrac{\sqrt{3}}{3}\pi$}. (结果中保留 $\pi$ )

    \item (2024$\cdot$崇明区一模$\cdot$6)已知圆锥的母线与底面所成角为 ${45}^{ \circ  }$ ，高为 1，则该圆锥的母线长为\blkda[2]{$\sqrt{2}$}.

    \item (2024$\cdot$黄浦区一模$\cdot$6)若棱长为 2 的正方体的所有顶点都在同一球面上，则该球的体积为\blkda[2]{$4\sqrt{3}\pi$}.

    \item (2024$\cdot$青浦区一模$\cdot$6)已知球的表面积是 ${4\pi }$ ，则该球的体积是\blkda[2]{$\dfrac{4}{3}\pi$}.

    \item (2025$\cdot$松江区一模$\cdot$6)已知一个圆锥的底面半径为 3，其侧面积为 ${15\pi }$ ，则该圆锥的高为\blkda[2]{$4$}.

    \item (2025$\cdot$嘉定区一模$\cdot$6)某圆锥的母线长为 2，底面半径为 1，则该圆锥的侧面积为\blkda[2]{${2\pi }$}.

    \item (2025$\cdot$虹口区一模$\cdot$6)若某圆锥的底面半径为 1，高为 1，则该圆锥的侧面积为\blkda[2]{$\sqrt{2}\pi$}.(结果保留 $\pi )$

    \item (2024$\cdot$金山区二模$\cdot$6)若长方体 ${ABCD} - {A}_{1}{B}_{1}{C}_{1}{D}_{1}$ 的体积是 ${120},E$ 为棱 $C{C}_{1}$ 的中点，则三棱锥 $E - {BCD}$ 的体积是\blkda[2]{$10$}.

    \item (2025$\cdot$杨浦区一模$\cdot$7)已知一个正四棱锥的每一条棱长都为 2，则该四棱锥的体积为\blkda[2]{$\dfrac{4\sqrt{2}}{3}$}.
    \begin{jieda}
        $V = \dfrac{1}{3} \times 4 \times \sqrt{2} = \dfrac{4\sqrt{2}}{3}$
    \end{jieda}

    \item (2025$\cdot$闵行区二模$\cdot$7)已知圆锥底面半径为 1，高为 $\sqrt{3}$ ，则过圆锥母线的截面面积的最大值为\blkda[2]{$\sqrt{3}$}.
    \begin{jieda}
        母线长为$2$，母线最大夹角为$\dfrac{\pi}{3}$，故轴截面为面积最大截面，$S = \dfrac{1}{2} \times 2 \times 2 \times \dfrac{\sqrt{3}}{2} = \sqrt{3}$.
    \end{jieda}

    \item (2024$\cdot$浦东新区一模$\cdot$7)已知圆锥的母线与底面所成的角为 $\dfrac{\pi }{3}$ ，体积为 ${3\pi }$ ，则圆锥的底面半径为\blkda[2]{$\sqrt{3}$}.
    \begin{jieda}
        $V = \dfrac{1}{3}\pi {r}^{2}h = \dfrac{1}{3}\pi {r}^{2} \times  \sqrt{3}r = {3\pi } \Longrightarrow  r = \sqrt{3}$.
    \end{jieda}

    \item (2024$\cdot$金山区一模$\cdot$7)已知某圆台上底面和下底面的半径分别为1 和 2， 母线长为 3，则该圆台的高为\blkda[2]{$2\sqrt{2}$}.
    \begin{jieda}
        作出圆台的轴截面, 如图示为等腰梯形, 梯形的高即为圆台的高, 即高为 $\sqrt{{3}^{2} - {\left( 2 - 1\right) }^{2}} = 2\sqrt{2}$.
        \begin{center}
            \includegraphics[width=0.15\textwidth]{bodyxb/src/107_664_274_155_176_0.jpg}
        \end{center}
    \end{jieda}

    \item (2025$\cdot$上海秋季卷$\cdot$7)如图所示，在正四棱柱 ${ABCD} - {A}_{1}{B}_{1}{C}_{1}{D}_{1}$ 中， ${BD} = 4\sqrt{2},{DB}_{1} = 9$ ，则该正四棱柱的体积为\blkda[2]{$112$}.
    \begin{center}
    \includegraphics[width=0.2\textwidth]{bodyxb/src/144_1221_389_249_289_0.jpg}
    \end{center}
    \begin{jieda}
        因为 ${BD} = 4\sqrt{2}$ 且四边形 ${ABCD}$ 为正方形,故 ${BA} = 4$，而 $D{B}_{1} = 9$ ,故 $B{B}_{1}^{2} + B{D}^{2} = {81}$ ,故 $B{B}_{1} = 7$ ,故所求体积为 $7 \times  {16} = {112}$.
    \end{jieda}

    \item (2025$\cdot$宝山区二模$\cdot$8)已知圆柱的底面积为$9\pi$，侧面积为$18\pi$，则该圆柱的体积为\blkda[2]{${27\pi }$}.
    \begin{jieda}
        设圆柱的底面圆的半径为 $r$ ,高为 $h$，故有 $\left\{ \begin{array}{l} \pi {r}^{2} = {9\pi } \\  {2\pi rh} = {18\pi } \end{array}\right.$，解得 $\left\{ \begin{array}{l} r = 3 \\  h = 3 \end{array}\right.$，所以圆柱的体积$V = \pi {r}^{2}h = {27\pi }.$
    \end{jieda}

    \item (2025$\cdot$普陀区二模$\cdot$8)若一个圆锥的高为 $\sqrt{2}$ ，侧面积为 $2\sqrt{2}\pi$ ，则该圆锥侧面展开图中扇形的中心角的大小为\blkda[2]{$\sqrt{2}\pi$}.
    \begin{jieda}
        设底面半径为 $r$ ,母线长为 $l$ ,由 ${\pi rl} = 2\sqrt{2}\pi$ ,得 ${rl} =  2\sqrt{2}$ ,又 $h = \sqrt{2}$ ,由勾股定理 $l = \sqrt{{r}^{2} + 2}$ ,所以 $r\sqrt{{r}^{2} + 2} =  2\sqrt{2}$ ,解得 $r = \sqrt{2},l = 2$ ,底面圆周长 $L = {2\pi r} = 2\sqrt{2}\pi$ ,扇形中心角 $\theta  = \dfrac{L}{l} = \dfrac{2\sqrt{2}\pi }{2} = \sqrt{2}\pi$.
    \end{jieda}

    \item (2024$\cdot$普陀区一模$\cdot$8)设圆锥的底面中心为 $O$ ， ${PB}$ 、 ${PC}$ 是它的两条母线，且 $\left| {BC}\right|  = 2$ ，若棱锥 $O - {PBC}$ 是正三棱锥，则该圆锥的侧面积为\blkda[2]{$2\sqrt{2}\pi$}.
    \begin{jieda}
        由棱锥 $O - {PBC}$ 为正三棱锥，得 ${PB} = {PC} = {BC} = 2,{OB} = {OC} =  {OP}$ ,而 ${PO} \bot  {OB},{PO} \bot  {OC}$ ,由勾股定理得 ${OB} = {OC} = {OP} = \sqrt{2}$ ,即圆锥的底面圆半径 $r = \sqrt{2}$ ,母线长 $l = 2$ ,则该圆锥的侧面积为 ${\pi rl} = 2\sqrt{2}\pi$.
        \begin{center}
        \includegraphics[width=0.2\textwidth]{bodyxb/src/107_605_1091_213_191_0.jpg}
        \end{center}
    \end{jieda}

    \item (2025$\cdot$普陀区一模$\cdot$8)若圆锥 $P - O$ 的体积为 $\dfrac{2\sqrt{2}\pi }{3}$ ，它的母线与底面所成的角的余弦值为 $\dfrac{1}{3}$ ， 则圆锥 $P - O$ 的表面积为\blkda[2]{${4\pi }$}.
    \begin{jieda}
        设圆锥底面半径 ${AO} = {OB} = r$ ,则母线长 $l = {PA} =  {3r}$ ,高 ${PO} = 2\sqrt{2}r$ ,则 $V = \dfrac{1}{3}\pi {r}^{2} \cdot  2\sqrt{2}r = \dfrac{2\sqrt{2}}{3}\pi$ ,解得 $r = 1$ , $l = {PA} = 3,{PO} = 2\sqrt{2}$ ,故该圆锥的表面积为 $S = {\pi rl} + \pi {r}^{2} =  {4\pi }$.
    \end{jieda}

    \item (2025$\cdot$青浦区一模$\cdot$8)已知圆柱 $M$ 的底面半径为3，高为 $\sqrt{3}$ ，圆锥 $N$ 的底面直径和母线长相等. 若圆柱 $M$ 和圆锥 $N$ 的体积相同，则圆锥 $N$ 的底面半径为\blkda[2]{3}.
    \begin{jieda}
        圆柱 $M$ 的体积为 $\pi  \cdot  {3}^{2} \times  \sqrt{3} = 9\sqrt{3}\pi$ ,设圆锥 $N$ 的底面半径为 $r$ ,则母线长为 ${2r}$ ,故圆锥的高为 $\sqrt{{\left( 2r\right) }^{2} - {r}^{2}} =  \sqrt{3}r$ ,则 $\dfrac{1}{3}\pi  \cdot  {r}^{2} \cdot  \sqrt{3}r = \dfrac{\sqrt{3}}{3}\pi {r}^{3}$ ,故 $\dfrac{\sqrt{3}}{3}\pi {r}^{3} = 9\sqrt{3}\pi$ ,解得 $r = 3$.
    \end{jieda}

    \item (2025$\cdot$奉贤区二模$\cdot$8)抛物线 ${x}^{2} = {4y}$ 的准线与圆 ${x}^{2} + {y}^{2} = {r}^{2}$ 相切，将圆绕直径所在直线旋转一周形成一个几何体，则该几何体的表面积为\blkda[2]{${4\pi }$}.
    \begin{jieda}
        因为抛物线方程为 ${x}^{2} = {4y}$ ,可知准线方程为 $y =$ -1,又由圆 ${x}^{2} + {y}^{2} = {r}^{2}$ 与准线相切,可知: $r = 1$ ,将圆绕直径旋转一周所成的球的表面积为: $S = {4\pi }$.
    \end{jieda}

    \item (2025$\cdot$杨浦区一模$\cdot$9)将一个半径为 1 的球形石材加工成一个圆柱形摆件，则该圆柱形摆件侧面积的最大值为\blkda[2]{${2\pi }$}.
    \begin{jieda}
        设圆柱形工件的高为 $h$ ,底面半径为 $r$，$\left\{\begin{array}{l}
            h = 2cos\theta \\
            r = sin\theta
        \end{array}\right., \theta \in \left(0, \dfrac{\pi}{2}\right)$，$S = {2\pi rh} = 4\pi sin\theta cos\theta = 2\pi sin2\theta \leq 2\pi$，当且仅当 $\theta = \dfrac{\pi}{4}$ 时等号成立.
    \end{jieda}

    \item (2024$\cdot$长宁区二模$\cdot$9)用铁皮制作一个有底无盖的圆柱形容器，若该容器的容积为 $\pi {{m}}^{3}$ ，则至少需要\blkda[2]{${3\pi }$} ${{m}}^{2}$ 铁皮.
    \begin{jieda}
        设圆柱形容器的底面半径为 $r$ ,高为 $h$ ,所以圆柱形容器的体积为 $V = \pi {r}^{2} \cdot  h = \pi$ ,所以 ${r}^{2} \cdot  h = 1$ ,所以圆柱形容器的表面积为: $S = \pi {r}^{2} + {2\pi rh} = \pi \left( {{r}^{2} + \dfrac{2}{r}}\right)  \geq  {3\pi }$.
    \end{jieda}

    \item (2024$\cdot$宝山区一模$\cdot$9)如图所示，在圆锥 $S - O$ 中， ${AC}$ 为底面圆 $O$ 的直径， ${SO} = {OC} = 1$ ，点 $B$ 在底面圆周上，且 ${AB} = {BC}$. 若 $E$ 为线段 ${AB}$ 上的动点，则 $\bigtriangleup  {SEC}$ 的周长最小值为\blkda[2]{$\sqrt{3} + \sqrt{2} + 1$}.
    \begin{center}
    \includegraphics[width=0.2\textwidth]{bodyxb/src/144_1141_1392_329_259_0.jpg}
    \end{center}
    \begin{jieda}
        连接 ${OB}$ ,依题意 ${SO} \bot$ 平面 ${ABC}$ ,而 ${OA}\text{ 、 }{OB}\text{ 、 }{OC} \subset$ 平面 ${ABC}$ ,所以 ${SO} \bot  {OA},{SO} \bot  {OB},{SO} \bot  {OC},{AB} = {BC},O$ 是 ${AC}$ 的中点,则 ${OB} \bot  {AC}$ ,由于 ${SO} = {OC} = 1$ ,所以 ${SA} =  {SC} = {SB} = {AB} = \sqrt{2}$ ,则 $\bigtriangleup {SAB}$ 是等边三角形， $\bigtriangleup {ABC}$ 是等腰直角三角形，将 $\bigtriangleup {SAB}$ 和 $\bigtriangleup {ABC}$ 展开在同一个平面，如图所示，连接 ${SC}$ ，交 ${AB}$ 于 $E$ ，在 $\bigtriangleup  {SAC}$ 中, 由余弦定理得${SC} = \sqrt{2 + 4 - 2 \times  \sqrt{2} \times  2 \times  \cos \left( {{60}^{ \circ  } + {45}^{ \circ  }}\right)} = \sqrt{6 - 4\sqrt{2}}\left( {\cos {60}^{ \circ  }\cos {45}^{ \circ  } - \sin {60}^{ \circ  }\sin {45}^{ \circ  }}\right) = \sqrt{4 + 2\sqrt{3}} = \sqrt{{\left( \sqrt{3} + 1\right) }^{2}} = \sqrt{3} + 1$

        \begin{center}
        \includegraphics[width=0.2\textwidth]{bodyxb/src/107_1350_710_180_134_0.jpg}
        \end{center}
        

        所以 $\bigtriangleup  {SEC}$ 的周长最小值为 $\sqrt{3} + \sqrt{2} + 1$ .
    \end{jieda}

    \item (2024$\cdot$静安区二模$\cdot$9)正四棱锥 $P - {ABCD}$ 底面边长为 2，高为 3， 则点 $A$ 到不经过点 $A$ 的侧面的距离为\blkda[2]{$\dfrac{3\sqrt{10}}{5}$}.
    \begin{jieda}
        设底面正方形 ${ABCD}$ 的中心为 $O,{BC}$ 的中点为 $E$ ,如图所示: 易知 ${EO} = 1$ , 高 ${PO} = 3$ ,所以其斜高 ${PE} = \sqrt{{3}^{2} + {1}^{2}} =  \sqrt{10}$ ,由对称性可知点 $A$ 到侧面 ${PCD}$ 与侧面 ${PBC}$ 的距离相等,易知侧面 ${PBC}$ 和侧面 ${PCD}$ 与的面积 ${S}_{\bigtriangleup {PDC}} = {S}_{\bigtriangleup {PBC}} = \dfrac{1}{2} \times  2 \times  \sqrt{10} = \sqrt{10}$ , 正四棱锥 $P - {ABCD}$ 的体积为 $V = \dfrac{1}{3}{S}_{ABCD} \cdot  {PO} = \dfrac{1}{3} \times  2 \times  2 \times  3 = 4$ ,设点 $A$ 到侧面 ${PBC}$ 的距离为 $d$ ,由等体积法可得 $V = {V}_{A - {PBC}} + {V}_{A - {PDC}} = \dfrac{1}{3}{S}_{\bigtriangleup {PDC}}d + \dfrac{1}{3}{S}_{\bigtriangleup {PBC}}d = \dfrac{1}{3}d \times  2\sqrt{10} = 4$ ,解得 $d = \dfrac{3\sqrt{10}}{5}$.

        \begin{center}
        \includegraphics[width=0.2\textwidth]{bodyxb/src/107_1341_1215_189_180_0.jpg}
        \end{center}
    \end{jieda}

    \item (2024$\cdot$奉贤区二模$\cdot$10)学生到工厂劳动实践，利用3D打印技术制作模型,如图所示. 该模型为长方体 ${ABCD} - {A}_{1}{B}_{1}{C}_{1}{D}_{1}$ 中挖去一个四棱锥 $O - {EFGH}$ ,其中 $O$ 为长方体的中心, $E\text{ 、 }F\text{ 、 }G\text{ 、 }H$ 分别为所在棱的中点, ${AB} = {BC} = 4{{cm}},{A{A}_{1}} = 2{{cm}},$ 3D 打印所用原料密度为 ${0.9}{g}/{{{cm}}}^{3}$. 不考虑打印损耗,制作该模型所需原料的质量为\blkda[2]{$\dfrac{132}{5}$}$g$.

    \begin{center}
    \includegraphics[width=0.2\textwidth]{bodyxb/src/144_1151_1759_319_287_0.jpg}
    \end{center}
    \begin{jieda}
        易知四棱锥 $O - {EFGH}$ 的底面积 ${S}_{EFGH} = {S}_{{BC}{C}_{1}{B}_{1}} -  4{S}_{\bigtriangleup {ECF}} = 4 \times  2 - 4 \times  \dfrac{1}{2} \times  2 \times  1 = 4$ ,高为 $h = 2{{cm}}$ ,所以四棱锥 $O - {EFGH}$ 的体积为 ${V}_{O - {EFGH}} = \dfrac{1}{3} \times  4 \times  2 = \dfrac{8}{3}{{{cm}}}^{3}$ ,长方体 ${ABCD} - {A}_{1}{B}_{1}{C}_{1}{D}_{1}$ 体积为 $V = 4 \times  4 \times  2 = {32}{{{cm}}}^{3}$ ,因此该模型的体积为 ${32} - \dfrac{8}{3} = \dfrac{88}{3}{{{cm}}}^{3}$ ,所以该模型所需原料的质量为 $\dfrac{88}{3} \times  {0.9} = \dfrac{132}{5}{g}$.
    \end{jieda}


    \item (2025$\cdot$宝山区一模$\cdot$10)将棱长为 2 的正四面体绕着它的某一条棱旋转一周所得的几何体的体积为\blkda[2]{${2\pi }$}.
    \begin{jieda}
        如图为棱长为 2 的正四面体, 作 ${BD} \bot  {AC},{BD} \cap  {AC}$ 于点 $D$ ,在 $\bigtriangleup {ABC}$ 中, ${AB} = 2,\angle {BAC} = {60}^{ \circ  }$ , ${BD} = {AB} \cdot  \sin {60}^{ \circ  } = \sqrt{3}$ ，根据题意，正四面体旋转后得到两个同底的圆锥，底面半径等于 ${BD} = \sqrt{3}$ ,高等于 $\dfrac{\left| AC\right| }{2} = 1$ ,所以旋转后几何体的体积为: $V = 2 \times  \dfrac{1}{3} \cdot  \pi  \cdot  {\left( \sqrt{3}\right) }^{2} = {2\pi }$.

        \begin{center}
            \includegraphics[width=0.15\textwidth]{bodyxb/src/108_657_196_211_190_0.jpg}
        \end{center}
    \end{jieda}

    \item (2024$\cdot$浦东新区二模$\cdot$10)如图所示，有一底面半径为1、高为3的圆柱，光源点 $A$ 沿着上底面圆周作匀速运动，射出的光线始终经过圆柱轴截面的中心，当光源点 $A$ 沿着上底面圆周运动半周时，其射出的光线在圆柱内部“扫过”的面积为\blkda[2]{$\dfrac{\sqrt{13}}{2}\pi$}.
    \begin{center}
    \includegraphics[width=0.15\textwidth]{bodyxb/src/145_364_1059_203_279_0.jpg}
    \end{center}
    \begin{jieda}
        如图所示光线在圆柱内部“扫过”的轨迹为两个半圆锥的侧面, $r = 1,h = \dfrac{3}{2},l =  \sqrt{1 + \dfrac{9}{4}} = \dfrac{\sqrt{13}}{2} \Longrightarrow  S = 2 \times  \dfrac{1}{2}{\pi rl} = \dfrac{\sqrt{13}}{2}\pi.$
        \begin{center}
        \includegraphics[width=0.15\textwidth]{bodyxb/src/108_725_579_143_198_0.jpg}
        \end{center}
    \end{jieda}
    \item (2025$\cdot$杨浦区二模$\cdot$10)如图，点 $D$ 、 $E$ 分别是直角三角形 ${ABC}$ 的边 ${AB}$ 、 ${BC}$ 上的点，斜边 ${AC}$ 与扇形的弧 $\overset{\frown}{DE}$ 相切,已知 ${AC} = 4,{BC} = 2$ ,则阴影部分绕直线 ${AB}$ 旋转一周所形成的几何体的体积为\blkda[2]{$\dfrac{2\sqrt{3}}{3}\pi$}.
    \begin{center}
    \includegraphics[width=0.15\textwidth]{bodyxb/src/145_702_989_214_349_0.jpg}
    \end{center}
    \begin{jieda}
        在$Rt\bigtriangleup {ABC}$ 中, $\angle B = {90}^{ \circ  }$ ,则 ${AB} = \sqrt{A{C}^{2} - B{C}^{2}} =  \sqrt{{4}^{2} - {2}^{2}} = 2\sqrt{3}$ ,由斜边 ${AC}$ 与扇形的弧 $\overset{\frown}{DE}$ 相切,扇形半径 $r = \dfrac{{AB} \cdot  {BC}}{AC} = \sqrt{3}$ ,阴影部分绕直线 ${AB}$ 旋转一周所形成的几何体是 ${{Rt}} \bigtriangleup  {ABC}$ 绕直线 ${AB}$ 旋转一周所得圆锥，挖去扇形弧 $\overset{\frown}{DE}$ 绕直线 ${AB}$ 旋转一周所得半球,所以所求体积为 $\dfrac{1}{3}\pi  \times  {2}^{2} \times  2\sqrt{3} - \dfrac{1}{2} \times  \dfrac{4\pi }{3} \times  {\left( \sqrt{3}\right) }^{3} = \dfrac{2\sqrt{3}}{3}\pi$.
    \end{jieda}

    \item (2024$\cdot$崇明区一模$\cdot$10)用易拉罐包装的饮料是超市和自动售卖机里的常见商品. 图为某品牌的易拉罐包装的饮料. 在满足容积要求的情况下，饮料生产商总希望包装材料的成本最低，也就是易拉罐本身的质量最小. 某数学兴趣小组对此想法通过数学建模进行验证. 为了建立数学模型, 他们提出以下 3 个假设: \ding{172}  易拉罐容积相同； \ding{173}  易拉罐是一个上下封闭的空心圆柱体； \ding{174}  易拉罐的罐顶、罐体和罐底的厚度和材质都相同. 你认为以此 3 个假设所建立的数学模型与实际情况相符吗？若相符，请在以下横线上填写“相符”； 若不相符，请选择其中的一个假设给出你的修改意见，并将修改意见填入横线:\blkda[2]{}。

    \begin{center}
    \includegraphics[width=0.1\textwidth]{bodyxb/src/145_1057_1042_157_275_0.jpg}
    \end{center}
    \begin{jieda}
        由题意知, 某品牌的易拉罐包装的饮料, 在满足容积要求的情况下, 饮料生产商总希望包装材料的成本最低, 也就是易拉罐本身的质量最小，所以假设 2 不合理，应为 “易拉罐的顶部类似于圆台”；假设 3 不合理，应为 “易拉罐的罐顶和罐底材质比罐体的材质厚”. 故答案为: 假设 2 中, 易拉罐的顶部类似于圆台; 假设 3 中, 易拉罐的罐顶和罐底材质比罐体的材质厚.
    \end{jieda}

    \item (2025$\cdot$徐汇区一模$\cdot$11)徐汇滨江作为 2024 年上海国际鲜花展的三个主会场之一, 吸引了广大市民前往观展并拍照留念. 图中的花盆是种植鲜花的常见容器，它可视作两个圆台的组合体，上面圆台的上、下底面直径分别为 ${30}{{cm}}$ 和 26 cm，下面圆台的上、下底面直径分别为 ${24}{{cm}}$ 和 ${18}{{cm}}$ ，且两个圆台侧面展开图的圆弧所对的圆心角相等. 若上面圆台的高为 $8{{cm}}$ ，则该花盆上、下两部分母线长的总和为\blkda[2]{$5\sqrt{17}$}$cm$.

    \begin{center}
    \includegraphics[width=0.15\textwidth]{bodyxb/src/145_1241_1439_178_214_0.jpg}
    \end{center}
    \begin{jieda}
        设上面圆台的母线长为 ${l}_{1}$ ,上面半径为 ${r}_{1} = {15}{{cm}}$ ,下半圆半径为 ${r}_{2} = {13}{{cm}}$ ,高为 $h = 8{{cm}}$ ,根据圆台的母线长公式 $l = \sqrt{{h}^{2} + {\left( {r}_{1} - {r}_{2}\right) }^{2}}$ ,带入数值计算得到 ${l}_{1} = \sqrt{{8}^{2} + {\left( {15} - {13}\right) }^{2}} =  \sqrt{68} = 2\sqrt{17}{{cm}}$ ; 设下面圆台的母线长为 ${l}_{2}$ ,上面半径为 ${r}_{3} =  {12}{{cm}}$ ,下半圆半径为 ${r}_{4} = 9{{cm}}$ ,由于两个圆台侧面展开图的圆弧所对的圆心角相等,可以得到 $\dfrac{{r}_{1} - {r}_{2}}{{l}_{1}} = \dfrac{{r}_{3} - {r}_{4}}{{l}_{2}}$ ,带入数值计算得到 ${l}_{2} = \dfrac{\left( {{r}_{3} - {r}_{4}}\right) {l}_{1}}{{r}_{1} - {r}_{2}} = \dfrac{\left( {{12} - 9}\right)  \times  2\sqrt{17}}{{15} - {13}} = 3\sqrt{17}{{cm}}$ ; 所以该花盆上、下两部分母线长的总和为 $2\sqrt{17} + 3\sqrt{17} =  5\sqrt{17}{{cm}}$. 
    \end{jieda}

    \item (2025$\cdot$奉贤区一模$\cdot$11)上海市奉贤区奉城镇的古建筑万佛阁(图 1)的屋檐下常系挂风铃 (图 2), 风吹铃动, 悦耳清脆, 亦称惊鸟铃. 惊鸟铃一般由铜铸造而成, 包括铃身和铃舌. 为了知道一个惊鸟铃的质量, 可以通过计算该惊鸟铃的体积, 然后由物理学知识计算出该惊鸟铃的质量, 因此我们需要作出一些合理的假设:

    假设 1: 铃身可近似看作由一个较大的圆锥挖去一个较小的圆锥;

    假设 2: 两圆锥的轴在同一条直线上;

    假设 3: 铃身内部有一个挂铃舌的部位的体积忽略不计. 截面图如下，其中 ${O}_{1}{O}_{3} = {20}{{cm}},{O}_{2}{O}_{3} = {18}{{cm}},{AB} = {16}{{cm}}$ ，则制作 100 个这样的惊鸟铃的铃身至少需要\blkda[2]{$120$} ${{kg}}$ 铜. (铜的密度为 ${8.9}{g}/{{{cm}}}^{3}$ ，结果精确到个位)
    \twotu{\includegraphics[width=0.9\textwidth]{bodyxb/src/146_229_349_933_332_0.jpg}}{\includegraphics[width=0.3\textwidth]{bodyxb/src/146_1194_308_252_330_0.jpg}}
    \begin{jieda}
        由题意知一个惊鸟铃的体积为 $\dfrac{1}{3}\pi  \times  {\left( \dfrac{16}{2}\right) }^{2} \times  \left( {{20} - }\right.  {18}) = \dfrac{128\pi }{3}$ ,则制作 100 个这样的惊鸟铃的铃身需要的铜的重量为 $\dfrac{\dfrac{128\pi }{3} \times  {8.9} \times  {100}}{1000} \approx  {119.3}$ ,则制作 100 个这样的惊鸟铃的铃身至少需要 120 千克铜.
    \end{jieda}

    \item (2024$\cdot$徐汇区一模$\cdot$11)已知一个棱长为 $a$ 的正方体木块可以在一个封闭的圆锥形容器内任意转动，若圆锥的底面半径为 3、母线长为 6，则实数 $a$ 的最大值为\blkda[2]{2}.
    \begin{jieda}
        棱长为 $a$ 的正方体绕其中心任意转动时,其顶点均在正方体外接球上,此正方体外接球半径为 $R =  \dfrac{1}{2}\sqrt{{a}^{2} + {a}^{2} + {a}^{2}} = \dfrac{\sqrt{3}a}{2}$ ,故问题转化为半径为 $\dfrac{\sqrt{3}a}{2}$ 的球体置于底面半径为 2,母线长为 4 的圆锥内,求球体半径的最大值. 当且仅当球体内切于圆锥时,球体的半径最大, 画出圆锥的轴截面图, 球的截面圆内切于圆锥, 如图所示, 由于圆锥底面半径为 3 , 母线长为 6 , 故其轴截面为一边长 6 的正三角形,其内切圆半径为 $\sqrt{3}$ ,即 $R = \sqrt{3} = \dfrac{\sqrt{3}a}{2} \Longrightarrow  a = 2$.
        \begin{center}
        \includegraphics[width=0.2\textwidth]{bodyxb/src/108_1362_532_214_209_0.jpg}
        \end{center}
    \end{jieda}

    \item (2024$\cdot$嘉定区一模$\cdot$12)正四棱台 ${ABCD} - {A}_{1}{B}_{1}{C}_{1}{D}_{1},{AB} = 3,{A}_{1}{B}_{1} = 1,{A{A}_{1}} = 2,M$ 是 ${C}_{1}{D}_{1}$ 的中点,在直线 $A{A}_{1}\text{ 、 }{BC}$ 上各取一个点 $P\text{ 、 }Q$ ,使得 $M\text{ 、 }P\text{ 、 }Q$ 三点共线,则线段 ${PQ}$ 的长度为\blkda[2]{}.

    \item (2024$\cdot$杨浦区一模$\cdot$12)我国古代数学著作《九章算术》中研究过一种叫“鳖(biē)臑(nào)”的几何体，它指的是由四个直角三角形围成的四面体，那么在一个长方体的八个顶点中任取四个，所组成的四面体中“鳖臑”的个数是\blkda[2]{}.

    \item (2024$\cdot$闵行区二模$\cdot$12)已知空间中有 2 个相异的点，现每增加一个点使得其与原有的点连接成尽可能多的等边三角形. 例如，空间中 3 个点最多可连接成 1 个等边三角形，空间中 4 个点最多可连接成 4 个等边三角形. 当增加到 8 个点时，空间中这 8 个点最多可连接成\blkda[2]{}个等边三角形.

    \item (2023$\cdot$上海秋季卷$\cdot$12)空间中有三个点 $A$ 、 $B$ 、 $C$ ，且 ${AB} = {BC} = {CA} = 1$ ，在空间中任取 2 个不同的点，使得它们与 $A$ 、 $B$ 、 $C$ 恰好成为一个正四棱锥的五个顶点，则不同的取法有\blkda[2]{}种.

    \item (2024$\cdot$闵行区一模$\cdot$12)已知点 $P$ 在正方体 ${ABCD} - {A}_{1}{B}_{1}{C}_{1}{D}_{1}$ 的表面上， $P$ 到三个平面 ${ABCD}\text{ 、 }{AD}{D}_{1}{A}_{1}\text{ 、 }{AB}{B}_{1}{A}_{1}$ 中的两个平面的距离相等，且 $P$ 到剩下一个平面的距离与 $P$ 到此正方体的中心的距离相等,则满足条件的点 $P$ 的个数为\blkda[2]{}.

    \item (2024 新清区二模$\cdot$12 )如图所示，在棱长为 1 的正方体 ${ABCD} - {A}_{1}{B}_{1}{C}_{1}{D}_{1}$ 中， $P$ 、 $Q$ 、 $R$ 在棱${AB}\text{ 、 }{BC}\text{ 、 }B{B}_{1}$ 上,且 ${PB} = \dfrac{1}{2},{QB} = \dfrac{1}{3},{RB} = \dfrac{1}{4}$ ,以 $\bigtriangleup {PQR}$ 为底面作一个三棱柱 ${PQR} -  {P}_{1}{Q}_{1}{R}_{1}$ ,使点 ${P}_{1}\text{ 、 }{Q}_{1}\text{ 、 }{R}_{1}$ 分别在平面 ${A}_{1}{AD}{D}_{1}\text{ 、 }{D}_{1}{DC}{C}_{1}\text{ 、 }{A}_{1}{B}_{1}{C}_{1}{D}_{1}$ 上,则这个三棱柱的侧棱长为\blkda[2]{}.

\begin{center}
\includegraphics[width=0.2\textwidth]{bodyxb/src/146_293_1737_303_259_0.jpg}
\end{center}


第 46 题图

\begin{center}
\includegraphics[width=0.2\textwidth]{bodyxb/src/146_697_1686_323_310_0.jpg}
\end{center}


第 48 题图

\begin{center}
\includegraphics[width=0.2\textwidth]{bodyxb/src/146_1124_1682_319_312_0.jpg}
\end{center}


第 49 题图
\end{enumerate}
\subsubsection*{二、选择题}

50.(2025$\cdot$崇明区二模$\cdot$14)已知一个圆锥的轴截面是边长为 2 的等边三角形，则这个圆锥的侧面积为 \dotfill{(\kh{})}

\xx{}{}{}{}${2\pi }$ B. ${3\pi }$ C. ${4\pi }$

D. $\dfrac{\sqrt{3}}{3}\pi$

51.(2024$\cdot$普陀区二模$\cdot$14)若一个圆锥的体积为 $\dfrac{2\sqrt{2}\pi }{3}$ ，用通过该圆锥的轴的平面截此圆锥，得到的截面三角形的顶角为 $\dfrac{\pi }{2}$ ,则该圆锥的侧面积为 \dotfill{(\kh{})}

\xx{}{}{}{}$\sqrt{2}\pi$ B. ${2\pi }$ C. $2\sqrt{2}\pi$ D. $4\sqrt{2}\pi$

52.(2024$\cdot$长宁区一模$\cdot$16)豆腐发酵后表面长出一层白绒绒的长毛就成了毛豆腐，将三角形豆腐 ${ABC}$ 悬空挂在发酵空间内，记发酵后毛豆腐所构成的几何体为 $T$. 若忽略三角形豆腐 ${ABC}$ 的厚度,设 ${AB} = 3\text{ 、 }{BC} = 4\text{ 、 }{AC} = 5$ ,点 $P$ 在 $\bigtriangleup {ABC}$ 内部. 假设对于任意点 $P$ ,满足 ${PQ} \leq  1$ 的点 $Q$ 都在 $T$ 内,且对于 $T$ 内任意一点 $Q$ ,都存在点 $P$ ,满足 ${PQ} \leq  1$ ,则 $T$ 的体积为 \dotfill{(\kh{})}

\xx{}{}{}{}${12} + {7\pi }$ B. ${12} + \dfrac{22\pi }{3}$ C. ${14} + {7\pi }$

D. ${14} + \dfrac{22\pi }{3}$

53.(2024$\cdot$虹口区一模$\cdot$15)“阿基米德多面体”也称为半正多面体，是由边数不全相同的正多边形为面围成的多面体, 它体现了数学的对称美. 如图所示, 将正方体沿同顶点出发的三条棱的中点截去一个三棱锥, 共可截去 8 个三棱锥, 得到 8 个面为正三角形、6 个面为正方形的一种半正多面体. 若 ${AB} = \sqrt{2}$ ,则此半正多面体外接球的表面积为 \dotfill{(\kh{})}

\begin{center}
\includegraphics[width=0.2\textwidth]{bodyxb/src/147_1169_1245_281_270_0.jpg}
\end{center}


第 53 题图

\xx{}{}{}{}$4\sqrt{3}\pi$ B. ${12\pi }$

C. $\dfrac{8\sqrt{2}\pi }{3}$ D. ${8\pi }$

54.(2024$\cdot$浦东新区一模$\cdot$15)已知棱长均为 1 的正 $n$ 棱柱有 2n 个顶点，从中任取两个顶点作为向量 $\vv{a}$ 的起点与终点,设底面的一条棱为 ${AB}$. 若集合 ${A}_{n} = \{ x \mid  x = \vv{a} \cdot  \vv{AB}\}$ ,则当 ${A}_{n}$ 中的元素个数最少时, $n$ 的值为 \dotfill{(\kh{})}

\xx{}{}{}{}3 B. 4 C. 6 D. 8

55.(2024$\cdot$嘉定区一模$\cdot$16)已知四面体 ${ABCD},{AB} = {BC},{AD} = {CD}$. 分别对于下列三个条件:  \ding{172}  ${AD}\bot {BC}$ ； \ding{173}  ${AC} = {BD}$ ； \ding{174}  $A{B}^{2} + C{D}^{2} = A{C}^{2} + B{D}^{2}$ ，是 ${AB}\bot {CD}$ 的充要条件的共有几个 \dotfill{(\kh{})}

\xx{}{}{}{}0 B. 1 C. 2 D. 3

56.(2024$\cdot$徐汇区二模$\cdot$16)三棱锥 $P - {ABC}$ 各顶点均在半径为 $2\sqrt{2}$ 的球 $O$ 的表面上， ${AB} = {AC} =  2\sqrt{2},\angle {BAC} = {90}^{ \circ  }$ ，二面角 $P - {BC} - A$ 的大小为 ${45}^{ \circ  }$ ，则对以下两个命题，判断正确的是\blkda[2]{}\dotfill{(\kh{})}  \ding{172}  三棱锥 $O - {ABC}$ 的体积为 $\dfrac{8}{3}$ ； \ding{173}  点 $P$ 形成的轨迹长度为 $2\sqrt{6}\pi$.

\xx{}{}{}{} \ding{172}  \ding{173} 都是真命题 B.  \ding{172} 是真命题， \ding{173} 是假命题

C.  \ding{172} 是假命题， \ding{173} 是真命题 D.  \ding{172}  \ding{173} 都是假命题

57.(2025$\cdot$虹口区二模$\cdot$16)在空间中，点 $O$ 为定点，设集合 $S = \left\{  {P\left| {{\left| \vv{OP}\right| }^{2} - 2\vv{OA} \cdot  \vv{OP} \leq  1}\right. }\right.$ ， $\left| \vv{OA}\right|  = 1\}$ ,则以下说法正确的是 \dotfill{(\kh{})}

 \ding{172}  若 $\vv{OP}$ 在 $\vv{OA}$ 上的数量投影为 $- \dfrac{1}{5}$ ，则线段 ${OP}$ 在运动过程中所形成的几何体体积为 $\dfrac{14}{375}\pi$ ；

 \ding{173}  对于任意的 ${P}_{i} \in  S$ 以及任意的正实数 ${a}_{i}$ ，设 $\vv{OQ} = \mathop{\sum }\limits_{{i = 1}}^{4}{a}_{i}\vv{O{P}_{i}}$ ，若 $\mathop{\sum }\limits_{{i = 1}}^{4}{a}_{i} = 1$ ，则 $Q \in  S$.

\xx{}{}{}{} \ding{172}  \ding{173} 都是真命题 B.  \ding{172} 是真命题， \ding{173} 是假命题

C.  \ding{172} 是假命题， \ding{173} 是真命题 D.  \ding{172}  \ding{173} 都是假命题

\subsubsection*{三、解答题}

58.(2024$\cdot$上海秋季卷$\cdot$17)如图为正四棱锥 $P - {ABCD}$ ， $O$ 为底面 ${ABCD}$ 的中心.

(1)若 ${AP} = 5,{AD} = 3\sqrt{2}$ ，求 $\bigtriangleup  {POA}$ 绕 ${PO}$ 旋转一周形成的几何体的体积；

(2)若 ${AP} = {AD},E$ 为 ${PB}$ 的中点，求直线 ${BD}$ 与平面 ${AEC}$ 所成角的大小.

\begin{center}
\includegraphics[width=0.2\textwidth]{bodyxb/src/148_1095_1085_328_298_0.jpg}
\end{center}


第 58 题图

59.(2025$\cdot$徐汇区二模$\cdot$17)如图， ${ABCD} - {A}_{1}{B}_{1}{C}_{1}{D}_{1}$ 是一块正四棱台形铁料，上、下底面的边长分别为 ${20}{{cm}}$ 和 ${40}{{cm}}$ ,高 ${30}{{cm}}$.

\begin{center}
\includegraphics[width=0.5\textwidth]{bodyxb/src/149_487_303_637_276_0.jpg}
\end{center}


第 59 题图

(1)求正四棱台 ${ABCD} - {A}_{1}{B}_{1}{C}_{1}{D}_{1}$ 的侧面 ${BC}{C}_{1}{B}_{1}$ 与底面 ${ABCD}$ 所成二面角的大小；

(2)现削去部分铁料(不计损耗)，将原正四棱台打磨为一个圆台，使得该圆台的上、下底面分别为原正四棱台上、下底面正方形的内切圆及其内部. 求削去部分与原正四棱台的体积之比.

60. (2025$\cdot$青浦区二模$\cdot$18)如图所示，已知四棱锥 $S - {ABCD}$ 的底面为菱形， $\angle {BAD} = \dfrac{\pi }{3},{AS} = {CS}$.

\begin{center}
\includegraphics[width=0.3\textwidth]{bodyxb/src/149_1064_1371_375_329_0.jpg}
\end{center}


第 60 题图

(1)求证: ${AC}\bot$ 平面 ${BDS}$ ；

(2)若 ${AB} = 2,{BS} = \sqrt{3},{DS} = 1$ ，求四棱锥 $S - {ABCD}$ 的体积.

61. (2025$\cdot$上海秋季卷$\cdot$18)如图所示， $P$ 是圆锥的顶点， $O$ 是底面圆心， ${AB}$ 是底面直径，且 ${AB} = 2$.

(1)若直线 ${PA}$ 与圆锥底面的所成角为 $\dfrac{\pi }{3}$ ，求圆锥的侧面积；

(2)已知 $Q$ 是母线 ${PA}$ 的中点，点 $C$ 、 $D$ 在底面圆周上，且弧 ${AC}$ 的长为 $\dfrac{\pi }{3},{CD}//{AB}$. 设点 $M$ 在线段 ${OC}$ 上，证明:直线 ${QM}//$ 平面 ${PBD}$.

\begin{center}
\includegraphics[width=0.2\textwidth]{bodyxb/src/150_1194_349_270_290_0.jpg}
\end{center}


第 61 题图

62. (2025$\cdot$静安区一模$\cdot$19)如图所示，正三棱锥 $A - {BCD}$ 的侧面是边长为 2 的正三角形.

(1)求正三棱锥 $A - {BCD}$ 的体积 $V$ ；

(2)设 $E\text{ 、 }F\text{ 、 }G$ 分别是线段 ${AC}\text{ 、 }{AD}\text{ 、 }{BC}$ 的中点.

求证: \ding{172}  ${CD}//$ 平面 ${EFG}$ ； \ding{173}  若平面 ${EFG}$ 交 ${BD}$ 于点 $H$ ，则四边形 ${EFHG}$ 是正方形.

\begin{center}
\includegraphics[width=0.2\textwidth]{bodyxb/src/150_1131_1293_335_315_0.jpg}
\end{center}


第 62 题图

\subsection*{考点2 空间向量及其应用}

\subsubsection*{一、填空题}

1.(2025$\cdot$徐汇区一模$\cdot$4)已知向量 $\vv{a} = \left( {2,5,1}\right) ,\vv{b} = \left( {4,m,5}\right)$ ，若 $\vv{a} \cdot  \vv{b} = 3$ ，则实数 $m$ 的值为\blkda[2]{$-2$}.

2.(2024$\cdot$长宁区一模$\cdot$9)若向量 $\vv{a} = \left( {1,0,2}\right) ,\vv{b} = \left( {0,1, - 1}\right)$ ，则 $\vv{a}$ 在 $\vv{b}$ 方向上的投影向量的坐标为\blkda[2]{$\left( {0, - 1,1}\right)$}.
\begin{jieda}
    $\vv{a} \cdot  \vv{b} =  - 2$ , $\left| \vv{b}\right|  = \sqrt{2}$ ,所以 $\vv{a}$ 在 $\vv{b}$ 方向上的投影向量为 $\dfrac{\vv{a} \cdot  \vv{b}}{\left| \vv{b}\right| } \cdot  \dfrac{\vv{b}}{\left| \vv{b}\right| } =  - \vv{b} = \left( {0, - 1,1}\right)$.
\end{jieda}

3.(2025$\cdot$浦东新区二模$\cdot$11)已知 $\vv{a}$ 、 $\vv{b}$ 、 $\vv{c}$ 为空间中三个单位向量，且 $\vv{a} \cdot  \vv{b} = \vv{b} \cdot  \vv{c} = \vv{c} \cdot  \vv{a} = 0$ ，若向量 $\vv{p}$ 满足 $\left| {\vv{p} - 2\vv{a}}\right|  = \dfrac{3}{2},\left| {\vv{p} - 2\vv{b}}\right|  = \dfrac{3}{2}$ ，则向量 $\vv{p}$ 与向量 $\vv{c}$ 夹角的最小值为\blkda[2]{$\arccos \dfrac{\sqrt{2}}{4}$}. (用反三角表示)
\begin{jieda}
    可设 $\vv{a} = \left( {1,0,0}\right) ,\vv{b} = \left( {0,1,0}\right) ,\vv{c} = \left( {0,0,1}\right)$ ,设 $\vv{p} =  \left( {x,y,z}\right)$ ,则 $\vv{p} - 2\vv{a} = \left( {x - 2,y,z}\right) ,\vv{p} - 2\vv{b} = \left( {x,y - 2,z}\right)$ , 所以 $\left\{  \begin{array}{l} {\left( x - 2\right) }^{2} + {y}^{2} + {z}^{2} = \dfrac{9}{4} \\  {x}^{2} + {\left( y - 2\right) }^{2} + {z}^{2} = \dfrac{9}{4} \end{array}\right.$ ,两式相减可得: $x = y$ ,再代入第一个式子,可得: $2{x}^{2} + {z}^{2} = {4x} - \dfrac{7}{4}$ ,设向量 $\vv{p}$ 与向量 $\vv{c}$ 夹角为 $\theta$ ,则 $\cos \theta  = \dfrac{z}{\sqrt{{x}^{2} + {y}^{2} + {z}^{2}}} = \dfrac{\sqrt{{4x} - \dfrac{7}{4} - 2{x}^{2}}}{\sqrt{{4x} - \dfrac{7}{4}}} =  \sqrt{1 - \dfrac{2{x}^{2}}{{4x} - \dfrac{7}{4}}} = \sqrt{1 - \dfrac{2}{\dfrac{4}{x} - \dfrac{7}{4{x}^{2}}}}$ ,易知对于 $y =  - \dfrac{7}{4{x}^{2}} + \dfrac{4}{x}$ 当 $\dfrac{1}{x} = \dfrac{8}{7}$ 即 $x = \dfrac{7}{8}$ 取得最大值 $\dfrac{16}{7}$ ,此时 $\sqrt{1 - \dfrac{2}{\dfrac{4}{x} - \dfrac{7}{4{x}^{2}}}}$ 取得最大值 $\dfrac{\sqrt{2}}{4}$ ，即 $\cos \theta$ 的最大值为 $\dfrac{\sqrt{2}}{4}$ ， $x = \dfrac{7}{8}$ 时取得，再由余弦函数的单调性可知 $\theta$ 的最小值为 $\arccos \dfrac{\sqrt{2}}{4}$.
\end{jieda}
1.(2025$\cdot$普陀区二模$\cdot$11)在棱长为 4 的正方体 ${ABCD} - {A}_{1}{B}_{1}{C}_{1}{D}_{1}$ 中， $\vv{AN} = 2\vv{AM} = \dfrac{1}{2}\vv{A{C}_{1}}$ ，若一动点 $G$ 满足 $\left| \vv{GN}\right|  = \sqrt{2}\left| \vv{GM}\right|$ ，则三棱锥 $G - {A}_{1}{B}_{1}{D}_{1}$ 体积的最大值为\blkda[2]{}.

2.(2024$\cdot$黄浦区二模$\cdot$12)在四面体 ${PABC}$ 中， $2\vv{PD} = \vv{PA} + \vv{PB},5\vv{PE} = 2\vv{PB} + 3\vv{PC},2\vv{PF} =  - \vv{PC} + 3\vv{PA}$ ，设四面体 ${PABC}$ 与四面体 ${PDEF}$ 的体积分别为 ${V}_{1}$ 、 ${V}_{2}$ ，则 $\dfrac{{V}_{2}}{{V}_{1}}$ 的值为\blkda[2]{}.

3. (2024$\cdot$普陀区一模$\cdot$12 )体积为 $\dfrac{\sqrt{2}}{12}{a}^{3}$ 的正四面体内有一个球 $O$ ，球 $O$ 与该正四面体的各面均有且只有一个公共点, $M\text{ 、 }N$ 是球 $O$ 的表面上的两动点,点 $P$ 在该正四面体的表面上运动,当 $\left| \vv{MN}\right|$ 最大时， $\vv{PM}$ $\cdot$ $\vv{PN}$ 的最大值为\blkda[2]{}.

\subsubsection*{二、选择题}

\begin{center}
\includegraphics[width=0.2\textwidth]{bodyxb/src/151_1177_1042_256_304_0.jpg}
\end{center}


第 4 题图

4.(2025$\cdot$奉贤区二模$\cdot$14)如图所示，在平行六面体 ${ABCD} - {A}_{1}{B}_{1}{C}_{1}{D}_{1}$ 中，点 $N$ 在对角线 ${A}_{1}C$ 上，点 $M$ 在对角线 ${A}_{1}B$ 上， $\vv{{A}_{1}N} = \dfrac{1}{3}\vv{NC}$ ， $\vv{{A}_{1}M} = \dfrac{1}{2}\vv{MB}$ ，以下命题正确的是 \dotfill{(\kh{})}

\xx{}{}{}{}$\vv{MN}//\vv{BC}$

B. ${D}_{1}\text{ 、 }N\text{ 、 }M$ 三点共线

C. ${D}_{1}M$ 与 ${A}_{1}C$ 是异面直线

\begin{center}
\includegraphics[width=0.2\textwidth]{bodyxb/src/151_1173_1408_262_236_0.jpg}
\end{center}


第 5 题图

D. $\vv{{D}_{1}N} = \dfrac{1}{2}\vv{NM}$

5.(2025$\cdot$黄浦区二模$\cdot$14)如图所示，在平行六面体 ${ABCD} - {A}_{1}{B}_{1}{C}_{1}{D}_{1}$ 中，设 $\vv{a} = \vv{A{A}_{1}},\vv{b} = \vv{D{B}_{1}}$ ，若 $\vv{a}$ 、 $\vv{b}$ 、 $\vv{c}$ 组成空间向量的一个基底，则 $\vv{c}$ 可以是\dotfill{(\kh{})}。

\xx{}{}{}{}$\vv{B{B}_{1}}$ B. $\vv{B{C}_{1}}$

C. $\vv{BD}$ D. $\vv{B{D}_{1}}$

\begin{center}
\includegraphics[width=0.3\textwidth]{bodyxb/src/151_1047_1698_388_337_0.jpg}
\end{center}


第 6 题图

6.(2025$\cdot$虹口区一模$\cdot$15)已知边长为 2 的正四面体 $A - {BCD}$ 的内切球 (球面与四面体四个面都相切的球) 的球心为 $O$ ，若空间中的动点 $P$ 满足 $\vv{OP} = x\vv{OC} + y\vv{OB} + z\vv{OD},x\text{ 、 }y\text{ 、 }z \in  \left\lbrack  {0,1}\right\rbrack$ ,则点 $P$ 的轨迹所形成的几何体的体积为 \dotfill{(\kh{})}

\xx{}{}{}{}$\sqrt{2}$ B. $\dfrac{\sqrt{2}}{3}$

C. $2\sqrt{3}$ D. $\dfrac{\sqrt{3}}{3}$


7. (2025$\cdot$闵行区二模$\cdot$16)设 $n$ 为正整数，空间中 $n$ 个单位向量构成集合 ${A}_{n} = \left\{  {{\vv{a}}_{1},{\vv{a}}_{2},\cdots ,{\vv{a}}_{n}}\right\}$ ，若存在实数 $t$ ,满足对任意 ${\vv{a}}_{i} \in  {A}_{n},{\vv{a}}_{j} \in  {A}_{n},{\vv{a}}_{i} \neq  {\vv{a}}_{j}$ ,都有 ${\vv{a}}_{i} \cdot  {\vv{a}}_{j} = t$ ,则当 $n$ 取得最大值时, $t$ 的值为 \dotfill{(\kh{})}
\xx{$- \dfrac{1}{2}$}{$\dfrac{1}{2}$}{$- \dfrac{1}{3}$}{$\dfrac{1}{3}$}
\begin{jieda}
    类比二维，直观可得，$n$为$4$，故解正四面体的两个外接球半径和棱构成的三角形可得.
\end{jieda}

\section*{直线、平面平行的判定与性质}

\subsubsection*{一、选择题}

1.(2024$\cdot$静安区二模$\cdot$14)设m、n是两条不同的直线，α、β是两个不同的平面，下列命题中真命题是 \dotfill{(\kh{})}

\xx{}{}{}{}若 $m//\alpha ,n//\alpha$ ,则 $m//n$ B. 若 $m \subset  \alpha ,n \subset  \beta ,m//n$ ,则 $\alpha //\beta$

C. 若 $m \bot  \alpha ,n//\alpha$ ,则 $m \bot  n$ D. 若 $m \subset  \alpha ,n \subset  \alpha ,m//\beta ,n//\beta$ ,则 $\alpha //\beta$

2.(2024$\cdot$长宁区二模$\cdot$14)已知直线 $a$ 、 $b$ 和平面 $\alpha$ ，则下列判断中正确的是 \dotfill{(\kh{})}

\begin{center}
\includegraphics[width=0.2\textwidth]{bodyxb/src/152_1139_688_294_456_0.jpg}
\end{center}


第 3 题图

\xx{}{}{}{}若 $a//\alpha ,b//\alpha$ ，则 $a//b$ B. 若 $a//b,b//a$ ，则 $a//\alpha$

C. 若 $a//\alpha ,b \bot  \alpha$ ，则 $\vv{a} \bot  \vv{b}$ D. 若 $a \bot  b,b//\alpha$ ，则 $a \bot  \alpha$

3.(2023$\cdot$上海春季卷$\cdot$15)如图所示， $P$ 是正方体 ${ABCD} - {A}_{1}{B}_{1}{C}_{1}{D}_{1}$ 边 ${A}_{1}{C}_{1}$ 上的动点，下列哪条边与 ${BP}$ 边始终异面 \dotfill{(\kh{})}

\xx{}{}{}{}$D{D}_{1}$

B. ${AC}$

C. $A{D}_{1}$

D. ${B}_{1}C$

4.(2025$\cdot$浦东新区一模$\cdot$14)设m、n为两条直线，α、β为两个平面，且 α、β $= n$. 下述四个命题中为假命题的是 \dotfill{(\kh{})}

\xx{}{}{}{}若 $m \bot  a$ ，则 $m \bot  n$ B. 若 $m//a$ ，则 $m//n$

C. 若 $m//\alpha$ 且 $m//\beta$ ，则 $m//n$ D. 若 $m//n$ ，则 $m//\alpha$ 或 $m//\beta$

5.(2024$\cdot$宝山区二模$\cdot$15)已知直线 $l$ 、 $m$ 、 $n$ 与平面 $\alpha$ 、 $\beta$ ，则下列命题中正确的是 \dotfill{(\kh{})}

\xx{}{}{}{}若 $\alpha //\beta ,l \subset  \alpha ,n \subset  \beta$ ,则 $l//n$ B. 若 $\alpha  \bot  \beta ,l \subset  \alpha$ ,则 $l \bot  \beta$

C. 若 $l \bot  \alpha ,l//\beta$ ，则 $\alpha  \bot  \beta$ D. 若 $l \bot  n,m \bot  n$ ,则 $l//m$

6.(2024$\cdot$青浦区一模$\cdot$15)已知直线 $m$ 、 $n$ ，平面 $\alpha$ 、 $\beta$. 给出下列命题:

 \ding{172}  若 $m\bot \alpha ,n\bot \beta$ ，且 $m\bot n$ ，则 $\alpha \bot \beta$ ；  \ding{173}  若 $m//\alpha ,n//\beta$ ，且 $m//n$ ，则 $\alpha //\beta$ ；

 \ding{174}  若 $m\bot \alpha ,n//\beta$ ，且 $m\bot n$ ，则 $\alpha \bot \beta$ ；  \ding{175}  若 $m\bot \alpha ,n//\beta$ ，且 $m//n$ ，则 $\alpha //\beta$.

其中正确的命题的个数是 \dotfill{(\kh{})}

\xx{}{}{}{}1 B. 2 C. 3 D. 4

\begin{center}
\includegraphics[width=0.2\textwidth]{bodyxb/src/152_1072_1681_335_315_0.jpg}
\end{center}


第 7 题图

7.(2024$\cdot$金山区一模$\cdot$15)如图，在正方体 ${ABCD} - {A}_{1}{B}_{1}{C}_{1}{D}_{1}$ 中， $E$ 、 $F$ 为正方体内 (含边界) 不重合的两个动点,下列结论错误的是 \dotfill{(\kh{})}

\xx{}{}{}{}若 $E \in  B{D}_{1},F \in  {BD}$ ,则 ${EF} \bot  {AC}$

B. 若 $E \in  B{D}_{1},F \in  {BD}$ ,则平面 ${BEF} \bot$ 平面 ${A}_{1}B{C}_{1}$

C. 若 $E \in  {AC},F \in  C{D}_{1}$ ,则 ${EF}//$ 平面 ${A}_{1}B{C}_{1}$

D. 若 $E \in  {AC},F \in  C{D}_{1}$ ,则 ${EF}//A{D}_{1}$

\subsubsection*{二、解答题}

8.(2023$\cdot$上海秋季卷$\cdot$17)直四棱柱 ${ABCD} - {A}_{1}{B}_{1}{C}_{1}{D}_{1},{AB}//{DC},{AB} \bot  {AD},{AB} = 2,{AD} = 3$ ， ${DC} = 4$.

(1)求证: ${A}_{1}B//$ 平面 ${DC}{C}_{1}{D}_{1}$ ；

(2)若四棱柱 ${ABCD} - {A}_{1}{B}_{1}{C}_{1}{D}_{1}$ 的体积为 36，求二面角 ${A}_{1} -  {BD} - A$ 的大小.

\begin{center}
\includegraphics[width=0.2\textwidth]{bodyxb/src/153_1071_437_341_278_0.jpg}
\end{center}


第 8 题图

9.(2024$\cdot$静安区二模$\cdot$19)如图1所示， ${ABCD}$ 是水平放置的矩形， ${AB} = {2\sqrt{3}},{BC} = {2\text{. 如图 2 所示， }}$ 将 ${ABD}$ 沿矩形的对角线 ${BD}$ 向上翻折，使得平面 ${ABD} \bot$ 平面 ${BCD}$.

(1)求四面体 ${ABCD}$ 的体积 $V$ ；

(2)试判断与证明以下两个问题:

 \ding{172}  在平面 ${BCD}$ 上是否存在经过点 $C$ 的直线 $l$ ，使得 $l\bot {AD}$ ?

 \ding{173}  在平面 ${BCD}$ 上是否存在经过点 $C$ 的直线 $l$ ，使得 $l//{AD}$ ?

\begin{center}
\includegraphics[width=0.4\textwidth]{bodyxb/src/153_846_1364_569_251_0.jpg}
\end{center}


第 9 题图

10.(2024$\cdot$浦东新区一模$\cdot$18)如图所示，在四棱锥 $P - {ABCD}$ 中， ${DP}\bot$ 底面 ${ABCD},{AB}//{CD}$ ， $\angle {ADC} = {90}^{ \circ  },{DP} = 3,{CD} = 2,{AB} = {AD} = 1$ ,点 $F$ 为 ${PD}$ 中点.

\begin{center}
\includegraphics[width=0.3\textwidth]{bodyxb/src/154_1065_288_395_279_0.jpg}
\end{center}


第 10 题图

(1)求证:直线 ${AF}//$ 平面 ${PBC}$ ;

(2)求点 $D$ 到平面 ${PBC}$ 的距离.

11. (2025$\cdot$虹口区一模$\cdot$18)如图所示，已知在四棱柱 ${ABCD} - {EFGH}$ 中， ${EA}\bot$ 平面 ${ABCD}$ ， $N$ 、 $M$ 分别是 ${EF}$ 、 ${HD}$ 的中点.

(1)求证: ${HN}//$ 平面 ${AFM}$ ；

(2)若底面 ${ABCD}$ 为梯形， ${AB}//{CD},{AB} = {EA} = 2,{AD} = {DC} = 1$ ，异面直线 ${AB}$ 与 ${EH}$ 所成角为 $\dfrac{\pi }{2}$. 求直线 ${AN}$ 与平面 ${AFM}$ 所成角的正弦值.

\begin{center}
\includegraphics[width=0.2\textwidth]{bodyxb/src/154_1166_1016_295_312_0.jpg}
\end{center}


第 11 题图

12.(2025$\cdot$金山区一模$\cdot$18)如图所示，在四棱锥 $P - {ABCD}$ 中，底面 ${ABCD}$ 是直角梯形， $\angle {BAD} =  \angle {CDA} = {90}^{ \circ  },{PA}\bot$ 平面 ${ABCD}$ ， $Q$ 是 ${PB}$ 的中点， ${PA} = {AD} = {DC} = 1$ ， ${AB} = {2.}$

\begin{center}
\includegraphics[width=0.3\textwidth]{bodyxb/src/154_1073_1636_388_239_0.jpg}
\end{center}


第 12 题图

(1)证明: ${CQ}//$ 平面 ${PAD}$ ；

(2)求点 $D$ 到平面 ${PAC}$ 的距离.

13.(2025$\cdot$嘉定区一模$\cdot$17)如图所示，在三棱柱 ${ABC} - {A}_{1}{B}_{1}{C}_{1}$ 中， ${AB} = {AC}$ ，侧面 $B{B}_{1}{C}_{1}C$ \_上底面 ${ABC}$ ,点 $E\text{ 、 }F$ 分别为棱 ${BC}$ 和 ${A}_{1}{C}_{1}$ 的中点.

(1)若底面 $\bigtriangleup  {ABC}$ 为边长为 2 的正三角形，且 $C{C}_{1} = {BC}$ ，侧棱 $C{C}_{1}$ 与底面 ${ABC}$ 所成的角为 ${60}^{ \circ  }$ ，求三棱柱 ${ABC} - {A}_{1}{B}_{1}{C}_{1}$ 的体积；

(2)求证: ${EF}//$ 平面 $A{A}_{1}{B}_{1}B$.

\begin{center}
\includegraphics[width=0.2\textwidth]{bodyxb/src/155_1075_390_320_275_0.jpg}
\end{center}


第 13 题图

\subsection*{考点4 直线、平面垂直的判定与性质}

\subsubsection*{一、选择题}

1.(2025$\cdot$长宁区一模$\cdot$14)已知非零空间向量 $\vv{a}$ 、 $\vv{b}$ 和 $\vv{c}$ ，则下列说法正确的是 \dotfill{(\kh{})}

\xx{}{}{}{}若 $\vv{a}\bot \vv{b},\vv{a}\bot \vv{c}$ ,则 $\vv{b}//\vv{c}$ B. 若 $\vv{a}\bot \vv{b},\vv{a}\bot \vv{c}$ ,则 $\vv{b}\bot \vv{c}$

C. 若 $\vv{a}\bot \vv{b},\vv{a}//\vv{c}$ ,则 $\vv{b}//\vv{c}$ D. 若 $\vv{a}\bot \vv{b},\vv{a}//\vv{c}$ ,则 $\vv{b}\bot \vv{c}$

2.(2025$\cdot$嘉定区二模$\cdot$14)已知平面 $\alpha$ 和平面 $\beta$ ，直线 $m \subset  \alpha$ ，直线 $n \subset  \beta$ ，则下列结论一定成立的是 \dotfill{(\kh{})}

\xx{}{}{}{}若 $m//n$ ，则 $\alpha //\beta$ B. 若 $m$ 与 $n$ 为异面直线，则 $\alpha //\beta$

C. 若 $m \bot  n$ ，则 $\alpha  \bot  \beta$ D. 若 $n \bot  \alpha$ ，则 $m \bot  n$

3.(2024$\cdot$上海春季卷$\cdot$14)空间中有两个不同的平面 $\alpha$ 、 $\beta$ ，两条不同的直线 $m$ 、 $n$ ，则下列说法正确的是 \dotfill{(\kh{})}

\xx{}{}{}{}若 $\alpha  \bot  \beta ,m \bot  \alpha ,n \bot  \beta$ ,则 $m \bot  n$ B. 若 $\alpha  \bot  \beta ,m \bot  \alpha ,m \bot  n$ ,则 $n \bot  \beta$

C. 若 $\alpha //\beta ,m//\alpha ,n//\beta$ ,则 $m//n$ D. 若 $\alpha //\beta ,m//\alpha ,m//n$ ,则 $n//\beta$

4.(2025$\cdot$崇明区一模$\cdot$14)已知直线 $l$ 和平面 $\alpha$ ，则“$l$ 垂直于 $\alpha$ 内的两条直线”是“$l\bot \alpha$ ”的\dotfill{(\kh{})}

\xx{}{}{}{}充分非必要条件 B. 必要非充分条件

C. 充要条件 D. 非充分非必要条件

5.(2025$\cdot$嘉定区一模$\cdot$14)已知 $\alpha$ 、 $\beta$ 是两个不同的平面， $a$ 、 $b$ 是两条不同的直线，下列条件中，一定得到直线 $l \bot  \alpha$ 的是 \dotfill{(\kh{})}

\xx{}{}{}{}$\alpha  \bot  \beta ,l//\beta$ B. $l \bot  a,a//a$

C. $l//a,a \bot  \alpha$ D. $l \bot  a,l \bot  b,a \subset  \alpha ,b \subset  \alpha$

6.(2024$\cdot$奉贤区一模$\cdot$13)已知两个不同的平面 $\alpha$ 和 $\beta$ ， $m$ 为平面 $\alpha$ 内的一条直线，则“$\alpha$ ⊥ $\beta$”是“$m$ 上 $\beta$ ”的 \dotfill{(\kh{})}

\xx{}{}{}{}充分非必要条件 B. 必要非充分条件

C. 充要条件 D. 既非充分也非必要条件.

7.(2024$\cdot$金山区二模$\cdot$15)如图，点 $N$ 为正方形 ${ABCD}$ 的中心， $\bigtriangleup  {ECD}$ 为正三角形，平面 ${ECD}\bot$ 平面 ${ABCD},M$ 是线段 ${ED}$ 的中点,则以下命题中正确的是 \dotfill{(\kh{})}

\begin{center}
\includegraphics[width=0.2\textwidth]{bodyxb/src/156_1134_457_295_197_0.jpg}
\end{center}


第 7 题图

\xx{}{}{}{}${BM} = {EN}$

B. ${CD} \bot  {MN}$

C. $A\text{ 、 }M\text{ 、 }N$ 三点共线

D. 直线 ${BM}$ 与 ${EN}$ 相交

\subsubsection*{二、解答题}

8.(2025$\cdot$宝山区一模$\cdot$17)如图所示，四棱锥 $P - {ABCD}$ 中，底面 ${ABCD}$ 为矩形， ${PA} = {PB} = {AD} =$ 3, ${AB} = 4$ ，且该四棱锥的体积为 $4\sqrt{5}$.

\begin{center}
\includegraphics[width=0.2\textwidth]{bodyxb/src/156_1099_907_329_285_0.jpg}
\end{center}


第 8 题图

(1)证明:平面 ${PAB} \bot$ 底面 ${ABCD}$ ；

(2)求异面直线 ${PC}$ 和 ${AB}$ 所成角的余弦值.

9.(2025$\cdot$黄浦区一模$\cdot$17)如图所示，在正方体 ${ABCD} - {A}_{1}{B}_{1}{C}_{1}{D}_{1}$ 中， $E$ 是 $B{C}_{1}$ 的中点.

\begin{center}
\includegraphics[width=0.2\textwidth]{bodyxb/src/156_1102_1511_331_317_0.jpg}
\end{center}


第 9 题图

(1)求证: $B{C}_{1} \bot$ 平面 ${CDE}$ ；

(2)求直线 ${DE}$ 与平面 ${ABCD}$ 所成角的大小.

10.(2025$\cdot$嘉定区二模$\cdot$17)如图所示，在四棱锥 $P - {ABCD}$ 中， ${PA}\bot$ 平面 ${ABCD},{PA} = {AC} = 2$ ， ${BC} = 1,{AB} = \sqrt{3}$.

(1)若 ${AD}//$ 平面 ${PBC}$ ，证明: ${AD}\bot {PB}$ ；

(2)在我国古代数学典籍《九章算术》中，记载了一种特殊的三棱锥——鳖臑，其四个面均为直角三角形，找出本题图中的一个鳖臑，并计算它的体积和表面积.

\begin{center}
\includegraphics[width=0.2\textwidth]{bodyxb/src/157_1189_402_272_271_0.jpg}
\end{center}


第 10 题图

11. (2025$\cdot$崇明区二模$\cdot$17)如图所示，在四棱锥 $P - {ABCD}$ 中，底面 ${ABCD}$ 是边长为 2 的正方形， 且 ${CB}\bot {BP},{CD}\bot {DP},{PA} = 2$ ，点 $E$ 、 $F$ 分别为 ${PB}$ 、 ${PD}$ 的中点.

\begin{center}
\includegraphics[width=0.2\textwidth]{bodyxb/src/157_1168_953_295_295_0.jpg}
\end{center}


第 11 题图

(1)求证: ${PA}\bot$ 平面 ${ABCD}$ ；

(2)求点 $P$ 到平面 ${AEF}$ 的距离.

12.(2025$\cdot$青浦区一模$\cdot$18)如图所示，在三棱锥 $P - {ABC}$ 中，平面 ${PAB}\bot$ 平面 ${ABC},{AB} = 6,{BC} =  2\sqrt{3},{AC} = 2\sqrt{6},D\text{ 、 }E$ 分别为线段 ${AB}\text{ 、 }{BC}$ 上的点,且 ${AD} = {2DB},{CE} = {2EB},{PD} \bot  {AC}$.

\begin{center}
\includegraphics[width=0.2\textwidth]{bodyxb/src/157_1211_1565_253_282_0.jpg}
\end{center}


第 12 题图

(1)求证: ${DE}//$ 平面 ${PAC}$ ；

(2)求证: ${PD} \bot$ 平面 ${ABC}$.

13.(2025$\cdot$长宁区一模$\cdot$18)如图所示，四棱柱 ${ABCD} - {A}_{1}{B}_{1}{C}_{1}{D}_{1}$ 的底面 ${ABCD}$ 是正方形， $O$ 是底面的中心， ${A}_{1}O \bot$ 平面 ${ABCD},{AB} = {A{A}_{1}} = \sqrt{2}$.

\begin{center}
\includegraphics[width=0.3\textwidth]{bodyxb/src/158_1047_307_372_226_0.jpg}
\end{center}


第 13 题图

(1)求证: ${A}_{1}C \bot$ 平面 ${BD}{D}_{1}{B}_{1}$ ；

(2)求直线 $O{A}_{1}$ 与平面 $A{A}_{1}B$ 所成角的正弦值.

14.(2025$\cdot$奉贤区一模$\cdot$19)如图所示，为正四棱锥 $P - {ABCD}$ ， $O$ 为底面 ${ABCD}$ 的中心.

(1)求证: ${CD}//$ 平面 ${PAB}$ ，平面 ${PAC}\bot$ 平面 ${PBD}$ ；

(2)设 $E$ 为 ${PB}$ 上的一点， $\vv{BE} = \dfrac{2}{3}\vv{BP}$.

在下面两问中选一个. 若都选，只按第 \ding{172} 问阅卷，第 1 问满分 5 分，第 2 问满分 7 分.

 \ding{172}  若 ${AD} = {AP} = 3\sqrt{2}$ ，求直线 ${EC}$ 与平面 ${BED}$ 所成角的大小；

 \ding{173}  已知平面 ${ECD}$ 与平面 ${ABCD}$ 所成锐二面角的大小为 $\arctan \dfrac{\sqrt{2}}{2}$ ，若 ${AD} = {3\sqrt{2}}$ ，求 ${AP}$ 的长.

\begin{center}
\includegraphics[width=0.5\textwidth]{bodyxb/src/158_811_1338_612_320_0.jpg}
\end{center}


第 14 题图

\subsection*{考点5 空间角与距离}

\subsubsection*{一、填空题}

1. (2024$\cdot$杨浦区二模$\cdot$7)正方体 ${ABCD} = {A}_{1}{B}_{1}{C}_{1}{D}_{1}$ 中，异面直线 ${AB}$ 与 $D{C}_{1}$ 所成角的大小为\blkda[2]{}.

\begin{center}
\includegraphics[width=0.2\textwidth]{bodyxb/src/159_1132_527_256_261_0.jpg}
\end{center}


第 3 题图

2.(2024$\cdot$崇明区一模$\cdot$7)在空间直角坐标系中，点 $P\left( {1, - 2,3}\right)$ 到 ${xOy}$ 平面的距离为\blkda[2]{}.

3.(2025$\cdot$上海春季卷$\cdot$9)已知 $P$ 是一个圆锥的顶点， ${PA}$ 是母线， ${PA} = \dfrac{PA}{Q\text{ ， }Q}$ 圆锥的底面半径是 $1.B\text{ 、 }C$ 分别在圆锥的底面上，则异面直线 ${PA}$ 与 ${BC}$ 所成角的最小值为\blkda[2]{}.

4.(2025$\cdot$金山区一模$\cdot$7)已知某圆锥的侧面展开图是圆心角为 $\sqrt{2}\pi$ 、半径为 2 的扇形，则该圆锥的母线与底面所成角的大小为\blkda[2]{}.

5.(2024$\cdot$奉贤区一模$\cdot$5)在四面体 $P - {ABC}$ 中，若底面 ${ABC}$ 的一个法向量为 $\vv{n} = \left( {1,1,0}\right)$ ，且 $\vv{CP} = \left( {2,2, - 1}\right)$ ，则顶点 $P$ 到底面 ${ABC}$ 的距离为\blkda[2]{}.

6.(2024$\cdot$崇明区二模$\cdot$8)已知底面半径为 1 的圆柱， $O$ 是其上底面圆心， $A$ 、 $B$ 是下底面圆周上两个不同的点， ${BC}$ 是母线. 若直线 ${OA}$ 与 ${BC}$ 所成角的大小为 $\dfrac{\pi }{3}$ ，则 ${BC} =$ \blkda[2]{}.

7.(2025$\cdot$崇明区一模$\cdot$9)在空间直角坐标系中，点(1，2，3)关于 ${xOy}$ 平面的对称点的坐标是\blkda[2]{}.

8.(2025$\cdot$黄浦区二模$\cdot$9)某商场要悬挂一个棱长为2m的正方体物件作为装饰，如图所示，A、B、C、D 为该正方体的顶点， $B{B}_{1}$ 、 $C{C}_{1}$ 、 $D{D}_{1}$ 为三根直绳索，且均垂直于屋顶所在平面 $\alpha$. 若平面 ${BCD}$ 与平面 $\alpha$ 平行，且点 $A$ 到 $\alpha$ 的距离为 $2{m}$ ，则直绳索 $B{B}_{1}$ 的长度约为\blkda[2]{}m.(结果精确到 ${0.01}{m}$ )

9.(2025$\cdot$徐汇区二模$\cdot$7)已知 ${PA}\bot$ 平面 ${ABC},{\bigtriangleup {ABC}}$ 是直角三角形，且 ${AB} = {AC} = 2,{PA} = 4$ ，则点 $P$ 到直线 ${BC}$ 的距离是\blkda[2]{}.

10.(2025$\cdot$虹口区一模$\cdot$9)如图所示，已知正三角形 ${ABC}$ 和正方形 ${BCDE}$ 的边长均为 2，且二面角 $A - {BC} - D$ 的大小为 $\dfrac{\pi }{6}$ ，则 $\vv{AC} \cdot  \vv{BD} =$ \blkda[2]{}.

11. (2025$\cdot$普陀区一模$\cdot$9)设 $\lambda  \in  \mathbf{R}$ ，在如图所示的平行六面体 ${ABCD} - {A}_{1}{B}_{1}{C}_{1}{D}_{1}$ 中， $\angle {A}_{1}{AB} =  \angle {A}_{1}{AD} = \angle {BAD} = \dfrac{\pi }{3},A{A}_{1} = 2,{AB} = {AD} = 1$ ,点 $M$ 是棱 ${C}_{1}{D}_{1}$ 的中点， $\vv{{A}_{1}N} = \lambda \vv{{A}_{1}{D}_{1}}$ ，若 $\vv{AM} \cdot  \vv{CN} = 2$ ,则 $\lambda$ 的值为\blkda[2]{}.

\begin{center}
\includegraphics[width=0.2\textwidth]{bodyxb/src/159_218_1720_294_292_0.jpg}
\end{center}


第 8 题图

\begin{center}
\includegraphics[width=0.2\textwidth]{bodyxb/src/159_619_1771_296_236_0.jpg}
\end{center}


第 10 题图

\begin{center}
\includegraphics[width=0.2\textwidth]{bodyxb/src/159_1017_1734_298_264_0.jpg}
\end{center}


第 11 题图

12.(2024$\cdot$静安区一模$\cdot$9)如图所示，在四棱锥 $P - {ABCD}$ 中， ${PA}\bot$ 平面 ${ABCD}$ ，底面 ${ABCD}$ 是矩形， $\left| {AP}\right|  = \left| {AB}\right|  = 2,\left| {AD}\right|  = 4,E$ 是 ${BC}$ 上的点， 直线 ${PB}$ 与平面 ${PDE}$ 所成的角是 $\arcsin \dfrac{\sqrt{3}}{6}$ ，则 ${BE}$ 的长为\blkda[2]{}.

\begin{center}
\includegraphics[width=0.2\textwidth]{bodyxb/src/160_1182_172_315_227_0.jpg}
\end{center}


第 12 题图

13.(2025$\cdot$嘉定区一模$\cdot$10)已知空间向量 $\vv{O{B}_{1}}$ 、 $\vv{O{B}_{2}}$ 、 $\vv{O{B}_{3}}$ 两两垂直，若空间点 $A$ 满足 $\left| \vv{A{B}_{1}}\right|  = \left| \vv{A{B}_{2}}\right|  = \left| \vv{A{B}_{3}}\right|  = 1$ ，记 $\vv{OP} = \vv{O{B}_{1}} + \vv{O{B}_{2}} + \vv{O{B}_{3}}$ ， 且 $\left| \vv{AP}\right|  \leq  1$ ，则 $\left| \vv{OA}\right|$ 的取值范围是\blkda[2]{}.

14.(2025$\cdot$黄浦区一模$\cdot$8)在正四面体 ${ABCD}$ 中，点 $N$ 是 $\bigtriangleup  {ABC}$ 的中心，若 $\vv{DN} = \lambda \vv{DA} + \mu \vv{DB} +  \nu \vv{BC}\left( {\lambda ,\mu ,\nu  \in  \mathbf{R}}\right)$ ，则 $\lambda  + \mu  + \nu  =$ \blkda[2]{}.

15.(2024$\cdot$上海春季卷$\cdot$10)在棱柱 ${ABCD} - {A}_{1}{B}_{1}{C}_{1}{D}_{1}$ 中，底面 ${ABCD}$ 为平行四边形， $\left| {A{A}_{1}}\right|  = 3$ ， $\left| {BD}\right|  = 4$ ，且 $\vv{A{B}_{1}} \cdot  \vv{BC} - \vv{A{D}_{1}} \cdot  \vv{DC} = 5$ ，则异面直线 $A{A}_{1}$ 与 ${BD}$ 的夹角为\blkda[2]{}.

16.(2025$\cdot$浦东新区一模$\cdot$11)已知空间中三个单位向量 $\vv{O{A}_{1}}$ 、 $\vv{O{A}_{2}}$ 、 $\vv{O{A}_{3}}$ ， $\vv{O{A}_{1}} \cdot  \vv{O{A}_{2}} = \vv{O{A}_{2}}$ 。 $\vv{O{A}_{3}} = \vv{O{A}_{3}} \cdot  \vv{O{A}_{1}} = 0,P$ 为空间中一点，且满足 $\left| {\vv{OP} \cdot  \vv{O{A}_{1}}}\right|  = 1,\left| {\vv{OP} \cdot  \vv{O{A}_{2}}}\right|  = 2,\left| {\vv{OP} \cdot  \vv{O{A}_{3}}}\right|  = 3$ , 则点 $P$ 个数的最大值为\blkda[2]{}.

17.(2024$\cdot$虹口区二模$\cdot$11)如图所示，在直四棱柱 ${ABCD} - {A}_{1}{B}_{1}{C}_{1}{D}_{1}$ 中，底面 ${ABCD}$ 为菱形，且 $\angle {BAD} = {60}^{ \circ  }$. 若 ${AB} = {A{A}_{1}} = 2$ ，点 $M$ 为棱 $C{C}_{1}$ 的中点，点 $P$ 在 ${A}_{1}B$ 上，则线段 ${PA}$ 、 ${PM}$ 的长度和的最小值为\blkda[2]{}.

18.(2025$\cdot$长宁区一模$\cdot$12)点 $P$ 、 $M$ 、 $N$ 分别位于正方体 ${ABCD} - {A}^{\prime }{B}^{\prime }{C}^{\prime }{D}^{\prime }$ 的面上， ${AB} = 1$ ，则 $\vv{PM} \cdot  \vv{PN}$ 的最小值是\blkda[2]{}.

19.(2025$\cdot$青浦区二模$\cdot$12)如图，正方体 ${ABCD} - {A}_{1}{B}_{1}{C}_{1}{D}_{1}$ 绕直线 $D{B}_{1}$ 旋转 $\dfrac{\pi }{3}$ ，直线 ${AB}$ 旋转至直线 ${A}^{\prime }{B}^{\prime }$ ，则直线 ${AB}$ 与直线 ${A}^{\prime }{B}^{\prime }$ 所成角的大小为\blkda[2]{}.

20.(2024$\cdot$浦东新区二模$\cdot$12)正三棱锥 S-ABC 中，底面边长 ${AB} = 2$ ，侧棱 ${AS} = 3$ ，向量 $\vv{a}$ 、 $\vv{b}$ 满足 $\vv{a} \cdot  \left( {\vv{a} + \vv{AC}}\right)  = \vv{a} \cdot  \vv{AB},\vv{b} \cdot  \left( {\vv{b} + \vv{AC}}\right)  = \vv{b} \cdot  \vv{AS}$ ，则 $\left| {\vv{a} - \vv{b}}\right|$ 的最大值为\blkda[2]{}.

\begin{center}
\includegraphics[width=0.2\textwidth]{bodyxb/src/160_170_1508_320_230_0.jpg}
\end{center}


第 15 题图

\begin{center}
\includegraphics[width=0.2\textwidth]{bodyxb/src/160_529_1447_309_293_0.jpg}
\end{center}


第 17 题图

\begin{center}
\includegraphics[width=0.2\textwidth]{bodyxb/src/160_880_1445_304_292_0.jpg}
\end{center}


第 19 题图

\begin{center}
\includegraphics[width=0.2\textwidth]{bodyxb/src/160_1228_1463_278_272_0.jpg}
\end{center}


第 20 题图

21. (2025$\cdot$宝山区二模$\cdot$12)空间中有相互垂直的两条异面直线 ${l}_{1}$ 、 ${l}_{2}$ ，点 $A$ 、 $B \in  {l}_{1}$ ， $C$ 、 $D \in  {l}_{2}$ ，且 ${AB} =$ 4， ${CD} = 1$ ，若 ${DA}\bot {DB}$ ，且 ${AC} = {BC} + 2$ ，则二面角 $D - {AB} - C$ 平面角的余弦值最小为\blkda[2]{}.

22.(2024$\cdot$宝山区二模$\cdot$12)在空间直角坐标系中，从原点出发的两个向量 $\vv{a}$ 、 $\vv{b}$ 满足: $\vv{a} \cdot  \vv{b} = 2$ ， $\left| \vv{b}\right|  = 1$ ， 且存在实数 $t$ ，使得 $\left| \vv{a}\right|  - 2\left| {\vv{a} + t\vv{b}}\right|  \geq  0$ 成立，则由 $\vv{a}$ 构成的空间几何体的体积是\blkda[2]{}.

23.(2024$\cdot$松江区一模$\cdot$12)已知正四面体 $A - {BCD}$ 的棱长为 2√2. 空间内任意点 $P$ 满足 $\left| {\vv{PB} + \vv{PC}}\right|  = 2$ ， 则 $\vv{AP} \cdot  \vv{AD}$ 的取值范围是\blkda[2]{}.

24.(2024$\cdot$奉贤区一模$\cdot$12)已知正方体 ${ABCD} - {A}_{1}{B}_{1}{C}_{1}{D}_{1}$ 的棱长为 $1,{\lambda }_{i} \in  \{  - 1,1\} (i = 1,2,3$ , 4),则 $\left| {{\lambda }_{1}\vv{AB} + {\lambda }_{2}\vv{BC} + {\lambda }_{3}\vv{A{C}_{1}} + {\lambda }_{4}\vv{B{D}_{1}}}\right|$ 的最大值是\blkda[2]{}.

25.(2024$\cdot$黄浦区一模$\cdot$12)若正三棱锥 $A - {BCD}$ 的底面边长为 6，高为 $\sqrt{13}$ ，动点 $P$ 满足(DA)+ $\vv{CB})\bot \left( {\vv{PA} + \vv{PB} + \vv{PC} + \vv{PD}}\right)$ ，则 $\left| {\vv{PA} + \vv{PB}}\right|  + 2\left| \vv{PA}\right|$ 的最小值为\blkda[2]{}.

26.(2023$\cdot$上海春季卷$\cdot$12)设 $\vv{OA}$ ， $\vv{OB}$ ， $\vv{OC}$ 为空间中三组单位向量，且 $\vv{OA}\bot \vv{OB}$ ， $\vv{OA}\bot \vv{OC}$ ， $\vv{OB}$ 与 $\vv{OC}$ 夹角为 ${60}^{ \circ  },P$ 为空间任意一点,且 $\left| \vv{OP}\right|  = 1$ ,满足 $\left| {\vv{OP} \cdot  \vv{OC}}\right|  \leq  \left| {\vv{OP} \cdot  \vv{OB}}\right|  \leq  \left| {\vv{OP} \cdot  \vv{OA}}\right|$ ,则 $\left| {\vv{OP} \cdot  \vv{OC}}\right|$ 的最大值为\blkda[2]{}.

\begin{center}
\includegraphics[width=0.2\textwidth]{bodyxb/src/161_1127_623_323_216_0.jpg}
\end{center}


第 27 题图

\subsubsection*{二、选择题}

27.(2025$\cdot$上海春季卷$\cdot$13)如图所示， ${ABCD} - {A}_{1}{B}_{1}{C}_{1}{D}_{1}$ 是正四棱台, 则下列各组直线中属于异面直线的是 \dotfill{(\kh{})}

\xx{}{}{}{}${AB}$ 和 ${C}_{1}{D}_{1}$ B. $A{A}_{1}$ 和 $C{C}_{1}$

C. $B{D}_{1}$ 和 ${B}_{1}D$ D. ${A}_{1}{D}_{1}$ 和 ${AB}$

28.(2025$\cdot$青浦区一模$\cdot$14)若点 $P\left( {x,y,z}\right)$ ( ${xyz} \neq  0$ )关于 ${xOy}$ 的对称点为 $A$ ，关于 $z$ 轴的对称点为 $B$ ,则 $A\text{ 、 }B$ 两点的对称是 \dotfill{(\kh{})}

\xx{}{}{}{}关于 ${xOz}$ 平面对称 B. 关于 $x$ 轴对称

C. 关于 $y$ 轴对称 D. 关于坐标原点对称

29.(2025$\cdot$奉贤区一模$\cdot$15)在四棱锥 S-ABCD 中，若 $\vv{SA} = x\vv{SB} + y\vv{SC} + z\vv{SD}$ ，则实数组(x, y, z) 可能是 \dotfill{(\kh{})}

\xx{}{}{}{}$\left( {1, - 1,1}\right)$ B. $\left( {1,0, - 1}\right)$ C. $\left( {1,0,0}\right)$ D. $\left( {1, - 1, - 1}\right)$

\begin{center}
\includegraphics[width=0.2\textwidth]{bodyxb/src/161_1163_1263_289_267_0.jpg}
\end{center}


第 30 题图

30. (2025$\cdot$宝山区一模$\cdot$15)如图所示，正四棱柱 ${ABCD} - {A}_{1}{B}_{1}{C}_{1}{D}_{1}$ 的底面 ${ABCD}$ 边长为 $1,E$ 为 ${AD}$ 上任意一点， $F$ 为 $C{C}_{1}$ 中点，若棱 ${C}_{1}{D}_{1}$ 上至少存在一点 $P$ 使得 ${PE} \bot  {PF}$ ,则棱长 $A{A}_{1}$ 的最大值为 \dotfill{(\kh{})}

\xx{}{}{}{}$\dfrac{\sqrt{2}}{2}$ B. 1

C. $\sqrt{2}$ D. 2

31.(2025$\cdot$黄浦区一模$\cdot$14)若从正方体八个顶点中任取四个顶点分别记为 $A\text{ 、 }B\text{ 、 }C\text{ 、 }D$ ,则直线 ${AB}$ 与 ${CD}$ 所成角的大小不可能为

\dotfill{(\kh{})}

\xx{}{}{}{}30 \textdegree  B. 45 \textdegree  C. 60 \textdegree  D. 90 \textdegree 

32.(2025$\cdot$静安区一模$\cdot$15)我国古代数学著作《九章算术》中将四个面都是直角三角形的空间四面体叫做“鳖臑”. 如图是一个水平放置的 $\bigtriangleup {ABC},{CD} \bot  {AB},\angle A = {30}^{ \circ  },\angle B = {45}^{ \circ  }$. 现将 ${{Rt}}\bigtriangleup {ACD}$ 沿 ${CD}$ 折起,使点 $A$ 移动到点 ${A}^{\prime }$ ,使得空间四面体 ${A}^{\prime }{BCD}$ 恰好是一个“鳖臑”，则二面角 ${A}^{\prime } - {CD} - B$ 的大小为 \dotfill{(\kh{})}

\begin{center}
\includegraphics[width=0.2\textwidth]{bodyxb/src/161_1125_1731_329_135_0.jpg}
\end{center}


第 32 题图

\xx{}{}{}{}60 \textdegree  B. ${90}^{ \circ  }$ C. arctan2

D. $\arccos \dfrac{\sqrt{3}}{3}$

33.(2025$\cdot$静安区一模$\cdot$16)在四棱锥 $P - {ABCD}$ 中， $\vv{AB} = \left( {4, - 2,3}\right) ,\vv{AD} = \left( {-4,1,0}\right) ,\vv{AP} =  \left( {-6,2, - 8}\right)$ ,则该四棱锥的高为 \dotfill{(\kh{})}

\xx{}{}{}{}4 B. 3 C. 2 D. 1

34.(2025$\cdot$黄浦区二模$\cdot$16)给定四面体 ${ABCD}$. 平面 $\alpha$ 满足: \ding{172}  $A$ 、 $B$ 、 $C$ 、 $D$ 四个点均不在平面 $\alpha$ 上，也不在 $\alpha$ 的同侧； \ding{173}  若平面 $\alpha$ 与四面体 ${ABCD}$ 的棱有公共点，则该公共点一定是此棱的中点或两个三等分点之一. 设 $A$ 、 $B$ 、 $C$ 、 $D$ 四个点到平面 $\alpha$ 的距离分别为 ${d}_{i}\left( {i = 1,2,3,4}\right)$ ，那么 ${d}_{i}$ 的所有不同值的个数组成的集合为 \dotfill{(\kh{})}

\xx{}{}{}{}$\{ 1,2,3,4\}$ B. $\{ 1,2,3\}$ C. $\{ 1,2\}$ D. $\{ 1\}$

35.(2024$\cdot$嘉定区一模$\cdot$16)已知四面体 ${ABCD},{AB} = {BC},{AD} = {CD}$. 分别对于下列三个条件:  \ding{172}  ${AD}\bot {BC}$ ； \ding{173}  ${AC} = {BD}$ ； \ding{174}  $A{B}^{2} + C{D}^{2} = A{C}^{2} + B{D}^{2}$ ，是 ${AB}\bot {CD}$ 的充要条件的共有几个 \dotfill{(\kh{})}

\xx{}{}{}{}0 B. 1 C. 2 D. 3

\subsubsection*{三、解答题}

36. (2023$\cdot$上海春季卷$\cdot$17)已知 ${AB}\bot {AC},{PA}\bot$ 平面 ${ABC},{PA} = {AB} = 3,{AC} = 4,M$ 为 ${BC}$ 中点， 过点 $M$ 分别作平行于平面 ${PAB}$ 的直线交 ${AC}$ 、 ${PC}$ 于点 $E$ 、 $F$.

\begin{center}
\includegraphics[width=0.2\textwidth]{bodyxb/src/162_1125_1007_335_290_0.jpg}
\end{center}


第 36 题图

(1)求直线 ${PM}$ 与平面 ${ABC}$ 所成的角的大小；

(2)证明:平面 ${MEF}//$ 平面 ${PAB}$ ,并求直线 ${ME}$ 到平面 ${PAB}$ 的距离.

37.(2024$\cdot$杨浦区二模$\cdot$17)如图所示， $P$ 为圆锥顶点， $O$ 为底面中心， $A$ 、 $B$ 、 $C$ 均在底面圆周上，且 $\bigtriangleup  {ABC}$ 为等边三角形.

\begin{center}
\includegraphics[width=0.2\textwidth]{bodyxb/src/163_1123_298_298_322_0.jpg}
\end{center}


第 37 题图

(1)求证:平面 ${POA} \bot$ 平面 ${PBC}$ ；

(2)若圆锥底面半径为 2，高为 $2\sqrt{2}$ ，求点 $A$ 到平面 ${PBC}$ 的距离.

38. (2025$\cdot$虹口区二模$\cdot$17)如图所示，在四棱锥 $P - {ABCD}$ 中， ${PA}\bot$ 平面 ${ABCD},{AB}\bot {AD},{AD}//  {BC},{AB} = {BC} = 2,{AD} = 4$.

(1)求证:平面 ${PAC} \bot$ 平面 ${PCD}$ ；

(2)若异面直线 ${PB}$ 和 ${CD}$ 所成角为 $\dfrac{\pi }{3}$ ，求点 $B$ 到平面 ${PCD}$ 的距离.

\begin{center}
\includegraphics[width=0.2\textwidth]{bodyxb/src/163_1093_1259_333_253_0.jpg}
\end{center}


第 38 题图

39. (2025$\cdot$普陀区二模$\cdot$17)如图所示，在三棱柱 ${ABC} - {{A}_{1}{B}_{1}{C}_{1}}$ 中， ${AC} = {A{A}_{1}} = {AB} = 2,\angle {{BA}{A}_{1}} =  \dfrac{\pi }{3}$ ,且 $C{A}_{1} = {CB},C{A}_{1} \bot  A{B}_{1}$.

\begin{center}
\includegraphics[width=0.2\textwidth]{bodyxb/src/164_1106_332_337_273_0.jpg}
\end{center}


第 39 题图

(1)求证:平面 ${CA}{B}_{1} \bot$ 平面 $A{A}_{1}{B}_{1}B$ ；

(2)求直线 $C{A}_{1}$ 与平面 $B{B}_{1}{C}_{1}C$ 所成的角的正弦值.

40. (2024$\cdot$杨浦区一模$\cdot$17)如图所示，在四棱锥 $P - {ABCD}$ 中， ${PA}\bot$ 平面 ${ABCD}$ ，底面 ${ABCD}$ 是正方形.

(1)求证:平面 ${PBD} \bot$ 平面 ${PAC}$ ；

(2)设 ${AB} = 2$ ，若四棱锥 $P - {ABCD}$ 的体积为 $\dfrac{8}{3}$ ，求点 $A$ 到平面 ${PBD}$ 的距离.

\begin{center}
\includegraphics[width=0.2\textwidth]{bodyxb/src/164_1172_1249_274_284_0.jpg}
\end{center}


第 40 题图

41. (2024$\cdot$松江区一模$\cdot$17)如图所示，在四棱锥 $P - {ABCD}$ 中， ${PA}\bot$ 底面 ${ABCD},{AB}\bot {AD}$ ，点 $E$ 在线段 ${AD}$ 上，且 ${CE}//{AB}$.

(1)求证: ${CE} \bot$ 平面 ${PAD}$ ；

(2)若四棱锥 $P - {ABCD}$ 的体积为 $\dfrac{5}{6}{AB} = 1,{AD} = 3,{CD} = \sqrt{2}$ ， $\angle {CDA} = {45}^{ \circ  }$ ,求二面角 $P - {CE} - A$ 的大小.

\begin{center}
\includegraphics[width=0.3\textwidth]{bodyxb/src/165_1060_344_356_249_0.jpg}
\end{center}


第 41 题图

42.(2025$\cdot$崇明区一模$\cdot$17)如图所示，在直三棱柱 ${ABC} - {A}_{1}{B}_{1}{C}_{1}$ 中， $E$ 、 $F$ 分别为 ${A}_{1}{C}_{1}$ 、 ${BC}$ 的中点， $A{A}_{1} = {AB} = {BC} = 2$ ， ${C}_{1}F\bot {AB}$.

\begin{center}
\includegraphics[width=0.2\textwidth]{bodyxb/src/165_1125_1239_292_332_0.jpg}
\end{center}


第 42 题图

(1)求证: ${C}_{1}F//$ 平面 ${ABE}$ ；

(2)求点 $C$ 到平面 ${ABE}$ 的距离.

43.(2024$\cdot$青浦区一模$\cdot$17)已知四棱锥 $P - {ABCD}$ ，底面 ${ABCD}$ 为正方形，边长为3， ${PD}\bot$ 平面 ${ABCD}$.

\begin{center}
\includegraphics[width=0.2\textwidth]{bodyxb/src/166_1155_295_308_233_0.jpg}
\end{center}


第 43 题图

(1)求证: ${BC}\bot$ 平面 ${CDP}$ ；

(2)若直线 ${AD}$ 与 ${BP}$ 所成的角大小为 ${60}^{ \circ  }$ ，求 ${DP}$ 的长.

44.(2024$\cdot$闵行区一模$\cdot$17)如图所示，在四棱锥 $P - {ABCD}$ 中，底面 ${ABCD}$ 是边长为 $a$ 的正方形， 侧面 ${PAD} \bot$ 底面 ${ABCD}$ ,且 ${PA} = {PD} = \dfrac{\sqrt{2}}{2}a$ ,设 $E$ 、 $F$ 分别为 ${PC}$ 、 ${BD}$ 的中点.

\begin{center}
\includegraphics[width=0.2\textwidth]{bodyxb/src/166_1139_1260_324_228_0.jpg}
\end{center}


第 44 题图

(1)证明:直线 ${EF}//$ 平面 ${PAD}$ ；

(2)求直线 ${PB}$ 与平面 ${ABCD}$ 所成的角的正切值.

45.(2024$\cdot$崇明区一模$\cdot$17)如图所示，四棱锥 $P - {ABCD}$ 中， ${PA}\bot$ 平面 ${ABCD}$ ， ${AB}//{CD}$ ， ${PA} =  {AB} = {AD} = 2,{CD} = 1,\angle {ADC} = {90}^{ \circ  },E\text{ 、 }F$ 分别为 ${PB}\text{ 、 }{AB}$ 的中点.

\begin{center}
\includegraphics[width=0.2\textwidth]{bodyxb/src/167_1096_287_329_267_0.jpg}
\end{center}


第 45 题图

(1)求证: ${CE}//$ 平面 ${PAD}$ ；

(2)求点 $B$ 到平面 ${PCF}$ 的距离.

46. (2024$\cdot$崇明区二模$\cdot$17)如图所示，在三棱锥 $P - {ABC}$ 中， ${AB} = {BC} = {2\sqrt{2}},{PA} = {PB} = {PC} = {AC} = 4$ ， $O$ 为 ${AC}$ 的中点.

\begin{center}
\includegraphics[width=0.2\textwidth]{bodyxb/src/167_1140_1163_290_304_0.jpg}
\end{center}


第 46 题图

(1)证明: ${PO}\bot$ 平面 ${ABC}$ ；

(2)若点 $M$ 在棱 ${BC}$ 上，且 ${MC} = {2MB}$ ，求点 $C$ 到平面 ${POM}$ 的距离.

47.(2025$\cdot$闵行区一模$\cdot$17)在直三棱柱 ${ABC} - {{A}_{1}{B}_{1}{C}_{1}}$ 中， ${AB} = {AC} = 2,{A{A}_{1}} = 3,\angle {BAC} = {90}^{ \circ  }$ ，连接 ${A}_{1}C,M\text{ 、 }E$ 分别为 ${A}_{1}C$ 和 ${BC}$ 的中点.

\begin{center}
\includegraphics[width=0.2\textwidth]{bodyxb/src/168_1162_300_290_341_0.jpg}
\end{center}


第 47 题图

(1)证明:直线 ${EM}//$ 平面 ${A}_{1}{AB}{B}_{1}$ ；

(2)求二面角 ${A}_{1} - {BC} - A$ 的大小.

48.(2024$\cdot$嘉定区二模$\cdot$17)如图所示，在三棱柱 ${ABC} - {A}_{1}{B}_{1}{C}_{1}$ 中， ${A}_{1}A\bot$ 平面 ${ABC},D$ 是 ${BC}$ 的中点, ${AC} = \sqrt{2},{A}_{1}A = {AB} = {BC} = 1$.

\begin{center}
\includegraphics[width=0.2\textwidth]{bodyxb/src/168_1138_1158_315_275_0.jpg}
\end{center}


第 48 题图

(1)求证: ${A}_{1}B//$ 平面 ${AD}{C}_{1}$ ；

(2)求直线 $D{C}_{1}$ 与 ${A}_{1}B$ 所成的角的大小.

49. (2024$\cdot$普陀区二模$\cdot$17)如图所示，在四棱锥 $S - {ABCD}$ 中，底面 ${ABCD}$ 是边长为 1 的正方形， ${SA} = {SB} = 2,E$ 、 $F$ 分别是 ${SC}$ 、 ${BD}$ 的中点.

(1)求证: ${EF}//$ 平面 ${SAB}$ ；

(2)若二面角 $S - {AB} - D$ 的大小为 $\dfrac{\pi }{2}$ ，求直线 ${SD}$ 与平面 ${ABCD}$ 所成角的大小.

\begin{center}
\includegraphics[width=0.2\textwidth]{bodyxb/src/169_1138_347_290_261_0.jpg}
\end{center}


第 49 题图

50. (2025$\cdot$宝山区二模$\cdot$17)如图所示，在四面体 ${ABCD}$ 中， $\bigtriangleup  {BCD}$ 是边长为 2 的正三角形，且 ${AB} =  {AD} = \sqrt{2}$.

(1)证明: ${BD} \bot  {AC}$ ；

(2)若 $E$ 是 ${BC}$ 的中点，且二面角 $A - {BD} - C$ 的大小为 $\dfrac{\pi }{2}$ ，求 ${AE}$ 与平面 ${BCD}$ 所成角的大小.

\begin{center}
\includegraphics[width=0.2\textwidth]{bodyxb/src/169_1137_1293_295_208_0.jpg}
\end{center}


第 50 题图

51.(2024$\cdot$虹口区二模$\cdot$18)如图所示，在三棱柱 ${ABC} - {A}_{1}{B}_{1}{C}_{1}$ 中， ${CA}\bot {CB}$ ， $D$ 为 ${AB}$ 的中点， ${CA} = {CB} = 2,C{C}_{1} = 3$.

\begin{center}
\includegraphics[width=0.2\textwidth]{bodyxb/src/170_1170_298_290_340_0.jpg}
\end{center}


第 51 题图

(1)求证: $A{C}_{1}//$ 平面 ${B}_{1}{CD}$ ；

(2)若 $C{C}_{1} \bot$ 平面 ${ABC}$ ，点 $P$ 在棱 $A{A}_{1}$ 上，且 ${PD} \bot$ 平面 ${B}_{1}{CD}$ ，求直线 ${CP}$ 与平面 ${B}_{1}{CD}$ 所成角的正弦值.

52. (2024$\cdot$黄浦区二模$\cdot$18)如图所示，在四棱锥 $P - {ABCD}$ 中，底面 ${ABCD}$ 为矩形，点 $E$ 是棱 ${PD}$ 上的一点， ${PB}//$ 平面 ${AEC}$.

(1)求证:点 $E$ 是棱 ${PD}$ 的中点；

(2)若 ${PA}\bot$ 平面 ${ABCD},{AP} = 2,{AD} = {2\sqrt{3}},{PC}$ 与平面 ${ABCD}$ 所成角的正切值为 $\dfrac{1}{3}$ ,求二面角 $D - {AE} - C$ 的大小.

\begin{center}
\includegraphics[width=0.2\textwidth]{bodyxb/src/170_1117_1246_339_253_0.jpg}
\end{center}


第 52 题图

53.(2025$\cdot$浦东新区一模$\cdot$18)如图所示，已知 ${AB}$ 为圆柱 $O{O}_{1}$ 底面圆 $O$ 的直径， ${OA} = 2$ ，母线 $A{A}_{1}$ 长为 3,点 $P$ 为底面圆 $O$ 的圆周上一点.

\begin{center}
\includegraphics[width=0.2\textwidth]{bodyxb/src/171_1172_291_291_271_0.jpg}
\end{center}


第 53 题图

(1)若 $\angle {BOP} = {90}^{ \circ  }$ ，求三棱锥 $A - {PB}{A}_{1}$ 的体积；

(2)若 $\angle {BOP} = {60}^{ \circ  }$ ，求异面直线 ${A}_{1}B$ 与 ${AP}$ 所成的角的余弦值.

54.(2025$\cdot$松江区一模$\cdot$18)如图所示，已知 ${AB}\bot$ 平面 ${ACD},{AB}//{DE}$ ， $\bigtriangleup  {ACD}$ 为等边三角形， ${AD} = {DE} = {2AB}$ ,点 $F$ 为 ${CD}$ 的中点.

\begin{center}
\includegraphics[width=0.2\textwidth]{bodyxb/src/171_1162_1145_299_262_0.jpg}
\end{center}


第 54 题图

(1)求证: ${AF}//$ 平面 ${BCE}$ ；

(2)求直线 ${BF}$ 和平面 ${ABC}$ 所成角的正弦值.

55. (2025$\cdot$徐汇区一模$\cdot$18)如图所示，在四棱锥 $P - {ABCD}$ 中， ${AD}//{BC}$ ， $\angle {ADC} = \angle {PAB} = \dfrac{\pi }{2}$ ， ${BC} = {CD} = \dfrac{1}{2}{AD},E$ 为棱 ${AD}$ 的中点，异面直线 ${PA}$ 与 ${CD}$ 所成角的大小为 $\dfrac{\pi }{2}$.

(1)求证: ${CD}//$ 平面 ${PBE}$ ；

( 2 )若二面角 $P - {CD} - A$ 的大小为 $\dfrac{\pi }{4}$ ，求直线 ${PA}$ 与平面 ${PCE}$ 所成角的正弦值.

\begin{center}
\includegraphics[width=0.2\textwidth]{bodyxb/src/172_1168_417_291_232_0.jpg}
\end{center}


第 55 题图

56.(2025$\cdot$黄浦区二模$\cdot$18)在四面体 ${ABCD}$ 中， ${DB} = {DC} = 2$ ， ${DB}\bot {DC}$.

(1)若 $\bigtriangleup  {ABC}$ 为正三角形，平面 ${ABC} \bot$ 平面 ${DBC}$ ，求四面体 ${ABCD}$ 的体积；

(2)若 ${AB} = {AC} = 4,{AD} = 3$ ，求二面角 $A - {BC} - D$ 的大小.

\begin{center}
\includegraphics[width=0.3\textwidth]{bodyxb/src/172_1100_1256_359_222_0.jpg}
\end{center}


第 56 题图

57.(2025$\cdot$杨浦区二模$\cdot$18)坐落于杨浦滨江的世界技能博物馆由百年历史文化保护建筑改建而成， 其中的支柱保留了原有的正八棱柱, 既考虑了结构力学优势, 又体现了对历史建筑的尊重和传承. 如图所示, ${O}_{1}O$ 分别为正八棱柱的上下两个底面的中心,已知 ${OA} = 1,A{A}_{1} = 4$.

\begin{center}
\includegraphics[width=0.2\textwidth]{bodyxb/src/173_1165_343_247_356_0.jpg}
\end{center}


第 57 题图

(1)求证: ${BC}\bot {C}_{1}F$ ；

(2)求点 $O$ 到平面 ${AC}{G}_{1}$ 的距离.

58.(2024$\cdot$徐汇区一模$\cdot$18)如图所示，某多面体的底面 ${ABCD}$ 为正方形， ${MA}//{PB},{MA}\bot {BC}$ ， ${AB} \bot  {PB},{MA} = 1,{AB} = {PB} = 2$.

\begin{center}
\includegraphics[width=0.2\textwidth]{bodyxb/src/173_1124_1193_301_329_0.jpg}
\end{center}


第 58 题图

(1)求四棱锥 $P - {ABCD}$ 的体积；

(2)求二面角 $B - {PM} - D$ 的平面角的正弦值.

59. (2024$\cdot$松江区二模$\cdot$18)如图所示，在四棱锥 $P - {ABCD}$ 中，底面 ${ABCD}$ 为菱形， ${PD}\bot$ 平面 ${ABCD}$ ， $E$ 为 ${PD}$ 的中点.

(1)设平面 ${ABE}$ 与直线 ${PC}$ 相交于点 $F$ ，求证: ${EF}//{CD}$ ；

(2)若 ${AB} = 2$ ， ${\angle {DAB}} = {60}^{ \circ  },{PD} = 4\sqrt{2}$ ，求直线 ${BE}$ 与平面 ${PAD}$ 所成角的大小.

\begin{center}
\includegraphics[width=0.2\textwidth]{bodyxb/src/174_1114_357_332_268_0.jpg}
\end{center}


第 59 题图

60. (2024$\cdot$宝山区二模$\cdot$18)如图所示，已知点 $P$ 在圆柱 $O{O}_{1}$ 的底面圆 $O$ 的圆周上， ${AB}$ 为圆 $O$ 的直径.

(1)求证: ${BP}\bot {A}_{1}P$ ；

(2)若 ${OA} = 2,\angle {BOP} = {60}^{ \circ  }$ ，圆柱的体积为 ${16}\sqrt{2}\pi$ ，求异面直线 ${AP}$ 与 ${A}_{1}B$ 所成角的大小.

\begin{center}
\includegraphics[width=0.2\textwidth]{bodyxb/src/174_1152_1265_291_340_0.jpg}
\end{center}


第 60 题图

61. (2024$\cdot$黄浦区一模$\cdot$18)如图所示，平面 ${ABCD} \bot$ 平面 ${ADEF}$ ，四边形 ADEF 是正方形， ${BC}//{AD}$ ， $\angle {BAD} = \angle {CDA} = {45}^{ \circ  },{CD} = 2,{AD} = 4\sqrt{2}$.

(1)证明: ${CD} \bot$ 平面 ${ABF}$ ；

(2)求二面角 $B - {EF} - A$ 的正切值.

\begin{center}
\includegraphics[width=0.3\textwidth]{bodyxb/src/175_984_362_398_198_0.jpg}
\end{center}


第 61 题图

62. (2025$\cdot$浦东新区二模$\cdot$18)如图所示，四边形 ${ABCD}$ 为长方形， ${PA}\bot$ 平面 ${ABCD},{AB} = {AP} = 2$ ， ${AD} = 3$.

(1)若 $E$ 、 $F$ 分别是 ${PB}$ 、 ${CD}$ 的中点，求证: ${EF}//$ 平面 ${PAD}$ ；

(2)边 ${BC}$ 上是否存在点 $G$ ，使得直线 ${PG}$ 与平面 ${PAD}$ 所成的角的大小为 ${30}^{ \circ  }$ ? 若存在，求 ${BG}$ 长；若不存在，说明理由.

\begin{center}
\includegraphics[width=0.2\textwidth]{bodyxb/src/175_1072_1208_323_338_0.jpg}
\end{center}


第 62 题图

63. (2024$\cdot$闵行区二模$\cdot$18)如图所示，已知 ${ABCD}\text{ 为 }$ ，等腰梯形， ${AD}//{BC},\angle {BAD} = {120}^{ \circ  },{PA}\bot$ 平面 ${ABCD},{AB} = {AD} = {AP} = 2$.

(1)求证: ${PC}\bot {AB}$ ；

(2)求二面角 $C - {BP} - A$ 的大小.

\begin{center}
\includegraphics[width=0.2\textwidth]{bodyxb/src/176_1104_334_321_243_0.jpg}
\end{center}


第 63 题图

64.(2024$\cdot$长宁区一模$\cdot$18)如图所示，在三棱锥 $A - {BCD}$ 中，平面 ${ABD}\bot$ 平面 ${BCD},{AB} = {AD}$ ， $O$ 为 ${BD}$ 的中点.

(1)求证: ${AO}\bot {CD}$ ；

(2)若 ${BD}\bot {DC},{BD} = {DC},{AO} = {BO}$ ，求异面直线 ${BC}$ 与 ${AD}$ 所成的角的大小.

\begin{center}
\includegraphics[width=0.2\textwidth]{bodyxb/src/176_1114_1187_312_259_0.jpg}
\end{center}


第 64 题图

65.(2024$\cdot$虹口区一模$\cdot$18)如图所示，在三棱柱 ${ABC} - {A}_{1}{B}_{1}{C}_{1}$ 中，侧面 $A{A}_{1}{C}_{1}C$ 为正方形， ${AB} =  {AC} = 4$ ,设 $M$ 是 $C{C}_{1}$ 的中点,满足 ${AM} \bot  {A}_{1}{B}_{1},N$ 是 ${BC}$ 的中点, $P$ 是线段 ${A}_{1}{B}_{1}$ 上的一点.

(1)证明: ${AM}\bot$ 平面 ${A}_{1}{PN}$ ;

(2)若 ${BC} = {4\sqrt{2}},{{A}_{1}P} = 1$ ，求直线 $A{B}_{1}$ 与平面 ${PMN}$ 所成角的大小.

\begin{center}
\includegraphics[width=0.2\textwidth]{bodyxb/src/177_1069_373_323_284_0.jpg}
\end{center}


第 65 题图

66. (2024$\cdot$上海春季卷$\cdot$18)如图所示， ${PA}$ 、 ${PB}$ 、 ${PC}$ 为圆锥的三条母线，且 ${AB} = {AC}$.

(1)证明: ${PA} \bot  {BC}$ ；

(2)若圆锥的侧面积为 $\sqrt{3}\pi ,{BC}$ 为底面直径， ${BC} = 2$ ，求二面角 $B - {PA} - C$ 的大小.

\begin{center}
\includegraphics[width=0.2\textwidth]{bodyxb/src/177_1117_1182_278_302_0.jpg}
\end{center}


第 66 题图

67.(2024$\cdot$长宁区二模$\cdot$18)如图所示，在长方体 ${ABCD} - {A}_{1}{B}_{1}{C}_{1}{D}_{1}$ 中， ${AB} = {AD} = 2,$ ， ${A{A}_{1}} = 1.$

(1)求二面角 ${D}_{1} - {AC} - D$ 的大小；

(2)若点 $P$ 在直线 ${A}_{1}{C}_{1}$ 上，求证:直线 ${BP}//$ 平面 ${D}_{1}{AC}$.

\begin{center}
\includegraphics[width=0.2\textwidth]{bodyxb/src/178_1125_303_335_279_0.jpg}
\end{center}


第 67 题图

68.(2024$\cdot$金山区二模$\cdot$18)如图所示，几何体是圆柱的一部分，它是由矩形 ${ABCD}$ (及其内部)以 ${AB}$ 边所在直线为旋转轴旋转 ${120}^{ \circ  }$ 得到的,点 $G$ 是 $\overset{\frown}{DF}$ 的中点,点 $P$ 在 $\overset{\frown}{CE}$ 上,异面直线 ${BP}$ 与 ${AD}$ 所成的角是 ${30}^{ \circ  }$.

\begin{center}
\includegraphics[width=0.2\textwidth]{bodyxb/src/178_1168_1104_291_256_0.jpg}
\end{center}


第 68 题图

(1)求证: ${AE}\bot {BP}$ ；

(2)若 ${AB} = 3,{AD} = 2$ ，求二面角 $E - {AG} - C$ 的大小.

69.(2024$\cdot$浦东新区二模$\cdot$18)在四棱锥 $P - {ABCD}$ 中，底面 ${ABCD}$ 为等腰梯形，平面 ${PAD} \bot$ 底面 ${ABCD}$ ，其中 ${AD}//{BC},{AD} = {2BC} = 4,{AB} = 3,{PA} = {PD} = 2\sqrt{3}$ ，点 $E$ 为 ${PD}$ 中点.

(1)证明: ${EC}//$ 平面 ${PAB}$ ；

(2)求二面角 $P - {AB} - D$ 的大小.

\begin{center}
\includegraphics[width=0.2\textwidth]{bodyxb/src/179_1085_356_303_287_0.jpg}
\end{center}


第 69 题图

70. (2024$\cdot$青浦区二模$\cdot$18)如图所示，三棱柱 ${ABC} - {A}_{1}{B}_{1}{C}_{1}$ 是所有棱长均为 2 的直三棱柱， $D$ 、 $E$ 分别为棱 ${AB}$ 和棱 $A{A}_{1}$ 的中点.

(1)求证: 平面 ${B}_{1}{CD} \bot$ 平面 ${AB}{B}_{1}{A}_{1}$ ；

(2)求二面角 ${B}_{1} - {CD} - E$ 的余弦值大小.

\begin{center}
\includegraphics[width=0.2\textwidth]{bodyxb/src/179_1134_1210_261_329_0.jpg}
\end{center}


第 70 题图

71. (2024$\cdot$徐汇区二模$\cdot$18)如图所示， $D$ 为圆锥的顶点， $O$ 是圆锥底面圆的圆心， ${AE}$ 为圆 $O$ 的直径,且 ${AE} = {AD} = 4$ ， $\bigtriangleup  {ABC}$ 是底面圆 $O$ 的内接正三角形， $P$ 为线段 ${DO}$ 上一点，且 ${DO} = \sqrt{6}{PO}$.

\begin{center}
\includegraphics[width=0.2\textwidth]{bodyxb/src/180_1121_300_293_314_0.jpg}
\end{center}


第 71 题图

(1)证明: ${PA} \bot$ 平面 ${PBC}$ ；

(2)求直线 ${PB}$ 与平面 ${PCE}$ 所成角的正弦值.

72. (2025$\cdot$闵行区二模$\cdot$18)如图所示，在四棱锥 $P - {ABCD}$ 中， ${PA}\bot$ 平面 ${PBC},{AB} = {2DC} = 4$ ， ${BC} = 2\sqrt{2},{AB} \bot  {BC},{DC}//{AB}$.

\begin{center}
\includegraphics[width=0.2\textwidth]{bodyxb/src/180_1118_1156_301_242_0.jpg}
\end{center}


第 72 题图

(1)证明:平面 ${ABCD} \bot$ 平面 ${PAB}$ ；

(2)若 $\angle {ABP} = \dfrac{\pi }{3}$ ，求点 $C$ 到平面 ${PAD}$ 的距离.



73. (2025$\cdot$静安区二模$\cdot$19)如图所示，在正三棱柱 ${ABC} - {A}_{1}{B}_{1}{C}_{1}$ 中， ${A{A}_{1}} = {2AC} = 4$ ，延长 ${CB}$ 至 $D$ ,使 ${CB} = {BD},E$ 是线段 ${BC}$ 的中点.

\begin{center}
\includegraphics[width=0.2\textwidth]{bodyxb/src/181_1232_295_217_337_0.jpg}
\end{center}


第 73 题图

(1)求证: \ding{172} 直线 ${C}_{1}B//$ 平面 $A{B}_{1}D$ ； \ding{173}  ${AE}\bot {D{B}_{1}}$ ；

(2)求二面角 ${B}_{1} - {AD} - C$ 的正弦值.

74.(2024$\cdot$宝山区一模$\cdot$19)如图所示，在直三棱柱 ${ABC} - {{A}_{1}{B}_{1}{C}_{1}}$ 中， ${AB} = {BC} = \sqrt{2},{AC} = {A{A}_{1}} = 2$ ， 且 $D$ 、 $E$ 分别是 ${AC}$ 、 ${A}_{1}{C}_{1}$ 的中点.

(1)证明: ${AC}\bot {BE}$ ；

(2)求三棱锥 $D - {ABE}$ 的体积；

(3)求直线 ${BD}$ 与平面 ${ABE}$ 所成角的大小.(结果用反三角函数值表示)

\begin{center}
\includegraphics[width=0.2\textwidth]{bodyxb/src/181_1161_1235_284_339_0.jpg}
\end{center}


第 74 题图

\subsection*{考点6 空间几何体的截面、交线及动态问题}

\subsubsection*{一、解答题}

1.(2025$\cdot$普陀区一模$\cdot$17)如图 1 所示的平行四边形 ${ABCD}$ 中， ${CA} = {CB} = 1,{CD} = \sqrt{2}$ ，现将 $\bigtriangleup  {DAC}$ 沿 ${AC}$ 折起,得到如图 2 所示的三棱锥 $P - {ABC}$ ，记棱 ${PC}$ 的中点为 $M$ ，且 ${PB} = \sqrt{3}$.

(1)求证: ${AM}\bot {BC}$ ；

(2)记棱 ${AB}$ 的中点为 $E$ ，在直线 ${CE}$ 上作出点 $N$ ，使得 ${PN}//$ 平面 ${MAB}$ ，请说明理由，并求出二面角 $P - {NB} - A$ 的大小.

\begin{center}
\includegraphics[width=0.4\textwidth]{bodyxb/src/182_936_684_510_284_0.jpg}
\end{center}


第 1 题图

2. (2024$\cdot$普陀区一模$\cdot$18)如图 1 所示等腰梯形 ${ABCD},{AB}//{DC},{AB} = {2DC} = 4,\angle {ABC} = \dfrac{\pi }{3}$ ， ${DE}\bot {AB}$ 于 $E$ 点，现将 $\bigtriangleup  {ADE}$ 沿直线 ${DE}$ 折起到 $\bigtriangleup  {PDE}$ 的位置，形成一个四棱锥 $P - {EBCD}$ ， 如图 2 所示.

(1)若 ${PC} = {2\sqrt{2}}$ ，求证: ${PE}\bot$ 平面 ${EBCD}$ ；

(2)若直线 ${PE}$ 与平面 ${EBCD}$ 所成的角为 $\dfrac{\pi }{3}$ ，求二面角 $P - {BC} - E$ 的大小.

\begin{center}
\includegraphics[width=0.4\textwidth]{bodyxb/src/182_955_1534_484_301_0.jpg}
\end{center}


第 2 题图

3.(2024$\cdot$奉贤区二模$\cdot$19)如图1是由两个三角形组成的图形,其中 $\angle {APC} = {90}^{ \circ  },\angle {PAC} = {30}^{ \circ  }$ ， ${AC} = {2AB},\angle {BCA} = {30}^{ \circ  }$. 将 $\bigtriangleup  {ABC}$ 沿 ${AC}$ 折起,使得平面 ${PAC}\bot$ 平面 ${ABC}$ ，如图 2 所示. 设 $O$ 是 ${AC}$ 的中点, $D$ 是 ${AP}$ 的中点.

(1)求直线 ${BD}$ 与平面 ${PAC}$ 所成角的大小；

(2)连接 ${PB}$ ，设平面 ${DBO}$ 与平面 ${PBC}$ 的交线为直线 $l$ ，判别 $l$ 与 ${PC}$ 的位置关系，并说明理由.

\begin{center}
\includegraphics[width=0.7\textwidth]{bodyxb/src/183_573_441_829_292_0.jpg}
\end{center}


第 3 题图

4.(2024$\cdot$嘉定区一模$\cdot$19)中国历史悠久，积累了许多房屋建筑的经验. 房梁为柱体，或取整根树干而制为圆柱形状，或作适当裁切而制为长方体形状，如图片所示. 材质确定的梁的承重能力取决于截面形状，现代工程科学常用抗弯截面系数 $W$ 来刻画梁的承重能力. 对于两个截面积相同的梁， 称 $W$ 较大的梁的截面形状更好. 三种不同截面形状的梁的抗弯截面系数公式，如下表所列:

\begin{center}
\includegraphics[width=0.3\textwidth]{bodyxb/src/183_167_1191_414_209_0.jpg}
\end{center}


图片引自梁思成《〈营造法式〉注释》卷五

\begin{center}

\begin{tabular}{|c|c|c|c|}
\hline
\phantom{X} & 圆形截面 & 正方形截面 & 矩形截面 \\
\cline{1-4}
条件 & $r$ 为圆半径 & $a$ 为正方形边长 & $h$ 为矩形的长, $b$ 为矩形的宽, $h > b$ \\
\cline{1-4}
抗弯截面系数 & ${W}_{1} = \dfrac{\pi }{4}{r}^{3}$ & ${W}_{2} = \dfrac{1}{6}{a}^{3}$ & ${W}_{3} = \dfrac{1}{6}b{h}^{2}$ \\
\cline{1-4}
\hline
\end{tabular}
\end{center}

(1)假设上表中的三种梁的截面面积相等，请问哪一种梁的截面形状最好？并具体说明；

(2)北宋学者李诫在《营造法式》中提出了矩形截面的梁的截面长宽之比应定为 $3 : 2$ 的观点. 考虑梁取材于圆柱形的树木,设矩形截面的外接圆的直径为常数 $D$ ，如图所示，请问 $h : b$ 为何值时，其抗弯截面系数取得最大值? 并据此分析李诫的观点是否合理.

\begin{center}
\includegraphics[width=0.2\textwidth]{bodyxb/src/183_1102_1530_303_166_0.jpg}
\end{center}


第 4 题图

\section*{第 11 章 导数及其应用}

\begin{center}

\begin{tabular}{|p{4cm}|p{12cm}|}
\hline
考点 & 课标要求与命题趋势 \\
\cline{1-2}
1. 导数的基本计算及其应用 & \ding{172}  借助物理背景与几何背景理解导数的意义，并能给出导数的严格数学定义 \ding{173}  通过导函数的概念，掌握二阶导数的概念，了解二阶导数的物理意义与几何意义；了解复合函数的求导公式 \ding{174}  掌握基本函数的导数求解，会导数的基本计算，会求切线方程，该内容是上海高考的命题热点,要熟练掌握 \\
\cline{1-2}
2. 求切线方程及其应用 & \ding{172}  掌握导数几何意义，能求曲线在某点或过某点的切线方程，并解决参数求解、切线应用等综合问题 \ding{173}  高考命题聚焦切线方程的多元考查: \ding{174} 基础题型侧重切点已知或未知的方程求解； \ding{175} 自高频结合参数讨论(如存在性、范围问题)；  \ding{176} 自创新融入公切线、切点弦及跨章节综合(如与函数零点、不等式结合)，强调逻辑推理与运算能力 \\
\cline{1-2}
3. 公切线问题 & \ding{172}  掌握导数几何意义，能解决两曲线的公切线方程求解及参数问题 \ding{173}  高考命题趋势聚焦: \ding{174} 基础题考查公切线方程的代数推导; \ding{175} 自高频结合参数范围、存在性探究(如斜率、截距条件)； \ding{176} 自创新融入多曲线(指数、对数、多项式)及跨章节综合(如与极值、零点结合)，强化逻辑推理与运算能力 \\
\cline{1-2}
4. 利用导数判断函数单调性及其应用 & \ding{172}  掌握导数与单调性的关系，能求单调区间(多项式不超过三次)，并应用于判断大小、求参等 \ding{173}  高考命题趋势显示，导数作为压轴题高频考查，题型涵盖含参讨论、不等式证明、零点问题等,综合度提升且贴近实际情境,强调数学建模与分类讨论能力 \\
\cline{1-2}
5. 求极值与最值及其应用 & \ding{172}  掌握极值的必要条件(导数为零)与充分条件(左右导数异号)，熟练运用导数求函数极值与闭区间最值(多项式≤3次) \ding{173}  高考命题聚焦: \ding{174}台基础题考查极值点判断及参数求解(如存在性、范围问题); \ding{175} 自高频结合参数讨论与分类思想 (如极值点个数、最值范围); \ding{176}- 创新融入实际优化问题及跨章节综合 (如与不等式、零点结合), 强化逻辑推理与运算能力 \\
\cline{1-2}
6. 利用导数研究函数的极值点及其应用 & \ding{172}  掌握极值点的定义、求法(通过一阶导数)及二阶导数判断极值类型(极大/极小)，能解决含参极值问题与简单优化问题(如成本最小化) \ding{173} 高考命题趋势显示，极值点问题作为压轴题高频考查，题型涵盖含参讨论、极值点偏移 (如比较极值点与中点关系)、实际情境优化(如利润最大化)，强调分类讨论、数形结合及数学建模能力，近年更注重与单调性、零点问题的综合考查 \\
\cline{1-2}
7. 导数与函数的基本性质结合问题 & \ding{172}  掌握导数与函数单调性、极值、最值的关联，强调利用导数解决实际优化问题 \ding{173}  高考命题趋势聚焦导数与函数性质的深度结合, 高频考查单调区间求解、极值存在性判断及参数范围探究, 常与不等式、方程等综合设问, 突出逻辑推理与数学建模能力 \\
\cline{1-2}
8. 利用导数研究函数的零点及其应用 & \ding{172}  掌握利用导数研究函数零点的方法，包括判断零点存在性、个数及参数范围，需结合单调性、极值与图像分析(如三次函数、复合函数) \ding{173} 高考命题趋势显示，零点问题作为压轴题高频考查，题型涵盖含参讨论、不等式证明、隐零点分析等，强调分类讨论、数形结合及数学建模能力，近年更注重与实际情境融合 \\
\cline{1-2}
\hline
\end{tabular}
\end{center}
\clearpage
\subsection*{考点1 分 导数的基本计算及其应用}

\subsubsection*{一、填空题}

1.(2024$\cdot$长宁区一模$\cdot$6)物体位移 $s$ 和时间 $t$ 满足函数关系 $s = {100t} - 5{t}^{2}\left( {0 < t < {20}}\right)$ ，则当 $t = 2$ 时， 物体的瞬时速度\blkda[2]{}.

2.(2024$\cdot$静安区二模$\cdot$6)已知物体的位移 $d$ (单位: ${m}$ )与时间 $t$ (单位:s)满足函数关系 $d = 2\sin t$ ，则

\begin{center}
\includegraphics[width=0.2\textwidth]{bodyxb/src/185_1287_575_158_267_0.jpg}
\end{center}


第 4 题图

在时间段 $t \in  \left( {2,6}\right)$ 内，物体的瞬时速度为 $1{m}/{s}$ 的时刻 $t =$ \blkda[2]{}(单位: ${s}$ ).

3.(2025$\cdot$静安区一模$\cdot$7)已知物体的位移 $d$ (单位: ${m}$ )与时间 $t$ (单位:s)满足函数关系 $d = 5\sin t - 2\cos t$ ，则该物体在 $t = \dfrac{\pi }{2}$ (s)时刻的瞬时速度为\blkda[2]{} $\left( {{m}/{s}}\right)$.

4.(2024$\cdot$青浦区二模$\cdot$11)如图所示，某酒杯上半部分的形状为倒立的圆锥，杯深 $8{{cm}}$ ，上口宽 $6{{cm}}$ ，若以 ${30}{{{cm}}}^{3}/{s}$ 的匀速往杯中注水，当水深为 $4{{cm}}$ 时，酒杯中水升高的瞬时变化率 $v =$ \blkda[2]{} ${{cm}}/{s}$.

5.(2025$\cdot$虹口区二模$\cdot$11)1798 年，人口学家马尔萨斯提出人口理论:从口数 $x\left( t\right)$ 是随着时间 $t$ 连续变化的函数; 人口增长率 $r$ 为常数,且单位时间内的人口增长量 ${x}^{\prime }\left( t\right)$ 与 $x\left( t\right)$ 成正比,进而建立了人口增长模型. 19 世纪中叶,生物学家发现由于人类生存条件的限制, $r$ 不是常数,因此改进了马尔萨斯的模型,并添加了新的假设:  \ding{172}  $r$ 是随着时间 $t$ 连续变化的函数,存在人口最大瞬时增长率 ${r}_{0}$ ,使 $r  = {r}_{0} - k \cdot  x\left( t\right) \left( {k > 0}\right)$ ,且 ${x}^{\prime }\left( t\right)$ 仅与 $r$ 和 $x\left( t\right)$ 有关;  \ding{173}  存在最大人口数 $N$ ,当人口数达到 $N$ 时, $r = 0$. 那么在这些假设下建立的人口增长模型 ${x}^{\prime }\left( t\right)  =$ \blkda[2]{}. [用含有 $x\left( t\right)$ 、 ${r}_{0}$ 、 $N$ 的式子表示]

\subsubsection*{二、选择题}

6.(2024$\cdot$青浦区一模$\cdot$14)若函数 $y = f\left( x\right)$ 在 $x = {x}_{0}$ 处的导数等于 $a$ ，则 $\mathop{\lim }\limits_{{{\Delta x} \rightarrow  0}}\dfrac{f\left( {{x}_{0} + {\Delta x}}\right)  - f\left( {{x}_{0} - {\Delta x}}\right) }{\Delta x}$ 的值为 \dotfill{(\kh{})}

\xx{}{}{}{}0 B. $a$ C. ${2a}$ D. ${3a}$

\begin{center}
\includegraphics[width=0.3\textwidth]{bodyxb/src/185_1090_1516_355_183_0.jpg}
\end{center}


第 7 题图

7.(2024$\cdot$松江区一模$\cdot$15)函数 $y = f\left( x\right)$ 的图像如图所示， $y = {f}^{\prime }\left( x\right)$ 为函数 $y = f\left( x\right)$ 的导函数，则不等式 $\dfrac{{f}^{\prime }\left( x\right) }{x} < 0$ 的解集为\dotfill{(\kh{})}

\xx{}{}{}{}$\left( {-3, - 1}\right)$ B. $\left( {0,1}\right)$

C. $\left( {-3, - 1}\right)  \cup  \left( {0,1}\right)$ D. $\left( {-\infty , - 3}\right)  \cup  \left( {1, + \infty }\right)$

8.(2025$\cdot$徐汇区二模$\cdot$16)已知函数 $y = f\left( x\right)$ 的定义域和值域都为 $\mathbf{R}$ ，且图像是一条连续不断的曲线,其导函数 $y = {f}^{\prime }\left( x\right)$ 的值如下表:

设 $D \subseteq  \mathbf{R}$ ,若集合 $\{ y \mid  y = f\left( x\right) ,x \in  D\}  =  \{ a,b,c\}$ ,其中 $a\text{ 、 }b\text{ 、 }c$ 为常数,则符合要求的集合 $D$ 的个数不可能是 \dotfill{(\kh{})}

\begin{center}

\begin{tabular}{|c|c|c|c|c|c|}
\hline
$x$ & $\left( {-\infty ,{x}_{1}}\right)$ & ${x}_{1}$ & $\left( {{x}_{1},{x}_{2}}\right)$ & ${x}_{2}$ & ( ${x}_{2}, + \infty$ ) \\
\cline{1-6}
${f}^{\prime }\left( x\right)$ & + & 0 & - & 0 & + \\
\cline{1-6}
\hline
\end{tabular}
\end{center}

C. 63 D. 343

\xx{}{}{}{}3 B. 27

9.(2024$\cdot$上海春季卷$\cdot$16)若函数 $y = f\left( x\right) \left( {x \in  \left( {n,n + 1}\right) ,n \in  \mathbf{N}}\right)$ 满足 $f\left( {x + 1}\right)  = {f}^{\prime }\left( x\right)$ ，则称函数 $y = f\left( x\right)$ 为延展函数. 已知延展函数 $y = g\left( x\right)$ 和函数 $y = h\left( x\right)$ ,满足当 $x \in  \left( {0,1}\right)$ 时, $g\left( x\right)  = {{e}}^{x}$ , $h\left( x\right)  = {x}^{10}$. 给定以下两个命题:

 \ding{172}  存在函数 $y = {kx} + b\left( {k\text{ 、 }b \in  \mathbf{R},k \neq  0}\right)$ 与 $y = g\left( x\right)$ 有无穷多个交点；

 \ding{173}  存在函数 $y = {kx} + b\left( {k\text{ 、 }b \in  \mathbf{R},k \neq  0}\right)$ 与 $y = h\left( x\right)$ 有无穷多个交点.

则正确的选项是 \dotfill{(\kh{})}

\xx{}{}{}{} \ding{172}  \ding{173} 都是真命题

B.  \ding{172}  \ding{173} 都是假命题

C.  \ding{172} 是真命题， \ding{173} 是假命题

D.  \ding{172} 是假命题， \ding{173} 是真命题

10.(2024$\cdot$虹口区二模$\cdot$16)已知定义在 $\mathbf{R}$ 上的函数 $f\left( x\right)$ 、 $g\left( x\right)$ 的导数满足 $\left| {{f}^{\prime }\left( x\right) }\right|  \leq  {g}^{\prime }\left( x\right)$ ，给出两个命题:

 \ding{172}  对任意 ${x}_{1}\text{ 、 }{x}_{2} \in  \mathbf{R}$ ，都有 $\left| {f\left( {x}_{1}\right)  - f\left( {x}_{2}\right) }\right|  \leq  \left| {g\left( {x}_{1}\right)  - g\left( {x}_{2}\right) }\right|$ ；

 \ding{173}  若 $g\left( x\right)$ 的值域为 $\left\lbrack  {m,M}\right\rbrack  ,f\left( {-1}\right)  = m,f\left( 1\right)  = M$ ,则对任意 $x \in  \mathbf{R}$ 都有 $f\left( x\right)  = g\left( x\right)$. 则下列判断正确的是 \dotfill{(\kh{})}

\xx{}{}{}{} \ding{172}  \ding{173} 都是假命题

B.  \ding{172}  \ding{173} 都是真命题

C.  \ding{172} 是假命题， \ding{173} 是真命题

D.  \ding{172} 是真命题， \ding{173} 是假命题

11. (2025$\cdot$奉贤区二模$\cdot$15)函数 $y = f\left( x\right)$ 的导函数为 $y = g\left( x\right)$ ，若存在实数 ${x}_{0}$ ，使得 $g\left( {x}_{0}\right) f\left( {-{x}_{0}}\right)  = 1$ 成立,则称函数 $y = f\left( x\right)$ 具有 ${x}_{0}$ 性质,下列函数 $y = f\left( x\right)$ 具有 ${x}_{0}$ 性质的函数是 \dotfill{(\kh{})}

\xx{}{}{}{}$y = {{e}}^{-x}$ B. $y = \sin x$

C. $y = {{e}}^{x} + {{e}}^{-x}$ D. $y = \dfrac{1}{2}\ln \left( {{x}^{2} + 1}\right)$

\subsection*{考点2 > 求切线方程及其应用}

\subsubsection*{一、填空题}

1.(2025$\cdot$长宁区一模$\cdot$3)曲线 $y = \ln x$ 在点 $\left( {1,0}\right)$ 处的切线方程为\blkda[2]{}.

2.(2025$\cdot$闵行区二模$\cdot$6)若直线 $l$ 是曲线 $y = \dfrac{2}{x - 1}$ 在 $x = 3$ 处的切线，则 $l$ 的斜率为\blkda[2]{}.

3.(2024$\cdot$金山区二模$\cdot$7)设 $f\left( x\right)  = {x}^{3} + a{x}^{2} + x\left( {a \in  \mathbf{R}}\right)$ ，若 $y = f\left( x\right)$ 为奇函数，则曲线 $y = f\left( x\right)$ 在点 $\left( {0,0}\right)$ 处的切线方程为\blkda[2]{}.

4.(2024$\cdot$闵行区二模$\cdot$8)函数 $y = \dfrac{2}{x} - x$ 在 $x = 1$ 处的切线方程为\blkda[2]{}.

5.(2024$\cdot$嘉定区二模$\cdot$8)已知曲线 $y = \dfrac{1}{3}{x}^{3}$ 上有一点 $P\left( {2,\dfrac{8}{3}}\right)$ ，则过点 $P$ 的切线的斜率为\blkda[2]{}.

\subsubsection*{二、解答题}

6.(2025$\cdot$嘉定区二模$\cdot$18)已知函数 $y = f\left( x\right)$ ，其中 $f\left( x\right)  = a \cdot  {{e}}^{x} + b \cdot  {{e}}^{-x},a$ 、 $b$ 为实常数且 ${ab} \neq  0$.

(1)若 $y = f\left( x\right)$ 为偶函数，且其最小值为 4，求实数 $a$ 与 $b$ 的值；

(2)若 $a = 1,g\left( x\right)  = {{e}}^{x} - x$ ，对任意实数 $x$ 均满足 $f\left( x\right)  \geq  g\left( x\right)$ ，求实数 $b$ 的取值范围.

7.(2024$\cdot$黄浦区一模$\cdot$20)设 $a$ 为实数， ${\Gamma }_{1}$ 是以点 $O\left( {0,0}\right)$ 为顶点、以点 $F\left( {0,\dfrac{1}{4}}\right)$ 为焦点的抛物线， ${\Gamma }_{2}$ 是以点 $A\left( {0,a}\right)$ 为圆心、半径为 1 的圆位于 $y$ 轴右侧且在直线 $y = a$ 下方的部分.

(1)求 ${\Gamma }_{1}$ 与 ${\Gamma }_{2}$ 的方程；

(2)若直线 $y = x + 2$ 被 ${\Gamma }_{1}$ 所截得的线段的中点在 ${\Gamma }_{2}$ 上,求 $a$ 的值；

(3)是否存在 $a$ ，满足: ${\Gamma }_{2}$ 在 ${\Gamma }_{1}$ 的上方，且 ${\Gamma }_{2}$ 有两条不同的切线被 ${\Gamma }_{1}$ 所截得的线段长相等? 若存在，求出 $a$ 的取值范围；若不存在, 请说明理由.

\begin{center}
\includegraphics[width=0.2\textwidth]{bodyxb/src/187_1085_1239_338_414_0.jpg}
\end{center}


第 7 题图

8.(2024$\cdot$崇明区二模$\cdot$21)已知 $f\left( x\right)  = 2\sqrt{x} - a\ln x - {ax} - 1$.

(1)若 $a =  - 1$ ，求曲线 $y = f\left( x\right)$ 在点 $P\left( {1,2}\right)$ 处的切线方程；

(2)若函数 $y = f\left( x\right)$ 存在两个不同的极值点 ${x}_{1}\text{ 、 }{x}_{2}$ ，求证: $f\left( {x}_{1}\right)  + f\left( {x}_{2}\right)  > 0$ ；

(3)若 $a = 1,g\left( x\right)  = f\left( x\right)  + x$ ，数列 $\left\{  {a}_{n}\right\}$ 满足 ${a}_{1} \in  \left( {0,1}\right)$ ， ${a}_{n + 1} = g\left( {a}_{n}\right)$. 求证:当 $n \geq  2$ 时， ${a}_{n} +  {a}_{n + 2} > 2{a}_{n + 1}$.

9.(2024$\cdot$奉贤区二模$\cdot$21)已知定义域为 $\mathbf{R}$ 的函数 $y = f\left( x\right)$ ，其图像是连续的曲线，且存在定义域也为 $\mathbf{R}$ 的导函数 $y = {f}^{\prime }\left( x\right)$.

(1)求函数 $f\left( x\right)  = {{e}}^{x} + {{e}}^{-x}$ 在点 $\left( {0,f\left( 0\right) }\right)$ 的切线方程；

(2)已知 $f\left( x\right)  = a\cos x + b\sin x$ ，当 $a$ 与 $b$ 满足什么条件时，存在非零实数 $k$ ，对任意的实数 $x$ 使得 $f\left( {-x}\right)  =  - k{f}^{\prime }\left( x\right)$ 恒成立?

(3)若函数 $y = f\left( x\right)$ 是奇函数，且满足 $f\left( x\right)  + f\left( {2 - x}\right)  = 3$. 试判断 ${f}^{\prime }\left( {x + 2}\right)  = {f}^{\prime }\left( {2 - x}\right)$ 对任意的实数 $x$ 是否恒成立，请说明理由.

10.(2025$\cdot$崇明区二模$\cdot$21)已知函数 $y = f\left( x\right)$ ， $P$ 为坐标平面上一点. 若函数 $y = f\left( x\right)$ 的图像上存在与 $P$ 不同的一点 $Q$ ,使得直线 ${PQ}$ 是函数 $y = f\left( x\right)$ 在点 $Q$ 处的切线,则称点 $P$ 具有性质 ${M}_{f}$. (1)若 $f\left( x\right)  = {x}^{2}$ ，判断点 $P\left( {1,0}\right)$ 是否具有性质 ${M}_{f}$ ，并说明理由；

(2)若 $f\left( x\right)  = 2{x}^{3} - 4{x}^{2} + {2x}$ ，证明:线段 $x = \dfrac{1}{2}\left( {-1 \leq  y \leq  1}\right)$ 上的所有点均具有性质 ${M}_{f}$ ；

(3)若 $f\left( x\right)  = {{e}}^{x}$ ，证明: “点 $P\left( {x,y}\right)$ 具有性质 ${M}_{f}$ ”的充要条件是“$y < {{e}}^{x}$ ”.

11. (2024$\cdot$宝山区二模$\cdot$21)函数 $y = g\left( x\right)$ 的表达式为 $g\left( x\right)  = \sin \left( {\omega x}\right) \left( {\omega  > 0}\right)$.

(1)若 $\omega  = 1$ ，直线 $l$ 与曲线 $y = g\left( x\right)$ 相切于点 $\left( {\dfrac{\pi }{2},1}\right)$ ，求直线 $l$ 的方程；

(2)函数 $y = g\left( x\right)$ 的最小正周期是 ${2\pi }$ ,令 $h\left( x\right)  = x \cdot  g\left( x\right)  - \ln x$ ,将函数 $y = h\left( x\right)$ 的零点由小到大依次记为 ${x}_{1},{x}_{2},\cdots ,{x}_{n},\cdots \left( {n \geq  1,n \in  \mathbf{N}}\right)$ ,证明: 数列 $\left\{  {\sin {x}_{n}}\right\}$ 是严格减数列;

(3)已知定义在 $\mathbf{R}$ 上的奇函数 $y = f\left( x\right)$ 满足 $f\left( {x + {2a}}\right)  =  - f\left( x\right) \left( {a > 0}\right)$ ，对任意 $x \in  \left\lbrack  {0,{2a}}\right\rbrack$ ，当 $x \neq  a$ 时,都有 $f\left( x\right)  < f\left( a\right)$ 且 $f\left( a\right)  = 1$. 记 $F\left( x\right)  = f\left( x\right)  + g\left( x\right) ,G\left( x\right)  = f\left( x\right)  +  g\left( {x + \dfrac{1}{2}}\right)$. 当 $\omega  = \pi$ 时,是否存在 ${x}_{1}\text{ 、 }{x}_{2} \in  \mathbf{R}$ ,使得 $F\left( {x}_{1}\right)  = G\left( {x}_{2}\right)  + 4$ 成立? 若存在,求出符合题意的 ${x}_{1}\text{ 、 }{x}_{2}$ ; 若不存在,请说明理由.

12.(2025$\cdot$崇明区一模$\cdot$21)定义:若曲线 ${C}_{1}$ 和曲线 ${C}_{2}$ 有公共点 $P$ ，且曲线 ${C}_{1}$ 在点 $P$ 处的切线与曲线 ${C}_{2}$ 在点 $P$ 处的切线重合,则称 ${C}_{1}$ 与 ${C}_{2}$ 在点 $P$ 处 “一线切”.

(1)已知圆 ${\left( x - a\right) }^{2} + {y}^{2} = {r}^{2}\left( {r > 0}\right)$ 与曲线 $y = {x}^{2}$ 在点 $\left( {1,1}\right)$ 处“一线切”,求实数 $a$ 的值；

(2)设 $f\left( x\right)  = {x}^{2} + {2x} + a,g\left( x\right)  = \ln \left( {x + 1}\right)$ ,若曲线 $y = f\left( x\right)$ 与曲线 $y = g\left( x\right)$ 在点 $P$ 处“一线切”,求实数 $a$ 的值;

(3)定义在 $\mathbf{R}$ 上的函数 $y = f\left( x\right)$ 的图像为连续曲线，函数 $y = f\left( x\right)$ 的导函数为 $y = {f}^{\prime }\left( x\right)$ ，对任意的 $x \in  \mathbf{R}$ ,都有 $\left\{  \begin{array}{l} \left| {{f}^{\prime }\left( x\right) }\right|  \geq  \left| {f\left( x\right) }\right| \\  \left| {f\left( x\right) }\right|  < \sqrt{2} \end{array}\right.$ 成立. 是否存在点 $P$ 使得曲线 $y = f\left( x\right) \sin x$ 和曲线 $y = 1$ 在点 $P$ 处“一线切”? 若存在,请求出点 $P$ 的坐标; 若不存在,请说明理由.

13.(2025$\cdot$奉贤区一模$\cdot$21)若函数 $y = f\left( x\right)$ 的图像上存在 $k$ 个不同点 ${P}_{1}$ 、 ${P}_{2}$ 、 $\cdots$ 、 ${P}_{k}\left( {k \geq  2,k \in  \mathbf{N}}\right)$ 处的切线重合,则称该切线为函数 $y = f\left( x\right)$ 的一条 $k$ 点切线,该系数具有 $k$ 点切线性质.

(1)判断函数 $y = {x}^{2} - 2\left| x\right| ,x \in  \mathbf{R}$ 的奇偶性并写出它的一条 2 点切线方程(无需理由);

(2)设 $f\left( x\right)  = {{e}}^{x} - \ln x$ ，判断函数 $y = f\left( x\right)$ 是否具有 $k$ 点切线性质，并说明理由；

(3)设 $g\left( x\right)  = \cos x + {2x}$ ，证明:对任意的 $m \geq  3$ ， $m \in  \mathbf{N}$ ，函数 $y = g\left( x\right)$ 具有 $m$ 点切线性质，并求出所有相应的切线方程.

14.(2024$\cdot$宝山区一模$\cdot$21)已知函数 $f\left( x\right)  = {{e}}^{x} - x,g\left( x\right)  = {{e}}^{-x} + x$ ，其中 ${e}$ 为自然对数的底数.

(1)求函数 $y = f\left( x\right)$ 的图像在点 $\left( {1,f\left( 1\right) }\right)$ 处的切线方程；

(2)设函数 $F\left( x\right)  = {af}\left( x\right)  - g\left( x\right)$ ，

 \ding{172}  若 $a = {e}$ ，求函数 $y = F\left( x\right)$ 的单调区间，并写出函数 $y = F\left( x\right)  - m$ 有三个零点时实数 $m$ 的取值范围;

 \ding{173}  当 $0 < a < 1$ 时， ${x}_{1}$ 、 ${x}_{2}$ 分别为函数 $y = F\left( x\right)$ 的极大值点和极小值点，且不等式 $F\left( {x}_{1}\right)  +  {tF}\left( {x}_{2}\right)  > 0$ 对任意 $a \in  \left( {0,1}\right)$ 恒成立,求实数 $t$ 的取值范围.

15.(2024$\cdot$黄浦区一模$\cdot$21)设函数 $f\left( x\right)$ 与 $g\left( x\right)$ 的定义域均为 $D$ ，若存在 ${x}_{0} \in  D$ ，满足 $f\left( {x}_{0}\right)  =  g\left( {x}_{0}\right)$ 且 ${f}^{\prime }\left( {x}_{0}\right)  = {g}^{\prime }\left( {x}_{0}\right)$ ,则称函数 $f\left( x\right)$ 与 $g\left( x\right)$ “局部趋同”.

(1)判断函数 ${f}_{1}\left( x\right)  = {{5x} + 1}$ 与 ${f}_{2}\left( x\right)  = {x}^{3} + {2x}$ 是否“局部趋同”，并说明理由;

(2)已知函数 ${g}_{1}\left( x\right)  =  - {x}^{2} + {ax}\left( {x > 0}\right) ,{g}_{2}\left( x\right)  = b{{e}}^{x}\left( {x > 0}\right)$. 求证:对任意的正数 $a$ ，都存在正数 $b$ ,使得函数 ${g}_{1}\left( x\right)$ 与 ${g}_{2}\left( x\right)$ “局部趋同”;

(3)对于给定的实数 $m$ ，若存在实数 $n$ ，使得函数 ${h}_{1}\left( x\right)  = {mx} + \dfrac{n}{x}\left( {x > 0}\right)$ 与 ${h}_{2}\left( x\right)  = \ln x$ “局部趋同”,求实数 $m$ 的取值范围.

\subsection*{考点3 分 公切线问题}

\subsubsection*{一、选择题}

1.(2024$\cdot$青浦区二模$\cdot$16)如图所示，已知直线 $y = {kx} + m$ 与函数 $y = f\left( x\right)$ ， $x \in  \left( {0, + \infty }\right)$ 的图像相切于两点,则函数 $y = f\left( x\right)  - {kx}$ 有 \dotfill{(\kh{})}

\begin{center}
\includegraphics[width=0.3\textwidth]{bodyxb/src/192_1000_534_435_196_0.jpg}
\end{center}


第 1 题图

\xx{}{}{}{}2 个极大值点，1 个极小值点

B. 3 个极大值点，2 个极小值点

C. 2 个极大值点，无极小值点

D. 3 个极大值点，无极小值点

\subsubsection*{二、解答题}

2.(2024$\cdot$黄浦区二模$\cdot$21)若函数 $y = f\left( x\right)$ 的图像上的两个不同点处的切线互相重合，则称该切线为函数 $y = f\left( x\right)$ 的图像的“自公切线”，称这两点为函数 $y = f\left( x\right)$ 的图像的一对“同切点”.

(1)分别判断函数 ${f}_{1}\left( x\right)  = \sin x$ 与 ${f}_{2}\left( x\right)  = \ln x$ 的图像是否存在“自公切线”，并说明理由；

(2)若 $a \in  \mathbf{R}$ ，求证:函数 $g\left( x\right)  = \tan x - x + a\left( {x \in  \left( {-\dfrac{\pi }{2},\dfrac{\pi }{2}}\right) }\right)$ 有唯一零点且该函数的图像不存在 “自公切线”;

(3)设 $n \in  {\mathbf{N}}^{ * }$ ， $h\left( x\right)  = \tan x - x + {n\pi }\left( {x \in  \left( {-\dfrac{\pi }{2},\dfrac{\pi }{2}}\right) }\right)$ 的零点为 ${x}_{n}$ ， $t \in  \left( {-\dfrac{\pi }{2},\dfrac{\pi }{2}}\right)$ ，求证: “存在 $s \in  \left( {{2\pi }, + \infty }\right)$ ,使得点 $\left( {s,\sin s}\right)$ 与 $\left( {t,\sin t}\right)$ 是函数 $y = \sin x$ 的图像的一对 “同切点” ”的充要条件是“$t$ 是数列 $\left\{  {x}_{n}\right\}$ 中的项”.

3. (2024$\cdot$闵行区一模$\cdot$21)已知 $a \in  \mathbf{R}$ ， $f\left( x\right)  = \left( {a - 2}\right) {x}^{3} - {x}^{2} + {5x} + \left( {1 - a}\right) \ln x$.

(1)若 1 为函数 $y = f\left( x\right)$ 的驻点，求实数 $a$ 的值；

(2)若 $a = 0$ ，试问曲线 $y = f\left( x\right)$ 是否存在切线与直线 $x - y - 1 = 0$ 互相垂直？请说明理由；

(3)若 $a = 2$ ，是否存在等差数列 ${x}_{1},{x}_{2},{x}_{3}\left( {0 < {x}_{1} < {x}_{2} < {x}_{3}}\right)$ ，使得曲线 $y = f\left( x\right)$ 在点 $\left( {x}_{2}\right.$ ， $\left. {f\left( {x}_{2}\right) }\right)$ 处的切线与过两点 $\left( {{x}_{1},f\left( {x}_{1}\right) }\right) \text{ 、 }\left( {{x}_{3},f\left( {x}_{3}\right) }\right)$ 的直线互相平行? 若存在,求出所有满足条件的等差数列; 若不存在, 请说明理由.

4.(2024$\cdot$松江区一模$\cdot$21)已知函数 $y = f\left( x\right)$ ，记 $f\left( x\right)  = x + \sin x$ ， $x \in  D$.

(1)若 $D = \left\lbrack  {0,{2\pi }}\right\rbrack$ ，判断函数的单调性；

(2)若 $D = \left( {0,\dfrac{\pi }{2}}\right\rbrack$ ，不等式 $f\left( x\right)  > {kx}$ 对任意 $x \in  D$ 恒成立，求实数 $k$ 的取值范围；

(3)若 $D = \mathbf{R}$ ，则曲线 $y = f\left( x\right)$ 上是否存在三个不同的点 $A$ 、 $B$ 、 $C$ ，使得曲线 $y = f\left( x\right)$ 在 $A$ 、 $B$ 、 $C$ 三点处的切线互相重合? 若存在,求出所有符合要求的切线的方程; 若不存在,请说明理由.

\subsection*{考点4 为 利用导数判断函数单调性及其应用}

\subsubsection*{一、填空题}

\begin{center}
\includegraphics[width=0.2\textwidth]{bodyxb/src/194_1199_421_264_264_0.jpg}
\end{center}


第 1 题图

1.(2025$\cdot$上海春季卷$\cdot$11)如图所示，正方形 ${ABCD}$ 是一块边长为 4 的工程用料，阴影部分所示是被腐蚀的区域，其余部分完好，曲线 ${MN}$ 为以 ${AD}$ 为对称轴的抛物线的一部分， ${DM} = {DN} = 3$. 工人师傅现要从完好的部分中截取一块矩形原料 BQPR，当其面积取得最大值时， ${AQ}$ 的长为\blkda[2]{}.

2.(2024$\cdot$长宁区一模$\cdot$11)若函数 $f\left( x\right)  = \sin x + a\cos x$ 在 $\left( {\dfrac{2\pi }{3},\dfrac{7\pi }{6}}\right)$ 上是严格单调函数，则实数 $a$ 的取值范围为\blkda[2]{}.

3. (2024$\cdot$普陀区一模$\cdot$11)设函数 $f\left( x\right)  = a{{e}}^{x} - 2{x}^{2}$ ，若对任意 ${x}_{0} \in  \left( {0,1}\right)$ ，皆有 $\mathop{\lim }\limits_{{x \rightarrow  {x}_{0}}}\dfrac{f\left( x\right)  - f\left( {x}_{0}\right)  - x + {x}_{0}}{x - {x}_{0}}$ >0 成立，则实数 $a$ 的取值范围是\blkda[2]{}.

4.(2024$\cdot$徐汇区二模$\cdot$11)如图所示，两条足够长且互相垂直的轨道 ${l}_{1}\text{ 、 }{l}_{2}$ 相交于点 $O$ ,一根长度为 8 的直杆 ${AB}$ 的两端点 $A\text{ 、 }B$ 分别在 ${l}_{1}\text{ 、 }{l}_{2}$ 上滑动 $\left( {A\text{ 、 }B}\right.$ 两点不与点 $O$ 重合,轨道与直杆的宽度等因素均可忽略不计),直杆上的点 $P$ 满足 ${OP} \bot  {AB}$ ,则 $\bigtriangleup {OAP}$ 面积的取值范围是\blkda[2]{}.

\begin{center}
\includegraphics[width=0.2\textwidth]{bodyxb/src/194_1130_896_329_256_0.jpg}
\end{center}


第 4 题图

5.(2024$\cdot$静安区一模$\cdot$12)记 $f\left( x\right)  = \ln x + {x}^{2} - {2kx} + {k}^{2}$ ，若存在实数 $a\text{ 、 }b$ ,满足 $\dfrac{1}{2} \leq  a < b \leq  2$ ,使得函数 $y = f\left( x\right)$ 在区间 $\left\lbrack  {a,b}\right\rbrack$ 上是严格增函数,则实数 $k$ 的取值范围是\blkda[2]{}.

6. (2024$\cdot$普陀区二模$\cdot$12)已知 $a \in  \mathbf{R}$ ，若关于 $x$ 的不等式 $a\left( {x - 2}\right) {{e}}^{-x} - x > 0$ 的解集中有且仅有一个负整数，则 $a$ 的取值范围是\blkda[2]{}.

7.(2025$\cdot$闵行区二模$\cdot$12)定义 $D = \left\lbrack  {a,b}\right\rbrack$ 的区间长度为 $b - a$. 若 $m < 0$ 且关于 $x$ 的不等式 $\mid  {\left( x - 1\right) }^{3} +  m\left( {x - 1}\right)  \leq  {16}$ 的解集的区间长度之和为 $T$ ，则当 $T$ 取最大值时，实数 $m$ 的值为\blkda[2]{}.

\subsubsection*{二、选择题}

8.(2025$\cdot$徐汇区一模$\cdot$15)已知函数 $y = f\left( x\right)$ 与它的导函数 $y = {f}^{\prime }\left( x\right)$ 的定义域均为\blkda[2]{}，若函数 $y = f\left( x\right)$ 是偶函数且 $y = {f}^{\prime }\left( x\right)$ 在 $\left( {-\infty ,0}\right)$ 上是严格增函数，则下列各表中，可能成为 $y = f\left( x\right)$ 取值的是\dotfill{(\kh{})}

\begin{center}

\begin{tabular}{|p{4cm}|p{12cm}|}
\hline
$x$ & $f\left( x\right)$ \\
\cline{1-2}
1 & 2.8188 \\
\cline{1-2}
2 & 1.0000 \\
\cline{1-2}
3 & 0.3644 \\
\cline{1-2}
4 & 0.2468 \\
\cline{1-2}
\hline
\end{tabular}
\end{center}

A.

\begin{center}

\begin{tabular}{|p{4cm}|p{12cm}|}
\hline
$x$ & $f\left( x\right)$ \\
\cline{1-2}
1 & 0.7580 \\
\cline{1-2}
2 & 1.0000 \\
\cline{1-2}
3 & 1.3188 \\
\cline{1-2}
4 & 1.7979 \\
\cline{1-2}
\hline
\end{tabular}
\end{center}

B.

\begin{center}

\begin{tabular}{|p{4cm}|p{12cm}|}
\hline
$x$ & $f\left( x\right)$ \\
\cline{1-2}
1 & 2.4132 \\
\cline{1-2}
2 & 1.0000 \\
\cline{1-2}
3 & 1.5885 \\
\cline{1-2}
4 & 4.1116 \\
\cline{1-2}
\hline
\end{tabular}
\end{center}

C.

\begin{center}

\begin{tabular}{|p{4cm}|p{12cm}|}
\hline
$x$ & $f\left( x\right)$ \\
\cline{1-2}
1 & 0.8664 \\
\cline{1-2}
2 & 1.0000 \\
\cline{1-2}
3 & 1.1188 \\
\cline{1-2}
4 & 1.2240 \\
\cline{1-2}
\hline
\end{tabular}
\end{center}

D.

9.(2025$\cdot$黄浦区一模$\cdot$16)设函数 $y = f\left( x\right)$ 在区间 $I$ 上有导函数 $y = {f}^{\prime }\left( x\right)$ ，且 ${f}^{\prime }\left( x\right)  < 0$ 在区间 $I$ 上恒成立,对任意的 $x \in  I$ ,有 $f\left( x\right)  \in  I$. 对于各项均不相同的数列 $\left\{  {a}_{n}\right\}  ,{a}_{1} \in  I,{a}_{n + 1} = f\left( {a}_{n}\right)$ ,下列结论正确的是 \dotfill{(\kh{})}

\xx{}{}{}{}数列 $\left\{  {a}_{{2n} - 1}\right\}$ 与 $\left\{  {a}_{2n}\right\}$ 均是严格增数列

B. 数列 $\left\{  {a}_{{2n} - 1}\right\}$ 与 $\left\{  {a}_{2n}\right\}$ 均是严格减数列

C. 数列 $\left\{  {a}_{{2n} - 1}\right\}$ 与 $\left\{  {a}_{2n}\right\}$ 中的一个是严格增数列,另一个是严格减数列

D. 数列 $\left\{  {a}_{{2n} - 1}\right\}$ 与 $\left\{  {a}_{2n}\right\}$ 均既不是严格增数列也不是严格减数列

10. (2024$\cdot$闵行区一模$\cdot$16)已知函数 $y = f\left( x\right)$ 的导函数为 $y = {f}^{\prime }\left( x\right)$ ， $x \in  \mathbf{R}$ ，且 $y = {f}^{\prime }\left( x\right)$ 在 $\mathbf{R}$ 上为严格增函数，关于下列两个命题的判断，说法正确的是 \dotfill{(\kh{})}

 \ding{172} “${x}_{1} > {x}_{2}$”是“$f\left( {{x}_{1} + 1}\right)  + f\left( {x}_{2}\right)  > f\left( {x}_{1}\right)  + f\left( {{x}_{2} + 1}\right)$ ”的充要条件;

 \ding{173}  “对任意 $x < 0$ 都有 $f\left( x\right)  < f\left( 0\right)$ ”是“$y = f\left( x\right)$ 在 $\mathbf{R}$ 上为严格增函数”的充要条件.

\xx{}{}{}{} \ding{172} 真命题， \ding{173} 假命题 B.  \ding{172} 假命题， \ding{173} 真命题

C.  \ding{172}  \ding{173} 均为真命题 D.  \ding{172}  \ding{173} 均为假命题

\subsubsection*{三、解答题}

11. (2025$\cdot$静安区一模$\cdot$17) 设 $f\left( x\right)  = x + \dfrac{4}{x},x \in  \left( {-\infty ,0}\right)  \cup  \left( {0, + \infty }\right)$.

(1)求函数 $y = f\left( x\right)$ 单调区间；

(2)求不等式 $f\left( x\right)  < {2x}$ 的解集.

12.(2025$\cdot$青浦区一模$\cdot$21)已知函数 $y = f\left( x\right)$ ，其中 $f\left( x\right)  = {{e}}^{x - 1} - {2\ln }x + x$.

(1)求函数 $y = f\left( x\right)$ 的单调区间；

(2)设 $g\left( x\right)  = f\left( x\right)  + 2\ln x$ ，问:函数 $y = g\left( x\right)$ 的图像上是否存在三点 $A$ 、 $B$ 、 $C$ ，使得它们的横坐标成等差数列，且直线 ${AC}$ 的斜率等于 $y = g\left( x\right)$ 在点 $B$ 处的切线的斜率？若存在，求出所有满足条件的点 $B$ 的坐标；若不存在，请说明理由；

(3)证明:函数 $y = f\left( x\right)$ 图像上任意一点都不落在函数 $y = {\left( x - 2\right) }^{3} - 3\left( {x - 2}\right)$ 图像的下方.

13.(2025$\cdot$宝山区二模$\cdot$21)定义在 $D$ 上的可导函数 $y = f\left( x\right)$ ，集合 ${A}_{\left( k,m\right) } = \left\{  {f\left( x\right)  \mid  F\left( {x}_{i}\right)  = k}\right.$ ， ${x}_{i} \in  D,i = 1\text{ 、 }2\text{ 、 }\cdots \text{ 、 }m,m$ 为正整数 $\}$ ,其中 $F\left( x\right)  = f\left( x\right)  + {f}^{\prime }\left( x\right)$ 称为 $f\left( x\right)$ 的自和函数, ${x}_{i}$ 称为 $y = f\left( x\right)$ 的固着点. 已知 $f\left( x\right)  = a{e}^{x} + {bx} + c\sin x\left( {a\text{ 、 }b\text{ 、 }c \in  \mathbf{R}}\right)$.

(1)若 $a = c = 0,b = 2,D = \mathbf{R},f\left( x\right)  \in  {A}_{\left( 1,m\right) }$ ，求 $m$ 的值及 $y = f\left( x\right)$ 的固着点；

(2)若 $a = 0,b = 1,c = 1,D = \left\lbrack  {s,t}\right\rbrack  \left( {s > 0}\right) ,F\left( x\right)$ 是 $f\left( x\right)$ 的自和函数，且 $F\left( x\right)$ 在 $D$ 上是严格增函数，求 $t - s$ 的最大值；

(3)若 $b =  - 1,c = 0,D = \left( {0, + \infty }\right) ,f\left( x\right)  \in  {A}_{\left( 0,1\right) }$ ，且 $t$ 是 $y = f\left( x\right)$ 的固着点，求 $a$ 的取值范围， 并证明: $\dfrac{1}{2a} < {{e}}^{t} < \dfrac{1}{{a}^{2}}$.

14.(2025$\cdot$黄浦区二模$\cdot$21)设 $D$ 是 $\mathbf{R}$ 的一个非空子集，函数 $y = f\left( x\right)$ 的定义域为 $D$ ，若 $y = f\left( x\right)$ 在 $D$ 上不是单调函数，且存在常数 $b$ ，使得 $f\left( x\right)  \geq  b$ 对任意的 $x \in  D$ 成立，则称函数 $y = f\left( x\right)$ 具有性质 $H$ ,称 $b$ 为该函数的一个下界.

(1)设 $f\left( x\right)  = x + \dfrac{1}{x},D = \left( {-\infty ,0}\right)$ ，判断函数 $y = f\left( x\right) ,x \in  D$ 是否具有性质 $H$ ；

(2)设 $m$ 为常数， $f\left( x\right)  = \dfrac{1}{3}{x}^{3} - x + 1,D = \left( {m,2}\right)$ ，当且仅当 $m$ 满足什么条件时，函数 $y =  f\left( x\right) ,x \in  D$ 具有性质 $H$ ,且 $b = \dfrac{1}{3}$ 是该函数的一个下界;

(3)设 $0 < a \leq  1$ ， $f\left( x\right)  = \ln \left( {x + 1}\right)  + {ax}\left( {x - 2}\right)$ ， $D = \left( {0,1}\right)$ ，若函数 $y = f\left( x\right)$ ， $x \in  D$ 具有性质 $H$ ， 求 $a$ 的取值范围；当 $a$ 在上述范围内变化时，若 $b$ 总是该函数的下界，求 $b$ 的取值范围.

15.(2024$\cdot$上海秋季卷$\cdot$21)对于一个函数 $y = f\left( x\right)$ 和一个点 $M\left( {a,b}\right)$ ，令 $s\left( x\right)  = {\left( x - a\right) }^{2} +  {\left( f\left( x\right)  - b\right) }^{2}$ ,若 $P\left( {{x}_{0},f\left( {x}_{0}\right) }\right)$ 是 $s\left( x\right)$ 取到最小值的点,则称 $P$ 是 $M$ 在 $f\left( x\right)$ 的“最近点”.

(1)对于 $f\left( x\right)  = \dfrac{1}{x}\left( {x > 0}\right)$ ，求证:对于点 $M\left( {0,0}\right)$ ，存在点 $P$ ，使得点 $P$ 是 $M$ 在 $f\left( x\right)$ 的“最近点”；

(2)对于 $f\left( x\right)  = {{e}}^{x},M\left( {1,0}\right)$ ，请判断是否存在一个点 $P$ ，它是 $M$ 在 $f\left( x\right)$ 的“最近点”，且直线 ${MP}$ 与 $y = f\left( x\right)$ 在点 $P$ 处的切线垂直；

(3)已知 $y = f\left( x\right)$ 在定义域 $\mathbf{R}$ 上存在导函数 $y = {f}^{\prime }\left( x\right)$ ，且函数 $y = g\left( x\right)$ 在定义域 $\mathbf{R}$ 上恒正，设点 ${M}_{1}\left( {t - 1,f\left( t\right)  - g\left( t\right) }\right)$ 、 ${M}_{2}\left( {t + 1,f\left( t\right)  + g\left( t\right) }\right)$. 若对任意的 $t \in  \mathbf{R}$ ，存在点 $P$ 同时是 ${M}_{1}$ 、 ${M}_{2}$ 在 $f\left( x\right)$ 的“最近点”，试判断 $y = f\left( x\right)$ 的单调性.

16.(2025$\cdot$杨浦区一模$\cdot$21)已知 $y = f\left( x\right)$ 是定义域为 $\left\lbrack  {0,1}\right\rbrack$ 的函数，实数 $p \in  \left( {0,1}\right)$ ，称函数 $y =  \left( {1 - p}\right) f\left( 0\right)  + {pf}\left( x\right)  - f\left( {px}\right) ,x \in  \left\lbrack  {0,1}\right\rbrack$ 为函数 $y = f\left( x\right)$ 的“ $p$ -生成函数”，记作 $y = {F}_{p}\left( x\right) ,x \in  \left\lbrack  {0,1}\right\rbrack$.

(1)若 $f\left( x\right)  = \cos {2\pi x}$ ，求函数 $y = {F}_{\dfrac{1}{2}}\left( x\right)$ 的值域；

( 2 )若 $f\left( x\right)  = a{x}^{2} + \ln \left( {1 + x}\right)$ ，函数 $y = {F}_{\dfrac{1}{3}}\left( x\right)$ 满足 ${F}_{\dfrac{1}{3}}\left( x\right)  \geq  0$ 对任意的 $0 \leq  x \leq  1$ 恒成立，求实数 $a$ 的取值范围;

(3)若 $y = f\left( x\right)$ 满足: \ding{172}  $f\left( 0\right)  = 0$ ;  \ding{173}  $y = f\left( x\right)$ 在 $\left( {0,1}\right)$ 上存在导函数 $y = {f}^{\prime }\left( x\right)$ ，且 $y = {f}^{\prime }\left( x\right)$ 在 $\left( {0,1}\right)$ 上是严格增函数;  \ding{174}  对于任意 $p \in  \left( {0,1}\right) ,y = f\left( x\right)$ 的“ $p$ -生成函数” $y = {F}_{p}\left( x\right)$ , $x \in  \left\lbrack  {0,1}\right\rbrack$ 的图像是一段连续曲线,求证: 函数 $y = \dfrac{f\left( x\right) }{x}$ 在 $\left( {0,1}\right)$ 上是严格增函数.

17.(2025$\cdot$虹口区一模$\cdot$21)设 $a \in  \mathbf{R}$ ， ${F}_{a}\left( x\right)  = \dfrac{f\left( x\right)  - f\left( a\right) }{x - a}$ ， $x \in  \left( {a - 1,a}\right)  \cup  \left( {a,a + 1}\right)$. 若函数 $y =  f\left( x\right)$ 满足 ${F}_{a}\left( x\right)  > 0$ 恒成立,则称函数 $y = f\left( x\right)$ 具有性质 $P\left( a\right)$.

(1)判断 $y = \sin x$ 是否具有性质 $P\left( 0\right)$ ，并说明理由;

(2)设 $f\left( x\right)  = {{e}}^{x} - x$ ，若函数 $y = f\left( x\right)$ 具有性质 $P\left( a\right)$ ，求实数 $a$ 的取值范围；

(3)设函数 $y = f\left( x\right)$ 的定义域为 $\mathbf{R}$ ，且对任意 $a \in  \mathbf{R}$ 以及 $t \in  \left( {0,1}\right)$ ，都有 ${F}_{a}\left( {a - t}\right)  < {F}_{a}\left( {a + t}\right)$. 若当 $x < 0$ 时,恒有 $f\left( x\right)  < 0$. 求证: 函数 $y = f\left( x\right)$ 对任意实数 $a$ 均具有性质 $P\left( a\right)$.

18.(2024$\cdot$虹口区一模$\cdot$21)已知 $y = f\left( x\right)$ 与 $y = g\left( x\right)$ 都是定义在 $\left( {0, + \infty }\right)$ 上的函数，若对任意 ${x}_{1}$ 、 ${x}_{2} \in  \left( {0, + \infty }\right)$ ,当 ${x}_{1} < {x}_{2}$ 时,都有 $g\left( {x}_{1}\right)  \leq  \dfrac{f\left( {x}_{1}\right)  - f\left( {x}_{2}\right) }{{x}_{1} - {x}_{2}} \leq  g\left( {x}_{2}\right)$ ,则称 $y = g\left( x\right)$ 是 $y = f\left( x\right)$ 的一个“控制函数”.

(1)判断 $y = {2x}$ 是否为函数 $y = {x}^{2}\left( {x > 0}\right)$ 的一个控制函数，并说明理由;

(2)设 $f\left( x\right)  = \ln x$ 的导数为 ${f}^{\prime }\left( x\right) ,0 < a < b$ ，求证:关于 $x$ 的方程 $\dfrac{f\left( b\right)  - f\left( a\right) }{b - a} = {f}^{\prime }\left( x\right)$ 在区间 $\left( {a,b}\right)$ 上有实数解;

(3)设 $f\left( x\right)  = x\ln x$ ，函数 $y = f\left( x\right)$ 是否存在控制函数？若存在，请求出 $y = f\left( x\right)$ 的所有控制函数; 若不存在, 请说明理由.

19.(2024$\cdot$静安区二模$\cdot$21)已知 $k \in  \mathbf{R}$ ，记 $f\left( x\right)  = {a}^{x} + k \cdot  {a}^{-x}$ ( $a > 0$ 且 $a \neq  1$ ).

(1)当 $a = {e}$ ( ${e}$ 是自然对数的底)时，试讨论函数 $y = f\left( x\right)$ 的单调性和最值;

(2)试讨论函数 $y = f\left( x\right)$ 的奇偶性；

(3)拓展与探究:

 \ding{172}  当 $k$ 在什么范围取值时，函数 $y = f\left( x\right)$ 的图像在 $x$ 轴上存在对称中心？请说明理由；

 \ding{173}  请提出函数 $y = f\left( x\right)$ 的一个新性质，并用数学符号语言表达出来. (不必证明)

20.(2024$\cdot$奉贤区一模$\cdot$21)若函数 $y = f\left( x\right)$ 满足:对任意的实数 $s$ 、 $t \in  \left( {0, + \infty }\right)$ ，有 $f\left( {s + t}\right)  >  f\left( s\right)  + f\left( t\right)$ 恒成立，则称函数 $y = f\left( x\right)$ 为“ $\sum$ 增函数”.

(1)求证:函数 $y = \sin x$ 不是“ $\sum$ 增函数”；

(2)若函数 $y = {2}^{x - 1} - x - a$ 是“ $\sum$ 增函数”，求实数 $a$ 的取值范围；

(3)设 $g\left( x\right)  = {{e}}^{x}\ln \left( {1 + x}\right)$ ，若曲线 $y = g\left( x\right)$ 在 $x = {x}_{0}$ 处的切线方程为 $y = x$ ，求 ${x}_{0}$ 的值，并证明函数 $y = g\left( x\right)$ 是“$\sum$ 增函数”.

\subsection*{考点5 > 求极值、最值及其应用}

\subsubsection*{一、填空题}

1. (2024$\cdot$徐汇区二模$\cdot$8)已知函数 $f\left( x\right)  = {x}^{3} + {2m}{x}^{2} - {nx} + m$ 在 $x = 1$ 处有极值 0，则 $m + n =$ \blkda[2]{}.

2.(2025$\cdot$静安区二模$\cdot$9)用总长为 ${14.8}{m}$ 的钢条制作一个长方体容器的框架，且容器底面的长边比短边长 0.5m(不计损耗). 若要使该容器的容积最大，则容器的高为\blkda[2]{}m。

3. (2024$\cdot$静安区二模$\cdot$11)已知实数 $a \in  \left( {0,6}\right)$ ，记 $f\left( x\right)  = \sqrt{x}\left( {x - a}\right)$. 若函数 $y = f\left( x\right)$ 在区间 $\left\lbrack  {0,2}\right\rbrack$ 上的最小值为 -2,则 $a$ 的值为\blkda[2]{}.

4.(2024$\cdot$闵行区二模$\cdot$11)对于任意的 ${x}_{1}$ 、 ${x}_{2} \in  \mathbf{R}$ ，且 ${x}_{2} > 0$ ，不等式 $\left| {{{e}}^{{x}_{1}} - {x}_{1}}\right|  + \left| {\ln {x}_{2} - {x}_{2}}\right|  > a$ 恒成立，则实数 $a$ 的取值范围\blkda[2]{}.

5.(2025$\cdot$嘉定区一模$\cdot$11)某公园为了美化环境，计划建造一座拱桥 DACBE. 已知该桥的剖面如图所示,其包括一段圆弧形桥面 $\overset{\frown}{ACB}$ 和两段长度相等的直线型桥面 ${AD}$ 、 ${BE}$ ,圆弧形桥面 $\overset{\frown}{ACB}$ 所在圆的半径为 $4{m}$ ,圆心 $O$ 在 ${DE}$ 上,且 ${AD}$ 和 ${BE}$ 所在直线与圆 $O$ 分别在连接点 $A$ 和 $B$ 处相切. 已知直线型桥面的修建费用是每米 0.4 万元，弧形桥面 $\overset{\frown}{ACB}$ 的修建费用是每米 2.5 万元,设 $\angle {ADO} = \theta$ ,根据空间限制及桥面坡度的限制, $\theta$ 的范围为 $\arcsin \dfrac{1}{3} \leq  \theta  \leq  \dfrac{\pi }{6}$ ，则当桥面修建总费用最低时 $\theta$ 的值为\blkda[2]{}.

\begin{center}
\includegraphics[width=0.2\textwidth]{bodyxb/src/200_1095_902_329_219_0.jpg}
\end{center}


第 5 题图

6.(2024$\cdot$黄浦区二模$\cdot$11)如图是某公园局部的平面示意图，图中的实线部分 (它由线段 ${CE}\text{ 、 }{DF}$ 与分别以 ${OC}\text{ 、 }{OD}$ 为直径的半圆弧组成) 表示一条步道. 其中的点 $C\text{ 、 }D$ 是线段 ${AB}$ 上的动点,点 $O$ 为线段 ${AB}\text{ 、 }{CD}$ 的中点,点 $E\text{ 、 }F$ 在以 ${AB}$ 为直径的半圆弧上,且 $\angle {OCE}$ 、 $\angle {ODF}$ 均为直角. 若 ${AB} = 1$ 百米,则此步道的最大长度为\blkda[2]{}百米.

\begin{center}
\includegraphics[width=0.2\textwidth]{bodyxb/src/200_1097_1211_329_198_0.jpg}
\end{center}


第 6 题图

7.(2025$\cdot$嘉定区二模$\cdot$11)某建筑公司欲设计一个正四棱锥形纪念碑，要求其顶点位于容积为 ${36\pi }{{m}}^{3}$ 的球形景观灯所在球面上. 考虑到抗风、抗震等结构安全需求,侧棱长度 $l$ 需满足 $2\sqrt{3} \leq  l \leq  3\sqrt{3}$. 当纪念碑体积取得最大值时，正四棱锥的侧棱长约为\blkda[2]{} ${m}$ (精确到 ${0.01}{m}$ ).

\subsubsection*{二、解答题}

8.(2024$\cdot$杨浦区一模$\cdot$18)设函数 $f\left( x\right)  = {{e}}^{x},x \in  \mathbf{R}$.

(1)求方程 ${\left( f\left( x\right) \right) }^{2} = f\left( x\right)  + 2$ 的实数解;

(2)若不等式 $x + b \leq  f\left( x\right)$ 对于一切 $x \in  \mathbf{R}$ 都成立，求实数 $b$ 的取值范围.

上海市高中数学“三年高考两年模拟”分类汇编

9.(2025$\cdot$上海秋季卷$\cdot$19)已知 $f\left( x\right)  = {x}^{2} - \left( {m + 2}\right) x + m\ln x,m \in  \mathbf{R}$.

(1)若 $f\left( 1\right)  = 0$ ，求不等式 $f\left( x\right)  \leq  {x}^{2} - 1$ 的解集；

(2)若函数 $y = f\left( x\right)$ 满足在 $\left( {0, + \infty }\right)$ 上存在极大值，求 $m$ 的取值范围.

10.(2024$\cdot$金山区一模$\cdot$19)网络购物行业日益发达，各销售平台通常会配备送货上门服务. 小金正在配送客户购买的电冰箱, 并获得了客户所在小区门户以及建筑转角处的平面设计示意图.

(1)为避免冰箱内部制冷液逆流，要求运送过程中发生倾斜时，外包装的底面与地面的倾斜角 $\alpha$ 不能超过 $\dfrac{\pi }{4}$ ，且底面至少有两个顶点与地面接触. 外包装看作长方体，如图 1 所示，记长方体的纵截面为矩形 ${ABCD},{AD} = {0.8}{m},{AB} = {2.4}{m}$ ,而客户家门高度为 ${2.3}{m}$ ,其他过道高度足够. 若以倾斜角 $\alpha  = \dfrac{\pi }{4}$ 的方式进客户家门,小金能否将冰箱运送入客户家中? 计算并说明理由;

(2)由于客户选择以旧换新服务，小金需要将客户长方体形状的旧冰箱进行回收. 为了省力，小金选择将冰箱水平推运 (冰箱背面水平放置于带滚轮的平板车上, 平板车长宽均小于冰箱背面). 推运过程中遇到一处直角过道,如图 2 所示,过道宽为 ${1.8}{m}$. 记此冰箱水平截面为矩形 ${EFGH},{EH} = {1.2}{m}$. 设 $\angle {PHG} = \beta$ ，当冰箱被卡住时 (即点 $H$ 、 $G$ 分别在射线 ${PR}$ 、 ${PQ}$ 上，点 $O$ 在线段 ${EF}$ 上)，尝试用 $\beta$ 表示冰箱高度 ${EF}$ 的长，并求出 ${EF}$ 的最小值，最后请帮助小金得出结论:按此种方式推运的旧冰箱，其高度的最大值是多少？(结果精确到 ${0.1}{m}$ )

\begin{center}
\includegraphics[width=0.5\textwidth]{bodyxb/src/201_749_1409_655_371_0.jpg}
\end{center}


第 10 题图

11. (2025$\cdot$金山区二模$\cdot$21)若函数 $y = f\left( x\right)$ 和 $y = g\left( x\right)$ 同时满足下列条件: \ding{172} 对任意 $x \in  \mathbf{R}$ ，都有 $f\left( x\right)  \leq  g\left( x\right)$ 成立； \ding{173}  存在 ${x}_{0} \in  \mathbf{R}$ ，使得 $f\left( {x}_{0}\right)  = g\left( {x}_{0}\right)$ ，则称函数 $y = g\left( x\right)$ 为 $y = f\left( x\right)$ 的“ $W$ 函数”,其中 ${x}_{0}$ 称为“ $W$ 点”.

(1)已知图像为一条直线的函数 $y = g\left( x\right)$ 是 $y = \sin x$ 的“ $W$ 函数”,请求出所有的“ $W$ 点”;

(2)设函数 $y = g\left( x\right)$ 为 $y = f\left( x\right)$ 的“ $W$ 函数”，其“ $W$ 点”组成集合 $M$ ；函数 $y = h\left( x\right)$ 为 $y = g\left( x\right)$ 的“ $W$ 函数”，其“ $W$ 点”组成集合 $N$. 试证明: “函数 $y = h\left( x\right)$ 为 $y = f\left( x\right)$ 的“$W$ 函数，”的一个充分必要条件是“$M \cap  N \neq  \varnothing$ ”；

(3)记 $f\left( x\right)  = \dfrac{x}{{{e}}^{x}}$ (e 为自然对数的底数)， $g\left( x\right)  = {kx} + m\left( {k\text{ 、 }m \in  \mathbf{R}}\right)$ ，若 $y = g\left( x\right)$ 为 $y = f\left( x\right)$ 的“$W$ 函数”，且“ $W$ 点” ${x}_{0} > 0$ ，求实数 $m$ 的最大值.

12.(2025$\cdot$宝山区一模$\cdot$21)已知 $y = f\left( x\right)$ 、 $y = g\left( x\right)$ 都是定义在实数集上的可导函数. 对于正整数 $k$ ,当 $m\text{ 、 }n$ 分别是 $y = f\left( x\right)$ 和 $y = g\left( x\right)$ 的驻点时,记 ${\Delta x} = \left| {m - n}\right|$ ,若 ${\Delta x} \leq  k$ ,则称 $f\left( x\right)$ 和 $g\left( x\right)$ 满足 $P\left( k\right)$ 性质; 当 ${x}_{1}\text{ 、 }{x}_{2} \in  \mathbf{R}$ ,且 $g\left( {x}_{1}\right)  \neq  g\left( {x}_{2}\right)$ 时,记 ${\Delta y} = \left| \dfrac{f\left( {x}_{1}\right)  - f\left( {x}_{2}\right) }{g\left( {x}_{1}\right)  - g\left( {x}_{2}\right) }\right|$ ,若 ${\Delta y} \geq  k$ , 则称 $f\left( x\right)$ 和 $g\left( x\right)$ 满足 $Q\left( k\right)$ 性质.

(1)若 $f\left( x\right)  = {2x} + 1,g\left( x\right)  = x$ ，判断 $f\left( x\right)$ 和 $g\left( x\right)$ 是否满足 $Q\left( 2\right)$ 性质，并说明理由；

(2)若 $f\left( x\right)  = {\left( x - 1\right) }^{2},g\left( x\right)  = \dfrac{{ax} + 1}{{{e}}^{x}}$ ，且 $f\left( x\right)$ 和 $g\left( x\right)$ 满足 $P\left( 1\right)$ 性质，求实数 $a$ 的取值范围；

(3)若 $y = f\left( x\right)$ 的最小正周期为 4，且 $g\left( {-1}\right)  = f\left( {-1}\right) ,g\left( 1\right)  = f\left( 1\right)$. 当 $x \in  \left\lbrack  {-1,3}\right\rbrack$ 时， $y =  f\left( x\right)$ 的驻点与其两侧区间的部分数据如下表所示:

\begin{center}

\begin{tabular}{|c|c|c|c|c|c|}
\hline
$x$ & -1 & (-1,1) & 1 & (1,3) & 3 \\
\cline{1-6}
${f}^{\prime }\left( x\right)$ & 0 & + & 0 & - & 0 \\
\cline{1-6}
$f\left( x\right)$ & 极小值 -1 & \phantom{X} & 极大值 1 & \phantom{X} & 极小值 -1 \\
\cline{1-6}
\hline
\end{tabular}
\end{center}

已知 $f\left( x\right)$ 和 $g\left( x\right)$ 满足 $Q\left( k\right)$ 性质,请写出 $f\left( x\right)  = g\left( x\right)$ 的充要条件,并说明理由.

13.(2025$\cdot$黄浦区一模$\cdot$21)函数 $y = f\left( x\right)$ 的定义域为 $D$ ，在 $D$ 上仅有一个极值点 ${x}_{0}$ ，方程 $f\left( x\right)  = 0$ 在 $D$ 上仅有两解,分别为 ${x}_{1}\text{ 、 }{x}_{2}$ ,且 ${x}_{1} < {x}_{0} < {x}_{2}$. 若 $\dfrac{{x}_{1} + {x}_{2}}{2} > {x}_{0}$ ,则称函数 $y = f\left( x\right)$ 在 $D$ 上的极值点左偏移; 若 $\dfrac{{x}_{1} + {x}_{2}}{2} < {x}_{0}$ ,则称函数 $y = f\left( x\right)$ 在 $D$ 上的极值点右偏移.

(1)设 $f\left( x\right)  = {x}^{2} - 1,D = \mathbf{R}$ ，判断函数 $y = f\left( x\right)$ 在 $D$ 上的极值点是否左偏移或右偏移；

(2)设 $m > 0$ 且 $m \neq  1$ ， $f\left( x\right)  = {x}^{3} - m{x}^{2} - x + m,D = \left( {0, + \infty }\right)$ ，求证:函数 $y = f\left( x\right)$ 在 $D$ 上的极值点右偏移;

(3)设 $a \in  \mathbf{R},f\left( x\right)  = \ln x - {ax},D = \left( {0, + \infty }\right)$ ，求证:当 $0 < a < {{e}}^{-1}$ 时，函数 $y = f\left( x\right)$ 在 $D$ 上的极值点左偏移.

14. (2024$\cdot$浦东新区一模$\cdot$21)设 $y = f\left( x\right)$ 是定义在 $\mathbf{R}$ 上的函数，若存在区间 $\left\lbrack  {a,b}\right\rbrack$ 和 ${x}_{0} \in  \left( {a,b}\right)$ ， 使得 $y = f\left( x\right)$ 在 $\left\lbrack  {a,{x}_{0}}\right\rbrack$ 上严格减，在 $\left\lbrack  {{x}_{0},b}\right\rbrack$ 上严格增，则称 $y = f\left( x\right)$ 为“含谷函数”， ${x}_{0}$ 为“谷点”, $\left\lbrack  {a,b}\right\rbrack$ 称为 $y = f\left( x\right)$ 的一个“含谷区间”.

(1)判断下列函数中，哪些是含谷函数？若是，请指出谷点；若不是，请说明理由.

 \ding{172}  $y = 2\left| x\right|$ ， \ding{173}  $y = x + \cos x$ ；

(2)已知实数 $m > 0$ ， $y = {x}^{2} - {2x} - m\ln \left( {x - 1}\right)$ 是含谷函数，且[2,4]是它的一个含谷区间，求 $m$ 的取值范围;

(3) 设 $p\text{ 、 }q \in  \mathbf{R},h\left( x\right)  =  - {x}^{4} + p{x}^{3} + q{x}^{2} + \left( {4 - {3p} - {2q}}\right) x$. 设函数 $y = h\left( x\right)$ 是含谷函数， $\left\lbrack  {a,b}\right\rbrack$ 是它的一个含谷区间,并记 $b - a$ 的最大值为 $L\left( {p,q}\right)$. 若 $h\left( 1\right)  \leq  h\left( 2\right)$ ,且 $h\left( 1\right)  \leq  0$ ,求 $L\left( {p,q}\right)$ 的最小值.

15.(2024$\cdot$嘉定区一模$\cdot$21)已知 $f\left( x\right)  = \dfrac{x}{{{e}}^{x}}$ ， $g\left( x\right)  = \dfrac{\ln x}{x}$.

(1)求函数 $y = f\left( x\right)$ 、 $y = g\left( x\right)$ 的单调区间和极值；

(2)请严格证明曲线 $y = f\left( x\right)$ 、 $y = g\left( x\right)$ 有唯一交点;

(3)对于常数 $a \in  \left( {0,\dfrac{1}{{e}}}\right)$ ，若直线 $y = a$ 和曲线 $y = f\left( x\right)$ 、 $y = g\left( x\right)$ 共有三个不同交点 $\left( {{x}_{1},a}\right)$ 、 $\left( {{x}_{2},a}\right) \text{ 、 }\left( {{x}_{3},a}\right)$ ,其中 ${x}_{1} < {x}_{2} < {x}_{3}$ ,求证: ${x}_{1},{x}_{2},{x}_{3}$ 成等比数列.

\subsection*{考点6 )利用导数研究函数的极值点及其应用}

\subsubsection*{一、填空题}

1.(2023$\cdot$上海秋季卷$\cdot$11)某公园欲建设一段斜坡，斜坡是一条直线，斜坡顶点距水平地面的高度为 $4{m}$ ，坡面与水平面所成角为 $\theta$ ，行人每沿着斜坡向上走 $1{m}$ 消耗的体力为 $\left( {{1.025} - \cos \theta }\right)$ ，欲使行人走上斜坡所消耗的总体力最小，则角 $\theta  =$ \blkda[2]{}.

2.(2025$\cdot$徐汇区一模$\cdot$9)设 $a \in  \mathbf{R}$ ， $f\left( x\right)  = {x}^{2} + {ax} + \ln x$ ，若函数 $y = f\left( x\right)$ 存在两个不同的极值点， 则 $a$ 的取值范围为\blkda[2]{}.

\subsubsection*{二、选择题}

3.(2025$\cdot$静安区二模$\cdot$15)设 $y = f\left( x\right)$ 是一个三次函数， $y = {f}^{\prime }\left( x\right)$ 为其导函数. 如图所示的是 $y =  x{f}^{\prime }\left( x\right)$ 的图像的一部分. 则 $f\left( x\right)$ 的极大值与极小值分别是 \dotfill{(\kh{})}

\begin{center}
\includegraphics[width=0.2\textwidth]{bodyxb/src/204_1097_1783_333_214_0.jpg}
\end{center}


第 3 题图

\xx{}{}{}{}$f\left( 1\right)$ 与 $f\left( {-1}\right)$

B. $f\left( {-1}\right)$ 与 $f\left( 1\right)$

C. $f\left( {-2}\right)$ 与 $f\left( 2\right)$

D. $f\left( 2\right)$ 与 $f\left( {-2}\right)$

\subsubsection*{三、解答题}

4.(2025$\cdot$嘉定区一模$\cdot$21)设 $A$ 为非空集合，函数 $y = f\left( x\right)$ 的定义域为 $D$. 若存在 ${x}_{0} \in  D$ 使得对任意的 $x \in  D$ 均有 $f\left( x\right)  - f\left( {x}_{0}\right)  \in  A$ ,则称 $f\left( {x}_{0}\right)$ 为函数 $y = f\left( x\right)$ 的一个 $A$ 值， ${x}_{0}$ 为相应的 $A$ 值点.

(1)若 $A = \left\lbrack  {-2,0}\right\rbrack  ,f\left( x\right)  = \sin x$. 证明: ${x}_{0} = {2k\pi } + \dfrac{1}{2}\pi ,k \in  \mathbf{Z}$ 是函数 $y = f\left( x\right)$ 的一个 $A$ 值点，并写出相应的 $A$ 值；

(2)若 $A = \lbrack 0, + \infty ),f\left( x\right)  =  - x,g\left( x\right)  = {x}^{2} + x + 1$. 分别判断函数 $y = f\left( x\right) \text{ 、 }y = g\left( x\right)$ 是否存在 $A$ 值? 若存在,求出相应的 $A$ 值点; 若不存在,说明理由;

(3)若 $A = ( - \infty ,0\rbrack$ ，且函数 $f\left( x\right)  = \ln x + a{x}^{2}\left( {a \in  \mathbf{R}}\right)$ 存在 $A$ 值，求函数 $y = f\left( x\right)$ 的 $A$ 值，并指出相应的 $A$ 值点.

\subsection*{考点7 )导数与函数的基本性质结合问题}

\subsubsection*{一、选择题}

1.(2024$\cdot$崇明区一模$\cdot$16)若存在实数 $a$ 、 $b$ ，对任意实数 $x \in  \left\lbrack  {0,1}\right\rbrack$ ，使得不等式 ${x}^{3} - m \leq  {ax} + b \leq  {x}^{3} + m$ 恒成立，则实数 $m$ 的取值范围是 \dotfill{(\kh{})}

\xx{}{}{}{}$\left\lbrack  {\dfrac{\sqrt{3}}{9}, + \infty }\right)$ B. $\left\lbrack  {\dfrac{8\sqrt{3}}{9}, + \infty }\right)$

C. $\left\lbrack  {\dfrac{\sqrt{3}}{3}, + \infty }\right)$ D. $\left\lbrack  {\dfrac{\sqrt{3}}{2}, + \infty }\right)$

\subsubsection*{二、解答题}

2. (2024$\cdot$嘉定区二模$\cdot$21)已知常数 $m \in  \mathbf{R}$ ,设 $f\left( x\right)  = \ln x + \dfrac{m}{x}$ ,

(1)若 $m = 1$ ，求函数 $y = f\left( x\right)$ 的最小值；

(2)是否存在 $0 < {x}_{1} < {x}_{2} < {x}_{3}$ ，且 ${x}_{1},{x}_{2},{x}_{3}$ 依次成等比数列，使得 $f\left( {x}_{1}\right)$ ， $f\left( {x}_{2}\right)$ ， $f\left( {x}_{3}\right)$ 依次成等差数列? 请说明理由;

(3)求证:“ $m \leq  0$ ”是“对任意 ${x}_{1}$ 、 ${x}_{2} \in  \left( {0, + \infty }\right)$ ， ${x}_{1} < {x}_{2}$ ，都有 $\dfrac{{f}^{\prime }\left( {x}_{1}\right)  + {f}^{\prime }\left( {x}_{2}\right) }{2} > \dfrac{f\left( {x}_{1}\right)  - f\left( {x}_{2}\right) }{{x}_{1} - {x}_{2}}$ ， 的充要条件.

\subsection*{考点8 > 利用导数研究函数的零点及其应用}

\subsubsection*{一、选择题}

1.(2024$\cdot$浦东新区二模$\cdot$16) ${f}_{0}\left( x\right)  = {a}_{m}{x}^{m} + {a}_{m - 1}{x}^{m - 1} + \cdots  + {a}_{1}x + {a}_{0}\left( {{a}_{m} \neq  0,m \geq  {10},m \in  \mathbf{Z}}\right)$ ，记 ${f}_{n}\left( x\right)  = {f}^{\prime }{}_{n - 1}\left( x\right) \left( {n = 1,2,\cdots ,m - 1}\right)$ ,令有穷数列 ${b}_{n}$ 为 ${f}_{n}\left( x\right)$ 零点的个数 $\left( {n = 1,2,\cdots ,m - 1}\right)$ , 则有以下两个结论:

 \ding{172}  存在 ${f}_{0}\left( x\right)$ ，使得 ${b}_{n}$ 为常数列； \ding{173}  存在 ${f}_{0}\left( x\right)$ ，使得 ${b}_{n}$ 为公差不为零的等差数列. 那么().

\xx{}{}{}{} \ding{172} 正确， \ding{173} 错误 B.  \ding{172} 错误， \ding{173} 正确

C.  \ding{172}  \ding{173} 都正确 D.  \ding{172}  \ding{173} 都错误

2. (2024$\cdot$普陀区一模$\cdot$16) 已知函数 $f\left( x\right)  = \left\{  \begin{array}{l} x + 1,x \in  \lbrack  - 1,0), \\  \dfrac{{3x} - 2}{1 - x},x \in  \lbrack 0,1). \end{array}\right.$ 若函数 $g\left( x\right)  = f\left( x\right)  - {mx} + \dfrac{2}{3}m$ 在 $\lbrack  - 1,1)$ 内有且仅有两个零点,则实数 $m$ 的取值范围是 \dotfill{(\kh{})}

\xx{}{}{}{}$\left( {-\dfrac{3}{2},0}\right\rbrack   \cup  \lbrack 3,8)$ B. $\left\lbrack  {-\dfrac{1}{2},0}\right\rbrack   \cup  \left\lbrack  {3,9}\right\rbrack$

C. $\left\lbrack  {-\dfrac{1}{2},0}\right\rbrack   \cup  \left\lbrack  {3,8}\right\rbrack$ D. $\left( {-\dfrac{3}{2},0}\right\rbrack   \cup  \lbrack 3,9)$

\subsubsection*{二、解答题}

3. (2025$\cdot$奉贤区二模$\cdot$21)函数 $y = f\left( x\right)$ ，其中 $f\left( x\right)  = {x}^{3} - {x}^{2} + \dfrac{1}{2}x + \dfrac{1}{4}$ ，定义域是一切实数.

(1)计算 $\mathop{\lim }\limits_{{h \rightarrow  0}}\dfrac{f\left( {2 + h}\right)  - f\left( 2\right) }{h}$ 的值并指出其几何意义;

(2)当 $x \in  \left( {0,\dfrac{1}{2}}\right)$ 时，方程 $f\left( x\right)  = a + x$ 只有一个解,求实数 $a$ 的取值范围；

(3) 设 ${x}_{1} = 0,{x}_{n + 1} = f\left( {x}_{n}\right) ,{y}_{1} = \dfrac{1}{2},{y}_{n + 1} = f\left( {y}_{n}\right) ,n \geq  1,n \in  \mathbf{N},{b}_{n} = {y}_{n} - {x}_{n}$. 求证: $\mathop{\sum }\limits_{{i = 1}}^{n}{b}_{n} \in  \left( {0,1}\right)$.

4.(2025$\cdot$虹口区二模$\cdot$21)对于定义在 $\mathbf{R}$ 上的函数 $y = f\left( x\right)$ 和 $y = g\left( x\right) ,a \in  \mathbf{R}$ ,设 ${M}_{a} = \{ t \mid  t =  \left. {f\left( x\right)  - g\left( a\right) ,x \geq  a}\right\}$.

(1)若 $f\left( x\right)  = {2}^{x} - 1,g\left( x\right)  = \cos x$ ，求 ${M}_{0}$ ；

(2)若 $f\left( x\right)  = {x}^{3} - 3{x}^{2},g\left( x\right)  =  - x,{M}_{a} \subseteq  \lbrack 0, + \infty )$ ，求实数 $a$ 的取值范围；

(3)已知对任意 $a \in  \mathbf{R}$ ，均有 ${M}_{a} = \lbrack 0, + \infty )$ ，记 $h\left( x\right)  = g\left( x\right)  - a$ ，求证:“对任意 $a \in  \mathbf{R}$ ，函数 $y =  h\left( x\right)$ 零点个数均有限”的充要条件是“$y = f\left( x\right)$ 在 $\mathbf{R}$ 上是严格增函数”.

5.(2024$\cdot$闵行区二模$\cdot$21)已知定义在 $\left( {0, + \infty }\right)$ 上的函数 $y = f\left( x\right)$ 的表达式为 $f\left( x\right)  = \sin x -  x\cos x$ ,其所有的零点按从小到大的顺序组成数列 $\left\{  {x}_{n}\right\}  \left( {n \geq  1,n \in  \mathbf{N}}\right)$.

(1)求函数 $y = f\left( x\right)$ 在区间 $\left( {0,\pi }\right)$ 上的值域；

(2)求证:函数 $y = f\left( x\right)$ 在区间 $\left( {{n\pi },\left( {n + 1}\right) \pi }\right) \left( {n \geq  1,n \in  \mathbf{N}}\right)$ 上有且仅有一个零点；

(3)求证: $\pi  < {x}_{n + 1} - {x}_{n} < \dfrac{\left( {n + 1}\right) \pi }{n}$.

6.(2025$\cdot$金山区一模$\cdot$21)对于函数 $y = f\left( x\right)$ 图像上不同的三点 $A\left( {{x}_{1},{y}_{1}}\right)$ 、 $B\left( {{x}_{2},{y}_{2}}\right)$ 、 $M\left( {{x}_{0},{y}_{0}}\right)$ [其中 ${x}_{0} \in  \left( {{x}_{1},{x}_{2}}\right)$ ],记点 $M$ 处的切线为 $l$ ,若 $l//{AB}$ ,则称 $M$ 为函数 $y = f\left( x\right)$ 在区间 $\left( {{x}_{1},{x}_{2}}\right)$ 上的“ $T$ 点”. 特别地,当 ${x}_{0} = \dfrac{{x}_{1} + {x}_{2}}{2}$ ,则称 $M$ 为函数 $y = f\left( x\right)$ 在区间 $\left( {{x}_{1},{x}_{2}}\right)$ 上的“和谐 $T$ 点”.

(1)设 $f\left( x\right)  = {x}^{2},M\left( {{x}_{0},{y}_{0}}\right)$ 是函数 $y = f\left( x\right)$ 在区间 $\left( {0,n}\right)$ 上的“ $T$ 点”，若 ${f}^{\prime }\left( {x}_{0}\right)  = 1$ ，求实数 $n$ 的值;

(2)设 $f\left( x\right)  = a\sin {2x} + \cos x + x - 1$ ，若函数 $y = f\left( x\right)$ 在区间 $\left( {0,{2\pi }}\right)$ 上恰有 3 个“$T$ 点”，求所有满足条件的实数 $a$ 的值组成的集合;

(3)设 $f\left( x\right)  = \ln x + b{x}^{2}\left( {b \in  \mathbf{R}}\right)$ ，试探究函数 $y = f\left( x\right)$ 的定义域内是否存在一个包含“和谐 $T$ 点” 的区间 $\left( {{x}_{1},{x}_{2}}\right)$ ? 若存在,求出该区间 $\left( {{x}_{1},{x}_{2}}\right)$ ; 若不存在,请说明理由.

7.(2025$\cdot$徐汇区一模$\cdot$21)已知定义域为 $D$ 的函数 $y = f\left( x\right)$ ，其导函数为 $y = {f}^{\prime }\left( x\right)$ ，若点 $\left( {{x}_{0},{y}_{0}}\right)$ 在导函数 $y = {f}^{\prime }\left( x\right)$ 图像上，且满足 ${f}^{\prime }\left( {x}_{0}\right)  \cdot  {f}^{\prime }\left( {y}_{0}\right)  \geq  0$ ，则称 ${x}_{0}$ 为函数 $y = f\left( x\right)$ 的一个“ $T$ 类数”,函数 $y = f\left( x\right)$ 的所有“ $T$ 类数”构成的集合称为“ $T$ 类集”.

(1)若 $f\left( x\right)  = \sin x$ ，分别判断 $\dfrac{\pi }{2}$ 和 $\dfrac{3\pi }{4}$ 是否为函数 $y = f\left( x\right)$ 的“ $T$ 类数”，并说明理由；

(2)设 $y = {f}^{\prime }\left( x\right)$ 的图像在 $\mathbf{R}$ 上连续不断,集合 $M = \left\{  {x \mid  {f}^{\prime }\left( x\right)  = 0}\right\}$. 记函数 $y = f\left( x\right)$ 的“ $T$ 类集” 为集合 $S$ ,若 $S \subset  \mathbf{R}$ ,求证: $M \neq  \varnothing$ ;

(3)已知 $f\left( x\right)  =  - \dfrac{1}{\omega }\cos \left( {{\omega x} + \varphi }\right) \left( {\omega  > 0}\right)$ ，若函数 $y = f\left( x\right)$ 的“ $T$ 类集”为 $\mathbf{R}$ 时 $\varphi$ 的取值构成集合 $A$ ,求当 $\varphi  \in  A$ 时 $\omega$ 的最大值.

8.(2025$\cdot$松江区一模$\cdot$21)定义在 $D$ 上的函数 $y = f\left( x\right)$ ，若对任意不同的两点 $A\left( {{x}_{1},f\left( {x}_{1}\right) }\right)$ ， $B({x}_{2},  f\left( {x}_{2}\right) )\left( {{x}_{1} < {x}_{2}}\right)$ ，都存在 ${x}_{0} \in  \left( {{x}_{1},{x}_{2}}\right)$ ，使得函数 $y = f\left( x\right)$ 在 ${x}_{0}$ 处的切线 $l$ 与直线 ${AB}$ 平行，则称函数 $y = f\left( x\right)$ 在 $D$ 上处处相依,其中 $l$ 称为直线 ${AB}$ 的相依切线, $\left( {{x}_{1},{x}_{2}}\right)$ 为函数 $y = f\left( x\right)$ 在 ${x}_{0}$ 的相依区间. 已知 $f\left( x\right)  =  - \left( {a + 1}\right) {x}^{2} + {ax}$.

(1)当 $a = 2$ 时，函数 $F\left( x\right)  = {x}^{3} + f\left( x\right)$ 在 $\mathbf{R}$ 上处处相依，证明:导函数 $y = {F}^{\prime }\left( x\right)$ 在 $\left( {0,1}\right)$ 上有零点；

(2)若函数 $G\left( x\right)  = \ln x + \dfrac{f\left( x\right) }{{x}^{2}}$ 在 $\left( {0, + \infty }\right)$ 上处处相依，且对任意实数 $m$ 、 $n$ ， $m > n > 0$ ，都有 $\dfrac{G\left( m\right)  - G\left( n\right) }{m - n} \leq  1$ 恒成立，求实数 $a$ 的取值范围；

(3)当 $a = 0$ 时， $H\left( x\right)  = \dfrac{{{e}}^{x}}{\sqrt{-f\left( x\right) }}\left( {x > 0}\right)$ ， $\left( {{x}_{1},{x}_{2}}\right)$ 为函数 $y = H\left( x\right)$ 在 ${x}_{0} = 1$ 的相依区间，证明: ${x}_{1} + {x}_{2} > 2$.

\section*{第 12 章 排列、组合与二项式定理}

\begin{center}

\begin{tabular}{|p{4cm}|p{12cm}|}
\hline
考点 & 课标要求与命题趋势 \\
\cline{1-2}
1. 排列组合 & \ding{172}  分类加法计数原理和分步乘法计数原理是解决计数问题的基础，称为基本计数原理. 通过实例，了解分类加法计数原理、分步乘法计数原理及其意义. 理解两个基本计数原理，运用计数原理探索排列、组合、二项式定理等问题 \ding{173}  通过实例，理解、掌握排列与组合的概念，能利用计数原理推导排列数公式、组合数公式. 会计算排列数与组合数, 熟练掌握排列组合的解题方法. 排列组合是高考的常考内容,一般会和分类加法原理与分步乘法原理结合在小题中考查,需重点复习 \\
\cline{1-2}
2. 二项式定理 & \ding{172}  能用多项式运算法则和计数原理证明二项式定理,会用二项式定理解决与二项展开式有关的简单问题 \ding{173}  理解、掌握二项式定理的通项公式. 二项式定理是高考的常考内容，需重点强化复习 \\
\cline{1-2}
\hline
\end{tabular}
\end{center}
\clearpage
\subsection*{考点1}

\subsubsection*{一、填空题}

1.(2024$\cdot$长宁区一模$\cdot$5)将 4 个人排成一排，若甲和乙必须排在一起，则共有\blkda[2]{}种不同排法.

2.(2025$\cdot$奉贤区一模$\cdot$5)若A、B、C、D、E 五人站成一排，如果 A、B 必须相邻，那么排法共\blkda[2]{}种.

3.(2024$\cdot$虹口区二模$\cdot$6)3 个男孩和 3 个女孩站成一排做游戏，3 个女孩不相邻的站法种数为\blkda[2]{}.

4.(2024$\cdot$闵行区二模$\cdot$7)现有 5 个工程队承建某工程的 5 个不同项目，每个工程队承建一个项目， 其中甲工程队不能承建 $A$ 项目,则不同的承建方案有\blkda[2]{}种. (用数字作答)

5.(2025$\cdot$宝山区一模$\cdot$7)已知关于正整数 $x$ 的方程 ${{C}}_{12}^{x - 1} = {{C}}_{12}^{{5x} - 5}$ ，则该方程的解为\blkda[2]{}.

6.(2025$\cdot$闵行区一模$\cdot$8)从 10 名数学老师中选出 3 人安排在 3 天的假期中值班，每天有且只有一人值班. 若老师甲必须参加且不安排在假期第一天值班, 则不同的值班安排方法种数为\blkda[2]{}.

7.(2024$\cdot$徐汇区一模$\cdot$8)要排出高一某班一天上午 5 节课的课表，其中语文、数学、英语、艺术、体育各一节，若要求语文、数学选一门第一节课上，且艺术、体育不相邻上课，则不同的排法种数是\blkda[2]{}.

8.(2024$\cdot$闵行区一模$\cdot$9)今年中秋和国庆共有连续 8 天小长假，某单位安排甲、乙、丙三名员工值班，每天都需要有人值班. 任选两名员工各值 3 天班，剩下的一名员工值 2 天班，且每名员工值班的日期都是连续的，则不同的安排方法数为\blkda[2]{}.

9.(2025$\cdot$上海秋季卷$\cdot$9)4 个家长和 2 个儿童去爬山，6 个人需要排成一条队列，要求队列的头和尾均是家长，则不同的排列个数有\blkda[2]{}种.

10.(2025$\cdot$虹口区二模$\cdot$10)已知 9 个小球的编号为 1、2、...、9，从中有放回地摸取小球三次，并依次记录其编号，若这三个编号成等差数列，则共有\blkda[2]{}种不同的摸取方法.

11.(2024$\cdot$上海秋季卷$\cdot$10)设集合 $A$ 中的元素皆为无重复数字的三位正整数，且元素中任意两者之积皆为偶数，求集合中元素个数的最大值\blkda[2]{}.

\begin{center}
\includegraphics[width=0.2\textwidth]{bodyxb/src/211_1135_463_290_243_0.jpg}
\end{center}


第 13 题图

12.(2024$\cdot$杨浦区二模$\cdot$10)有 5 名志愿者报名参加周六、周日的公益活动，若每天从这 5 人中安排 2 人参加，则恰有 1 人在这两天都参加的不同安排方式共有\blkda[2]{}种

13.(2024$\cdot$徐汇区二模$\cdot$10)如图所示，将一个四棱锥的每一个顶点染一种颜色，并使同一条棱上的两端点异色，如果只有 4 种颜色可供使用，则不同的染色方法有\blkda[2]{}种.

14.(2025$\cdot$青浦区一模$\cdot$10)已知函数 $y = f\left( x\right)$ 的定义域为 $\{  - 2, - 1,1,2\}$ ，值域为 $\{  - 2,2\}$ ，则满足条件的函数 $y = f\left( x\right)$ 最多有\blkda[2]{}个.

15.(2024$\cdot$上海春季卷$\cdot$12)已知 ${a}_{1} = 2,{a}_{2} = 4,{a}_{3} = 8,{a}_{4} = {16}$ ，若对任意实数 ${b}_{1}$ 、 ${b}_{2}$ 、 ${b}_{3}$ 、 ${b}_{4}$ ，均满足 $\left\{  {{a}_{i} + {a}_{j} \mid  1 \leq  i < j \leq  4}\right\}   = \left\{  {{b}_{i} + {b}_{j} \mid  1 \leq  i < j \leq  4}\right\}$ ,则有序数组 $\left\{  {{b}_{1},{b}_{2},{b}_{3},{b}_{4}}\right\}$ 有\blkda[2]{}对.

16.(2025$\cdot$闵行区二模$\cdot$12)设 ${x}_{1}$ 、 ${x}_{2}$ 、 ${x}_{3}$ 、 ${x}_{4}$ 、 ${x}_{5}$ 均是正整数，且 $\left\{  {x \mid  x = {x}_{m}{x}_{n}{x}_{p}{x}_{q},1 \leq  m < n < p < }\right.  q \leq  5\}  = \{ {108},{144},{288},{432}\}$ ，则 ${x}_{1} + {x}_{2} + {x}_{3} + {x}_{4} + {x}_{5}$ 的值为\blkda[2]{}.

17.(2024$\cdot$松江区二模$\cdot$12)某校高一数学兴趣小组一共有 30 名学生，学号分别为 1、2、3、...、30，老师要随机挑选 3 名学生参加某项活动, 要求任意两人的学号之差绝对值大于等于 5 , 则有\blkda[2]{}种不同的选择方法.

18.(2025$\cdot$虹口区一模$\cdot$12)已知项数为 10 的数列 $\left\{  {a}_{n}\right\}$ 中任一项均为集合 $\{ x \mid  1 \leq  x \leq  {10},x \in  \mathbf{N}\}$ 中的元素,且相邻两项满足 ${a}_{n} < {a}_{n + 1} + 3,n = 1,2,\cdots ,9$. 若 $\left\{  {a}_{n}\right\}$ 中任意两项都不相等,则满足条件的数列 $\left\{  {a}_{n}\right\}$ 有\blkda[2]{}个.

\subsubsection*{二、选择题}

19.(2025$\cdot$杨浦区二模$\cdot$14)3 名同学报名参加社团活动，有 4 个社团可以报名，这些社团招收人数不限，但每位同学只能报名其中 1 个社团，则这 3 位同学可能的报名结果共有\blkda[2]{}种\blkda[2]{}. \dotfill{(\kh{})}

\xx{}{}{}{}6 B. 24 C. 64 D. 81

20.(2024$\cdot$黄浦区二模$\cdot$13)某学校为了解学生参加体育运动的情况，用分层抽样的方法作抽样调查, 拟从初中部和高中部两层共抽取 40 名学生, 已知该校初中部和高中部分别有 500 名和 300 名学生，则不同的抽样结果的种数为 \dotfill{(\kh{})}

\xx{}{}{}{}${{C}}_{500}^{25} + {{C}}_{300}^{15}$ B. ${{C}}_{500}^{25} \cdot  {{C}}_{300}^{15}$ C. ${{C}}_{500}^{20} + {{C}}_{300}^{20}$ D. ${{C}}_{500}^{20} \cdot  {{C}}_{300}^{20}$

21.(2025$\cdot$奉贤区二模$\cdot$13)下列有关排列组合数的计算公式，错误的是 \dotfill{(\kh{})}

\xx{}{}{}{}${{P}}_{n}^{m} + {{P}}_{n}^{m - 1} = {{P}}_{n + 1}^{m}\left( {m\text{ 、 }n}\right.$ 是正整数，且 $\left. {m \leq  n}\right)$ B. ${{P}}_{n}^{m} = {{C}}_{n}^{m}{{P}}_{m}^{m}\left( {m\text{ 、 }n}\right.$ 是正整数，且 $m \leq  n$

C. ${{P}}_{n}^{m} = n{{P}}_{n - 1}^{m - 1}\left( {m\text{ 、 }n\text{ 是正整数,且 }m \leq  n}\right)$ D. ${{C}}_{n}^{m} + {{C}}_{n}^{m - 1} = {{C}}_{n + 1}^{m}\left( {m\text{ 、 }n}\right.$ 是正整数,且 $\left. {m \leq  n}\right)$

考点2 又 二项式定理

\subsubsection*{一、填空题}

1.(2024$\cdot$静安区一模$\cdot$2)在( $x + \dfrac{2}{x}{)}^{3}$ 的二项展开式中， $x$ 的系数为\blkda[2]{}.

2.(2025$\cdot$青浦区二模$\cdot$3)在(1-ax) ${}^{6}$ 的二项展开式中， ${x}^{3}$ 项的系数是 20，则实数 $a$ 的值是\blkda[2]{}.

3. (2024$\cdot$长宁区二模$\cdot$3)在( $x + \dfrac{1}{x}$ )的展开式中， ${x}^{2}$ 的系数为\blkda[2]{}.

4.(2025$\cdot$徐汇区一模$\cdot$3)在(1+x) ${}^{n}$ 的二项展开式中，若各项系数和为 32，则正整数 $n$ 的值为\blkda[2]{}.

5.(2025$\cdot$宝山区一模$\cdot$4)在( $x + 2{)}^{5}$ 的二项展开式中， ${x}^{3}$ 的系数为\blkda[2]{}.

6.(2025$\cdot$崇明区一模$\cdot$4)在(x-1) ${ }^{7}$ 的二项展开式中， ${x}^{3}$ 的系数为\blkda[2]{}.

7.(2025$\cdot$奉贤区一模$\cdot$6)在( ${x}^{6} + \dfrac{1}{x\sqrt{x}}{)}^{5}$ 的二项展开式中，常数项为\blkda[2]{}.(用数字作答)

8.(2025$\cdot$普陀区一模$\cdot$4)若( $x + {x}^{2}{)}^{5} = {a}_{0}{x}^{5} + {a}_{1}{x}^{6} + {a}_{2}{x}^{7} + {a}_{3}{x}^{8} + {a}_{4}{x}^{9} + {a}_{5}{x}^{10}$ ，则 ${a}_{3}$ 的值为\blkda[2]{}.

9.(2024$\cdot$上海春季卷$\cdot$4)在(x-1) ${}^{6}$ 的二项展开式中， ${x}^{4}$ 项的系数为\blkda[2]{}.

10.(2025$\cdot$虹口区一模$\cdot$4)在(x-2) ${}^{6}$ 的二项展开式中， ${x}^{3}$ 项的系数为\blkda[2]{}.

11.(2025$\cdot$上海秋季卷$\cdot$4)在二项式( ${2x} - 1{)}^{5}$ 的展开式中，常数项为\blkda[2]{}.

12.(2025$\cdot$奉贤区二模$\cdot$4)在( $\sqrt{x} - \dfrac{2}{x}{)}^{9}$ 的二项展开式中，常数项为\blkda[2]{}.(用数字作答)

13.(2025$\cdot$浦东新区 一模$\cdot$4)在( ${x}^{2} + \dfrac{1}{x}{)}^{6}$ 的展开式中， ${x}^{3}$ 的系数为\blkda[2]{}.(用数字作答)

14.(2024$\cdot$闵行区一模$\cdot$4)已知 ${\left( x - 1\right) }^{4} = {a}_{0} + {a}_{1}x + {a}_{2}{x}^{2} + {a}_{3}{x}^{3} + {a}_{4}{x}^{4}$ ，则 ${a}_{2} =$ \blkda[2]{}.

15.(2024$\cdot$嘉定区二模$\cdot$4)在二项式 ${\left( x + 1\right) }^{5}$ 的展开式中， ${x}^{2}$ 项的系数是\blkda[2]{}.

16.(2024$\cdot$闵行区二模$\cdot$4)在( ${2x} - 1{)}^{6}$ 的二项展开式中， ${x}^{3}$ 项的系数为\blkda[2]{}.

17.(2024$\cdot$清东新区二模$\cdot$4)在( $3{x}^{2} + \dfrac{1}{x}{)}^{5}$ 的二项展开式中， ${x}^{4}$ 项的系数为\blkda[2]{}.(用数值回答)

18.(2024$\cdot$青浦区二模$\cdot$4)( $\sqrt{x} + \dfrac{2}{\sqrt{x}}{)}^{6}$ 的二项展开式中的常数项为\blkda[2]{}.

19.(2024$\cdot$虹口区一模$\cdot$5)在( $x - \dfrac{2}{\sqrt{x}}$ )的二项展开式中， $x$ 项的系数为\blkda[2]{}.

20.(2024$\cdot$黄浦区二模$\cdot$5)若( $a{x}^{2} + \dfrac{1}{x}{\text{ ) }}^{5}$ 的展开式中 ${x}^{4}$ 的系数是 -80，则实数 $a =$ \blkda[2]{}.

21.(2024$\cdot$崇明区一模$\cdot$5)在( $x + \dfrac{2}{{x}^{2}}$ )5的展开式中， ${x}^{2}$ 的系数为\blkda[2]{}.(用数字作答)

22.(2024$\cdot$松江区二模$\cdot$5)已知 ${x}^{7} = {a}_{0} + {a}_{1}\left( {x - 1}\right)  + {a}_{2}{\left( x - 1\right) }^{2} + \cdots  + {a}_{7}{\left( x - 1\right) }^{7}$ ，则 ${a}_{5} =$ \blkda[2]{}.

23.(2024$\cdot$青浦区一模$\cdot$5)在(3-2x) ${}^{6}$ 的二项展开式中， ${x}^{3}$ 项的系数是\blkda[2]{}.(用数值作答)

24.(2024$\cdot$普陀区一模$\cdot$5)设(1-2x) $x = {a}_{0} + {a}_{1}x + {a}_{2}{x}^{2} + {a}_{3}{x}^{3} + {a}_{4}{x}^{4} + \cdots  + {a}_{n}{x}^{n}$ ，若 ${a}_{0} + {a}_{4} = {17}$. 则 $n =$ \blkda[2]{}.

25.(2024$\cdot$杨浦区二模$\cdot$5)已知二项式(1+x) ${1}^{10}$ ，其展开式中含 ${x}^{2}$ 项的系数为\blkda[2]{}.

26.(2025$\cdot$闵行区二模$\cdot$5)在( ${2x} + \dfrac{1}{x}{)}^{6}$ 的二项展开式中，常数项是\blkda[2]{}.(用数值作答)

27.(2025$\cdot$闵行区二模$\cdot$5)在( ${ax} + \dfrac{1}{x}{)}^{5}$ 展开式中 $x$ 的系数为 80，则实数 $a$ 的值为\blkda[2]{}.

28.(2025$\cdot$嘉定区二模$\cdot$5)在( ${2x} - \dfrac{1}{\sqrt{x}}$  ${)}^{6}$ 的二项展开式中，常数项的值为\blkda[2]{}.

29.(2025$\cdot$金山区一模$\cdot$5)在(3x-1) ${}^{6}$ 的二项展开式中， ${x}^{2}$ 项的系数为\blkda[2]{}.

30.(2025$\cdot$黄浦区一模$\cdot$5)在( $x + \dfrac{1}{x}{)}^{6}$ 的二项展开式中，常数项为\blkda[2]{}.

31.(2025$\cdot$闵行区一模$\cdot$6)在( $x + \dfrac{1}{x}$ )的二项展开式中， ${x}^{4}$ 项的系数为\blkda[2]{}.

32.(2025$\cdot$长宁区一模$\cdot$6)在( $x - \dfrac{1}{x}{)}^{6}$ 的二项展开式中的常数项是\blkda[2]{}.

33.(2024$\cdot$上海秋季卷$\cdot$6)在(x+1)”的二项展开式中，若各项系数和为32，则 ${x}^{2}$ 项的系数为\blkda[2]{}.

34.(2025$\cdot$上海春季卷$\cdot$6)已知 ${\left( x + \dfrac{a}{x}\right) }^{6}$ 的展开式中常数项是 20，则 $a =$ \blkda[2]{}.

35.(2025$\cdot$宝山区二模$\cdot$6)在( $x + \dfrac{2}{x}{)}^{6}$ 的二项展开式中， ${x}^{2}$ 项的系数为\blkda[2]{}.

36.(2025$\cdot$杨浦区二模$\cdot$6)在( $x + \dfrac{1}{\sqrt{x}}{\text{ ) }}^{6}$ 的二项展开式中，常数项的值为\blkda[2]{}.

37.(2025$\cdot$普陀区二模$\cdot$6)设 $t \in  \mathbf{R}$ ，若 $\left( {1 + \dfrac{t}{x}}\right) {\left( 1 - x\right) }^{6}$ 的展开式中 ${x}^{3}$ 项的系数为10，则 $t =$ \blkda[2]{}.

38.(2024$\cdot$普陀区二模$\cdot$6)设(1+x) ${x}^{n} = {a}_{0} + {a}_{1}x + {a}_{2}{x}^{2} + \cdots  + {a}_{n}{x}^{n}\left( {n \geq  1,n \in  \mathbf{N}}\right)$ ，若 ${a}_{5} > {a}_{4}$ ，且 ${a}_{5} > {a}_{6}$ ,则 $\mathop{\sum }\limits_{{i = 1}}^{n}{a}_{i} =$ \blkda[2]{}.

39.(2024$\cdot$徐汇区二模$\cdot$7)已知 ${\left( {x}^{3} + \dfrac{1}{{x}^{2}}\right) }^{n}$ 的二项展开式中各项系数和为 1024，则展开式中常数项的值为\blkda[2]{}.

40.(2024$\cdot$徐汇区一模$\cdot$7)已知 ${\left( 1 - x\right) }^{10} = {a}_{0} + {a}_{1}x + {a}_{2}{x}^{2} + \cdots  + {a}_{10}{x}^{10}$ ，则 ${a}_{1} + {a}_{3} + {a}_{5} + {a}_{7} + {a}_{9} =$ \blkda[2]{}.

41.(2024$\cdot$松江区一模$\cdot$7)在二项式(3+x) $x$ 函数用式中， ${x}^{2}$ 项的系数是常数项的 5 倍，则 $n =$ \blkda[2]{}.

42.(2025$\cdot$嘉定区一模$\cdot$7)在( $x - \dfrac{1}{x}$ )的二项展开式中 ${x}^{3}$ 项的系数为\blkda[2]{}.

43.(2024$\cdot$崇明区二模$\cdot$7)若(x+a) ${}^{5}$ 的二项式展开式中 ${x}^{2}$ 的系数为 10，则 $a =$ \blkda[2]{}.

44.(2025$\cdot$浦东新区二模$\cdot$7)在( $x - \dfrac{1}{x}{)}^{10}$ 的二项展开式中，常数项为\blkda[2]{}.

45.(2025$\cdot$松江区一模$\cdot$7)已知 ${\left( x + 2\right) }^{4} = {a}_{0} + {a}_{1}x + {a}_{2}{x}^{2} + {a}_{3}{x}^{3} + {a}_{4}{x}^{4}$ ，则 ${a}_{1} + {a}_{2} + {a}_{3} + {a}_{4} =$ \blkda[2]{}.

46.(2024$\cdot$嘉定区一模$\cdot$8)在(1+-2x) ${}^{6}$ 的二项展开式中，系数最大的项为\blkda[2]{}.

47.(2024$\cdot$宝山区一模$\cdot$8)若对于任意实数 $x$ ，都有 ${x}^{4} = {a}_{0} + {a}_{1}\left( {x + 2}\right)  + {a}_{2}{\left( x + 2\right) }^{2} + {a}_{3}{\left( x + 2\right) }^{3} +  {a}_{4}{\left( x + 2\right) }^{4}$ ，则 ${a}_{3}$ 的值为\blkda[2]{}.

48.(2023$\cdot$上海春季卷$\cdot$8)设(1-2x) $= {a}_{0} + {a}_{1}x + {a}_{2}{x}^{2} + {a}_{3}{x}^{3} + {a}_{4}{x}^{4}$ ，则 ${a}_{0} + {a}_{4} =$ \blkda[2]{}.

49.(2025$\cdot$崇明区二模$\cdot$9)若( $x + 1{)}^{10} = {a}_{0} + {a}_{1}\left( {x - 1}\right)  + {a}_{2}{\left( x - 1\right) }^{2} + \cdots  + {a}_{10}{\left( x - 1\right) }^{10}$ ，则 ${a}_{0} +  {a}_{1} + {a}_{2} + \cdots  + {a}_{10} =$ \blkda[2]{}.

50.(2024$\cdot$奉贤区二模$\cdot$9)已知多项式(1+x)(1−-x) ${}^{5} = {a}_{0} + {a}_{1}x + {a}_{2}{x}^{2} + {a}_{3}{x}^{3} + {a}_{4}{x}^{4} + {a}_{5}{x}^{5} +  {a}_{6}{x}^{6}$ 对一切实数 $x$ 恒成立，则 ${a}_{0} + {a}_{3} =$ \blkda[2]{}.

51.(2025$\cdot$青浦区一模$\cdot$9)( $x + y$ )( $x - y{\text{ ) }}^{6}$ 的展开式中， ${x}^{4}{y}^{3}$ 项的系数为\blkda[2]{}.

52.(2023$\cdot$上海秋季卷$\cdot$10)已知(1+2023x) ${}^{100} + {\left( {2023} - x\right) }^{100} = {a}_{0} + {a}_{1}x + {a}_{2}{x}^{2} + \cdots  + {a}_{100}{x}^{100}$ ，且 $k \in  \{ 0,1,2,3,\cdots ,{100}\}$ ，当 ${a}_{k} < 0$ 时，正整数 $k$ 的最大值为\blkda[2]{}.

53.(2024$\cdot$金山区一模$\cdot$10)若(10x + 6y) ${}^{3} = a{x}^{3} + b{x}^{2}y + {cx}{y}^{2} + d{y}^{3}$ ，则 $- a + {2b} - {4c} + {8d} =$ \blkda[2]{}.

54.(2024$\cdot$金山区二模$\cdot$10)在(1+x) ${}^{5}{\left( 1 + y\right) }^{3}$ 的展开式中，记 ${x}^{m}{y}^{n}$ 项的系数为 $f\left( {m,n}\right)$ ，则 $f\left( {3,0}\right)  + f\left( {2,1}\right)  =$ \blkda[2]{}.

55.(2024$\cdot$杨浦区一模$\cdot$9)已知 ${\left( 1 + x\right) }^{m} + {\left( 1 + x\right) }^{n} = {a}_{0} + {a}_{1}x + {a}_{2}{x}^{2} + \cdots  + {a}_{m + n}{x}^{m + n}$ ( $m$ 、 $n$ 为正整数) 对任意实数 $x$ 都成立,若 ${a}_{1} = {12}$ ,则 ${a}_{2}$ 的最小值为\blkda[2]{}.

\subsubsection*{二、选择题}

56.(2024$\cdot$奉贤区一模$\cdot$15)若 ${\left( x - \dfrac{2}{{x}^{2}}\right) }^{n} + {\left( x + \dfrac{1}{x}\right) }^{n}$ ( $2 \leq  n \leq  {10}$ ， $n \in  \mathbf{N}$ )的展开式中存在常数项，则下列选项中 $n$ 的取值不可能是 \dotfill{(\kh{})}

\xx{}{}{}{}3 B. 4 C. 5 D. 6

\section*{第 13 章 概率与统计}

\begin{center}

\begin{tabular}{|p{4cm}|p{12cm}|}
\hline
考点 & 课标要求与命题趋势 \\
\cline{1-2}
1. 随机抽样 & \ding{172}  简单随机抽样; 通过实例,了解简单随机抽样的含义及其解决问题的过程,掌握两种简单随机抽样方法:抽签法和随机数法. 会计算样本均值和样本方差，了解样本与总体的关系 \ding{173}  分层随机抽样: 通过实例, 了解分层随机抽样的特点和适用范围, 了解分层随机抽样的必要性, 掌握各层样本量比例分配的方法. 结合具体实例, 掌握分层随机抽样的样本均值和样本方差 \ding{174}  抽样方法的选择:在简单的实际情境中，能根据实际问题的特点，设计恰当的抽样方法解决问题. 理解、掌握简单随机抽样、分层抽样定义及计算, 理解、掌握总体样本估计的定义及计算. 统计与数字特征是高考的常考内容, 需重点强化复习 \\
\cline{1-2}
2. 图表类统计图综合 & \ding{172}  知道获取数据的基本途径，包括:统计报表和年鉴、社会调查、试验设计、普查和抽样、 互联网等。了解总体、样本、样本量的概念，了解数据的随机性 \ding{173}  统计图表:如根据实际问题的特点，选择恰当的统计图表对数据进行可视化描述，体会合理使用统计图表的重要性 \\
\cline{1-2}
3. 样本的数字特征 & \ding{172}  结合具体实例，理解样本点和有限样本空间的含义，理解随机事件与样本点的关系 \ding{173}  通过实例，理解概率的性质，掌握随机事件概率的运算法则 \ding{174}  结合实例，会用频率估计概率 \\
\cline{1-2}
4. 变量间的相关关系 & 了解随机事件的并、交与互斥的含义, 能结合实例进行随机事件的并、交运算 \\
\cline{1-2}
5. 互斥事件的概率计算 & \ding{172}  结合具体实例，理解古典概型，能计算古典概型中简单随机事件的概率。 \ding{173}  掌握互斥事件的概率计算 \\
\cline{1-2}
6. 古典概率 & \ding{172}  理解、掌握古典概率的定义，并会相关计算. 古典概率是高考的常考内容，一般考查古典概型的概率计算及互斥、对立事件的辨析及计算，需强化训练 \ding{173}  掌握对立事件、相互独立事件的概率求解，会求古典概率、条件概率、全概率，是高考命题热点,会概率统计的综合运算及知识杂糅问题 \\
\cline{1-2}
7. 条件概率 & \ding{172}  结合古典概型，了解条件概率，能计算简单随机事件的条件概率 \ding{173}  会条件概率和全概率的计算. 该内容是高考必考内容, 一般结合条件概率、全概率综合考查, 需重点强化复习 \\
\cline{1-2}
8. 全概率公式 & 结合古典概型, 了解条件概率与独立性的关系; 会利用乘法公式计算概率; 会利用全概率公式计算概率 \\
\cline{1-2}
9. 正态分布指定区间的概率 & \ding{172}  通过误差模型，了解服从正态分布的随机变量. 通过具体实例、借助频率直方图的几何直观，了解正态分布的特征；了解正态分布的均值、方差及其含义 \ding{173}  理解、掌握正态分布的定义及指定区间的概率计算 \\
\cline{1-2}
10. 独立性检验为载体及其应用 & \ding{172}  结合有限样本空间, 了解两个随机事件独立性的含义. 结合古典概型, 利用独立性计算概率 \ding{173}  熟练掌握独立性检验和线性回归直线方程的求解，该内容会继续作为载体内容命题 \ding{174}  熟练掌握二项分布、超几何分布及其他类别的分布与期望方差问题，这是高考命题热点 \\
\cline{1-2}
11. 离散型分布及期望、方差 & \ding{172}  结合实例、能用样本估计总体的集中趋势参数(平均数、中位数、众数)，理解集中趋势参数的统计含义 \ding{173}  结合实例，能用样本估计总体的离散程度参数(标准差、方差、极差)，理解离散程度参数的统计含义 \ding{174}  结合实例，能用样本估计总体的取值规律 \ding{175}  结合实例，能用样本估计百分位数，理解百分位数的统计含义 \\
\cline{1-2}
\hline
\end{tabular}
\end{center}

考点1 》 随机抽样

\subsubsection*{一、填空题}

1.(2024$\cdot$静安区一模$\cdot$6)某果园种植了 222 棵苹果树，现从中随机抽取了 20 棵苹果树，算得这 20 棵苹果树平均每棵产量为 ${28}{{kg}}$ ,则预估该果园的苹果产量为\blkda[2]{} ${{kg}}$.

2.(2025$\cdot$虹口区二模$\cdot$6)某公司为了解用电量 $y$ (单位: ${{kW}} \cdot  {h}$ )与气温 $x$ (单位: ${}^{ \circ  }{C}$ )之间的关系,随机统计了 4 天的用电量与当天气温,绘制了如右表格,由表中数据可得回归方程 $y =  - {2x} + {59.5}$ ，则实数 $m =$ \blkda[2]{}.

\begin{center}

\begin{tabular}{|c|c|c|c|c|}
\hline
$x$ & -1 & 10 & 13 & 18 \\
\cline{1-5}
$y$ & 62 & 38 & 34 & $m$ \\
\cline{1-5}
\hline
\end{tabular}
\end{center}

3.(2024$\cdot$黄浦区一模$\cdot$7)某城市 30 天的空气质量指数如下:29、26、28、29、38、29、26、40、31、35、 44、33、28、30、86、65、53、70、34、36、41、31、38、63、60、56、34、74、34. 则这组数据的第 75 百分位数为\blkda[2]{}.

4.(2024$\cdot$长宁区一模$\cdot$7)现利用随机数表法从编号为 00、01、02、...、18、19 的 20 支水笔中随机选取 6 支, 选取方法是从下列随机数表第 1 行的第 9 个数字开始由左到右依次选取两个数字, 则选出来的第 6 支水笔的编号为\blkda[2]{}.

$\begin{array}{llllll} {95226000} & {49840128} & {66175168} & {39682927} & {43772366} & {27096623} \end{array}$

92580956 43890890 06482834 59741458 29778149 64608925

\subsubsection*{二、选择题}

5.(2024$\cdot$普陀区二模$\cdot$13)从放有两个红球、一个白球的袋子中一次任意取出两个球，两个红球分别标记为 $A\text{ 、 }B$ ,白球标记为 $C$ ,则它的一个样本空间可以是 \dotfill{(\kh{})}

\xx{}{}{}{}$\{ {AB},{BC}\}$ B. $\{ {AB},{AC},{BC}\}$

C. $\{ {AB},{BA},{BC},{CB}\}$ D. $\{ {AB},{BA},{AC},{CA},{CB}\}$

6.(2024$\cdot$闵行区一模$\cdot$14)某校读书节期间，共120名同学获奖(分金、银、铜三个等级)，从中随机抽取 24 名同学参加交流会，若按高一、高二、高三分层随机抽样，则高一年级需抽取 6 人；若按获奖等级分层随机抽样, 则金奖获得者需抽取 4 人. 下列说法正确的是 \dotfill{(\kh{})}

\xx{}{}{}{}高二和高三年级获奖同学共 80 人 B. 获奖同学中金奖所占比例一定最低

C. 获奖同学中金奖所占比例可能最高 D. 获金奖的同学可能都在高一年级

\subsection*{考点2 > 图表类统计图综合}

\begin{center}
\includegraphics[width=0.2\textwidth]{bodyxb/src/216_1301_1880_144_152_0.jpg}
\end{center}


第 1 题图

\subsubsection*{一、填空题}

1.(2025$\cdot$青浦区二模$\cdot$4)如图是 6 株果树植株挂果个数(两位数)的茎叶图，则 6 株果数植株挂果个数的中位数为\blkda[2]{}.

2.(2024$\cdot$宝山区一模$\cdot$5)在一次为期30天的博览会上，主办方统计了每天的参观人数(单位:千人)，得到样本的茎叶图(如图)，则该样本的第 70 百分位数是\blkda[2]{}.

\begin{center}
\includegraphics[width=0.2\textwidth]{bodyxb/src/217_1154_170_306_196_0.jpg}
\end{center}


第 2 题图

3.(2024$\cdot$徐汇区一模$\cdot$5)某学校组织全校学生参加网络安全知识竞赛，成绩 (单位: 分) 的频率分布直方图如图所示, 数据的分组依次为[20,40)、 [40,60)、[60,80)、[80,100]，若该校的学生总人数为 1000，则成绩低于 60 分的学生人数为\blkda[2]{}。

4.(2025$\cdot$崇明区二模$\cdot$7)某次数学考试后，随机选取 14 名学生的成绩，得到如下茎叶图，其中个数部分作为 “叶”，百位数和十位数作为 “茎”，若该组数据的第 25 百分位数是 87，则 $x$ 的值为\blkda[2]{}.

5.(2025$\cdot$黄浦区一模$\cdot$7)从 $A$ 校高一年级学生中抽取 66 名学生测量他们的身高，其中最大值为 184 cm，最小值 152 cm，绘制身高频率分布直方图，若组距为 3，且第一组下限为 151.5，则组数为\blkda[2]{}.

6.(2025$\cdot$徐汇区一模$\cdot$7)某景点对 30 天内每天的游客人数(单位:万人)进行统计，得到样本的茎叶图(如图所示)，则该样本的第 75 百分位数是\blkda[2]{}.

\begin{center}
\includegraphics[width=0.3\textwidth]{bodyxb/src/217_290_877_457_275_0.jpg}
\end{center}


第 3 题图

\begin{center}
\includegraphics[width=0.2\textwidth]{bodyxb/src/217_820_876_180_277_0.jpg}
\end{center}


第 4 题图

\begin{center}
\includegraphics[width=0.2\textwidth]{bodyxb/src/217_1069_875_293_279_0.jpg}
\end{center}


第 6 题图

7.(2023$\cdot$上海春季卷$\cdot$7)某校抽取 100 名学生测身高，其中身高最大值为 186 cm，最小值 154 cm，根据身高数据绘制频率组距分布直方图，组距为 5，且第一组下限为 153.5，则组数为\blkda[2]{}.

8.(2024$\cdot$崇明区一模$\cdot$8)如图是小王同学在篮球赛中得分记录的茎叶图，则他平均每场得分为\blkda[2]{}.

9.(2024$\cdot$黄浦区一模$\cdot$9)某校共有 400 名学生参加了趣味知识竞赛(满分:150 分)，且每位学生的竞赛成绩均不低于 90 分. 将这 400 名学生的竞赛成绩分组如下: [90,100), [100,110), [110,120)、[120,130)、[130,140)、[140,150]，得到的频率分布直方图，如图所示，则这 400 名学生中竞赛成绩不低于 120 分的人数为\blkda[2]{}.

10.(2024$\cdot$浦东新区一模$\cdot$9)小明为了解自己每天花在体育锻炼上的时间(单位:min)，连续记录了 7 天的数据并绘制成如图所示的茎叶图，则这组数据的第 60 百分位数是\blkda[2]{}.

\begin{center}
\includegraphics[width=0.2\textwidth]{bodyxb/src/217_312_1886_229_115_0.jpg}
\end{center}


第 8 题图

\begin{center}
\includegraphics[width=0.4\textwidth]{bodyxb/src/217_612_1740_508_258_0.jpg}
\end{center}


第 9 题图

\begin{center}
\includegraphics[width=0.2\textwidth]{bodyxb/src/217_1189_1793_152_205_0.jpg}
\end{center}


第 10 题图

\subsubsection*{二、选择题}

\begin{center}
\includegraphics[width=0.4\textwidth]{bodyxb/src/218_922_233_473_285_0.jpg}
\end{center}


第 11 题图

11.(2023$\cdot$上海春季卷$\cdot$14)如图所示，下面是出口，上面是进口，下列选项叙述错误的是 \dotfill{(\kh{})}

\xx{}{}{}{}从 2018 年开始, 2021 年进出口总额增长率最大

B. 从 2018 年开始, 进出口总额逐年增大

C. 从 2018 年开始, 进口总额逐年增大

D. 从 2018 年开始, 2020 年进出口总额增长率最小

12.(2025$\cdot$普陀区一模$\cdot$14)某机构对 2014 年至 2023 年的中国新能源汽车的年销售量进行了统计， 结果如图所示(单位: 万辆), 则下列结论中正确的是 \dotfill{(\kh{})}

\begin{center}
\includegraphics[width=0.6\textwidth]{bodyxb/src/218_375_695_776_374_0.jpg}
\end{center}


第 12 题图

\xx{}{}{}{}这十年中国新能源汽车年销售量的中位数为 123

B. 这十年中国新能源汽车年销售量的极差为 721

C. 这十年中国新能源汽车年销售量的第 70 百分位数为 136.6

D. 这十年中的前五年的年销售量的方差小于后五年的年销售量的方差

13.(2024$\cdot$虹口区一模$\cdot$14)空气质量指数 AQI 是反映空气质量状况的指数，其对应关系如下表: 为监测某化工厂排放废气对周边空气质量指数的影响, 某科学兴趣小组在工厂附近某处测得 10 月1日～20 日 AQI 的数据并绘成折线图如下:

\begin{center}
\includegraphics[width=0.5\textwidth]{bodyxb/src/218_103_1526_698_321_0.jpg}
\end{center}


第 13 题图

\begin{center}

\begin{tabular}{|c|c|c|c|c|c|c|}
\hline
AQI 指数值 & 0\textasciitilde 50 & 51\textasciitilde  100 & 101〜 150 & 151～ 200 & 201～ 300 & >300 \\
\cline{1-7}
空气质量 & 优 & 良 & 轻度污染 & 中度污染 & 重度污染 & 严重污染 \\
\cline{1-7}
\hline
\end{tabular}
\end{center}

下列叙述正确的是 \dotfill{(\kh{})}

\xx{}{}{}{}这 20 天中 AQI 的中位数略大于 150 B. 10 月 4 日到 10 月 11 日，空气质量越来越好

C. 这 20 天中的空气质量为优的天数占 25\% D. 10 月上旬 AQI 的极差大于中旬 AQI 的极差

\subsubsection*{三、解答题}

14.(2025$\cdot$上海春季卷$\cdot$19)甲、乙是两个体育社团的小组. 如图所示的是甲、乙两组组员身高分布的茎叶图 (单位:cm)，以身高的百位数和十位数作为“茎”排列在中间、个位数作为“叶”分列在两边.

(1)求甲、乙两组组员身高的第 60 百分位数；

(2)从甲、乙两组各选取一个组员，求两人身高均在 ${170}{{cm}}$ 以上的概率；

(3)为使两组人数相同，从甲组中调派一个队员到乙组. 是否存在甲组的一个组员, 将他调派至乙组后, 甲、乙两组的平均身高都增大?

\begin{center}
\includegraphics[width=0.4\textwidth]{bodyxb/src/219_972_536_501_252_0.jpg}
\end{center}


第 14 题图

15. (2024$\cdot$静安区二模$\cdot$18) 某高中随机抽取 100 名学生，测得他们的身高(单位:cm)，按照区间 $\lbrack {160},{165})$ 、 $\lbrack {165},{170})$ 、 $\lbrack {170},{175})$ 、 $\lbrack {175},{180})$ 、 $\left\lbrack  {{180},{185}}\right\rbrack$ 分组，得到样本身高的频率分布直方图 (如下图所示).

(1)求身高不低于 ${170}{{cm}}$ 的学生人数；

(2)将身高在 $\lbrack {170},{175})$ 、 $\lbrack {175},{180})$ 、 $\left\lbrack  {{180},{185}}\right\rbrack$ 区间内的学生依次记为 $A$ 、 $B$ 、 $C$ 三个组,用分层抽样的方法从三个组中抽取 6 人.

\begin{center}
\includegraphics[width=0.4\textwidth]{bodyxb/src/219_922_1443_551_441_0.jpg}
\end{center}


第 15 题图

 \ding{172}  求从这三个组分别抽取的学生人数；

 \ding{173}  若要从 6 名学生中抽取 2 人，求 B 组中至少有 1 人被抽中的概率.

考点3 解 样本的数字特征

\subsubsection*{一、填空题}

1.(2023$\cdot$上海秋季卷$\cdot$9)国内生产总值(GDP)是衡量该地区经济状况的最佳指标. 根据统计数据显示，某市在 2020 年间经济高质量增长，GDP 稳定增长，第一个季度和第四季度的 GDP 分别为 232 亿元、241 亿元，且四个季度的 GDP 逐季度增长，若这四个季度的 GDP 的中位数与平均数相等， 则该市 2020 年的 GDP 总额为\blkda[2]{}.

2.(2025$\cdot$奉贤区二模$\cdot$3)假设生产某产品的一个部件来自三个供应商，供货占比分别是 $\dfrac{1}{2}\text{ 、 }\dfrac{1}{6}\text{ 、 }\dfrac{1}{3}$ ， 而它们的良品率分别是 0.96、0.90、0.93，则该部件的总体良品率是\blkda[2]{}.

3.(2025$\cdot$宝山区一模$\cdot$6)某运动员在某次男子10m气手枪射击比赛中的得分数据(单位:环)为:9.6、 9.9、9.2、9.4、9.9、10.1、10.2、9.7、9.6、9.3、10.0、10.4，则这组数据的第 25 百分位数为\blkda[2]{}.

\begin{center}
\includegraphics[width=0.2\textwidth]{bodyxb/src/220_1210_865_219_152_0.jpg}
\end{center}


第 5 题图

4.(2024$\cdot$宝山区二模$\cdot$6)有一组数从小到大排列为:3、5、x、8、9、10. 若其极差与平均数相等，则这组数据的中位数为\blkda[2]{}.

5.(2024$\cdot$金山区一模$\cdot$6)已知甲、乙两组数据的茎叶图如图所示，若它们的中\blkda[2]{} $9m3248$ 位数相同，则图中 $m$ 的值为\blkda[2]{}.

6.(2024$\cdot$嘉定区二模$\cdot$7)数据1、2、、3、4、5 的方差为 ${s}_{1}^{2}$ ，数据3、6、9、12、15 的方差为 ${s}_{2}^{2}$ ，则 $\dfrac{{s}_{2}^{2}}{{s}_{1}^{2}} =$ \blkda[2]{}.

7.(2025$\cdot$闵行区二模$\cdot$7)已知数据 ${x}_{1}$ 、 ${x}_{2}$ 、 $\cdots$ 、 ${x}_{100}$ 的平均数为 2，方差为 5，则 ${x}_{1}^{2}$ 、 ${x}_{2}^{2}$ 、 $\cdots$ 、 ${x}_{100}^{2}$ 的平均数为\blkda[2]{}.

8.(2024$\cdot$青浦区一模$\cdot$7)某家大型超市统计了八次节假日的客流量(单位:百人)分别为29、30、38、 25、37、40、42、32，那么这组数据的第 75 百分位数为\blkda[2]{}.

9.(2025$\cdot$杨浦区一模$\cdot$8)某次杨浦区高三质量调研数学试卷中的填空题第八题，答对得 5 分，答错或不答得 0 分，全区共 4000 人参加调研，该题的答题正确率是 60\%，则该次调研中全区间学该题得分的方差为\blkda[2]{}.

10.(2025$\cdot$浦东新区二模$\cdot$9)李老师在整理建模小组 10 名学生的成绩时不小心遗失了一位学生的成绩,且剩余学生的成绩数据如下: $5\text{ 、 }6\text{ 、 }6\text{ 、 }7\text{ 、 }7\text{ 、 }7\text{ 、 }8\text{ 、 }9\text{ 、 }9$ ,但李老师记得这名学生的成绩恰好是本组学生成绩的第 25 百分位数，则这 10 名学生的成绩的方差为\blkda[2]{}.

11.(2025$\cdot$崇明区一模$\cdot$10)某校四个植树小队，在植树节这天种下柏树的棵数分别为10、x、10、8， 若这组数据的中位数和平均数相等，那么 $x =$ \blkda[2]{}.

\subsubsection*{二、选择题}

12.(2025$\cdot$普陀区二模$\cdot$13)某市职业技能大赛的移动机器人比赛项目有 19 名同学参赛，他们在预赛中所得的积分互不相同，只有积分在前 10 名的同学才能进入决赛. 若该比赛项目中的某同学知道自己的积分后，要判断自己能否进入决赛，则他只需要知道这 19 名同学的预赛积分的\blkda[2]{} \dotfill{(\kh{})}

\xx{}{}{}{}平均数 B. 众数 C. 中位数 D. 方差

13.(2024$\cdot$浦东新区一模$\cdot$14)某组样本数据由10个互不相同的数组成，去掉其中的最小数和最大数后, 得到一组新的样本数据, 则下列选项一定成立的是 \dotfill{(\kh{})}.

\xx{}{}{}{}两组样本数据的平均数相同 B. 两组样本数据的方差相同

C. 两组样本数据的中位数相同 D. 两组样本数据的极差相同

\begin{center}
\includegraphics[width=0.2\textwidth]{bodyxb/src/221_1161_388_295_253_0.jpg}
\end{center}


第 14 题图

14.(2024$\cdot$杨浦区一模$\cdot$14)在一次男子 ${10}{m}$ 气手枪射击比赛中，甲运动员的成绩(单位:环)为7、5、7、、、10.9；乙运动员的成绩为8、、3、、4、、、 10.1, 如茎叶图所示. 从这组数据来看, 下列说法正确的是 \dotfill{(\kh{})}

\xx{}{}{}{}甲的平均成绩和乙一样, 且甲更稳定

B. 甲的平均成绩和乙一样, 但乙更稳定

C. 甲的平均成绩高于乙，且甲更稳定

D. 乙的平均成绩高于甲，且乙更稳定

15.(2025$\cdot$闵行区二模$\cdot$14)某人统计了甲、乙两家零售商店在周一到周五的营业额(单位:百元)情况，得到了如下的茎叶图(其中茎表示十位数，叶表示个位数)，关于这 5 天的营业额情况，下列结论正确的是 \dotfill{(\kh{})}

\begin{center}

\begin{tabular}{|c|c|c|}
\hline
甲 & \phantom{X} & 乙 \\
\cline{1-3}
8 7 & 2 & 6 8 \\
\cline{1-3}
7 6 2 & 3 & 0 2 4 \\
\cline{1-3}
\hline
\end{tabular}
\end{center}

第 15 题图

\xx{}{}{}{}甲、乙两家商店营业额的极差相同

B. 甲、乙两家商店营业额的中位数相同

C. 从营业额超过 3000 元的天数所占比例来看, 甲商店较高

D. 甲商店营业额的方差小于乙商店营业额的方差

16.(2024$\cdot$崇明区二模$\cdot$14)某单位有 $A$ 、 $B$ 两个部门，1月份进行服务满意度问卷调查，得到两部门服务满意度得分的频率分布条形图如图所示. 设 $A$ 、 $B$ 两部门的服务满意度得分的第 75 百分位数分别为 ${n}_{1}$ 、 ${n}_{2}$ ， 方差分别为 ${s}_{1}^{2}\text{ 、 }{s}_{2}^{2}$ ,则下列说法正确的是 \dotfill{(\kh{})}

\begin{center}
\includegraphics[width=0.4\textwidth]{bodyxb/src/221_957_1134_500_284_0.jpg}
\end{center}


第 16 题图

\xx{}{}{}{}${n}_{1} > {n}_{2},{s}_{1}^{2} > {s}_{2}^{2}$

B. ${n}_{1} > {n}_{2},{s}_{1}^{2} < {s}_{2}^{2}$

C. ${n}_{1} < {n}_{2},{s}_{1}^{2} < {s}_{2}^{2}$

D. ${n}_{1} < {n}_{2},{s}_{1}^{2} > {s}_{2}^{2}$

17.(2024$\cdot$嘉定区一模$\cdot$14)两位跳水运动员甲和乙，某次比赛中的得分如下表所示，则正确的选项为 \dotfill{(\kh{})}

\begin{center}

\begin{tabular}{|c|c|c|c|c|c|}
\hline
\phantom{X} & 第一跳 & 第二跳 & 第三跳 & 第四跳 & 第五跳 \\
\cline{1-6}
甲 & 85.5 & 96 & 86.4 & 75.9 & 94.4 \\
\cline{1-6}
乙 & 79.5 & 80 & 95.7 & 94.05 & 86.4 \\
\cline{1-6}
\hline
\end{tabular}
\end{center}

\xx{}{}{}{}甲和乙的中位数相等, 甲的均分小于乙

B. 甲的均分大于乙, 甲的方差大于乙

C. 甲的均分大于乙, 甲的方差等于乙

D. 甲的均分大于乙，甲的方差小于乙

18.(2024$\cdot$徐汇区一模$\cdot$14)跳水比赛共有 7 位评委分别给出某选手的原始评分，评定该选手的成绩时, 从 7 个原始评分中去掉 1 个最高分和 1 个最低分, 得到 5 个有效评分. 5 个有效评分与 7 个原始评分相比，一定不变的数字特征是 \dotfill{(\kh{})}

\xx{}{}{}{}中位数 B. 平均数 C. 方差 D. 极差

19.(2024$\cdot$宝山区一模$\cdot$14)下列说法中错误的是 \dotfill{(\kh{})}

\xx{}{}{}{}一组数据的平均数、中位数可能相同

B. 一组数据中比中位数大的数和比中位数小的数一样多

C. 平均数、众数和中位数都是描述一组数据的集中趋势的统计量

D. 极差、方差、标准差都是描述一组数据的离散程度的统计量

\begin{center}
\includegraphics[width=0.2\textwidth]{bodyxb/src/222_1184_617_208_235_0.jpg}
\end{center}


第 20 题图

20. (2024$\cdot$松江区一模$\cdot$14)如图所示的茎叶图记录了甲、乙两支篮球队各 6 名队员某场比赛的得分数据(单位:分)，则下列说法中正确的是 \dotfill{(\kh{})}

\xx{}{}{}{}甲队数据的中位数大于乙队数据的中位数

B. 甲队数据的平均值小于乙队数据的平均值

C. 甲队数据的标准差大于乙队数据的标准差

D. 乙队数据的第 75 百分位数为 27

21.(2025$\cdot$浦东新区一模$\cdot$15)对一组数据3、3、3、1、1、5、5、2、4，若任意去掉其中一个数据，剩余数据的统计量一定会发生变化的为 \dotfill{(\kh{})}

\xx{}{}{}{}中位数 B. 众数 C. 平均数 D. 方差

22.(2024$\cdot$普陀区一模$\cdot$15)已知一组数据3、1、5、3、2，现加入 $p$ 、 $q$ 两数对该组数据进行处理，若经过处理后的这组数据的极差为 $p - q$ ,则经过处理后的这组数据与之前的那组数据相比,一定会变大的数字特征是 \dotfill{(\kh{})}

\xx{}{}{}{}平均数 B. 方差 C. 众数 D. 中位数

23.(2025$\cdot$宝山区二模$\cdot$15)甲、乙两名篮球运动员在 8 场比赛中的单场得分用茎叶图表示如下左图, 茎叶图中甲的得分有部分数据丢失, 但甲得分的折线图完好 (下右图), 则下列结论正确的是 \dotfill{(\kh{})}

\begin{center}
\includegraphics[width=0.7\textwidth]{bodyxb/src/222_323_1467_880_326_0.jpg}
\end{center}


第 23 题图

\xx{}{}{}{}甲得分的极差小于乙得分的极差

B. 甲得分第 25 百分位数大于乙得分的第 75 百分位数

C. 甲得分的平均数大于乙得分的平均数

D. 甲得分的方差小于乙得分的方差

\subsubsection*{三、解答题}

24.(2025$\cdot$静安区二模$\cdot$18)某校高三学生体检后，为了解高三学生的视力情况，该校从高三六个班的 300 名学生中以班为单位(每班学生 50 人)，每班按随机抽样方法抽取了 8 名学生的视力数据. 其中高三(1)班抽取的 8 名学生的视力数据与人数见下表:

\begin{center}

\begin{tabular}{|c|c|c|c|c|c|c|c|c|c|c|c|c|c|c|}
\hline
视力数据 & 4.0 & 4.1 & 4.2 & 4.3 & 4.4 & 4.5 & 4.6 & 4.7 & 4.8 & 4.9 & 5.0 & 5.1 & 5.2 & 5.3 \\
\cline{1-15}
人数 & \phantom{X} & \phantom{X} & \phantom{X} & \phantom{X} & 2 & \phantom{X} & 2 & \phantom{X} & 2 & 1 & \phantom{X} & 1 & \phantom{X} & \phantom{X} \\
\cline{1-15}
\hline
\end{tabular}
\end{center}

(1)用上述样本数据估计高三(1)班学生视力的平均值；

(2)已知其余五个班学生视力的平均值分别为4.3,4.4,4.5,4.6,4.8. 若从这六个班中任意抽取两个班学生视力的平均值作比较，求抽取的两个班学生视力的平均值之差的绝对值不小于 0.2 的概率.

25.(2024$\cdot$上海春季卷$\cdot$19)从某果园中采摘某种水果共 136 箱，每箱均装有相同个数的此种水果， 此种水果分为一级果和二级果, 其中一级果 102 箱, 二级果 34 箱.

(1)从中随机挑出两箱水果，求恰好选到一级果、二级果各一箱的概率；

(2)若采用分层随机抽样的方法从中抽出 8 箱水果，求一级果、二级果各几箱；

(3)如果抽出若干箱水果，其中，一级果120个，单果质量平均数为303.45 g，方差为603.46；二级果 48 个，单果质量平均数为 240.41 g，方差为 648.21；求 168 个此种水果的单果质量平均数与方差，并预估果园中水果的单果质量. (结果保留两位小数)

26.(2025$\cdot$金山区一模$\cdot$19)某高中举行了一次知识竞赛. 为了了解本次竞赛成绩情况，从中抽取了部分学生的成绩作为样本进行统计. 将成绩进行整理后, 依次分为五组: $\lbrack {50},{60})$ 、[60,70)、 $\lbrack {70},{80})$ 、[80,90)、[90,100)，其中第 1 组的频率为第 2 组和第 4 组频率的等比中项. 请根据下面的频率分布直方图如图所示)解决下列问题:

(1)求a、b的值；

(2)从样本数据在 $\lbrack {50},{60})$ 、 $\lbrack {70},{80})$ 两个小组内的学生中，用分层抽样的方法抽取 7 名学生，再从这 7 名学生中随机选出 2 人，求选出的两人恰好来自不同小组的概率；

(3)某老师在此次竞赛成绩中抽取了 10 名学生的分数: ${x}_{1}\text{ 、 }{x}_{2}\text{ 、 }{x}_{3}\text{ 、 }\cdots \text{ 、 }{x}_{10}$ ,已知这 10 个分数的平均数 $\overline{x} = {88}$ ,方差 ${s}^{2} = {25}$ ,若剔除其中的 95 和 81 两个分数，求剩余 8 个分数的平均数与方差.

\begin{center}
\includegraphics[width=0.4\textwidth]{bodyxb/src/224_886_542_543_396_0.jpg}
\end{center}


第 26 题图

27.(2025$\cdot$青浦区一模$\cdot$19)第七届中国国际进口博览会于2024 年 11 月 5 日至 10 日在上海举办，某公司生产的 $A$ 、 $B$ 、 $C$ 三款产品在博览会上亮相，每一种产品均有普通装和精品装两种款式，该公司每天产量如下表:(单位:个)

\begin{center}

\begin{tabular}{|c|c|c|c|}
\hline
\phantom{X} & 产品 $A$ & 产品 $B$ & 产品 $C$ \\
\cline{1-4}
普通装 & $n$ & 180 & 400 \\
\cline{1-4}
精品装 & 300 & 420 & 600 \\
\cline{1-4}
\hline
\end{tabular}
\end{center}

现采用分层抽样的方法在某一天生产的产品中抽取 100 个,其中 $B$ 款产品有 30 个.

(1)求 $n$ 的值；

(2)用分层抽样的方法在 $C$ 款产品中抽取一个容量为 5 的样本, 从样本中任取 2 个产品, 求其中至少有一个精品装产品的概率;

(3)对抽取到的 $B$ 款产品样本中某种指标进行统计，普通装产品的平均数为 10，方差为 2，精品装产品的平均数为 12 , 方差为 1.8 , 试估计这天生产的 $B$ 款产品的某种指标的总体方差. (精确到 0.01)

28.(2025$\cdot$崇明区一模$\cdot$19)王老师将全班 40 名学生的高一数学期中考试(满分 100 分)成绩分成 5 组,绘制成如图所示的频率分布直方图,现将 $\lbrack {50},{60})$ 记作第一组, $\lbrack {60},{70})\text{ 、 }\lbrack {70},{80})\text{ 、 }\lbrack {80},{90})$ 、 [90,100]分别记作第二、三、四、五组. 已知第一组、第二组的频率之和为 0.3,第一组和第五组的频率相同.

(1)估计此次考试成绩的平均值;(同一组数据用该组数据的中点值代替)

(2)王老师将测试成绩在 $\lbrack {80},{90})$ 和 $\left\lbrack  {{90},{100}}\right\rbrack$ 内的试卷进行分析，再从中选 2 人的试卷进行优秀答卷展示，求被选中进行优秀答卷展示的这 2 人的测试成绩至少 1 个在 $\left\lbrack  {{90},{100}}\right\rbrack$ 内的概率；

(3) 已知第二组考生成绩的平均数和方差分别为 65 和 40，第四组考生成绩的平均数和方差分别为 83 和 70，据此计算第二组和第四组所有学生成绩的方差.

\begin{center}
\includegraphics[width=0.4\textwidth]{bodyxb/src/225_895_540_545_439_0.jpg}
\end{center}


第 28 题图

29.(2025$\cdot$黄浦区一模$\cdot$19)A 校高一年级共有学生 330 名，为了解该校高一年级学生的身高情况， 学校采用分层随机抽样的方法抽取 66 名学生，其中女生 32 名，男生 34 名，测量他们的身高.

(1)该校高一学生中男、女生各有多少名？

(2)若从这 66 名学生中随机抽取两名，求这两名都是男生的概率；

(3)在 32 名女生身高的数据中，其中一个数据记录有误，错将 165 cm 记录为 156 cm，由错误数据求得这 32 个数据的平均数为 ${161}{{cm}}$ ，方差为 23.6875 ，求原始数据的平均数及方差.(平均数结果保留精确值, 方差结果精确到 0.01)

30.(2025$\cdot$虹口区一模$\cdot$19)2024 年巴黎奥运会落下帷幕. 某平台为了解观众对本次奥运会的满意度, 随机调查了本市 1000 名观众, 得到他们对本届奥运会的满意度评分(满分 100 分), 平台将评分分为 $\lbrack {50},{60})$ 、 $\lbrack {60},{70})$ 、 $\lbrack {70},{80})$ 、 $\left\lbrack  {{80},{90})}\right\rbrack$ 、 $\left\lbrack  {{90},{100}}\right\rbrack$ 共 5 层，绘制成频率分布直方图(如图 1 所示). 并在这些评分中以分层抽样的方式从这 5 层中再抽取了共 20 名观众的评分, 绘制成茎叶图，但由于某种原因茎叶图受到了污损，可见部分信息，如图 2 所示.

\begin{center}
\includegraphics[width=0.6\textwidth]{bodyxb/src/226_431_450_784_480_0.jpg}
\end{center}


第 30 题图

(1)求图 2 中这 20 名观众的满意度评分的第 35 百分位数；

(2)若从图 2 中的 20 名观众中再任选取 3 人做深度采访，求其中至少有 1 名观众的评分大于等于 90 分的概率;

(3)已知这 1000 名观众的评分位于 $\left\lbrack  {{50},{80}}\right)$ 上的均值为 67，方差为 64.7，位于 $\left\lbrack  {{50},{100}}\right\rbrack$ 上的均值为 73 , 方差为 134.6 , 求这 1000 名观众的评分位于[80,100]上的均值与方差.

\subsection*{考点中 1 )变量间的相关关系}

\subsubsection*{一、填空题}

1.(2025$\cdot$徐汇区二模$\cdot$5)如下是一个 $2 \times  2$ 列联表，则 $s =$ \blkda[2]{}.

\begin{center}

\begin{tabular}{|c|c|c|c|}
\hline
\phantom{X} & ${y}_{1}$ & ${y}_{2}$ & 总计 \\
\cline{1-4}
${x}_{1}$ & $a$ & 35 & 45 \\
\cline{1-4}
${x}_{2}$ & 7 & $b$ & $n$ \\
\cline{1-4}
总计 & $m$ & 73 & $S$ \\
\cline{1-4}
\hline
\end{tabular}
\end{center}

2.(2024$\cdot$长宁区二模$\cdot$7)收集数据，利用 $2 \times  2$ 列联表,分析学习成绩好与上课注意力集中是否有关时，提出的零假设为:学习成绩好与上课注意力集中\blkda[2]{}(填:“有关”或“无关”)

3.(2024$\cdot$普陀区二模$\cdot$7)为了提高学生参加体育锻炼的积极性，某校本学期依据学生特点，有针对性地组建了五个特色运动社团. 学校为了了解学生参与运动的情况，对每个特色运动社团的参与人数进行了统计，其中一个特色运动社团开学第 1 周至第 5 周参与运动的人数统计数据如表所示.

\begin{center}

\begin{tabular}{|c|c|c|c|c|c|}
\hline
周次 $x$ & 1 & 2 & 3 & 4 & 5 \\
\cline{1-6}
参与运动的人数 $y$ & 35 & 36 & 40 & 39 & 45 \\
\cline{1-6}
\hline
\end{tabular}
\end{center}

若表中数据可用回归方程 $y = {2.3x} + \widehat{b}\left( {1 \leq  x \leq  {18},x \in  \mathbf{N}}\right)$ 来预测，则本学期第 11 周参与该特色运动社团的人数约为\blkda[2]{}. (精确到整数)

4.(2025$\cdot$杨浦区二模$\cdot$9)植物社团的同学观察一株植物的生长情况，为了解植物高度 $y$ (单位: ${{cm}}$ ) 与生长期 $x$ (单位: 天)之间的关系，随机统计了某 4 天的植物高度，并制作了如右对照表:

\begin{center}

\begin{tabular}{|c|c|c|c|c|}
\hline
生长期 $x$ & 3 & 9 & 11 & 17 \\
\cline{1-5}
植物高度 $y$ & 2.4 & 3.4 & 3.8 & 5.2 \\
\cline{1-5}
\hline
\end{tabular}
\end{center}

由表中数据可得回归方程 $y = \widehat{a}x + \widehat{b}$ 中 $\widehat{a} = 0$ 。2,试预测生长期是 30 天时,植物高度约为\blkda[2]{} ${{cm}}.\left( {\widehat{a} = \dfrac{\mathop{\sum }\limits_{{i = 1}}^{n}{x}_{i}{y}_{i} - n\overline{x}\overline{y}}{\mathop{\sum }\limits_{{i = 1}}^{n}{x}_{i}^{2} - n{\overline{x}}^{2}},\widehat{b} = \overline{y} - \widehat{a}\overline{x}}\right)$

5.(2024$\cdot$宝山区二模$\cdot$9)某公司为了了解某商品的月销售量 $y$ (单位:万件)与月销售单价 $x$ (单位: 元/件)之间的关系, 随机统计了 5 个月的销售量与销售单价, 并制作了如下对照表:

\begin{center}

\begin{tabular}{|c|c|c|c|c|c|}
\hline
月销售单价 $x$ (元/件) & 10 & 15 & 20 & 25 & 30 \\
\cline{1-6}
月销售量 $y$ (万件) & 11 & 10 & 8 & 6 & 5 \\
\cline{1-6}
\hline
\end{tabular}
\end{center}

由表中数据可得回归方程 $y = \widehat{a}x + \widehat{b}$ 中 $\widehat{a} =  - {0.32}$ ,试预测当月销售单价为 40 元/件时,月销售量为\blkda[2]{}万件.

6.(2024$\cdot$宝山区一模$\cdot$10)随着我国国民教育水平的提高，越来越多的有志青年报考研究生. 现阶段,我国研究生入学考试科目为思政、外语和专业课三门，录取工作将这样进行:在每门课均及格 (60 分) 的考生中, 按总分进行排序, 择优录取. 振华同学刚刚完成报名, 尚有 11 周复习时间, 下表是他每门课的复习时间和预计得分. 设思政、外语和专业课分配到的周数分别为 $x\text{ 、 }y\text{ 、 }z$ ,则自然数数组 $\left( {x,y,z}\right)  =$ \blkda[2]{}时，振华被录取的可能性最大.

\begin{center}

\begin{tabular}{|c|c|c|c|c|c|c|c|c|c|c|c|}
\hline
\multirow{2}{*}{科目} & \multicolumn{11}{|c|}{周数} \\
\cline{2-12}
 & 0 & 1 & 2 & 3 & 4 & 5 & 6 & 7 & 8 & 9 & 10 \\
\cline{1-12}
思政 & 20 & 40 & 55 & 65 & 72 & 78 & 80 & 82 & 83 & 84 & 85 \\
\cline{1-12}
外语 & 30 & 45 & 53 & 58 & 62 & 65 & 68 & 70 & 72 & 74 & 75 \\
\cline{1-12}
专业课 & 50 & 70 & 85 & 90 & 93 & 95 & 96 & 96 & 96 & 96 & 96 \\
\cline{1-12}
\hline
\end{tabular}
\end{center}

7.(2025$\cdot$奉贤区二模$\cdot$9)通过随机抽样，获得某种商品消费者年需求量与该商品每千克价格之间的一组数据调查，如下表所示:

\begin{center}

\begin{tabular}{|c|c|c|c|c|c|c|c|c|c|c|}
\hline
价格(百元) & ${x}_{1}$ & ${x}_{2}$ & ${x}_{3}$ & ${x}_{4}$ & ${x}_{5}$ & ${x}_{6}$ & ${x}_{7}$ & ${x}_{8}$ & ${x}_{9}$ & ${x}_{10}$ \\
\cline{1-11}
\phantom{X} & 4 & 4 & 4.6 & 5 & 5.2 & 5.6 & 6 & 6.6 & 7 & 10 \\
\cline{1-11}
需求量(kg) & ${y}_{1}$ & ${y}_{2}$ & ${y}_{3}$ & ${y}_{4}$ & ${y}_{5}$ & ${y}_{6}$ & ${y}_{7}$ & ${y}_{8}$ & ${y}_{9}$ & ${y}_{10}$ \\
\cline{1-11}
\phantom{X} & 3.5 & 3 & 2.7 & 2.4 & 2.5 & 2 & 1.5 & 1.2 & 1.2 & 1 \\
\cline{1-11}
\hline
\end{tabular}
\end{center}

那么线性相关系数 $r =$ \blkda[2]{}. (精确到 0.001)

线性相关系数公式: $r = \dfrac{\mathop{\sum }\limits_{{i = 1}}^{n}\left( {{x}_{i} - \overline{x}}\right) \left( {{y}_{i} - \overline{y}}\right) }{\sqrt{\mathop{\sum }\limits_{{i = 1}}^{n}{\left( {x}_{i} - \overline{x}\right) }^{2} \cdot  \mathop{\sum }\limits_{{i = 1}}^{n}{\left( {y}_{i} - \overline{y}\right) }^{2}}}$.

\subsubsection*{二、选择题}

8.(2025$\cdot$黄浦区二模$\cdot$13)如果两种证券在一段时间内收益数据的相关系数为 0.8，那么表明\dotfill{(\kh{})}

\xx{}{}{}{}两种证券的收益有反向变动的倾向

B. 两种证券的收益有同向变动的倾向

C. 两种证券的收益之间存在完全反向的联动关系, 即涨或跌是相反的

D. 两种证券的收益之间存在完全同向的联动关系, 即同时涨或同时跌

9.(2025$\cdot$闵行区二模$\cdot$13)两个变量: $x$ 与 $y$ 之间的回归方程 \dotfill{(\kh{})}

\xx{}{}{}{}表示 $x$ 与 $y$ 之间的函数关系 B. 表示 $x$ 与 $y$ 之间的不确定关系

C. 反映 $x$ 与 $y$ 之间的真实关系 D. 是反映 $x$ 与 $y$ 之间的真实关系的一种最佳拟合

10.(2024$\cdot$上海秋季卷$\cdot$13)已知气候温度和海水表层温度相关，且相关系数为正数，对此描述正确的是 \dotfill{(\kh{})}

\xx{}{}{}{}气候温度高，海水表层温度就高

B. 气候温度高，海水表层温度就低

C. 随着气候温度由低到高，海水表层温度呈上升趋势

D. 随着气候温度由低到高，海水表层温度呈下降趋势

11.(2024$\cdot$奉贤区二模$\cdot$13)设某大学的女生体重 $y$ (单位: ${{kg}}$ )与身高 $x$ (单位: ${{cm}}$ )具有线性相关关系，根据一组样本数据 $\left( {{x}_{i},{y}_{i}}\right) \left( {i = 1,2,\cdots ,n}\right)$ ，用最小二乘法建立的回归方程为 $y = {0.85x} - {85.71}$ ， 则下列结论中不正确的是 \dotfill{(\kh{})}

\xx{}{}{}{}$y$ 与 $x$ 具有正的线性相关关系

B. 回归直线过样本点的中心 $\left( {\overline{x},\overline{y}}\right)$

C. 若该大学某女生身高增加 $1{{cm}}$ ，则其体重约增加0.85 ${{kg}}$

D. 若该大学某女生身高为 ${170}{{cm}}$ ，则可断定其体重必为 ${58.79}{{kg}}$

12.(2023$\cdot$上海秋季卷$\cdot$14)为了解某校高中男生身高与体重的关系，随机选取 50 名男生的身高与

\begin{center}
\includegraphics[width=0.4\textwidth]{bodyxb/src/229_880_619_551_192_0.jpg}
\end{center}


第 12 题图

体重的数据绘制散点图，如图所示，则下列说法正确的是 \dotfill{(\kh{})}

\xx{}{}{}{}身高越高，体重越重

B. 身高越高，体重越轻

C. 身高与体重呈正相关

D. 身高与体重呈负相关

13.(2024$\cdot$徐汇区二模$\cdot$14)为了研究 $y$ 关于 $x$ 的线性相关关系，收集了 5 组样本数据(见下表):

\begin{center}

\begin{tabular}{|c|c|c|c|c|c|}
\hline
$x$ & 1 & 2 & 3 & 4 & 5 \\
\cline{1-6}
$y$ & 0.5 & 0.9 & 1 & 1.1 & 1.5 \\
\cline{1-6}
\hline
\end{tabular}
\end{center}

若已求得一元线性回归方程为 $y = {\text{ â }x} + {0.34}$ ，则下列选项中正确的是 \dotfill{(\kh{})}

\xx{}{}{}{}$\widehat{a} = {0.21}$

B. 当 $x = 8$ 时， $y$ 的预测值为 2.2

C. 样本数据 $y$ 的第 40 百分位数为 1

D. 去掉样本点 (3,1) 后, $x$ 与 $y$ 的样本相关系数 $r$ 不会改变

14.(2024$\cdot$松江区二模$\cdot$14)某小区为了倡导居民对生活垃圾进行分类，对垃圾分类后处理垃圾 $x$ (kg)所需的费用 $y$ (角)的情况作了调研，并统计得到下表中几组对应数据，同时用最小二乘法得到 $y$ 关于 $x$ 的线性回归方程为 $y = {0.7x} + {0.4}$ ,则下列说法错误的是 \dotfill{(\kh{})}

\begin{center}

\begin{tabular}{|c|c|c|c|c|}
\hline
$x$ & 2 & 3 & 4 & 5 \\
\cline{1-5}
$y$ & 2 & 2.3 & 3.4 & $m$ \\
\cline{1-5}
\hline
\end{tabular}
\end{center}

\xx{}{}{}{}变量 $x\text{ 、 }y$ 之间呈正相关关系

B. 可以预测当 $x = 8$ 时， $y$ 的值为 6

C. $m = {3.9}$

D. 由表格中数据知样本中心点为 $\left( {{3.5},{2.85}}\right)$

15.(2025$\cdot$徐汇区二模$\cdot$14)在研究线性回归模型时，若样本数据( \ding{55} ) 点都在直线 $y =  - \dfrac{1}{3}x + 2$ 上，则两组数据 ${x}_{i}$ 和 ${y}_{i}\left( {i = 1,2,3,\cdots ,n}\right)$ 的线性相关系数为\dotfill{(\kh{})}

\xx{}{}{}{}-1 B. 1

C. $- \dfrac{1}{3}$ D. 2

16.(2024$\cdot$浦东新区二模$\cdot$15)通过随机抽样，我们绘制了如图所示的某种商品每千克价格(单位: 百元)与该商品消费者年需求量(单位: ${{kg}}$ )的散点图. 若去掉图中右下方的点 $A$ 后，下列说法正确的是 \dotfill{(\kh{})}

\begin{center}
\includegraphics[width=0.5\textwidth]{bodyxb/src/230_813_322_609_306_0.jpg}
\end{center}


第 16 题图

\xx{}{}{}{}“每千克价格”与“年需求量”这两个变量由负相关变为正相关

B. “每千克价格”与“年需求量”这两个变量的线性相关程度不变

C. “每千克价格”与“年需求量”这两个变量的线性相关系数变大

D. “每千克价格”与“年需求量”这两个变量的线性相关系数变小

17.(2025$\cdot$浦东新区二模$\cdot$15)研究变量 $x$ 、 $y$ 得到一组成对数据( \ding{55} ； 线性回归分析,接着增加一个数据 $\left( {{x}_{n + 1},{y}_{n + 1}}\right)$ ,其中 ${x}_{n + 1} = \dfrac{1}{n}\mathop{\sum }\limits_{{i = 1}}^{n}{x}_{i},{y}_{n + 1} = \dfrac{1}{n}\mathop{\sum }\limits_{{i = 1}}^{n}{y}_{i}$ ，再重新进行一次线性回归分析, 则下列说法正确的是 \dotfill{(\kh{})}

\xx{}{}{}{}变量 $x$ 与变量 $y$ 的相关性变强 B. 相关系数 $r$ 的绝对值变小

C. 线性回归方程 $y = \widehat{a}x + \widehat{b}$ 不变 D. 拟合误差 $Q$ 变大

\subsubsection*{三、解答题}

18.(2025$\cdot$上海秋季卷$\cdot$17)2024 年东京奥运会,中国获得了男子4×100 m混合泳接力金牌. 以下是历届奥运会男子 4×100 m 混合泳接力项目冠军成绩记录(单位:s)，数据按照升序排列.

\begin{center}

\begin{tabular}{|c|c|c|c|c|}
\hline
206.78 & 207.46 & 207.95 & 209.34 & 209.35 \\
\cline{1-5}
210.68 & 213.73 & 214.84 & 216.93 & 216.93 \\
\cline{1-5}
\hline
\end{tabular}
\end{center}

(1)求这组数据的极差与中位数；

(2)从这 10 个数据中任选 3 个，求恰有 2 个数据在 211 以上的概率；

(3)若比赛成绩 $y$ 关于年份 $x$ 的回归方程为 $y =  - {0.311x} + \widehat{b}$ ，年份 $x$ 的平均数为 2006，预测 2028 年冠军队的成绩. (精确到 0.01 s).

\subsection*{考点5 > 互斥事件的概率计算}

\subsubsection*{一、填空题}

1.(2024$\cdot$浦东新区一模$\cdot$3)已知事件 $A$ 与事件 $B$ 互斥，且 $P\left( A\right)  = {0.3}$ ， $P\left( B\right)  = {0.4}$ ，则 $P\left( {A \cup  B}\right)  =$ \blkda[2]{}.

2. (2025$\cdot$普陀区二模$\cdot$3)已知事件 $A$ 与事件 $B$ 相互独立，若 $P\left( {A \cap  \overline{B}}\right)  = \dfrac{3}{25}$ ， $P\left( B\right)  = \dfrac{3}{5}$ ，则 $P\left( A\right)  =$ \blkda[2]{}.

3.(2024$\cdot$青浦区一模$\cdot$4)2023 年 10 月 25 日至 11 月 12 日，青浦曲水园推出以“曲水流觞$\cdot$花趣水乡”为主题的菊花展. 花展结束后，园方挑选数百盆菊花免费赠送给市民. 其中有红色、黄色、橙色菊花各 1 盆，分别赠送给甲、乙、丙 3 人，每人 1 盆，则甲没有拿到橙色菊花的概率是\blkda[2]{}.

4.(2023$\cdot$上海春季卷$\cdot$5)已知事件 $A$ 发生的概率为 $P\left( A\right)  = {0.5}$ ，则它的对立事件 $P\left( \overline{A}\right)  =$ \blkda[2]{}.

5.(2024$\cdot$嘉定区一模$\cdot$6)已知事件 $A$ 和 $B$ 相互独立， $P\left( A\right)  = \dfrac{1}{4},P\left( B\right)  = \dfrac{1}{13}$ ，则 $P\left( {A \cap  B}\right)  =$ \blkda[2]{}.

6.(2025$\cdot$奉贤区一模$\cdot$9)A、B 两人下棋，每局两人获胜的可能性一样. 某一天两人要进行一场三局两胜的比赛,最终胜者赢得 100 元奖金. 第一局比赛 $A$ 胜,后因为有其他要事中止比赛. 为求公平，则 $A$ 应该分得\blkda[2]{}元奖金.

7.(2024$\cdot$崇明区一模$\cdot$9)已知事件 $A$ 与事件 $B$ 相互独立，如果 $P\left( A\right)  = {0.4}$ ， $P\left( B\right)  = {0.7}$ ，则 $P\left( {\overline{A} \cap  B}\right)  =$ \blkda[2]{}.

8.(2024$\cdot$崇明区二模$\cdot$10)某学习小组共有 10 名学生，其中至少有 2 名学生在同一月份出生的概率是\blkda[2]{}. (默认每月天数相同,结果精确到 0.001)

\subsubsection*{二、选择题}

9.(2025$\cdot$上海秋季卷$\cdot$13)已知事件 $A$ 、 $B$ 相互独立，事件 $A$ 发生的概率为 $P\left( A\right)  = \dfrac{1}{2}$ ，事件 $B$ 发生的概率为 $P\left( B\right)  = \dfrac{1}{2}$ ,则事件 $A \cap  B$ 发生的概率 $P\left( {A \cap  B}\right)$ 为 \dotfill{(\kh{})}

\xx{}{}{}{}$\dfrac{1}{8}$ B. $\dfrac{1}{4}$ C. $\dfrac{1}{2}$ D. 0

10.(2025$\cdot$徐汇区二模$\cdot$13)已知 $A$ 、 $B$ 为两个随机事件，则“ $A$ 、 $B$ 为互斥事件”是“$A$ 、 $B$ 为对立事件” 的 \dotfill{(\kh{})}

\xx{}{}{}{}充分非必要条件 B. 必要非充分条件

C. 充要条件 D. 非充分非必要条件

11.(2025$\cdot$黄浦区一模$\cdot$13)掷一颗质地均匀的骰子，观察朝上面的点数. 设事件 $E$ :点数是奇数. 事件 $F$ : 点数是偶数. 事件 $G$ : 点数是 3 的倍数. 事件 $H$ : 点数是 4. 下列每对事件中,不是互斥事件的为\dotfill{(\kh{})}

\xx{}{}{}{}$E$ 与 $F$ B. $F$ 与 $G$ C. $E$ 与 $H$ D. $G$ 与 $H$

12.(2024$\cdot$金山区一模$\cdot$14)已知事件 $A$ 和 $B$ 相互独立，且 $P\left( A\right)  = \dfrac{1}{3}$ ， $P\left( B\right)  = \dfrac{3}{7}$ ，则 $P\left( {A \cap  B}\right)$ 等于 \dotfill{(\kh{})}

\xx{}{}{}{}$\dfrac{1}{7}$ B. $\dfrac{2}{21}$ C. $\dfrac{2}{7}$ D. $\dfrac{16}{21}$

13.(2025$\cdot$虹口区一模$\cdot$14)已知事件 $A$ 和事件 $B$ 满足 $A \cap  B = \varnothing$ ，则下列说法正确的是\blkda[2]{}C. 当(C)

\xx{}{}{}{}事件 $\overline{A}$ 和事件 $\overline{B}$ 互斥 B. 事件 $A$ 和事件 $B$ 互斥

C. 事件 $A$ 和事件 $B$ 对立 D. 事件 $\overline{A}$ 和事件 $B$ 对立

14.(2024$\cdot$上海春季卷$\cdot$15)有四个礼品盒. 已知前三个礼品盒中分别只装了一支钢笔、一本书以及一个笔袋，第四个盒子中钢笔、书、笔袋都有. 现随机抽取一个盒子，事件 $A$ 为抽中的盒子里面有钢笔， 事件 $B$ 为抽中的盒子里面有书,事件 $C$ 为抽中的盒子里面有笔袋. 则下面正确的选项是\dotfill{(\kh{})}

\xx{}{}{}{}$A$ 与 $B$ 互斥 B. $A$ 与 $B$ 相互独立 C. $A$ 与 $B \cup  C$ 互斥 D. $A$ 与 $B \cap  C$ 独立

15.(2025$\cdot$崇明区一模$\cdot$15)抛掷一红一绿两颗质地均匀的骰子，记录骰子朝上面的点数，若用 $x$ 表示红色骰子的点数，用 $y$ 表示绿色骰子的点数，用 $\left( {x,y}\right)$ 表示一次试验结果，设事件 $E : x + y =$ 8;事件 $F$ : 至少有一颗点数为6;事件 $G : x > 4$ ;事件 $H : y < 4$. 则下列说法正确的是 \dotfill{(\kh{})}

\xx{}{}{}{}事件 $E$ 与事件 $F$ 为互斥事件 B. 事件 $F$ 与事件 $G$ 为互斥事件

C. 事件 $E$ 与事件 $G$ 相互独立 D. 事件 $G$ 与事件 $H$ 相互独立

16.(2024$\cdot$奉贤区二模$\cdot$15)有 6 个相同的球，分别标有数字1、2、3、4、5、6，从中有放回地随机取两次, 每次取 1 个球. 甲表示事件“第一次取出的球的数字是 1 ”, 乙表示事件“第二次取出的球的数字是 2 ”，丙表示事件“两次取出的球的数字之和是 5 ”，丁表示事件“两次取出的球的数字之和是 6”, 则下列选项正确的是 \dotfill{(\kh{})}

\xx{}{}{}{}甲与乙相互独立 B. 乙与丙相互独立 C. 甲与丙相互独立 D. 乙与丁相互独立

17.(2025$\cdot$松江区一模$\cdot$15)抛掷三枚硬币，若记“出现三个正面”“两个正面一个反面”和“两个反面一个正面”分别为事件 $A$ 、 $B$ 和 $C$ ，则下列说法错误的是 \dotfill{(\kh{})}

\xx{}{}{}{}事件 $A\text{ 、 }B$ 和 $C$ 两两互斥 B. $P\left( A\right)  + P\left( B\right)  + P\left( C\right)  = \dfrac{7}{8}$

C. 事件 $A$ 与事件 $B \cup  C$ 是对立事件 D. 事件 $A \cup  B$ 与 $B \cup  C$ 相互独立

18.(2025$\cdot$嘉定区一模$\cdot$15)假定生男生女是等可能的，设事件A:一个家庭中既有男孩又有女孩； 事件 $B$ : 一个家庭中最多有一个女孩. 针对下列两种情形: \ding{172}  家庭中有 2 个小孩； \ding{173}  家庭中有 3 个小孩, 下面说法正确是 \dotfill{(\kh{})}

\xx{}{}{}{} \ding{172} 中事件 $A$ 与事件 $B$ 相互独立， \ding{173} 中的事件 $A$ 与事件 $B$ 相互独立

B.  \ding{172} 中事件 $A$ 与事件 $B$ 不相互独立， \ding{173} 中的事件 $A$ 与事件 $B$ 相互独立

C.  \ding{172} 中事件 $A$ 与事件 $B$ 相互独立， \ding{173} 中的事件 $A$ 与事件 $B$ 不相互独立

D.  \ding{172} 中事件 $A$ 与事件 $B$ 不相互独立， \ding{173} 中的事件 $A$ 与事件 $B$ 不相互独立

19.(2025$\cdot$青浦区二模$\cdot$15)一个质地均匀的正四面体，四个面上分别标有数字1、2、3、4. 任意掷一次该四面体，观察它与地面接触面上的数字，得到样本空间 $\Omega  = \{ 1,2,3,4\}$ ，记事件 $A = \{ 1,2\}$ ，事件 $B = \{ 1,3\}$ ,事件 $C = \{ 1,4\}$ ,则 \dotfill{(\kh{})}

\xx{}{}{}{}事件 $A\text{ 、 }B\text{ 、 }C$ 两两独立,事件 $A\text{ 、 }B\text{ 、 }C$ 相互独立

B. 事件 $A\text{ 、 }B\text{ 、 }C$ 两两独立，事件 $A\text{ 、 }B\text{ 、 }C$ 不相互独立

C. 事件 $A\text{ 、 }B\text{ 、 }C$ 不两两独立，事件 $A\text{ 、 }B\text{ 、 }C$ 相互独立

D. 事件 $A\text{ 、 }B\text{ 、 }C$ 不两两独立，事件 $A\text{ 、 }B\text{ 、 }C$ 不相互独立

20.(2025$\cdot$崇明区二模$\cdot$15)抛掷一枚质地均匀的硬币 $n$ 次(其中 $n$ 为大于等于 2 的整数)，设事件 $A$ 表示“ $n$ 次中既有正面朝上又有反面朝上”，事件 $B$ 表示“ $n$ 次中至多有一次正面朝上”，若事件 $A$ 与事件 $B$ 是独立的,则 $n$ 的值为 \dotfill{(\kh{})}

\xx{}{}{}{}5 B. 4 C. 3 D. 2

\subsubsection*{三、解答题}

21.(2024$\cdot$宝山区一模$\cdot$17)一个盒子中装有 4 张卡片，卡片上分别写有数字1、2、3、4. 现从盒子中随机抽取卡片.

(1)若一次抽取 3 张卡片，事件 $A$ 表示“3 张卡片上数字之和大于 7”，求 $P\left( A\right)$ ；

(2)若第一次抽取 1 张卡片，放回后再抽取 1 张卡片，事件 $B$ 表示“两次抽取的卡片上数字之和大于 6 ”，求 $P\left( B\right)$ ；

(3)若一次抽取 2 张卡片，事件 $C$ 表示“ 2 张卡片上数字之和是 3 的倍数”，事件 $D$ 表示“ 2 张卡片上数字之积是 4 的倍数”. 验证 $C$ 、 $D$ 是独立的.

22.(2024$\cdot$静安区一模$\cdot$18)甲、乙两人每下一盘棋，甲获胜的概率是 0.4，甲不输的概率为 0.9.

(1)若甲、乙两人下一盘棋，求他们下成和棋的概率；

(2)若甲、乙两人连下两盘棋，假设两盘棋之间的胜负互不影响，求甲至少获胜一盘的概率.

23.(2025$\cdot$徐汇区一模$\cdot$19)某企业招聘员工，指定“英语听说”“信息技术”“逻辑推理”作为三门考试课程，有两种考试方案.

方案一:参加三门课程的考试，至少有两门及格为通过；

方案二:在三门课程中，随机选取两门，并参加这两门课程的考试，两门都及格为通过.

假设某应聘者参加三门指定课程考试及格的概率分别是 ${p}_{1}\text{ 、 }{p}_{2}\text{ 、 }{p}_{3}\left( {{p}_{i} \in  \left( {0,1}\right) ,i = 1,2,3}\right)$ ,且三门课程考试是否及格相互之间没有影响.

(1)分别求该应聘者选方案一考试通过的概率 ${T}_{1}$ 和选方案二考试通过的概率 ${T}_{2}$ ；

(2)试比较该应聘者在上述两种方案下考试通过的概率的大小，并说明理由.

考点6 > 古典概率

\subsubsection*{一、填空题}

1.(2025$\cdot$长宁区一模$\cdot$5)投掷两枚质地均匀的骰子，观察掷得的点数，则掷得的点数之和为 7 的概率是\blkda[2]{}.

2.(2024$\cdot$浦东新区一模$\cdot$8)在 100 件产品中有 90 件一等品、10 件二等品，从中随机抽取 3 件产品， 则恰好含 1 件二等品的概率为\blkda[2]{}. (结果精确到 0.01)

3.(2024$\cdot$松江区一模$\cdot$8)有 5 名同学报名参加暑期区科技馆志愿者活动，共服务两天，每天需要两人参加活动，则恰有 1 人连续参加两天志愿者活动的概率为\blkda[2]{}.

4.(2024$\cdot$金山区一模$\cdot$8)从1、2、3、4、5这五个数中随机抽取两个不同的数，则所抽到的两个数的和大于 6 的概率为\blkda[2]{}(结果用数值表示)

5.(2025$\cdot$闵行区二模$\cdot$9)投篮测试中，每人投 3 次，至少投中 2 次才能通过测试. 已知某同学每次投篮投中的概率为 0.6 ，且各次投篮是否投中相互独立，则该同学通过测试的概率为\blkda[2]{}。

6.(2025$\cdot$嘉定区二模$\cdot$9)在由1、2、3、4、5 这五个数组成的无重复数字的四位数中，其能被 3 整除的概率为\blkda[2]{}.

7.(2024$\cdot$黄浦区二模$\cdot$9)某校高三年级举行演讲比赛，共有 5 名选手参加. 若这 5 名选手甲、乙、丙、 丁、戊通过抽签来决定上场顺序, 则甲、乙两位选手上场顺序不相邻的概率为\blkda[2]{}.

8.(2025$\cdot$黄浦区二模$\cdot$10)若从 2025 的所有正约数中任取一个数，则这个数是一个完全平方数的概率为\blkda[2]{}.

9.(2023$\cdot$上海春季卷$\cdot$10)已知有 4 名男生、6 名女生，从10人中任选3人，则恰有 1 名男生 2 名女生的概率为\blkda[2]{}.

10.(2024$\cdot$虹口区一模$\cdot$10)将甲、乙等 8 人安排在 4 天值班，若每天安排两人，则甲、乙两人安排在同一天的概率为\blkda[2]{}. (结果用分数表示)

11.(2024$\cdot$嘉定区一模$\cdot$10)已知 11 个大小相同的球，其中 3 个是红球，3 个是黑球，5 个是白球，从中随机取出 4 个形成一组，其中三种颜色都有的概率为\blkda[2]{}.

12.(2025$\cdot$崇明区一模$\cdot$12)已知函数 $y = f\left( x\right)$ 的定义域 $D = \{ 1,2,3,4\}$ ，值域 $A = \{ 5,6,7\}$ ，则函数 $y = f\left( x\right)$ 为增函数的概率是\blkda[2]{}.

13.(2025$\cdot$静安区二模$\cdot$11)从 $m$ 个男生和 $n$ 个女生(10≥m>n之4)中任选 2 个人当队长，假设事件 $A$ 表示选出的 2 人性别相同，事件 $B$ 表示选出的 2 人性别不同. 如果事件 $A$ 的概率和事件 $B$ 的概率相等,那么 $m - n$ 的可能值为\blkda[2]{}.

14.(2025$\cdot$长宁区一模$\cdot$11)设 $O$ 为坐标原点，从集合\{1,2,3,4,5,6,7,8,9\}中任取两个不同的元素 $x\text{ 、 }y$ ，组成 $A\text{ 、 }B$ 两点的坐标 $\left( {x,y}\right) \text{ 、 }\left( {y,x}\right)$ ，则 ${S}_{\bigtriangleup {AOB}} \leq  {10}$ 的概率为\blkda[2]{}.

\subsubsection*{二、选择题}

15.(2024$\cdot$黄浦区一模$\cdot$14)从 3 名男同学和 2 名女同学中任选 2 名同学参加志愿者服务，则选出的 2 名同学中至少有 1 名女同学的概率是 \dotfill{(\kh{})}

\xx{}{}{}{}$\dfrac{7}{20}$ B. $\dfrac{7}{10}$

C. $\dfrac{3}{10}$ D. $\dfrac{3}{5}$

16.(2025$\cdot$徐汇区一模$\cdot$14)一个不透明的盒子中装有若干个红球和 5 个黑球，这些球除颜色外均相同. 每次将球充分搅匀后, 任意摸出 1 个球记下颜色后再放回盒子. 经过重复摸球足够多次试验后发现, 摸到黑球的频率稳定在 0.1 左右, 则据此估计盒子中红球的个数约为 \dotfill{(\kh{})}

\xx{}{}{}{}40 个 B. 45 个

C. 50 个 D. 55 个

17.(2025$\cdot$上海春季卷$\cdot$15)有一四边形 ABCD,对于其四边 ${AB}$ 、 ${BC}$ 、 ${CD}$ 、 ${DA}$ ，按顺序分别抛掷一枚质量均匀的硬币:如硬币正面朝上，则将其擦去；如硬币反面朝上，则不擦去. 最后，以 $A$ 为起点沿着尚未擦去的边出发，可以到达点 $C$ 的概率为 \dotfill{(\kh{})}

\xx{}{}{}{}$\dfrac{1}{2}$ B. $\dfrac{7}{16}$

C. $\dfrac{1}{4}$ D. $\dfrac{3}{16}$

18.(2024$\cdot$长宁区二模$\cdot$15)某运动员 8 次射击比赛的成绩为:9.6、9.7、9.5、9.9、9.4、9.8、9.3、 10.0 ; 已知这组数据的第 $x$ 百分位为 $m$ ,若从这组数据中任取一个数,这个数比 $m$ 大的概率为 0.25,则 $x$ 的取值不可能是 \dotfill{(\kh{})}

\xx{}{}{}{}65 B. 70

C. 75 D. 80

\subsubsection*{三、解答题}

19.(2025$\cdot$松江区一模$\cdot$17)某日用品按行业质量标准分成五个等级，等级系数X 依次为 1、2、、3、4、 5，现从一批该日用品中随机抽取 20 件，对其等级系数进行统计分析，得到频率分布表如下:

\begin{center}

\begin{tabular}{|c|c|c|c|c|c|}
\hline
$X$ & 1 & 2 & 3 & 4 & 5 \\
\cline{1-6}
$f$ & $a$ & 0.2 & 0.45 & $b$ & $C$ \\
\cline{1-6}
\hline
\end{tabular}
\end{center}

(1)若所抽取的 20 件日用品中，等级系数为 4 的恰有 3 件，等级系数为 5 的恰有 2 件，求 $a$ 、 $b$ 、 $c$ 的值;

(2)在(1)的条件下，将等级系数为 4 的 3 件日用品记为 ${x}_{1}$ 、 ${x}_{2}$ 、 ${x}_{3}$ ，等级系数为 5 的 2 件日用品记为 ${y}_{1}\text{ 、 }{y}_{2}$ ,现从 ${x}_{1}\text{ 、 }{x}_{2}\text{ 、 }{x}_{3}\text{ 、 }{y}_{1}\text{ 、 }{y}_{2}$ 这 5 件日用品中任取两件(假定每件日用品被取出的可能性相同), 写出所有可能的结果, 并求这两件日用品的等级系数恰好相等的概率.

20. (2025$\cdot$奉贤区一模$\cdot$18)某芯片代工厂生产甲、乙两种型号的芯片，为了解芯片的某项指标，从这两种芯片中各抽取 100 件进行检测，获得该项指标的频率分布直方图，如图所示:

\begin{center}
\includegraphics[width=0.8\textwidth]{bodyxb/src/236_272_1117_981_406_0.jpg}
\end{center}


第 20 题图

假设数据在组内均匀分布, 以样本估计总体, 以事件发生的频率作为相应事件发生的概率.

(1)求频率分布直方图中 $x$ 的值并估计乙型芯片该项指标的平均值; (同一组中的数据用该组区间的中点值为代表)

(2)已知甲型芯片指标在 $\lbrack {80},{100})$ 为航天级芯片，乙型芯片指标在 $\lbrack {60},{70})$ 为航天级芯片. 现分别采用分层抽样的方式,从甲型芯片指标在 $\lbrack {70},{90})$ 内取 2 件,乙型芯片指标在 $\lbrack {50},{70})$ 内取 4 件, 再从这 6 件中任取 2 件, 求至少有一件为航天级芯片的概率.

21. (2024$\cdot$崇明区一模$\cdot$19)交通拥堵指数(TPI)是表征交通拥堵程度的客观指标，用 TPI 表示，TPI 越大，代表拥堵程度越高. 某平台计算 TPI 的公式为:TPI = 纺锤行程时间，并按 TPI 的大小将城市道路拥堵程度划分如下表所示的 4 个等级:

\begin{center}

\begin{tabular}{|c|c|c|c|c|}
\hline
TPI & [1,1.5) & [1.5,2) & $\lbrack 2,4)$ & 不低于 4 \\
\cline{1-5}
拥堵等级 & 畅通 & 缓行 & 拥堵 & 严重拥堵 \\
\cline{1-5}
\hline
\end{tabular}
\end{center}

某市 2023 年元旦及前后共 7 天与 2022 年同期的交通高峰期城市道路 TPI 的统计数据如下图:

\begin{center}
\includegraphics[width=0.7\textwidth]{bodyxb/src/237_392_602_882_391_0.jpg}
\end{center}


第 21 题图

(1)从 2022 年元旦及前后共 7 天中任取 1 天，求这一天交通高峰期城市道路拥堵程度为 “拥堵” 的概率;

(2)从 2023 年元旦及前后共 7 天中任取 3 天，将这 3 天中交通高峰期城市道路 TPI 比 2022 年同日 TPI 高的天数记为 $X$ ,求所有 $X$ 的可能值及其发生的概率.

22.(2024$\cdot$闵行区一模$\cdot$19)2023 年 9 月 23 日至 10 月 8 日，第19届亚运会在杭州成功举办，杭州亚运会的志愿者被称为“小青荷”，某运动场馆内共有小青荷 36 名，其中男生 12 名，女生 24 名，这些小青荷中会说日语和会说韩语的人数统计如表所示: (其中 $m\text{ 、 }n$ 均为正整数, $6 \leq  m \leq  8$.)

\begin{center}

\begin{tabular}{|c|c|c|}
\hline
\phantom{X} & 男生小青荷 & 女生小青荷 \\
\cline{1-3}
会说日语 & 8 & 12 \\
\cline{1-3}
会说韩语 & $m$ & $n$ \\
\cline{1-3}
\hline
\end{tabular}
\end{center}

(1)从这 36 名小青荷中随机抽取两名作为某活动主持人， 求抽取的两名小青荷中至少有一名会说日语的概率;

(2)从这些小青荷中随机抽取一名去接待外宾,用 $A$ 表示事件“抽到的小青荷是男生”，用 $B$ 表示事件“抽到的小青荷会说韩语”. 试给出一组 $m\text{ 、 }n$ 的值,使得事件 $A$ 与 $B$ 相互独立,并说明理由.

23.(2025$\cdot$嘉定区一模$\cdot$19)在一场盛大的电竞比赛中，有两支实力强劲的队伍甲和乙进行对决. 比赛采用 5 局 3 胜制, 最终的胜者将赢得 10 万元奖金. 比赛过程中, 每局比赛双方获胜的概率相互独立且甲队每局获胜概率为 0.4 ，乙队每局获胜概率为 0.6. 比赛开始后，甲队先连胜两局，此时， 主办方记录了两队队员在这两局比赛中的一些数据. 甲队队员的击杀数(单位:个)数据如下: 24、31、31、36、36、37、39、44、49、50；乙队队员的击杀数(单位:个)数据如下:8、13、14、16、23、26、 28、33、38、39. 然而此时比赛场地突发技术故障，比赛不得不中止. 请回答以下问题:

(1)根据目前情况(甲队已连胜两局)，写出甲、乙两队“采用 5 局 3 胜制”的比赛结果的样本空间；

(2)根据所给数据，绘制甲、乙两队队员的击杀数分布的茎叶图；

(3)在目前情况下(甲队已连胜两局)，估算甲、乙两队获胜概率，并据此分配 10 万元奖金。

24.(2025$\cdot$闵行区一模$\cdot$19)为了解某市高三学生的睡眠时长，从该市 6.6 万名高三学生中随机抽取 600 人，统计他们的日均睡眠时长及分布人数如下表所示:

\begin{center}

\begin{tabular}{|c|c|c|c|}
\hline
睡眠时长(h) & $\lbrack 4,6)$ & [6,8) & [8, 10] \\
\cline{1-4}
人数 & 150 & 270 & 180 \\
\cline{1-4}
\hline
\end{tabular}
\end{center}

其中睡眠时长在 $\left\lbrack  {8,{10}}\right\rbrack$ 的为睡眠充足,在 $\lbrack 6,8)$ 的为睡眠良好，在 $\lbrack 4,6)$ 的为睡眠不足.

(1)估计该市 6.6 万名高三学生中日均睡眠时长大于等于 6 小时的人数约为多少？

(2)估计该市高三学生日均睡眠时长；

(3)若从这 600 名学生中利用分层抽样的方法抽取 20 人，再从这 20 人中随机抽取 4 人做进一步访谈调查, 求这 4 人中既有睡眠充足, 又有睡眠良好, 也有睡眠不足学生的概率. 25.(2025$\cdot$宝山区一模$\cdot$19)甲、乙两人轮流掷质地均匀的骰子，每人每次掷两颗.

(1)甲掷一次，求两颗骰子点数不同的概率；

(2)甲、乙各掷一次，求甲的点数和恰好比乙的点数和大 7 的概率；

(3)若第一次掷出点数之和大于 6 的人为胜者，同时比赛结束；否则，由另一人继续投掷，直到比赛结束. 例如, 甲乙先后轮流掷出的点数之和为:5、4、3、7, 此时乙为胜者. 设甲先投掷, 求甲最终获胜的概率.

26.(2025$\cdot$杨浦区一模$\cdot$19)为加强学生睡眠监测督导，学校对高中三个年级学生的日均睡眠时间进行调查. 根据分层随机抽样法,学校在高一、高二和高三年级中共抽取了 100 名学生的日均睡眠时间作为样本, 其中高一 35 人, 高二 33 人. 已知该校高三年级一共 512 人.

(1)学校高中三个年级一共有多少个学生？

(2)若抽取 100 名学生的样本极差为 2，数据如下表所示:(其中 $x < {10}$ ， $n$ 是正整数)

\begin{center}

\begin{tabular}{|c|c|c|c|c|c|}
\hline
日均睡眠时间(h) & $x$ & 8.5 & 9 & 9.5 & 10 \\
\cline{1-6}
学生数量 & $n$ & 32 & 13 & 11 & 4 \\
\cline{1-6}
\hline
\end{tabular}
\end{center}

求该样本的第 40 百分位数;

(3)从这 100 名学生的样本中随机抽取三个学生的日均睡眠时间，求其中至少有 1 个数据来自高三学生的概率.

27.(2024$\cdot$普陀区一模$\cdot$17)我国随着人口老龄化程度的加剧，劳动力人口在不断减少，“延迟退休” 已成为公众关注的热点话题之一. 为了了解公众对“延迟退休”的态度，某研究机构对属地所在的一社区进行了调查, 并将随机抽取的 50 名被调查者的年龄制成如图所示的茎叶图.

\begin{center}
\includegraphics[width=0.4\textwidth]{bodyxb/src/240_892_349_551_297_0.jpg}
\end{center}


第 27 题图

(1)经统计发现，投赞成票的人均年龄恰好是这 50 人年龄的第 60 百分位数，求此百分位数；

(2)经统计年龄在 $\lbrack {50},{59})$ 的被调查者中,投赞成票的男性有 3 人、女性有 2 人，现从该组被调查者中随机选取男女各 2 人进行跟踪调查, 求被选中的 4 人中至少有 3 人投赞成票的概率. (结果用最简分数表示)

考点7 条件概率

\subsubsection*{一、填空题}

1.(2025$\cdot$普陀区二模$\cdot$7)在一个不透明的盒中装着标有数字1、2、、3、4 的大小与质地都相同的小球各2个，现从该盒中一次取出 2 个球，设事件A为“取出 2 个球的数字之和大于 5”，事件B 为“取出的 2 个球中最小数字是 2 ”，则 $P\left( {B \mid  A}\right)  =$ \blkda[2]{}.

2.(2025$\cdot$静安区二模$\cdot$7)设一个罐子中有大小与质地相同的黑、白、红三个球，不放回的每次摸一个球,设第一次没有摸到黑球是事件 $A$ ,第二次没有摸到黑球是事件 $B$ ,则 $P\left( {B \mid  A}\right)$ 的值为\blkda[2]{}.

3.(2024$\cdot$嘉定区二模$\cdot$9)小张、小王两家计划假期去嘉定游玩，他们分别从古猗园、秋霞圃、州桥老街三个景点中随机选择一个游玩，记事件 $A$ 表示“两家至少有一家选择古猗园”，事件 $B$ 表示“两家选择景点不同”，则概率 $P\left( {B \mid  A}\right)  =$ \blkda[2]{}.

4.(2025$\cdot$徐汇区二模$\cdot$9)已知两个随机事件 $A$ 、 $B$ ，若 $P\left( A\right)  = \dfrac{1}{5},P\left( B\right)  = \dfrac{1}{4},P\left( {B \mid  A}\right)  = \dfrac{2}{3}$ ，则 $P\left( {\overline{A} \mid  B}\right)  =$

5.(2025$\cdot$奉贤区二模$\cdot$10)盒子中有大小与质地均相同的 $a$ 个红球和 $b$ 个白球，从中随机取 1 个球， 观察其颜色后放回，并同时放入与其相同颜色的球 $c$ 个(大小与质地均相同)，再从中随机取 1 个球，计算此次取到白球的概率是\blkda[2]{}.

6.(2024$\cdot$青浦区二模$\cdot$10)从1、2、3、4、5中任取2个不同的数字，设“取到的2个数字之和为偶数” 为事件 $A$ ，“取到的 2 个数字均为奇数”为事件 $B$ ，则 $P\left( {B \mid  A}\right)  =$ \blkda[2]{}.

7.(2024$\cdot$静安区二模$\cdot$10)某工厂生产的产品以 100 个为一批. 在进行抽样检查时，只从每批中抽取 10 个来检查, 如果发现其中有次品, 则认为这批产品是不合格的. 假定每一批产品中的次品最多不超过 2 个，并且其中恰有 $i\left( {i = 0,1,2}\right)$ 个次品的概率如下: 则各批产品通过检查的概率为\blkda[2]{}.

\begin{center}

\begin{tabular}{|c|c|c|c|}
\hline
一批产品中有次品的个数 $i$ & 0 & 1 & 2 \\
\cline{1-4}
概率 & 0.3 & 0.5 & 0.2 \\
\cline{1-4}
\hline
\end{tabular}
\end{center}

(精确到 0.01)

\subsubsection*{二、解答题}

8.(2025$\cdot$宝山区二模$\cdot$19)某游乐园的活动项目共有三类，分别是“过山车”等 10 个体验类项目、“海豚之舞”等 4 个表演类项目、“智力闯关”等 3 个互动类项目. 因设备维护需要，项目并非每日都全部开放. 以下数据是项目开放的数量 $x$ (个)和游客平均等待时间 $y$ (min/个)的关系:

\begin{center}

\begin{tabular}{|c|c|c|c|c|c|c|c|c|c|}
\hline
项目类别 & \multicolumn{5}{|c|}{体验类} & \multicolumn{2}{|c|}{演出类} & \multicolumn{2}{|c|}{互动类} \\
\cline{1-10}
开放数量 $x$ (个) & 4 & 5 & 6 & 7 & 8 & 2 & \phantom{X} & 2 & 3 \\
\cline{1-10}
平均等待时间 $y$ (分钟/个) & 76 & 73 & 67 & $m$ & 60 & 53 & \phantom{X} & \phantom{X} & 30 \\
\cline{1-10}
\hline
\end{tabular}
\end{center}

(1)体验类项目中,若 $y$ 关于 $x$ 的回归方程为 $y =  - {4.3x} + {93.8}$ ，请计算 $m$ 的值，并依据该模型预测所有体验类项目均开放时的平均等待时间; (精确到整数)

(2)小王游玩当日，体验类、演出类、互动类项目分别开放了 8 个、 4 个、 3 个，他计划随机游玩其中的 3 个项目,已知他选择的项目中至少包含 1 个互动类项目,求他的等待总时间恰为 ${120}{{min}}$ 的概率;

(3)为提高游客的参与度，园方在互动类项目“智力闯关”中设计了两关，通过第一关的游客奖励 20 个游园币, 游客可以选择结束或继续闯关. 若继续闯关, 则必须完成第二关的所有题目. 第二关包含 2 道相互独立的选择题, 每答对 1 题可再奖励 20 个游园币, 每答错 1 题则要扣除 10 个游园币. 每个游园币可兑换园区内任意一个项目的 1 min 等待时间. 小王已通过第一关，假设他在第二关中每道题答对的概率均为 $p$ ,为了获得更多项目等待时间的兑换奖励,小王是否应该继续闯关？请你帮他做出决策.

9.(2025$\cdot$长宁区一模$\cdot$19)2024 年第七届中国国际进口博览会(简称进博会)于 11 月 5 日至 10 日在上海国家会展中心举行. 为了解进博会参会者的年龄结构, 某机构随机抽取了年龄在 15-75 岁之间的 200 名参会者进行调查,并按年龄绘制了频率分布直方图,分组区间为 $\lbrack {15},{25})\text{ 、 }\lbrack {25},{35})$ 、 [35,45)、[45,55)、[55,65)、[65,75]. 把年龄落在区间 $\lbrack {15},{35})$ 内的人称为“青年人”，把年龄落在区间[35,65)内的人称为“中年人”，把年龄落在[65,75]内的人称为“老年人”.

(1)求所抽取的“青年人”的人数；

(2)以分层抽样的方式从“青年人”“中年人”“老年人”中抽取 10 名参会者做进一步访谈，发现其中女性共 4 人，这 4 人中有 3 人是“中年人”。再用抽签法从所抽取的 10 名参会者中任选 2 人.

\begin{center}
\includegraphics[width=0.4\textwidth]{bodyxb/src/242_950_595_486_359_0.jpg}
\end{center}


第 9 题图

 \ding{172}  简述如何采用抽签法任选 2 人；

 \ding{173}  设事件 $A : 2$ 人均为“中年人”，事件 $B : 2$ 人中至少有 1 人为男性，判断事件 $A$ 与事件 $B$ 是否独立,并说明理由.

10.(2024$\cdot$崇明区二模$\cdot$19)某疾病预防中心随机调查了 340 名 50 岁以上的公民，研究吸烟习惯与慢性气管炎患病的关系，调查数据如表所示.

\begin{center}

\begin{tabular}{|c|c|c|c|}
\hline
\phantom{X} & 不吸烟者 & 吸烟者 & 总计 \\
\cline{1-4}
不患慢性气管炎者 & 120 & 160 & 280 \\
\cline{1-4}
患慢性气管炎者 & 15 & 45 & 60 \\
\cline{1-4}
总计 & 135 & 205 & 340 \\
\cline{1-4}
\hline
\end{tabular}
\end{center}

(1)是否有 95\% 的把握认为患慢性气管炎与吸烟有关？

(2)常用 $L\left( {B \mid  A}\right)  = \dfrac{P\left( {B \mid  A}\right) }{P\left( {\overline{B} \mid  A}\right) }$ 表示在事件 $A$ 发生的条件下事件 $B$ 发生的优势,在统计中称为似然比. 现从 340 人中任选一人, $A$ 表示“选到的人是吸烟者”, $B$ 表示“选到的人患慢性气管炎者”. 请利用样本数据,估计 $L\left( {B \mid  A}\right)$ 的值.

(3)现从不患慢性气管炎者的样本中，按分层抽样的方法选出 7 人，从这 7 人里再随机选取 3 人， 求这 3 人中，不吸烟者的人数 $X$ 的数学期望.

附: ${\chi }^{2} = \dfrac{n{\left( ad - bc\right) }^{2}}{\left( {a + b}\right) \left( {c + d}\right) \left( {a + c}\right) \left( {b + d}\right) },P\left( {{\chi }^{2} \geq  {3.841}}\right)  \approx  {0.05}$.

考点8 二 全概率公式

\subsubsection*{一、填空题}

1.(2024$\cdot$浦东新区二模$\cdot$7)某校面向高一全体学生共开设 3 门体育类选修课，每人限选一门. 已知这三门体育类选修课的选修人数之比为 $6 : 3 : 1$ ，考核优秀率分别为 ${20}\%$ 、 ${16}\%$ 和 ${12}\%$ ，现从该年级所有选择体育类选修课的同学中任取一名，其成绩是优秀的概率为\blkda[2]{}.

2.(2024$\cdot$杨浦区一模$\cdot$8)甲和乙两射手射击同一目标，命中的概率分别为0.7 和0.8，两人各射击一次，假设事件“甲命中”与“乙命中”是独立的，则至少一人命中目标的概率为\blkda[2]{}.

3.(2024$\cdot$上海秋季卷$\cdot$8)某校举办科学竞技比赛，有 $A$ 、 $B$ 、 $C$ 3 种题库， $A$ 题库有 5000 道题、 $B$ 题库有 4000 道题、 $C$ 题库有 3000 道题. 小申已完成所有题,已知他完成 $A$ 题库的正确率是 0.92, $B$ 题库的正确率是 0.86 ，C 题库的正确率是 0.72 。现他从所有的题中随机选一题，正确率是\blkda[2]{}.

4.(2025$\cdot$金山区一模$\cdot$11)抛掷一枚质地均匀的硬币 n 次(其中 n 为大于等于 2 的整数)，设事件 $A$ : $n$ 次中既有正面朝上又有反面朝上，事件 $B : n$ 次中至多有一次正面朝上，若事件 $A$ 与事件 $B$ 是独立的，则 $n$ 的值为\blkda[2]{}.

\subsubsection*{二、选择题}

5.(2025$\cdot$杨浦区一模$\cdot$14)如果 $A$ 、 $B$ 是独立事件， $\overline{A}$ 、 $\overline{B}$ 分别是 $A$ 、 $B$ 的对立事件，那么以下等式不一定成立的是 \dotfill{(\kh{})}

\xx{}{}{}{}$P\left( {A \cap  B}\right)  = P\left( A\right) P\left( B\right)$ B. $P\left( {\overline{A} \cap  B}\right)  = P\left( \overline{A}\right) P\left( B\right)$

C. $P\left( {A \cup  B}\right)  = P\left( A\right)  + P\left( B\right)$ D. $P\left( {\overline{A} \cap  \overline{B}}\right)  = \left\lbrack  {1 - P\left( A\right) }\right\rbrack  \left\lbrack  {1 - P\left( B\right) }\right\rbrack$

6.(2025$\cdot$虹口区二模$\cdot$15)春节期间，小明和弟弟玩起了一种自定义游戏，规定先由弟弟掷一颗质地均匀的骰子,若弟弟掷出的点数为 6 , 则吃 1 颗花生; 若掷出其他点数,则记下这个点数,然后由小明开始两个人轮流掷这颗骰子，直至任意一方掷出这个记下的点数或者 6 ，一次游戏结束. 若掷出的是这个记下的点数，则弟弟吃 1 颗花生；若是 6 ，则小明吃 3 颗花生. 任意一次游戏中弟弟能吃到 1 颗花生的概率为 \dotfill{(\kh{})}

\xx{}{}{}{}$\dfrac{5}{24}$ B. $\dfrac{5}{12}$ C. $\dfrac{3}{8}$ D. $\dfrac{7}{12}$

\subsubsection*{三、解答题}

7.(2025$\cdot$黄浦区二模$\cdot$19)一盒子中有大小与质地均相同的 20 个小球,其中白球 $n\left( {3 \leq  n \leq  {13}}\right)$ 个， 其余为黑球.

(1)当盒中的白球数 $n = 6$ 时,从盒中不放回地随机取两次,每次取一个球,用 $A$ 表示事件“第一次取到白球”,用 $B$ 表示事件“第二次取到白球”,求 $P\left( {B \mid  A}\right)$ 和 $P\left( B\right)$ ,并判断事件 $A$ 与 $B$ 是否相互独立.

(2)某同学要策划一个抽奖活动，参与者从盒中一次性随机取 10 个球. 若其中恰有 3 个白球，则获奖， 否则不获奖. 要使参与者获奖的可能性最大、最小,该同学应该分别如何放置白球的数量 $n$ ?

8.(2025$\cdot$普陀区一模$\cdot$19)机器人竞技是继电子竞技之后热门的科技竞技项目. 某区为了参加市机器人竞技总决赛，开展了区内选拔赛，其中 $A$ 、 $B$ 、 $C$ 、 $D$ 四人进入区内个人组决赛，按照规则每人与其他三人各进行一场比赛,且这三场比赛互相独立. 下表统计的是 $A$ 在近期热身中分别与 $B$ 、 $C$ 、 $D$ 三人比赛的情况.

(1)根据表格中的数据，试估计在区内决赛中 $A$ 至少获胜一场的概率；

(2)根据表格中的数据，请给 $B$ 、 $C$ 、 $D$ 三人设计一个出场顺序，使得 $A$ 在这三场比赛中连胜两场的概率最大,并说明理由.

\begin{center}

\begin{tabular}{|c|c|c|c|}
\hline
\phantom{X} & $B$ & $C$ & $D$ \\
\cline{1-4}
比赛的次数 & 12 & 10 & 15 \\
\cline{1-4}
$A$ 获胜的次数 & 4 & 5 & 12 \\
\cline{1-4}
\hline
\end{tabular}
\end{center}

9.(2025$\cdot$浦东新区一模$\cdot$19)申辉中学为期两周的高一、高二年级校园篮球赛告一段落. 高一小 A、 高二小 $B$ 分别荣获了高一年级和高二年级比赛的年级 MVP (最有价值球员). 以下是他们在各自 8 场比赛的二分球和三分球出手次数及其命中率.

\begin{center}

\begin{tabular}{|c|c|c|c|c|}
\hline
\phantom{X} & 二分球出手 & 二分球命中率 & 三分球出手 & 三分球命中率 \\
\cline{1-5}
小 $A$ & 100 次 & 80\% & 100 次 & 40\% \\
\cline{1-5}
小 $B$ & 190 次 & 70\% & 10 次 & 30\% \\
\cline{1-5}
\hline
\end{tabular}
\end{center}

现以两人的总投篮命中率(二分球+三分球)较高者评为校 MVP(总投篮命中率= 总命中次数÷ 总出手次数)

(1)小 C 认为，目测小 $A$ 的二分球命中率和三分球命中率均高于小 $B$ ，此次必定能评为校 MVP， 试通过计算判断小 $C$ 的想法是否准确？

(2)小D是游戏爱好者，设置了一款由游戏人物小 a、小 $b$ 轮流投篮对战游戏，游戏规则如下:  \ding{172}  游戏中小 a 的命中率始终为 0.4 ，小 $b$ 的命中率始终为 0.3 ； \ding{173}  游戏中投篮总次数最多为 $k  \left( {3 \leq  k \leq  {20},k \in  \mathbf{Z}}\right)$ 次，且同一个游戏人物不允许连续投篮； \ding{174}  游戏中若投篮命中，则游戏结束，投中者获得胜利，若直至第 $k$ 次投篮都没有命中，则规定第二次投篮者获胜. 若每次游戏对战前必须设置 “第一次投篮人物” 和“$k$”的值，请解答以下两个问题.

(i)若小 $a$ 第一次投篮，请证明小 $a$ 获胜概率大；

(ii)若小 $b$ 第一次投篮，试问谁的获胜概率大？并说明理由.

10.(2025$\cdot$嘉定区二模$\cdot$19)某学校对学生的课外阅读时间进行调查，随机抽取了 150 名学生，得到如下样本数据频率分布直方图.

(1)估计该校学生的平均课外阅读时间；(同一组数据用该区间的中点值作代表)

(2)估计该校学生课外阅读时间位于区间 $\lbrack {30},{60})$ (单位: ${h}/{\text{ 月 }}$ ) 的概率；

(3)已知该校喜欢阅读的学生占比为 18\%，初一年级学生占该校总学生数的 28\%，且初一年级学生中喜欢阅读的占 40\%，求其他年级学生中喜欢阅读的比例.

\begin{center}
\includegraphics[width=0.5\textwidth]{bodyxb/src/246_809_503_599_310_0.jpg}
\end{center}


第 10 题图

\subsubsection*{一、填空题}

1.(2024$\cdot$松江区二模$\cdot$3)已知随机变量 $X$ 服从正态分布 $N\left( {3,{\sigma }^{2}}\right)$ ，且 $P\left( {3 \leq  X \leq  5}\right)  = {0.3}$ ，则 $P\left( {X > 5}\right)  =$ \blkda[2]{}.

2.(2024$\cdot$普陀区二模$\cdot$4)已知 $X \sim  N\left( {4,{2}^{2}}\right)$ ，若 $P\left( {X < 0}\right)  = 0$.02，则 $P\left( {4 < X < 8}\right)  =$ \blkda[2]{}.

3.(2024$\cdot$奉贤区一模$\cdot$4)某公司生产的糖果每包标识质量是 500g，但公司承认实际质量存在误差. 已知糖果的实际质量 $X$ 服从 $\mu  = {500}$ 的正态分布. 若随意买一包糖果,假设质量误差超过 $5{g}$ 的可能性为 $p$ ,则 $P\left( {{495} \leq  X \leq  {500}}\right)$ 的值为\blkda[2]{}. (用含 $p$ 的代数式表达)

4.(2024$\cdot$浦东新区二模$\cdot$5)已知随机变量 $X$ 服从正态分布 $N\left( {{95},{\sigma }^{2}}\right)$ ，若 $P\left( {{75} \leq  X \leq  {115}}\right)  = {0.4}$ ，则 $P\left( {X > {115}}\right)  =$ \blkda[2]{}.

5.(2024$\cdot$青浦区二模$\cdot$5)设随机变量 $\xi$ 服从正态分布 $N\left( {2,1}\right)$ ，若 $P\left( {\xi  < a - 3}\right)  = P\left( {\xi  > 1 - {2a}}\right)$ ，则实数 $a =$ \blkda[2]{}.

6.(2024$\cdot$奉贤区二模$\cdot$5)已知随机变量 $X$ 服从正态分布 $N\left( {2,{\sigma }^{2}}\right)$ ，且 $P\left( {2 < X \leq  {2.5}}\right)  = {0.36}$ ，则 $P\left( {X > {2.5}}\right)  =$ \blkda[2]{}.

7.(2024$\cdot$静安区二模$\cdot$5)某地区高三年级 2000 名学生参加了地区教学质量调研测试，已知数学测试成绩 $X$ 服从正态分布 $N\left( {{100},{\sigma }^{2}}\right)$ (试卷满分 150 分)，统计结果显示，有 320 名学生的数学成绩低于 80 分，则数学分数属于闭区间[80,120]的学生人数约为\blkda[2]{}.

8.(2024$\cdot$黄浦区二模$\cdot$7)随机变量 $X$ 服从正态分布 $N\left( {2,{\sigma }^{2}}\right)$ ，若 $P\left( {2 < X \leq  {2.5}}\right)  = {0.36}$ ，则 $P\left( {\left| {X - 2}\right|  > {0.5}}\right)  =$ \blkda[2]{}.

9.(2025$\cdot$青浦区二模$\cdot$8)已知随机变量 $\xi  \sim  N\left( {3,{\sigma }^{2}}\right) \left( {\sigma  > 0}\right)$ ，若 $P\left( {\xi  \geq  1}\right)  = {0.9}$ ，则 $P\left( {3 \leq  \xi  \leq  5}\right)  =$ \blkda[2]{}.

10.(2025$\cdot$闵行区二模$\cdot$9)某公司生产的糖果每包的标识质量是 ${500}{g}$ ，但公司承认实际质量存在误差. 已知每包糖果的实际质量服从正态分布 $N\left( {{500},{\sigma }^{2}}\right)$ ,且任意一包的糖果质量介于 ${495}{g}$ 到 ${505}{g}$ 之间的可能性为 ${95.4}\%$ ,则随意买一包该公司生产的糖果,其质量超过 ${505}{g}$ 的可能性约为\blkda[2]{}. (精确到 ${0.1}\%$ )

11. (2025$\cdot$虹口区二模$\cdot$9) 某工厂生产的零件长度 $X$ (单位: ${{mm}}$ ) 服从正态分布 $N\left( {3,{\sigma }^{2}}\right)$ ,且 $P\left( {\left| {X - 3}\right|  \leq  {0.5}}\right)  = {0.8}$ ,若对该工厂同批生产的 4 个零件逐一检查,则仅有 1 个零件的长度大于 ${3.5}{{mm}}$ 的概率为\blkda[2]{}.

12.(2025$\cdot$嘉定区二模$\cdot$10)已知某次数学的测试成绩 $X$ 服从 $\mu  = {75}\text{ 、 }{\sigma }^{2} = {64}$ 的正态分布，若小明的成绩不低于 91 分，那么他的成绩大约超过了\blkda[2]{}\% 的学生(精确到 0.1\%). (参考数据: $P\left( {\left| X\right|  < \sigma }\right)  \approx  {68.3}\% ,P\left( {\left| X\right|  < {2\sigma }}\right)  \approx  {95.4}\% ,P\left( {\left| X\right|  < {3\sigma }}\right)  \approx  {99.7}\% )$

\subsubsection*{二、选择题}

13.(2024$\cdot$宝山区二模$\cdot$14)已知随机变量 $X$ 服从正态分布 $N\left( {0,{\sigma }^{2}}\right)$. 若 $P\left( {X \leq  a}\right)  = \dfrac{5}{6}$ ，则 $P\left( {\left| X\right|  \leq  a}\right)  =$ \dotfill{(\kh{})}

\xx{}{}{}{}$\dfrac{2}{3}$ B. $\dfrac{1}{2}$

C. $\dfrac{1}{3}$ D. $\dfrac{1}{6}$

14.(2024$\cdot$金山区二模$\cdot$14)下列说法不正确的是 \dotfill{(\kh{})}

\xx{}{}{}{}一组数据 ${10}\text{ 、 }{11}\text{ 、 }{11}\text{ 、 }{12}\text{ 、 }{13}\text{ 、 }{14}\text{ 、 }{16}\text{ 、 }{18}\text{ 、 }{20}\text{ 、 }{22}$ 的第 60 百分位数为 14

B. 若随机变量 $X$ 服从正态分布 $N\left( {3,{\sigma }^{2}}\right)$ ，且 $P\left( {X \leq  4}\right)  = {0.7}$ ，则 $P\left( {3 < X < 4}\right)  = {0.2}$

C. 若线性相关系数 $\left| r\right|$ 越接近 1,则两个变量的线性相关程度越高

D. 对具有线性相关关系的变量 $x$ 、 $y$ ，且回归方程为 $y = {0.3x} - m$ ，若样本点的中心为 $\left( {m,2,8}\right)$ ， 则实数 $m$ 的值是 -4

15.(2024$\cdot$杨浦区二模$\cdot$15)某区高三年级 3 200 名学生参加了区统一考试，考试成绩 $X$ 服从正态分布 $N\left( {{100},{\sigma }^{2}}\right)$ (试卷满分为 150 分). 统计结果显示,考试成绩在 80 分到 120 分之间的人数约为总人数的 $\dfrac{3}{4}$ ，则此次考试中成绩不低于 120 分的学生人数约为 \dotfill{(\kh{})}

\xx{}{}{}{}350 B. 400

C. 450 D. 500

\subsubsection*{三、解答题}

16.(2025$\cdot$徐汇区二模$\cdot$19)某公司生产的糖果每包标识“净含量 500g”，但公司承认实际的净含量存在误差. 已知每包糖果的实际净含量 $\xi$ (单位: ${g}$ ) 服从正态分布 $N\left( {{500},{2.5}^{2}}\right)$.

(1)随机抽取一包该公司生产的糖果，求其净含量误差超过 $5{g}$ 的概率(精确到 0.001)；

(2)随机抽取 3 包该公司生产的糖果，记其中净含量小于 497.5 g 的包数为 $X$. 求 $X$ 的分布和期望 (精确到 0.001 ). (参考数据: $\Phi \left( 1\right)  \approx  {0.8413},\Phi \left( 2\right)  \approx  {0.9772},\Phi \left( 3\right)  \approx  {0.9987}$ ,其中 $y =  \Phi \left( x\right)$ 为标准正态分布函数.)

\subsubsection*{一、填空题}

1.(2024$\cdot$金山区二模$\cdot$9)为了考察某种药物预防疾病的效果，进行动物试验，得到结果如列联表所示:

\begin{center}

\begin{tabular}{|c|c|c|c|}
\hline
\multirow{2}{*}{药物} & \multicolumn{2}{|c|}{疾病} & \multirow{2}{*}{合计} \\
\cline{2-3}
 & 未患病 & 患病 &  \\
\cline{1-4}
服用 & $m$ & 50-m & 50 \\
\cline{1-4}
未服用 & 80-m & m-30 & 50 \\
\cline{1-4}
合计 & 80 & 20 & 100 \\
\cline{1-4}
\hline
\end{tabular}
\end{center}

取显著性水平 $\alpha  = {0.05}$ ,若本次考察结果支持 “药物对疾病预防有显著效果”,则 $m\left( {m \geq  {40}}\right.$ , $m \in  \mathbf{N}$ ) 的最小值为\blkda[2]{}.

(参考公式: ${\chi }^{2} = \dfrac{n{\left( ad - bc\right) }^{2}}{\left( {a + b}\right) \left( {c + d}\right) \left( {a + c}\right) \left( {b + d}\right) }$ ; 参考值: $P\left( {{\chi }^{2} \geq  {3.841}}\right)  \approx  {0.05})$

\subsubsection*{二、选择题}

2.(2024$\cdot$闵行区二模$\cdot$15)某疾病预防中心随机调查了 339 名 50 岁以上的公民，研究吸烟习惯与慢性气管炎患病的关系，调查数据如下表:

\begin{center}

\begin{tabular}{|c|c|c|c|}
\hline
\phantom{X} & 不吸烟者 & 吸烟者 & 总计 \\
\cline{1-4}
不患慢性气管炎者 & 121 & 162 & 283 \\
\cline{1-4}
患慢性气管炎者 & 13 & 43 & 56 \\
\cline{1-4}
总计 & 134 & 205 & 339 \\
\cline{1-4}
\hline
\end{tabular}
\end{center}

假设 ${H}_{0}$ : 患慢性气管炎与吸烟没有关系,即它们相互独立. 通过计算统计量 ${\chi }^{2}$ ,得 ${\chi }^{2} \approx  {7.468}$ ,根据 ${\chi }^{2}$ 分布概率表: $P\left( {{\chi }^{2} \geq  {6.635}}\right)  \approx  {0.01},P\left( {{\chi }^{2} \geq  {5.024}}\right)  \approx  {0.025},P\left( {{\chi }^{2} \geq  {3.841}}\right)  \approx  {0.05},P\left( {{\chi }^{2} \geq  }\right.$ 2.706)≈0.1. 给出下列 3 个命题, 其中正确的个数是 \dotfill{(\kh{})}

 \ding{172}  “患慢性气管炎与吸烟没有关系”成立的可能性小于 5\%；

 \ding{173}  有 99\% 的把握认为患慢性气管炎与吸烟有关；

 \ding{174}  ${\chi }^{2}$ 分布概率表中的 0.05、0.01 等小概率值在统计上称为显著性水平，小概率事件一般认为不太可能发生.

\xx{}{}{}{}0 个 B. 1 个 C. 2 个 D. 3 个

\subsubsection*{三、解答题}

3.(2025$\cdot$奉贤区二模$\cdot$17)某疾病预防中心随机调查了 339 名 50 岁以上的公民，研究吸烟习惯与慢性气管炎患病的关系，测得数据如表所示:

\begin{center}

\begin{tabular}{|c|c|c|c|}
\hline
\phantom{X} & 不吸烟者 & 吸烟者 & 总计 \\
\cline{1-4}
不患慢性气管炎者 & 121 & $b$ & 283 \\
\cline{1-4}
患慢性气管炎者 & C & $d$ & \phantom{X} \\
\cline{1-4}
总计 & 134 & \phantom{X} & 339 \\
\cline{1-4}
\hline
\end{tabular}
\end{center}

(1)估算样本中吸烟者中患慢性支气管炎的百分比；

(2)有多少把握认为患慢性支气管炎与吸烟有关？

附: ${\chi }^{2} = \dfrac{n{\left( ad - bc\right) }^{2}}{\left( {a + b}\right) \left( {c + d}\right) \left( {a + c}\right) \left( {b + d}\right) }$ ，其中 $n = a + b + c + d$ ， $P\left( {{\chi }^{2} \geq  {6.635}}\right)  \approx  {0.01}$ ， $P\left( {{\chi }^{2} \geq  }\right.$

${5.024}) \approx  {0.025},P\left( {{\chi }^{2} \geq  {3.841}}\right)  \approx  {0.05},P\left( {{\chi }^{2} \geq  {2.706}}\right)  \approx  {0.1}$.

4.(2024$\cdot$奉贤区二模$\cdot$18)某学生兴趣小组随机调查了某市 100 天中每天的空气质量等级和当天到某公园锻炼的人次, 整理数据如下表所示 (单位:天):

\begin{center}

\begin{tabular}{|c|c|c|c|}
\hline
\makecell{ 锻炼人次 \\ 空气质量等级 } & [0,200] & (200,400) & (400,600] \\
\cline{1-4}
1(优) & 3 & 18 & 25 \\
\cline{1-4}
2(良) & 6 & $x$ & 14 \\
\cline{1-4}
3(轻度污染) & 5 & 5 & 6 \\
\cline{1-4}
4(中度污染) & 6 & 3 & 0 \\
\cline{1-4}
\hline
\end{tabular}
\end{center}

(1)求一天中到该公园锻炼的平均人次的估计值; (同一组中的数据用该组区间的中点值为代表)

(2)若某天的空气质量等级为 1 或 2，则称这天“空气质量好”；若某天的空气质量等级为 3 或 4，则称这天“空气质量不好”. 根据所给数据,完成下面的 $2 \times  2$ 列联表,请根据表中的数据判断: 一天中到该公园锻炼的人次是否与该市当天的空气质量有关? (规定显著性水平 $\alpha  = {0.05}$ )

\begin{center}

\begin{tabular}{|c|c|c|c|}
\hline
\phantom{X} & 人次<400 & 人次>400 & 总计 \\
\cline{1-4}
空气质量好 & \phantom{X} & \phantom{X} & \phantom{X} \\
\cline{1-4}
空气质量不好 & \phantom{X} & \phantom{X} & \phantom{X} \\
\cline{1-4}
总计 & \phantom{X} & \phantom{X} & \phantom{X} \\
\cline{1-4}
\hline
\end{tabular}
\end{center}

附: ${\chi }^{2} = \dfrac{n{\left( ad - bc\right) }^{2}}{\left( {a + b}\right) \left( {c + d}\right) \left( {a + c}\right) \left( {b + d}\right) },P\left( {{\chi }^{2} \geq  {3.841}}\right)  \approx  {0.05}$.

5.(2024$\cdot$杨浦区二模$\cdot$19)某工厂为提高生产效率，开展技术创新活动，提出了完成某项生产任务的两种新的生产方式. 为比较两种生产方式的效率, 选取 40 名工人, 将他们随机分成两组, 每组 20 人，第一组工人用第一种生产方式，第二组工人用第二种生产方式. 完成生产任务的工作时间不超过 70 min 的工人为“优秀”，否则为“合格”. 根据工人完成生产任务的工作时间(单位:min)绘制了如图所示茎叶图:

\begin{center}

\begin{tabular}{|c|c|c|c|c|c|c|c|c|c|}
\hline
第一种生产方式 & \phantom{X} & \multicolumn{8}{|c|}{第二种生产方式} \\
\cline{1-10}
3 1 & 6 & 1 & 2 & 3 & 5 & \phantom{X} & 7 & 8 & 9 \\
\cline{1-10}
5 2 1 & 7 & 0 & 0 & 2 & \phantom{X} & 4 & 7 & 8 & \phantom{X} \\
\cline{1-10}
4 4 3 2 2 1 & 8 & 4 & 7 & 7 & \phantom{X} & \phantom{X} & \phantom{X} & \phantom{X} & \phantom{X} \\
\cline{1-10}
4 4 3 2 2 1 1 0 & 9 & 0 & 1 & \phantom{X} & \phantom{X} & \phantom{X} & \phantom{X} & \phantom{X} & \phantom{X} \\
\cline{1-10}
\hline
\end{tabular}
\end{center}

第 5 题图

(1)求 40 名工人完成生产任务所需时间的第 75 百分数;

(2)独立地从两种生产方式中各选出一个人，求选出的两个人均为优秀的概率;

(3)根据工人完成生产任务的工作时间，两种生产方式优秀与合格的人数填入下面的 $2 \times  2$ 列联表.

\begin{center}

\begin{tabular}{|c|c|c|c|}
\hline
\phantom{X} & 第一种生产方式 & 第二种生产方式 & 总计 \\
\cline{1-4}
优秀 & \phantom{X} & \phantom{X} & \phantom{X} \\
\cline{1-4}
合格 & \phantom{X} & \phantom{X} & \phantom{X} \\
\cline{1-4}
总计 & \phantom{X} & \phantom{X} & \phantom{X} \\
\cline{1-4}
\hline
\end{tabular}
\end{center}

根据上面的 $2 \times  2$ 列联表,判断能否有 ${95}\%$ 的把握认为两种生产方式的工作效率有显著差异. $\left( {{\chi }^{2} = \dfrac{n{\left( ad - bc\right) }^{2}}{\left( {a + b}\right) \left( {c + d}\right) \left( {a + c}\right) \left( {b + d}\right) }}\right.$. 其中 $n = a + b + c + d,P\left( {{\chi }^{2} \geq  {3.841}}\right)  \approx  {0.05})$.

6.(2024$\cdot$上海秋季卷$\cdot$19)为了解某地初中学生体育锻炼时长与学业成绩的关系，从该地区 29000 名学生中抽取 580 人，得到日均体育锻炼时长与学业成绩的数据如下表所示:

\begin{center}

\begin{tabular}{|c|c|c|c|c|c|}
\hline
时间范围学业成绩 & $\lbrack 0,{0.5})$ & [0.5,1) & [1,1.5) & $\lbrack {1.5},2)$ & $\lbrack 2,{2.5})$ \\
\cline{1-6}
优秀 & 5 & 44 & 42 & 3 & 1 \\
\cline{1-6}
不优秀 & 134 & 147 & 137 & 40 & 27 \\
\cline{1-6}
\hline
\end{tabular}
\end{center}

(1)该地区 29000 名学生中体育锻炼时长不少于 1 h 人数约为多少？

(2)估计该地区初中学生日均体育锻炼的时长;(精确到 0.1 )

(3)是否有 95\%的把握认为学业成绩优秀与日均体育锻炼时长不小于 1 h 且小于 2 h 有关？(附: ${\chi }^{2} = \dfrac{n{\left( ad - bc\right) }^{2}}{\left( {a + b}\right) \left( {c + d}\right) \left( {a + c}\right) \left( {b + d}\right) }$ ,其中 $n = a + b + c + d,P\left( {{\chi }^{2} \geq  {3.841}}\right)  \approx  {0.05}.)$

7.(2025$\cdot$虹口区二模$\cdot$19)已知某区组建了一支 120 人的志愿者队伍，并由其中 72 人组成“志愿模范队”. 经过一年的实践，全队共有 72 人的周平均服务时长超过 2 h，其中有 54 人来自“志愿模范队”，如下表所示.

\begin{center}

\begin{tabular}{|c|c|c|c|}
\hline
\phantom{X} & 是“志愿模范队”成员 & 不是“志愿模范队”成员 & 总计 \\
\cline{1-4}
周平均服务时长超过 $2{h}$ & 54 & \phantom{X} & 72 \\
\cline{1-4}
周平均服务时长不超过 $2{h}$ & \phantom{X} & \phantom{X} & \phantom{X} \\
\cline{1-4}
总计 & 72 & \phantom{X} & 120 \\
\cline{1-4}
\hline
\end{tabular}
\end{center}

(1)已知一名志愿者是“志愿模范队”成员，求其周平均服务时长超过 2 h 的概率；

(2)请完成 $2 \times  2$ 列联表,并根据表中数据回答: 是否有 99.9\% 的把握认为 “是‘志愿模范队’成员” 与“周平均服务时长超过 $2{h}$ ”有关系?

(3)现从周平均服务时长超过 $2{h}$ 的人员中按照是否为“志愿模范队”成员进行分层抽样，选取 8 人组建“志愿突击队”，并从这 8 人中再随机选取 2 人做深度访谈，记随机变量 $X$ 为这 2 人中来自于“志愿模范队”的人数，求 $X$ 的分布与方差.

附录: ${\chi }^{2} = \dfrac{n{\left( ad - bc\right) }^{2}}{\left( {a + b}\right) \left( {c + d}\right) \left( {a + c}\right) \left( {b + d}\right) }$ ,其中 $n = a + b + c + d$.

\begin{center}

\begin{tabular}{|c|c|c|c|c|}
\hline
$P\left( {{\chi }^{2} \geq  k}\right)$ & 0.100 & 0.050 & 0.010 & 0.001 \\
\cline{1-5}
$k$ & 2.706 & 3.841 & 6.635 & 10.828 \\
\cline{1-5}
\hline
\end{tabular}
\end{center}

8.(2024$\cdot$奉贤区一模$\cdot$19)某连锁便利店从 2014 年到 2018 年销售商品品种为 2000 种，从 2019 年开始, 该便利店进行了全面升级, 销售商品品种为 3000 种. 下表中列出了从 2014 年到 2023 年的利润额.

\begin{center}

\begin{tabular}{|c|c|c|c|c|c|c|c|c|c|c|}
\hline
年份 $x$ & 2014 & 2015 & 2016 & 2017 & 2018 & 2019 & 2020 & 2021 & 2022 & 2023 \\
\cline{1-11}
利润额 $y/$ 万元 & 27.6 & 42.0 & 38.4 & 48.0 & 63.6 & 63.7 & 72.8 & 80.1 & 60.5 & 99.3 \\
\cline{1-11}
\hline
\end{tabular}
\end{center}

(1)若某年的利润额超过 45.0 万元，则该便利店当年会被评选为示范店；若利润额不超过 45.0 万元,则该便利店当年不会被评选为示范店. 试完成 $2 \times  2$ 列联表,并判断商品品种数量与便利店是否为示范店有关? (显著性水平 $\alpha  = {0.05},P\left( {{\chi }^{2} \geq  {3.841}}\right)  \approx  {0.05}$ )

\begin{center}

\begin{tabular}{|c|c|c|c|}
\hline
\phantom{X} & 品种为 2000 种 & 品种为 3000 种 & 总计 \\
\cline{1-4}
被评为示范店次数 & \phantom{X} & \phantom{X} & \phantom{X} \\
\cline{1-4}
未被评为示范店次数 & \phantom{X} & \phantom{X} & \phantom{X} \\
\cline{1-4}
总计 & \phantom{X} & \phantom{X} & \phantom{X} \\
\cline{1-4}
\hline
\end{tabular}
\end{center}

(2)请根据 2014 年至 2023 年(剔除 2022 年的数据)的数据建立 $y$ 与 $x$ 的线性回归模型 \ding{172} ；根据 2019 年至 2023 年的数据建立 $y$ 与 $x$ 的线性回归模型 \ding{173}. 分别用这两个模型，预测 2024 年该便利店的利润额并说明这样的预测值是否可靠? (回归系数精确到 0.001 , 利润精确到 0.1 万元) 回归系数 $\widehat{a}$ 与 $\widehat{b}$ 的公式如下:

\[
\widehat{a} = \dfrac{\mathop{\sum }\limits_{{i = 1}}^{n}\left( {{x}_{i} - \overline{x}}\right) \left( {{y}_{i} - \overline{y}}\right) }{\mathop{\sum }\limits_{{i = 1}}^{n}{\left( {x}_{i} - \overline{x}\right) }^{2}} = \dfrac{\mathop{\sum }\limits_{{i = 1}}^{n}{x}_{i}{y}_{i} - n\overline{x}\overline{y}}{\mathop{\sum }\limits_{{i = 1}}^{n}{x}_{i}^{2} - n{\overline{x}}^{2}},\widehat{b} = \overline{y} - \widehat{a}\overline{x} = \dfrac{\mathop{\sum }\limits_{{i = 1}}^{n}{y}_{i} - \widehat{a}\mathop{\sum }\limits_{{i = 1}}^{n}{x}_{i}}{n}
\]

\subsection*{考点11、> 离散型分布及期望、方差}

\subsubsection*{一、填空题}

1. (2025$\cdot$静安区二模$\cdot$4)已知随机变量 $X$ 服从二项分布 $B \sim  \left( {n,p}\right)$ ，若 $E\left\lbrack  X\right\rbrack   = {30}$ ， $D\left\lbrack  X\right\rbrack   = {20}$ ，则 $p$ 的值为\blkda[2]{}.

2.(2024$\cdot$虹口区二模$\cdot$5)已知随机变量 $X \sim  B\left( {{50},p}\right)$ ，且 $E\left\lbrack  X\right\rbrack   = {20}$ ，则 $D\left\lbrack  X\right\rbrack   =$ \blkda[2]{}.

3.(2025$\cdot$上海秋季卷$\cdot$6)已知随机变量 $X$ 的分布为 $\left( \begin{matrix} 5 & 6 & 7 \\  {0.2} & {0.3} & {0.5} \end{matrix}\right)$ ，则期望 $E\left\lbrack  X\right\rbrack   =$ \blkda[2]{}.

4.(2024$\cdot$徐汇区二模$\cdot$9)同时抛掷三枚相同的质地均匀硬币，设随机变量X=1 表示结果中有正面朝上， $X = 0$ 表示结果中没有正面朝上，则 $D\left\lbrack  X\right\rbrack   =$ \blkda[2]{}.

\subsubsection*{二、选择题}

5.(2025$\cdot$黄浦区二模$\cdot$15)设 ${x}_{1} < {x}_{2} < {x}_{3} < {x}_{4}$ ，随机变量 $X$ 取值 ${x}_{1}$ 、 ${x}_{2}$ 、 ${x}_{3}$ 、 ${x}_{4}$ 的概率均为 0.25，随机变量 ${X}_{1}$ 取值 $\dfrac{{x}_{1} + {x}_{2}}{2}\text{ 、 }\dfrac{{x}_{2} + {x}_{3}}{2}\text{ 、 }\dfrac{{x}_{3} + {x}_{4}}{2}\text{ 、 }\dfrac{{x}_{4} + {x}_{1}}{2}$ 的概率也均为 0.25,随机变量 ${X}_{2}$ 取值 $2{x}_{1}$ - ${x}_{2}\text{ 、 }2{x}_{2} - {x}_{3}\text{ 、 }2{x}_{3} - {x}_{4}\text{ 、 }2{x}_{4} - {x}_{1}$ 的概率也均为 0.25. 若记 $D\left\lbrack  {X}_{1}\right\rbrack  \text{ 、 }D\left\lbrack  {X}_{2}\right\rbrack$ 分别为 ${X}_{1}\text{ 、 }{X}_{2}$ 的方差, 则 \dotfill{(\kh{})}

\xx{}{}{}{}$D\left\lbrack  {X}_{1}\right\rbrack   < D\left\lbrack  {X}_{2}\right\rbrack$

B. $D\left\lbrack  {X}_{1}\right\rbrack   = D\left\lbrack  {X}_{2}\right\rbrack$

C. $D\left\lbrack  {X}_{1}\right\rbrack   > D\left\lbrack  {X}_{2}\right\rbrack$

D. $D\left\lbrack  {X}_{1}\right\rbrack$ 与 $D\left\lbrack  {X}_{2}\right\rbrack$ 的大小关系与 ${x}_{1}\text{ 、 }{x}_{2}\text{ 、 }{x}_{3}\text{ 、 }{x}_{4}$ 的取值有关

6.(2024$\cdot$虹口区二模$\cdot$15)给出下列 4 个命题:

 \ding{172}  若事件 $A$ 和事件 $B$ 互斥，则 $P\left( {A \cap  B}\right)  = P\left( A\right) P\left( B\right)$ ；

 \ding{173}  数据2、3、6、7、8、10、11、13的第 70 百分位数为 10 ;

 \ding{174} 已知 $y$ 关于 $x$ 的回归方程为 $y =  - {0.5x} + {0.7}$ ，则样本点 $\left( {2, - 1}\right)$ 的离差为 -0.7；

 \ding{175}  随机变量 $X$ 的分布为 $\left( \begin{matrix} 0 & 1 & 2 & 3 \\  {0.2} & {0.2} & {0.3} & {0.3} \end{matrix}\right)$ ,则其数学期望 $E\left\lbrack  X\right\rbrack   = {1.6}$.

其中正确命题的序号为 \dotfill{(\kh{})}

\xx{}{}{}{} \ding{172}  \ding{173} 

B.  \ding{172}  \ding{174} 

C.  \ding{173}  \ding{174} 

D.  \ding{173}  \ding{175} 

\subsubsection*{三、解答题}

7.(2024$\cdot$宝山区二模$\cdot$19)在课外活动中，甲、乙两名同学进行投篮比赛，每人投 3 次，每投进一次得 2 分,否则得 0 分. 已知甲每次投进的概率为 $\dfrac{1}{2}$ ,且每次投篮相互独立; 乙第一次投篮,投进的概率为 $\dfrac{1}{2}$ ，从第二次投篮开始，若前一次投进，则该次投进的概率为 $\dfrac{3}{5}$ ，若前一次没投进，则该次投进的概率为 $\dfrac{2}{5}$.

(1)求甲投篮 3 次得 2 分的概率；

(2)若乙投篮 3 次得分为 $X$ ，求 $X$ 的分布和期望；

(3)比较甲、乙的比赛结果.

8.(2024$\cdot$闵行区二模$\cdot$19)ChatGPT 是 OpenAI 研发的一款聊天机器人程序，是人工智能技术驱动的自然语言处理工具，它能够基于在预训练阶段所见的模式和统计规律来生成回答，但它的回答可能会受到训练数据信息的影响，不一定完全正确. 某科技公司在使用 ChatGPT 对某一类问题进行测试时发现，如果输入的问题没有语法错误，它回答正确的概率为 0.98 ；如果出现语法错误，它回答正确的概率为 0.18. 假设每次输入的问题出现语法错误的概率为 0.1 ，且每次输入问题， ChatGPT 的回答是否正确相互独立. 该公司科技人员小张想挑战一下 ChatGPT, 小张和 ChatGPT 各自从给定的 10 个问题中随机抽取 9 个作答, 已知在这 10 个问题中, 小张能正确作答其中的 9 个.

(1)求小张能全部回答正确的概率；

(2)求一个问题能被 ChatGPT 回答正确的概率；

(3)在这轮挑战中，分别求出小张和 ChatGPT 答对题数的期望与方差.

9.(2024$\cdot$徐汇区二模$\cdot$19)为了解中草药甲对某疾病的预防效果，研究人员随机调查了 100 名人员， 调查数据如表所示. (单位:个)

\begin{center}

\begin{tabular}{|c|c|c|c|}
\hline
\phantom{X} & 未患病者 & 患病者 & 合计 \\
\cline{1-4}
未服用中草药甲 & 29 & 16 & 45 \\
\cline{1-4}
服用中草药甲 & 46 & 9 & 55 \\
\cline{1-4}
合计 & 75 & 25 & 100 \\
\cline{1-4}
\hline
\end{tabular}
\end{center}

(1)若规定显著性水平 $\alpha  = {0.05}$ ，试分析中草药甲对预防此疾病是否有效；

(2)已知中草药乙对该疾病的治疗有效率数据如下:对未服用过中草药甲的患者治疗有效率为 $\dfrac{1}{2}$ ， 对服用过中草药甲的患者治疗有效率为 $\dfrac{3}{4}$. 若用频率估计概率，现从患此疾病的人员中随机选取 2 人(分两次选取，每次 1 人，两次选取的结果独立)使用中草药乙进行治疗，记治疗有效的人数为 $X$ ,求 $X$ 的分布和数学期望.

附: ${\chi }^{2} = \dfrac{n{\left( ad - bc\right) }^{2}}{\left( {a + b}\right) \left( {c + d}\right) \left( {a + c}\right) \left( {b + d}\right) },n = a + b + c + d$.

\begin{center}

\begin{tabular}{|c|c|c|c|c|}
\hline
$P\left( {{\chi }^{2} \geq  k}\right)$ & 0.100 & 0.050 & 0.010 & 0.001 \\
\cline{1-5}
$k$ & 2.706 & 3.841 & 6.635 & 10.828 \\
\cline{1-5}
\hline
\end{tabular}
\end{center}

10.(2024$\cdot$虹口区二模$\cdot$19)某企业监控汽车零件的生产过程，现从汽车零件中随机抽取 100 件作为样本，测得质量差(零件质量与标准质量之差的绝对值)的样本数据如下表:

\begin{center}

\begin{tabular}{|c|c|c|c|c|c|}
\hline
质量差(单位:mg) & 54 & 57 & 60 & 63 & 66 \\
\cline{1-6}
件数(单位:件) & 5 & 21 & 46 & 25 & 3 \\
\cline{1-6}
\hline
\end{tabular}
\end{center}

(1)求样本质量差的平均数 $\overline{x}$ ; 假设零件的质量差 $X \sim  N\left( {\mu ,{\sigma }^{2}}\right)$ ，其中 ${\sigma }^{2} = {16}$ ，用 $\overline{x}$ 作为 $\mu$ 的近似值，求 $P\left( {{56} < X < {68}}\right)$ 的值；

(2)已知该企业共有两条生产汽车零件的生产线，其中全部零件的 $\dfrac{3}{4}$ 来自第 1 条生产线. 若两条生产线的废品率分别为 0.016 和 0.012 , 且这两条生产线是否产出废品是相互独立的. 现从该企业生产的汽车零件中随机抽取一件.

(i) 求抽取的零件为废品的概率;

(ii) 若抽取出的零件为废品,求该废品来自第 1 条生产线的概率.

参考数据: 若随机变量 $X \sim  N\left( {\mu ,{\sigma }^{2}}\right)$ ,则 $P\left( {\mu  - \sigma  < X \leq  \mu  + \sigma }\right)  \approx  {0.6827},P\left( {\mu  - {2\sigma } < X \leq  \mu  + {2\sigma }}\right)  \approx  {0.9545},P\left( {\mu  - {3\sigma } < X \leq  \mu  + {3\sigma }}\right)  \approx  {0.9973}.$

11.(2024$\cdot$嘉定区二模$\cdot$19)据文化和旅游部发布的数据显示，2023 年国内出游人次达 48.91 亿次， 总花费 4.91 万亿元. 人们选择的出游方式不尽相同，有自由行，也有跟团游. 为了了解年龄因素是否影响出游方式的选择，我们按年龄将成年人群分为青壮年组(大于等于 14 岁，小于 40 岁)和中老年组(大于等于 40 岁). 现在 $S$ 市随机抽取 170 名成年市民进行调查，得到如表所示的数据.

\begin{center}

\begin{tabular}{|c|c|c|c|}
\hline
\phantom{X} & 青壮年 & 中老年 & 合计 \\
\cline{1-4}
自由行 & 60 & 40 & \phantom{X} \\
\cline{1-4}
跟团游 & 20 & 50 & \phantom{X} \\
\cline{1-4}
合计 & \phantom{X} & \phantom{X} & \phantom{X} \\
\cline{1-4}
\hline
\end{tabular}
\end{center}

(1)请补充 $2 \times  2$ 列联表，并判断能否有 95\%的把握认为年龄与出游方式的选择有关；

(2)用分层抽样的方式从跟团游中抽取 14 个人，再从 14 个人中随机抽取 7 个人，用随机变量 $X$ 表示这 7 个人中中老年与青壮年人数之差的绝对值,求 $X$ 的分布和数学期望.

附: ${\chi }^{2} = \dfrac{n{\left( ad - bc\right) }^{2}}{\left( {a + b}\right) \left( {c + d}\right) \left( {a + c}\right) \left( {b + d}\right) }$ ，其中 $n = a + b + c + d$.

\begin{center}

\begin{tabular}{|c|c|c|c|}
\hline
$P\left( {{\chi }^{2} \geq  k}\right)$ & 0.10 & 0.05 & 0.025 \\
\cline{1-4}
$k$ & 2.706 & 3.841 & 5.024 \\
\cline{1-4}
\hline
\end{tabular}
\end{center}

12.(2024$\cdot$金山区二模$\cdot$19)有标号依次为1、2、...、n(n≥2，n∈N)的n个盒子，标号为1号的盒子里有 3 个红球和 3 个白球, 其余盒子里都是 1 个红球和 1 个白球. 现从 1 号盒子里取出 2 个球放入 2 号盒子,再从 2 号盒子里取出 2 个球放入 3 号盒子,……,依次进行到从 $n - 1$ 号盒子里取出 2 个球放入 $n$ 号盒子为止.

(1)当 $n = 2$ 时,求 2 号盒子里有 2 个红球的概率;

(2)设 $n$ 号盒子中红球个数为随机变量 ${X}_{n}$ ，求 ${X}_{3}$ 的分布及 $E\left\lbrack  {X}_{3}\right\rbrack$ ，并猜想 $E\left\lbrack  {X}_{n}\right\rbrack$ 的值. (见面)的证明此猜想)

13.(2023$\cdot$上海秋季卷$\cdot$19)已知某汽车模型公司共有 25 个汽车模型, 其外观和内饰的颜色分布如下表所示:

(1)若小明从这些模型中随机拿一个模型，记事件 $A$ 为 “小明取到红色外观的模型”,事件 $B$ 为 “小明取到米色内饰的模型”. 求 $P\left( B\right)$ 和 $P\left( {B \mid  A}\right)$ ,并据此判断事件 $A$ 和事件 $B$ 是否独立;

\begin{center}

\begin{tabular}{|c|c|c|}
\hline
\phantom{X} & 红色外观 & 蓝色外观 \\
\cline{1-3}
棕色内饰 & 8 & 12 \\
\cline{1-3}
米色内饰 & 2 & 3 \\
\cline{1-3}
\hline
\end{tabular}
\end{center}

(2)该公司举行了一个抽奖活动，规定在一次抽奖中，每人可以一次性从这些模型中掌网个代车模型，给出以下假设:

假设1: 拿到的两个模型会出现三种结果，即外观和内饰均为同色、外观和内饰都异色，以及仅外观或仅内饰同色;

假设 2: 按结果的可能性大小,概率越小奖项越高;

假设 3: 该抽奖活动的奖金额为:一等奖 600 元、二等奖 300 元、三等奖 150 元;

请你分析奖项对应的结果,设 $X$ 为奖金额,写出 $X$ 的分布,并求出 $X$ 的数字期望.

14.(2024$\cdot$松江区二模$\cdot$19)某素质训练营设计了一项闯关比赛. 规定:三人组队参赛，每次只派一个人，且每人只派一次；如果一个人闯关失败，再派下一个人重新闯关；三人中只要有人闯关成功即视作比赛胜利，无需继续闯关. 现有甲、乙、丙三人组队参赛，他们各自闯关成功的概率分别为 ${p}_{1}\text{ 、 }{p}_{2}\text{ 、 }{p}_{3}$ ,假定 ${p}_{1}\text{ 、 }{p}_{2}\text{ 、 }{p}_{3}$ 互不相等,且每人能否闯关成功的事件相互独立.

(1)计划依次派甲、乙、丙进行闯关，若 ${p}_{1} = \dfrac{3}{4}$ 、 ${p}_{2} = \dfrac{2}{3}$ 、 ${p}_{3} = \dfrac{1}{2}$ ，求该小组比赛胜利的概率；

(2)若依次派甲、乙、丙进行闯关，则写出所需派出的人员数目 $X$ 的分布，并求 $X$ 的期望 $E\left\lbrack  X\right\rbrack$ ；

(3)已知 $1 > {p}_{1} > {p}_{2} > {p}_{3}$ ，若乙只能安排在第二个派出，要使派出人员数目的期望较小，试确定甲、丙谁先派出.

15.(2024$\cdot$普陀区二模$\cdot$19)张先生每周有 5 个工作日，工作日出行采用自驾方式，必经之路上有一个十字路口,直行车道有三条,直行车辆可以随机选择一条车道通行,记事件 $A$ 为 “张先生驾车从左侧直行车道通行”.

(1)某日张先生驾车上班接近路口时，看到自己车前是一辆大货车，遂选择不与大货车从同一车道通行,记事件 $B$ 为 “大货车从中间直行车道通行”,求 $P\left( {A \cap  B}\right)$ ;

(2)用 $X$ 表示张先生每周工作日出行事件 $A$ 发生的次数，求 $X$ 的分布及期望 $E\left\lbrack  X\right\rbrack$.

16.(2025$\cdot$崇明区二模$\cdot$19)某区 2025 年 3 月 31 日至 4 月 13 日的天气预报如图所示.

\begin{center}

\begin{tabular}{|c|c|c|c|c|c|c|}
\hline
\makecell{ 31 \\ \includegraphics[width=0.2\textwidth]{bodyxb/src/262_223_302_92_87_0.jpg} \\ 大部分晴 17/9  \textdegree C } & \makecell{ 04月01日 \\ \includegraphics[width=0.2\textwidth]{bodyxb/src/262_387_309_94_83_0.jpg} \\ 阵雨 18/9  \textdegree C } & \makecell{ 02 \\ \includegraphics[width=0.2\textwidth]{bodyxb/src/262_554_307_95_85_0.jpg} \\ 阵雨 18/10  \textdegree C } & \makecell{ 03 \\ \includegraphics[width=0.2\textwidth]{bodyxb/src/262_722_305_95_85_0.jpg} \\ 阵雨 19/10  \textdegree C } & \makecell{ 04 \\ \includegraphics[width=0.2\textwidth]{bodyxb/src/262_898_300_80_79_0.jpg} \\ 多云 18/9  \textdegree C } & \makecell{ 05 \\ \includegraphics[width=0.2\textwidth]{bodyxb/src/262_1059_301_95_87_0.jpg} \\ 间歇性多云 19/10  \textdegree C } & \makecell{ 06 \\ \includegraphics[width=0.2\textwidth]{bodyxb/src/262_1234_297_81_79_0.jpg} \\ 多云 20/10  \textdegree C } \\
\cline{1-7}
\makecell{ 07 \\ \includegraphics[width=0.2\textwidth]{bodyxb/src/262_220_561_95_82_0.jpg} \\ 阵雨 19/10  \textdegree C } & \makecell{ 08 \\ \includegraphics[width=0.2\textwidth]{bodyxb/src/262_393_554_82_76_0.jpg} \\ 大部分多云 19/10  \textdegree C } & \makecell{ 09 \\ \includegraphics[width=0.2\textwidth]{bodyxb/src/262_555_553_95_86_0.jpg} \\ 大部分晴 20/8  \textdegree C } & \makecell{ 10 \\ \includegraphics[width=0.2\textwidth]{bodyxb/src/262_733_551_80_76_0.jpg} \\ 多云 18/8  \textdegree C } & \makecell{ 11 \\ \includegraphics[width=0.2\textwidth]{bodyxb/src/262_895_551_92_84_0.jpg} \\ 大部分晴 18/10  \textdegree C } & \makecell{ 12 \\ \includegraphics[width=0.2\textwidth]{bodyxb/src/262_1055_553_94_83_0.jpg} \\ 阵雨 17/8  \textdegree C } & \makecell{ 13 \\ \includegraphics[width=0.2\textwidth]{bodyxb/src/262_1224_552_94_82_0.jpg} \\ 阵雨 17/9  \textdegree C } \\
\cline{1-7}
\hline
\end{tabular}
\end{center}

第 16 题图

(1)从3月31日至4月13日某天开始，连续统计三天，求这三天中至少有两天是阵雨的概率；

(2)根据天气预报，该区 4 月 14 日的最低气温是 ${9}^{ \circ  }{C}$ ，温差是指一段时间内最高温度与最低温度之间的差值，例如 3 月 31 日的最高温度为 ${17}^{ \circ  }{C}$ ，最低温度为 ${9}^{ \circ  }{C}$ ，当天的温差为 ${8}^{ \circ  }{C}$. 记4月 1 日至 4 日这 4 天温差的方差为 ${s}_{1}^{2},4$ 月 11 日至 14 日这 4 天温差的方差为 ${s}_{2}^{2}$ ,若 ${s}_{2}^{2} = \dfrac{4}{3}{s}_{1}^{2}$ ,求 4 月 14 日天气预报的最高气温;

(3)从3月31日至4月13日中随机抽取两天,用 $X$ 表示一天温差不高于 ${9}^{ \circ  }{C}$ 的天数，求 $X$ 的分布及期望.

17.(2024$\cdot$青浦区二模$\cdot$19)垃圾分类能减少有害垃圾对环境的破坏，同时能提高资源循环利用的效率. 目前上海社区的垃圾分类基本采用四类分类法, 即干垃圾、湿垃圾、可回收垃圾与有害垃圾. 某校为调查学生对垃圾分类的了解程度,随机抽取 100 名学生作为样本,按照了解程度分为 $A$ 等级和 $B$ 等级,得到如下列联表:

(1)根据表中的数据回答:学生对垃圾分类的了解程度是否与性别有关？(规定:显著性水平 $\alpha  = {0.05})$ 附: ${\chi }^{2} = \dfrac{n{\left( ad - bc\right) }^{2}}{\left( {a + b}\right) \left( {c + d}\right) \left( {a + c}\right) \left( {b + d}\right) }$ , 其中 $n = a + b + c + d,P\left( {{\chi }^{2} \geq  {3.841}}\right)  \approx  {0.05}$.

\begin{center}

\begin{tabular}{|c|c|c|c|}
\hline
\phantom{X} & 男生 & 女生 & 总计 \\
\cline{1-4}
$A$ 等级 & 40 & 20 & 60 \\
\cline{1-4}
$B$ 等级 & 20 & 20 & 40 \\
\cline{1-4}
总计 & 60 & 40 & 100 \\
\cline{1-4}
\hline
\end{tabular}
\end{center}

(2)为进一步加强垃圾分类的宣传力度，学校特举办垃圾分类知识问答比赛. 每局比赛由二人参加,主持人 $A$ 和 $B$ 轮流提问,先赢 3 局者获得奖项并结束比赛. 甲、乙两人参加比赛,已知主持人 $A$ 提问甲赢的概率为 $\dfrac{2}{3}$ ，主持人 $B$ 提问甲赢的概率为 $\dfrac{1}{2}$ ，每局比赛互相独立，且每局都分输赢. 现抽签决定第一局由主持人 A 提问.

(i) 求比赛只进行 3 局就结束的概率;

(ii) 设 $X$ 为结束比赛时甲赢的局数,求 $X$ 的分布和数学期望 $E\left\lbrack  X\right\rbrack$.

18.(2025$\cdot$金山区二模$\cdot$19)为了研究高三学生每天整理数学错题的情况，某校数学建模兴趣小组的同学在本校高三年级学生中采用随机抽样的方法抽取了 40 名学生，调查他们平时的数学成绩和整理数学错题的情况，现统计得到的部分数据如下表所示:

\begin{center}

\begin{tabular}{|c|c|c|c|}
\hline
\phantom{X} & 数学成绩总评优秀人数 & 数学成绩总评非优秀人数 & 合计 \\
\cline{1-4}
每天都整理数学错题人数 & 14 & \phantom{X} & \phantom{X} \\
\cline{1-4}
不是每天都整理数学错题人数 & \phantom{X} & 15 & 20 \\
\cline{1-4}
合计 & \phantom{X} & \phantom{X} & 40 \\
\cline{1-4}
\hline
\end{tabular}
\end{center}

(1)完成上述样本数据的 $2 \times  2$ 列联表，并计算:每天都整理数学错题且数学成绩总评优秀的经验概率;

(2)是否有 99\% 的把握认为“数学成绩总评优秀与每天都整理数学错题有关”？

附: ${\chi }^{2} = \dfrac{n{\left( ad - bc\right) }^{2}}{\left( {a + b}\right) \left( {c + d}\right) \left( {a + c}\right) \left( {b + d}\right) }$ ;

\begin{center}

\begin{tabular}{|c|c|c|c|}
\hline
$P\left( {{\chi }^{2} \geq  k}\right)$ & 0.10 & 0.01 & 0.001 \\
\cline{1-4}
$k$ & 2.706 & 6.635 & 10.828 \\
\cline{1-4}
\hline
\end{tabular}
\end{center}

(3)从不是每天都整理数学错题的学生中随机抽取 3 名学生做进一步访谈，设恰好抽取到数学成绩总评优秀的人数为 $X$ ,求 $X$ 的分布和期望.

19.(2025$\cdot$闵行区二模$\cdot$19)某社团共有 12 名成员，其中高一男生2人、女生4人，高二男生3人、女生 3 人. 现从中随机抽选 2 人参加数学知识问答.

(1)若逐个抽选，求恰好第一个抽选的是男生的概率；

(2)若恰好抽选了 1 名男生与 1 名女生，求这 2 人都是高二学生的概率；

(3)若恰好抽选了 1 名高一学生与 1 名高二学生，记抽选出来的男生与女生的人数之差的绝对值为 $X$ ,求 $X$ 的分布与数学期望 $E\left\lbrack  X\right\rbrack$.

20. (2025$\cdot$浦东新区二模$\cdot$19)为测试 A、B 两款人工智能软件解答数学问题的能力，将 100 道难度相当的数学试题从 1 到 100 编号后随机分配给这两款软件测试. 每道试题只被一款软件解答一次,并记录结果如下:

\begin{center}

\begin{tabular}{|c|c|c|c|c|}
\hline
\multirow{2}{*}{试题类别} & \multicolumn{2}{|c|}{$A$ 软件} & \multicolumn{2}{|c|}{$B$ 软件} \\
\cline{2-5}
 & 测试试题数量 & 正确解答的数量 & 测试试题数量 & 正确解答的数量 \\
\cline{1-5}
几何试题 & 20 & 16 & 30 & 20 \\
\cline{1-5}
函数试题 & 30 & 24 & 20 & 18 \\
\cline{1-5}
\hline
\end{tabular}
\end{center}

(1)分别估计 $A$ 软件、 $B$ 软件能正确解答数学问题的概率；

(2)小浦准备用这两款软件来解决某次数学测试中的第 12 题(假设其难度和测试的 100 道题基本相同)，但该题内容还未知，从以往情况来看，该题是几何题的概率为 $\dfrac{1}{3}$ ，是函数题的概率为 $\dfrac{2}{3}$. 将频率视为概率,试通过计算来说明小浦应该用哪款软件解决这道试题?

(3)小浦决定采用这两款软件解答 6 道类似试题，其中几何、函数各 3 道，每道试题只用其中一款软件解答一次. 将频率视为概率, 小浦比较了这两款软件在解答几何和函数题上的正确率, 决定用表现较好的那款软件解决其擅长的题型. 用 ${X}_{1}\text{ 、 }{X}_{2}$ 分别表示这 3 道几何试题与 3 道函数试题被正确解答的个数,求随机变量 ${X}_{1} + {X}_{2}$ 的数学期望和方差.

21.(2025$\cdot$普陀区二模$\cdot$19)某区为推进教育数字化转型, 通过聚合区域学校的教育资源, 依托 AI 技术搭建了区域智慧题库系统，形成了“ $A$ 通识过关一 $B$ 综合拓展一 $C$ 创新提升”三层动态题库，且 $A$ 、 $B$ 、 $C$ 三层题量之比为 $7 : 3 : 2$ ，设该题库中任意 1 道题被选到的可能性都相同.

(1)现有 4 人参加一项比赛，若每人分别独立地从该题库中随机选取一道题作答，求这 4 人中至少有 2 人的选题来自 $B$ 层的概率;

(2)现采用分层随机抽样的方法, 使用智能组卷系统从该题库中选取 12 道题生成试卷, 若某老师要从生成的这份 12 道题的试卷中随机选取 3 道题做进一步改编，记该老师选到 $A$ 层题的题数为 $X$ ,求 $X$ 的分布与期望 $E\left\lbrack  X\right\rbrack$.

22.(2024$\cdot$浦东新区二模$\cdot$19)某商店随机抽取了当天 100 名客户的消费金额，并分组如下: $\lbrack 0,{200})$ 、 $\lbrack {200},{400})$ 、 $\lbrack {400},{600})$ 、 $\cdots$ 、 $\left\lbrack  {{1000},{1200}}\right\rbrack$ (单位:元)，得到如图所示的频率分布直方图.

(1)若该店当天总共有 1350 名客户进店消费，试估计其中有多少客户的消费额不少于 800 元？

(2)若利用分层随机抽样的方法从消费不少于 800 元的客户中共抽取 6 人，再从这 6 人中随机抽取 2 人做进一步调查, 则抽到的 2 人中至少有 1 人的消费金额不少于 1000 元的概率是多少?

(3)为吸引顾客消费，该商店考虑两种促销方案:

方案一:消费金额每满 300 元可立减 50 元，并可叠加使用；

方案二:消费金额每满 1000 元即可抽奖三次，每次中奖的概率均为 $\dfrac{1}{3}$ ，且每次抽奖互不影响, 中奖 1 次当天消费金额可打 9 折, 中奖 2 次当天消费金额可打 6 折, 中奖 3 次当天消费金额可打 3 折.

若两种方案只能选择其中一种, 小王准备购买的商品又恰好标价 1000 元,请帮助他选择合适的促销方案并说明理由.

\begin{center}
\includegraphics[width=0.4\textwidth]{bodyxb/src/266_812_1487_583_345_0.jpg}
\end{center}


第 22 题图

23.(2025$\cdot$杨浦区二模$\cdot$19)为弘扬中华民族传统文化、增强民族自豪感，某学校开展中华古诗词背诵比赛，分为初赛和复赛. 全校同学都参加了初赛，并随机抽取一个班级进行初赛成绩统计，已知该班级共有 40 名学生，他们的初赛分数的频率分布直方图如图所示:

\begin{center}
\includegraphics[width=0.4\textwidth]{bodyxb/src/267_934_344_506_396_0.jpg}
\end{center}


第 23 题图

(1)计算 $b$ 的值，并估计该校这次初赛的平均分数；

(2)初赛分数达到 80 及以上的同学，称为优秀参赛选手，现从班级中随机选出 2 名同学，用 $X$ 代表其中的优秀参赛选手人数，求 $X$ 的分布；

(3)为增加比赛的趣味性，复赛规则如下:复赛试题将从题库中随机抽取，每位参赛选手将有机会回答填空、选择和简答各 1 题; 每答对 1 题得 1 分, 答错或不答得 0 分, 每位选手可以自行选择回答问题的顺序,若答对一题可继续答下一题,直到 3 题全部答完; 若答错或不答则比赛结束. 例如: 选手甲可自行按“简答一填空一选择”顺序答题，甲答对第一题得 1 分，并继续回答第二题且答错得 0 分，结束比赛，总分为 1 分.

小杨作为优秀参赛选手，代表班级参加复赛. 根据他初赛的答题正确频率，可估计他填空、选择和简答的答题正确概率如下表所示:

\begin{center}

\begin{tabular}{|c|c|c|c|}
\hline
题型 & 填空 & 选择 & 简答 \\
\cline{1-4}
答题正确概率 & 80\% & 90\% & 80\% \\
\cline{1-4}
\hline
\end{tabular}
\end{center}

若小杨每次答题的结果都相互独立，那么为尽量在比赛中获得较高分数，小杨应该采用怎样的答题顺序? 请说明理由.

\section*{第 14 章 数学建模问题}

\begin{center}

\begin{tabular}{|p{4cm}|p{12cm}|}
\hline
考点 & 课标要求与命题趋势 \\
\cline{1-2}
1. 圆锥曲线模型 &  新课标强调圆锥曲线的实际应用与跨学科融合, 要求学生掌握椭圆、双曲线、抛物线的定义、标准方程及几何性质, 并能将其应用于天体运动、工程设计等实际问题. 重点培养数学建模素养，如通过参数方程法解决最优化问题，或结合向量、导数等知识进行综合分析高考命题趋势:   \ding{172}  题型分布:圆锥曲线稳居压轴题(第 20 题)，占分约 18 分，常与直线、圆、向量综合考查 \ding{173}  核心考点:椭圆为主(占比超 60\%)，涉及轨迹方程、定值定点、存在性问题，近年弱化复杂计算，强化几何直观与参数方程应用 \ding{174}  创新方向:结合实际情境建模，如利用抛物线定义解决光线反射问题，或通过双曲线性质分析物理运动轨迹  复习策略:   \ding{172}  基础夯实:熟练掌握教材定义、方程及几何性质，重点突破椭圆第二定义、双曲线渐近线等拓展内容 \ding{173}  技巧提升:针对性训练“直曲联立转韦达”“参数方程消参”等方法，熟练运用硬解定理、 点差法简化运算 \ding{174}  真题精练:精研三年高考两年模拟真题，总结压轴题结构(如“第一问送分，第二问联立，第三问探索”)，强化存在性问题的逻辑推导 \\
\cline{1-2}
2. 函数模型 & 新课标强调函数模型的实际应用与跨学科整合，要求学生掌握指数函数、对数函数、幂函数的增长特征, 能将其应用于经济、物理等领域的最优化问题. 重点培养数学建模素养, 如通过导数分析函数极值，或结合分段函数解决复杂情境问题(如体形系数计算). 教材中专门设置数学建模模块，强化从实际问题抽象函数关系的能力  高考命题趋势:   \ding{172}  题型分布:函数应用题固定出现在解答题第 19 题(约 14 分)，填空题第 11 题(约 5 分)，常结合导数求最值 \ding{173}  核心考点:以指数函数(增长率)、对数函数(衰减模型)、分段函数(阶梯费用)为主，近年新增污染指数、体育竞技等真实情境建模 \ding{174}  创新方向:新定义问题(如“延展函数”)与导数结合，考查函数性质的灵活应用  复习策略:   \ding{172}  教材为本:梳理教材中函数模型案例(如利润最大化、人口增长)，掌握“审题一建模一求解一反馈”四步法 \ding{173}  专题突破:针对导数应用题(如容器体积优化、污染指数分析)，强化“求导一临界点分析一验证”流程 \ding{174}  真题精练:精研三年高考两年模拟真题，总结题型规律(如 2023 年体形系数问题的变量替换技巧)，重点训练分段函数的图像与方程转化 \\
\cline{1-2}
3. 不等式模型 & 新课标强调不等式模型的实际应用与跨学科整合，要求学生掌握基本不等式、一元二次不等式、分式不等式的解法及几何意义，能将其应用于资源分配、成本优化等现实问题. 重点培养数学建模素养，如通过导数分析不等式恒成立问题，或结合新定义(如 “延展函数”)探究不等式性质. 教材新增实际情境案例(如环保监测数据处理), 强化从实际问题抽象不等式关系的能力  高考命题趋势:   \ding{172}  题型分布:不等式固定出现在填空题第 1-2 题(约 5 分)、解答题第 19 题(约 14 分)， 常与导数、函数综合考查 \ding{173}  核心考点:基本不等式求最值(如 2025 年真题“$a + 1/b = 1$ ，求 $b + 1/a$ 最小值”)、一元二次不等式恒成立问题、分式不等式与函数定义域结合. 近年新增导数与不等式结合的存在性问题(如“存在极大值时 $m$ 的取值范围”) \ding{174}  创新方向:以实际数据为背景(如体育竞技得分、碳排放约束)，构建非线性不等式模型，考查参数范围求解与几何直观分析  复习策略:   \ding{172}  基础强化:熟练掌握教材中不等式案例(如利润最大化、仓储费用优化)，重点突破“1” 的代换、配凑法等技巧 \ding{173}  专题突破: 针对导数与不等式结合题(如“ $f\left( x\right)  \leq  {x}^{2} - 1$ 解集”),强化“构造函数 $\rightarrow$ 求导一分析极值”流程 \ding{174}  真题精练:精研三年高考两年模拟真题，总结“不等式十导数”压轴题规律(如 2025 年第 19 题分步设问),重点训练含参不等式的分类讨论 \\
\cline{1-2}
\hline
\end{tabular}
\end{center}
\clearpage
\subsection*{考点1 >> 圆锥曲线模型}

\subsubsection*{一、填空题}

1.(2025$\cdot$杨浦区二模$\cdot$11)如图所示，阿基米德椭圆规是由基座、带孔的横杆、两条互相垂直的空槽、两个可动滑块 $A$ 、 $B$ 组成的一种绘图工具，横杆的一端 $C$ 上装有铅笔，假设两条互相垂直的空槽和带孔的横杆都足够长， 将滑块 $A\text{ 、 }B$ 固定在带孔的横杆上,令滑块 $A$ 在其中一条空槽上滑动,滑块 $B$ 在另一条空槽上滑动,铅笔 $C$ 随之运动就能画出椭圆. 当 $A\text{ 、 }B$ 之间的距离为 ${14}{{cm}}$ 时,若需要画出一个离心率为 $\dfrac{4}{5}$ 的椭圆,则 $B\text{ 、 }C$ 之间的距离为\blkda[2]{}cm.

\begin{center}
\includegraphics[width=0.2\textwidth]{bodyxb/src/269_1111_1029_322_257_0.jpg}
\end{center}


第 1 题图

\subsubsection*{二、解答题}

2.(2025$\cdot$普陀区二模$\cdot$21)已知 $a$ 、 $b \in  \mathbf{R}$ ，对于函数 $y = f\left( x\right)$ ， $x \in  D$ ，设集合 $A = \{ \left( {x,y}\right)  \mid  y =  f\left( x\right) \} ,B = \left\{  {\left( {x,y}\right)  \mid  {\left( x - a\right) }^{2} + {\left( y - b\right) }^{2} \leq  1}\right\}$ ,记 ${M}_{f}\left( {a,b}\right)  = A \cap  B$.

(1)若 $f\left( x\right)  = \dfrac{1}{x} - \dfrac{3x}{4}$ ，请判断 ${M}_{f}\left( {0,0}\right)$ 中元素的个数，并说明理由；

(2)设 $k > 0,f\left( x\right)  = \sqrt{x + k}$ ，若 ${M}_{f}\left( {\dfrac{1}{2},0}\right)  = \left\{  \left( {{x}_{0},{y}_{0}}\right) \right\}$ ，求 $k$ 的值以及曲线 $y = f\left( x\right)$ 在点 $P\left( {{x}_{0},{y}_{0}}\right)$ 处的切线方程;

(3)设 $m \in  \mathbf{R},f\left( x\right)  = {{e}}^{x} - \dfrac{\ln x}{x} - \dfrac{1}{x} + m$ ，若对于任意的 $a$ ，皆有 ${M}_{f}\left( {a,0}\right)  = \varnothing$ 成立，求 $m$ 的取值范围.

3.(2025$\cdot$静安区一模$\cdot$20)如图所示的封闭图形的边缘由抛物线 $\Gamma$ 和垂直于抛物线对称轴的线段 ${AB}$ 组成. 已知 ${AB} = 4$ ，抛物线的顶点到线段 ${AB}$ 所在直线的距离为 2.

(1)请用数学符号语言表达这个封闭图形的边缘；

(2)在该封闭图形上截取一个矩形 ${CDEF}$ ，其中点 $C$ 、 $D$ 在线段 ${AB}$ 上，点 $E$ 、 $F$ 在抛物线 $\Gamma$ 上. 求以矩形 ${CDEF}$ 为侧面， ${CF}$ 为母线的圆柱的体积最大值；

(3)求证:抛物线 $\Gamma$ 的任何两条相互垂直的切线的交点都在同一条直线上.

\begin{center}
\includegraphics[width=0.2\textwidth]{bodyxb/src/270_1052_453_349_207_0.jpg}
\end{center}


第 3 题图

4.(2024$\cdot$静安区二模$\cdot$20)江南某公园内正在建造一座跨水拱桥，如平面图所示，现已经在地平面以上造好了一个外沿直径为 ${20}{m}$ 的半圆形拱桥洞,地平面与拱桥洞外沿交于点 $A$ 与点 $B$. 现在准备以地平面上的点 $C$ 与点 $D$ 为起点建造上、下桥坡道,要求:  \ding{172}  $\left| {BD}\right|  = \left| {AC}\right|$ ;  \ding{173}  在拱桥洞左侧建造平面图为直线的坡道,坡度为 $1 : 2\sqrt{2}$ (坡度为坡面的垂直高度和水平方向的距离的比);  \ding{174}  在拱桥洞右侧建造平面图为圆弧的坡道;  \ding{175}  在过桥的路面上骑车不颠簸.

\begin{center}
\includegraphics[width=0.5\textwidth]{bodyxb/src/270_446_1162_678_141_0.jpg}
\end{center}


第 4 题图

(1)请你设计一条过桥道路，画出大致的平面图，并用数学符号语言刻画与表达出来；

(2)按你的方案计算过桥道路的总长度；(精确到 ${0.1}{m}$ )

(3)若整个过桥坡道的路面宽为 ${10}{m}$ ，且铺设坡道全部使用混凝土. 请设计出所铺设路面的相关几何体，提出一个实际问题，写出解决该问题的方案，并说明理由. (如果需要，可通过假设的运算结果列式说明，不必计算)

\subsection*{考点2 > 函数模型}

\subsubsection*{一、填空题}

1. (2025$\cdot$徐汇区一模$\cdot$12)已知定义域为 $A = \{ 1,2,3\}$ 的函数 $y = f\left( x\right)$ 的值域也是 $A$ ，所有这样的函数 $y = f\left( x\right)$ 形成全集 $B$. 设非空集合 $C \subseteq  B$ 且 $\overline{C}$ 中的每一个函数都是 $C$ 中的两个函数(可以相同)的复合函数，则集合 $C$ 的元素个数的最小值为\blkda[2]{}.

2.(2024$\cdot$长宁区二模$\cdot$11)甲、乙、丙三辆出租车 2023 年运营的相关数据如下表所示:

\begin{center}

\begin{tabular}{|c|c|c|c|}
\hline
\phantom{X} & 甲 & 乙 & 丙 \\
\cline{1-4}
接单量 $t$ (单) & 7831 & 8225 & 8338 \\
\cline{1-4}
油费 $s$ (元) & 107 150 & 110264 & 110376 \\
\cline{1-4}
平均每单里程 $k\left( {{km}}\right)$ & 15 & 15 & 15 \\
\cline{1-4}
平均每公里油费 $a$ (元) & 0.7 & 0.7 & 0.7 \\
\cline{1-4}
\hline
\end{tabular}
\end{center}

出租车空驶率=地租车没有载客行驶的里程;依据上述数据，小明建立了求解三辆车的空驶率的模型 $u = f\left( {s,t,k,a}\right)$ ,并求得甲、乙、丙的空驶率分别为 ${23.26}\%$ 、 ${21.68}\%$ 、 $x\%$ ,则 $x =$ \blkda[2]{}. (精确到 0.01)

\subsubsection*{二、选择题}

3.(2025$\cdot$杨浦区一模$\cdot$15)小李研究数学建模“雨中行”问题，在作出“降雨强度保持不变”“行走速度保持不变”“将人体视作一个长方体”等合理假设的前提下，他设了如下表所示的变量:

\begin{center}

\begin{tabular}{|c|c|c|c|c|c|c|c|}
\hline
人的身高 & 人体宽度 & 人体厚度 & 降雨速度 & 雨滴密度 & 行走距离 & 风速 & 行走速度 \\
\cline{1-8}
$h$ & $w$ & $d$ & ${v}_{r}$ & $p$ & $D$ & ${v}_{w}$ & $v$ \\
\cline{1-8}
\hline
\end{tabular}
\end{center}

并构建模型如下:

当人迎风行走时，人体总的淋雨量为 $T = \dfrac{pwD}{v}\left\lbrack  {d{v}_{r} + h\left( {{v}_{w} + v}\right) }\right\rbrack$.

根据模型,小李对“雨中行”作出如下解释:

 \ding{172}  若两人结伴迎风行走，则体型较高大魁梧的人淋雨量较大；

 \ding{173}  若某人迎风行走，则走得越快淋雨量越小；若背风行走，则走得越慢淋雨量越小；

 \ding{174}  若某人迎风行走了 ${10}{s}$ ，则行走距离越长淋雨量越大.

这些解释合理的个数为 \dotfill{(\kh{})}

\xx{}{}{}{}0 B. 1

C. 2 D. 3

\subsubsection*{三、解答题}

4.(2025$\cdot$奉贤区二模$\cdot$19)将一块边长为 10 cm 的正方形铁片制作一个正四棱锥的容器罩. 同学们设计了甲、乙、丙三个不同的方案，各自裁下阴影部分，用余下的制作成正四棱锥容器罩，形如最右边的图. 甲和丙是去制作有盖的容器罩, 乙是去制作无盖的容器罩. 假设加工过程中铁片损失忽略不计. 设甲、乙、丙中白色的四个等腰三角形的底边分别是 $x\text{ 、 }m\text{ 、 }n$.

\begin{center}
\includegraphics[width=0.9\textwidth]{bodyxb/src/272_189_491_1116_265_0.jpg}
\end{center}


第 4 题图

(1)请你选择其中的某一个方案，而且只需选一个方案(选择超过一个方案的，按第一个方案处理). 你选择的方案是\blkda[2]{}，求解以下 2 个问题:

 \ding{172}  求出所选方案相对应的棱锥的侧面积 $S\left( x\right)$ 、 $S\left( m\right)$ 、 $S\left( n\right)$ ；

 \ding{173}  求出所选方案相对应棱锥的体积 $V\left( x\right)$ 、 $V\left( m\right)$ 、 $V\left( n\right)$ 的最大值.

(2)假设三个方案中相应的体积最大值分别记作 $V{\left( x\right) }_{\max }$ 、 $V{\left( m\right) }_{\max }$ 、 $V{\left( n\right) }_{\max }$ ，请直接写出三者的大小关系. (不写判断理由与过程)

5.(2025$\cdot$上海春季卷$\cdot$21)已知函数 $y = f\left( x\right)$ 的定义域是 $D$. 对于 $t \in  D$ ，定义集合 ${S}_{f\left( t\right) } = \{ x \mid  f\left( x\right)  \geq  f\left( t\right) \}$.

(1)若 $f\left( x\right)  = {\log }_{2}x$ ，求 ${S}_{f\left( {16}\right) }$ ；

(2)对于集合 $A$ ，若对任意 $x \in  A$ 都有 $- x \in  A$ ，则称 $A$ 是对称集. 若 $D$ 是对称集，证明: “函数 $y = f\left( x\right)$ 是偶函数”的充要条件是“对任意 $t \in  D$ ,都有 ${S}_{f\left( t\right) } = {S}_{f\left( {-t}\right) }$ ”;

(3)若 $f\left( x\right)  = {{e}}^{x} - \dfrac{1}{2}m{x}^{2},x \in  \mathbf{R}$ ，对任意 ${t}_{1} < {t}_{2} \in  D$ ，都有 ${S}_{f\left( {t}_{2}\right) } \subseteq  {S}_{f\left( {t}_{1}\right) }$ ，求 $m$ 的取值范围.

考点3 一、不等式模型

\subsubsection*{一、填空题}

1.(2024$\cdot$青浦区一模$\cdot$12)已知三个互不相同的实数 $a$ 、 $b$ 、 $c$ 满足 $a + b + c = 1$ ， ${a}^{2} + {b}^{2} + {c}^{2} = 3$ ，则 ${abc}$ 的取值范围为\blkda[2]{}.

\subsubsection*{二、解答题}

2.(2024$\cdot$青浦区一模$\cdot$19)上海各中学都定期进行紧急疏散演习:当警报响起，建筑物内师生马上有组织、尽快地疏散撤离. 对于一个特定的建筑物, 管理人员关心房间内所有人疏散完毕 (房间最后一个人到达安全出口处)所用时间. 数学建模小组准备对某教学楼第一层楼两间相同的教室展开研究. 为此, 他们提出如下模型假设:

 \ding{172}  疏散时所有人员有秩序地撤离建筑物;  \ding{173}  所有人员排成单列行进撤离;  \ding{174}  队列中人员的间隔是均匀的;  \ding{175}  队列匀速地撤离建筑物.

(1)上述模型假设是否合理，请任选两个模型假设说明理由；

(2)如图所示，设第一间教室的人数为 ${n}_{1} + 1$ ，第二间教室的人数为 ${n}_{2} + 1$ ，每间教室的长度为 $l$ ， 其中 ${n}_{1}\text{ 、 }{n}_{2}$ 都是正整数， $l > 0$ ，忽略教室门的宽度及忽略教室内人群到教室门口的时间. 请再引入适当的变量，建立两个教室内的人员完全撤离所用时间的数学模型.

\begin{center}
\includegraphics[width=0.3\textwidth]{bodyxb/src/273_960_1641_456_287_0.jpg}
\end{center}


第 2 题图

3.(2025$\cdot$静安区二模$\cdot$21)若存在实常数 $k$ 、 $m$ ，对任意 $x \in  D$ ，不等式 $f\left( x\right)  \geq  {kx} + m \geq  g\left( x\right)$ 恒成立， 则称直线 $y = {kx} + m$ 是函数 $y = f\left( x\right)$ 和函数 $y = g\left( x\right)$ 在 $D$ 上的分界线.

(1)请写出函数 $y = \dfrac{1}{2}{x}^{2}$ 和函数 $y = \ln x$ 在 $\left( {0, + \infty }\right)$ 上的一条斜率为 1 的分界线; (不必证明)

(2)求证:函数 $y = x + \dfrac{1}{x}$ 和函数 $y = 2 - \dfrac{1}{x}$ 在 $\left( {0, + \infty }\right)$ 上过坐标原点的分界线有且只有一条；

(3)试探究函数 $y = {{e}}^{x}\left( {x + 1}\right)$ ( ${e}$ 为自然对数的底数) 和函数 $y =  - {x}^{2} + {2x} + 1$ 在 $\mathbf{R}$ 上是否存在分界线? 若存在, 求出分界线方程; 若不存在, 请说明理由.

\section*{第 15 章 新颖题型}

\begin{center}

\begin{tabular}{|p{4cm}|p{12cm}|}
\hline
考点 & 课标要求与命题趋势 \\
\cline{1-2}
1. 数列新定义 & 课标要求:   \ding{172}  核心素养导向. 根据《普通高中数学课程标准(2017 年版 2020 年修订)》(以下简称《新课标》)，数列教学需强化数学抽象与逻辑推理能力. 数列新定义题型要求学生通过阅读理解新情境，将实际问题转化为数学模型 \ding{173}  知识迁移能力. 课程标准强调数列作为特殊函数的本质，新定义题型常需学生将函数、 不等式等知识迁移至数列情境. 如:函数与数列的关联；跨学科融合，如“斐波那契数列”与生物繁殖模型结合, 体现数学在实际问题中的应用 \ding{174}  综合应用能力. 数列新定义题型常与其他模块交叉考查, 例如: 将等比数列证明与双曲线性质结合; 概率统计, 部分新定义题型涉及数列与概率的综合, 进行概率分布分析  命题趋势:   \ding{172}  题型创新与多样化. 高阶数列与特殊结构, 近年高频出现“差分数列”“对称数列”“非典型新定义数列”等; 数学文化渗透, 如 “杨辉三角数列” “牛顿数列” 等历史背景题型, 考查学生对数学史的理解与应用 \ding{173}  难度分层与梯度设计. 基础题:侧重定义理解与简单应用；压轴题:强调逻辑推理与创新性 \ding{174}  跨学科与实际情境结合. 经济与金融模型;物理与工程背景, 考查学生将数学工具迁移至实际问题的能力 \ding{175}  开放探究与存在性证明  备考策略:   \ding{172}  强化定义理解与信息提取. 精准翻译，将新定义转化为数学语言；类比迁移:通过熟悉数列(如等差、等比数列)的研究方法，探索新数列性质 \ding{173}  掌握核心解题方法. 特殊化验证, 通过具体数值代入，观察并发现、表述规律; 数学归纳法；构造辅助数列 \ding{174}  提升综合应用能力. 知识串联，将数列与导数、概率等模块结合; 实际问题建模，训练数学建模能力 \\
\cline{1-2}
2. 函数新定义 & 课标要求:   \ding{172}  核心素养导向. 根据《新课标》，函数教学需强化数学抽象与逻辑推理能力. 新定义题型要求学生通过阅读理解新情境，将实际问题转化为数学模型 \ding{173}  知识迁移能力. 课程标准强调函数作为变量关系刻画工具的本质，新定义题型常需学生将导数、不等式等知识迁移至函数情境. 例如导数与函数的关联; 跨学科融合, 如 “斐波那契函数”与生物繁殖模型结合, 体现数学在实际问题中的应用 \ding{174}  综合应用能力. 函数新定义题型常与其他模块交叉考查  命题趋势:  题型创新与多样化;难度分层与梯度设计;跨学科与实际情境结合等  备考策略:  强化定义理解与信息提取, 掌握核心解题方法, 提升综合应用能力 \\
\cline{1-2}
\hline
3. 集合新定义 & 课标要求:  \ding{172}  核心素养导向. 根据《新课标》，集合教学需强化数学抽象与逻辑推理能力. 新定义题型要求学生通过阅读理解新情境，将实际问题转化为数学模型 \ding{173} 知识迁移能力，课程标准强调集合作为数学基础工具的本质，新定义题型常需学生将函数、数列等知识迁移至集合情境. 如:函数与集合的关联；跨学科融合等 \ding{174}  综合应用能力. 集合新定义题型常与其他模块交叉考查  命题趋势:  题型创新与多样化; 难度分层与梯度设计; 跨学科与实际情境结合等 备考策略: 强化定义理解与信息提取,掌握核心解题方法,提升综合应用能力 \\
\cline{1-2}
4. 其他新定义 & \ding{172}  新定义题型的特点:通过给出一个新概念，或约定一种新运算，或给出几个新模型来创设全新的问题情境，要求考生在阅读理解的基础上，依据题目提供的信息，联系所学的知识和方法，实现信息的迁移，达到灵活解题的目的；遇到新定义问题，应耐心读题，分析新定义的特点，弄清新定义的性质，按新定义照章办事，逐条分析、验证、运算，使问题得以解决, 难度较高, 需重点特训 。上海高考数学新结构体系下，新定义类试题更综合性地考查学生的思维能力和推理能力，以问题为抓手，创新设问方式，搭建思维平台，引导考生思考，在思维过程中领悟，数学方法. 题目更加注重综合性、应用性、创新性, 该类题分值最高, 试题容量明显增大, 对学科核心素养的考查也更深入 氏轴颞俞颞钉破了试题颞理、命题方式、试卷结构的固有模式，增强试题的灵活性，米取多样的形式、多角度的提问，考查学生的数学能力 \\
\cline{1-2}
\hline
\end{tabular}
\end{center}
\clearpage
\subsection*{考点1}

\subsubsection*{一、填空题}

1.(2025$\cdot$杨浦区二模$\cdot$12)由若干个多边形所覆盖的区域，称为这些多边形的并集,如图所示,梯形 ${ACDE}$ 是 $\bigtriangleup  {ACE}$ 与矩形 ${BCDE}$ 的并集. 已知 $n$ 是正整数，在平面直角坐标系 ${xOy}$ 中，直线 ${l}_{n}$ 的方程为 $y =  - {2}^{n}x + n$ ,若直线 ${l}_{n}$ 交 $x$ 轴于点 ${A}_{n}$ ,交 $y$ 轴于点 ${B}_{n}$ ,则 $\bigtriangleup {A}_{1}O{B}_{1}$ 、 $\bigtriangleup  {A}_{2}O{B}_{2}\text{ 、 }\cdots \text{ 、 } \bigtriangleup  {A}_{10}O{B}_{10}$ 的并集，其面积为\blkda[2]{}.

\begin{center}
\includegraphics[width=0.2\textwidth]{bodyxb/src/276_1088_1404_325_227_0.jpg}
\end{center}


第 1 题图

\subsubsection*{二、选择题}

2.(2025$\cdot$普陀区二模$\cdot$16)设 $k \geq  2,n \geq  1,k\text{ 、 }n \in  \mathbf{N},{S}_{n}$ 是数列 $\left\{  {a}_{n}\right\}$ 的前 $n$ 项和，且满足 $2{S}_{n} = 3{a}_{n} - 1$ ， 数列 $\left\{  {b}_{n}\right\}$ 是由 $k$ 个大于 -2 的整数组成的有穷数列,若 ${b}_{1} \neq  0,\mathop{\sum }\limits_{{i = 1}}^{k}{b}_{i}{a}_{k - i + 1} = 0$ ,则称数列 $\left\{  {b}_{n}\right\}$ 是数列 $\left\{  {a}_{n}\right\}$ 的“ $T$ 数列”. 对于数列 $\left\{  {b}_{n}\right\}$ 有如下两个命题:  \ding{172} 若 ${b}_{1} > 0$ ，则数列 $\left\{  {b}_{n}\right\}$ 不是数列 $\left\{  {a}_{n}\right\}$ 的“ $T$ 数列”;  \ding{173} 若 $k \geq  3$ ，则数列 $\left\{  {a}_{n}\right\}$ 的“$T$ 数列”至少有 5 个. 则下列结论中正确的是 \dotfill{(\kh{})}

\xx{}{}{}{} \ding{172}  \ding{173} 均为真命题 B.  \ding{172} 为真命题， \ding{173} 为假命题

C.  \ding{172} 为假命题， \ding{173} 为真命题 D.  \ding{172}  \ding{173} 均为假命题



3.(2025$\cdot$青浦区二模$\cdot$16)数学上用符号 $\mathop{\prod }\limits_{{i = 1}}^{1}{a}_{i}$ 表示 $n$ 个实数 ${a}_{1}$ 、 ${a}_{2}$ 、 $\cdots$ 、 ${a}_{n}$ 的积. 设 ${x}_{1}$ 、 ${x}_{2}$ 、、、 ${x}_{200}$ 、 ${y}_{1}$ 、 ${y}_{2}\text{ 、 }\cdots \text{ 、 }{y}_{200}$ 为互不相同的实数,已知 $\mathop{\prod }\limits_{{j = 1}}^{{200}}\left( {{x}_{i} + {y}_{j}}\right)  = {2025}\left( {i = 1,2,\cdots ,{200}}\right)$ ,则 $\mathop{\prod }\limits_{{i = 1}}^{{200}}\left( {{x}_{i} + {y}_{j}}\right)$ 等于 \dotfill{(\kh{})}

\xx{}{}{}{}-2025 B. 2025 C. -2026 D. 2026

\subsubsection*{三、解答题}

4.(2024$\cdot$普陀区一模$\cdot$21)若存在常数 $t$ ，使得数列 $\left\{  {a}_{n}\right\}$ 满足 ${a}_{n + 1} - {a}_{1}{a}_{2}{a}_{3}\cdots {a}_{n} = t\left( {n \geq  1,n \in  \mathbf{N}}\right)$ ，则称数列 $\left\{  {a}_{n}\right\}$ 为“ $H\left( t\right)$ 数列”.

(1)判断数列:1,2,3,8,49 是否为“ $H\left( 1\right)$ 数列”，并说明理由;

(2)若数列 $\left\{  {a}_{n}\right\}$ 是首项为 2 的“$H\left( t\right)$ 数列”，数列 $\left\{  {b}_{n}\right\}$ 是等比数列，且 $\left\{  {a}_{n}\right\}$ 与 $\left\{  {b}_{n}\right\}$ 满足 $\mathop{\sum }\limits_{{i = 1}}^{n}{a}_{i}^{2} =  {a}_{1}{a}_{2}{a}_{3}\cdots {a}_{n} + {\log }_{2}{b}_{n}$ ,求 $t$ 的值和数列 $\left\{  {b}_{n}\right\}$ 的通项公式;

(3)若数列 $\left\{  {a}_{n}\right\}$ 是“ $H\left( t\right)$ 数列”， ${S}_{n}$ 为数列 $\left\{  {a}_{n}\right\}$ 的前 $n$ 项和， ${a}_{1} > 1$ ， $t > 0$ ，试比较 $\ln {a}_{n}$ 与 ${a}_{n} - 1$ 的大小，并证明 $t > {S}_{n + 1} - {S}_{n} - {{e}}^{{S}_{n} - n}$.

考点2 > 函数新定义

\subsubsection*{一、选择题}

1.(2024$\cdot$浦东新区一模$\cdot$16)对于函数 $y = f\left( x\right)$ 和 $y = g\left( x\right)$ ，及区间 $D$ ，若存在实数 $k$ 、 $b$ ，使得 $f\left( x\right)  \geq  {kx} + b \geq  g\left( x\right)$ 对任意 $x \in  D$ 恒成立,则称 $y = f\left( x\right)$ 在区间 $D$ 上 “优于” $y = g\left( x\right)$. 有以下两个结论:

 \ding{172}  $f\left( x\right)  = {\log }_{2}x$ 在区间 $D = \left\lbrack  {1,2}\right\rbrack$ 上优于 $g\left( x\right)  = {x}^{2} - {2x} + 1$ ;

 \ding{173}  当 $m \leq   - 2$ 时， $f\left( x\right)  = {x}^{3}$ 在区间 $D = \left\lbrack  {-1,1}\right\rbrack$ 上优于 $g\left( x\right)  = {{e}}^{x} + m$. 则下列判断正确的是 \dotfill{(\kh{})}

\xx{}{}{}{} \ding{172}  \ding{173} 均正确 B.  \ding{172} 正确， \ding{173} 错误

C.  \ding{172} 错误， \ding{173} 正确 D.  \ding{172}  \ding{173} 均错误

\subsubsection*{二、解答题}

2.(2025$\cdot$闵行区二模$\cdot$21)已知函数 $y = f\left( x\right)$ 在定义域 $D$ 上存在导函数 $y = {f}^{\prime }\left( x\right)$. 对于给定的一个有序实数对 $\left( {k,m}\right)$ ,若存在 ${x}_{1}\text{ 、 }{x}_{2} \in  D$ ,使得 $\left\lbrack  {k{x}_{1} - f\left( {x}_{1}\right)  + m}\right\rbrack   \cdot  \left\lbrack  {k{x}_{2} - f\left( {x}_{2}\right)  + m}\right\rbrack   < 0$ ,则称 $\left( {k,m}\right)$ 为 $y = f\left( x\right)$ 在定义域 $D$ 上的一个“分割数对”.

(1)已知 $f\left( x\right)  = {x}^{2},D = \mathbf{R}$ ，判断数对 $\left( {1,0}\right)$ 是否为 $y = f\left( x\right)$ 在 $D$ 上的“分割数对”，并说明理由；

(2)已知 $f\left( x\right)  = \ln x,D = \left( {1,2}\right)$ ，若 $\left( {\ln 2,m}\right)$ 为 $y = f\left( x\right)$ 在区间 $D$ 上的“分割数对”，求实数 $m$ 的取值范围;

(3)已知 $f\left( x\right)  = \left( {{x}^{2} + {ax} + b}\right)  \cdot  {{e}}^{x},D = \mathbf{R}$ ，若有且仅有一个实数 $a$ 满足对任意 $t \in  \mathbf{R},\left( {{f}^{\prime }\left( t\right) }\right.$ , $f\left( t\right)  - t{f}^{\prime }\left( t\right) )$ 都不是 $y = f\left( x\right)$ 在 $D$ 上的“分割数对”,求实数 $b$ 的值.

3.(2025$\cdot$青浦区二模$\cdot$21)函数的导函数有很多有趣的性质，例如:函数 $y = c$ (实数 $c$ 为常数)的导函数为 $y = 0$ ; 反之,若函数 $y = \varphi \left( x\right)$ 的导函数为 ${\varphi }^{\prime }\left( x\right)  = 0$ ,则 $\varphi \left( x\right)  = c$ (实数 $c$ 为常数). 已知函数 $y = f\left( x\right)$ 与 $y = g\left( x\right)$ 定义域都是 $\mathbf{R}$ ,导函数分别为 $y = {f}^{\prime }\left( x\right)$ 和 $y = {g}^{\prime }\left( x\right)$. 若 ${f}^{\prime }\left( x\right)  = f\left( x\right)$ , 则称 $y = f\left( x\right)$ 是 “自导函数”; 若 ${f}^{\prime }\left( x\right)  = g\left( x\right)$ 且 ${g}^{\prime }\left( x\right)  =  - f\left( x\right)$ ,则称 $y = f\left( x\right)$ 与 $y = g\left( x\right)$ 是 “共轭互导函数”.

(1)请判断函数 $y = {{e}}^{{ax} + b}\left( {a\text{ 、 }b \in  \mathbf{R},a \neq  0}\right)$ 是否是“自导函数”，并说明理由；

(2)若函数 $y = f\left( x\right)$ 是“自导函数”，且满足 $f\left( 0\right)  = 1$ ，求证: $f\left( x\right) f\left( {-x}\right)  = 1$ ；

(3)若函数 $y = f\left( x\right)$ 与 $y = g\left( x\right)$ 是 “共轭互导函数”，满足 $f\left( 0\right)  = 0$ 、 $g\left( 0\right)  = 1$ ，求证: ${f}^{2}\left( x\right)  +  {g}^{2}\left( x\right)  = 1$ ,进而证明 $f\left( x\right)  = \sin x$ 且 $g\left( x\right)  = \cos x$.

4.(2024$\cdot$徐汇区一模$\cdot$21)若函数 $y = f\left( x\right) ,x \in  \mathbf{R}$ 的导函数 $y = {f}^{\prime }\left( x\right) ,x \in  \mathbf{R}$ 是以 $T\left( {T \neq  0}\right)$ 为周期的函数,则称函数 $y = f\left( x\right) ,x \in  \mathbf{R}$ 具有“$T$ 性质”.

(1)试判断函数 $y = {x}^{2}$ 和 $y = \sin x$ 是否具有“${2\pi }$ 性质”，并说明理由；

(2)已知函数 $y = h\left( x\right)$ ，其中 $h\left( x\right)  = a{x}^{2} + {bx} + 2\sin {bx}\left( {0 < b < 3}\right)$ 具有“$\pi$ 性质”，求函数 $y = h\left( x\right)$ 在 $\left\lbrack  {0,\pi }\right\rbrack$ 上的极小值点;

(3)若函数 $y = f\left( x\right) ,x \in  \mathbf{R}$ 具有“ $T$ 性质”，且存在实数 $M > 0$ 使得对任意 $x \in  \mathbf{R}$ 都有 $\left| {f\left( x\right) }\right|  < M$ 成立,求证: $y = f\left( x\right) ,x \in  \mathbf{R}$ 为周期函数.

[可用结论: 若函数 $y = f\left( x\right) ,x \in  \mathbf{R}$ 的导函数满足 ${f}^{\prime }\left( x\right)  = 0,x \in  \mathbf{R}$ ,则 $f\left( x\right)  = C$ (常数).]

5.(2023$\cdot$上海春季卷$\cdot$21)设 $f\left( x\right)  = a{x}^{3} - \left( {a + 1}\right) {x}^{2} + x,g\left( x\right)  = {kx} + m$ ，其中 $a \geq  0,k$ 、 $m \in  \mathbf{R}$ ，若对任意 $x \in  \left\lbrack  {0,1}\right\rbrack$ 均有 $f\left( x\right)  \leq  g\left( x\right)$ ,则称函数 $y = g\left( x\right)$ 是函数 $y = f\left( x\right)$ 的“控制函数”,且对所有的函数 $y = g\left( x\right)$ 取最小值定义为 $\overline{f}\left( x\right)$.

(1)若 $a = 2,g\left( x\right)  = x$ ，试问 $y = g\left( x\right)$ 是否为函数 $y = f\left( x\right)$ 的“控制函数”；

(2)若 $a = 0$ ，使得直线 $y = h\left( x\right)$ 是曲线 $y = f\left( x\right)$ 在 $x = \dfrac{1}{4}$ 处的切线. 证明:函数 $y = h\left( x\right)$ 为函数 $y = f\left( x\right)$ 的“控制函数”，并求 $\overline{f}\left( \dfrac{1}{4}\right)$ 的值；

(3)若曲线 $y = f\left( x\right)$ 在 $x = {x}_{0},{x}_{0} \in  \left( {0,1}\right)$ 处的切线过点 $\left( {1,0}\right)$ ，且 $c \in  \left\lbrack  {{x}_{0},1}\right\rbrack$. 证明:当且仅当 $c =  {x}_{0}$ 或 $c = 1$ 时, $\overline{f}\left( c\right)  = f\left( c\right)$.

6.(2024$\cdot$金山区二模$\cdot$21)已知函数 $y = f\left( x\right)$ 与 $y = g\left( x\right)$ 有相同的定义域 $D$. 若存在常数 $a(a \in$ R),使得对于任意的 ${x}_{1} \in  D$ ,都存在 ${x}_{2} \in  D$ ,满足 $f\left( {x}_{1}\right)  + g\left( {x}_{2}\right)  = a$ ,则称函数 $y = g\left( x\right)$ 是函数 $y = f\left( x\right)$ 关于 $a$ 的“ $S$ 函数”.

(1)若 $f\left( x\right)  = \ln x,g\left( x\right)  = {{e}}^{x}$ ，试判断函数 $y = g\left( x\right)$ 是否是 $y = f\left( x\right)$ 关于 0 的“$S$ 函数”，并说明理由;

(2)若函数 $y = f\left( x\right)$ 与 $y = g\left( x\right)$ 均存在最大值与最小值，且函数 $y = g\left( x\right)$ 是 $y = f\left( x\right)$ 关于 $a$ 的“$S$ 函数”, $y = f\left( x\right)$ 又是 $y = g\left( x\right)$ 关于 $a$ 的“ $S$ 函数”,证明: ${\left\lbrack  f\left( x\right) \right\rbrack  }_{\min } + {\left\lbrack  g\left( x\right) \right\rbrack  }_{\max } = a$ ;

(3)已知 $f\left( x\right)  = \left| {x - 1}\right| ,g\left( x\right)  = \sqrt{x}$ ，其定义域均为 $\left\lbrack  {0,t}\right\rbrack$. 给定正实数 $t$ ，若存在唯一的 $a$ ，使得 $y = g\left( x\right)$ 是 $y = f\left( x\right)$ 关于 $a$ 的“ $S$ 函数”，求 $t$ 的所有可能值.

7.(2025$\cdot$浦东新区一模$\cdot$21)过曲线 $y = f\left( x\right)$ 上一点 $P$ 作其切线，若恰有两条，则称 $P$ 为 $f\left( x\right)$ 的“$A$ 类点”；过曲线 $y = f\left( x\right)$ 外一点 $Q$ 作其切线，若恰有三条，则称 $Q$ 为 $f\left( x\right)$ 的“ $B$ 类点”；若点 $R$ 为 $f\left( x\right)$ 的“ $A$ 类点”或“ $B$ 类点”，且过 $R$ 存在两条相互垂直的切线，则称 $R$ 为 $f\left( x\right)$ 的“ $C$ 类点”.

(1)设 $f\left( x\right)  = \dfrac{1}{{x}^{2}}$ ，判断点 $P\left( {1,1}\right)$ 是否为 $f\left( x\right)$ 的“ $A$ 类点”，并说明理由；

(2)设 $f\left( x\right)  = {x}^{3} - {mx}$ ，若点 $Q\left( {2,0}\right)$ 为 $f\left( x\right)$ 的“ $B$ 类点”，且过点 $Q$ 的三条切线的切点横坐标可构成等差数列，求实数 $m$ 的值；

(3)设 $f\left( x\right)  = \dfrac{x + 1}{{{e}}^{x}}$ ，证明: $y$ 轴上不存在 $f\left( x\right)$ 的“ $C$ 类点”.

8.(2024$\cdot$金山区一模$\cdot$21)设函数 $y = f\left( x\right)$ 的定义域为 $D$ ，给定区间 $\left\lbrack  {a,b}\right\rbrack   \subseteq  D$ ，若存在 ${x}_{0} \in  \left( {a,b}\right)$ ， 使得 $f\left( {x}_{0}\right)  = \dfrac{f\left( b\right)  - f\left( a\right) }{b - a}$ ,则称函数 $y = f\left( x\right)$ 为区间 $\left\lbrack  {a,b}\right\rbrack$ 上的“均值函数”, ${x}_{0}$ 为函数 $y =  f\left( x\right)$ 的“均值点”.

(1)试判断函数 $y = {x}^{2}$ 是否为区间 $\left\lbrack  {1,2}\right\rbrack$ 上的“均值函数”，如果是，请求出其“均值点”；如果不是， 请说明理由;

(2)已知函数 $y =  - {2}^{{2x} - 1} + m \cdot  {2}^{x - 1} - {12}$ 是区间 $\left\lbrack  {1,3}\right\rbrack$ 上的“均值函数”,求实数 $m$ 的取值范围;

(3)若函数 $y = \dfrac{{x}^{2} + a}{2\left( {{x}^{2} - {2x} + 2}\right) }$ (常数 $a \in  \mathbf{R}$ )是区间 $\left\lbrack  {-2,2}\right\rbrack$ 上的“均值函数”，且 $\dfrac{2}{3}$ 为其“均值点”. 将区间 $\left\lbrack  {-2,0}\right\rbrack$ 任意划分成 $m + 1\left( {m \in  \mathbf{N}}\right)$ 份,设分点的横坐标从小到大依次为 ${t}_{1},{t}_{2},\cdots ,{t}_{m}$ , 记 ${t}_{0} =  - 2,{t}_{m + 1} = 0,G = \mathop{\sum }\limits_{{i = 0}}^{m}\left| {f\left( {t}_{i + 1}\right)  - f\left( {t}_{i}\right) }\right|$. 再将区间 $\left\lbrack  {0,2}\right\rbrack$ 等分成 ${2}^{n} + 1\left( {n \in  \mathbf{N}}\right)$ 份,设等分点的横坐标从小到大依次为 ${x}_{1},{x}_{2},\cdots ,{x}_{{2}^{n}}$ ,记 $H = \mathop{\sum }\limits_{{i = 1}}^{{2}^{n}}f\left( {x}_{i}\right)$. 求使得 $H \cdot  G > {2023}$ 的最小整数 $n$ 的值.

9.(2025$\cdot$长宁区一模$\cdot$21)双曲余弦函数 $\cosh x = \dfrac{{{e}}^{x} + {{e}}^{-x}}{2}$ ，双曲正弦函数 $\sinh x = \dfrac{{{e}}^{x} - {{e}}^{-x}}{2}$.

(1)求函数 $\cos {hx} = \dfrac{{{e}}^{x} + {{e}}^{-x}}{2}$ 的单调增区间；

(2)若函数 $y = \cosh {2x} - a\sinh x$ 在 $\left( {0, + \infty }\right)$ 上的最小值是 $\dfrac{1}{4}$ ，求实数 $a$ 的值；

(3)对任意 $x \in  \mathbf{R},\cos h\left( x\right)  \geq  \cos x + m{x}^{2}$ 恒成立，求实数 $m$ 的取值范围.

考点3 > 集合新定义

\subsubsection*{一、填空题}

1. (2025$\cdot$奉贤区一模$\cdot$12)已知集合 $M = \left\{  {{P}_{0},{P}_{1},{P}_{2},\cdots ,{P}_{n}}\right\}  \left( {n \geq  2,n \in  \mathbf{N}}\right)$ 是由函数 $y = \cos x,x \in  \left\lbrack  {0,{2\pi }}\right\rbrack$ 的图像上两两不相同的点构成的点集,集合 $S = \left\{  {a\left| {a = {\vv{OP}}_{0} \cdot  {\vv{OP}}_{i}}\right. ,i = 0,1,2,\cdots ,n,n \geq  2}\right.$ , $n \in  \mathbf{N}\}$ ,其中 ${P}_{0}\left( {0,1}\right) \text{ 、 }{P}_{1}\left( {\pi , - 1}\right)$. 若集合 $S$ 中的元素按照从小到大的顺序排列能构成公差为 $d$ 的等差数列,当 $d \in  \left\{  {\dfrac{1}{2},1}\right\}$ 时,则符合条件的点集 $M$ 的个数为\blkda[2]{}.

2.(2025$\cdot$崇明区二模$\cdot$12)已知集合 $M$ 中的任一个元素都是整数，当存在整数 $a$ 、 $c \in  M$ ， $b \notin  M$ 且 $a < b < c$ 时，称 $M$ 为 “间断整数集”. 集合 $\{ x \mid  1 \leq  x \leq  {10},x \in  \mathbf{Z}\}$ 的所有子集中，是 “间断整数集” 的个数为\blkda[2]{}.

3.(2025$\cdot$虹口区二模$\cdot$12)记|A|为有限集合 $A$ 中的元素个数. 设 $\omega  > 0,{S}_{\omega } = \left\{  {\theta  \mid  {2}^{2025} + \omega  \cdot  \theta }\right.$ 能被 7 整除 4,若对于任意实数 $a$ 和任意正整数 $n$ ,恒有 $\left| {{S}_{\omega } \cap  \left( {a,a + n{{e}}^{-{0.5n}}}\right) }\right|  \leq  3$ ,则实数 $\omega$ 的取值范围是\blkda[2]{}.

\subsubsection*{二、选择题}

4.(2025$\cdot$金山区一模$\cdot$16)已知三棱锥 ${A}_{1} - {A}_{2}{A}_{3}{A}_{4}$ 的侧棱长相等，且侧棱两两垂直. 设 $P$ 为该三棱锥表面 (含棱) 上异于顶点 ${A}_{1}$ 、 ${A}_{2}$ 、 ${A}_{3}$ 、 ${A}_{4}$ 的点,记 $D = \left\{  {d\left| {d = }\right| P{A}_{i}\mid ,i = 1,2,3,4}\right\}$. 若集合 $D$ 中有且只有 2 个元素,则符合条件的点 $P$ 有\blkda[2]{}个. \dotfill{(\kh{})}

\xx{}{}{}{}3 B. 6 C. 7 D. 10

5.(2024$\cdot$全山区一模$\cdot$16)设集合 $A = \{ 1,2,\cdots ,{100}\} ,X$ 、 $Y$ 均为 $A$ 的非空子集(允许 $X = Y$ ). $X$ 中的最大元素与 $Y$ 中的最小元素分别记为 $M$ 、 $m$ ，则满足 $M > m$ 的有序集合对 $\left( {X,Y}\right)$ 的个数为

\dotfill{(\kh{})}

\xx{}{}{}{}${2}^{200} - {100} \cdot  {2}^{100}$ B. ${2}^{200} - {101} \cdot  {2}^{100}$ C. ${2}^{201} - {100} \cdot  {2}^{100}$ D. ${2}^{201} - {101} \cdot  {2}^{100}$

6.(2024$\cdot$静安区二模$\cdot$16)如果一个非空集合 $G$ 上定义了一个运算 *，满足如下性质，则称 $G$ 关于运算 * 构成一个群.

(1)封闭性，即对于任意的 $a$ 、 $b \in  G$ ，有 $a * b \in  G$ ；

(2)结合律，即对于任意的 $a$ 、 $b$ 、 $c \in  G$ ，有 $\left( {a * b}\right)  * c = a * \left( {b * c}\right)$ ；

(3)对于任意的 $a$ 、 $b \in  G$ ，方程 $x * a = b$ 与 $a * y = b$ 在 $G$ 中都有解.

例如，整数集 $\mathbf{Z}$ 关于整数的加法 $\left( +\right)$ 构成群，因为任意两个整数的和还是整数，且满足加法结合律,对于任意的 $a\text{ 、 }b \in  \mathbf{Z}$ ,方程 $x + a = b$ 与 $a + y = b$ 都有整数解; 而实数集 $\mathbf{R}$ 关于实数的乘法 ( $\times$ ) 不构成群,因为方程 $0 \times  y = 1$ 没有实数解. 以下关于 “群” 的真命题个数有 \dotfill{(\kh{})}

 \ding{172}  自然数集 ${N}$ 关于自然数的加法 (+) 构成群;  \ding{173}  有理数集 ${Q}$ 关于有理数的乘法 ( $\times$ ) 构成群;

 \ding{174}  平面向量集关于向量的数量积 ($\cdot$) 构成群； \ding{175}  复数集 ${C}$ 关于复数的加法 (+) 构成群.

\xx{}{}{}{}0 个 B. 1 个 C. 2 个 D. 3 个

\subsection*{考点4 > 其他新定义}

\begin{center}
\includegraphics[width=0.2\textwidth]{bodyxb/src/283_1116_348_284_271_0.jpg}
\end{center}


第 1 题图

\subsubsection*{一、填空题}

1. (2024$\cdot$奉贤区二模$\cdot$11)点 $P$ 是棱长为 1 的正方体 ${ABCD} - {A}_{1}{B}_{1}{C}_{1}{D}_{1}$ 棱上一点，则满足 $\left| {PA}\right|  + \left| {P{C}_{1}}\right|  = 2$ 的点 $P$ 的个数为\blkda[2]{}.

\subsubsection*{二、选择题}

2.(2025$\cdot$徐汇区二模$\cdot$15)在桌面上有一个质地均匀的正四面体 D 一

${ABC}$. 从该正四面体与桌面贴合的面上的三条棱中等可能地选取一条棱,沿其翻转正四面体至正四面体的另一个面与桌面贴合, 如此翻转称为一次操作. 如图, 开始时, 正四面体与桌面贴合的面为 ${ABC}$ ,操作 $n\left( {n = 1,2,3,\cdots }\right)$ 次后,正四面体与桌面贴合的面是

\begin{center}
\includegraphics[width=0.2\textwidth]{bodyxb/src/283_1089_798_315_275_0.jpg}
\end{center}


第 2 题图

${ABC}$ 的概率记为 ${P}_{n}$. 现有下列两个结论:  \ding{172}  ${P}_{2} = \dfrac{1}{3}$ ;  \ding{173}  ${P}_{25} < {P}_{24}$. 则下列说法正确的是 \dotfill{(\kh{})}

\xx{}{}{}{} \ding{172} 正确， \ding{173} 错误 B.  \ding{172} 错误， \ding{173} 正确

C.  \ding{172}  \ding{173} 都正确 D.  \ding{172}  \ding{173} 都错误

\subsubsection*{三、解答题}

3. (2025$\cdot$浦东新区二模$\cdot$21)定义域为 $\mathbf{R}$ 的可导函数 $y = f\left( x\right)$ 满足，曲线 $y = f\left( x\right)$ 上存在三个不同的点 $A\left( {{x}_{1},{y}_{1}}\right)$ 、 $B\left( {{x}_{2},{y}_{2}}\right)$ 、 $C\left( {{x}_{3},{y}_{3}}\right)$ 和 $\left( {{x}_{1} < {x}_{2} < {x}_{3}}\right)$ ，使得直线 ${AC}$ 与曲线 $y = f\left( x\right)$ 在点 $B$ 处的切线平行 (或重合). 若 ${x}_{1},{x}_{2},{x}_{3}$ 成等差数列,则称 $f\left( x\right)$ 为“等差函数”; 若 ${x}_{1},{x}_{2},{x}_{3}$ 成等差数列且 ${x}_{1}\text{ 、 }{x}_{2}\text{ 、 }{x}_{3}$ 均为正整数,则称 $f\left( x\right)$ 为“整数等差函数”.

(1)设 $f\left( x\right)  = {x}^{2} + x,g\left( x\right)  = \sin x$ ，分别判断 $f\left( x\right)$ 和 $g\left( x\right)$ 是否为“整数等差函数”，直接写出结论;

(2)若 $f\left( x\right)  = \dfrac{1}{{x}^{2} + m}$ 为 “整数等差函数”，求实数 $m$ 的最小值；

(3)已知 $y = f\left( x\right)$ 的导函数 $y = {f}^{\prime }\left( x\right)$ 在 $\mathbf{R}$ 上为增函数，且存在一个正常数 $T$ ，使得对任意 $x \in  \mathbf{R},f\left( {x + T}\right)  = {f}^{\prime }\left( x\right)$ 成立，证明: $f\left( x\right)$ 为“等差函数”的充要条件是 $f\left( x\right)$ 为常值函数.